% !TEX program = xelatex
% 中文自动译本：Stephen B. Pope, Turbulent Flows (Cambridge University Press, 2000).
\documentclass[11pt,a4paper,openany,fontset=none]{ctexbook}
\usepackage[margin=2.55cm,headheight=15pt]{geometry}
\usepackage{microtype}
\usepackage{fancyhdr}
\usepackage{hyperref}
\usepackage{xurl}
\usepackage{titlesec}
\usepackage{amsmath,amssymb,amsfonts,bm,mathtools}
\usepackage[export]{adjustbox}
\usepackage{enumitem}
\usepackage{graphicx}
\usepackage{tikz}
\usepackage{longtable,booktabs,array}
\setmainfont[AutoFakeBold=2,AutoFakeSlant=.2]{Arial Unicode MS}
\setCJKmainfont{Songti SC}
\setCJKsansfont{Songti SC}
\setCJKmonofont{Songti SC}
\setlength{\parindent}{2em}
\setlength{\parskip}{0.65em}
\setlength{\emergencystretch}{3em}
\linespread{1.30}\selectfont
\xeCJKsetup{CJKglue = \hskip 0.075em plus 0.04em minus 0.03em}
\setlist[itemize]{leftmargin=2.8em,itemsep=0.3em,topsep=0.45em}
\titleformat{\chapter}[display]{\bfseries\Huge\centering}{第\thechapter 章}{0.75em}{}
\titleformat{\section}{\large\bfseries}{}{0pt}{}
\titleformat{\subsection}{\normalsize\bfseries}{}{0pt}{}
\titlespacing*{\section}{0pt}{1.45em}{0.7em}
\titlespacing*{\subsection}{0pt}{1.2em}{0.55em}
\pagestyle{fancy}
\fancyhf{}
\fancyhead[C]{\small\nouppercase{\leftmark}}
\fancyfoot[C]{\thepage}
\renewcommand{\headrulewidth}{0.3pt}
\hypersetup{colorlinks=true,linkcolor=black,urlcolor=blue,
  pdftitle={湍流},pdfauthor={Stephen B. Pope}}
\graphicspath{{./}}
\title{湍流\\[0.5em]\Large 中文自动译文}
\author{Stephen B. Pope 著\\\small 译文由自动化工具生成，供学习与检索使用}
\date{原著：Cambridge University Press，2000\\生成日期：\today}
% 统一译本版式：避免公式、图片和题注跨页或侵入页眉、页脚。
\usepackage{needspace}
\clubpenalty=10000
\widowpenalty=10000
\displaywidowpenalty=10000
\raggedbottom
\setlength{\abovedisplayskip}{0.8\baselineskip plus 0.2\baselineskip minus 0.1\baselineskip}
\setlength{\belowdisplayskip}{0.8\baselineskip plus 0.2\baselineskip minus 0.1\baselineskip}
\setlength{\abovedisplayshortskip}{0.55\baselineskip}
\setlength{\belowdisplayshortskip}{0.55\baselineskip}
\newenvironment{sourceblock}
  {\par\addvspace{0.65\baselineskip}\noindent\begin{minipage}{\linewidth}\centering}
  {\end{minipage}\par\addvspace{0.45\baselineskip}}
\begin{document}
\frontmatter
\maketitle
\chapter*{说明}
\markboth{说明}{说明}
本文件是 Stephen B. Pope 所著 \emph{Turbulent Flows} 的中文自动译本。
正文按章节与小节的正常阅读顺序重新排版，不与原书逐页对齐。

为保证数学内容及实验数据可靠，本译本将显示公式、插图、复杂表格、符号表和参考文献
作为原 PDF 的高清裁片嵌入；正文、标题和图表题注重新排版为中文。
专业术语按湍流理论、RANS、PDF 方法和 LES 的常用中文译法统一。
涉及推导、图表数据和模型常数时，请同时核对原著。原著版权归作者与 Cambridge University Press 所有。

\begingroup\small\tableofcontents\endgroup
\clearpage
\makeatletter\@mainmattertrue\makeatother
\pagenumbering{arabic}
\clearpage\chapter*{前言}\addcontentsline{toc}{chapter}{前言}
\markboth{前言}{前言}
概率论知识，因此文中提供了必要的材料（例如，第 3.2-3.5 节）。

如果节奏要求不高，第一部分和第二部分可以在两个学期内完成-材料充足。或者，如果不需要建模，第一部分本身就提供了对湍流的相当完整的介绍。

许多练习要求读者“证明……”，从而引入额外的结果和观察结果。因此，建议阅读所有练习，即使没有进行。本书被设计为独立的文本，但提供了足够的参考文献以提供研究文献的入口。

无论这种性质的书在准备时多么小心，第一次印刷时难免会出现错误。已知更正的列表在以下位置给出。请读者通过 pope@mae.cornell.edu 向作者报告任何进一步的更正。

我非常感谢许多人在本书的准备过程中提供的帮助。我感谢康奈尔大学的同事 David Caughey、Sidney Leibovich、John Lumley、Dietmar Rempfer 和 Zellman Warhaft 的支持和技术投入。我感谢 Peter Bradshaw、Paul Durbin、Rodney Fox、Kemo Hanjalic、Charles Meneveau、Robert Moser、Blair Perot、Ugo Piomelli、P. K. Yeung 和 Norman Zabusky 在阅读章节草稿的基础上提出的宝贵建议。同样，我也感谢以下康奈尔大学毕业生对本书草稿的反馈：Bertrand Delarue、Thomas Dreeben、Matthew Overholt、Paul Van Slooten、Jun Xu、Cem Albukrek、Dawn Chamberlain、Timothy Fisher、Laurent Mydlarski、Gad Reinhorn、Shankar Subramaniam 和 Walter Welton。还感谢前五位在制作这些数据时提供的帮助。大部分打字稿是由琼·迈耶曼 (June Meyermann) 编写的，我们非常感谢她的耐心、准确性和热情。经过莎拉·波普 (Sarah Pope) 的仔细检查，参考书目的准确性得到了很大提高。最重要的是，我要感谢我的妻子 Linda，感谢她在这个项目期间以及多年来的耐心、支持和鼓励。

\clearpage\chapter*{符号表}\addcontentsline{toc}{chapter}{符号表}
\markboth{符号表}{符号表}
下表给出正文中最常用符号的中文释义。符号本身、上下标和张量阶数保持原书约定；
其后的完整原书符号表继续保留，便于逐项核对较少使用的模型常数。

\begingroup\small
\renewcommand{\arraystretch}{1.18}
\begin{longtable}{>{\centering\arraybackslash}p{0.22\textwidth}p{0.68\textwidth}}
\toprule 符号 & 中文释义 \\ \midrule
\endfirsthead
\toprule 符号 & 中文释义 \\ \midrule
\endhead
$\boldsymbol{U}(\boldsymbol{x},t)$ & 瞬时速度矢量 \\
$\langle\boldsymbol{U}\rangle$ & 系综平均或统计平均速度 \\
$\boldsymbol{u}=\boldsymbol{U}-\langle\boldsymbol{U}\rangle$ & 速度脉动 \\
$k=\tfrac12\langle u_i u_i\rangle$ & 湍动能 \\
$\varepsilon$ & 湍动能耗散率 \\
$\rho,\ p,\ \nu,\ \mu$ & 密度、压力、运动黏度和动力黏度 \\
$R_{ij}=\langle u_i u_j\rangle$ & 雷诺应力张量（按单位质量计） \\
$S_{ij},\ \Omega_{ij}$ & 平均应变率张量和平均旋转率张量 \\
$E(\kappa)$ & 三维能谱函数；$\kappa$ 为波数 \\
$D_{LL}(r)$ & 二阶纵向速度结构函数 \\
$f_U(V),\ F_U(V)$ & 随机变量 $U$ 的概率密度函数和累积分布函数 \\
$\eta,\ u_\eta,\ \tau_\eta$ & 柯尔莫哥洛夫长度、速度和时间尺度 \\
$L,\ T_L$ & 积分长度尺度和拉格朗日积分时间尺度 \\
$Re,\ Re_\lambda$ & 雷诺数和泰勒微尺度雷诺数 \\
$C_\mu$ & $k$--$\varepsilon$ 模型中的湍黏度常数 \\
$C_S$ & Smagorinsky 亚格子尺度模型系数 \\
$\kappa$ & 冯·卡门常数（壁面律语境） \\
$D/Dt$ & 随体导数（物质导数） \\
\bottomrule
\end{longtable}
\endgroup
\clearpage
\begin{sourceblock}
\includegraphics[page=15,pagebox=mediabox,trim={70.83bp 198.97bp 70.83bp 198.98bp},clip,width=\textwidth,max height=0.82\textheight]{Pope.pdf}
\end{sourceblock}
\clearpage
\begin{sourceblock}
\includegraphics[page=16,pagebox=mediabox,trim={70.83bp 198.97bp 70.83bp 198.98bp},clip,width=\textwidth,max height=0.82\textheight]{Pope.pdf}
\end{sourceblock}
\clearpage
\begin{sourceblock}
\includegraphics[page=17,pagebox=mediabox,trim={70.83bp 198.97bp 70.83bp 198.98bp},clip,width=\textwidth,max height=0.82\textheight]{Pope.pdf}
\end{sourceblock}
\clearpage
\begin{sourceblock}
\includegraphics[page=18,pagebox=mediabox,trim={70.83bp 198.97bp 70.83bp 198.98bp},clip,width=\textwidth,max height=0.82\textheight]{Pope.pdf}
\end{sourceblock}
\clearpage
\begin{sourceblock}
\includegraphics[page=19,pagebox=mediabox,trim={70.83bp 198.97bp 70.83bp 198.98bp},clip,width=\textwidth,max height=0.82\textheight]{Pope.pdf}
\end{sourceblock}
\clearpage
\begin{sourceblock}
\includegraphics[page=20,pagebox=mediabox,trim={70.83bp 198.97bp 70.83bp 198.98bp},clip,width=\textwidth,max height=0.82\textheight]{Pope.pdf}
\end{sourceblock}
\clearpage
\begin{sourceblock}
\includegraphics[page=21,pagebox=mediabox,trim={70.83bp 198.97bp 70.83bp 198.98bp},clip,width=\textwidth,max height=0.82\textheight]{Pope.pdf}
\end{sourceblock}
\clearpage
\begin{sourceblock}
\includegraphics[page=22,pagebox=mediabox,trim={70.83bp 198.97bp 70.83bp 198.98bp},clip,width=\textwidth,max height=0.82\textheight]{Pope.pdf}
\end{sourceblock}
\clearpage
\begin{sourceblock}
\includegraphics[page=23,pagebox=mediabox,trim={70.83bp 198.97bp 70.83bp 198.98bp},clip,width=\textwidth,max height=0.82\textheight]{Pope.pdf}
\end{sourceblock}
\clearpage
\begin{sourceblock}
\includegraphics[page=24,pagebox=mediabox,trim={70.83bp 198.97bp 70.83bp 198.98bp},clip,width=\textwidth,max height=0.82\textheight]{Pope.pdf}
\end{sourceblock}
\clearpage
\begin{sourceblock}
\includegraphics[page=25,pagebox=mediabox,trim={70.83bp 198.97bp 70.83bp 198.98bp},clip,width=\textwidth,max height=0.82\textheight]{Pope.pdf}
\end{sourceblock}
\clearpage
\begin{sourceblock}
\includegraphics[page=26,pagebox=mediabox,trim={70.83bp 198.97bp 70.83bp 198.98bp},clip,width=\textwidth,max height=0.82\textheight]{Pope.pdf}
\end{sourceblock}
\clearpage
\begin{sourceblock}
\includegraphics[page=27,pagebox=mediabox,trim={70.83bp 198.97bp 70.83bp 198.98bp},clip,width=\textwidth,max height=0.82\textheight]{Pope.pdf}
\end{sourceblock}
\clearpage
\begin{sourceblock}
\includegraphics[page=28,pagebox=mediabox,trim={70.83bp 198.97bp 70.83bp 198.98bp},clip,width=\textwidth,max height=0.82\textheight]{Pope.pdf}
\end{sourceblock}
\part{基础}
\chapter{引言}
\section*{1.1 湍流的性质}\phantomsection\addcontentsline{toc}{section}{1.1 湍流的性质}
在我们的日常环境中，有很多机会观察湍流，无论是烟囱中的烟雾、河流或瀑布中的水，还是强风的冲击。在观察瀑布时，我们立即看到水流是不稳定的、不规则的、看似随机且混乱的，当然每个漩涡或水滴的运动都是不可预测的。在固体火箭发动机形成的羽流中（参见图 1.1），可以观察到多种尺度的湍流运动，从大小与羽流宽度相当的涡流和凸起，到相机可以解析的最小尺度。这两个示例中提到的特征对于所有湍流都是常见的。

可以在实验室实验中进行更详细和仔细的观察。图 1.2 显示了两个不同雷诺数下湍流射流的平面图像。同样，浓度场是不规则的，并且可以观察到大范围的长度尺度。

正如上述讨论所暗示的，湍流的一个基本特征是流体速度场在位置和时间上都显着且不规则地变化。速度场（在 2.1 节中正确介绍）用 U(x, t) 表示，其中 x 是位置，t 是时间。

图 1.3 显示了在湍流射流中心线上测得的速度轴向分量的时程 Ul(t)（与图 1.2 类似）。水平线（图 1.3 中）显示由 (U1) 表示的平均速度，并在第 3.2 节中定义。可以观察到，速度 Ul (t) 显示出显着的波动（大约为 (U1) 的 25\%），并且时间历史远非周期性，而是在大范围的时间尺度上表现出变化。非常重要的是，我们观察到 U l( t) 及其平均值 (U1) 在某种意义上是“稳定的”：没有观察到 U1(t) 的巨大变化； Ul(t) 也不会长时间处于与 (Ul) 不同的值附近。

\begin{sourceblock}
\includegraphics[page=32,pagebox=mediabox,trim={64.54bp 785.80bp 64.54bp 109.48bp},clip,width=\textwidth,max height=0.72\textheight]{Pope.pdf}
\par\vspace{0.45\baselineskip}
\small 图 1.1。泰坦四号火箭发动机地面测试中的湍流羽流照片。喷嘴的出口直径为 3 m，估计羽流高度为 1,500 m，估计雷诺数为 200 x 10°。欲了解更多详细信息，请参阅 Mungal 和 Hollingsworth (1989)，经 San Jose Mercury \& News 许可。
\end{sourceblock}
图 1.4 显示了在相似湍流射流中测得的平均速度〈U1〉的剖面
该剖面随横流坐标变化。与瞬时速度 U1 的剧烈变化形成鲜明对比，平均速度〈U1〉的剖面十分平滑，不含精细结构；事实上，其形状与层流射流相差很小。

在工程应用中，湍流很常见，但不太容易看到。在用泵、压缩机、管道等处理液体或气体时，流动通常是湍流。类似地，飞机、汽车、船舶和潜艇等交通工具周围的流动也是湍流。

发动机、锅炉和熔炉中燃料与空气的混合，以及化学反应器中反应物的混合，都发生在湍流中。

\begin{sourceblock}
\includegraphics[page=33,pagebox=mediabox,trim={64.39bp 1151.86bp 64.39bp 109.44bp},clip,width=\textwidth,max height=0.72\textheight]{Pope.pdf}
\par\vspace{0.45\baselineskip}
\small 图 1.2。湍流射流中浓度的平面图像：（a）$Re=5{,}000$；（b）$Re=20{,}000$。引自 Dahm 和 Dimotakis（1990）。
\end{sourceblock}
\begin{sourceblock}
\includegraphics[page=33,pagebox=mediabox,trim={64.39bp 533.61bp 64.39bp 904.62bp},clip,width=\textwidth,max height=0.55\textheight]{Pope.pdf}
\par\vspace{0.45\baselineskip}
\small 图 1.3。湍流射流中心线上轴向速度分量 $U_1(t)$ 的时间历程。数据来自 Tong 和 Warhaft（1995）的实验。
\end{sourceblock}

湍流的一个重要特征是它比类似的层流更有效地输送和混合流体。 Osborne Reynolds (1883) 首次报道的实验很好地证明了这一点。染料被稳定地注入有水流动的长管道的中心线上。正如雷诺兹 (Reynolds) (1 894) 后来建立的那样，该流动的特征是单个无量纲参数，现在称为雷诺数 Re。

\begin{sourceblock}
\includegraphics[page=34,pagebox=mediabox,trim={64.54bp 1279.68bp 64.54bp 109.48bp},clip,width=\textwidth,max height=0.72\textheight]{Pope.pdf}
\par\vspace{0.45\baselineskip}
\small 图 1.4。湍流射流中的平均轴向速度分布。平均速度 (U;) 通过其在中心线上的值 (U;)y; 进行归一化；横流（径向）坐标x2通过距喷嘴x1的距离标准化。雷诺数为 95,500。改编自 Hussein、Capp 和 George (1994)。
\end{sourceblock}
流动的速度和长度尺度，v 是流体的运动粘度。 （对于管道流，2.4 和 L 分别为面积平均轴向速度和管道直径。）在雷诺管道流实验中，如果 Re 小于约 2,300，则流动为层流，流体速度不随时间变化，并且所有流线均平行于管道的轴线。在这种（层流）情况下，注入到中心线上的染料形成一条长条纹，其直径仅随着下游距离的增加而略有增加。另一方面，如果 Re 超过约 4,000，则流动为湍流。在靠近注射器的地方，染料条纹会因湍流运动而摇晃；随着下游距离的增加，它逐渐变得不那么明显；最终与周围的水混合，将峰值染料浓度降低到不再可见的程度。

（范戴克 (Van Dyke，1982) 中包含了雷诺兹实验复制品和其他规范湍流的可视化效果。研究这组照片具有很大的教育价值。）

在许多应用中，湍流对于输送和混合流体的有效性至关重要。当不同的流体流聚集在一起进行混合时，通常希望这种混合尽可能快地发生。对于污染物流来说确实是这样
随着雷诺数的增加，从层流到湍流的转变发生在
Re的范围，这个范围取决于实验的细节。

\section*{1.2 湍流研究}\phantomsection\addcontentsline{toc}{section}{1.2 湍流研究}
释放到大气中或水体中，以及用于在燃烧装置和化学反应器中混合不同反应物。

湍流对于“混合”流体的动量也很有效。因此，在飞机机翼和船体上，壁面剪应力（以及阻力）比层流时大得多。同样，与层流相比，固-液和液-气界面的传热传质速率在湍流中大大提高。

研究湍流的主要动机是前面三个观察结果的结合：绝大多数流动是湍流；物质、动量和热量在流动中的传输和混合具有重要的实际意义；湍流极大地提高了这些过程的速率。

1.2 湍流的研究 许多不同的技术已被用来解决有关湍流和湍流的许多不同问题。对这些研究进行分类的第一步是区分湍流中的小规模湍流和大规模运动。

正如第 6 章详细讨论的那样，在高雷诺数时会出现尺度分离。大尺度运动受到流动几何形状（例如边界条件）的强烈影响，并且它们控制传输和混合。另一方面，小尺度运动的行为几乎完全取决于它们从大尺度接收能量的速率以及粘度。因此，这些小尺度运动具有通用特征，与流动几何形状无关。人们很自然会问小尺度运动的特征是什么。它们可以从控制流体运动的方程中预测出来吗？这些是湍流理论的问题，在 Batchelor (1953)、Monin 和 Yaglom (1975)、Panchev (197 I)、Lesieur (19 90)、McComb (1990) 和其他人的书中都有讨论，而在本书中只稍微涉及到这些问题（在第 1 章）。

本书的重点是湍流，其研究可分为三类。

(i) 发现：旨在提供的实验（或模拟）研究
有关特定流量的定性或定量信息。

(ii) 建模：理论（或建模）研究，旨在开发
易于处理的数学模型，可以准确预测属性
(iii) 对照：研究（通常涉及实验和理论）

组件）旨在操纵或控制流量或
以有益的方式产生湍流-例如改变边界
增强混合的几何形状；或使用主动控制来减少阻力。

本书第一部分的其余部分主要基于第一类的研究，目的是让读者了解简单湍流的重要特征、主要物理过程以及它们与流体运动方程的关系。第一部分中对湍流的描述并不全面：其他材料可以在 Monin 和 Yaglom (1971)、Townsend (1976)、Hinze (1975) 和 Schlichting (1979) 的书中找到。

对于第二类旨在开发易于处理的数学模型的研究来说，“易于处理”这个词至关重要。对于流体流动，无论是层流还是湍流，其控制定律都体现在纳维-斯托克斯方程中，该方程已为人所知一个多世纪了。 （这些方程将在第二章中进行回顾。）考虑到流体流动的多样性和复杂性，相对简单的纳维-斯托克斯方程能够准确且完整地描述它们，这是非常值得注意的。然而，在湍流的背景下，它们的力量也是它们的弱点：方程描述了湍流速度场的每个细节，从最大到最小的长度和时间尺度。速度场中包含的信息量非常大，因此（通常）直接求解纳维-斯托克斯方程是不可能的。因此，虽然纳维-斯托克斯方程准确地描述了湍流，但它们并没有为湍流提供易于处理的模型。

求解湍流纳维-斯托克斯方程的直接方法称为直接数值模拟（DNS），在第 9 章中进行了描述。虽然 DNS 对于实际感兴趣的高雷诺数流动来说很棘手，但它仍然是研究中等雷诺数下简单湍流的强大研究工具。在第 7 章对墙界流的描述中，DNS 结果被广泛用于研究所涉及的物理过程。

对于应用中普遍存在的高雷诺数流动，自然的选择是采用统计方法。也就是说，描述湍流时，不是用速度而是用一些统计量，最简单的是平均速度场（U(x, t)）。基于此类统计数据的模型可以得出一组易于处理的方程，因为统计场在位置和时间上平滑变化（如果有的话）。在第 3 章中我们介绍
流场；在第一部分中，我们描述了可用于计算湍流特性的统计模型。所描述的方法包括： 湍流黏度模型，例如 k-E 模型（第 10 章）；雷诺应力模型（第 11 章）；基于速度概率密度函数 (PDF) 的模型（第 12 章）；和大涡模拟 (LES)（第 13 章）。

第 I1 部分中描述的统计模型可用于上述第三类的一些研究，即旨在操纵或控制流动或湍流的研究。然而，这里没有明确讨论此类研究。

广泛的数学技术用于描述和模拟湍流。提供了其中一些技术的附录作为简短的教程和摘要。第一个是笛卡尔张量，它被广泛使用。读者在继续阅读之前可能希望先阅读本材料（附录 A）。整本书都有练习，为读者提供了练习所用数学技巧的机会。大多数这些练习还包含额外的结果和观察结果。第 xxi 页上提供了术语和缩写的列表。

运动方程

在本章中，我们简要回顾一下控制恒性质牛顿流体流动的纳维-斯托克斯方程。更全面的描述可以在 Batchelor (1967)、Panton (1984) 和 Tritton (1988) 的文本中找到。在湍流研究中，两个重要的主题没有在这些文本中广泛讨论，即压力的泊松方程（Secti0.n 2.5）和纳维-

\chapter{流体运动方程}
采用向量表示法或笛卡尔张量表示法，附录 A 对此进行了回顾。

\section*{2.1 连续介质流体性质}\phantomsection\addcontentsline{toc}{section}{2.1 连续介质流体性质}
将流体视为连续介质的想法既自然又熟悉。然而，值得回顾一下连续统假说，该假说将流体的离散分子性质与连续统观点相协调，以避免在引入“流体颗粒”和“无穷小物质元素”等数量时产生混淆。

与人类的尺度相比，分子运动的长度和时间尺度非常小。以大气条件下的空气为例，分子间的平均间距为3 x m，平均自由程为6 x 10-8 m，分子连续碰撞的平均时间为10 s。相比之下，最小的几何长度尺度
流中的 -lo，l 很少小于 0.1 mm = 10- 4 m，对于高达 100 m s 的流速，会产生大于 的流时间尺度。 因此，即使对于这个具有较小长度和时间尺度的流的示例，这些流尺度也超过分子尺度三个或更多数量级。

长度尺度的分离通过努森数来量化

\begin{sourceblock}
\includegraphics[page=39,pagebox=mediabox,trim={64.39bp 1720.78bp 64.39bp 227.97bp},clip,width=0.94\textwidth,max height=0.38\textheight]{Pope.pdf}
\end{sourceblock}
在上面的例子中，K n 小于，而一般来说，连续统方法适用于 K n 1。

对于非常小的Kn，由于尺度的分离，存在中间长度尺度t'，使得t'与分子尺度相比较大，但与流动尺度相比较小( e , 1 << l' << l)。粗略地说，连续介质流体性质可以被认为是在尺寸为 V l 的体积上平均的分子性质。 让 V 表示以点 x 为中心的体积 V 的球形区域。然后，在时间 t 时，流体的密度 p(x, t) 为 V 中的分子质量除以 V。

类似地，流体的速度 U(x, t) 是 V 内分子的平均速度。由于尺度分离，连续介质特性对 l' 选择的依赖性可以忽略不计。

（虽然上一段中提出的方法是标准的（参见，例如，Batchelor (19 67) 和 Panton (1984)），如练习 2.1 所示，需要更加小心地根据尺度 l' 的平均来提供连续介质性质的正确定义。事实上，连续介质场最好定义为分子性质的期望，请参阅第 12 章。）

重要的是要认识到，一旦我们调用连续统假设来获得连续场，例如 p(x, t) 和 U(x, t)，我们就可以抛弃流体离散分子性质的所有概念，并且分子尺度不再相关。我们可以有意义地谈论“x，t 处的密度”，即使（从微观角度来看）（x，t）处很可能不存在任何物质。类似地，我们可以考虑比分子尺度更小的距离上的属性差异：实际上，当我们定义梯度时，我们就是这样做的，

\begin{sourceblock}
\includegraphics[page=39,pagebox=mediabox,trim={64.39bp 544.93bp 64.39bp 1351.83bp},clip,width=0.94\textwidth,max height=0.38\textheight]{Pope.pdf}
\end{sourceblock}
在理想气体的流动中，设 m i 、 x ( t ) 和 u(')(t) 为质量，
第 i 个分子的位置和速度。作为一般化的
标准连续统假设，考虑定义

\begin{sourceblock}
\includegraphics[page=39,pagebox=mediabox,trim={64.39bp 220.97bp 64.39bp 1665.79bp},clip,width=0.94\textwidth,max height=0.38\textheight]{Pope.pdf}
\end{sourceblock}
其中 uci) x ( \textasciitilde{}-) 和 K ( r ) 是平滑核，例如

\begin{sourceblock}
\includegraphics[page=39,pagebox=mediabox,trim={64.39bp 119.39bp 64.39bp 1849.77bp},clip,width=0.94\textwidth,max height=0.38\textheight]{Pope.pdf}
\end{sourceblock}
\begin{sourceblock}
\includegraphics[page=40,pagebox=mediabox,trim={64.39bp 1740.78bp 64.39bp 117.98bp},clip,width=0.94\textwidth,max height=0.38\textheight]{Pope.pdf}
\end{sourceblock}
其中 K'(r) dK(r)/dr。

（显然，方程（2.3）定义的连续域继承了
核的数学连续性特性。如果按照标准
处理时，K ( r ) 是分段常数，即

\begin{sourceblock}
\includegraphics[page=40,pagebox=mediabox,trim={64.39bp 1418.82bp 64.39bp 458.94bp},clip,width=0.94\textwidth,max height=0.38\textheight]{Pope.pdf}
\end{sourceblock}
那么 U(x, t) 是分段常数，因此不是连续差分
可以的。）

\section*{2.2 欧拉场与拉格朗日场}\phantomsection\addcontentsline{toc}{section}{2.2 欧拉场与拉格朗日场}
连续密度场和速度场 p(x, t) 和 U(x, t) 是欧拉场，因为它们由惯性系中的位置 x 索引。另一种拉格朗日描述的起点是流体粒子的定义，它是一个连续体概念。根据定义，流体质点是一个以局部流体速度运动的点：X+(t,Y)表示在指定的固定参考时间t 处位于Y 处的流体质点在时间t 的位置，见图2.1。从数学上讲，流体质点的位置

\begin{sourceblock}
\includegraphics[page=40,pagebox=mediabox,trim={64.39bp 270.97bp 64.39bp 1160.85bp},clip,width=\textwidth,max height=0.72\textheight]{Pope.pdf}
\par\vspace{0.45\baselineskip}
\small 图 2.1。 x-1 空间中流体粒子的轨迹 X(t, Y) 的草图，显示其在参考时间 to 和稍后时间的位置 Y
\end{sourceblock}
\textasciitilde{} + ( Yt ) 由两个方程定义。首先，参考时间to处的位置定义为

\begin{sourceblock}
\includegraphics[page=41,pagebox=mediabox,trim={64.39bp 1687.79bp 64.39bp 254.97bp},clip,width=0.94\textwidth,max height=0.38\textheight]{Pope.pdf}
\end{sourceblock}
二、方程

\begin{sourceblock}
\includegraphics[page=41,pagebox=mediabox,trim={64.39bp 1521.81bp 64.39bp 386.95bp},clip,width=0.94\textwidth,max height=0.38\textheight]{Pope.pdf}
\end{sourceblock}
表达了流体粒子随局部流体速度运动的事实。给定欧拉速度场 ( x ,t)，则；对于任意 Y，方程(2.8) 可以在时间上向前和向后积分以获得所有 t 的 X+(t,Y )。

例如，密度和速度的拉格朗日场是根据欧拉对应项定义的：

\begin{sourceblock}
\includegraphics[page=41,pagebox=mediabox,trim={64.39bp 1219.85bp 64.39bp 724.91bp},clip,width=0.94\textwidth,max height=0.38\textheight]{Pope.pdf}
\end{sourceblock}
\begin{sourceblock}
\includegraphics[page=41,pagebox=mediabox,trim={64.39bp 1135.86bp 64.39bp 807.90bp},clip,width=0.94\textwidth,max height=0.38\textheight]{Pope.pdf}
\end{sourceblock}
请注意，拉格朗日场 p+ 和 U+ 不是由流体粒子的当前位置索引，而是由其在参考时间 t0 处的位置 Y 索引。因此，Y 称为拉格朗日坐标或材料坐标。

对于固定 Y，Xf (t, Y ) 定义了一条轨迹（在 x-t 空间中），即 J 流体粒子路径，类似地，p+(t, Y ) 是流体粒子密度。偏导数 ap+(t, Y ) / a t 是固定 Y 处的密度变化率，即跟随流体粒子。弗罗恩方程。 (2.9) 我们得到

\begin{sourceblock}
\includegraphics[page=41,pagebox=mediabox,trim={64.39bp 409.95bp 64.39bp 1386.83bp},clip,width=0.94\textwidth,max height=0.38\textheight]{Pope.pdf}
\end{sourceblock}
其中实质性导数或随体导数由下式定义：

\begin{sourceblock}
\includegraphics[page=41,pagebox=mediabox,trim={64.39bp 244.97bp 64.39bp 1662.79bp},clip,width=0.94\textwidth,max height=0.38\textheight]{Pope.pdf}
\end{sourceblock}
因此，跟随流体粒子的密度变化率由下式给出
apf /a;) 拉格朗日场 ( e ) 的偏导数
类似地，对于固定的 Y，U+(t, Y ) 是流体粒子速度，并且
是流体质点速度的变化率，即流体质点加速度。

流体粒子也称为质点，我们已经看到它是由其在时间 t 时的位置 Y 及其随局部流体速度的运动来定义的（方程（2.8））。材料线、表面和体积的定义类似。例如，有时考虑一个简单的闭合曲面 So，它包含体积 Vo。相应的材料表面 S(t)i 被定义为在时间 t0 处与 So 重合，并且由 S(t) 上的每个点都以局部流体速度移动的性质定义。因此 S(t)i 由流体粒子 Y) 组成，它构成表面 So：

\begin{sourceblock}
\includegraphics[page=42,pagebox=mediabox,trim={64.39bp 1175.85bp 64.39bp 768.90bp},clip,width=0.94\textwidth,max height=0.38\textheight]{Pope.pdf}
\end{sourceblock}
由于材料表面随流体运动，因此表面与流体之间的相对速度为零。因此，流体粒子不能穿过材料表面；材料表面也不存在质量通量。

考虑两个流体粒子，它们在参考时间位于
在 Y 和 Y d Y 处，其中 d Y 是无穷小的位移。时
t，两个粒子之间的直线构成无穷小线元

\begin{sourceblock}
\includegraphics[page=42,pagebox=mediabox,trim={64.39bp 746.91bp 64.39bp 1195.85bp},clip,width=0.94\textwidth,max height=0.38\textheight]{Pope.pdf}
\end{sourceblock}
证明 s(t) 的演化过程为

\begin{sourceblock}
\includegraphics[page=42,pagebox=mediabox,trim={64.39bp 590.93bp 64.39bp 1315.83bp},clip,width=0.94\textwidth,max height=0.38\textheight]{Pope.pdf}
\end{sourceblock}
（提示：在泰勒级数中展开 U+(t, Y+dY) = U(X+( t, Y)+s(t), t)。

由于 s 无穷小，因此仅需要保留前导项。）

\section*{2.3 连续性方程}\phantomsection\addcontentsline{toc}{section}{2.3 连续性方程}
质量守恒或连续性方程是

\begin{sourceblock}
\includegraphics[page=42,pagebox=mediabox,trim={64.39bp 199.97bp 64.39bp 1707.78bp},clip,width=0.94\textwidth,max height=0.38\textheight]{Pope.pdf}
\end{sourceblock}
该方程在控制体积方面的推导和解释

在这里。进一步有用的解释是根据流体的比体积 9(x, t) = l/p(x, t)。方程的操纵(2.17) 产量

\begin{sourceblock}
\includegraphics[page=43,pagebox=mediabox,trim={64.39bp 1653.79bp 64.39bp 254.97bp},clip,width=0.94\textwidth,max height=0.38\textheight]{Pope.pdf}
\end{sourceblock}
左边是比容的对数增长率，而（如练习 2.3 和 2.4 所示）膨胀 V U 给出了无穷小的材料体积的对数增长率。因此，连续性方程可以被视为流体粒子的比容变化与无穷小材料体积元的体积变化之间的一致性条件。

在本书中，我们考虑常密度流（例如，p 独立于 x 和 t 的流）。在这种情况下，进化方程Eq。 (2.17) 退化为速度场为螺线管或无散度的运动条件：

\begin{sourceblock}
\includegraphics[page=43,pagebox=mediabox,trim={64.39bp 1154.85bp 64.39bp 792.90bp},clip,width=0.94\textwidth,max height=0.38\textheight]{Pope.pdf}
\end{sourceblock}
令 V(t) 为由材料表面 S(t) 界定的材料体积。

从几何表明 V(t) 内的流体体积 V(t) 的演变

\begin{sourceblock}
\includegraphics[page=43,pagebox=mediabox,trim={64.39bp 874.89bp 64.39bp 1072.86bp},clip,width=0.94\textwidth,max height=0.38\textheight]{Pope.pdf}
\end{sourceblock}
其中 dA 是 S(t) 上的面积元素，a 和
是向外指向

正常。利用散度定理得到

\begin{sourceblock}
\includegraphics[page=43,pagebox=mediabox,trim={64.39bp 654.92bp 64.39bp 1251.84bp},clip,width=0.94\textwidth,max height=0.38\textheight]{Pope.pdf}
\end{sourceblock}
证明，对于无穷小体积 dV(t)

材料体积dV，

\begin{sourceblock}
\includegraphics[page=43,pagebox=mediabox,trim={64.39bp 468.94bp 64.39bp 1476.81bp},clip,width=0.94\textwidth,max height=0.38\textheight]{Pope.pdf}
\end{sourceblock}
雅可比矩阵的行列式

\begin{sourceblock}
\includegraphics[page=43,pagebox=mediabox,trim={64.39bp 293.96bp 64.39bp 1593.80bp},clip,width=0.94\textwidth,max height=0.38\textheight]{Pope.pdf}
\end{sourceblock}
给出无穷小材料体积之间的体积比

t 时刻的 dV(t)，以及 t 时刻的体积 dV(to)。第一次订购
无穷小的 dt，表明

\begin{sourceblock}
\includegraphics[page=43,pagebox=mediabox,trim={64.39bp 119.39bp 64.39bp 1847.77bp},clip,width=0.94\textwidth,max height=0.38\textheight]{Pope.pdf}
\end{sourceblock}
\begin{sourceblock}
\includegraphics[page=44,pagebox=mediabox,trim={64.39bp 1753.78bp 64.39bp 146.98bp},clip,width=0.94\textwidth,max height=0.38\textheight]{Pope.pdf}
\end{sourceblock}
\begin{sourceblock}
\includegraphics[page=44,pagebox=mediabox,trim={64.39bp 1684.79bp 64.39bp 258.97bp},clip,width=0.94\textwidth,max height=0.38\textheight]{Pope.pdf}
\end{sourceblock}
因此表明

\begin{sourceblock}
\includegraphics[page=44,pagebox=mediabox,trim={64.39bp 1524.81bp 64.39bp 369.95bp},clip,width=0.94\textwidth,max height=0.38\textheight]{Pope.pdf}
\end{sourceblock}
练习 2.3 中定义的体积 V(t) 可以写成

\begin{sourceblock}
\includegraphics[page=44,pagebox=mediabox,trim={64.39bp 1362.83bp 64.39bp 528.93bp},clip,width=0.94\textwidth,max height=0.38\textheight]{Pope.pdf}
\end{sourceblock}
将该方程中的第一个和最后一个表达式微分为

相对于时间，并将结果与方程进行比较。 (2.21) 得到

\begin{sourceblock}
\includegraphics[page=44,pagebox=mediabox,trim={64.39bp 1172.85bp 64.39bp 735.91bp},clip,width=0.94\textwidth,max height=0.38\textheight]{Pope.pdf}
\end{sourceblock}
因此认为，在恒定密度流中，J(t,Y ) 是统一的。

\section*{2.4 动量方程}\phantomsection\addcontentsline{toc}{section}{2.4 动量方程}
动量方程基于牛顿第二定律，将流体粒子加速度 DUIDt 与流体所经历的表面力和体力联系起来。一般来说，分子起源的表面力由对称的应力张量 zij(x,t ) 描述，即
感兴趣的身体力是重力。当 Y 为引力 zij r j i a 势（ e，与重力相关的每单位质量的势能）时，每单位质量的体积力为

\begin{sourceblock}
\includegraphics[page=44,pagebox=mediabox,trim={64.39bp 594.92bp 64.39bp 1350.83bp},clip,width=0.94\textwidth,max height=0.38\textheight]{Pope.pdf}
\end{sourceblock}
（对于恒定的重力场，势为 Y = gz，其中 g 是重力加速度，z 是垂直坐标。）这些力导致流体根据动量方程加速

\begin{sourceblock}
\includegraphics[page=44,pagebox=mediabox,trim={64.39bp 342.96bp 64.39bp 1554.80bp},clip,width=0.94\textwidth,max height=0.38\textheight]{Pope.pdf}
\end{sourceblock}
现在，我们专门研究恒性牛顿流体的动量方程，这是本书中考虑的基本流动类别。在这种情况下，应力张量为

\begin{sourceblock}
\includegraphics[page=44,pagebox=mediabox,trim={64.39bp 119.39bp 64.39bp 1798.77bp},clip,width=0.94\textwidth,max height=0.38\textheight]{Pope.pdf}
\end{sourceblock}
其中 P 是压力，p 是（常数）粘度系数。回想一下（对于所考虑的恒定密度流）速度场是螺线管的（ e d Ui / a x , - - 0），我们观察到方程： (2.32) 将应力表示为各向同性 (-Pdij) 和偏向贡献之和。

通过将应力张量的表达式（方程（2.32））代入一般动量方程方程（2.32）。 (2.31)（并利用 p 和 y 一致且 V U = 0 的事实），我们得到纳维-斯托克斯方程
dt 轴 轴 ax,- dx,
此外，通过以下方式定义修正压力 p：

\begin{sourceblock}
\includegraphics[page=45,pagebox=mediabox,trim={64.39bp 1286.84bp 64.39bp 660.92bp},clip,width=0.94\textwidth,max height=0.38\textheight]{Pope.pdf}
\end{sourceblock}
这个方程简化为

其中 v p/p 是运动粘度。总之：恒定性质牛顿流体的流动由纳维-斯托克斯方程控制。 (2.35) 以及源自质量守恒的螺线管条件 V U 0。

在具有单位法线的静止实体壁上，速度满足的边界条件是不渗透条件

\begin{sourceblock}
\includegraphics[page=45,pagebox=mediabox,trim={64.39bp 755.90bp 64.39bp 1192.85bp},clip,width=0.94\textwidth,max height=0.38\textheight]{Pope.pdf}
\end{sourceblock}
和无滑移条件

\begin{sourceblock}
\includegraphics[page=45,pagebox=mediabox,trim={64.39bp 606.92bp 64.39bp 1331.83bp},clip,width=0.94\textwidth,max height=0.38\textheight]{Pope.pdf}
\end{sourceblock}
（它们一起产生 U 0）。

有时考虑理想（无粘性）流体的假设情况很有用，该流体被定义为具有各向同性应力张量

\begin{sourceblock}
\includegraphics[page=45,pagebox=mediabox,trim={64.39bp 376.95bp 64.39bp 1565.80bp},clip,width=0.94\textwidth,max height=0.38\textheight]{Pope.pdf}
\end{sourceblock}
动量守恒由欧拉方程给出

\begin{sourceblock}
\includegraphics[page=45,pagebox=mediabox,trim={64.39bp 207.97bp 64.39bp 1701.79bp},clip,width=0.94\textwidth,max height=0.38\textheight]{Pope.pdf}
\end{sourceblock}
由方程式得出。 （2.3I）、（2.34）和（2.38）。因为欧拉方程
边界条件比纳维-斯托克斯方程的边界条件要小，例如静止的固体壁，只能应用不渗透条件，并且一般情况下速度的切向分量不为零。

虽然最好直接从各向同性的定义获得欧拉方程（以及从它们导出的其他方程）

Ti，（方程（2.38）），但是可以观察到，通过将 v 设置为零，可以从纳维-斯托克斯方程获得欧拉方程。然而，重要的是要认识到 v = 0 是一个奇异极限：

0) 是消失粘度极限下的纳维-斯托克斯方程（v 与欧拉方程的解不同。一方面，即使在这个极限下，方程也需要不同的边界条件。

\section*{2.5 压力的作用}\phantomsection\addcontentsline{toc}{section}{2.5 压力的作用}
压力在（恒定密度）纳维-S 托克斯方程中的作用需要进一步评论。首先，我们观察到各向同性应力和保守体力具有相同的效果，这由修正的压力梯度表示。因此，体积力对速度场和修正压力场没有影响。 （当然，这与变密度流相反，在变密度流中浮力可能很重要。）从此以后，我们将 p 简称为“压力”。

我们可能习惯于将压力视为热力学变量，通过状态方程与密度和温度相关。然而，对于恒定密度流，压力和密度之间没有联系，需要对压力有不同的理解。

为此，我们取纳维-斯托克斯方程的散度。 （2.35），不假设速度场是螺线管的，而是用 A 表示膨胀率（ i . A V U ）。结果是

\begin{sourceblock}
\includegraphics[page=46,pagebox=mediabox,trim={64.39bp 515.93bp 64.39bp 1388.82bp},clip,width=0.94\textwidth,max height=0.38\textheight]{Pope.pdf}
\end{sourceblock}
\begin{sourceblock}
\includegraphics[page=46,pagebox=mediabox,trim={64.39bp 354.95bp 64.39bp 1495.81bp},clip,width=0.94\textwidth,max height=0.38\textheight]{Pope.pdf}
\end{sourceblock}
考虑方程的解。 (2.40)，初始条件和边界条件为 A 0。当且仅当 R 处处为零时，解为 A O，这在 tu.m 中意味着（根据方程 (2.41)）p 满足泊松方程

\begin{sourceblock}
\includegraphics[page=46,pagebox=mediabox,trim={64.39bp 125.98bp 64.39bp 1815.77bp},clip,width=0.94\textwidth,max height=0.38\textheight]{Pope.pdf}
\end{sourceblock}
\textasciitilde{} hu, swe 的结论是

满足这个泊松方程是必要的

磁场和螺线管速度保持螺线管速度的充分条件。

在静止的平面固体表面，纳维-斯托克斯方程(2.35)

\begin{sourceblock}
\includegraphics[page=47,pagebox=mediabox,trim={64.39bp 1580.80bp 64.39bp 287.96bp},clip,width=0.94\textwidth,max height=0.38\textheight]{Pope.pdf}
\end{sourceblock}
其中 n 是壁法线方向的坐标，U 是法线于壁的速度分量。该方程为泊松方程提供了诺伊曼边界条件，方程： （2.42）。给定这种形式的诺依曼条件，泊松方程(2.42) 根据同一时刻的速度场确定压力场 p(x, (t 在常数内)。因此，Vp 由当前速度场唯一确定，与流动的历史无关。

泊松方程的解(2.42) 可以根据格林函数明确地写出。 C.考虑泊松方程

\begin{sourceblock}
\includegraphics[page=47,pagebox=mediabox,trim={64.39bp 1107.86bp 64.39bp 836.89bp},clip,width=0.94\textwidth,max height=0.38\textheight]{Pope.pdf}
\end{sourceblock}
在域 V 中。源 S(x) 可以写成

\begin{sourceblock}
\includegraphics[page=47,pagebox=mediabox,trim={64.39bp 936.88bp 64.39bp 958.88bp},clip,width=0.94\textwidth,max height=0.38\textheight]{Pope.pdf}
\end{sourceblock}
其中 y 是 V 中的点，6(x y) 是 y 处的三维狄拉克δ函数'。泊松方程的一个解

\begin{sourceblock}
\includegraphics[page=47,pagebox=mediabox,trim={64.39bp 769.90bp 64.39bp 1171.85bp},clip,width=0.94\textwidth,max height=0.38\textheight]{Pope.pdf}
\end{sourceblock}
\begin{sourceblock}
\includegraphics[page=47,pagebox=mediabox,trim={64.39bp 640.92bp 64.39bp 1241.84bp},clip,width=0.94\textwidth,max height=0.38\textheight]{Pope.pdf}
\end{sourceblock}
（正如符号所暗示的，解取决于 x 和 delta 函数 y 的位置。）当它乘以 S(y) 并在方程上积分时， (2.46) 变为 V2f S ( e. 方程 (2.44))，因此方程 (2.44) (2.47) 成为解：

压力方程泊松方程的解。因此，(2.42) 是

\begin{sourceblock}
\includegraphics[page=47,pagebox=mediabox,trim={64.39bp 158.98bp 64.39bp 1716.78bp},clip,width=0.94\textwidth,max height=0.38\textheight]{Pope.pdf}
\end{sourceblock}
其中 p(h) 是取决于边界条件的调和函数 (V2p(h) 0)。 （可以用 V 边界上的表面积分来表示 p(h)i，例如参见 Kellogg (1967)）

表明（远离原点）

\begin{sourceblock}
\includegraphics[page=48,pagebox=mediabox,trim={64.39bp 1462.82bp 64.39bp 441.94bp},clip,width=0.94\textwidth,max height=0.38\textheight]{Pope.pdf}
\end{sourceblock}
求解纳维-斯托克斯方程的简单数值方法
对于恒定属性流，以小时间步长推进解决方案
At，从初始条件U(x,O)开始。在第 n 步
数值解用 U (") (x) 表示，它是一个近似值
到 U(x, n At)。每个时间步骤由两个子步骤组成，第一个子步骤
产生由下式定义的中间结果 U (x)

\begin{sourceblock}
\includegraphics[page=48,pagebox=mediabox,trim={64.39bp 1012.87bp 64.39bp 874.89bp},clip,width=0.94\textwidth,max height=0.38\textheight]{Pope.pdf}
\end{sourceblock}
第二个子步骤是

\begin{sourceblock}
\includegraphics[page=48,pagebox=mediabox,trim={64.39bp 819.90bp 64.39bp 1088.86bp},clip,width=0.94\textwidth,max height=0.38\textheight]{Pope.pdf}
\end{sourceblock}
其中 \$ ( ") (x) 是标量场。

（a）评论方程之间的联系。 （2.51）和
Navier-S 方程。

(b) 假设无散度，由等式得到： (2.51)

表达式（以 u(")fo 表示）U 的散度。

(c) 从等式获得(2.52) 的散度表达式
(d) 因此表明满足要求 V U 0 如果：

并且仅当 \$(")(x) 满足泊松方程

\begin{sourceblock}
\includegraphics[page=48,pagebox=mediabox,trim={64.39bp 235.97bp 64.39bp 1654.79bp},clip,width=0.94\textwidth,max height=0.38\textheight]{Pope.pdf}
\end{sourceblock}
(e) \$(")(x)a 与压力之间有什么关系？

\section*{2.6 守恒无源标量}\phantomsection\addcontentsline{toc}{section}{2.6 守恒无源标量}
除了速度 U(x, t) 之外，我们还考虑用 \$(x, t) 表示的守恒无源标量。在恒定特性流中，守恒定律

\begin{sourceblock}
\includegraphics[page=49,pagebox=mediabox,trim={64.39bp 1563.80bp 64.39bp 300.96bp},clip,width=0.94\textwidth,max height=0.38\textheight]{Pope.pdf}
\end{sourceblock}
其中 T 是（恒定且均匀）扩散率。标量 \$ 是守恒的，因为等式中没有源项或汇项。 （2.54）。它是被动的，因为（根据假设）它的值对材料属性没有影响（ e m p 、 ? v 和
因此它对流量没有影响。

标量\$可以表示各种物理属性。它可以是温度的小幅超额，足够小，以至于对材料性能的影响可以忽略不计。在这种情况下，T 是热扩散率，v/T 比率是普朗特数 Pr。或者，\$ 可以是痕量物质的浓度，在这种情况下，T 是分子扩散率，v/T 是施密特数 Sc。

标量的一个重要属性是它的有界性。如果 \$ 的初始值和边界值在给定范围内

\begin{sourceblock}
\includegraphics[page=49,pagebox=mediabox,trim={64.39bp 958.88bp 64.39bp 987.88bp},clip,width=0.94\textwidth,max height=0.38\textheight]{Pope.pdf}
\end{sourceblock}
那么所有 (x, t) 的 \$(x, t) 也位于此范围内：大于、\$、或小于的值不会出现。

为了显示这个结果，我们检查标量场中的局部最大值。假设此时存在局部最大值，并且我们选择一个坐标系，使得 d2\$/(dxi axj) 位于该处的主轴上。最大值的数学属性意味着

\begin{sourceblock}
\includegraphics[page=49,pagebox=mediabox,trim={64.39bp 602.92bp 64.39bp 1342.83bp},clip,width=0.94\textwidth,max height=0.38\textheight]{Pope.pdf}
\end{sourceblock}
并且二阶导数 d2\$/dxi、d2\$/dx; 和 d2\$/dx2 为负或零。因此，对于他们的总和，拉普拉斯算子，我们有

\begin{sourceblock}
\includegraphics[page=49,pagebox=mediabox,trim={64.39bp 425.95bp 64.39bp 1522.81bp},clip,width=0.94\textwidth,max height=0.38\textheight]{Pope.pdf}
\end{sourceblock}
然后，根据 \$ 的守恒方程（Eq. (2.54)），我们得到
对于每个向量 V；表明，遵循从局部最大值开始的任何轨迹， 的值不会增加。因此，没有办法
和边界条件。显然，类似的论点也适用于较低的

\section*{2.7 涡量方程}\phantomsection\addcontentsline{toc}{section}{2.7 涡量方程}
湍流的一个基本特征是它们是旋转的：也就是说，它们具有非零涡度。涡度 w(x, t) 是速度的旋度

\begin{sourceblock}
\includegraphics[page=50,pagebox=mediabox,trim={64.39bp 1466.82bp 64.39bp 479.94bp},clip,width=0.94\textwidth,max height=0.38\textheight]{Pope.pdf}
\end{sourceblock}
它等于 (x, t) 处流体旋转率的两倍。

涡度演化方程是通过纳维-斯托克斯方程的旋度获得的（2.35）：

\begin{sourceblock}
\includegraphics[page=50,pagebox=mediabox,trim={64.39bp 1203.85bp 64.39bp 696.91bp},clip,width=0.94\textwidth,max height=0.38\textheight]{Pope.pdf}
\end{sourceblock}
对于恒定密度流，压力项 (-V x Vp/p) 消失。

材料 s(t) 的无穷小线元的演化方程（见方程（2.16））为

\begin{sourceblock}
\includegraphics[page=50,pagebox=mediabox,trim={64.39bp 992.87bp 64.39bp 933.88bp},clip,width=0.94\textwidth,max height=0.38\textheight]{Pope.pdf}
\end{sourceblock}
除了粘性项外，它与涡度方程相同。因此，在无粘流中，涡量矢量的行为与无穷小的物质线元素相同（亥姆霍兹定理）。如果速度梯度产生的应变率作用于拉伸与 w 对齐的材料线单元，则 w 的大小相应地增加。这就是涡旋拉伸现象，涡旋拉伸是湍流中的一个重要过程，w VU 被称为涡旋拉伸项。

对于二维流，涡旋拉伸项消失，涡量的一个非零分量演化为守恒标量。由于不存在涡旋拉伸，二维湍流（在特殊情况下可能发生）与三维湍流有本质上的不同。

使用后缀表示法来验证关系：

\begin{sourceblock}
\includegraphics[page=50,pagebox=mediabox,trim={64.39bp 189.98bp 64.39bp 1758.78bp},clip,width=0.94\textwidth,max height=0.38\textheight]{Pope.pdf}
\end{sourceblock}
\begin{sourceblock}
\includegraphics[page=50,pagebox=mediabox,trim={64.39bp 119.39bp 64.39bp 1840.77bp},clip,width=0.94\textwidth,max height=0.38\textheight]{Pope.pdf}
\end{sourceblock}
\section*{2.8 应变率与旋转率}\phantomsection\addcontentsline{toc}{section}{2.8 应变率与旋转率}
\begin{sourceblock}
\includegraphics[page=51,pagebox=mediabox,trim={64.39bp 1788.77bp 64.39bp 152.98bp},clip,width=0.94\textwidth,max height=0.38\textheight]{Pope.pdf}
\end{sourceblock}
\begin{sourceblock}
\includegraphics[page=51,pagebox=mediabox,trim={64.39bp 1706.78bp 64.39bp 224.97bp},clip,width=0.94\textwidth,max height=0.38\textheight]{Pope.pdf}
\end{sourceblock}
是等式中的表达式。 （2.64）和（2.65）张量？

证明纳维-斯托克斯方程（Eq. (2.35)）可以写成
斯托克斯形式

\begin{sourceblock}
\includegraphics[page=51,pagebox=mediabox,trim={64.39bp 1464.82bp 64.39bp 440.94bp},clip,width=0.94\textwidth,max height=0.38\textheight]{Pope.pdf}
\end{sourceblock}
因此得到伯努利定理：对于稳定的、无粘性的、常数-
密度流、伯努利积分、

\begin{sourceblock}
\includegraphics[page=51,pagebox=mediabox,trim={64.39bp 1296.84bp 64.39bp 650.92bp},clip,width=0.94\textwidth,max height=0.38\textheight]{Pope.pdf}
\end{sourceblock}
是常数

(a) 沿着路线，
(b) 沿着涡线（ e 线平行于 o ），a nd
(c) 无旋流中的任何地方 (o 0)。

表明涡度平方或熵 w 2 o o 的演变

\begin{sourceblock}
\includegraphics[page=51,pagebox=mediabox,trim={64.39bp 910.88bp 64.39bp 1026.87bp},clip,width=0.94\textwidth,max height=0.38\textheight]{Pope.pdf}
\end{sourceblock}
2.8 应变率和旋转率 速度梯度 d U i / d x j 是二阶张量的分量，其一般性质在附录 B 中描述。 将 d U i / d x j 分解为各向同性、对称偏量和反对称张量

\begin{sourceblock}
\includegraphics[page=51,pagebox=mediabox,trim={64.39bp 455.94bp 64.39bp 1396.82bp},clip,width=0.94\textwidth,max height=0.38\textheight]{Pope.pdf}
\end{sourceblock}
其中，恒定密度流的膨胀 A V U 为零，Sij 是对称的偏应变率张量

\begin{sourceblock}
\includegraphics[page=51,pagebox=mediabox,trim={64.39bp 252.97bp 64.39bp 1641.79bp},clip,width=0.94\textwidth,max height=0.38\textheight]{Pope.pdf}
\end{sourceblock}
Qij 是反对称转速张量

\begin{sourceblock}
\includegraphics[page=51,pagebox=mediabox,trim={64.39bp 119.39bp 64.39bp 1803.77bp},clip,width=0.94\textwidth,max height=0.38\textheight]{Pope.pdf}
\end{sourceblock}
\section*{2.9 变换性质}\phantomsection\addcontentsline{toc}{section}{2.9 变换性质}
可以看出，牛顿应力定律（方程（2.32））可以重新表达为

\begin{sourceblock}
\includegraphics[page=52,pagebox=mediabox,trim={64.39bp 1690.79bp 64.39bp 248.97bp},clip,width=0.94\textwidth,max height=0.38\textheight]{Pope.pdf}
\end{sourceblock}
表明粘性应力与应变率线性相关，与旋转率无关。

涡度和旋转率的关系为

\begin{sourceblock}
\includegraphics[page=52,pagebox=mediabox,trim={64.39bp 1383.83bp 64.39bp 555.93bp},clip,width=0.94\textwidth,max height=0.38\textheight]{Pope.pdf}
\end{sourceblock}
其中 是交替符号。因此 Qij 和 包含相同的
Eijk oi 信息，但是（如附录 A 中所述）Qv 是张量，而不是。

从方程。 (2.16)，推导出长度演化方程
无穷小的物质线元素。显示增长率
线的长度与应变率线性相关，并且是独立的
的旋转率，

表明涡度方程（方程（2.60））也可以是
重新表达压力泊松方程中的源（方程（2.42））

就 Sij 和 Qij 而言。

在简单剪切流中，除了
ul/ax2。对于这种情况，写下 Sij 的分量和（如
流动特性 通过研究纳维-斯托克斯方程在进行各种变换时的行为，我们能够推断出它们所描述的流体流动的重要特性。这些属性中最重要的是雷诺数相似性、固定旋转和坐标轴反射下的不变性、伽利略不变性以及缺乏

\begin{sourceblock}
\includegraphics[page=53,pagebox=mediabox,trim={64.39bp 461.94bp 64.39bp 376.75bp},clip,width=\textwidth,max height=0.72\textheight]{Pope.pdf}
\par\vspace{0.45\baselineskip}
\small 图 2.2。用于研究纳维-斯托克斯方程变换性质的实验草图：（a）参考实验（指E坐标系）； (b)\{h)其他实验（参考E坐标系）。
\end{sourceblock}
考虑在实验室进行的特定流体力学实验，并考虑第二个实验，该实验与第一个实验类似，但在某些方面有所不同。例如：第二个实验可以在不同的时间进行；该装置可以放置在不同的位置；它
可以使用不同的流体；或者可以建造第二个装置，其几何形状与第一个装置相似，但尺寸不同。对于这些差异中的每一个，我们可以询问两个实验中的速度场是否相似。也就是说，当速度场被适当缩放并参考适当的坐标系时，速度场是否相同？这些问题可以通过研究纳维-斯托克斯方程的变换性质（也称为不变性或对称性）来回答。

这些是湍流建模中的重要考虑因素。除非模型的变换属性与纳维-斯托克斯方程的变换属性一致，否则模型在质量上将是错误的。

图 2.2(a) 是第一个（参考）实验中考虑的装置的草图。装置的尺寸由长度尺度 L 表征，速度的初始条件和边界条件由速度尺度 U 表征。坐标系（用 E 表示，正交基向量 ei）的原点和轴相对于装置固定，装置静止在惯性系中。

长度尺度 L 和速度尺度 U 用于定义无量纲自变量

\begin{sourceblock}
\includegraphics[page=54,pagebox=mediabox,trim={64.39bp 989.88bp 64.39bp 948.88bp},clip,width=0.94\textwidth,max height=0.38\textheight]{Pope.pdf}
\end{sourceblock}
第三因变量

\begin{sourceblock}
\includegraphics[page=54,pagebox=mediabox,trim={64.39bp 852.89bp 64.39bp 1089.86bp},clip,width=0.94\textwidth,max height=0.38\textheight]{Pope.pdf}
\end{sourceblock}
将这些简单的缩放变换应用于连续性方程；q。 (2.19)，纳维-斯托克斯方程(2.35) 和泊松方程2.42），我们得到

\begin{sourceblock}
\includegraphics[page=54,pagebox=mediabox,trim={64.39bp 618.92bp 64.39bp 1329.83bp},clip,width=0.94\textwidth,max height=0.38\textheight]{Pope.pdf}
\end{sourceblock}
\begin{sourceblock}
\includegraphics[page=54,pagebox=mediabox,trim={64.39bp 465.94bp 64.39bp 1437.82bp},clip,width=0.94\textwidth,max height=0.38\textheight]{Pope.pdf}
\end{sourceblock}
\begin{sourceblock}
\includegraphics[page=54,pagebox=mediabox,trim={64.39bp 342.96bp 64.39bp 1547.80bp},clip,width=0.94\textwidth,max height=0.38\textheight]{Pope.pdf}
\end{sourceblock}
.这里雷诺数是

\begin{sourceblock}
\includegraphics[page=54,pagebox=mediabox,trim={64.39bp 211.97bp 64.39bp 1729.78bp},clip,width=0.94\textwidth,max height=0.38\textheight]{Pope.pdf}
\end{sourceblock}
显然，雷诺数是这些中出现的唯一参数
雷诺数相似性 图 2.2(b) 所示的实验具有不同的长度尺度 £,、速度尺度 U, 以及流体特性 vy 和 pp。如果尺度变量以类似的方式定义（\% = x/Lp, U = 等），则边界条件（在两个实验中用 表示）是相同的，并且变换后的纳维-斯托克斯方程与方程 1 相同。 (2.78)-(2.80)，除了 Re 被替换为
因此，如果雷诺数相同 (Re = Re,)，则缩放速度场 0 (\%,7) 也相同，因为它们由具有相同初始条件和边界条件的相同方程控制。这是雷诺数相似性的性质。

缩放后的欧拉方程与方程相同。 (2.79)，但省略了 Re 中的术语。因此，欧拉方程给出的缩放速度场是相同的，与 Ly、U4 和 py 无关：它们表现出尺度相似性，并且据说欧拉方程对于尺度变换是不变的。

时间和空间不变性

纳维-斯托克斯方程最简单的不变性是它们相对于时间和空间变化的不变性。如图2.2(c)所示，我们认为第二个实验比参考实验晚了T时间，设备平移了X量。第二个实验中的速度场参考图2.2(c)所示的E坐标系，其具有正交基向量@；。缩放自变量定义为

\begin{sourceblock}
\includegraphics[page=55,pagebox=mediabox,trim={64.39bp 566.93bp 64.39bp 1381.83bp},clip,width=0.94\textwidth,max height=0.38\textheight]{Pope.pdf}
\end{sourceblock}
\begin{sourceblock}
\includegraphics[page=55,pagebox=mediabox,trim={64.39bp 480.94bp 64.39bp 1461.82bp},clip,width=0.94\textwidth,max height=0.38\textheight]{Pope.pdf}
\end{sourceblock}
证明变换后的方程与方程相同是微不足道的。 (2.78)-(2.80)。

旋转和反射不变性

图2.2（d）显示了与参考实验不同方向的装置；通过旋转参考 (E) 坐标轴可以获得适当的 E 坐标系。图2.2(e)显示了
装置。此时，通过坐标轴的反射得到合适的E坐标系。

这些坐标变换（轴的旋转和反射）正是在笛卡尔张量中考虑的（参见附录 A）。与; = e;* 2;是方向余弦，缩放变量是

\begin{sourceblock}
\includegraphics[page=56,pagebox=mediabox,trim={64.39bp 1560.80bp 64.39bp 383.95bp},clip,width=0.94\textwidth,max height=0.38\textheight]{Pope.pdf}
\end{sourceblock}
纳维-斯托克斯方程可以用笛卡尔张量表示法书写，因此变换后的方程与参考系中的方程相同（方程（2.78）-（2.80））。因此，纳维-斯托克斯方程对于坐标轴的旋转和反射是不变的。

在这些考虑中，区分两种“旋转”很重要。这里我们考虑通过 E 坐标轴的固定旋转获得的 E 坐标系。 “固定”是指方向余弦 aj；不依赖于时间。相反，我们考虑下面的旋转框架，因此方向余弦与时间相关。

关于反射的不变性具有物理意义和数学结果，这在附录 A 中进行了更详细的讨论。物理意义是纳维-斯托克斯方程不包含右手或左手运动的偏差。当然，这种偏差可能会发生在流动中-在龙卷风中最为明显-但它是由初始条件或边界条件，或由坐标系旋转引起的，而不是由运动方程（以惯性坐标系表示）引起的。

任何用笛卡尔张量表示法编写的方程都可以确保坐标轴旋转和反射下的不变性。相反，以向量表示法编写并涉及伪向量（例如，涡量）的方程，或使用交替符号 ¢, 以后缀表示法编写的方程，不能确保这些不变性。

时间反转 类似于坐标轴的反射（例如， = 我们可以
通过定义时间的逆转

\begin{sourceblock}
\includegraphics[page=56,pagebox=mediabox,trim={64.39bp 283.96bp 64.39bp 1654.79bp},clip,width=0.94\textwidth,max height=0.38\textheight]{Pope.pdf}
\end{sourceblock}
\begin{sourceblock}
\includegraphics[page=56,pagebox=mediabox,trim={64.39bp 210.97bp 64.39bp 1736.78bp},clip,width=0.94\textwidth,max height=0.38\textheight]{Pope.pdf}
\end{sourceblock}
t 很容易证明相应的变换纳维-斯托克斯等于
项（与 \textasciitilde{}e-' 成比例）已更改。因此，纳维-S 托克斯方程在时间反转下不是不变的；但欧拉方程是。

伽利略不变性 本节的其余主题与移动框架有关。我们首先考虑如图2.2（f）所示，装置以*固定速度V运动，使得两个坐标系（E 和 ）都在惯性系中。坐标系之间的变换为

\begin{sourceblock}
\includegraphics[page=57,pagebox=mediabox,trim={64.39bp 1420.82bp 64.39bp 526.93bp},clip,width=0.94\textwidth,max height=0.38\textheight]{Pope.pdf}
\end{sourceblock}
\begin{sourceblock}
\includegraphics[page=57,pagebox=mediabox,trim={64.39bp 1340.83bp 64.39bp 600.92bp},clip,width=0.94\textwidth,max height=0.38\textheight]{Pope.pdf}
\end{sourceblock}
量表示同一不同惯性系伽利略不变量。从方程式。 (2.89) 和 (2.90) 我们得到

\begin{sourceblock}
\includegraphics[page=57,pagebox=mediabox,trim={64.39bp 1143.86bp 64.39bp 774.90bp},clip,width=0.94\textwidth,max height=0.38\textheight]{Pope.pdf}
\end{sourceblock}
\begin{sourceblock}
\includegraphics[page=57,pagebox=mediabox,trim={64.39bp 1011.87bp 64.39bp 918.88bp},clip,width=0.94\textwidth,max height=0.38\textheight]{Pope.pdf}
\end{sourceblock}
\begin{sourceblock}
\includegraphics[page=57,pagebox=mediabox,trim={64.39bp 896.89bp 64.39bp 1051.87bp},clip,width=0.94\textwidth,max height=0.38\textheight]{Pope.pdf}
\end{sourceblock}
表明速度梯度和流体加速度是伽利略不变的，而速度及其偏时间导数则不是。其他伽利略不变量包括标量，例如 4(x,t) 和压力 p(x,t)，以及与速度梯度相关的量，例如 Sjj7R jj 和涡度 w。

简单地表明，变换后的纳维-斯托克斯方程（用 f x/L 等形式表示为 U = U/U 等）与方程 1 相同。 (2.78)-(2.80)，因此是伽利略不变量。就像经典力学描述的所有现象一样，流体流动的行为在所有惯性框架中都是相同的。

下列哪些是伽利略不变量：

(a) 流线（根据定义，它是一条无处不在的曲线
平行于速度矢量），
(b) 涡线（根据定义，它是一条无处不在的曲线
(c) 螺旋度，定义为 U
(d) 熵，定义为 o
(e) 材料线、表面和体积，
(f) 对于标量场； a\$/dt、a\$/dxi 和 D\$/Dt？

扩展伽利略不变性 纳维-斯托克斯方程的一个特殊性质是它们在框架的直线加速度下保持不变。我们认为，如图 2.2(g) 所示，第二个实验是在以可变速度 V(l) 移动的平台上进行的，但框架没有旋转，因此坐标方向（例如 el 和 El ）保持平行。变换后的变量 、 7 和 U 由等式定义。 （2.89）和（2.90），变换后的纳维-斯托克斯方程为

\begin{sourceblock}
\includegraphics[page=58,pagebox=mediabox,trim={64.39bp 1119.86bp 64.39bp 774.90bp},clip,width=0.94\textwidth,max height=0.38\textheight]{Pope.pdf}
\end{sourceblock}
其中右侧的附加项是框架的加速度 A d V/dt。最后两项可以写成

\begin{sourceblock}
\includegraphics[page=58,pagebox=mediabox,trim={64.39bp 916.88bp 64.39bp 981.88bp},clip,width=0.94\textwidth,max height=0.38\textheight]{Pope.pdf}
\end{sourceblock}
表明框架加速度可以被修改后的压力吸收。因此，变换变量的纳维-斯托克斯方程

\begin{sourceblock}
\includegraphics[page=58,pagebox=mediabox,trim={64.39bp 753.90bp 64.39bp 1182.85bp},clip,width=0.94\textwidth,max height=0.38\textheight]{Pope.pdf}
\end{sourceblock}
与等式相同。 (2.76)-(2.80)。因此，具有任意直线加速度的实验中的缩放速度和修正压力 j!i 场与惯性火焰中的相同。这是扩展的伽利略 incariarzce（仅适用于恒定密度流）。

框架旋转 最后，我们考虑在非惯性旋转框架中进行第二个实验，图 2.2(h)。在坐标系中，与时间相关的基向量 gj(t) 演化为

\begin{sourceblock}
\includegraphics[page=58,pagebox=mediabox,trim={64.39bp 252.97bp 64.39bp 1655.79bp},clip,width=0.94\textwidth,max height=0.38\textheight]{Pope.pdf}
\end{sourceblock}
其中 Qij(t) 是框架的旋转率。请注意，在这个
= -Qjj(t) 情况下，方向余弦 aij(t) ei zj(t) 与时间相关。

与等式相同（2.94），但框架加速度 -Aj 被虚拟力取代

\begin{sourceblock}
\includegraphics[page=59,pagebox=mediabox,trim={64.39bp 1653.79bp 64.39bp 266.97bp},clip,width=0.94\textwidth,max height=0.38\textheight]{Pope.pdf}
\end{sourceblock}
（参见练习 2.15）。对 F 的三个贡献分别是离心力、科里奥利力和角加速度力。离心力可以被吸收到修改后的压力中，但其余两个力则不能。正如气象学和涡轮机械中众所周知的那样，科里奥利力会对旋转框架中的流动产生重大影响。

旋转坐标系和非旋转坐标系中相同的量被认为具有材料坐标系无关性。显然，纳维-斯托克斯方程不具备这个性质。

框架旋转的影响在涡度方程中也很明显。在非惯性坐标系中，涡度演化方程

\begin{sourceblock}
\includegraphics[page=59,pagebox=mediabox,trim={64.39bp 1065.87bp 64.39bp 881.89bp},clip,width=0.94\textwidth,max height=0.38\textheight]{Pope.pdf}
\end{sourceblock}
从 Navier-Stokes 方程获得 ( e . E q. (2.94) with F, in place

\begin{sourceblock}
\includegraphics[page=59,pagebox=mediabox,trim={64.39bp 801.90bp 64.39bp 1090.86bp},clip,width=0.94\textwidth,max height=0.38\textheight]{Pope.pdf}
\end{sourceblock}
显然，由于最后两项对应于科里奥利力和角加速度力，旋转坐标系中的涡量方程与惯性坐标系中的涡量方程不同（方程（2.60））。

设 X(t) 为移动点相对于原点的位置
惯性系中的 E 坐标系。令 Y(r) ai(t)x-(t)b e
同一点相对于非惯性系E的位置。

坐标系原点以速度 V(t) 移动，其基础
向量根据方程演化。 （2.97）。如果起源重合
在时间 t = 0 时，则

\begin{sourceblock}
\includegraphics[page=59,pagebox=mediabox,trim={64.39bp 165.98bp 64.39bp 1734.78bp},clip,width=0.94\textwidth,max height=0.38\textheight]{Pope.pdf}
\end{sourceblock}
\begin{sourceblock}
\includegraphics[page=60,pagebox=mediabox,trim={64.39bp 1738.78bp 64.39bp 157.98bp},clip,width=0.94\textwidth,max height=0.38\textheight]{Pope.pdf}
\end{sourceblock}
\begin{sourceblock}
\includegraphics[page=60,pagebox=mediabox,trim={64.39bp 1660.79bp 64.39bp 278.96bp},clip,width=0.94\textwidth,max height=0.38\textheight]{Pope.pdf}
\end{sourceblock}
其中过点表示相对于时间的差异。

二维流 纳维-斯托克斯方程的另一个特殊性质是，对于二维流（例如，在 XI-x2 平面中），它们相对于流平面中框架的稳定旋转（e.绕 x3 轴的旋转）是不变的，参见 Speziale (198 1)。对于二维流，有时用流函数和涡度重新表达纳维-斯托克斯方程是有用的。流函数 y,(xl,x 2, t) 的速度为

\begin{sourceblock}
\includegraphics[page=60,pagebox=mediabox,trim={64.39bp 1080.86bp 64.39bp 779.90bp},clip,width=0.94\textwidth,max height=0.38\textheight]{Pope.pdf}
\end{sourceblock}
涡度的唯一非零分量是

\begin{sourceblock}
\includegraphics[page=60,pagebox=mediabox,trim={64.39bp 921.88bp 64.39bp 969.88bp},clip,width=0.94\textwidth,max height=0.38\textheight]{Pope.pdf}
\end{sourceblock}
对于框架的稳定旋转，方程中的最后一项。 (2.100) 为零，并且对于所考虑的二维流，对倒数第二项的显式求值表明它也为零（参见练习 2.17）。因此，对于这种特殊情况，涡量不受坐标系旋转的影响，因此纳维-斯托克斯方程表现出材料-坐标系无差异（在这种限制意义上）。

对于二维流，U3 0 和 UI 和 U2 由下式给出
等式。 （2.104），表明速度散度对于所有
流函数。表明流函数和涡度是重新
由泊松方程得出

参考等式中的倒数第二项。 （2.100），考虑

\begin{sourceblock}
\includegraphics[page=60,pagebox=mediabox,trim={64.39bp 119.39bp 64.39bp 1749.78bp},clip,width=0.94\textwidth,max height=0.38\textheight]{Pope.pdf}
\end{sourceblock}
对于二维流（在平面中），以及对于框架旋转
系统蒸发散在同一平面上。 d \& / d \textasciitilde{}a\textasciitilde{}nd 的组成部分是
零？显示 Q；和Q；为零。获取结果

\begin{sourceblock}
\includegraphics[page=61,pagebox=mediabox,trim={64.39bp 1596.80bp 64.39bp 300.96bp},clip,width=0.94\textwidth,max height=0.38\textheight]{Pope.pdf}
\end{sourceblock}
因此认为对于所考虑的流类别，R' 为零。

\chapter{湍流的统计描述}
湍流的统计描述

\section*{3.1 湍流的随机性质}\phantomsection\addcontentsline{toc}{section}{3.1 湍流的随机性质}
在湍流中，速度场是随机的。这是什么作用
U(x,t)语句是什么意思？为什么会这样呢？

作为第一步，我们需要理解“随机”这个词。考虑在一组指定条件 C 下可以重复多次的流体流动实验，并考虑事件 A，例如 A \{U < 10 m s- 1，其中 U 是指定位置和时间的速度的指定分量（从实验开始时测量）。如果事件A不可避免地发生，则A是确定的或确定的。如果事件A不能发生，那么它是不可能的。第三种可能性是A可能发生或者可能但不一定发生。在这种情况下，事件 A 是随机的。那么，在示例中 A \{U < 10 ms \} U 是随机变量。

有时会犯的一个错误是错误地赋予“随机”这一名称额外的意义，然后对湍流是一种随机现象这一事实提出质疑。事件 A 是随机的仅意味着它既不是确定的也不是不可能的。 U 是一个随机变量，仅意味着每次在同一组条件 C 下重复实验时，它不具有相同的唯一值。图 3.1 说明了随机变量 U 在 40 次重复实验中所取的值 U (n = 1, 2,..., 40)。

下一个要解决的问题是湍流的随机性质与纳维-斯托克斯方程中体现的经典力学的确定性性质之间的一致性。如果运动方程是确定性的，为什么解是随机的？答案在于两个观察结果的结合。

(i) 在任何湍流中，初始阶段都不可避免地存在扰动
条件、边界条件和材料属性。

\begin{sourceblock}
\includegraphics[page=63,pagebox=mediabox,trim={64.39bp 1299.84bp 64.39bp 109.44bp},clip,width=\textwidth,max height=0.72\textheight]{Pope.pdf}
\par\vspace{0.45\baselineskip}
\small 图 3.1。湍流实验第 n 次重复时随机速度变量 U 的值 U 的草图。
\end{sourceblock}
(ii) 湍流场对这种扰动表现出敏锐的敏感性。

在讨论随机性时，我们考虑了“在一组指定的条件 C 下可以重复多次的流体流动实验”。一个例子是纯水在 20°C 下流过光滑的直管。应当理解，如此定义的条件 C 是不完整的：在实践中，不可避免地存在来自这些标称条件的扰动。边界条件中可能存在扰动，例如，通过设备的振动，或来自标称光滑表面的详细光洁度。由于温度的微小不均匀性或杂质的存在，可能会导致流体特性的扰动，并且可能存在通过小心和努力，这些扰动可以减少，但无法消除，因此，标称条件 C 是不完整的，因此不能唯一地确定湍流的演变。

扰动的存在本身并不能解释湍流的随机性质，因为事实上，这种扰动也存在于层流中。然而，在高雷诺数湍流下，流场的演化对初始条件、边界条件和材料特性的微小变化极其敏感。这种敏感性在动力系统的研究中得到了很好的理解，并且在混沌书籍中得到了普及（例如，Gleick (1988) 和 Moon (1992)）-现在可以使用洛伦兹方程来证明。

Lorenz（1963）研究了一个由三个常微分方程描述的瞬态系统。

\begin{sourceblock}
\includegraphics[page=64,pagebox=mediabox,trim={64.39bp 1005bp 64.39bp 110bp},clip,width=\textwidth,max height=0.60\textheight]{Pope.pdf}
\par\vspace{0.45\baselineskip}
\small 图 3.2。洛伦兹方程组（3.1）的时间历程：（a）由初始条件（3.2）得到的 $x(t)$；
（b）由略有差异的初始条件（3.3）得到的 $\widehat{x}(t)$；（c）二者之差 $\widehat{x}(t)-x(t)$。
\end{sourceblock}

常微分方程

\begin{sourceblock}
\includegraphics[page=64,pagebox=mediabox,trim={64.39bp 636.92bp 64.39bp 1302.84bp},clip,width=0.94\textwidth,max height=0.38\textheight]{Pope.pdf}
\end{sourceblock}
其中系数为 a -- 1 0 , P = ,, a n d p = 28。对于初始条件

\begin{sourceblock}
\includegraphics[page=64,pagebox=mediabox,trim={64.39bp 507.94bp 64.39bp 1430.82bp},clip,width=0.94\textwidth,max height=0.38\textheight]{Pope.pdf}
\end{sourceblock}
图 3.2(a) 显示了从方程 1 的数值积分获得的时间历史 x(t)。 ( 3 ) 初始条件稍有不同时得到的结果用 ?(t) 表示

\begin{sourceblock}
\includegraphics[page=64,pagebox=mediabox,trim={64.39bp 290.96bp 64.39bp 1647.79bp},clip,width=0.94\textwidth,max height=0.38\textheight]{Pope.pdf}
\end{sourceblock}
如图 3+2(b)+I t 所示，可以观察到（如预期）.s(t) 和 d ?(t) 最初无法区分，但到 t 35 它们就完全不同了-这一观察在猪身上变得更加清晰。 3.2(c)，显示了差异

\section*{3.2 随机变量的表征}\phantomsection\addcontentsline{toc}{section}{3.2 随机变量的表征}
这种对初始条件极度敏感的结果是，超过某个点后，系统的状态就无法预测。在这个例子中，如果初始状态的已知程度仅在 10\% 之内，那么图 3.2 清楚地表明，在 t = 35 之外无法做出有用的预测。

这个例子旨在证明一组简单的确定性方程（比纳维-斯托克斯方程简单得多）可以表现出对初始条件的敏锐敏感性，从而表现出不可预测性。

洛伦兹系统的定性行为取决于系数。特别是，对于固定值 o = 10 且 p= g，行为取决于 p。如果 p 小于临界值 \textasciitilde{} 24.74，则系统达到稳定不动点，即状态变量 [x(t), y(t), z(t)] 渐近趋于固定值。然而，对于 p > p”（例如，图 3.2 中的 9 = 28），混沌行为随之而来。同样，与纳维-斯托克斯方程有相似之处，该方程（具有稳定边界条件）在足够低的雷诺数下具有稳定解，但在高 Re 下具有混沌、湍流解。关于洛伦兹方程、动力系统和方程、动力系统和混沌的进一步讨论包含在古根海默和福尔摩斯 (1983)、穆恩 (1992) 和格莱克 (1988)。

3.2 随机变量的表征
对于层流，我们可以使用理论（即纳维-斯托克斯方程）来计算 U（在指定位置和时间的速度的特定分量），并且我们可以进行实验来测量 U。根据一个世纪的经验，我们非常有信心 U 的计算值和测量值会一致（在较小的数值和实验误差范围内）。

纳维-斯托克斯方程同样适用于湍流，但这里理论的目标必须不同。由于 U 是一个随机变量，它的值本质上是不可预测的：预测 U 的特定值的理论几乎肯定是错误的。然而，理论可以旨在确定诸如 A=\{U
在本节中，我们开发用于表征随机变量（例如 U）的概念和工具。特别是 U 完全由其概率密度函数 (PDF) 表征。湍流中的随机速度场 U(x, ¢) 是一个比单个随机变量 U 复杂得多的数学对象。在后续章节中，我们将介绍一些用于表征随机变量集（例如 U;、U2 和 U;）、时间随机函数（例如 U(t)）和位置随机函数的量。

\begin{sourceblock}
\includegraphics[page=66,pagebox=mediabox,trim={64.39bp 1297.84bp 64.39bp 109.44bp},clip,width=\textwidth,max height=0.72\textheight]{Pope.pdf}
\par\vspace{0.45\baselineskip}
\small 图 3.3。 U 样本空间的草图，显示与事件对应的区域 (a) B= \{U < Vp\} 和 (bl) C= \{Yas U
\end{sourceblock}
样本空间 为了能够讨论比 A r \{ U < 10 m s-I\} 更一般的事件，我们引入一个独立的速度变量 V，称为
对应于 U 的变量。如图 3.3 所示，不同的样本空间事件，例如

\begin{sourceblock}
\includegraphics[page=66,pagebox=mediabox,trim={64.39bp 895.89bp 64.39bp 1046.87bp},clip,width=0.94\textwidth,max height=0.38\textheight]{Pope.pdf}
\end{sourceblock}
\begin{sourceblock}
\includegraphics[page=66,pagebox=mediabox,trim={64.39bp 813.90bp 64.39bp 1128.86bp},clip,width=0.94\textwidth,max height=0.38\textheight]{Pope.pdf}
\end{sourceblock}
对应于样本空间的不同区域。

例如，事件 B 的概率可写为

\begin{sourceblock}
\includegraphics[page=66,pagebox=mediabox,trim={64.39bp 543.93bp 64.39bp 1398.82bp},clip,width=0.94\textwidth,max height=0.38\textheight]{Pope.pdf}
\end{sourceblock}
目前，读者对概率的直观理解就足够了：p是一个实数（0 - < p 1），表示事件发生的可能性。因为不可能的事件为零；对于确定事件 p 是统一的。 （概率将在 3.8 节中进一步讨论。）

累积分布函数 任何事件的概率可以由下式确定
累积分布

(CDF)，由bution函数定义

\begin{sourceblock}
\includegraphics[page=66,pagebox=mediabox,trim={64.39bp 119.39bp 64.39bp 1839.77bp},clip,width=0.94\textwidth,max height=0.38\textheight]{Pope.pdf}
\end{sourceblock}
例如，我们有

\begin{sourceblock}
\includegraphics[page=67,pagebox=mediabox,trim={64.39bp 1713.78bp 64.39bp 228.97bp},clip,width=0.94\textwidth,max height=0.38\textheight]{Pope.pdf}
\end{sourceblock}
\begin{sourceblock}
\includegraphics[page=67,pagebox=mediabox,trim={64.39bp 1557.80bp 64.39bp 387.95bp},clip,width=0.94\textwidth,max height=0.38\textheight]{Pope.pdf}
\end{sourceblock}
CDF 的三个基本属性是

\begin{sourceblock}
\includegraphics[page=67,pagebox=mediabox,trim={64.39bp 1408.82bp 64.39bp 537.93bp},clip,width=0.94\textwidth,max height=0.38\textheight]{Pope.pdf}
\end{sourceblock}
因为 \{ U < --a i \} 是不可能的；

\begin{sourceblock}
\includegraphics[page=67,pagebox=mediabox,trim={64.39bp 1258.84bp 64.39bp 687.91bp},clip,width=0.94\textwidth,max height=0.38\textheight]{Pope.pdf}
\end{sourceblock}
因为 \{ U < co\} 是确定的；并且，

\begin{sourceblock}
\includegraphics[page=67,pagebox=mediabox,trim={64.39bp 1108.86bp 64.39bp 837.89bp},clip,width=0.94\textwidth,max height=0.38\textheight]{Pope.pdf}
\end{sourceblock}
因为每个事件的概率都是非负的，即

\begin{sourceblock}
\includegraphics[page=67,pagebox=mediabox,trim={64.39bp 957.88bp 64.39bp 985.88bp},clip,width=0.94\textwidth,max height=0.38\textheight]{Pope.pdf}
\end{sourceblock}
第三个性质（方程（3.12））表达了 CDF 是非递减函数的事实。

概率密度函数 (PDF) 被定义为
概率密度函数CDF：

\begin{sourceblock}
\includegraphics[page=67,pagebox=mediabox,trim={64.39bp 575.93bp 64.39bp 1370.83bp},clip,width=0.94\textwidth,max height=0.38\textheight]{Pope.pdf}
\end{sourceblock}
从 CDF 的属性可以简单看出 P D F 是非负的

\begin{sourceblock}
\includegraphics[page=67,pagebox=mediabox,trim={64.39bp 380.95bp 64.39bp 1557.80bp},clip,width=0.94\textwidth,max height=0.38\textheight]{Pope.pdf}
\end{sourceblock}
它满足归一化条件

\begin{sourceblock}
\includegraphics[page=67,pagebox=mediabox,trim={64.39bp 204.97bp 64.39bp 1698.79bp},clip,width=0.94\textwidth,max height=0.38\textheight]{Pope.pdf}
\end{sourceblock}
f(-a) f ( m ) 0。 (3.13) 由此可见，概率

\begin{sourceblock}
\includegraphics[page=68,pagebox=mediabox,trim={64.39bp 837.90bp 64.39bp 109.44bp},clip,width=\textwidth,max height=0.72\textheight]{Pope.pdf}
\par\vspace{0.45\baselineskip}
\small 图 3.4。 (a) 随机变量 U 的 CDF 草图显示了事件 C = \{V, < U < V4\} 的概率，(b) 相应的 PDF。 (b)中的阴影区域是C的概率。
\end{sourceblock}
该间隔内的 PDF：

\begin{sourceblock}
\includegraphics[page=68,pagebox=mediabox,trim={64.39bp 480.94bp 64.39bp 1356.83bp},clip,width=0.94\textwidth,max height=0.38\textheight]{Pope.pdf}
\end{sourceblock}
图 3-4 提供了该方程的图形解释。

对于无穷小区间，方程(3.17) 变为

\begin{sourceblock}
\includegraphics[page=68,pagebox=mediabox,trim={64.39bp 257.97bp 64.39bp 1635.79bp},clip,width=0.94\textwidth,max height=0.38\textheight]{Pope.pdf}
\end{sourceblock}
因此 P D F f ( V ) 是样本空间中每单位距离的概率
因此，术语“概率密度函数”。 PDF f ( V ) 有
是无量纲的。在变量变化的情况下，密度（例如 PDF）的变换规则与函数的变换规则不同：参见第 49 页的练习 3.9。

需要强调的是，P D F (V)（或同等的CDF）完全表征了随机变量U。具有相同P D F 的两个或多个随机变量被认为是同分布的，或者等效地统计上相同的。

均值和特征 随机变量 U 的平均值（或期望）定义为

\begin{sourceblock}
\includegraphics[page=69,pagebox=mediabox,trim={64.39bp 1302.84bp 64.39bp 597.92bp},clip,width=0.94\textwidth,max height=0.38\textheight]{Pope.pdf}
\end{sourceblock}
它是 U 所有可能值的概率加权平均值。更一般地，如果 Q(U) 是 U 的函数，则 Q(U) 的平均值为

\begin{sourceblock}
\includegraphics[page=69,pagebox=mediabox,trim={64.39bp 1100.86bp 64.39bp 802.90bp},clip,width=0.94\textwidth,max height=0.38\textheight]{Pope.pdf}
\end{sourceblock}
即使没有明确说明条件，也应该理解（此处和下文），仅当方程 1 中的积分存在时，均值 (Q(U)) 才存在。 (3.20)绝对收敛。

获取手段的规则非常简单。如果 Q(U) 和 R(U) 是 U 的函数，并且 a 和 b 是常数，则

\begin{sourceblock}
\includegraphics[page=69,pagebox=mediabox,trim={64.39bp 788.90bp 64.39bp 1138.86bp},clip,width=0.94\textwidth,max height=0.38\textheight]{Pope.pdf}
\end{sourceblock}
这可以很容易地从等式中得到验证。 （3.20）。因此，尖括号 ( ) 的行为相当于线性运算符。虽然 U、Q(U) 和 R ( U ) 都是随机变量，但 (U)、(Q(U)) 和 (R(U)) 不是。因此，平均值的平均值是平均值：((U))(U)。

U 的波动定义为

\begin{sourceblock}
\includegraphics[page=69,pagebox=mediabox,trim={64.39bp 496.94bp 64.39bp 1449.82bp},clip,width=0.94\textwidth,max height=0.38\textheight]{Pope.pdf}
\end{sourceblock}
方差定义为均方波动：

\begin{sourceblock}
\includegraphics[page=69,pagebox=mediabox,trim={64.39bp 345.96bp 64.39bp 1587.80bp},clip,width=0.94\textwidth,max height=0.38\textheight]{Pope.pdf}
\end{sourceblock}
方差的平方根是标准偏差

\begin{sourceblock}
\includegraphics[page=69,pagebox=mediabox,trim={64.39bp 204.97bp 64.39bp 1737.78bp},clip,width=0.94\textwidth,max height=0.38\textheight]{Pope.pdf}
\end{sourceblock}
也用 u' 和 表示，也称为 r.m.s. （根
第 n 个中心矩定义为

\begin{sourceblock}
\includegraphics[page=70,pagebox=mediabox,trim={64.39bp 1695.79bp 64.39bp 209.97bp},clip,width=0.94\textwidth,max height=0.38\textheight]{Pope.pdf}
\end{sourceblock}
显然我们有 p(, = 11,p 1 = 0, p2 = 0,。

（相反，关于原点的第 n 个矩或第 n 个原始矩定义为 (Un)。）

Q 和 R 是随机变量，a 和 b 是常数，
使用方程式(3.20) 验证关系

\begin{sourceblock}
\includegraphics[page=70,pagebox=mediabox,trim={64.39bp 1301.84bp 64.39bp 642.92bp},clip,width=0.94\textwidth,max height=0.38\textheight]{Pope.pdf}
\end{sourceblock}
\begin{sourceblock}
\includegraphics[page=70,pagebox=mediabox,trim={64.39bp 1216.85bp 64.39bp 726.91bp},clip,width=0.94\textwidth,max height=0.38\textheight]{Pope.pdf}
\end{sourceblock}
\begin{sourceblock}
\includegraphics[page=70,pagebox=mediabox,trim={64.39bp 1133.86bp 64.39bp 809.90bp},clip,width=0.94\textwidth,max height=0.38\textheight]{Pope.pdf}
\end{sourceblock}
\begin{sourceblock}
\includegraphics[page=70,pagebox=mediabox,trim={64.39bp 1050.87bp 64.39bp 892.89bp},clip,width=0.94\textwidth,max height=0.38\textheight]{Pope.pdf}
\end{sourceblock}
令 Q 定义为

\begin{sourceblock}
\includegraphics[page=70,pagebox=mediabox,trim={64.39bp 874.89bp 64.39bp 1063.87bp},clip,width=0.94\textwidth,max height=0.38\textheight]{Pope.pdf}
\end{sourceblock}
其中 U 是随机变量，a 和 b 是常数。表明

\begin{sourceblock}
\includegraphics[page=70,pagebox=mediabox,trim={64.39bp 740.91bp 64.39bp 1203.85bp},clip,width=0.94\textwidth,max height=0.38\textheight]{Pope.pdf}
\end{sourceblock}
\begin{sourceblock}
\includegraphics[page=70,pagebox=mediabox,trim={64.39bp 652.92bp 64.39bp 1280.84bp},clip,width=0.94\textwidth,max height=0.38\textheight]{Pope.pdf}
\end{sourceblock}
\begin{sourceblock}
\includegraphics[page=70,pagebox=mediabox,trim={64.39bp 573.93bp 64.39bp 1368.83bp},clip,width=0.94\textwidth,max height=0.38\textheight]{Pope.pdf}
\end{sourceblock}
还表明

\begin{sourceblock}
\includegraphics[page=70,pagebox=mediabox,trim={64.39bp 445.94bp 64.39bp 1494.81bp},clip,width=0.94\textwidth,max height=0.38\textheight]{Pope.pdf}
\end{sourceblock}
标准化 使用标准化随机变量通常很方便，根据定义，标准化随机变量的均值和单位方差为零。 U对应的标准化随机变量为

\begin{sourceblock}
\includegraphics[page=70,pagebox=mediabox,trim={64.39bp 119.39bp 64.39bp 1834.77bp},clip,width=0.94\textwidth,max height=0.38\textheight]{Pope.pdf}
\end{sourceblock}
\section*{3.3 概率分布示例}\phantomsection\addcontentsline{toc}{section}{3.3 概率分布示例}
其 P D F U 的标准化 PDF 为

\begin{sourceblock}
\includegraphics[page=71,pagebox=mediabox,trim={64.39bp 1724.78bp 64.39bp 218.97bp},clip,width=0.94\textwidth,max height=0.38\textheight]{Pope.pdf}
\end{sourceblock}
fi U 的标准化矩为

\begin{sourceblock}
\includegraphics[page=71,pagebox=mediabox,trim={64.39bp 1562.80bp 64.39bp 327.96bp},clip,width=0.94\textwidth,max height=0.38\textheight]{Pope.pdf}
\end{sourceblock}
f i fiz 显然我们有 1、0 和 1。 第三个标准化矩 f i fi4
称为偏度，第四个是适应度或峰度。

证明 U 和 Q 的标准化矩（定义为
等式。 (3.30)) 是相同的。

特征函数 随机变量 U 的特征函数定义为

\begin{sourceblock}
\includegraphics[page=71,pagebox=mediabox,trim={64.39bp 1024.87bp 64.39bp 875.89bp},clip,width=0.94\textwidth,max height=0.38\textheight]{Pope.pdf}
\end{sourceblock}
可以认识到，方程中的积分。 (3.38) 是傅里叶逆变换：Y(s) 和 f ( V ) 形成傅里叶变换对，因此它们包含相同的信息。

特征函数是一种数学手段，有助于一些推导和证明。它的属性在附录 I 中描述。特征函数在第 12 章中被广泛使用，但之前没有。因此，可以推迟对附录一的研究。

3.3 概率分布的例子 为了巩固所发展的概念，并说明一些性质上不同的行为，我们现在给出一些概率分布的具体例子。这些分布将在后面的章节中遇到。

均匀分布 如果 U 在区间 a < V < b 内均匀分布，则 U 的 PDF

\begin{sourceblock}
\includegraphics[page=71,pagebox=mediabox,trim={64.39bp 155.98bp 64.39bp 1667.79bp},clip,width=0.94\textwidth,max height=0.38\textheight]{Pope.pdf}
\end{sourceblock}
图 3.5。均匀随机变量的 CDF (a) 和 PDF (b)（方程（3.39））。

对于均匀分布方程。 （3.39）表明
(c) ji3 0，以及
指数分布 如果 U 随参数 I 呈指数分布，则其 PDF（见图 3.6）

图 3.6。指数分布随机变量的 CDF (a) 和 PDF (b)

对于指数分布方程。 （3.40）表明
(a) 满足归一化条件，
(d) F ( V ) 1 exp(-V/?,)，对于 V > 0，
(e) Prob\{U 2 a?,) = e-"，对于 a 2 0。

概率论中最重要的正态分布是正态分布或高斯分布。如果 U 服从均值和标准差 a 的正态分布，则 U 的 PDF 为

\begin{sourceblock}
\includegraphics[page=73,pagebox=mediabox,trim={64.39bp 119.39bp 64.39bp 1807.77bp},clip,width=0.94\textwidth,max height=0.38\textheight]{Pope.pdf}
\end{sourceblock}
这里 ( V 2 ) - 或有时 N(p, a2) 表示平均值为 p 且方差为 a2 的正态分布。我们也可以写

\begin{sourceblock}
\includegraphics[page=74,pagebox=mediabox,trim={64.39bp 1684.79bp 64.39bp 263.97bp},clip,width=0.94\textwidth,max height=0.38\textheight]{Pope.pdf}
\end{sourceblock}
表示 U 是正态随机变量，即
U 的 PDF 分布相等，由式 (1) 给出。 （3.41）。

如果 U 根据方程呈正态分布： (3.41) 那么

\begin{sourceblock}
\includegraphics[page=74,pagebox=mediabox,trim={64.39bp 1466.82bp 64.39bp 475.94bp},clip,width=0.94\textwidth,max height=0.38\textheight]{Pope.pdf}
\end{sourceblock}
是具有 PDF 的标准化高斯随机变量

\begin{sourceblock}
\includegraphics[page=74,pagebox=mediabox,trim={64.39bp 1295.84bp 64.39bp 598.92bp},clip,width=0.94\textwidth,max height=0.38\textheight]{Pope.pdf}
\end{sourceblock}
此 PDF 和相应的 CDF

\begin{sourceblock}
\includegraphics[page=74,pagebox=mediabox,trim={64.39bp 1116.86bp 64.39bp 761.90bp},clip,width=0.94\textwidth,max height=0.38\textheight]{Pope.pdf}
\end{sourceblock}
如图3.7所示。

通过考虑数量

\begin{sourceblock}
\includegraphics[page=74,pagebox=mediabox,trim={64.39bp 838.89bp 64.39bp 1060.87bp},clip,width=0.94\textwidth,max height=0.38\textheight]{Pope.pdf}
\end{sourceblock}
获得标准化矩的递推关系

高斯分布。证明奇矩 ( , . . . ) 是
零，峰度为

\begin{sourceblock}
\includegraphics[page=74,pagebox=mediabox,trim={64.39bp 637.92bp 64.39bp 1308.83bp},clip,width=0.94\textwidth,max height=0.38\textheight]{Pope.pdf}
\end{sourceblock}
那就是

\begin{sourceblock}
\includegraphics[page=74,pagebox=mediabox,trim={64.39bp 504.94bp 64.39bp 1442.82bp},clip,width=0.94\textwidth,max height=0.38\textheight]{Pope.pdf}
\end{sourceblock}
对数正态分布 我们再次将 U 视为均值 p 和方差 a2 的正态分布。那么正随机变量

\begin{sourceblock}
\includegraphics[page=74,pagebox=mediabox,trim={64.39bp 213.97bp 64.39bp 1730.78bp},clip,width=0.94\textwidth,max height=0.38\textheight]{Pope.pdf}
\end{sourceblock}
根据定义，是对数正态分布。

\begin{sourceblock}
\includegraphics[page=75,pagebox=mediabox,trim={64.39bp 798.90bp 64.39bp 109.44bp},clip,width=\textwidth,max height=0.72\textheight]{Pope.pdf}
\par\vspace{0.45\baselineskip}
\small 图 3.7。标准化高斯随机变量的 CDF (a) 和 PDF (b)。
\end{sourceblock}
即 F ( V ) 和 f ( V ) 由方程式给出。 （3.41）。由于 Y 为正，因此样本空间可以取为正实线，即 y 2 0。从 CDF 的定义出发，我们得到

\begin{sourceblock}
\includegraphics[page=75,pagebox=mediabox,trim={64.39bp 465.94bp 64.39bp 1425.82bp},clip,width=0.94\textwidth,max height=0.38\textheight]{Pope.pdf}
\end{sourceblock}
然后通过对 y 求导即可获得 P D F：

\begin{sourceblock}
\includegraphics[page=75,pagebox=mediabox,trim={64.39bp 196.97bp 64.39bp 1608.80bp},clip,width=0.94\textwidth,max height=0.38\textheight]{Pope.pdf}
\end{sourceblock}
图 3.8 显示了 (Y) 1 和 (Y) 1 的 PDF f y ( y ) 和 CDF F y ( y )

\begin{sourceblock}
\includegraphics[page=76,pagebox=mediabox,trim={64.39bp 874.89bp 64.39bp 109.44bp},clip,width=\textwidth,max height=0.72\textheight]{Pope.pdf}
\par\vspace{0.45\baselineskip}
\small 图 3.8。对数范数的 CDF (a) 和 PDF (b)！随机变量 Y，其中 (¥) = 1 且 var(Y) = 4 和 5。
\end{sourceblock}
生成不同形状的 PDF。特别是，较大的 a' 值会导致 PDF 具有长尾，这一点在 CDF 缓慢的统一方法中最为明显。如 L xercise 3.7 所示，归一化方差 var( Y ( Y ) ) 随着 eu 增加。

方程(3.50)和(3.5 1 )说明了PDF和CDF的变换规则。这些在练习 3.9 中得到进一步发展。

证明 Y 的原始矩（由方程 (3.49) 定义）为

\begin{sourceblock}
\includegraphics[page=76,pagebox=mediabox,trim={64.39bp 289.96bp 64.39bp 1657.79bp},clip,width=0.94\textwidth,max height=0.38\textheight]{Pope.pdf}
\end{sourceblock}
（提示：评估 J" e 11 I> ' .. f ' (V) d V。）

表明规格

\begin{sourceblock}
\includegraphics[page=76,pagebox=mediabox,trim={64.39bp 119.39bp 64.39bp 1840.77bp},clip,width=0.94\textwidth,max height=0.38\textheight]{Pope.pdf}
\end{sourceblock}
导致 ( Y ) 为 1，并且 Y 的方差为

\begin{sourceblock}
\includegraphics[page=77,pagebox=mediabox,trim={64.39bp 1725.78bp 64.39bp 218.97bp},clip,width=0.94\textwidth,max height=0.38\textheight]{Pope.pdf}
\end{sourceblock}
随机变量 Z 定义为

\begin{sourceblock}
\includegraphics[page=77,pagebox=mediabox,trim={64.39bp 1597.80bp 64.39bp 342.96bp},clip,width=0.94\textwidth,max height=0.38\textheight]{Pope.pdf}
\end{sourceblock}
其中 Y 是对数正态随机变量，a 和 b 为正数
常数。证明 Z 也是对数正态分布

\begin{sourceblock}
\includegraphics[page=77,pagebox=mediabox,trim={64.39bp 1425.82bp 64.39bp 515.93bp},clip,width=0.94\textwidth,max height=0.38\textheight]{Pope.pdf}
\end{sourceblock}
随机变量U具有CDF F(V)和PDF(V)。的
随机变量 Y 定义为

\begin{sourceblock}
\includegraphics[page=77,pagebox=mediabox,trim={64.39bp 1209.85bp 64.39bp 736.91bp},clip,width=0.94\textwidth,max height=0.38\textheight]{Pope.pdf}
\end{sourceblock}
其中 Q(V) 是单调递增函数。继
方程式中的步骤。 (3.50) 和 (3.51) 表明 CDF Fy(v) 和 P D F
对于 Y 由下式给出

\begin{sourceblock}
\includegraphics[page=77,pagebox=mediabox,trim={64.39bp 825.90bp 64.39bp 957.88bp},clip,width=0.94\textwidth,max height=0.38\textheight]{Pope.pdf}
\end{sourceblock}
\begin{sourceblock}
\includegraphics[page=77,pagebox=mediabox,trim={64.39bp 763.90bp 64.39bp 1183.85bp},clip,width=0.94\textwidth,max height=0.38\textheight]{Pope.pdf}
\end{sourceblock}
表明 Q(V) 的相应结果是单调的
递减函数是

\begin{sourceblock}
\includegraphics[page=77,pagebox=mediabox,trim={64.39bp 585.93bp 64.39bp 1353.83bp},clip,width=0.94\textwidth,max height=0.38\textheight]{Pope.pdf}
\end{sourceblock}
\begin{sourceblock}
\includegraphics[page=77,pagebox=mediabox,trim={64.39bp 471.94bp 64.39bp 1425.82bp},clip,width=0.94\textwidth,max height=0.38\textheight]{Pope.pdf}
\end{sourceblock}
表明方程。 (3.59) 和 (3.62) 可以写成共同的形式

\begin{sourceblock}
\includegraphics[page=77,pagebox=mediabox,trim={64.39bp 349.96bp 64.39bp 1595.80bp},clip,width=0.94\textwidth,max height=0.38\textheight]{Pope.pdf}
\end{sourceblock}
其中 dV 和

\begin{sourceblock}
\includegraphics[page=77,pagebox=mediabox,trim={64.39bp 192.98bp 64.39bp 1697.79bp},clip,width=0.94\textwidth,max height=0.38\textheight]{Pope.pdf}
\end{sourceblock}
是相应的 inhnitesin\textasciitilde{}ali 区间。

\begin{sourceblock}
\includegraphics[page=78,pagebox=mediabox,trim={64.39bp 885.89bp 64.39bp 109.44bp},clip,width=\textwidth,max height=0.72\textheight]{Pope.pdf}
\par\vspace{0.45\baselineskip}
\small 图 3.9。均值为 1、方差为 o 的伽玛分布的 CDF (a) 和 PDF (b)？ = 4、4 和 5。
\end{sourceblock}
伽马分布 正随机变量 U（均值 p 和方差 a2）具有伽马分布，如果其 PDF 为
其中 a 定义为

\begin{sourceblock}
\includegraphics[page=78,pagebox=mediabox,trim={64.39bp 372.95bp 64.39bp 1536.81bp},clip,width=0.94\textwidth,max height=0.38\textheight]{Pope.pdf}
\end{sourceblock}
T(a) 是伽马函数

\begin{sourceblock}
\includegraphics[page=78,pagebox=mediabox,trim={64.39bp 204.97bp 64.39bp 1697.79bp},clip,width=0.94\textwidth,max height=0.38\textheight]{Pope.pdf}
\end{sourceblock}
对于 1，这成为指数分布，并且值
方差）PDF 在原点处为零，而对于较小的 a 值，它是无限的 - 如图 3.9 所示。

使用替换 x aV/p 来显示归一化的原始数据
伽马分布的矩是

\begin{sourceblock}
\includegraphics[page=79,pagebox=mediabox,trim={64.39bp 1374.83bp 64.39bp 430.95bp},clip,width=0.94\textwidth,max height=0.38\textheight]{Pope.pdf}
\end{sourceblock}
验证此结果对于 n 0、1 和 2 的一致性。

Delta 函数分布 假设 U 是一个随机变量，其取值 a 的概率为 p，取值 b (b > a) 的概率为 1 p。推导 U 的 CDF 很简单：

见图3.10。这可以用 Heavis,ide 函数写成

\begin{sourceblock}
\includegraphics[page=79,pagebox=mediabox,trim={64.39bp 779.90bp 64.39bp 1157.85bp},clip,width=0.94\textwidth,max height=0.38\textheight]{Pope.pdf}
\end{sourceblock}
相应的PDF（通过微分方程（3.70）获得）是

\begin{sourceblock}
\includegraphics[page=79,pagebox=mediabox,trim={64.39bp 650.92bp 64.39bp 1293.84bp},clip,width=0.94\textwidth,max height=0.38\textheight]{Pope.pdf}
\end{sourceblock}
见图3.10。 （狄拉克 δ 函数和赫维赛德函数的性质在附录 C 中进行了回顾。）

只能取有限数量值的随机变量是离散随机变量（与连续随机变量相对）。尽管本节中介绍的工具旨在描述连续随机变量，但显然（借助 Heaviside 和 Dirac delta 函数）也可以处理离散随机变量。此外，如果 U 是确定变量，其值为 a 的概率为 1，则其 CDF 和 P D F 始终由下式给出

\begin{sourceblock}
\includegraphics[page=79,pagebox=mediabox,trim={64.39bp 119.39bp 64.39bp 1749.78bp},clip,width=0.94\textwidth,max height=0.38\textheight]{Pope.pdf}
\end{sourceblock}
\begin{sourceblock}
\includegraphics[page=80,pagebox=mediabox,trim={64.39bp 782.90bp 64.39bp 109.44bp},clip,width=\textwidth,max height=0.72\textheight]{Pope.pdf}
\par\vspace{0.45\baselineskip}
\small 图 3.10。离散随机变量 U 的 CDF (a) 和 PDF (b)，方程： （3.69）。
\end{sourceblock}
\begin{sourceblock}
\includegraphics[page=80,pagebox=mediabox,trim={64.39bp 443.94bp 64.39bp 1470.81bp},clip,width=0.94\textwidth,max height=0.38\textheight]{Pope.pdf}
\end{sourceblock}
\begin{sourceblock}
\includegraphics[page=80,pagebox=mediabox,trim={64.39bp 291.96bp 64.39bp 1622.80bp},clip,width=0.94\textwidth,max height=0.38\textheight]{Pope.pdf}
\end{sourceblock}
方差的最大可能值发生在

采用双δ函数分布

\begin{sourceblock}
\includegraphics[page=81,pagebox=mediabox,trim={64.39bp 1676.79bp 64.39bp 254.97bp},clip,width=0.94\textwidth,max height=0.38\textheight]{Pope.pdf}
\end{sourceblock}
对于这个分布表明

\begin{sourceblock}
\includegraphics[page=81,pagebox=mediabox,trim={64.39bp 1525.81bp 64.39bp 395.95bp},clip,width=0.94\textwidth,max height=0.38\textheight]{Pope.pdf}
\end{sourceblock}
\begin{sourceblock}
\includegraphics[page=81,pagebox=mediabox,trim={64.39bp 1454.82bp 64.39bp 480.94bp},clip,width=0.94\textwidth,max height=0.38\textheight]{Pope.pdf}
\end{sourceblock}
注意：对于 0、1，这些结果是 p 和
柯西分布 均值、方差和其他矩被定义为 PDF 的积分（方程（3.20））。我们隐含地假设所有此类积分收敛；事实上，除了少数例外，湍流研究中遇到的 PDF 都是如此。有一个简单的反例很有用：这是由柯西分布提供的。

以 c 为中心、半宽为 w 的柯西分布的 PDF

\begin{sourceblock}
\includegraphics[page=81,pagebox=mediabox,trim={64.39bp 874.89bp 64.39bp 1027.87bp},clip,width=0.94\textwidth,max height=0.38\textheight]{Pope.pdf}
\end{sourceblock}
对于大 V f 随 V 变化，因此 Vf (V) 的积分发散为 In V 因此，尽管分布关于其中心 V = c 对称，但均值（由方程 (3.19) 定义）不存在。方差是

图 3.11 显示了柯西密度（方程（3.79））和相应的

\begin{sourceblock}
\includegraphics[page=81,pagebox=mediabox,trim={64.39bp 443.94bp 64.39bp 1371.83bp},clip,width=0.94\textwidth,max height=0.38\textheight]{Pope.pdf}
\end{sourceblock}
图 3.12 中描绘的 PDF 具有均值零和单位方差
它是标准化的）。表明草图中定义的变量是

\begin{sourceblock}
\includegraphics[page=81,pagebox=mediabox,trim={64.39bp 152.98bp 64.39bp 1719.78bp},clip,width=0.94\textwidth,max height=0.38\textheight]{Pope.pdf}
\end{sourceblock}
\begin{sourceblock}
\includegraphics[page=82,pagebox=mediabox,trim={64.39bp 845.89bp 64.39bp 109.44bp},clip,width=\textwidth,max height=0.72\textheight]{Pope.pdf}
\par\vspace{0.45\baselineskip}
\small 图 3.11。 c=0 时柯西分布（方程（3.79）和（3.80））的 CDF (a) 和 PDF (b)，
\end{sourceblock}
\section*{3.4 联合随机变量}\phantomsection\addcontentsline{toc}{section}{3.4 联合随机变量}
在本节中，将单个随机变量 U 获得的结果扩展到两个或多个随机变量。我们以湍流中特定位置和时间的速度分量 (U1, U2,U 3) 为例。

对应于随机变量 U = \{U1,U 2,U 3\}a 的样本空间变量重新表示为 V = \{ V , , V3\}。对于 U1 和 U2 两个分量， \textasciitilde{} i 3.1\textasciitilde{}3 显示了由 N = 100 个点组成的散点图
v2）。 = ( u , (n ) , u ( 2 n) ) ,n = 1 ? 2、**N*、？其中 (u,("), U2(n) ) 是第 n 次重复实验时的 (U1, U2) 值。联合随机变量 (U1, U2) 的 CDF 定义为

\begin{sourceblock}
\includegraphics[page=83,pagebox=mediabox,trim={64.39bp 1461.82bp 64.39bp 109.44bp},clip,width=\textwidth,max height=0.72\textheight]{Pope.pdf}
\par\vspace{0.45\baselineskip}
\small 图 3.12。练习 3.13 中标准化 PDF 的草图。
\end{sourceblock}
\begin{sourceblock}
\includegraphics[page=83,pagebox=mediabox,trim={64.39bp 800.90bp 64.39bp 563.93bp},clip,width=\textwidth,max height=0.72\textheight]{Pope.pdf}
\par\vspace{0.45\baselineskip}
\small 图 3.13。联合随机变量 (U1, U2) 的 100 个样本的样本空间中的散点图。 （在此示例中，U 和 U; 共同正规，(U;) = 2、(U2) = 1、(uj) = 1、uh) = 且 pra = 1/5。）
\end{sourceblock}
图 3.14 的阴影区域。显然，F 1 2V2) 是其每个参数的非减函数：

\begin{sourceblock}
\includegraphics[page=83,pagebox=mediabox,trim={64.39bp 492.94bp 64.39bp 1452.82bp},clip,width=0.94\textwidth,max height=0.38\textheight]{Pope.pdf}
\end{sourceblock}
CDF 的其他性质是

\begin{sourceblock}
\includegraphics[page=83,pagebox=mediabox,trim={64.39bp 350.96bp 64.39bp 1589.80bp},clip,width=0.94\textwidth,max height=0.38\textheight]{Pope.pdf}
\end{sourceblock}
因为\{U1 是不可能的；和

\begin{sourceblock}
\includegraphics[page=83,pagebox=mediabox,trim={64.39bp 162.98bp 64.39bp 1728.78bp},clip,width=0.94\textwidth,max height=0.38\textheight]{Pope.pdf}
\end{sourceblock}
图3.1-4。 V1-V2样本空间显示与事件对应的区域
U2（在方程（3.85）中定义）称为CDF。同样，边际
U1 的边际 CD F 为 F1(Vl) F12(Vl, a)。

U1 和 U2 的联合 PDF (JPDF) 定义为

\begin{sourceblock}
\includegraphics[page=84,pagebox=mediabox,trim={64.39bp 957.88bp 64.39bp 967.88bp},clip,width=0.94\textwidth,max height=0.38\textheight]{Pope.pdf}
\end{sourceblock}
其基本属性如图 3.1 5 所示，是

\begin{sourceblock}
\includegraphics[page=84,pagebox=mediabox,trim={64.39bp 782.90bp 64.39bp 1098.86bp},clip,width=0.94\textwidth,max height=0.38\textheight]{Pope.pdf}
\end{sourceblock}
可以很容易地推断出的其他属性是

\begin{sourceblock}
\includegraphics[page=84,pagebox=mediabox,trim={64.39bp 672.91bp 64.39bp 1273.84bp},clip,width=0.94\textwidth,max height=0.38\textheight]{Pope.pdf}
\end{sourceblock}
\begin{sourceblock}
\includegraphics[page=84,pagebox=mediabox,trim={64.39bp 561.93bp 64.39bp 1341.83bp},clip,width=0.94\textwidth,max height=0.38\textheight]{Pope.pdf}
\end{sourceblock}
\begin{sourceblock}
\includegraphics[page=84,pagebox=mediabox,trim={64.39bp 448.94bp 64.39bp 1453.82bp},clip,width=0.94\textwidth,max height=0.38\textheight]{Pope.pdf}
\end{sourceblock}
其中 f 是 U2 的 PDF。

如果 Q(U1,U 2) 是随机变量的函数，则其均值定义为

\begin{sourceblock}
\includegraphics[page=84,pagebox=mediabox,trim={64.39bp 270.97bp 64.39bp 1673.79bp},clip,width=0.94\textwidth,max height=0.38\textheight]{Pope.pdf}
\end{sourceblock}
\textasciitilde{}\textasciitilde{}\textasciitilde{} e(aUnl) sa 和 (U2)，以及方差 (u:) 和 (ut)，可以根据该方程确定，或者同样地，根据边际 PDF f ( Vl ) 和 f V2) 确定

\begin{sourceblock}
\includegraphics[page=85,pagebox=mediabox,trim={64.39bp 1281.84bp 64.39bp 109.44bp},clip,width=\textwidth,max height=0.72\textheight]{Pope.pdf}
\par\vspace{0.45\baselineskip}
\small 图 3.15。 Vi-V> 样本空间显示对应于事件 \{Vig < U1 < Viv) Vea < Uz < Voy\} 的区域，参见方程 1。 （3.87）。
\end{sourceblock}
U1和U2的协方差是混合二阶矩

\begin{sourceblock}
\includegraphics[page=85,pagebox=mediabox,trim={64.39bp 985.88bp 64.39bp 962.88bp},clip,width=0.94\textwidth,max height=0.38\textheight]{Pope.pdf}
\end{sourceblock}
相关系数为

\begin{sourceblock}
\includegraphics[page=85,pagebox=mediabox,trim={64.39bp 858.89bp 64.39bp 1065.87bp},clip,width=0.94\textwidth,max height=0.38\textheight]{Pope.pdf}
\end{sourceblock}
如图 3.13 中的散点图所示，当一个随机变量相对于均值的正偏移（例如 ul > 0）优先与另一个随机变量的正偏移（例如 u2 0）相关时，就会出现正相关系数。相反，如果 ul 的正偏移优先与 u2 的负偏移相关，如图 3.16 所示，则相关系数为负。一般来说，我们有柯西-施瓦茨不等式

\begin{sourceblock}
\includegraphics[page=85,pagebox=mediabox,trim={64.39bp 515.93bp 64.39bp 1431.82bp},clip,width=0.94\textwidth,max height=0.38\textheight]{Pope.pdf}
\end{sourceblock}
参见练习 3.16。

如果相关系数p12为零（这意味着协方差（u1u2）为零），则随机变量U1和U2不相关。相反，如果 pl2 为 1，则 U1 和 U2 完全相关；并且，如果 pl2 等于
1、它们完全负相关。练习 3.17 中给出了这些相关性的示例。

对于图 3.16 所示的散点图，很明显，具有 U1 V 的样本和具有 U1 Vlb 的样本很可能具有显着差异

\begin{sourceblock}
\includegraphics[page=86,pagebox=mediabox,trim={64.39bp 1057.87bp 64.39bp 109.44bp},clip,width=\textwidth,max height=0.72\textheight]{Pope.pdf}
\par\vspace{0.45\baselineskip}
\small 图 3.16。负相关随机变量的散点图 ((U;) = 1、(U2) = -1、= 2、() = 12 和 px = -V/2/3）。
\end{sourceblock}
\begin{sourceblock}
\includegraphics[page=86,pagebox=mediabox,trim={64.39bp 441.94bp 64.39bp 1011.87bp},clip,width=\textwidth,max height=0.72\textheight]{Pope.pdf}
\par\vspace{0.45\baselineskip}
\small 图 3.17。图 3.16 所示的分布的联合 PDF 是针对 VY = Vig = 1 且 Vy = Vip = S 的 V2 绘制的。
\end{sourceblock}
= V1 和 Vl -- V 对于固定 V f 12(Vla,V 2) 表示 U2 如何针对具有 U1 V 的样本 (U1,U 2) 进行分布 这些想法通过定义条件 PDF 得以精确化； U1 = V 条件下 U2 的 PDF 为

\begin{sourceblock}
\includegraphics[page=86,pagebox=mediabox,trim={64.39bp 119.39bp 64.39bp 1839.77bp},clip,width=0.94\textwidth,max height=0.38\textheight]{Pope.pdf}
\end{sourceblock}
这只是缩放后的联合 PDF，以便满足归一化

\begin{sourceblock}
\includegraphics[page=87,pagebox=mediabox,trim={64.39bp 1662.79bp 64.39bp 199.97bp},clip,width=0.94\textwidth,max height=0.38\textheight]{Pope.pdf}
\end{sourceblock}
对于给定，如果 fl (V l) 为零，则 f211(V2/2)Vi1s) 未定义。否则，很容易验证 f 2 1 1 ( V 2 V ls)a 满足 PDF 的所有条件（即，它是非负的，并且满足归一化条件，方程（3.96））。 （关于符号的一个词：V 是 U1 = V 的缩写，读作“以 U1 = V 为条件”或“给定 U1 = V 或“给定 Vl”。）

对于函数 Q(U1,U 2)，条件均值（以 Vl 为条件）（QI 定义为

\begin{sourceblock}
\includegraphics[page=87,pagebox=mediabox,trim={64.39bp 1233.84bp 64.39bp 666.92bp},clip,width=0.94\textwidth,max height=0.38\textheight]{Pope.pdf}
\end{sourceblock}
独立的概念至关重要。如果 U1 和 U2 是独立的，那么对其中一个的值的了解不会提供有关另一个的信息。因此，“条件”没有效果，并且条件概率密度函数和边际概率密度函数是相同的：

\begin{sourceblock}
\includegraphics[page=87,pagebox=mediabox,trim={64.39bp 971.88bp 64.39bp 975.88bp},clip,width=0.94\textwidth,max height=0.38\textheight]{Pope.pdf}
\end{sourceblock}
因此（根据方程（3.95））联合 PDF 是边际的乘积：

\begin{sourceblock}
\includegraphics[page=87,pagebox=mediabox,trim={64.39bp 830.89bp 64.39bp 1115.86bp},clip,width=0.94\textwidth,max height=0.38\textheight]{Pope.pdf}
\end{sourceblock}
独立随机变量不相关；但一般来说，反之则不然。

显示联合 PDF 方程的属性。 (3.87)-(3.90) 以下
低于 CDF（方程（3.82））和联合 PDF 的定义
证明，对于单独的 U1 的函数 R(U1)，
联合 PDF（方程（3.91））的平均值 (R(U1))i 是一致的
其定义为边际 PDF（方程（3.20））。

通过考虑数量 (ul /u', u2/u12)2，建立柯西-

\begin{sourceblock}
\includegraphics[page=87,pagebox=mediabox,trim={64.39bp 210.97bp 64.39bp 1677.79bp},clip,width=0.94\textwidth,max height=0.38\textheight]{Pope.pdf}
\end{sourceblock}
你在哪里；和你；是 U1 和 U2 的标准差，p12 是

\begin{sourceblock}
\includegraphics[page=88,pagebox=mediabox,trim={64.39bp 1689.79bp 64.39bp 162.98bp},clip,width=0.94\textwidth,max height=0.38\textheight]{Pope.pdf}
\end{sourceblock}
其中 a、b 和 c 是常数。证明相关系数

\begin{sourceblock}
\includegraphics[page=88,pagebox=mediabox,trim={64.39bp 1473.81bp 64.39bp 379.95bp},clip,width=0.94\textwidth,max height=0.38\textheight]{Pope.pdf}
\end{sourceblock}
因此证明 U1 和 U2 是
(a) 不相关 ( p 0) 如果 b 为零并且 c 不为零，
(b) 如果 c 为零且 b 为正，则完全相关 (p12 1)，并且
(c) 完全负相关 (p12 -1) 如果 为零且 b 为
对两个随机变量求和，得到结果

\begin{sourceblock}
\includegraphics[page=88,pagebox=mediabox,trim={64.39bp 1070.86bp 64.39bp 873.89bp},clip,width=0.94\textwidth,max height=0.38\textheight]{Pope.pdf}
\end{sourceblock}
对 N 个独立随机变量求和，得到结果

\begin{sourceblock}
\includegraphics[page=88,pagebox=mediabox,trim={64.39bp 862.89bp 64.39bp 1014.87bp},clip,width=0.94\textwidth,max height=0.38\textheight]{Pope.pdf}
\end{sourceblock}
设 U1 为标准化高斯随机变量，U2 为
由 U2 U1 定义 绘制 (U1,U 2) 的可能值
VI-V2样本空间。证明 U1 和 U2 不相关。争论
U2 的条件 PDF 为

\begin{sourceblock}
\includegraphics[page=88,pagebox=mediabox,trim={64.39bp 588.93bp 64.39bp 1357.83bp},clip,width=0.94\textwidth,max height=0.38\textheight]{Pope.pdf}
\end{sourceblock}
因此 U2 和 U1 不是独立的。

对于任何函数 R(U1)，从等式开始： (3.97)，验证结果

\begin{sourceblock}
\includegraphics[page=88,pagebox=mediabox,trim={64.39bp 393.95bp 64.39bp 1554.80bp},clip,width=0.94\textwidth,max height=0.38\textheight]{Pope.pdf}
\end{sourceblock}
证明无条件均值可以从条件中获得

概念上的意思是

\begin{sourceblock}
\includegraphics[page=88,pagebox=mediabox,trim={64.39bp 161.98bp 64.39bp 1743.78bp},clip,width=0.94\textwidth,max height=0.38\textheight]{Pope.pdf}
\end{sourceblock}
\section*{3.5 正态分布与联合正态分布}\phantomsection\addcontentsline{toc}{section}{3.5 正态分布与联合正态分布}
3.5 正态分布和联合正态分布 在本节中，我们介绍中心极限定理，该定理（除其他外）表明正态分布或高斯分布（方程（3.41））在概率论中起着核心作用。然后描述了联合正态分布及其特殊性质。给出的许多结果最容易通过特征函数获得（附录 I）。

我们首先检查整体平均值。令 U 表示可重复湍流实验中特定位置和时间的速度分量，并令 U”) 表示第 n 次重复的 U。每次重复都是在相同的标称条件下进行的，并且不同重复之间不存在依赖性。因此，随机变量 U, U®) 是独立的并且具有相同的分布（即 U 的分布）：它们被称为独立且同分布 (1.1.d.)。

整体平均值（超过 N 次重复）定义为

\begin{sourceblock}
\includegraphics[page=89,pagebox=mediabox,trim={64.39bp 1094.86bp 64.39bp 787.90bp},clip,width=0.94\textwidth,max height=0.38\textheight]{Pope.pdf}
\end{sourceblock}
集合平均值本身就是一个随机变量，很容易证明它的均值和方差为

\begin{sourceblock}
\includegraphics[page=89,pagebox=mediabox,trim={64.39bp 922.88bp 64.39bp 1022.87bp},clip,width=0.94\textwidth,max height=0.38\textheight]{Pope.pdf}
\end{sourceblock}
\begin{sourceblock}
\includegraphics[page=89,pagebox=mediabox,trim={64.39bp 808.90bp 64.39bp 1091.86bp},clip,width=0.94\textwidth,max height=0.38\textheight]{Pope.pdf}
\end{sourceblock}
因此（参见练习 3.22）\& 定义为

\begin{sourceblock}
\includegraphics[page=89,pagebox=mediabox,trim={64.39bp 674.91bp 64.39bp 1263.84bp},clip,width=0.94\textwidth,max height=0.38\textheight]{Pope.pdf}
\end{sourceblock}
是标准化随机变量（例如，(0) = 0，(0？) = 1）。

中心极限定理指出，当 N 趋于无穷大时，U,f(V) 的 PDF 趋于标准化正态分布

\begin{sourceblock}
\includegraphics[page=89,pagebox=mediabox,trim={64.39bp 411.95bp 64.39bp 1483.81bp},clip,width=0.94\textwidth,max height=0.38\textheight]{Pope.pdf}
\end{sourceblock}
（参见第 47 页的图 3.7 和第 709 页的练习 I.3），该结果取决于 \{U'),U®, UY?独立同分布。但它对基础随机变量 U 的唯一限制是它具有有限方差。

现在我们转向联合正态分布，它在概率论和湍流中都很重要。例如，在均匀湍流的实验中，速度分量和守恒被动
第 175 页）。联合正态分布的定义和性质

\{ul,u 2 arle 现在给出一组 D 随机变量 U UD \}

, 4 + * , 对于 D 2 或 3，U 可以被认为是湍流中的速度分量。

使用矩阵表示法很方便。的均值和波动
U 由随机向量表示

\begin{sourceblock}
\includegraphics[page=90,pagebox=mediabox,trim={64.39bp 1502.81bp 64.39bp 441.94bp},clip,width=0.94\textwidth,max height=0.38\textheight]{Pope.pdf}
\end{sourceblock}
\begin{sourceblock}
\includegraphics[page=90,pagebox=mediabox,trim={64.39bp 1403.82bp 64.39bp 530.93bp},clip,width=0.94\textwidth,max height=0.38\textheight]{Pope.pdf}
\end{sourceblock}
(对称 D x D) 协方差矩阵为

\begin{sourceblock}
\includegraphics[page=90,pagebox=mediabox,trim={64.39bp 1262.84bp 64.39bp 678.91bp},clip,width=0.94\textwidth,max height=0.38\textheight]{Pope.pdf}
\end{sourceblock}
如果 U \{U1,U 2, U3\} 是速度，则协方差矩阵是一个二阶张量，其分量为 Cij
如果 U \{U1,u 2 , ..., UD\} 是联合正态分布，则（根据定义）其联合 PDF 为

\begin{sourceblock}
\includegraphics[page=90,pagebox=mediabox,trim={64.39bp 980.88bp 64.39bp 955.88bp},clip,width=0.94\textwidth,max height=0.38\textheight]{Pope.pdf}
\end{sourceblock}
请注意，联合 PDF 的 V 相关性包含在二次方程中

\begin{sourceblock}
\includegraphics[page=90,pagebox=mediabox,trim={64.39bp 804.90bp 64.39bp 1140.86bp},clip,width=0.94\textwidth,max height=0.38\textheight]{Pope.pdf}
\end{sourceblock}
对于D 2，对应于恒定概率密度的恒定g值是V1-V2平面中的椭圆。对于 D 3，恒定概率表面是 V 空间中的椭球体。

我们现在更详细地检查联合正态随机变量对 \{U1,U 2\}（即 D 2）。图 3.18 显示了 p 和 C 的特定选择的散点图和常数概率密度线。

根据方差 (uf) 和 (u) 以及相关系数，联合正态 PDF（方程 (3.116)）为

\begin{sourceblock}
\includegraphics[page=90,pagebox=mediabox,trim={64.39bp 175.98bp 64.39bp 1771.78bp},clip,width=0.94\textwidth,max height=0.38\textheight]{Pope.pdf}
\end{sourceblock}
\begin{sourceblock}
\includegraphics[page=91,pagebox=mediabox,trim={64.39bp 1296.84bp 64.39bp 109.44bp},clip,width=\textwidth,max height=0.72\textheight]{Pope.pdf}
\par\vspace{0.45\baselineskip}
\small 图 3.18。平面中联合正态随机变量 (U;, U2) 的散点图和常概率密度线，其中 (Uy) = 2、(U2) = 1、(wt) = 1、(uj) = 4 且 py = 1/,/5。
\end{sourceblock}
(i) U1 和 U2 (f 和 ( V ) ) 的边际 PDF 呈高斯分布。

(ii) 如果 U1 和 U2 不相关 ( i e p12 0)，那么它们也处于-
依赖（从那时起 f 12(V1)，这是一个特殊的
联合正态分布的性质：一般来说，缺乏相关性
化并不意味着独立。

(iii) U1 的条件平均值为

\begin{sourceblock}
\includegraphics[page=91,pagebox=mediabox,trim={64.39bp 765.90bp 64.39bp 1132.86bp},clip,width=0.94\textwidth,max height=0.38\textheight]{Pope.pdf}
\end{sourceblock}
(iv) U1 的条件方差为

\begin{sourceblock}
\includegraphics[page=91,pagebox=mediabox,trim={64.39bp 627.92bp 64.39bp 1311.83bp},clip,width=0.94\textwidth,max height=0.38\textheight]{Pope.pdf}
\end{sourceblock}
(v) 条件 PDF f V2) 是高斯分布的。

回到 U \{U1,U 2, . 的一般情况。。 UD\} 是联合法线，通过考虑 U 的线性变换可以获得额外的见解。一个基本结果（参见附录 I）是，如果 U 是联合法线，则由 U 的一般线性变换形成的随机向量也是联合法线。

由于协方差矩阵 C 是对称的，因此可以通过酉变换对它进行对角化，酉变换由酉矩阵 A 定义。（酉矩阵的属性是

\begin{sourceblock}
\includegraphics[page=91,pagebox=mediabox,trim={64.39bp 167.98bp 64.39bp 1776.78bp},clip,width=0.94\textwidth,max height=0.38\textheight]{Pope.pdf}
\end{sourceblock}
\begin{sourceblock}
\includegraphics[page=92,pagebox=mediabox,trim={64.39bp 1734.78bp 64.39bp 158.98bp},clip,width=0.94\textwidth,max height=0.38\textheight]{Pope.pdf}
\end{sourceblock}
其中 A 是包含 C 特征值的对角矩阵

\begin{sourceblock}
\includegraphics[page=92,pagebox=mediabox,trim={64.39bp 1525.81bp 64.39bp 351.96bp},clip,width=0.94\textwidth,max height=0.38\textheight]{Pope.pdf}
\end{sourceblock}
因此变换后的随机向量

\begin{sourceblock}
\includegraphics[page=92,pagebox=mediabox,trim={64.39bp 1314.83bp 64.39bp 633.92bp},clip,width=0.94\textwidth,max height=0.38\textheight]{Pope.pdf}
\end{sourceblock}
有一个对角协方差矩阵 A：

\begin{sourceblock}
\includegraphics[page=92,pagebox=mediabox,trim={64.39bp 1166.85bp 64.39bp 770.90bp},clip,width=0.94\textwidth,max height=0.38\textheight]{Pope.pdf}
\end{sourceblock}
从这一转变中可以进行一些观察并推断出结果。

(i) 如果 U 是速度矢量，则 ii 是 a 中的脉动速度
特定坐标系，即 (uiuj)+ 的主轴
(ii) C、Ai 的特征值是

\begin{sourceblock}
\includegraphics[page=92,pagebox=mediabox,trim={64.39bp 815.90bp 64.39bp 1127.86bp},clip,width=0.94\textwidth,max height=0.38\textheight]{Pope.pdf}
\end{sourceblock}
（其中括号后缀被排除在求和约定之外
化）。因此，由于每个特征值都是非负的，因此 C 是对称的
正半负债 e.

(iii) 协方差矩阵是对角的表明反式
形成随机变量 \{\&, G2, .。 , G3\} 不相关。

无论 U 是否为联合正态，这三个观察结果都适用。此外我们还有以下
(iv) 如果 U 是联合法线，则 \{fi1,fi2,,4 ti+D+\} 是独立高斯分布
随机变量。

从集合平均值的定义（方程（3.108））可知

\section*{3.6 随机过程}\phantomsection\addcontentsline{toc}{section}{3.6 随机过程}
从而验证方程。 3 10 提示：

\begin{sourceblock}
\includegraphics[page=93,pagebox=mediabox,trim={64.39bp 1669.79bp 64.39bp 213.97bp},clip,width=0.94\textwidth,max height=0.38\textheight]{Pope.pdf}
\end{sourceblock}
获得峰度的显式表达式：

N 和 U 的峰度。根据以下结果对结果进行评论
中心极限定理。

表明，对于较大的 N，集合平均值（方程（3.108））可以是
其中 u' sdev(U) 是标准化高斯随机变量
令 U 为联合正态随机向量，均值为 p 且为正值
定协方差矩阵 C A A A \textasciitilde{} w，这里 A 是酉且 A 是
对角线。表明随机变量
是一个标准化的联合正态，即它的均值为零，恒等协方差
ance，以及联合PDF

\begin{sourceblock}
\includegraphics[page=93,pagebox=mediabox,trim={64.39bp 861.89bp 64.39bp 1038.87bp},clip,width=0.94\textwidth,max height=0.38\textheight]{Pope.pdf}
\end{sourceblock}
高斯随机数生成器产生一系列输入

相关标准化高斯随机数：( 1 ) , 5(3)

如何使用这些来生成联合正态随机向量 U
具有指定的均值 p 和协方差矩阵 C？

（提示：这可以通过多种方式实现，其中最好的方式是
涉及：Cholesky 分解，即对称半定
矩阵可以分解为 C LL"，其中 L 是下三角。）

3.6 随机过程 作为随机变量的一个例子，我们考虑（在第 3.2 节中）在可重复的湍流实验中，在特定位置和时间（相对于实验的开始）速度 U 的一个分量。随机变量 U 完全由其 PDF f (V) 表征。现在考虑相同的速度，但作为时间的函数，即 U(t)。这样的时刻-

\begin{sourceblock}
\includegraphics[page=94,pagebox=mediabox,trim={64.39bp 1330.83bp 64.39bp 109.44bp},clip,width=\textwidth,max height=0.72\textheight]{Pope.pdf}
\par\vspace{0.45\baselineskip}
\small 图 3.19。湍流实验的三次重复的 U(t) 的样本路径。
\end{sourceblock}
在不同的实验重复中获得的样本路径（U(t) 的 e 值）。

如何表征随机过程？在每个时间点，随机变量 U(t) 的特征在于其一次性 CDF

\begin{sourceblock}
\includegraphics[page=94,pagebox=mediabox,trim={64.39bp 1000.87bp 64.39bp 941.88bp},clip,width=0.94\textwidth,max height=0.38\textheight]{Pope.pdf}
\end{sourceblock}
或者，等效地，通过一次性 PDF

\begin{sourceblock}
\includegraphics[page=94,pagebox=mediabox,trim={64.39bp 834.89bp 64.39bp 1074.86bp},clip,width=0.94\textwidth,max height=0.38\textheight]{Pope.pdf}
\end{sourceblock}
然而，这些量不包含两次或多次有关 U(t) 的联合信息。为了说明这一限制，图 3.20 显示了五个不同随机过程的示例路径，每个过程都具有相同的一次性 PDF。显然，完全不同的行为（定性和定量）是可能的，但一次性 PDF 无法表示。过程 U(t) 的 N 次联合 CDF 定义为

\begin{sourceblock}
\includegraphics[page=94,pagebox=mediabox,trim={64.39bp 434.94bp 64.39bp 1512.81bp},clip,width=0.94\textwidth,max height=0.38\textheight]{Pope.pdf}
\end{sourceblock}
其中 \{tl,t 2, .。 , t \} 是指定的时间点，fN(Vl,t l; V2, t2; +,.V N 是相应的 N 时间联合 PDF。为了完全表征随机过程，需要知道所有时刻的联合 PDF，这通常是一项不可能的任务。

如果过程在统计上是平稳的，那么就会发生相当大的简化，许多（但肯定不是全部）湍流也是如此。一个过程是统计的

\begin{sourceblock}
\includegraphics[page=95,pagebox=mediabox,trim={64.39bp 606.92bp 64.39bp 231.77bp},clip,width=\textwidth,max height=0.72\textheight]{Pope.pdf}
\par\vspace{0.45\baselineskip}
\small 图 3.20。五个统计平稳随机过程的样本路径。每个的一次性 PDF 都是标准化高斯分布。 (a) 测量的湍流速度。 (b) 测量的湍流速度的频率高于 (a) 的频率。 (c) 具有与 (a) 相同频谱的高斯过程。 (d) 具有与 (a) 相同的积分时间尺度的 Ornstein-Uhlenbeck 过程（参见第 12 章）。 (e) 跳转过程
\end{sourceblock}
具有与(d)相同的光谱。

对于所有正时间间隔 T，以及 \{tl , t 2, ... 的所有选择。。。 , t \} , 我们有

\begin{sourceblock}
\includegraphics[page=95,pagebox=mediabox,trim={64.39bp 207.97bp 64.39bp 1737.78bp},clip,width=0.94\textwidth,max height=0.38\textheight]{Pope.pdf}
\end{sourceblock}
层流启动后，它可以通过初始层流

\begin{sourceblock}
\includegraphics[page=96,pagebox=mediabox,trim={64.39bp 1435.82bp 64.39bp 109.44bp},clip,width=\textwidth,max height=0.72\textheight]{Pope.pdf}
\par\vspace{0.45\baselineskip}
\small 图 3.21。过程的平均值 (U(t))（实线）和方差 (u(\#)*) 所示
\end{sourceblock}
与时间无关。湍流在经过初始瞬态期后可以达到统计静止状态，其中即使流变量（例如 U(t)）随时间变化，统计数据也与时间无关。图 3.19 所示的流程就是这种情况。均值(U(t))和方差(\textasciitilde{}(t)o\textasciitilde{}f)th过程如图3.21所示。显然，在 t 5 之后，统计数据变得与时间无关，即使过程本身 U(t) 继续发生显着变化。

对于统计平稳过程，可以考虑的最简单的多次统计量是自协方差

\begin{sourceblock}
\includegraphics[page=96,pagebox=mediabox,trim={64.39bp 863.89bp 64.39bp 1081.86bp},clip,width=0.94\textwidth,max height=0.38\textheight]{Pope.pdf}
\end{sourceblock}
或者，以标准化形式，自相关函数

\begin{sourceblock}
\includegraphics[page=96,pagebox=mediabox,trim={64.39bp 727.91bp 64.39bp 1213.85bp},clip,width=0.94\textwidth,max height=0.38\textheight]{Pope.pdf}
\end{sourceblock}
其中 u(t) U(t) - ( U ) 是波动。 （请注意，考虑到假定的统计平稳性，均值 (U)、方差 (u2)、R(s) 和 p(s) 不依赖于 t。）自相关函数是时间 t 和 t s 过程之间的相关系数。因此它具有以下属性

\begin{sourceblock}
\includegraphics[page=96,pagebox=mediabox,trim={64.39bp 474.94bp 64.39bp 1471.81bp},clip,width=0.94\textwidth,max height=0.38\textheight]{Pope.pdf}
\end{sourceblock}
\begin{sourceblock}
\includegraphics[page=96,pagebox=mediabox,trim={64.39bp 414.95bp 64.39bp 1531.81bp},clip,width=0.94\textwidth,max height=0.38\textheight]{Pope.pdf}
\end{sourceblock}
此外，放入 t' t s，我们得到

\begin{sourceblock}
\includegraphics[page=96,pagebox=mediabox,trim={64.39bp 226.97bp 64.39bp 1663.79bp},clip,width=0.94\textwidth,max height=0.38\textheight]{Pope.pdf}
\end{sourceblock}
即 p(s) 是偶函数。

如果 U(t) 是周期性的，周期为 T ( e mU (t T) U(t))，那么也是
P ( \textasciitilde{} () i . p (s T) p(s))。然而，对于在湍流中出现的过程，我们预计相关性会随着滞后时间 s 的增加而减小。通常 p(s) 下降得足够快，以至于积分

\begin{sourceblock}
\includegraphics[page=97,pagebox=mediabox,trim={64.39bp 1592.80bp 64.39bp 308.96bp},clip,width=0.94\textwidth,max height=0.38\textheight]{Pope.pdf}
\end{sourceblock}
收敛：则 7 是该过程的积分时间尺度。

图 3.22 显示了图 3.20 中给出的五个过程的自相关函数。特别注意，高频过程 (b) 具有较窄的自相关函数（因此比低频过程 (a) 更小。通过构造，过程 (c) 具有与 (a) 相同的自相关函数。过程 (d) 和 (e) 都具有自相关函数 p(s) exp(-Isl/\textasciitilde{})，与过程 (a) 具有相同的积分时间尺度。因此，除了 (b) 之外，所有过程都具有相同的积分时间尺度。

自协方差 R(s) (u(t)u(t s)) ( \textasciitilde{} ( t ) \textasciitilde{} ) \textasciitilde{}an(ds )(t wice) 频谱 E ( o ) 形成傅里叶变换对：

\begin{sourceblock}
\includegraphics[page=97,pagebox=mediabox,trim={64.39bp 884.89bp 64.39bp 1010.87bp},clip,width=0.94\textwidth,max height=0.38\textheight]{Pope.pdf}
\end{sourceblock}
\begin{sourceblock}
\includegraphics[page=97,pagebox=mediabox,trim={64.39bp 616.92bp 64.39bp 1284.84bp},clip,width=0.94\textwidth,max height=0.38\textheight]{Pope.pdf}
\end{sourceblock}
（傅里叶变换的定义和性质在附录 D 中给出。）显然，R(s) 和 E ( o ) 包含相同的信息，只是形式不同。因为 R(s) 是实数且偶数，所以 E(o) 也是实数且偶数。

正如附录 E 中更全面讨论的那样，速度波动 u(t) 具有频谱表示形式，即不同频率 o 的傅里叶模式的加权和，即 eio' cos ( a t ) 位于 (ot) 中。频谱的基本属性是（对于 o, ob）e 积分

\begin{sourceblock}
\includegraphics[page=97,pagebox=mediabox,trim={64.39bp 168.98bp 64.39bp 1728.78bp},clip,width=0.94\textwidth,max height=0.38\textheight]{Pope.pdf}
\end{sourceblock}
\begin{sourceblock}
\includegraphics[page=98,pagebox=mediabox,trim={64.39bp 180.98bp 64.39bp 657.72bp},clip,width=\textwidth,max height=0.72\textheight]{Pope.pdf}
\par\vspace{0.45\baselineskip}
\small 图 3.22。过程的自相关函数如图 3.20 所示。如插图所示，对于过程 (a) 和 (c)，自相关函数在原点处是平滑的。
\end{sourceblock}
\begin{sourceblock}
\includegraphics[page=99,pagebox=mediabox,trim={64.39bp 1188.85bp 64.39bp 109.44bp},clip,width=\textwidth,max height=0.72\textheight]{Pope.pdf}
\par\vspace{0.45\baselineskip}
\small 图 3.23。过程谱如图 3.20 所示。
\end{sourceblock}
w, o < 特别是方差是

\begin{sourceblock}
\includegraphics[page=99,pagebox=mediabox,trim={64.39bp 964.88bp 64.39bp 937.88bp},clip,width=0.94\textwidth,max height=0.38\textheight]{Pope.pdf}
\end{sourceblock}
从方程式可以看出(3.141) 为 0。

频谱和自相关之间的另一个简单联系是积分时间尺度由下式给出

\begin{sourceblock}
\includegraphics[page=99,pagebox=mediabox,trim={64.39bp 719.91bp 64.39bp 1203.85bp},clip,width=0.94\textwidth,max height=0.38\textheight]{Pope.pdf}
\end{sourceblock}
通过在方程式中设置 w = 0 可以很容易地验证这一点。 （3.140）。附录 bm 中给出了频谱表示和频谱解释的更完整解释

图 3.23 显示了以下给出的平稳随机过程的谱
过程。然而，应当理解，一次性PDF和自相关函数仅提供该过程的部分特征。这一点在过程(d)和(e)中得到了充分的证明。

function (p(s) e ) 重复一下，一般来说，一次性 P D F 和自相关函数不能完全表征随机过程。

高斯过程是一个重要但非常特殊的情况。如果一个过程是高斯过程，那么根据定义，一般的 N-ti:me PDF（方程（3.133））是联合正态的。现在，联合正态分布完全由其均值 ( ( t ) ) 及其协方差 (u(tn)u(tm)F) 来表征。或者是一个统计平稳过程，我们有

\begin{sourceblock}
\includegraphics[page=100,pagebox=mediabox,trim={64.39bp 1454.82bp 64.39bp 488.94bp},clip,width=0.94\textwidth,max height=0.38\textheight]{Pope.pdf}
\end{sourceblock}
因此，统计平稳高斯过程完全由其均值 ( \textasciitilde{} ( t ) )it,s 方差 ( \textasciitilde{} ( t ) \textasciitilde{}an) d、自相关函数 p(s) （或等效的频谱 E (o)）来表征。

在图 3.20 中，过程 (c) 被定义为高斯过程，其频谱与湍流速度过程 (a) 的频谱相同。过程（a）和（c）之间的一些差异可能是显而易见的；这些差异可以通过检查 \textasciitilde{} ( t ) d 2u (t)/dt2 的样本路径来清楚地揭示，见图 3.24。对于高斯过程 (c)，它遵循 \textasciitilde{} ( t i)s 也是高斯分布，因此峰度为 \textasciitilde{} ( t i)s 3。然而，对于湍流速度，过程 (a)， \textasciitilde{} ( tis) 远离高斯分布，并且具有峰度
由湍流引起的随机过程（例如过程 (a)）是可微分的，即对于每个样本路径，存在以下限制：

\begin{sourceblock}
\includegraphics[page=100,pagebox=mediabox,trim={64.39bp 737.91bp 64.39bp 1147.86bp},clip,width=0.94\textwidth,max height=0.38\textheight]{Pope.pdf}
\end{sourceblock}
在这种情况下，取均值并取极限通勤，以便
dU(t) U(t At) U(t)))

\begin{sourceblock}
\includegraphics[page=100,pagebox=mediabox,trim={64.39bp 339.96bp 64.39bp 1593.80bp},clip,width=0.94\textwidth,max height=0.38\textheight]{Pope.pdf}
\end{sourceblock}
综合过程（例如过程（d）和（e））不需要是可微的（极限方程（3.146）不存在）。可以观察到，这些过程的频谱在高频处以 E ( o ) 的形式衰减，并且（相应地）它们的自相关函数 p(s) e 不是

图 3.24。 \textasciitilde{} ( t f) 或过程 (a) 和 (c) 的示例路径如图 3.20 所示。

过程 (d) 是一个 Ornstein-Uhlenbeck 过程，它是扩散过程的典型示例。此类过程用于 PDF 方法，并在第 12 章和附录 J 中进行了描述。

在下面的练习中，u(t) 是一个零均值、统计平稳、可微的随机过程，具有自协方差 R(s)、自相关函数 p(s) 和谱 E (a)。

证明 u(t) 和 L(t) 不相关，并且 u(t) 和 u(t) 是
负相关。

表明

\begin{sourceblock}
\includegraphics[page=101,pagebox=mediabox,trim={64.39bp 212.97bp 64.39bp 1683.79bp},clip,width=0.94\textwidth,max height=0.38\textheight]{Pope.pdf}
\end{sourceblock}
证明原点 (s = 0) dR(s)/ ds 为零，且 d 2 \textasciitilde{} ( s )d/s 为
证明过程 u(t) 的自协方差 B(s) 为

\begin{sourceblock}
\includegraphics[page=102,pagebox=mediabox,trim={64.39bp 1711.78bp 64.39bp 210.97bp},clip,width=0.94\textwidth,max height=0.38\textheight]{Pope.pdf}
\end{sourceblock}
证明 u(t) 的积分时间尺度为零。

证明u(t) 的频谱为0 2 \textasciitilde{} ( o )。

如果 u(t) 是高斯过程，则证明

\begin{sourceblock}
\includegraphics[page=102,pagebox=mediabox,trim={64.39bp 1480.81bp 64.39bp 463.94bp},clip,width=0.94\textwidth,max height=0.38\textheight]{Pope.pdf}
\end{sourceblock}
\begin{sourceblock}
\includegraphics[page=102,pagebox=mediabox,trim={64.39bp 1399.82bp 64.39bp 540.93bp},clip,width=0.94\textwidth,max height=0.38\textheight]{Pope.pdf}
\end{sourceblock}
\section*{3.7 随机场}\phantomsection\addcontentsline{toc}{section}{3.7 随机场}
在湍流中，速度 U(x,t) 是一个与时间相关的随机向量 fie1.d。它可以被描述，即部分地通过前面部分中介绍的工具的扩展来表征。

一点统计 速度的一点一次性联合 CDF 为

\begin{sourceblock}
\includegraphics[page=102,pagebox=mediabox,trim={64.39bp 882.89bp 64.39bp 1059.87bp},clip,width=0.94\textwidth,max height=0.38\textheight]{Pope.pdf}
\end{sourceblock}
然后联合 PDF 是

\begin{sourceblock}
\includegraphics[page=102,pagebox=mediabox,trim={64.39bp 711.91bp 64.39bp 1212.85bp},clip,width=0.94\textwidth,max height=0.38\textheight]{Pope.pdf}
\end{sourceblock}
在每个点和时间，该 PDF 都充分表征了随机速度矢量，但它不包含两个或多个时间或位置的联合信息。就该 PDF 而言，平均速度场为

\begin{sourceblock}
\includegraphics[page=102,pagebox=mediabox,trim={64.39bp 444.94bp 64.39bp 1458.82bp},clip,width=0.94\textwidth,max height=0.38\textheight]{Pope.pdf}
\end{sourceblock}
\begin{sourceblock}
\includegraphics[page=102,pagebox=mediabox,trim={64.39bp 333.96bp 64.39bp 1569.80bp},clip,width=0.94\textwidth,max height=0.38\textheight]{Pope.pdf}
\end{sourceblock}
该方程的第二行引入了一个缩写符号： ) d V 写为
脉动速度场定义为

\begin{sourceblock}
\includegraphics[page=103,pagebox=mediabox,trim={64.39bp 1725.78bp 64.39bp 217.97bp},clip,width=0.94\textwidth,max height=0.38\textheight]{Pope.pdf}
\end{sourceblock}
速度的（一点、一次）协方差为 (ui(x,t )uj(x,t ))。由于下一章给出的原因，这些协方差被称为雷诺应力，并被写为 (uiuj)，并理解对 x 和 t 的依赖性。

关于符号的说明：f ( V ; x, t) 中的分号表示 f 是相对于出现在分号 e Vl、V2 和 V3 左侧的样本空间变量的密度，而 是相对于其余变量（ e xl 、 x2、x3 和 t）的函数。这种区别很有用，因为密度和函数具有不同的变换属性（参见第 49 页的练习 3.9）。

湍流速度场是可微的，并且（如第 3.6 节中讨论的）微分并取平均通程：

\begin{sourceblock}
\includegraphics[page=103,pagebox=mediabox,trim={64.39bp 1081.86bp 64.39bp 814.90bp},clip,width=0.94\textwidth,max height=0.38\textheight]{Pope.pdf}
\end{sourceblock}
\begin{sourceblock}
\includegraphics[page=103,pagebox=mediabox,trim={64.39bp 951.88bp 64.39bp 945.88bp},clip,width=0.94\textwidth,max height=0.38\textheight]{Pope.pdf}
\end{sourceblock}
N 点统计 N 点、N 次联合 P D F 可以定义为式（1）的简单扩展。 （3.132）。令 \{(x("),t(")),n 1,2,..., N\} 为一组指定的位置和时间。然后我们定义

\begin{sourceblock}
\includegraphics[page=103,pagebox=mediabox,trim={64.39bp 682.91bp 64.39bp 1257.84bp},clip,width=0.94\textwidth,max height=0.38\textheight]{Pope.pdf}
\end{sourceblock}
为 U(x, t) 在这 N 个时空点的联合 P D F。确定所有时空点的 N 点 P D F 显然是不可能的，因此实际上无法完全表征随机速度场。

发现湍流速度场不是高斯场：高斯场完全由均值 (U(x, t)) 和自协方差表征
统计平稳性和同质性 如果所有统计量在时间变化下保持不变，则随机场 U(x, t) 是统计平稳的。就 N 点 PDF 而言，这意味着
如果对于所有 N 个点 (x("),t (")) 被替换为 (x(")t, T)，则 fN 不变，其中 T 是时移。

\begin{sourceblock}
\includegraphics[page=104,pagebox=mediabox,trim={64.39bp 1539.81bp 64.39bp 109.44bp},clip,width=\textwidth,max height=0.72\textheight]{Pope.pdf}
\par\vspace{0.45\baselineskip}
\small 图 3.25。湍流槽道流装置的草图。
\end{sourceblock}
在位置的转变下。对于所有 N 个点，如果 (x,t) 替换为 (x + X,1)，则 fy 不变，其中 X 是位置偏移。如果速度场 U(x,t) 在统计上是均匀的，则平均速度 (U) 是均匀的；并且，通过适当选择框架，(U) 可以取为零。均匀湍流的定义限制较少：具体而言，在均匀湍流中，脉动速度场 u(x,t) 在统计上是均匀的。它与平均速度梯度 0(U;)/@x 的定义一致；非零，但均匀（参见第 5.4.5 节）。在风洞实验中可以很好地近似均匀湍流；均匀湍流是使用直接数值模拟研究的最简单的流动类别。

以类似的方式，湍流在统计上可以是二维的或一维的。例如，图3.25是槽道流动装置的草图。对于大纵横比 (b/h >> 1) 且远离端壁 (|x3|/b < 1)，流的统计数据在展向 (x3) 方向上变化很小。那么，在近似范围内，速度场 U(x,t) 在统计上是二维的 - 统计数据与 x3 无关。在通道足够远的地方 (x;/h >> 1)，流动变得（统计上）完全发展。那么速度场在统计上是一维的，统计数据与 x 和 x3 无关。类似地，管道中的湍流在统计上是轴对称的，因为（在极柱坐标中）所有统计数据对于圆周坐标都是独立的。

应该强调的是，即使流在统计上是均匀的或一维的，U(x,t) 的所有三个分量仍然在所有三个坐标方向和时间上发生变化。只是统计量与某些坐标方向无关。

各向同性湍流 根据定义，统计均匀场 U(x, t) 在平移（即坐标系原点的移动）下统计不变。如果
坐标系，那么它（统计上）是各向同性的。各向同性的概念在湍流中极其重要：已经对（近似）各向同性湍流进行了数百个风洞实验，并且许多湍流理论都以它为中心。就 N 点 P D F (Eq. (3.159)) 而言，在各向同性湍流中，如果将 U(x , t(")) 替换为 y u ( x t )，则各向同性湍流不变，
f N (") -(n) 其中，U 表示通过坐标轴的旋转和反射获得的任意坐标系中的位置和速度。

两点相关性包含随机场空间结构一些信息的最简单的统计量是两点一次性自协方差

\begin{sourceblock}
\includegraphics[page=105,pagebox=mediabox,trim={64.39bp 1266.84bp 64.39bp 675.91bp},clip,width=0.94\textwidth,max height=0.38\textheight]{Pope.pdf}
\end{sourceblock}
这通常被称为两点相关性。由此可以定义各种积分长度尺度，例如

\begin{sourceblock}
\includegraphics[page=105,pagebox=mediabox,trim={64.39bp 1039.87bp 64.39bp 856.89bp},clip,width=0.94\textwidth,max height=0.38\textheight]{Pope.pdf}
\end{sourceblock}
其中 el 是 XI 坐标方向的单位向量。

波数谱 对于均匀湍流，两点相关性 Rij(r,t ) 与 x 无关，并且它包含的信息可以用波数谱重新表示。空间傅里叶模式

\begin{sourceblock}
\includegraphics[page=105,pagebox=mediabox,trim={64.39bp 663.92bp 64.39bp 1278.84bp},clip,width=0.94\textwidth,max height=0.38\textheight]{Pope.pdf}
\end{sourceblock}
是在波数向量 r 的方向上呈正弦变化（波长为 2n/IrI）的函数，且 n d 在垂直于 K 的平面中恒定。速度谱张量 mij(r,t) 是两点相关性的傅立叶变换

\begin{sourceblock}
\includegraphics[page=105,pagebox=mediabox,trim={64.39bp 322.96bp 64.39bp 1558.80bp},clip,width=0.94\textwidth,max height=0.38\textheight]{Pope.pdf}
\end{sourceblock}
逆变换是

\begin{sourceblock}
\includegraphics[page=105,pagebox=mediabox,trim={64.39bp 119.39bp 64.39bp 1770.78bp},clip,width=0.94\textwidth,max height=0.38\textheight]{Pope.pdf}
\end{sourceblock}
其中 d r 和 d\textasciitilde{} 分别写为 drl dr2d r3 和 dlcl dK2d \textasciitilde{} res。该方程中的 Slettingr 0 产生

\begin{sourceblock}
\includegraphics[page=106,pagebox=mediabox,trim={64.39bp 1615.80bp 64.39bp 270.97bp},clip,width=0.94\textwidth,max height=0.38\textheight]{Pope.pdf}
\end{sourceblock}
因此 mij(u, t) 表示波数为 K 的速度模态对协方差 (uiuj) 的贡献。

两点相关性和频谱包含两种不同的方向信息。 Rij(r, t) 对 r 的依赖性，以及 , t) 对
给出相关性方向依赖性的信息； Rij 的分量给出了速度方向的信息。

一个有用的量，特别是对于定性讨论来说，是能谱函数：

\begin{sourceblock}
\includegraphics[page=106,pagebox=mediabox,trim={64.39bp 1032.87bp 64.39bp 862.89bp},clip,width=0.94\textwidth,max height=0.38\textheight]{Pope.pdf}
\end{sourceblock}
可以将其视为 (Dij(u,t ) 去除所有方向信息。对所有标量波数 rc 积分方程 (3.166)，得出

\begin{sourceblock}
\includegraphics[page=106,pagebox=mediabox,trim={64.39bp 782.90bp 64.39bp 1119.86bp},clip,width=0.94\textwidth,max height=0.38\textheight]{Pope.pdf}
\end{sourceblock}
因此，E(K,t ) d\textasciitilde{} 表示对湍动能的贡献 1 I K \textasciitilde{} I K I
来自 5 du 范围内的所有模式。湍流中的速度 -2 (uiui) K K 谱在第 6.5 节中进行了详细研究。

根据替换 x' x r 和两点的定义
相关性（方程（3.160）），表明

\begin{sourceblock}
\includegraphics[page=106,pagebox=mediabox,trim={64.39bp 348.96bp 64.39bp 1596.80bp},clip,width=0.94\textwidth,max height=0.38\textheight]{Pope.pdf}
\end{sourceblock}
因此，对于统计上均匀的场，

\begin{sourceblock}
\includegraphics[page=106,pagebox=mediabox,trim={64.39bp 185.98bp 64.39bp 1740.78bp},clip,width=0.94\textwidth,max height=0.38\textheight]{Pope.pdf}
\end{sourceblock}
\section*{3.8 概率与平均}\phantomsection\addcontentsline{toc}{section}{3.8 概率与平均}
相关性（方程（3.160））满足

\begin{sourceblock}
\includegraphics[page=107,pagebox=mediabox,trim={64.39bp 1679.79bp 64.39bp 215.97bp},clip,width=0.94\textwidth,max height=0.38\textheight]{Pope.pdf}
\end{sourceblock}
此外，如果 u(x, t) 在统计上是齐次的，那么

\begin{sourceblock}
\includegraphics[page=107,pagebox=mediabox,trim={64.39bp 1511.81bp 64.39bp 400.95bp},clip,width=0.94\textwidth,max height=0.38\textheight]{Pope.pdf}
\end{sourceblock}
3.8 概率和平均 在开发了描述随机变量、随机过程和随机场的工具之后，我们现在回到起点，以澄清概率的概念，一切都建立在概率的概念之上。诸如密度和速度之类的物理量是可操作地定义的（例如，在第 2. 1 节中，以便（至少在原则上）它们的值可以通过测量来确定。概率的可操作定义，例如时间平均值或集合平均值，尽管经常使用，但并不令人满意。相反，在现代处理中，概率论是公理化的。本节的目的是描述这种公理化方法，并解释与可测量量的联系（例如时间平均值）为了简单起见，我们从抛硬币实验开始讨论。

考虑一枚可以抛掷任意次数的硬币，有两种可能的结果：“正面”和“反面”。我们将变量 p 定义为“正面朝上”的概率。 （假设每次抛掷在统计上都是独立的，并且与其他抛掷没有区别。）

假设耐心地进行了一个将硬币抛掷 N 1,000,000 次的实验。抛掷产生正面的比例是一个随机变量，用 pN 表示。在这个特定的实验中，假设 p\textasciitilde{} 的测量值为 0.5024。

抛硬币实验是伯努利试验的一个例子，伯努利试验有一个完整的理论。例如，假设我们假设硬币是“fan, 9”。 p = 那么简单的统计计算表明（对于 N 1,000,000）p\textasciitilde{} 的概率为 99\% 位于范围内
由于测量值 pN = 0.5024 超出了这个范围，我们可以有
根据 pN 的观测值进行计算表明，在 99\% 的置信度下，p 位于以下范围内
总结一下

(i) p 被定义为“正面”的概率，
(ii) pN 是测量的“正面”频率，
(iii) 给定关于 p 的假设，可以预测 p\textasciitilde{} 的范围，并且
(iv) 给定 PN 的测量值，p 的置信区间可以为
最需要理解的两点是 p 无法测量
只能以一定的置信水平进行估计；并且，虽然当N趋于无穷大时p\textasciitilde{}趋于p，但这并不被视为p的定义。

将速度 U(t) 视为湍流时，我们定义 f ( V ; t) 来消除其 PDF，然后将平均值定义为
在湍流实验和模拟中，多种平均用于定义与 (U(t)) 相关的其他方法。对于统计平稳流，时间平均值（在时间间隔 T 内）定义为

\begin{sourceblock}
\includegraphics[page=108,pagebox=mediabox,trim={64.39bp 799.90bp 64.39bp 1096.86bp},clip,width=0.94\textwidth,max height=0.38\textheight]{Pope.pdf}
\end{sourceblock}
对于可以重复或复制 N 次的流，整体平均值定义为

\begin{sourceblock}
\includegraphics[page=108,pagebox=mediabox,trim={64.39bp 586.93bp 64.39bp 1295.84bp},clip,width=0.94\textwidth,max height=0.38\textheight]{Pope.pdf}
\end{sourceblock}
其中 \textasciitilde{} ( " ) ( ti) 是第 n 个实现的测量值。在 C 边立方域中的均匀湍流模拟中，U(x, t) 的空间平均值由下式定义：

\begin{sourceblock}
\includegraphics[page=108,pagebox=mediabox,trim={64.39bp 353.95bp 64.39bp 1592.80bp},clip,width=0.94\textwidth,max height=0.38\textheight]{Pope.pdf}
\end{sourceblock}
可以为统计上的一维流和二维流定义类似的空间平均值。

这些平均值 (U) 和 (U) La 涉及（如 pN）r 个随机变量。他们
( u ) 可以用来估计，但不能用来确定地测量。最重要的是

\begin{sourceblock}
\includegraphics[page=109,pagebox=mediabox,trim={64.39bp 1128.86bp 64.39bp 109.44bp},clip,width=\textwidth,max height=0.72\textheight]{Pope.pdf}
\par\vspace{0.45\baselineskip}
\small 图 3.26。 Durst 等人测量的速度分布。 (1974) 在矩形管道中对称膨胀下游的稳定层流中。几何形状和边界条件关于平面 y = 0 对称。符号：©，稳定状态 1； A、稳定状态2； e，轮廓 1 绕 y 轴的反射。
\end{sourceblock}
\begin{sourceblock}
\includegraphics[page=109,pagebox=mediabox,trim={64.39bp 161.98bp 64.39bp 1752.78bp},clip,width=0.94\textwidth,max height=0.38\textheight]{Pope.pdf}
\end{sourceblock}
矩形管道中对称扩展的下游流动。铝-
尽管几何形状和边界条件关于

平面 y 0 时，流动不对称。每次流量都是
从静止开始，在初始瞬态之后，流量达到两个之一
稳定的稳态。对于此流程，讨论之间的关系
期望 (U)、时间平均值 (U) T 和整体平均值

\chapter{平均流方程}
\section*{4.1 雷诺平均方程}\phantomsection\addcontentsline{toc}{section}{4.1 雷诺平均方程}
在上一章中，介绍了各种统计量来描述湍流速度场均值、PDF、两点相关性等。从控制底层湍流速度场 U(x, t) 的纳维-斯托克斯方程开始，可以导出所有这些量的演化方程。这些方程中最基本的方程（首先由 Reynolds (1894) 导出）是控制平均速度场 (U(x, t)) 的方程。

将速度 U(x, t) 分解为其平均值 (U(x, t)) 和波动

\begin{sourceblock}
\includegraphics[page=111,pagebox=mediabox,trim={64.39bp 928.88bp 64.39bp 1016.87bp},clip,width=0.94\textwidth,max height=0.38\textheight]{Pope.pdf}
\end{sourceblock}
被称为雷诺分解，即

\begin{sourceblock}
\includegraphics[page=111,pagebox=mediabox,trim={64.39bp 787.90bp 64.39bp 1157.85bp},clip,width=0.94\textwidth,max height=0.38\textheight]{Pope.pdf}
\end{sourceblock}
由连续性方程得出（方程（2.19））

\begin{sourceblock}
\includegraphics[page=111,pagebox=mediabox,trim={64.39bp 642.92bp 64.39bp 1288.84bp},clip,width=0.94\textwidth,max height=0.38\textheight]{Pope.pdf}
\end{sourceblock}
(U(x, t)) 和 u(x, t) 都是螺线管。对于这个方程的平均值

\begin{sourceblock}
\includegraphics[page=111,pagebox=mediabox,trim={64.39bp 473.94bp 64.39bp 1411.82bp},clip,width=0.94\textwidth,max height=0.38\textheight]{Pope.pdf}
\end{sourceblock}
然后通过减法我们得到

\begin{sourceblock}
\includegraphics[page=111,pagebox=mediabox,trim={64.39bp 329.96bp 64.39bp 1598.80bp},clip,width=0.94\textwidth,max height=0.38\textheight]{Pope.pdf}
\end{sourceblock}
（请注意，取平均值和微分交换，使得 (V U) V ( U ) 和 (V u) V a ( u) 0。）

取动量方程（方程（2.35））的平均值就不那么简单了
由于非线性对流项。第一步是以保守形式写出实质性导数，

\begin{sourceblock}
\includegraphics[page=112,pagebox=mediabox,trim={64.39bp 1647.79bp 64.39bp 276.96bp},clip,width=0.94\textwidth,max height=0.38\textheight]{Pope.pdf}
\end{sourceblock}
所以平均值是

\begin{sourceblock}
\includegraphics[page=112,pagebox=mediabox,trim={64.39bp 1481.81bp 64.39bp 407.95bp},clip,width=0.94\textwidth,max height=0.38\textheight]{Pope.pdf}
\end{sourceblock}
然后，将 Ui 和 Uj 的雷诺分解代入，非线性项变为

\begin{sourceblock}
\includegraphics[page=112,pagebox=mediabox,trim={64.39bp 1162.85bp 64.39bp 778.90bp},clip,width=0.94\textwidth,max height=0.38\textheight]{Pope.pdf}
\end{sourceblock}
由于很快就会给出的原因，速度协方差被称为
(tli\textasciitilde{}lj) Revnolds 强调。因此，根据前面的两个方程，我们得到

\begin{sourceblock}
\includegraphics[page=112,pagebox=mediabox,trim={64.39bp 823.90bp 64.39bp 1078.86bp},clip,width=0.94\textwidth,max height=0.38\textheight]{Pope.pdf}
\end{sourceblock}
第二步是 d(Ui)/?xi = 0（方程（4.4））。

最终结果可以通过定义平均实质性导数来有效地重新表达

D Q / D \textasciitilde{} 对于任何属性 Q(x, t)，表示其跟随以局部平均速度 (U(x, t )) 移动的点的变化率。就该导数而言，方程。 (4.9) 是

\begin{sourceblock}
\includegraphics[page=112,pagebox=mediabox,trim={64.39bp 303.96bp 64.39bp 1591.80bp},clip,width=0.94\textwidth,max height=0.38\textheight]{Pope.pdf}
\end{sourceblock}
显然，随体导数 (D u j / D t ) 的平均值不等于
D ( u j ) / \textasciitilde{} t * 均值的均值实质性导数
现在计算动量方程的平均值就很简单了

图 4.1。控制体 V 的草图，带有边界控制面 A，显示
向外指向单位法线n。

平均动量或雷诺方程

\begin{sourceblock}
\includegraphics[page=113,pagebox=mediabox,trim={64.39bp 1147.86bp 64.39bp 756.90bp},clip,width=0.94\textwidth,max height=0.38\textheight]{Pope.pdf}
\end{sourceblock}
从表面上看，雷诺方程（方程（4.12））和纳维-斯托克斯方程（方程（2.35））是相同的，除了雷诺方程中的术语有一个关键的区别。

与 p(x, t) 一样，平均压力场 (p(x, t ) ) 满足泊松方程。这可以通过取 V'P 的平均值（方程（2.42））或通过取雷诺方程的散度来获得：

\begin{sourceblock}
\includegraphics[page=113,pagebox=mediabox,trim={64.39bp 664.92bp 64.39bp 1135.86bp},clip,width=0.94\textwidth,max height=0.38\textheight]{Pope.pdf}
\end{sourceblock}
从雷诺方程（方程（4.12））获得方程
固定控制体积 V 中平均动量的变化率（参见

图 4.1）。在可能的情况下将项表达为边界上的积分
控制面

对于随机场 +(x,t )，获得结果

\begin{sourceblock}
\includegraphics[page=113,pagebox=mediabox,trim={64.39bp 119.39bp 64.39bp 1796.77bp},clip,width=0.94\textwidth,max height=0.38\textheight]{Pope.pdf}
\end{sourceblock}
和平均旋转率定义为

平均应变率

\begin{sourceblock}
\includegraphics[page=114,pagebox=mediabox,trim={64.39bp 1518.81bp 64.39bp 375.95bp},clip,width=0.94\textwidth,max height=0.38\textheight]{Pope.pdf}
\end{sourceblock}
获取结果

\begin{sourceblock}
\includegraphics[page=114,pagebox=mediabox,trim={64.39bp 1394.82bp 64.39bp 544.93bp},clip,width=0.94\textwidth,max height=0.38\textheight]{Pope.pdf}
\end{sourceblock}
\begin{sourceblock}
\includegraphics[page=114,pagebox=mediabox,trim={64.39bp 1119.86bp 64.39bp 827.90bp},clip,width=0.94\textwidth,max height=0.38\textheight]{Pope.pdf}
\end{sourceblock}
\section*{4.2 雷诺应力}\phantomsection\addcontentsline{toc}{section}{4.2 雷诺应力}
显然，雷诺应力 (uiuj) 在平均速度场 (U) 方程中起着至关重要的作用。如果 (uiuj) 为零，则 U(x, t) 和 (U(x, t)) 的方程将相同。因此，U(x, t) 和 (U(x, t)) 的非常不同的行为（参见第 6 页的图 1.4）可归因于雷诺应力的影响。现在描述它们的一些特性。

解释为应力 雷诺方程可以重写

\begin{sourceblock}
\includegraphics[page=114,pagebox=mediabox,trim={64.39bp 464.94bp 64.39bp 1428.82bp},clip,width=0.94\textwidth,max height=0.38\textheight]{Pope.pdf}
\end{sourceblock}
这是动量守恒方程的一般形式（参见方程 (2.3 I)），方括号中的项代表三个应力之和：粘性应力、平均压力场中的各向同性应力 -(p)Sij 以及脉动速度场产生的表观应力 -p (uiuj)。尽管该表观应力为 -p(u+um)，但将 (uiuj) 称为雷诺应力是方便且传统的。

粘性应力（即每单位面积的力）最终源于力矩
通过脉动速度场进行动量传递。参见图 4，由于流经边界表面 A，固定控制体积 V 内的动量增益率为

\begin{sourceblock}
\includegraphics[page=115,pagebox=mediabox,trim={64.39bp 1605.80bp 64.39bp 295.96bp},clip,width=0.94\textwidth,max height=0.38\textheight]{Pope.pdf}
\end{sourceblock}
（每单位体积的动量为 pU，通过 A 进入 V 的每单位面积的体积流量为 -U n。）该值的 j 分量的平均值

\begin{sourceblock}
\includegraphics[page=115,pagebox=mediabox,trim={64.39bp 1270.84bp 64.39bp 482.94bp},clip,width=0.94\textwidth,max height=0.38\textheight]{Pope.pdf}
\end{sourceblock}
发散定理的最后一步。因此，对于控制体积 V，雷诺方程中出现的雷诺应力（即 -p a(uiuj)/axi）是由边界 A (uiuj)ni 上的脉动速度引起的平均动量通量产生的。

闭合问题 对于一般的统计三维流，有四个独立的方程控制平均速度场；即雷诺方程（方程（4.12））的三个分量以及平均连续性方程（方程（4.4））或（p）的泊松方程（方程（4.13））。但是，这四个方程包含四个以上的未知数。除了 (U) 和 (p)（四个量）之外，还有雷诺应力。

这是闭合问题的表现。一般来说，一组统计量的演化方程（从纳维-斯托克斯方程获得）包含所考虑的集合中的附加统计量。因此，在缺乏单独信息来确定附加统计数据的情况下，无法求解该组方程。这样一组未知数比方程组多的方程组被称为未封闭的。雷诺方程是未封闭的：除非以某种方式确定雷诺应力，否则无法求解它们。

雷诺应力是二阶张量的分量，它显然是对称的，即 (uiuj) (ujui)+ 对角分量 ((u:) (ulu 1) 、 ( u ) 和 (u:)) 是法向应力，而非对角分量（例如 (ulu Z)）是剪应力。

湍动能 k(x, t) 定义为雷诺应力张量的迹的一半：

\begin{sourceblock}
\includegraphics[page=116,pagebox=mediabox,trim={64.39bp 1678.79bp 64.39bp 260.97bp},clip,width=0.94\textwidth,max height=0.38\textheight]{Pope.pdf}
\end{sourceblock}
它是脉动速度场中每单位质量的平均动能。

在雷诺应力张量的主轴上，剪应力为零，法向应力为特征值，其为非负 i。 , ( u f ) 0)。因此雷诺应力张量是对称正半定的。一般来说，所有特征值都严格为正；但是，在特殊或极端情况下，一个或多个特征值可以为零。

各向异性 剪应力和法向应力之间的区别取决于坐标系的选择。各向同性应力和各向异性应力之间存在本质区别。各向同性应力为 1k6-，则偏各向异性部分为
下面广泛使用的归一化各向异性张量定义为

\begin{sourceblock}
\includegraphics[page=116,pagebox=mediabox,trim={64.39bp 878.89bp 64.39bp 1014.87bp},clip,width=0.94\textwidth,max height=0.38\textheight]{Pope.pdf}
\end{sourceblock}
就这些各向异性张量而言，雷诺应力张量为
只有各向异性分量 aij 才能有效地传输动量，因为我们有

\begin{sourceblock}
\includegraphics[page=116,pagebox=mediabox,trim={64.39bp 433.94bp 64.39bp 1478.81bp},clip,width=0.94\textwidth,max height=0.38\textheight]{Pope.pdf}
\end{sourceblock}
表明各向同性分量 ( I k ) 可以在修改后的平均压力中被吸收。

无旋运动 湍流的一个基本特征是它们是旋转的。考虑一个无旋随机速度场，例如（在大约
平均涡度、脉动涡度和 dui/dxj duj/dxi 也为零。因此我们有

\begin{sourceblock}
\includegraphics[page=117,pagebox=mediabox,trim={64.39bp 1649.79bp 64.39bp 247.97bp},clip,width=0.94\textwidth,max height=0.38\textheight]{Pope.pdf}
\end{sourceblock}
由此得出 Corrsin-Kistler 方程（Corrsin 和 Kistler 1954）

\begin{sourceblock}
\includegraphics[page=117,pagebox=mediabox,trim={64.39bp 1495.81bp 64.39bp 424.95bp},clip,width=0.94\textwidth,max height=0.38\textheight]{Pope.pdf}
\end{sourceblock}
对于无旋流。在这种情况下，雷诺应力 (uiuj) 与各向同性应力具有相同的效果，可以在修改的压力中吸收。换句话说，无旋场 t) 产生的雷诺应力对平均速度场绝对没有影响。

对称性 对于某些流动，流动几何形状的对称性决定雷诺应力的属性。

考虑统计二维流，其中统计数据独立于 x3，并且在 x3 坐标轴的反射下统计不变。对于速度 f ( V ; x, t) 的 PDF，这两个条件

\begin{sourceblock}
\includegraphics[page=117,pagebox=mediabox,trim={64.39bp 863.89bp 64.39bp 979.88bp},clip,width=0.94\textwidth,max height=0.38\textheight]{Pope.pdf}
\end{sourceblock}
\begin{sourceblock}
\includegraphics[page=117,pagebox=mediabox,trim={64.39bp 753.90bp 64.39bp 1181.85bp},clip,width=0.94\textwidth,max height=0.38\textheight]{Pope.pdf}
\end{sourceblock}
在 x3 0 处，最后一个方程得出 (U3) -(Us)，即 (Us) 0；它类似地产生 0 和 (u2u3) 0。第一个方程 (Eq. (4.31)) 表明
( \textasciitilde{} 1 \textasciitilde{} =3 ) = 这些关系对于所有 x 都成立。因此，对于这种统计二维流，(U3) 为零，雷诺应力张量为

\begin{sourceblock}
\includegraphics[page=117,pagebox=mediabox,trim={64.39bp 414.95bp 64.39bp 1440.82bp},clip,width=0.94\textwidth,max height=0.38\textheight]{Pope.pdf}
\end{sourceblock}
除了统计上的二维之外，第 76 页图 3.25 中描绘的湍流槽道流在统计上关于平面 0 对称。这种对称性意味着

\begin{sourceblock}
\includegraphics[page=117,pagebox=mediabox,trim={64.39bp 204.97bp 64.39bp 1741.78bp},clip,width=0.94\textwidth,max height=0.38\textheight]{Pope.pdf}
\end{sourceblock}
由此可见 (U2) 和 (uIu2) 是 x2 的奇函数，而
以下每个方程都是不正确的。为什么？

在均匀湍流剪切流的实验中（其中
(Ul)/axzi 是唯一非零平均速度梯度）雷诺
应力（由 k 归一化）测量为

\begin{sourceblock}
\includegraphics[page=118,pagebox=mediabox,trim={64.39bp 665.92bp 64.39bp 1190.85bp},clip,width=0.94\textwidth,max height=0.38\textheight]{Pope.pdf}
\end{sourceblock}
(a) 确定相应的各向异性张量和 bij。

(b) 和 uz 之间的相关系数是多少？

(c) 方程形式的矩阵。 (4.35) 可以转化为
主轴由以下形式的酉矩阵表示

也就是说，对于角度 8 的特定值，Aki(ukuj)Aje 为
对角线。确定角度 8 = OR (0 I O R I n /2)

\section*{4.3 平均标量方程}\phantomsection\addcontentsline{toc}{section}{4.3 平均标量方程}
(d) 确定角度 0 0s (0
0s In /2) 转换
主轴的平均应变率。

(e) 确定 (uiuj) 的特征值。

4.3 平均标量方程 正如湍流速度场 U(X, t) 的最基本描述是由平均速度 (U(x, t)) 提供的一样，守恒无源标量场 +(x, t) 的最基本描述也是由其均值 (+(x,t )) 提供的。 (\$(x, t)) 的守恒方程通过与获取雷诺方程相同的过程获得。

涨落标量场定义为

\begin{sourceblock}
\includegraphics[page=119,pagebox=mediabox,trim={64.39bp 1198.85bp 64.39bp 744.91bp},clip,width=0.94\textwidth,max height=0.38\textheight]{Pope.pdf}
\end{sourceblock}
使得标量场的雷诺分解为

\begin{sourceblock}
\includegraphics[page=119,pagebox=mediabox,trim={64.39bp 1068.87bp 64.39bp 877.89bp},clip,width=0.94\textwidth,max height=0.38\textheight]{Pope.pdf}
\end{sourceblock}
+(x, t) 的守恒方程（方程（2.54））可以写成

\begin{sourceblock}
\includegraphics[page=119,pagebox=mediabox,trim={64.39bp 881.89bp 64.39bp 1060.87bp},clip,width=0.94\textwidth,max height=0.38\textheight]{Pope.pdf}
\end{sourceblock}
唯一的非线性项是涉及对流通量 U\$ 的项，其平均值为

\begin{sourceblock}
\includegraphics[page=119,pagebox=mediabox,trim={64.39bp 645.92bp 64.39bp 1245.84bp},clip,width=0.94\textwidth,max height=0.38\textheight]{Pope.pdf}
\end{sourceblock}
速度标量协方差 (u\$') 是一个向量，称为
标量通量：它表示由于脉动速度场而产生的标量通量（单位面积的流量）（参见练习4.6）。因此，取方程的平均值。 （4.39，我们得到

\begin{sourceblock}
\includegraphics[page=119,pagebox=mediabox,trim={64.39bp 353.95bp 64.39bp 1573.80bp},clip,width=0.94\textwidth,max height=0.38\textheight]{Pope.pdf}
\end{sourceblock}
或者，就平均实质性导数而言（方程（4.10）），

\begin{sourceblock}
\includegraphics[page=119,pagebox=mediabox,trim={64.39bp 191.98bp 64.39bp 1706.78bp},clip,width=0.94\textwidth,max height=0.38\textheight]{Pope.pdf}
\end{sourceblock}
显然，在这个平均标量方程中，标量通量起着类似的作用
它们引起了闭合问题：即使 (U) 已知，方程：如果没有 -(u\$') 的规定，则 (4.41) 无法求解 ( )。

令 @(t) 为守恒无源标量场 \$(x,
在固定控制体积 V 上（参见 \textasciitilde{} i g 4. . 1）。得到一个方程
对于 d(@(t))/dt.W 这里可以将每一项表示为积分
边界控制面 A，并描述其意义
考虑一个统计二维流，其中的统计量
速度场和标量场与 x3 无关且不变
在 x3 坐标轴的反射下。写出对称性
U(x, t) 和 \$(x, t) 的单点联合 PDF 满足的条件。

证明 a(\$)/ax3 和 为零。进一步如果 x2 0 是平面
统计对称性，表明 是 xz 的偶函数，
(u2cpr)是一个函数。

\section*{4.4 梯度扩散与湍流黏度假设}\phantomsection\addcontentsline{toc}{section}{4.4 梯度扩散与湍流黏度假设}
在“科学探究领域”的历史发展中，通常会提出一系列模型来描述正在研究的现象。通常，例如在湍流研究中，早期模型很简单，但随后发现缺乏物理内容和预测准确性。后来的模型可能在物理内容和预测准确性方面更出色，但缺乏简单性。尽管存在缺陷，但欣赏早期的简单模型还是很有价值的。原因之一是模型隐含的行为可以通过简单的推理或简单的分析来确定，而不是更复杂的模型通常需要的数值解。其次，简单的模型可以提供参考，可以比较正在研究的现象以及更复杂的模型。

正是本着这种精神，我们引入了梯度扩散假说、湍流湍流理论和相关思想。这些都是有价值的概念，但应始终牢记其局限性。

标量通量 (u\$') 向量给出了守恒标量 的湍流传输的方向和大小。根据梯度扩散假设，这种传输沿着平均标量梯度下降，即
正标量 r T ( x ,t ) 湍流扩散率，使得

\begin{sourceblock}
\includegraphics[page=121,pagebox=mediabox,trim={64.39bp 1723.78bp 64.39bp 219.97bp},clip,width=0.94\textwidth,max height=0.38\textheight]{Pope.pdf}
\end{sourceblock}
有效扩散系数定义为分子扩散系数和湍流扩散系数之和

\begin{sourceblock}
\includegraphics[page=121,pagebox=mediabox,trim={64.39bp 1563.80bp 64.39bp 330.96bp},clip,width=0.94\textwidth,max height=0.38\textheight]{Pope.pdf}
\end{sourceblock}
结合梯度扩散假设（方程（4.42））的平均标量守恒方程（方程（4.41））是

\begin{sourceblock}
\includegraphics[page=121,pagebox=mediabox,trim={64.39bp 1356.83bp 64.39bp 547.93bp},clip,width=0.94\textwidth,max height=0.38\textheight]{Pope.pdf}
\end{sourceblock}
那么可以看出，这个方程与守恒方程相同
( 4 ) 4， \$ 的方程 (Eq. (2.54)) 但用 (U) 代替 U，并且
在数学上，梯度扩散假设（方程（4.42））类似于傅立叶热传导定律和菲克分子扩散定律。类似地，Boussinesq 在 1877 年提出的湍流黏度 h v 假设在数学上类似于牛顿流体的应力-应变率关系（方程（2.32））。根据假设，偏雷诺应力 (-p(ujuj) Pkdij) 与平均应变率成正比，

\begin{sourceblock}
\includegraphics[page=121,pagebox=mediabox,trim={64.39bp 789.90bp 64.39bp 1054.87bp},clip,width=0.94\textwidth,max height=0.38\textheight]{Pope.pdf}
\end{sourceblock}
其中正标量系数是湍流黏度（也称为
VT 涡黏度）。

结合湍流黏度假设的平均动量方程（将 e 方程（4.45）代入方程（4.12））为

\begin{sourceblock}
\includegraphics[page=121,pagebox=mediabox,trim={64.39bp 370.95bp 64.39bp 1515.81bp},clip,width=0.94\textwidth,max height=0.38\textheight]{Pope.pdf}
\end{sourceblock}
是有效粘度。这与纳维-斯托克斯方程相同，用 (U) 和 v, 代替 U 和 v，并用 (p) i p k 修改平均压力。

迄今为止，梯度扩散和湍流黏度假说的引入还没有任何合理性或批评。对假设的彻底评估被推迟到第 10 章。在此阶段，以下内容
(i) 梯度扩散假设意味着标量通量向量为
与平均标量梯度向量对齐。即使在简单的动荡中
流动发现事实并非如此。例如，在一项实验中
均匀湍流剪切流（Tavoularis 和 Corrsin 1981）

和 ( u ) w 之间的角度经测量为 65"。

(ii) 类似地，湍流黏度假说意味着各向异性
张量与平均应变率张量对齐，即
由于是对称且有偏差的，并且 和 的平均速率
应变有五个独立的分量。根据动荡的情况-
粘度假设，这五个组成部分彼此相关
通过标量系数 vT。同样，即使在简单的剪切流中，
发现这种对齐不会发生（参见练习 4.5）。

(iii) 一类重要的流量由可描述的流量组成
通过二维湍流边界层方程（其中
见第 5 章）。在这些流动中，平均速度为
主要在 xl 坐标方向，而平均值的变化
数量主要集中在 x2 坐标方向。只有一个
标量通量的分量和一个雷诺应力 (ulu2)，
出现在边界层方程中。因此，梯度
扩散假说简化为

\begin{sourceblock}
\includegraphics[page=122,pagebox=mediabox,trim={64.39bp 647.92bp 64.39bp 1295.84bp},clip,width=0.94\textwidth,max height=0.38\textheight]{Pope.pdf}
\end{sourceblock}
和湍流黏度假说

\begin{sourceblock}
\includegraphics[page=122,pagebox=mediabox,trim={64.39bp 464.94bp 64.39bp 1458.82bp},clip,width=0.94\textwidth,max height=0.38\textheight]{Pope.pdf}
\end{sourceblock}
这两个方程都将单个协方差与单个梯度相关联。

假设协方差和梯度具有相反的符号

情况几乎总是如此，而不是假设
或假设，这些方程可以作为 TT 和
(iv) vT(x,t ) 和 rT(x,t ) 的规范解决了闭合问题。也就是说，
如果 和 TT 可以以某种方式指定，则平均流量方程
在高雷诺数且远离墙壁时，vy 和 Ty 是
发现与速度尺度 U/ 和长度尺度 £ 成比例
流量，与流体 v 和 I 的分子特性无关。

因此，比率和 I';/I 均随
雷诺数，因此（在给定情况下）分子
运输可以忽略不计。

湍流普朗特数 oy 定义为
在大多数简单的湍流中或被发现是有序的统一。

表明，根据梯度扩散假设，在一个静态
统计平稳流 (¢) 的最大值和最小值
发生在边界处。

（提示：修改 2.6 节中使用的有界参数以应用
其为等式。 （4.44）。）

证明，为了实现湍流黏度假设（方程（4.45））

为了产生非负正应力，有必要且充分
为使湍流黏度满足

\begin{sourceblock}
\includegraphics[page=123,pagebox=mediabox,trim={64.39bp 862.89bp 64.39bp 1038.87bp},clip,width=0.94\textwidth,max height=0.38\textheight]{Pope.pdf}
\end{sourceblock}
其中 Sy;是平均应变率张量的最大特征值。

\chapter{自由剪切流}
最常研究的湍流自由剪切流是射流、尾流和混合层。正如“自由”这个名字所暗示的那样，这些流动远离墙壁，并且由于平均速度差异而产生湍流。

我们首先检查圆射流。通过将实验观察（第 5.1 节）与雷诺方程（第 5.2 节）相结合，不仅可以了解圆射流，还可以了解一般湍流的行为。在5.3节中，我们研究了圆射流中的湍动能，以及能量产生和耗散的重要过程。其他自相似自由剪切流在 5.4 节中进行了简要描述；第 5.5 节对自由剪切流行为进行了进一步观察。

\section*{5.1 圆射流：实验观测}\phantomsection\addcontentsline{toc}{section}{5.1 圆射流：实验观测}
\subsection*{5.1.1 流程描述}\phantomsection\addcontentsline{toc}{subsection}{5.1.1 流程描述}
我们已经在第一章中遇到过圆射流，例如图 1 和 2.1.1-1.4。理想的实验配置和采用的坐标系如图 5.1 所示。牛顿流体稳定地流过直径为 d 的喷嘴，产生（近似）平顶速度分布，速度为 U；。来自喷嘴的射流流入同一流体的环境中，该流体在无穷远处静止。流动在统计上是平稳且轴对称的。因此，统计数据取决于轴向和径向坐标（x 和 r），但与时间和圆周坐标 0 无关。x、r 和 @ 坐标方向上的速度分量由 U、V 和 W 表示。

在理想实验中，流动完全由 U;,d 和 vy 定义，因此唯一的无量纲参数是雷诺数，定义

\begin{sourceblock}
\includegraphics[page=125,pagebox=mediabox,trim={64.39bp 1290.84bp 64.39bp 109.44bp},clip,width=\textwidth,max height=0.72\textheight]{Pope.pdf}
\end{sourceblock}
\subsection*{5.1.2 平均速度场}\phantomsection\addcontentsline{toc}{subsection}{5.1.2 平均速度场}
\begin{center}\small 图 5.1。圆形喷射实验的草图，显示了所采用的极柱坐标系。\end{center}
作者：Re UJd/v。 （实际上，喷嘴和周围环境的细节会产生一些影响，正如 Schneider (1985) 和 Hussein 等人 (1994) 所讨论的那样。）

平均速度场 正如从流动的视觉外观所预期的那样（第 4 页和第 5 页的图 1.1 和 1.2），平均速度主要在轴向方向。测量的平均轴向速度的径向分布如图 5.2 所示。 （请注意，0 是轴，(U) 的轮廓绕其对称。）图 5.2 中未显示初始开发区域（例如 0 I x l d 25），其中轮廓从（近似）方形变为图 5.2 中所示的圆形。平均圆周速度为零 ( e . , ( W ) 0 )，而如练习 5.5 所示，平均径向速度 ( V ) 比 ( U ) 小一个数量级。

轴向速度 根据平均轴向速度场 (U(X, 8))（与 8 无关），中心线速度为

\begin{sourceblock}
\includegraphics[page=125,pagebox=mediabox,trim={64.39bp 335.96bp 64.39bp 1609.80bp},clip,width=0.94\textwidth,max height=0.38\textheight]{Pope.pdf}
\end{sourceblock}
射流的半宽度定义为

\begin{sourceblock}
\includegraphics[page=125,pagebox=mediabox,trim={64.39bp 199.97bp 64.39bp 1739.78bp},clip,width=0.94\textwidth,max height=0.38\textheight]{Pope.pdf}
\end{sourceblock}
图 5.2 中两个明显的观察结果是，随着轴向距离的增加，

\begin{sourceblock}
\includegraphics[page=126,pagebox=mediabox,trim={64.39bp 1187.85bp 64.39bp 109.44bp},clip,width=\textwidth,max height=0.72\textheight]{Pope.pdf}
\par\vspace{0.45\baselineskip}
\small 图 5.2。湍流圆射流中平均轴向速度的径向分布，Re = 95,500。虚线表示轮廓的半宽。 （改编自 Hussein 等人 (1994) 的数据。）
\end{sourceblock}
\begin{sourceblock}
\includegraphics[page=126,pagebox=mediabox,trim={64.39bp 440.94bp 64.39bp 901.89bp},clip,width=\textwidth,max height=0.72\textheight]{Pope.pdf}
\par\vspace{0.45\baselineskip}
\small 图 5.3：湍流圆射流中平均轴向速度与径向距离的关系，Re = 10°； Wygnanski 和 Fiedler (1969) 的测量。符号：0，x/d=40； A、x/d=50； O，x/d=60； 0，x/d = 75； e、x/d
\end{sourceblock}
随着射流的衰减和扩散，平均速度分布会发生变化，如图 5.2 所示，但分布的 不会发生变化。超越
形状发展区域（xld 30，例如），(U)/Uo(x) 的轮廓，绘制
Wygnanski 和 Fiedler (1969) 以这种方式绘制了 x/d 在 40 到 100 之间的情况。重要的结论是平均速度分布变得自相似。

自相似性 自相似性是在湍流研究中的几个不同背景下出现的一个重要概念。为了探索一般思想，请考虑取决于两个自变量（ e 、 x 和 y）的量 Q(x,y)。 x 的 AS 函数、特征尺度 Qo(x) 和 6(x) 分别为因变量 Q 和自变量 y 定义。然后缩放变量定义为

\begin{sourceblock}
\includegraphics[page=127,pagebox=mediabox,trim={64.39bp 1269.84bp 64.39bp 641.92bp},clip,width=0.94\textwidth,max height=0.38\textheight]{Pope.pdf}
\end{sourceblock}
\begin{sourceblock}
\includegraphics[page=127,pagebox=mediabox,trim={64.39bp 1176.85bp 64.39bp 741.91bp},clip,width=0.94\textwidth,max height=0.38\textheight]{Pope.pdf}
\end{sourceblock}
如果缩放因变量独立于 x，即存在一个函数

\begin{sourceblock}
\includegraphics[page=127,pagebox=mediabox,trim={64.39bp 1000.87bp 64.39bp 885.89bp},clip,width=0.94\textwidth,max height=0.38\textheight]{Pope.pdf}
\end{sourceblock}
那么 Q(x, y) 是自相似的。在这种情况下，Q(x, y) 可以表示为
～（5）．单自变量 Qo(x)、6(x) 和的函数
以下是一些评论和资格：

(i) 必须适当选择尺度 Qo(x) 和 6(x)，它们通常
对 x 具有幂律依赖性；

(ii) 在某些情况下，需要更一般的转换，
(iii) 可以观察到自相似行为（在良好的近似范围内）

化）在 x 的范围内（但不是所有 x ）；和
(iv) 如果自相似量 Q(x,y) 受偏微分控制
tial 方程，则 Qo(x),S (x)，并且由普通方程控制
微分方程。

尺度的轴向变化 回到圆射流，为了完成这幅图，我们需要确定 Uo(x) 和 ( x ) 的轴向变化，图 5.4 显示了 Uo(x) 的逆向，特别是 UJ Uo(x)，相对于 x/d 绘制。显然，在 x/d 范围内

\begin{sourceblock}
\includegraphics[page=128,pagebox=mediabox,trim={64.39bp 1192.85bp 64.39bp 109.44bp},clip,width=\textwidth,max height=0.72\textheight]{Pope.pdf}
\par\vspace{0.45\baselineskip}
\small 图 5.4。湍流圆射流中沿中心线的平均速度随轴向距离的变化，Re = 95,500：符号，Hussein 等人的实验数据，(1994)；和线，方程。 (5.6) 其中 xo/d = 4 且 B = 5.8。
\end{sourceblock}
\begin{sourceblock}
\includegraphics[page=128,pagebox=mediabox,trim={64.39bp 847.89bp 64.39bp 1075.86bp},clip,width=0.94\textwidth,max height=0.38\textheight]{Pope.pdf}
\end{sourceblock}
\begin{sourceblock}
\includegraphics[page=128,pagebox=mediabox,trim={64.39bp 606.92bp 64.39bp 1316.83bp},clip,width=0.94\textwidth,max height=0.38\textheight]{Pope.pdf}
\end{sourceblock}
\begin{sourceblock}
\includegraphics[page=128,pagebox=mediabox,trim={64.39bp 452.94bp 64.39bp 1494.81bp},clip,width=0.94\textwidth,max height=0.38\textheight]{Pope.pdf}
\end{sourceblock}
\begin{sourceblock}
\includegraphics[page=128,pagebox=mediabox,trim={64.39bp 174.98bp 64.39bp 1770.78bp},clip,width=0.94\textwidth,max height=0.38\textheight]{Pope.pdf}
\end{sourceblock}
Panchapakesan 和 Hussein 等人。 (1994)，侯赛因等 a1。 （1994），
热线数据激光多普勒数据

雷诺数 在理想的圆射流实验中，唯一的无量纲参数是射流的雷诺数 Re。因此，我们应该问；自相似轮廓形状、速度衰减常数和铺展率 S 如何随 Re 变化。答案简单而深刻：不依赖Re。表 5.1 显示，对于 Re 相差近十倍的射流，B 和 S 测量值的微小差异在实验不确定性范围内。此外，从目视观察来看，Re 增大千倍的射流的扩散率是相同的（参见第 4 页的图 1.1 以及 Mungal 和 Hollingsworth (1989)）。从第 5 页的图 1.2 可以明显看出，雷诺数确实会影响流动：雷诺数较大时，小规模结构较小。然而，重复一遍，平均速度分布和扩散率与 Re 无关。

小结 在高雷诺数湍流射流 (Re lo4) 的自相似区域 ( x l d > 30) 中，中心线速度 Uo(x) 和半宽 rl12(x)v 由式(1) 确定。 （5.6）和（5.8）。这些定律中的经验常数与 Re 无关：为了确定性，我们将它们的值设为 = 5.8 和 S 0.094（见表 5.1）。跨流相似性变量可以取为

\begin{sourceblock}
\includegraphics[page=129,pagebox=mediabox,trim={64.39bp 425.95bp 64.39bp 1520.81bp},clip,width=0.94\textwidth,max height=0.38\textheight]{Pope.pdf}
\end{sourceblock}
\begin{sourceblock}
\includegraphics[page=129,pagebox=mediabox,trim={64.39bp 298.96bp 64.39bp 1595.80bp},clip,width=0.94\textwidth,max height=0.38\textheight]{Pope.pdf}
\end{sourceblock}
两者通过 q S t 相关。自相似平均速度剖面为

\begin{sourceblock}
\includegraphics[page=129,pagebox=mediabox,trim={64.39bp 147.98bp 64.39bp 1740.78bp},clip,width=0.94\textwidth,max height=0.38\textheight]{Pope.pdf}
\end{sourceblock}
\begin{sourceblock}
\includegraphics[page=130,pagebox=mediabox,trim={64.39bp 1087.86bp 64.39bp 109.44bp},clip,width=\textwidth,max height=0.72\textheight]{Pope.pdf}
\par\vspace{0.45\baselineskip}
\small 图 5.5。自相似圆射流中平均轴向速度的自相似轮廓：与 Hussein 等人 (1994) 的 LDA 数据的曲线拟合。
\end{sourceblock}
图 5.6。自相似圆射流的平均横向速度。来自LDA数据
侯赛因 (1994)。

等人。

横向荷载

在圆射流的自相似区域中，平均横向速度 (V) 可以通过连续性方程从 (U) 确定（参见练习 5.4 和 5.5）。图 5.6 显示了用这种方法得到的 (v)/Uo 的自相似轮廓。应该观察到 (V) 非常小，小于 Uo 的 40 倍。还请注意 (V) 在射流边缘处为负值，表明

\begin{sourceblock}
\includegraphics[page=131,pagebox=mediabox,trim={64.39bp 1571.80bp 64.39bp 205.97bp},clip,width=0.94\textwidth,max height=0.38\textheight]{Pope.pdf}
\end{sourceblock}
\begin{sourceblock}
\includegraphics[page=131,pagebox=mediabox,trim={64.39bp 1482.81bp 64.39bp 444.94bp},clip,width=0.94\textwidth,max height=0.38\textheight]{Pope.pdf}
\end{sourceblock}
从湍流圆射流中的自相似速度分布 f (q)

（方程（5.11）和（5.12））表明

\begin{sourceblock}
\includegraphics[page=131,pagebox=mediabox,trim={64.39bp 1274.84bp 64.39bp 625.92bp},clip,width=0.94\textwidth,max height=0.38\textheight]{Pope.pdf}
\end{sourceblock}
\begin{sourceblock}
\includegraphics[page=131,pagebox=mediabox,trim={64.39bp 1156.85bp 64.39bp 767.90bp},clip,width=0.94\textwidth,max height=0.38\textheight]{Pope.pdf}
\end{sourceblock}
（其中质数表示对 q 的微分）。

自相似速度分布的近似为

\begin{sourceblock}
\includegraphics[page=131,pagebox=mediabox,trim={64.39bp 986.88bp 64.39bp 960.88bp},clip,width=0.94\textwidth,max height=0.38\textheight]{Pope.pdf}
\end{sourceblock}
表明（根据 rl12 的定义）常数 a 由下式给出

\begin{sourceblock}
\includegraphics[page=131,pagebox=mediabox,trim={64.39bp 834.89bp 64.39bp 1088.86bp},clip,width=0.94\textwidth,max height=0.38\textheight]{Pope.pdf}
\end{sourceblock}
证明，根据这个近似，

\begin{sourceblock}
\includegraphics[page=131,pagebox=mediabox,trim={64.39bp 621.92bp 64.39bp 1226.85bp},clip,width=0.94\textwidth,max height=0.38\textheight]{Pope.pdf}
\end{sourceblock}
（请注意， ( U\}/ ar),=,1/2i 大约是 ( U) 的六倍
轴向位置。）

对于湍流圆射流，在极柱坐标中，平均值
连续性方程为

\begin{sourceblock}
\includegraphics[page=131,pagebox=mediabox,trim={64.39bp 237.97bp 64.39bp 1709.78bp},clip,width=0.94\textwidth,max height=0.38\textheight]{Pope.pdf}
\end{sourceblock}
证明，如果 ( U ) 自相似

\begin{sourceblock}
\includegraphics[page=131,pagebox=mediabox,trim={64.39bp 119.39bp 64.39bp 1849.77bp},clip,width=0.94\textwidth,max height=0.38\textheight]{Pope.pdf}
\end{sourceblock}
那么 (V) 也自相似

\begin{sourceblock}
\includegraphics[page=132,pagebox=mediabox,trim={64.39bp 1729.78bp 64.39bp 215.97bp},clip,width=0.94\textwidth,max height=0.38\textheight]{Pope.pdf}
\end{sourceblock}
其中 f ( q ) 和 h(q) 的关系为

\begin{sourceblock}
\includegraphics[page=132,pagebox=mediabox,trim={64.39bp 1604.80bp 64.39bp 340.96bp},clip,width=0.94\textwidth,max height=0.38\textheight]{Pope.pdf}
\end{sourceblock}
如果自相似轴向速度分布 f (7) 由方程给出： （5.17），
由方程式可知(5.24) 横向速度剖面 h(q) 为

\begin{sourceblock}
\includegraphics[page=132,pagebox=mediabox,trim={64.39bp 1429.82bp 64.39bp 506.94bp},clip,width=0.94\textwidth,max height=0.38\textheight]{Pope.pdf}
\end{sourceblock}
表明，根据该方程，横向速度
半宽是

\begin{sourceblock}
\includegraphics[page=132,pagebox=mediabox,trim={64.39bp 1259.84bp 64.39bp 675.91bp},clip,width=0.94\textwidth,max height=0.38\textheight]{Pope.pdf}
\end{sourceblock}
表明，对于较大的 r，方程： (5.25) 意味着

\begin{sourceblock}
\includegraphics[page=132,pagebox=mediabox,trim={64.39bp 1048.87bp 64.39bp 898.89bp},clip,width=0.94\textwidth,max height=0.38\textheight]{Pope.pdf}
\end{sourceblock}
令 f ( c , F, , , 8 ) 表示 U、V 和 W 的联合 PDF
湍流圆射流： 、 、 和 是样本空间变量。

流量在统计上是轴对称的：

\begin{sourceblock}
\includegraphics[page=132,pagebox=mediabox,trim={64.39bp 801.90bp 64.39bp 1105.86bp},clip,width=0.94\textwidth,max height=0.38\textheight]{Pope.pdf}
\end{sourceblock}
在圆周坐标的反射下保持不变

\begin{sourceblock}
\includegraphics[page=132,pagebox=mediabox,trim={64.39bp 633.92bp 64.39bp 1292.84bp},clip,width=0.94\textwidth,max height=0.38\textheight]{Pope.pdf}
\end{sourceblock}
表明，对于 8 0，等式： (5.29) 意味着 (W)、( U W ) 和 ( V W )

为零，并且（根据方程（5.28））这些量为零
画出 r-8 平面的草图并表明（对于给定的 x，
8) 对应于 V(x, r, e)、V (x, r, 8 n) 的方向，
W(x,r,8n/2)和W(x,r,8n/2)相同。因此争论
（对于统计轴对称流动）

\begin{sourceblock}
\includegraphics[page=132,pagebox=mediabox,trim={64.39bp 242.97bp 64.39bp 1682.79bp},clip,width=0.94\textwidth,max height=0.38\textheight]{Pope.pdf}
\end{sourceblock}
证明在轴上 和 ( W 2 ) 相等，并且 ( V ) 和
(UV) 为零。

\subsection*{5.1.3 雷诺应力}\phantomsection\addcontentsline{toc}{subsection}{5.1.3 雷诺应力}
x、r 和 8 坐标方向上的脉动速度分量用 u、v 表示，并且 在湍流圆射流中，雷诺应力

\begin{sourceblock}
\includegraphics[page=133,pagebox=mediabox,trim={64.39bp 1524.81bp 64.39bp 294.96bp},clip,width=0.94\textwidth,max height=0.38\textheight]{Pope.pdf}
\end{sourceblock}
也就是说，由于圆周对称，( u w ) 和 (vw) 为零（参见练习 5.6）。流动的几何形状还决定了法向应力是 r 的偶函数，而剪切应力 (uv) 是奇函数。当接近轴线r 0 时，速度的径向V分量和周向W分量变得难以区分。因此 (v2) 和 (w2) 在轴上相等。

考虑中心线上的 ramas.a 轴向速度

\begin{sourceblock}
\includegraphics[page=133,pagebox=mediabox,trim={64.39bp 1137.86bp 64.39bp 810.90bp},clip,width=0.94\textwidth,max height=0.38\textheight]{Pope.pdf}
\end{sourceblock}
u\&(x)v 如何与 x 相加？或者，就无量纲量而言，ub(x)/Uo(x)如何与x/d和Re变化？答案同样很简单，但却非常有启发性。在发展区域之后，ub(x)/Uo(x) 渐近趋于约 0.25 的恒定值（参见 Panchapakesan 和 Lumley (1993a)）。因此，像 Uo(x) 一样，ub(x) 也会随着 x-I 衰减。不同实验中 ub/Uo 存在一些差异，但尚未记录到对 Re 的系统依赖性。

正如从上述观察中可以预期的那样，我们发现雷诺应力变得自相似。也就是说，对于超出开发区域的所有 x，针对或 q r / ( x xo) 绘制的 ( u \textasciitilde{} \textasciitilde{} \textasciitilde{} \textasciitilde{} ) / u \textasciitilde{} ( x ) \textasciitilde{} 的剖面图会崩溃。图 5.7 显示了 Hussein 等人测量的自相似剖面。 （1994）。这些数据的一些重要观察结果是
(i) 在中心线上，均方根值速度约为平均值的 25\%。

(ii) 朝向射流边缘，尽管雷诺应力衰减（随着
增加 / r 2 )，拉玛斯与局部平均值的比率增加
无限制（见图 5.8）。

(iii) 雷诺应力表现出显着的各向异性，这表明
剪切应力和法向应力的差异。

(iv) 剪切应力的相对大小可以通过以下方式量化
比率 (uv)/k，以及 u-v 相关系数 p,,（见图 5.9）。

两条曲线形状相同，中心部分平坦，

\begin{sourceblock}
\includegraphics[page=134,pagebox=mediabox,trim={64.39bp 1217.85bp 64.39bp 109.44bp},clip,width=\textwidth,max height=0.72\textheight]{Pope.pdf}
\par\vspace{0.45\baselineskip}
\small 图 5.7。自相似圆射流中的雷诺应力分布：与 Hussein 等人的 LDA 数据的曲线拟合。 （1994）。
\end{sourceblock}
\begin{sourceblock}
\includegraphics[page=134,pagebox=mediabox,trim={64.39bp 461.94bp 64.39bp 838.90bp},clip,width=\textwidth,max height=0.72\textheight]{Pope.pdf}
\par\vspace{0.45\baselineskip}
\small 图 5.8。自相似圆射流中局部湍流强度的轮廓 - (u*)!/?/(U)。从曲线拟合到 Hussein 等人的实验数据。 （1994）。
\end{sourceblock}
(v) 剪应力为正，其中 d ( U ) / d r 为负，并且变为
零，其中 d ( U ) / d r 趋于零。因此，对于这个流程，有一个
正湍流黏度使得

\begin{sourceblock}
\includegraphics[page=134,pagebox=mediabox,trim={64.39bp 119.39bp 64.39bp 1819.77bp},clip,width=0.94\textwidth,max height=0.38\textheight]{Pope.pdf}
\end{sourceblock}
图 5.9。自相似中 (uu)/k 和 u-E 相关系数 p ,, 的轮廓
圆形喷射。从曲线拟合到 Hussein 等人的实验数据。 （1994）。

(vi) 由于 (uv) 和 d(U)/dr 的轮廓是自相似的，显然
由方程定义的湍流黏度曲线。 (5.33) 也是自相似的。

\begin{sourceblock}
\includegraphics[page=135,pagebox=mediabox,trim={64.39bp 798.90bp 64.39bp 1145.86bp},clip,width=0.94\textwidth,max height=0.38\textheight]{Pope.pdf}
\end{sourceblock}
其中 CiT 是归一化轮廓，如图 5.10 所示。它
可以观察到 CT 在整个射流上相当均匀
对于 0.1 r 1.5，在 0.028 的 15\% 以内，但它会减小到
零向边缘。

(vii) 湍流黏度的量纲为速度乘以长度。康-
随后，局部长度尺度 ( x , ) ,c 可以由下式定义

\begin{sourceblock}
\includegraphics[page=135,pagebox=mediabox,trim={64.39bp 421.95bp 64.39bp 1523.81bp},clip,width=0.94\textwidth,max height=0.38\textheight]{Pope.pdf}
\end{sourceblock}
其中 u'(x, r ) 是本地均方根。轴向速度 显然 I 是自
类似。 I/rl12的轮廓（图5.1 1）相当平坦，在15\%以内
大多数射流的值为 0.12 (0.1 r 2.1)。

长度尺度 I，在方程式中定义。 （5.35）是一个导出量，而不是可以直接测量并具有明确的物理意义。

\begin{sourceblock}
\includegraphics[page=136,pagebox=mediabox,trim={64.39bp 1190.85bp 64.39bp 109.44bp},clip,width=\textwidth,max height=0.72\textheight]{Pope.pdf}
\par\vspace{0.45\baselineskip}
\small 图 5.10。归一化湍流扩散系数\$； （方程（5.34））在自相似圆射流中。从曲线拟合到 Hussein 等人的实验数据。 （1994）。
\end{sourceblock}
图 5.12。湍流圆射流中整体长度尺度的自相似轮廓。

来自 Wygnanski 和 Fiedler (1969)。

\begin{sourceblock}
\includegraphics[page=137,pagebox=mediabox,trim={64.39bp 947.88bp 64.39bp 888.89bp},clip,width=0.94\textwidth,max height=0.38\textheight]{Pope.pdf}
\end{sourceblock}
然后相应的积分长度尺度是

\begin{sourceblock}
\includegraphics[page=137,pagebox=mediabox,trim={64.39bp 658.92bp 64.39bp 1243.84bp},clip,width=0.94\textwidth,max height=0.38\textheight]{Pope.pdf}
\end{sourceblock}
\begin{sourceblock}
\includegraphics[page=137,pagebox=mediabox,trim={64.39bp 543.93bp 64.39bp 1358.83bp},clip,width=0.94\textwidth,max height=0.38\textheight]{Pope.pdf}
\end{sourceblock}
图 5.12 显示了这些积分长度尺度的测量自相似轮廓。可以看出，L l1 和Lzz 通常分别为0.7r1i2 和0.3r1i2，分别显着大于i 0.lrll2.It 还应该理解，对于大于Lll 的间隔距离可能存在显着的相关性，如图5.13所示。自相似圆射流的一些进一步特征在 5.4 节中描述。

在 r 1 处，平均应变率张量的主轴为

\begin{sourceblock}
\includegraphics[page=138,pagebox=mediabox,trim={64.39bp 1358.83bp 64.39bp 109.44bp},clip,width=\textwidth,max height=0.72\textheight]{Pope.pdf}
\par\vspace{0.45\baselineskip}
\small 图 5.13。自相似圆射流中轴向速度的纵向自相关。来自 Wygnanski 和 Fiedler (1969)。
\end{sourceblock}
图 5.7) 雷诺应力张量的主轴位于
与 x 轴的角度小于 30"。

将 Lll 与图 2 和 3 中喷射流的可见宽度进行比较。 1.1 和 1.2。

在极柱坐标（x、r 和 B）中，连续性方程
（其中 U、V 和 W 是三个坐标方向上的速度
蒸发散）； Navier-S 托克斯方程为

\begin{sourceblock}
\includegraphics[page=138,pagebox=mediabox,trim={64.39bp 747.91bp 64.39bp 1149.85bp},clip,width=0.94\textwidth,max height=0.38\textheight]{Pope.pdf}
\end{sourceblock}
\begin{sourceblock}
\includegraphics[page=138,pagebox=mediabox,trim={64.39bp 524.93bp 64.39bp 1372.83bp},clip,width=0.94\textwidth,max height=0.38\textheight]{Pope.pdf}
\end{sourceblock}
\begin{sourceblock}
\includegraphics[page=138,pagebox=mediabox,trim={64.39bp 327.96bp 64.39bp 1521.81bp},clip,width=0.94\textwidth,max height=0.38\textheight]{Pope.pdf}
\end{sourceblock}
\begin{sourceblock}
\includegraphics[page=138,pagebox=mediabox,trim={64.39bp 166.98bp 64.39bp 1695.79bp},clip,width=0.94\textwidth,max height=0.38\textheight]{Pope.pdf}
\end{sourceblock}
\section*{5.2 圆射流：平均动量}\phantomsection\addcontentsline{toc}{section}{5.2 圆射流：平均动量}
(W)、(uw) 和 (vw) 为零。表明雷诺方程
这些流量是：

\begin{sourceblock}
\includegraphics[page=139,pagebox=mediabox,trim={64.39bp 1669.79bp 64.39bp 237.97bp},clip,width=0.94\textwidth,max height=0.38\textheight]{Pope.pdf}
\end{sourceblock}
\begin{sourceblock}
\includegraphics[page=139,pagebox=mediabox,trim={64.39bp 1548.80bp 64.39bp 342.96bp},clip,width=0.94\textwidth,max height=0.38\textheight]{Pope.pdf}
\end{sourceblock}
P DR AX R DR R

\begin{sourceblock}
\includegraphics[page=139,pagebox=mediabox,trim={64.39bp 1198.85bp 64.39bp 560.93bp},clip,width=0.94\textwidth,max height=0.38\textheight]{Pope.pdf}
\end{sourceblock}
3.2 圆射流：平均动量

\subsection*{5.2.1 边界层方程}\phantomsection\addcontentsline{toc}{subsection}{5.2.1 边界层方程}
在湍流圆射流中，存在占主导地位的平均流动方向 (x)，平均横向速度相对较小 (l(V)l 0.03I(U)I)，流动逐渐扩散 (drI12/dx 0 . 1 )，并且 s o (表示平均值) 轴向梯度与横向梯度相比较小。所有自由剪切流所共有的这些特征允许使用边界层方程代替完整的雷诺方程当然，湍流边界层方程也适用于湍流边界层，也适用于其他一些壁面流动，这些流动将在第 7 章中讨论。

我们首先考虑统计二维静止流，其中 x 是流的主导方向，均值梯度主要在 y 方向，并且统计数据在该方向上不变化。速度分量为 U、V 和 W，其中 ( W ) 为零。图 5.14 概述了此类流程的示例。我们考虑如图 5.14 所示的情况
其中上边界 (y co) 是静止的或非湍流的自由流。对于每个流，作为 x 的函数，我们可以定义 6(x) 为特征流宽度，U,(x) 为特征对流速度，U,(x) 为特征速度差。

对于这些流动，平均连续性和动量方程为

\begin{sourceblock}
\includegraphics[page=139,pagebox=mediabox,trim={64.39bp 119.39bp 64.39bp 1808.77bp},clip,width=0.94\textwidth,max height=0.38\textheight]{Pope.pdf}
\end{sourceblock}
喷气式飞机

\begin{sourceblock}
\includegraphics[page=140,pagebox=mediabox,trim={64.39bp 194.98bp 64.39bp 643.72bp},clip,width=\textwidth,max height=0.72\textheight]{Pope.pdf}
\par\vspace{0.45\baselineskip}
\small 图 5.14。平面二维剪切流草图，显示特征流宽度 5(x)、特征对流速度 U. 和特征速度差 U.。
\end{sourceblock}
这些方程也适用于层流，在这种情况下雷诺应力为零。在边界层近似中忽略大括号 (\{ \}) 中的项。

湍流边界层方程只需忽略括号中的项即可获得，其原因与在层流情况下忽略它们的原因相同，并且忽略雷诺应力的轴向导数，因为它们与横向梯度相比较小。横向动量方程（方程（5.51））则变为

\begin{sourceblock}
\includegraphics[page=141,pagebox=mediabox,trim={64.39bp 962.88bp 64.39bp 933.88bp},clip,width=0.94\textwidth,max height=0.38\textheight]{Pope.pdf}
\end{sourceblock}
在自由流 ( v m) 中，压力用 po(x) 表示，并且 (u2) 为零。因此，可以对方程进行积分，得到

\begin{sourceblock}
\includegraphics[page=141,pagebox=mediabox,trim={64.39bp 791.90bp 64.39bp 1149.85bp},clip,width=0.94\textwidth,max height=0.38\textheight]{Pope.pdf}
\end{sourceblock}
则轴向压力梯度为

\begin{sourceblock}
\includegraphics[page=141,pagebox=mediabox,trim={64.39bp 618.92bp 64.39bp 1274.84bp},clip,width=0.94\textwidth,max height=0.38\textheight]{Pope.pdf}
\end{sourceblock}
对于静态或均匀自由流的流动，压力梯度 dpo/dx 为零。一般来说，dpo/dx 由伯努利方程根据自由流速度给出。

在轴向动量方程中，Eq. （5.50），忽略大括号中的项并代入等式。 (5.54)，我们得到

\begin{sourceblock}
\includegraphics[page=141,pagebox=mediabox,trim={64.39bp 288.96bp 64.39bp 1619.80bp},clip,width=0.94\textwidth,max height=0.38\textheight]{Pope.pdf}
\end{sourceblock}
右侧的第一项和最后一项需要进一步讨论。

在湍流自由剪切流中，v d 2 ( \textasciitilde{} ) / d v 2 的阶数为 v u , / d 2，因此，与方程（1）中的主要项相比， （5.55），它的顺序是 \textasciitilde{} e - ' ,and
在边界层中，速度导数非常大，并且不与 Us 和 6 成比例。在这种情况下，粘性项 v a 2 ( U ) / a y 2 是方程 (1) 中的前导项。 （5.55）。

在层流边界层方程中，轴向扩散项 v
u / a x 2 的相对顺序为 \textasciitilde{} e - '，因此小得可以忽略不计。湍流边界层流中可比较的轴向应力梯度项是方程式中的最后一项。 （5.55）。忽略这个术语是一致的；但应该意识到，这并不是一个微不足道的近似值。正如练习 5.11 所示，在自由剪切流中，被忽略的项可能约为方程中主要项的 10\%。

综上所述：对于以静态流体或均匀流为界的统计二维静态流，湍流边界层方程由连续性方程（方程（5.49））和轴向动量方程组成

\begin{sourceblock}
\includegraphics[page=142,pagebox=mediabox,trim={64.39bp 1117.86bp 64.39bp 778.90bp},clip,width=0.94\textwidth,max height=0.38\textheight]{Pope.pdf}
\end{sourceblock}
除了靠近墙壁之外，粘性项可以忽略不计。平均压力分布由方程式给出。 （5.53）。

对于统计轴对称、静止非旋流（例如圆射流或球体后面的尾流），相应的湍流边界层方程为

\begin{sourceblock}
\includegraphics[page=142,pagebox=mediabox,trim={64.39bp 664.92bp 64.39bp 1116.86bp},clip,width=0.94\textwidth,max height=0.38\textheight]{Pope.pdf}
\end{sourceblock}
平均压力分布为

\begin{sourceblock}
\includegraphics[page=142,pagebox=mediabox,trim={64.39bp 500.94bp 64.39bp 1398.82bp},clip,width=0.94\textwidth,max height=0.38\textheight]{Pope.pdf}
\end{sourceblock}
和轴向应力梯度项，在右侧忽略

\begin{sourceblock}
\includegraphics[page=142,pagebox=mediabox,trim={64.39bp 297.96bp 64.39bp 1539.81bp},clip,width=0.94\textwidth,max height=0.38\textheight]{Pope.pdf}
\end{sourceblock}
从极柱坐标中的雷诺方程开始

（方程（5.45）-（5.47）），验证相应的边界层
对于自相似圆射流的中心线，可得：

对边界层中或省略的项进行估计

\subsection*{5.2.2 质量、动量和能量的流量}\phantomsection\addcontentsline{toc}{subsection}{5.2.2 质量、动量和能量的流量}
我们回到湍流圆射流，进行一些源于动量守恒的基本观察。忽略粘性项并乘以 r，动量方程（方程（5.58））变为

\begin{sourceblock}
\includegraphics[page=143,pagebox=mediabox,trim={64.39bp 992.87bp 64.39bp 951.88bp},clip,width=0.94\textwidth,max height=0.38\textheight]{Pope.pdf}
\end{sourceblock}
（连续性方程 Eq. (5.57) 用于以保守形式写出对流项。）对 r 进行积分，我们得到

\begin{sourceblock}
\includegraphics[page=143,pagebox=mediabox,trim={64.39bp 718.91bp 64.39bp 1140.86bp},clip,width=0.94\textwidth,max height=0.38\textheight]{Pope.pdf}
\end{sourceblock}
因为，对于较大的 r， (U V ) 趋于零的速度比 r 更快。平均流量的动量流量为

\begin{sourceblock}
\includegraphics[page=143,pagebox=mediabox,trim={64.39bp 498.94bp 64.39bp 1401.82bp},clip,width=0.94\textwidth,max height=0.38\textheight]{Pope.pdf}
\end{sourceblock}
我们从等式中可以看出。 (5.63) 动量流量守恒：\textasciitilde{}(x) 与 x 无关。同样的结论适用于所有喷射流（喷射到静止环境或均匀流中）和尾流（在均匀流中）。

自相似圆射流的平均速度分布可以写成

\begin{sourceblock}
\includegraphics[page=143,pagebox=mediabox,trim={64.39bp 226.97bp 64.39bp 1715.78bp},clip,width=0.94\textwidth,max height=0.38\textheight]{Pope.pdf}
\end{sourceblock}
\begin{sourceblock}
\includegraphics[page=143,pagebox=mediabox,trim={64.39bp 119.39bp 64.39bp 1849.77bp},clip,width=0.94\textwidth,max height=0.38\textheight]{Pope.pdf}
\end{sourceblock}
f ( < ) 5 并且 是相似性轮廓（作为函数而不是，方程（5.12））。然后可以重写动量流量（方程（5.64））

\begin{sourceblock}
\includegraphics[page=144,pagebox=mediabox,trim={64.39bp 1638.79bp 64.39bp 261.97bp},clip,width=0.94\textwidth,max height=0.38\textheight]{Pope.pdf}
\end{sourceblock}
积分是无量纲常数，由轮廓的形状决定，但与 x 无关。由于 M 独立于 x，因此乘积 r l , 2 ( \textasciitilde{} ) \textasciitilde{} Om(u\textasciitilde{}st )th 也独立于 x。鉴于实验观察到射流呈线性扩散（drllt/dx - S 常数），平均速度 Uo(x) 不可避免地会随着 x-I 衰减。

对于自相似圆射流，与平均速度场相关的质量 m(x) 和动能 B(x) 的流量为

\begin{sourceblock}
\includegraphics[page=144,pagebox=mediabox,trim={64.39bp 1108.86bp 64.39bp 701.91bp},clip,width=0.94\textwidth,max height=0.38\textheight]{Pope.pdf}
\end{sourceblock}
\begin{sourceblock}
\includegraphics[page=144,pagebox=mediabox,trim={64.39bp 870.89bp 64.39bp 937.88bp},clip,width=0.94\textwidth,max height=0.38\textheight]{Pope.pdf}
\end{sourceblock}
由于 f 的积分和乘积 rllzUo 与 x 无关，因此可以看出，质量流量与 rl/2 成线性比例，因此与 x 成线性比例，而能量流量与 rll2（和 x）成反比。

\subsection*{5.2.3 自相似性}\phantomsection\addcontentsline{toc}{subsection}{5.2.3 自相似性}
对于湍流圆射流，经验观察发现 ( U) Uo(x) 和 ( u u *) U \textasciitilde{} ( X ) a\textasciitilde{}s 函数的轮廓变得自相似
r/rl /-？ (x)（e 取决于 x）。 (U) 的自相似轮廓为等式： （5.65），我们将 定义为

\begin{sourceblock}
\includegraphics[page=144,pagebox=mediabox,trim={64.39bp 309.96bp 64.39bp 1637.79bp},clip,width=0.94\textwidth,max height=0.38\textheight]{Pope.pdf}
\end{sourceblock}
我们现在从边界层方程表明，这种自相似行为意味着射流线性扩散（drlI2/dx S 常数），因此 Uo(x) 随着 x-I 衰减，正如所观察到的那样。

边界层动量方程（方程（5.58））可以写成
其中质数表示关于 的微分。 （练习 5.12 中给出了推导该方程所涉及的步骤。） 平方项
( [ I ) 5，括号仅依赖于，而大括号 (\{ )) 中的项仅依赖于 由于右侧仅依赖于 (，因此左侧不可能有 x 依赖。因此，大括号中的项与 x 无关，即

\begin{sourceblock}
\includegraphics[page=145,pagebox=mediabox,trim={64.39bp 1324.83bp 64.39bp 572.93bp},clip,width=0.94\textwidth,max height=0.38\textheight]{Pope.pdf}
\end{sourceblock}
\begin{sourceblock}
\includegraphics[page=145,pagebox=mediabox,trim={64.39bp 1193.85bp 64.39bp 708.91bp},clip,width=0.94\textwidth,max height=0.38\textheight]{Pope.pdf}
\end{sourceblock}
其中 C 和 S 是常数。 （这个论证取决于这样一个事实：( 依赖项不全为零。)

通过从上述两个方程中消除 C，我们得到

\begin{sourceblock}
\includegraphics[page=145,pagebox=mediabox,trim={64.39bp 933.88bp 64.39bp 996.87bp},clip,width=0.94\textwidth,max height=0.38\textheight]{Pope.pdf}
\end{sourceblock}
表明射流的线性扩散速率是自相似性的必然结果。方程（5.72）意味着 Uo(x) 作为 x 的幂变化，但它并没有确定幂。然而，考虑到 rl/2 随 x 线性变化，我们已经观察到动量守恒要求 Uo(x) 随 x 变化由此得出常数 C 为

\begin{sourceblock}
\includegraphics[page=145,pagebox=mediabox,trim={64.39bp 582.93bp 64.39bp 1319.83bp},clip,width=0.94\textwidth,max height=0.38\textheight]{Pope.pdf}
\end{sourceblock}
从方程 (U) u o ( x ) f ( t ) 和 r 2 ( x ) 可知
自相似圆射流中 (U) 的导数为

\begin{sourceblock}
\includegraphics[page=145,pagebox=mediabox,trim={64.39bp 297.96bp 64.39bp 1602.80bp},clip,width=0.94\textwidth,max height=0.38\textheight]{Pope.pdf}
\end{sourceblock}
\begin{sourceblock}
\includegraphics[page=145,pagebox=mediabox,trim={64.39bp 160.98bp 64.39bp 1739.78bp},clip,width=0.94\textwidth,max height=0.38\textheight]{Pope.pdf}
\end{sourceblock}
方程（方程（5.57）），表明平均横向速度为

\begin{sourceblock}
\includegraphics[page=146,pagebox=mediabox,trim={64.39bp 1493.81bp 64.39bp 395.95bp},clip,width=0.94\textwidth,max height=0.38\textheight]{Pope.pdf}
\end{sourceblock}
其中 ia 和 是积分变量。通过替换这些关系
代入边界层动量方程（方程（5.58），忽略
粘性项），验证方程。 （5.71）。

\subsection*{5.2.4 均匀湍流黏度}\phantomsection\addcontentsline{toc}{subsection}{5.2.4 均匀湍流黏度}
湍流边界层方程呈现闭合问题：有两个方程（连续性和轴向动量），涉及三个因变量 (U)、(V) 和 (uv)。如果可以指定湍流黏度 vT(x,r )，则可以克服此闭合问题，因为剪切应力为

\begin{sourceblock}
\includegraphics[page=146,pagebox=mediabox,trim={64.39bp 875.89bp 64.39bp 985.88bp},clip,width=0.94\textwidth,max height=0.38\textheight]{Pope.pdf}
\end{sourceblock}
对于我们观察到的自相似圆射流，首先，它与
VT r1/2 和 Uo，即

\begin{sourceblock}
\includegraphics[page=146,pagebox=mediabox,trim={64.39bp 722.91bp 64.39bp 1224.85bp},clip,width=0.94\textwidth,max height=0.38\textheight]{Pope.pdf}
\end{sourceblock}
其次，整个射流 gT(q) 与值 0.028 的偏差在 15\% 以内。因此，研究边界层方程的解是合理的，其中 O T ( \textasciitilde{} ) 取为常数，与 q 无关。事实上，由于乘积 rl12(x)Uo(x)i 独立于 x，这对应于取
均匀独立于 x 和 r，因此边界层 VT - 动量方程变为
a r r dr （鉴于假设的高雷诺数，粘性项已被忽略，尽管可以简单地用 v,fi 替换它来保留它。）这是
VT 精确地包含 ( U ) ( ) 和 d 的层流边界层方程
U 的 VT 和 v *

\begin{sourceblock}
\includegraphics[page=147,pagebox=mediabox,trim={64.39bp 1184.85bp 64.39bp 109.44bp},clip,width=\textwidth,max height=0.72\textheight]{Pope.pdf}
\par\vspace{0.45\baselineskip}
\small 图 5.15。自相似圆射流中的平均速度分布：实线，曲线拟合 Hussein 等人的实验数据。 （1994）；虚线，均匀湍流黏度溶液（方程 5.82）。
\end{sourceblock}
最早由 Schlichting (1933) 获得。在这里我们给出解决方案，讨论其一些后果，然后给出推导。

就相似度分布 (U)/Uo 与 r/(x xo) 而言，

\begin{sourceblock}
\includegraphics[page=147,pagebox=mediabox,trim={64.39bp 769.90bp 64.39bp 1080.86bp},clip,width=0.94\textwidth,max height=0.38\textheight]{Pope.pdf}
\end{sourceblock}
其中系数 a 以传播速率 S 的形式给出

\begin{sourceblock}
\includegraphics[page=147,pagebox=mediabox,trim={64.39bp 558.93bp 64.39bp 1293.84bp},clip,width=0.94\textwidth,max height=0.38\textheight]{Pope.pdf}
\end{sourceblock}
该轮廓（S 0.094）与图 5.15 中的 (U)/Uo 测量值进行比较。除射流边缘外，轮廓之间有很好的一致性，其中凭经验确定的湍流黏度
衰减到零（见图 5.10）。铺展率由指定的标准化粘度根据以下公式确定

\begin{sourceblock}
\includegraphics[page=147,pagebox=mediabox,trim={64.39bp 241.97bp 64.39bp 1667.79bp},clip,width=0.94\textwidth,max height=0.38\textheight]{Pope.pdf}
\end{sourceblock}
铺展率 S 0.094 是通过 1T 0.028 获得的，这（毫不奇怪）是从测量中获得的平均值（图 5.10）。

\begin{sourceblock}
\includegraphics[page=148,pagebox=mediabox,trim={64.39bp 1710.78bp 64.39bp 159.98bp},clip,width=0.94\textwidth,max height=0.38\textheight]{Pope.pdf}
\end{sourceblock}
\textasciitilde{} h u 由于均匀湍流黏度近似，湍流圆射流中的平均速度场与雷诺数 35 的层流射流中的速度场相同。

均匀湍流的解

粘度 斯托克斯流函数 y(x, r ) 的引入为

\begin{sourceblock}
\includegraphics[page=148,pagebox=mediabox,trim={64.39bp 1353.83bp 64.39bp 554.93bp},clip,width=0.94\textwidth,max height=0.38\textheight]{Pope.pdf}
\end{sourceblock}
\begin{sourceblock}
\includegraphics[page=148,pagebox=mediabox,trim={64.39bp 1248.84bp 64.39bp 680.91bp},clip,width=0.94\textwidth,max height=0.38\textheight]{Pope.pdf}
\end{sourceblock}
从而自动满足连续性方程（方程（5.57））。 x 从虚拟原点测量（因此 rl x ），E q。 (5.86) 导致

\begin{sourceblock}
\includegraphics[page=148,pagebox=mediabox,trim={64.39bp 957.88bp 64.39bp 853.89bp},clip,width=0.94\textwidth,max height=0.38\textheight]{Pope.pdf}
\end{sourceblock}
由于 x2Uo(x)v 与 x 线性变化，显然存在自相似的标度流函数 F(q) 使得
（包含常数 vT，因此 F(q) 是无量纲的。）

从上面的方程我们得到

\begin{sourceblock}
\includegraphics[page=148,pagebox=mediabox,trim={64.39bp 581.93bp 64.39bp 1342.83bp},clip,width=0.94\textwidth,max height=0.38\textheight]{Pope.pdf}
\end{sourceblock}
\begin{sourceblock}
\includegraphics[page=148,pagebox=mediabox,trim={64.39bp 466.94bp 64.39bp 1431.82bp},clip,width=0.94\textwidth,max height=0.38\textheight]{Pope.pdf}
\end{sourceblock}
其中 F' dF/dq。为了满足轴上为零的条件，F

\begin{sourceblock}
\includegraphics[page=148,pagebox=mediabox,trim={64.39bp 311.96bp 64.39bp 1576.80bp},clip,width=0.94\textwidth,max height=0.38\textheight]{Pope.pdf}
\end{sourceblock}
边界层方程（方程（5.81））中的所有项都可以用F及其导数来表示。化简后的结果是

\begin{sourceblock}
\includegraphics[page=148,pagebox=mediabox,trim={64.39bp 119.39bp 64.39bp 1797.77bp},clip,width=0.94\textwidth,max height=0.38\textheight]{Pope.pdf}
\end{sourceblock}
左边是 ( - q )，因此方程可以积分为

\begin{sourceblock}
\includegraphics[page=149,pagebox=mediabox,trim={64.39bp 1688.79bp 64.39bp 203.97bp},clip,width=0.94\textwidth,max height=0.38\textheight]{Pope.pdf}
\end{sourceblock}
鉴于方程。 （5.92），积分常数为零。方程可以

\begin{sourceblock}
\includegraphics[page=149,pagebox=mediabox,trim={64.39bp 1515.81bp 64.39bp 370.95bp},clip,width=0.94\textwidth,max height=0.38\textheight]{Pope.pdf}
\end{sourceblock}
然后进行第二次积分，再次积分常数为

\begin{sourceblock}
\includegraphics[page=149,pagebox=mediabox,trim={64.39bp 1345.83bp 64.39bp 549.93bp},clip,width=0.94\textwidth,max height=0.38\textheight]{Pope.pdf}
\end{sourceblock}
\begin{sourceblock}
\includegraphics[page=149,pagebox=mediabox,trim={64.39bp 1187.85bp 64.39bp 676.91bp},clip,width=0.94\textwidth,max height=0.38\textheight]{Pope.pdf}
\end{sourceblock}
以积分常数 c 进行第三次积分，我们得到

\begin{sourceblock}
\includegraphics[page=149,pagebox=mediabox,trim={64.39bp 1042.87bp 64.39bp 882.89bp},clip,width=0.94\textwidth,max height=0.38\textheight]{Pope.pdf}
\end{sourceblock}
设置一个e2'，解决办法是

\begin{sourceblock}
\includegraphics[page=149,pagebox=mediabox,trim={64.39bp 843.89bp 64.39bp 1052.87bp},clip,width=0.94\textwidth,max height=0.38\textheight]{Pope.pdf}
\end{sourceblock}
通过对该解求微分，我们找到平均速度分布（方程（5.90））

\begin{sourceblock}
\includegraphics[page=149,pagebox=mediabox,trim={64.39bp 652.92bp 64.39bp 1206.85bp},clip,width=0.94\textwidth,max height=0.38\textheight]{Pope.pdf}
\end{sourceblock}
因此中心线速度为

\begin{sourceblock}
\includegraphics[page=149,pagebox=mediabox,trim={64.39bp 494.94bp 64.39bp 1413.82bp},clip,width=0.94\textwidth,max height=0.38\textheight]{Pope.pdf}
\end{sourceblock}
自相似轮廓是

\begin{sourceblock}
\includegraphics[page=149,pagebox=mediabox,trim={64.39bp 324.96bp 64.39bp 1572.80bp},clip,width=0.94\textwidth,max height=0.38\textheight]{Pope.pdf}
\end{sourceblock}
常数a和湍流黏度vT可以与铺展速率S rll2/x 相关。注意到 r rl/2 对应于 q S，根据 rl/2 的定义我们需要 f(S) - 这导致

\begin{sourceblock}
\includegraphics[page=149,pagebox=mediabox,trim={64.39bp 119.39bp 64.39bp 1823.77bp},clip,width=0.94\textwidth,max height=0.38\textheight]{Pope.pdf}
\end{sourceblock}
然后，从方程。 ( 5 1 0 1 ) 我们得到

\begin{sourceblock}
\includegraphics[page=150,pagebox=mediabox,trim={64.39bp 1662.79bp 64.39bp 206.97bp},clip,width=0.94\textwidth,max height=0.38\textheight]{Pope.pdf}
\end{sourceblock}
\section*{5.3 圆射流：动能}\phantomsection\addcontentsline{toc}{section}{5.3 圆射流：动能}
动能的分解 流体的动能（每单位质量）为

\begin{sourceblock}
\includegraphics[page=150,pagebox=mediabox,trim={64.39bp 1338.83bp 64.39bp 601.92bp},clip,width=0.94\textwidth,max height=0.38\textheight]{Pope.pdf}
\end{sourceblock}
E的均值可以分解为两部分：

\begin{sourceblock}
\includegraphics[page=150,pagebox=mediabox,trim={64.39bp 1210.85bp 64.39bp 732.91bp},clip,width=0.94\textwidth,max height=0.38\textheight]{Pope.pdf}
\end{sourceblock}
其中 E(x, t) 是平均流量的动能

\begin{sourceblock}
\includegraphics[page=150,pagebox=mediabox,trim={64.39bp 1070.86bp 64.39bp 860.89bp},clip,width=0.94\textwidth,max height=0.38\textheight]{Pope.pdf}
\end{sourceblock}
k(x, t) 是湍动能
（可以通过将雷诺分解 U ( U ) u 代入方程（5.105）并取平均值来验证该分解。）

湍动能 k 决定雷诺应力张量的各向同性部分 ( e 'k dij)；但我们也发现各向异性部分随 k 变化。例如，在大部分湍流圆射流上，我们观察到 (uv) 0.27k（图 5.9），剪切应力的数学界限为 Iuvl 5 k（练习 5,13）。因此，k 是一个相当重要的量。在本节中，我们考虑湍流中产生和耗散湍动能的过程。这也导致了对 E、(E) 和 E 的考虑。

根据 Cauchy-Schwartz 不等式和 k 的定义，显示

\begin{sourceblock}
\includegraphics[page=150,pagebox=mediabox,trim={64.39bp 258.97bp 64.39bp 1689.79bp},clip,width=0.94\textwidth,max height=0.38\textheight]{Pope.pdf}
\end{sourceblock}
表明（对于不可压缩流动）

\begin{sourceblock}
\includegraphics[page=150,pagebox=mediabox,trim={64.39bp 119.39bp 64.39bp 1812.77bp},clip,width=0.94\textwidth,max height=0.38\textheight]{Pope.pdf}
\end{sourceblock}
由动量方程可知

\begin{sourceblock}
\includegraphics[page=151,pagebox=mediabox,trim={64.39bp 1695.79bp 64.39bp 213.97bp},clip,width=0.94\textwidth,max height=0.38\textheight]{Pope.pdf}
\end{sourceblock}
其中 zij 是应力张量，证明

\begin{sourceblock}
\includegraphics[page=151,pagebox=mediabox,trim={64.39bp 1532.81bp 64.39bp 375.95bp},clip,width=0.94\textwidth,max height=0.38\textheight]{Pope.pdf}
\end{sourceblock}
其中 Sij '(aUi/8xj aUj/axi) 是应变率张量。对于一个
牛顿流体由下式给出

那么证明动能方程是

\begin{sourceblock}
\includegraphics[page=151,pagebox=mediabox,trim={64.39bp 1187.85bp 64.39bp 720.91bp},clip,width=0.94\textwidth,max height=0.38\textheight]{Pope.pdf}
\end{sourceblock}
瞬时动能 E 的演化方程由纳维-斯托克斯方程得出，为
其中 Sij 1 Ui/axj Uj/axi) 是应变率张量，并且

\begin{sourceblock}
\includegraphics[page=151,pagebox=mediabox,trim={64.39bp 617.92bp 64.39bp 1307.84bp},clip,width=0.94\textwidth,max height=0.38\textheight]{Pope.pdf}
\end{sourceblock}
是能量通量，参见练习 5.14。 E 方程在固定控制体积上的积分为

\begin{sourceblock}
\includegraphics[page=151,pagebox=mediabox,trim={64.39bp 409.95bp 64.39bp 1481.81bp},clip,width=0.94\textwidth,max height=0.38\textheight]{Pope.pdf}
\end{sourceblock}
表面积分考虑了流入、流出和控制面上所做的功：它表示 E 从一个区域到另一区域的转移。重要的是要观察到流动中不存在能量“来源”。分量平方和的数量 Sw*Sij 为正（如果 Sij 的所有分量均为零，则量为零）。因此，右侧是能量的“汇”：它代表粘性耗散，即
通过展开 V 2 ( 1 U U ) 可知

\begin{sourceblock}
\includegraphics[page=152,pagebox=mediabox,trim={64.39bp 1642.79bp 64.39bp 256.97bp},clip,width=0.94\textwidth,max height=0.38\textheight]{Pope.pdf}
\end{sourceblock}
使用该结果，从 Navier-S tokes 方程中获得（方程 (2.35)）

动能方程的另一种形式：

\begin{sourceblock}
\includegraphics[page=152,pagebox=mediabox,trim={64.39bp 1440.82bp 64.39bp 457.94bp},clip,width=0.94\textwidth,max height=0.38\textheight]{Pope.pdf}
\end{sourceblock}
\begin{sourceblock}
\includegraphics[page=152,pagebox=mediabox,trim={64.39bp 1344.83bp 64.39bp 560.93bp},clip,width=0.94\textwidth,max height=0.38\textheight]{Pope.pdf}
\end{sourceblock}
Sij 是应变率张量，Rij L ( \textasciitilde{} U \textasciitilde{} / \textasciitilde{} X \textasciitilde{}
Uj/axi) 是旋转率张量，表明

表明方程。 (5.1 14) 和(5.120) 相同。

表明

\begin{sourceblock}
\includegraphics[page=152,pagebox=mediabox,trim={64.39bp 827.90bp 64.39bp 1069.86bp},clip,width=0.94\textwidth,max height=0.38\textheight]{Pope.pdf}
\end{sourceblock}
因此（从方程（5.1 14））表明

\begin{sourceblock}
\includegraphics[page=152,pagebox=mediabox,trim={64.39bp 666.92bp 64.39bp 1229.84bp},clip,width=0.94\textwidth,max height=0.38\textheight]{Pope.pdf}
\end{sourceblock}
其中 E 和 由方程式定义。 （5.127）和（5.128）。

平均动能 平均动能 (E) 的方程可通过取方程 (Eq) 的平均值来简单获得。 （5.116）：

\begin{sourceblock}
\includegraphics[page=152,pagebox=mediabox,trim={64.39bp 299.96bp 64.39bp 1605.80bp},clip,width=0.94\textwidth,max height=0.38\textheight]{Pope.pdf}
\end{sourceblock}
参见练习 5.18。右边的两项是

\begin{sourceblock}
\includegraphics[page=152,pagebox=mediabox,trim={64.39bp 119.39bp 64.39bp 1845.77bp},clip,width=0.94\textwidth,max height=0.38\textheight]{Pope.pdf}
\end{sourceblock}
Sij 其中 s,i 是应变的平均率和波动率：

\begin{sourceblock}
\includegraphics[page=153,pagebox=mediabox,trim={64.39bp 1683.79bp 64.39bp 224.97bp},clip,width=0.94\textwidth,max height=0.38\textheight]{Pope.pdf}
\end{sourceblock}
\begin{sourceblock}
\includegraphics[page=153,pagebox=mediabox,trim={64.39bp 1539.81bp 64.39bp 357.95bp},clip,width=0.94\textwidth,max height=0.38\textheight]{Pope.pdf}
\end{sourceblock}
第一个贡献 是平均流引起的耗散：一般来说，与其他项相比，它的量级为 \textasciitilde{} e - '，因此可以忽略不计。正如我们将看到的，第二个贡献至关重要。

平均流和湍动能 ( U ) ( U ) 和 k (u u) 的方程可以写成

\begin{sourceblock}
\includegraphics[page=153,pagebox=mediabox,trim={64.39bp 1121.86bp 64.39bp 779.90bp},clip,width=0.94\textwidth,max height=0.38\textheight]{Pope.pdf}
\end{sourceblock}
参见练习 5.19 和 5.20（其中 T 和 由方程（5.136）和 5.140 定义）。数量

\begin{sourceblock}
\includegraphics[page=153,pagebox=mediabox,trim={64.39bp 797.90bp 64.39bp 1150.85bp},clip,width=0.94\textwidth,max height=0.38\textheight]{Pope.pdf}
\end{sourceblock}
通常为正，因此是 k 方程中的“源”：它被称为湍动能的产生或简称为产生。

从雷诺方程（方程（4.12））开始表明，平均
动能方程（对于 E ( U ) ( U ) ） 1s

\begin{sourceblock}
\includegraphics[page=153,pagebox=mediabox,trim={64.39bp 191.98bp 64.39bp 1721.78bp},clip,width=0.94\textwidth,max height=0.38\textheight]{Pope.pdf}
\end{sourceblock}
\begin{sourceblock}
\includegraphics[page=153,pagebox=mediabox,trim={64.39bp 119.39bp 64.39bp 1836.77bp},clip,width=0.94\textwidth,max height=0.38\textheight]{Pope.pdf}
\end{sourceblock}
通过从纳维方程中减去雷诺方程（方程（4.12））

斯托克斯方程（Eq.（2.35）），表明脉动速度u(x,t)

演变为

\begin{sourceblock}
\includegraphics[page=154,pagebox=mediabox,trim={64.39bp 1601.80bp 64.39bp 312.96bp},clip,width=0.94\textwidth,max height=0.38\textheight]{Pope.pdf}
\end{sourceblock}
\begin{sourceblock}
\includegraphics[page=154,pagebox=mediabox,trim={64.39bp 1467.81bp 64.39bp 406.95bp},clip,width=0.94\textwidth,max height=0.38\textheight]{Pope.pdf}
\end{sourceblock}
其中 p' 是脉动压力场 (p' = p - (p))。因此显示
湍动能的演化过程为

\begin{sourceblock}
\includegraphics[page=154,pagebox=mediabox,trim={64.39bp 1146.86bp 64.39bp 743.91bp},clip,width=0.94\textwidth,max height=0.38\textheight]{Pope.pdf}
\end{sourceblock}
生产 E 和 k 的方程清楚地表明了生产所发挥的重要作用。平均速度梯度对抗雷诺应力的作用消除了平均流中的动能（E 的方程（5.131）中的 -p）并将其传递到方程（5.131）中的脉动速度场。 (5.133) 对于 k)。

关于生产的一些观察结果如下。

(i) 只有速度梯度张量的对称部分影响亲

\begin{sourceblock}
\includegraphics[page=154,pagebox=mediabox,trim={64.39bp 593.92bp 64.39bp 1297.84bp},clip,width=0.94\textwidth,max height=0.38\textheight]{Pope.pdf}
\end{sourceblock}
(ii) 只有雷诺应力张量的各向异性部分影响亲

\begin{sourceblock}
\includegraphics[page=154,pagebox=mediabox,trim={64.39bp 424.95bp 64.39bp 1461.82bp},clip,width=0.94\textwidth,max height=0.38\textheight]{Pope.pdf}
\end{sourceblock}
哪里 (u-u 'kdije
(iii) 根据湍流黏度假设 ( e , -2VTSij,
等式。 (4.48)) 的产生式是
可以看出，这个表达式与
(iv) 在边界层近似中，所有平均速度梯度
除 d(U)/dy（或 d ( U )l d r ）生产外，忽略 dients

\begin{sourceblock}
\includegraphics[page=155,pagebox=mediabox,trim={64.39bp 1556.80bp 64.39bp 303.96bp},clip,width=0.94\textwidth,max height=0.38\textheight]{Pope.pdf}
\end{sourceblock}
(v) 根据湍流黏度假说，在边界层
产量的近似值是

\begin{sourceblock}
\includegraphics[page=155,pagebox=mediabox,trim={64.39bp 1362.83bp 64.39bp 534.93bp},clip,width=0.94\textwidth,max height=0.38\textheight]{Pope.pdf}
\end{sourceblock}
证明生产有界于

其中 Sn 是平均值特征值的最大绝对值
应变率张量。

耗散 在 k 方程中，汇是湍动能的耗散，或简称耗散。脉动速度梯度 (aui/axj) 对抗脉动偏应力 (2vsij)，将动能转化为内能。 （如练习 5.22\# 所示，由此产生的温度上升几乎总是可以忽略不计。）从其定义 2v (s4.s 4.) 可以看出，耗散是非负的。

伊吉伊吉？

为了理解 E 最重要的特性，我们回到自相似圆射流。我们已经看到 (U)/Uo 和 (uiuj)/u 的轮廓； （作为 r/r1l2 的函数）是自相似的并且独立于 Re（对于
k / \textasciitilde{} i 足够大的 x/d 和 Re)。因此和

\begin{sourceblock}
\includegraphics[page=155,pagebox=mediabox,trim={64.39bp 407.95bp 64.39bp 1497.81bp},clip,width=0.94\textwidth,max height=0.38\textheight]{Pope.pdf}
\end{sourceblock}
也是自相似且独立于Re。在 k 的平衡方程中，由于 D k / \textasciitilde{} 在 n d 尺度上都与 u,3/r1i2 相同，因此几乎不可避免地具有相同的尺度：即

\begin{sourceblock}
\includegraphics[page=155,pagebox=mediabox,trim={64.39bp 199.97bp 64.39bp 1740.78bp},clip,width=0.94\textwidth,max height=0.38\textheight]{Pope.pdf}
\end{sourceblock}
自相似，与 Re 无关。这通过测量得到证实（例如，
乍一看，这种行为令人费解。假设进行两个高雷诺数圆射流实验，分别用 a 和 b 表示，具有相同的喷嘴直径 d 和射流速度 U，但流体粘度不同，并且在给定 x（在自相似区域中）

两个实验中的 va v,,, 速度 Uo(x) 和半宽 r1l2(x)i 相同。因此，在给定 ( x , ) 时，两个实验中的耗散也相同，即

\begin{sourceblock}
\includegraphics[page=156,pagebox=mediabox,trim={64.39bp 1457.82bp 64.39bp 490.94bp},clip,width=0.94\textwidth,max height=0.38\textheight]{Pope.pdf}
\end{sourceblock}
然而，根据其定义 (sijsij)，与
2 v v，两次实验不一样！那么如何才能和相等呢？答案可以在第 5 页的图 1.2 中找到。可以看出，雷诺数较高的射流具有更精细的小结构规模，因此，梯度更陡，并且
L(aui/ax.aui/axi)。 S i j 2 J
考虑理想气体的自相似圆射流。由于耗散
加热时，中心线温度 To(x) 略高于
环境温度，TI。使用简单的能量平衡
获得粗略估计

从而获得最大温度超出的估计值

其中 M a 是基于 UJ 的马赫数。 (C, 是常数-
压力比热，y 是比热之比。）

柯尔莫哥洛夫尺度 在第六章中，我们将看到最小湍流运动的特征尺度是柯尔莫哥洛夫尺度。这些是由 和 v 形成的长度 (q)、时间 (z,) 和速度 (u,) 尺度：

\begin{sourceblock}
\includegraphics[page=156,pagebox=mediabox,trim={64.39bp 119.39bp 64.39bp 1827.77bp},clip,width=0.94\textwidth,max height=0.38\textheight]{Pope.pdf}
\end{sourceblock}
\begin{sourceblock}
\includegraphics[page=157,pagebox=mediabox,trim={64.39bp 1762.78bp 64.39bp 148.98bp},clip,width=0.94\textwidth,max height=0.38\textheight]{Pope.pdf}
\end{sourceblock}
\begin{sourceblock}
\includegraphics[page=157,pagebox=mediabox,trim={64.39bp 1710.78bp 64.39bp 234.97bp},clip,width=0.94\textwidth,max height=0.38\textheight]{Pope.pdf}
\end{sourceblock}
从这些定义和方程。 (5.147) 由此可见，与平均流量尺度 r 和 Uo 相比，柯尔莫哥洛夫尺度随雷诺数 Reo U Or l / 2 /a\textasciitilde{}的变化，根据

\begin{sourceblock}
\includegraphics[page=157,pagebox=mediabox,trim={64.39bp 1485.81bp 64.39bp 451.94bp},clip,width=0.94\textwidth,max height=0.38\textheight]{Pope.pdf}
\end{sourceblock}
\begin{sourceblock}
\includegraphics[page=157,pagebox=mediabox,trim={64.39bp 1398.82bp 64.39bp 538.93bp},clip,width=0.94\textwidth,max height=0.38\textheight]{Pope.pdf}
\end{sourceblock}
\begin{sourceblock}
\includegraphics[page=157,pagebox=mediabox,trim={64.39bp 1310.83bp 64.39bp 625.92bp},clip,width=0.94\textwidth,max height=0.38\textheight]{Pope.pdf}
\end{sourceblock}
（回想一下，2 是无量纲的并且与雷诺数无关。）因此，与视觉观察一致，我们看到（相对于平均流量尺度）最小运动的尺寸和时间尺度随着雷诺数的增加而减小。

源自柯尔莫哥洛夫量表定义的两个揭示性恒等式是

\begin{sourceblock}
\includegraphics[page=157,pagebox=mediabox,trim={64.39bp 950.88bp 64.39bp 980.88bp},clip,width=0.94\textwidth,max height=0.38\textheight]{Pope.pdf}
\end{sourceblock}
\begin{sourceblock}
\includegraphics[page=157,pagebox=mediabox,trim={64.39bp 810.90bp 64.39bp 1070.86bp},clip,width=0.94\textwidth,max height=0.38\textheight]{Pope.pdf}
\end{sourceblock}
第一个结果表明，无论 Reo 多大，基于柯尔莫哥洛夫尺度的雷诺数都是统一的，这表明这些尺度上的运动受到粘度的强烈影响。第二个表明速度梯度以独立于 v 的方式缩放（如 u v / q 1 ）。

因此，上述难题的解决方案是均方应变率 ( s i j s i j ) 与 v 成反比，即
- ( s - ) 和 = vb (Sijsij) 相等。剩下的问题为什么要这样做
- V, ASeb - \&ar j ij a 小尺度湍流运动是这样缩放的吗？已在章节中讨论
湍动能收支 对于自相似圆射流，湍动能收支如图 5.16 所示。绘制的数量是 k 方程（方程（5.132））中的四项，通过 \textasciitilde{} , 3 / r \textasciitilde{}T/h\textasciitilde{}e 标准化。贡献是产量，；耗散、n\textasciitilde{}ean-flowc 对流、-\textasciitilde{} k / D t 和湍流传输，

\begin{sourceblock}
\includegraphics[page=158,pagebox=mediabox,trim={64.39bp 1152.86bp 64.39bp 109.44bp},clip,width=\textwidth,max height=0.72\textheight]{Pope.pdf}
\par\vspace{0.45\baselineskip}
\small 图 5.16。自相似圆射流中的湍动能收支。数量按 Up 和 1/2 标准化。 （摘自 Panchapakesan 和 Lumley (1993a)。）
\end{sourceblock}
不同的调查之间存在一致性（比如 20\% 以内）。然而，其他两项存在相当大的不确定性，不同实验中的测量结果相差两倍或更多。）在整个射流中，耗散是一个主导项。产量在 r 0.6 处达到峰值，此时 P/Ei 之比约为 0.8。在中心线上， -(uv) a ( \textasciitilde{} ) / a irs 为零（并且随着 r 2 变化），因此那里的产生是由于项 -((u2) - (v2)) a ( U ) / a x （在边界层中被忽略）

PIE 近似）。在射流边缘趋于零，湍流传输平衡
评估和比较与平均流量和 k 相关的不同速率和时间尺度非常有用。这在表 5.2 和图 5.17 中完成。时间尺度 z 和 zp 提供了射流中湍流寿命的测量。以恒定速率耗散能量 k 需要时间 z，以速率 P 产生 k 需要时间 zp。这些时间尺度很大且近似相等：它们相当于在中心线上以速度移动的粒子从虚拟原点出发的飞行时间
TJ Uo(x);它们大约是施加剪切力时间尺度的三倍
动荡是长期存在的。

\begin{sourceblock}
\includegraphics[page=159,pagebox=mediabox,trim={64.54bp 1667.58bp 64.54bp 109.48bp},clip,width=\textwidth,max height=0.72\textheight]{Pope.pdf}
\par\vspace{0.45\baselineskip}
\small 图 5.17。自相似圆射流中的时间尺度，单位为 to。定义见表 5.2。
\end{sourceblock}
\begin{sourceblock}
\includegraphics[page=159,pagebox=mediabox,trim={64.54bp 1398.65bp 64.54bp 382.90bp},clip,width=\textwidth,max height=0.72\textheight]{Pope.pdf}
\par\vspace{0.45\baselineskip}
\small 图 5.18。自相似圆射流的长度尺度，单位为 1/2。 Lj; Ly 为纵向和横向积分尺度； L = k?/?/e;1 <= vr/u';在 r/ryj 进行评估。 \textasciitilde{}0.7。 （注意对数刻度。）
\end{sourceblock}
\begin{center}\small 表 5.2。自相似圆射流中的时间尺度、速率和比率：前四个条目根据 Uo(x) 和扩散速率 S\$ 进行评估；其余条目根据\textasciitilde{} 0.7 的实验数据估计，其中 \{uv) 和 |@(U)/ér|峰值\end{center}
\begin{sourceblock}
\includegraphics[page=159,pagebox=mediabox,trim={64.54bp 787.80bp 64.54bp 924.77bp},clip,width=\textwidth,max height=0.72\textheight]{Pope.pdf}
\end{sourceblock}
平均飞行时间从

虚拟原点

夹带率

轴向应变率

应变率

湍流衰减率

湍流产生率

产量比例

应变率与

衰减率

L I I 和 有直接物理意义，Z vT/ur 和

\begin{sourceblock}
\includegraphics[page=160,pagebox=mediabox,trim={64.39bp 1714.78bp 64.39bp 228.97bp},clip,width=0.94\textwidth,max height=0.38\textheight]{Pope.pdf}
\end{sourceblock}
不。

参考表 5.2，验证

\begin{sourceblock}
\includegraphics[page=160,pagebox=mediabox,trim={64.39bp 1406.82bp 64.39bp 532.93bp},clip,width=0.94\textwidth,max height=0.38\textheight]{Pope.pdf}
\end{sourceblock}
对于归一化速度分布 f 1/ ( I + av12)2s，如何
la(U)/arl 的最大值为
并发生在

赝耗散 赝耗散定义为
与真实耗散有关

（参见练习 5.25）。几乎在所有情况下，方程中的最后一项。 （5.160）很小（至多百分之几，因此两者之间的区别
E 很少很重要。事实上，许多作者将 E 称为“耗散”。正如练习 5.27 所示，当 E 为
获得耗散和功率之间的以下关系

赝耗散

\begin{sourceblock}
\includegraphics[page=161,pagebox=mediabox,trim={64.39bp 1491.81bp 64.39bp 309.96bp},clip,width=0.94\textwidth,max height=0.38\textheight]{Pope.pdf}
\end{sourceblock}
湍流动力学中的耗散项和粘性扩散项

能量方程由表达式得出 表明这
可以用其他形式重新表达

\begin{sourceblock}
\includegraphics[page=161,pagebox=mediabox,trim={64.39bp 1145.86bp 64.39bp 797.90bp},clip,width=0.94\textwidth,max height=0.38\textheight]{Pope.pdf}
\end{sourceblock}
\begin{sourceblock}
\includegraphics[page=161,pagebox=mediabox,trim={64.39bp 878.89bp 64.39bp 1029.87bp},clip,width=0.94\textwidth,max height=0.38\textheight]{Pope.pdf}
\end{sourceblock}
表明湍动能方程（Eq.（5.139））可以
或者写成

\begin{sourceblock}
\includegraphics[page=161,pagebox=mediabox,trim={64.39bp 635.92bp 64.39bp 1258.84bp},clip,width=0.94\textwidth,max height=0.38\textheight]{Pope.pdf}
\end{sourceblock}
在均匀各向同性湍流中，四阶张量
是各向同性的，因此可以写成
其中 a 、 p 和 y 是标量。鉴于连续性方程
0，表明标量之间的关系是

\begin{sourceblock}
\includegraphics[page=161,pagebox=mediabox,trim={64.39bp 119.39bp 64.39bp 1837.77bp},clip,width=0.94\textwidth,max height=0.38\textheight]{Pope.pdf}
\end{sourceblock}
通过考虑 (8/axj)(uia uj/axe) （由于 ho- 其为零
运动性）表明
为零，因此

\begin{sourceblock}
\includegraphics[page=162,pagebox=mediabox,trim={64.39bp 1548.80bp 64.39bp 399.95bp},clip,width=0.94\textwidth,max height=0.38\textheight]{Pope.pdf}
\end{sourceblock}
显示方程。 (5.165) 则变为

\begin{sourceblock}
\includegraphics[page=162,pagebox=mediabox,trim={64.39bp 1368.83bp 64.39bp 527.93bp},clip,width=0.94\textwidth,max height=0.38\textheight]{Pope.pdf}
\end{sourceblock}
表明

\section*{5.4 其他自相似流动}\phantomsection\addcontentsline{toc}{section}{5.4 其他自相似流动}
我们已经详细研究了自相似圆射流中的平均速度和雷诺应力。我们现在简要地考察一下其他经典的自由

\subsection*{5.4.1 喷气式飞机}\phantomsection\addcontentsline{toc}{subsection}{5.4.1 喷气式飞机}
尾流 - 以及均匀剪切流（即受到恒定且均匀的平均剪切的均匀湍流）。然后，在 5.5 节中，描述了这些流程的一些其他特征。

5.4.1 平面射流 理想的平面射流（如第 112 页图 5.14 所示）在统计学上是二维的。平均流的主导方向为 x，跨流坐标为 y，统计数据与展向坐标 z 无关。关于平面 y 0 存在统计对称性。在实验室实验中，有一个槽高为 d（y 方向）的矩形喷嘴，并且

\begin{sourceblock}
\includegraphics[page=163,pagebox=mediabox,trim={64.39bp 1366.83bp 64.39bp 109.44bp},clip,width=\textwidth,max height=0.72\textheight]{Pope.pdf}
\par\vspace{0.45\baselineskip}
\small 图 5.19。自相似平面射流中的平均速度剖面。符号，Heskestad（1965）的实验数据；线，均匀湍流黏度解决方案，方程。 (5.187)（经 ASME 许可）。
\end{sourceblock}
50) 因此，对于 0（达到良好的近似），流量在统计上是二维的并且没有末端效应，至少对于 x/w 不太大。

正如在圆射流中一样，中心线速度 Uo(x) 和半宽 Y \textasciitilde{} / \textasciitilde{} (aXre) 定义为

\begin{sourceblock}
\includegraphics[page=163,pagebox=mediabox,trim={64.39bp 978.88bp 64.39bp 966.88bp},clip,width=0.94\textwidth,max height=0.38\textheight]{Pope.pdf}
\end{sourceblock}
\begin{sourceblock}
\includegraphics[page=163,pagebox=mediabox,trim={64.39bp 876.89bp 64.39bp 1062.87bp},clip,width=0.94\textwidth,max height=0.38\textheight]{Pope.pdf}
\end{sourceblock}
在实验中（例如，Heskestad (1965)、Bradbury (1965) 以及 Gutmark 和 Wygnanski (1976)），发现当用 Uo(x) 和 ylI2(x) 缩放时，平均速度和雷诺应力分布变得自相似（超过 x/d 40）。这些轮廓如图 1 和 2 所示。 5.19 和 5.20。可以看出，轮廓形状和雷诺应力水平与在圆射流中观察到的相当。

发现 yli2(x) 的变化是线性的，即

\begin{sourceblock}
\includegraphics[page=163,pagebox=mediabox,trim={64.39bp 431.95bp 64.39bp 1475.81bp},clip,width=0.94\textwidth,max height=0.38\textheight]{Pope.pdf}
\end{sourceblock}
其中传播速率 S 是常数，S 0.10。然而，与圆射流相比，我们发现 Uo(xj 随着 x-'12 变化。如现在所示，这些变化是自相似性的结果。

在保守形式中，边界层方程（忽略粘性

\begin{sourceblock}
\includegraphics[page=163,pagebox=mediabox,trim={64.39bp 119.39bp 64.39bp 1755.78bp},clip,width=0.94\textwidth,max height=0.38\textheight]{Pope.pdf}
\end{sourceblock}
\begin{sourceblock}
\includegraphics[page=164,pagebox=mediabox,trim={64.39bp 1104.86bp 64.39bp 109.44bp},clip,width=\textwidth,max height=0.72\textheight]{Pope.pdf}
\par\vspace{0.45\baselineskip}
\small 图 5.20。自相似平面射流中的雷诺应力分布。根据 Heskestad (1965) 的测量（经 ASME 许可）。
\end{sourceblock}
对 y 积分，我们得到

\begin{sourceblock}
\includegraphics[page=164,pagebox=mediabox,trim={64.39bp 845.89bp 64.39bp 1057.87bp},clip,width=0.94\textwidth,max height=0.38\textheight]{Pope.pdf}
\end{sourceblock}
因为 ( U ) 和 (uv) 对于 y + +-co 为零。因此动量流量（每单位跨度）

\begin{sourceblock}
\includegraphics[page=164,pagebox=mediabox,trim={64.39bp 638.92bp 64.39bp 1263.84bp},clip,width=0.94\textwidth,max height=0.38\textheight]{Pope.pdf}
\end{sourceblock}
是守恒的（独立于 x）。

在自相似区域，平均轴向速度分布为

\begin{sourceblock}
\includegraphics[page=164,pagebox=mediabox,trim={64.39bp 462.94bp 64.39bp 1482.81bp},clip,width=0.94\textwidth,max height=0.38\textheight]{Pope.pdf}
\end{sourceblock}
\begin{sourceblock}
\includegraphics[page=164,pagebox=mediabox,trim={64.39bp 334.96bp 64.39bp 1612.80bp},clip,width=0.94\textwidth,max height=0.38\textheight]{Pope.pdf}
\end{sourceblock}
所以动量流量为

\begin{sourceblock}
\includegraphics[page=164,pagebox=mediabox,trim={64.39bp 168.98bp 64.39bp 1737.78bp},clip,width=0.94\textwidth,max height=0.38\textheight]{Pope.pdf}
\end{sourceblock}
我们从中推断出的观察结果

\begin{sourceblock}
\includegraphics[page=165,pagebox=mediabox,trim={64.39bp 1688.79bp 64.39bp 214.97bp},clip,width=0.94\textwidth,max height=0.38\textheight]{Pope.pdf}
\end{sourceblock}
将自相似轮廓替换为 (U) 时，E q。 （5.178），以及

\begin{sourceblock}
\includegraphics[page=165,pagebox=mediabox,trim={64.39bp 1545.81bp 64.39bp 393.95bp},clip,width=0.94\textwidth,max height=0.38\textheight]{Pope.pdf}
\end{sourceblock}
代入边界层方程，我们得到

\begin{sourceblock}
\includegraphics[page=165,pagebox=mediabox,trim={64.39bp 1353.83bp 64.39bp 541.93bp},clip,width=0.94\textwidth,max height=0.38\textheight]{Pope.pdf}
\end{sourceblock}
（这涉及与圆射流相同的操作，请参见练习 5.12。） (5.183)，由于右边和括号中的项与 x 无关，因此 d \textasciitilde{} \textasciitilde{} / \textasciitilde{}m/ udstx a 也与 x 无关。因此，自相似性要求传播速率 S dylI2/dx 恒定；然后动量守恒要求 Uo 变化为 x -1/2（方程（5.180））。

( U ) 和 ( u v ) 的自相似性意味着湍流的自相似性

\begin{sourceblock}
\includegraphics[page=165,pagebox=mediabox,trim={64.39bp 911.88bp 64.39bp 963.88bp},clip,width=0.94\textwidth,max height=0.38\textheight]{Pope.pdf}
\end{sourceblock}
可以看出，对于自相似平面射流，随着 x 的增加而增加
VT 当地雷诺数

\begin{sourceblock}
\includegraphics[page=165,pagebox=mediabox,trim={64.39bp 724.91bp 64.39bp 1219.85bp},clip,width=0.94\textwidth,max height=0.38\textheight]{Pope.pdf}
\end{sourceblock}
另一方面，湍流雷诺数

\begin{sourceblock}
\includegraphics[page=165,pagebox=mediabox,trim={64.39bp 570.93bp 64.39bp 1364.83bp},clip,width=0.94\textwidth,max height=0.38\textheight]{Pope.pdf}
\end{sourceblock}
独立于 x。

如果湍流黏度被认为在整个流动中是均匀的（即 0,' 常数），则可以求解边界层方程（方程（5.183））的自相似形式（见下文）以得出

\begin{sourceblock}
\includegraphics[page=165,pagebox=mediabox,trim={64.39bp 296.96bp 64.39bp 1649.79bp},clip,width=0.94\textwidth,max height=0.38\textheight]{Pope.pdf}
\end{sourceblock}
其中 ln(1 ) 该结果与实验数据进行比较
产生观察到的扩散速率 S 0.1 的缩放湍流黏度 OT
GV对应于

（见方程（5.200））。这可以与圆射流的 RT w 35 进行比较。

均匀湍流黏度的解 我们现在获得自相似边界层方程（方程（5.183））的解，剪切应力由湍流黏度假设给出，并且 OT 被视为均匀的，即独立于 5。方程为

\begin{sourceblock}
\includegraphics[page=166,pagebox=mediabox,trim={64.39bp 1283.84bp 64.39bp 609.92bp},clip,width=0.94\textwidth,max height=0.38\textheight]{Pope.pdf}
\end{sourceblock}
替换起来很方便

\begin{sourceblock}
\includegraphics[page=166,pagebox=mediabox,trim={64.39bp 1126.86bp 64.39bp 768.90bp},clip,width=0.94\textwidth,max height=0.38\textheight]{Pope.pdf}
\end{sourceblock}
由于 F(5) 是偶函数，因此 F(5) 是奇函数。特别是 F(0) F"(0) 0。通过这种替换，我们得到

\begin{sourceblock}
\includegraphics[page=166,pagebox=mediabox,trim={64.39bp 952.88bp 64.39bp 984.88bp},clip,width=0.94\textwidth,max height=0.38\textheight]{Pope.pdf}
\end{sourceblock}
注意方括号中的术语是

\begin{sourceblock}
\includegraphics[page=166,pagebox=mediabox,trim={64.39bp 821.90bp 64.39bp 1116.86bp},clip,width=0.94\textwidth,max height=0.38\textheight]{Pope.pdf}
\end{sourceblock}
我们可以积分两次得到

\begin{sourceblock}
\includegraphics[page=166,pagebox=mediabox,trim={64.39bp 694.91bp 64.39bp 1251.84bp},clip,width=0.94\textwidth,max height=0.38\textheight]{Pope.pdf}
\end{sourceblock}
由于F\textasciitilde{}和F'是偶函数，积分常数b为零，而边界条件F'(0) 1 确定a 9,。定义

\begin{sourceblock}
\includegraphics[page=166,pagebox=mediabox,trim={64.39bp 499.94bp 64.39bp 1424.82bp},clip,width=0.94\textwidth,max height=0.38\textheight]{Pope.pdf}
\end{sourceblock}
等式。 (5.192) 则变为

\begin{sourceblock}
\includegraphics[page=166,pagebox=mediabox,trim={64.39bp 367.95bp 64.39bp 1573.80bp},clip,width=0.94\textwidth,max height=0.38\textheight]{Pope.pdf}
\end{sourceblock}
可以积分得到

\begin{sourceblock}
\includegraphics[page=166,pagebox=mediabox,trim={64.39bp 119.39bp 64.39bp 1698.79bp},clip,width=0.94\textwidth,max height=0.38\textheight]{Pope.pdf}
\end{sourceblock}
根据我们的定义有

\begin{sourceblock}
\includegraphics[page=167,pagebox=mediabox,trim={64.39bp 1717.78bp 64.39bp 218.97bp},clip,width=0.94\textwidth,max height=0.38\textheight]{Pope.pdf}
\end{sourceblock}
从中我们得到

\begin{sourceblock}
\includegraphics[page=167,pagebox=mediabox,trim={64.39bp 1561.80bp 64.39bp 345.96bp},clip,width=0.94\textwidth,max height=0.38\textheight]{Pope.pdf}
\end{sourceblock}
这与方程式一起。 (5.193)，将扩散速率 S 与缩放的湍流黏度相关联：

\begin{sourceblock}
\includegraphics[page=167,pagebox=mediabox,trim={64.39bp 1389.82bp 64.39bp 556.93bp},clip,width=0.94\textwidth,max height=0.38\textheight]{Pope.pdf}
\end{sourceblock}
\begin{sourceblock}
\includegraphics[page=167,pagebox=mediabox,trim={64.39bp 1235.84bp 64.39bp 634.92bp},clip,width=0.94\textwidth,max height=0.38\textheight]{Pope.pdf}
\end{sourceblock}
取实验值S=0.1

\subsection*{5.4.2 平面混合层}\phantomsection\addcontentsline{toc}{subsection}{5.4.2 平面混合层}
正如第 112 页图 5.14 所示，混合层是在两个不同速度 Uh 和 U1 (Uh > Ul 0) 的均匀、几乎平行的流之间形成的湍流。这种混合层形成于流入静止环境（U1 0）的平面射流（Uh UJ）的边缘（以及初始区域）。或者，可以在风洞中创建流动，使用分流板将 x < 0 的两股流分开，然后在 x > 0 时形成混合层。

就像在平面射流中一样，流动的主要方向是 x；横流坐标是 y；统计数据与展向坐标 z 无关。与圆形和平面射流相比，对于混合层有两个施加的速度，Uh 和 Ul。因此，流量取决于无量纲参数 U1/ Uh，并且可以定义两个特征速度：特征对流速度

\begin{sourceblock}
\includegraphics[page=167,pagebox=mediabox,trim={64.39bp 391.95bp 64.39bp 1554.80bp},clip,width=0.94\textwidth,max height=0.38\textheight]{Pope.pdf}
\end{sourceblock}
和特征速度差

\begin{sourceblock}
\includegraphics[page=167,pagebox=mediabox,trim={64.39bp 252.97bp 64.39bp 1694.79bp},clip,width=0.94\textwidth,max height=0.38\textheight]{Pope.pdf}
\end{sourceblock}
所有提到的速度（Uh、U、U 和 Us）都是常数独立的

\begin{sourceblock}
\includegraphics[page=168,pagebox=mediabox,trim={64.39bp 1285.84bp 64.39bp 109.44bp},clip,width=\textwidth,max height=0.72\textheight]{Pope.pdf}
\par\vspace{0.45\baselineskip}
\small 图 5.21。平均速度 (U) 相对于 y 的草图，以及显示 yo1、yoo 和 6 定义的缩放平均速度剖面图。
\end{sourceblock}
基于平均速度剖面的方法（U(x,y，这是独立的
2 ) ) 的 z。对于 0 < cx < 1，我们定义跨流位置 y,(x) 使得

\begin{sourceblock}
\includegraphics[page=168,pagebox=mediabox,trim={64.39bp 1017.87bp 64.39bp 927.88bp},clip,width=0.94\textwidth,max height=0.38\textheight]{Pope.pdf}
\end{sourceblock}
然后取 6(x) 为

\begin{sourceblock}
\includegraphics[page=168,pagebox=mediabox,trim={64.39bp 880.89bp 64.39bp 1063.87bp},clip,width=0.94\textwidth,max height=0.38\textheight]{Pope.pdf}
\end{sourceblock}
见图 5.21。 IA 另外，参考横向位置\textasciitilde{}（xis）由'定义

\begin{sourceblock}
\includegraphics[page=168,pagebox=mediabox,trim={64.39bp 740.91bp 64.39bp 1197.85bp},clip,width=0.94\textwidth,max height=0.38\textheight]{Pope.pdf}
\end{sourceblock}
然后缩放的跨流坐标定义为

\begin{sourceblock}
\includegraphics[page=168,pagebox=mediabox,trim={64.39bp 607.92bp 64.39bp 1316.83bp},clip,width=0.94\textwidth,max height=0.38\textheight]{Pope.pdf}
\end{sourceblock}
和缩放速度

\begin{sourceblock}
\includegraphics[page=168,pagebox=mediabox,trim={64.39bp 471.94bp 64.39bp 1473.81bp},clip,width=0.94\textwidth,max height=0.38\textheight]{Pope.pdf}
\end{sourceblock}
从这些定义我们有 f (+a) 和 f k0.4 参见
请注意，练习 5.29-5.32 中使用了 \textasciitilde{} ( x i)s 的不同定义，其中控制方程

图 5.22t。层。

平面混合中的比例速度分布 符号、实验数据
Champagne et (1976) x 厘米； Ø，X厘米； o，x 厘米）；线，
误差函数曲线（方程所示为
y \textasciitilde{} 、ly、a s 和在平面混合层中的轴向变化，显示
传播。等人。 （1976）。

Champagne等的线性实验数据

至方程式(5.207)；图 5.23 清楚地表明混合层呈线性扩展。实验还表明雷诺应力是自相似的。应该注意的是，流动既不关于 y 0 对称，也不关于 y 0 对称。

。。 - \&.即使约为 0，它也会优先传播到低速流中。

正如圆射流和平面射流一样，混合层的线性扩展是自相似性的必然结果，如练习 5.29 所示，使用不同的 ji(x) 定义，并且 g(5) (uv)/u2 为缩放剪应力，自相似混合层的边界层方程为

\begin{sourceblock}
\includegraphics[page=169,pagebox=mediabox,trim={64.39bp 151.98bp 64.39bp 1745.78bp},clip,width=0.94\textwidth,max height=0.38\textheight]{Pope.pdf}
\end{sourceblock}
速率 d6 /dx 和参数

\begin{sourceblock}
\includegraphics[page=170,pagebox=mediabox,trim={64.39bp 1694.79bp 64.39bp 210.97bp},clip,width=0.94\textwidth,max height=0.38\textheight]{Pope.pdf}
\end{sourceblock}
正如所观察到的，必须是独立于 x 的常数。

对于以速度 U, 沿 x 方向行进的观察者，混合层的分数 g:生长速率为 U, d l n 6 l d x。如果该速率按本地时间尺度 6/Us 归一化，则得到的无量纲参数为

\begin{sourceblock}
\includegraphics[page=170,pagebox=mediabox,trim={64.39bp 1407.82bp 64.39bp 494.94bp},clip,width=0.94\textwidth,max height=0.38\textheight]{Pope.pdf}
\end{sourceblock}
因此可以预期，发现S近似独立于速度比AS，使得d6/dx（近似）随着Us/U变化。从一个实验到另一个实验，S 的测量值存在相当大的变化，这至少部分归因于离开分流板时的流动状态 (x 0)。报告值的范围为 S w 0.06 至 S w 0.11 (Dimotakis 1991)。在Champagne等人的实验中
(1976)，值为 S w 0.097。等人。

需要考虑的一个有趣的限制是 Us/ U, 0，对应于 CJl/ Uh 1。在此限制中，边界层方程 Eq. 中的 Us/ U 中的项。 （5.208）消失。该方程对应于

\begin{sourceblock}
\includegraphics[page=170,pagebox=mediabox,trim={64.39bp 864.89bp 64.39bp 1041.87bp},clip,width=0.94\textwidth,max height=0.38\textheight]{Pope.pdf}
\end{sourceblock}
其余对流项 ( ( U ) U,) d(U)/dx 和 ( V ) d(U)/dy 与 U、d(U)/dx 相比可忽略不计。观察者以速度 U 沿 x 方向行进，会看到两个流（y co 和 y -co）向右和向左移动，速度分别为 和 - U s。与 y 方向上的梯度相比，x 方向上的均值梯度非常小（Us/Uc 量级）。该层的厚度随时间以 S Us 的速率增长。因此，在移动坐标系中，当 Us/U 趋于零时，流量在统计上变得一维且与时间相关。它被称为时间混合层（与实验室坐标中的空间混合层相对）。时间混合层关于 y 0 统计对称。

Rogers 和 Moser (1994) 描述了时间混合层的直接数值模拟 (DNS)。纳维-斯托克斯方程通过具有 512 x 210 x 192 模式（在 x、y 和方向）的谱方法求解。在初始瞬态之后，混合层变得自相似，并且宽度6随时间线性增加。观察到的扩散参数 S w 0.062 接近实验观察到的低端

\begin{sourceblock}
\includegraphics[page=171,pagebox=mediabox,trim={64.39bp 1435.82bp 64.39bp 109.44bp},clip,width=\textwidth,max height=0.72\textheight]{Pope.pdf}
\par\vspace{0.45\baselineskip}
\small 图 5.24。自相似平面混合层中的缩放平均速度分布。符号，Bell 和 Mehta 的实验 (1990) = 0.6)；实线，时间混合层的 DNS 数据（Rogers 和 Moser 1994）；虚线，误差函数轮廓，其宽度选择为与层中心的数据匹配。
\end{sourceblock}
1对应于时间混合层。然而，Bell 和 Mehta (1990) 使用 Ul/Uh 0.6 进行的实验产生的结果与 Rogers 和 Moser (1994) 的时间混合层的结果相似。实验中扩散参数为S 0.069。

图 5.24 显示，U1/Uh 为 0.6 的时间混合层和空间混合层的缩放平均速度分布是无法区分的。图 5.24 还显示了误差函数曲线，它是时间混合层的恒定湍流黏度解（参见练习 5.33）。正如射流一样，与测量结果相比，恒定湍流黏度解给出的平均速度分布更趋向于自由流速度。图 5.25 显示了按比例缩放的雷诺应力分布，两个混合层的分布略有不同。

对于混合层，由于 Us 是固定的并且 6 随 x 线性变化，因此雷诺数 Reo@) U,6/v 和湍流黏度也随 x 线性增加。湍动能流速 k ( x ) (U)k dy
-co 随 U , U : \textasciitilde{} 缩放，因此也随着 x 线性增加。这与喷流和尾流相反，其中随着 x 减少。因为混合层中的 x 会随着流的平均而增加，所以产量必须超过耗散
在该层的中心 Rogers 和 Moser (1994) 观察到 w 1.4。

在以下练习中，开发了空间和时间混合层的相似方程。时间混合层是对称的（ i . , y 为零且 f 是奇函数），并且自由流是并行的（ ix. ,
(V)y=-m (V)y=m 0)。空间混合层不是对称的，它优先扩散到低速流中（dpldx为负），并且夹带
(V)y=m (V)y=-m 流体（ e . 为零时为正）。因此两个

\begin{sourceblock}
\includegraphics[page=172,pagebox=mediabox,trim={64.39bp 629.92bp 64.39bp 208.77bp},clip,width=\textwidth,max height=0.72\textheight]{Pope.pdf}
\par\vspace{0.45\baselineskip}
\small 图 5.25。自相似平面混合层中的缩放雷诺应力分布。符号，Bell 和 Mehta 的实验 (1990) = 0.6)；实线，时间混合层的 DNS 数据（Rogers 和 Moser 1994）。
\end{sourceblock}
通过调整风洞壁的倾斜度来保持大致均匀的速度。）

根据 和缩放速度的定义（方程（5.206）

和（5.207））表明

\begin{sourceblock}
\includegraphics[page=172,pagebox=mediabox,trim={64.39bp 119.39bp 64.39bp 1786.77bp},clip,width=0.94\textwidth,max height=0.38\textheight]{Pope.pdf}
\end{sourceblock}
在自相似混合层中，无论定义如何
因此，根据平均连续性方程，表明
选择坐标系，使 x 轴平行于
高速自由流中的速度，因此 ( V )

为零。然后证明横向平均速度为

\begin{sourceblock}
\includegraphics[page=173,pagebox=mediabox,trim={64.39bp 1218.85bp 64.39bp 676.91bp},clip,width=0.94\textwidth,max height=0.38\textheight]{Pope.pdf}
\end{sourceblock}
因此证明边界层方程可以写成

\begin{sourceblock}
\includegraphics[page=173,pagebox=mediabox,trim={64.39bp 1047.87bp 64.39bp 847.89bp},clip,width=0.94\textwidth,max height=0.38\textheight]{Pope.pdf}
\end{sourceblock}
\begin{sourceblock}
\includegraphics[page=173,pagebox=mediabox,trim={64.39bp 927.88bp 64.39bp 960.88bp},clip,width=0.94\textwidth,max height=0.38\textheight]{Pope.pdf}
\end{sourceblock}
上述结果适用于 y(x) 的任何规范。我们现在
明确规定 y(x) 是峰值的位置
通过考虑等式。 (5.21 5) 在 0 处，表明
剪切应力 (gJ.=

\begin{sourceblock}
\includegraphics[page=173,pagebox=mediabox,trim={64.39bp 569.93bp 64.39bp 1237.84bp},clip,width=0.94\textwidth,max height=0.38\textheight]{Pope.pdf}
\end{sourceblock}
表明，根据边界层方程的规定

\begin{sourceblock}
\includegraphics[page=173,pagebox=mediabox,trim={64.39bp 365.95bp 64.39bp 1481.81bp},clip,width=0.94\textwidth,max height=0.38\textheight]{Pope.pdf}
\end{sourceblock}
证明夹带速度为

\begin{sourceblock}
\includegraphics[page=173,pagebox=mediabox,trim={64.39bp 198.97bp 64.39bp 1700.79bp},clip,width=0.94\textwidth,max height=0.38\textheight]{Pope.pdf}
\end{sourceblock}
讨论时空混合的 (V) 符号
通过积分方程。 (5.218) 从 到 W，得到关系式
由此表明空间混合层不可能是对称的。

如果是通过误差函数，那么

\begin{sourceblock}
\includegraphics[page=174,pagebox=mediabox,trim={64.39bp 1336.83bp 64.39bp 549.93bp},clip,width=0.94\textwidth,max height=0.38\textheight]{Pope.pdf}
\end{sourceblock}
通过这个近似，得到结果

\begin{sourceblock}
\includegraphics[page=174,pagebox=mediabox,trim={64.39bp 1170.85bp 64.39bp 732.91bp},clip,width=0.94\textwidth,max height=0.38\textheight]{Pope.pdf}
\end{sourceblock}
对于 Uh/Ul = 2 的混合层，将铺展率 S 取为
为 0.09，计算等式右侧。 (5.222) 并表明，
在低速流中，流线与流线之间的夹角
轴约为 1。

根据湍流黏度假设，g = -OTff，并假设 OT
为了保持一致，证明动量方程（方程 5.208）

时间混合层减少为

\begin{sourceblock}
\includegraphics[page=174,pagebox=mediabox,trim={64.39bp 759.90bp 64.39bp 1188.85bp},clip,width=0.94\textwidth,max height=0.38\textheight]{Pope.pdf}
\end{sourceblock}
证明该方程的解（满足适当的
边界条件）是

\begin{sourceblock}
\includegraphics[page=174,pagebox=mediabox,trim={64.39bp 292.96bp 64.39bp 1655.79bp},clip,width=0.94\textwidth,max height=0.38\textheight]{Pope.pdf}
\end{sourceblock}
并且满足条件 f (+A) = + - 0 .4

\begin{sourceblock}
\includegraphics[page=174,pagebox=mediabox,trim={64.39bp 119.39bp 64.39bp 1799.77bp},clip,width=0.94\textwidth,max height=0.38\textheight]{Pope.pdf}
\end{sourceblock}
对于自相似时间混合层，从时刻开始
tum 方程（Eq. (5.218)），表明归一化剪切应力
层的中心由下式给出

\begin{sourceblock}
\includegraphics[page=175,pagebox=mediabox,trim={64.39bp 1597.80bp 64.39bp 304.96bp},clip,width=0.94\textwidth,max height=0.38\textheight]{Pope.pdf}
\end{sourceblock}
如果由误差函数轮廓近似（方程（5.224））

\begin{sourceblock}
\includegraphics[page=175,pagebox=mediabox,trim={64.39bp 1396.82bp 64.39bp 456.94bp},clip,width=0.94\textwidth,max height=0.38\textheight]{Pope.pdf}
\end{sourceblock}
g(0) 和 S 的测量值与此的吻合程度如何

\subsection*{5.4.3 飞机尾流}\phantomsection\addcontentsline{toc}{subsection}{5.4.3 飞机尾流}
如图 5.14(c)（第 112 页）所示，当均匀流（速度为 U，沿 x 方向）流过圆柱体（与 z 轴对齐）时，会形成平面尾流。流动在统计上是平稳的、二维的并且关于平面 y 0 对称。

特征对流速度为自由流速度 Uc，特征速度差为

\begin{sourceblock}
\includegraphics[page=175,pagebox=mediabox,trim={64.39bp 774.90bp 64.39bp 1170.85bp},clip,width=0.94\textwidth,max height=0.38\textheight]{Pope.pdf}
\end{sourceblock}
半宽 yli2(x) 的定义如下

\begin{sourceblock}
\includegraphics[page=175,pagebox=mediabox,trim={64.39bp 628.92bp 64.39bp 1314.83bp},clip,width=0.94\textwidth,max height=0.38\textheight]{Pope.pdf}
\end{sourceblock}
正如预期的那样，随着下游距离的增加，尾流会扩散（yllz 增加）并衰减（U,/U, 减小到零）。

正如混合层一样，有两种不同的速度尺度，Us 和 U,。在混合层中，它们具有恒定的比率，与 x 无关。随后，该比率随着 U/U 的衰减而变化。因此，流不可能完全自相似；但随着 U,/U 趋于零，它确实会在不久的将来变得渐近自相似。在实验中，当该速度比小于约 1 时，会观察到自相似行为
当 y/yl12(x) 为缩放的跨流变量时，自相似速度亏损 f 定义为

\begin{sourceblock}
\includegraphics[page=175,pagebox=mediabox,trim={64.39bp 119.39bp 64.39bp 1848.77bp},clip,width=0.94\textwidth,max height=0.38\textheight]{Pope.pdf}
\end{sourceblock}
平均速度是这样，

\begin{sourceblock}
\includegraphics[page=176,pagebox=mediabox,trim={64.39bp 1721.78bp 64.39bp 224.97bp},clip,width=0.94\textwidth,max height=0.38\textheight]{Pope.pdf}
\end{sourceblock}
从这些定义我们有 f(0) 1 和 f(f1)

动量不足流量（每单位跨度）\textasciitilde{} ( x ) 定义为

\begin{sourceblock}
\includegraphics[page=176,pagebox=mediabox,trim={64.39bp 1502.81bp 64.39bp 400.95bp},clip,width=0.94\textwidth,max height=0.38\textheight]{Pope.pdf}
\end{sourceblock}
并且，在自相似区域，这是

\begin{sourceblock}
\includegraphics[page=176,pagebox=mediabox,trim={64.39bp 1329.83bp 64.39bp 560.93bp},clip,width=0.94\textwidth,max height=0.38\textheight]{Pope.pdf}
\end{sourceblock}
将动量定理应用于圆柱体和尾流（参见 Batchelor (1967)）表明动量赤字流量是守恒的

\textasciitilde{} ( x ) （与 x 无关）并且等于圆柱体上的阻力（每单位跨度）（另请参见练习 5.35）。因此，从等式可以看出。 (5.234) 在远尾 (Us/ U, O) 中，乘积 U , ( X ) \textasciitilde{} \textasciitilde{} / \textasciitilde{} 取决于 (Xin)d

在远尾流中，平均对流项 D ( U ) 减少为 U, 3 ( U )/ ax，因此边界层方程变为

\begin{sourceblock}
\includegraphics[page=176,pagebox=mediabox,trim={64.39bp 899.89bp 64.39bp 999.87bp},clip,width=0.94\textwidth,max height=0.38\textheight]{Pope.pdf}
\end{sourceblock}
( 5 ) / 就相似变量 f 和 g ( t ) (uv) U: 而言，该方程为

\begin{sourceblock}
\includegraphics[page=176,pagebox=mediabox,trim={64.39bp 763.90bp 64.39bp 1181.85bp},clip,width=0.94\textwidth,max height=0.38\textheight]{Pope.pdf}
\end{sourceblock}
其中传播参数是

\begin{sourceblock}
\includegraphics[page=176,pagebox=mediabox,trim={64.39bp 591.92bp 64.39bp 1314.83bp},clip,width=0.94\textwidth,max height=0.38\textheight]{Pope.pdf}
\end{sourceblock}
（参见练习 5.36）。方程（5.236）表明，自相似性决定了 S 是常数。这与 U 的恒定性一起， ( x ) \textasciitilde{} \textasciitilde{} / \textasciitilde{} i(mxp) l,i 表明 U,(x) 和 y1/2(x) 分别随 x 和 x 变化。请注意等式。 (5.236) 可以积分得出剪切应力、平均速度和铺展率之间的简单关系：

\begin{sourceblock}
\includegraphics[page=176,pagebox=mediabox,trim={64.39bp 281.96bp 64.39bp 1665.79bp},clip,width=0.94\textwidth,max height=0.38\textheight]{Pope.pdf}
\end{sourceblock}
湍流黏度与 U,(x)y 成比例，因此是独立的
{\small 1 2 ( \textasciitilde{} )，x 的下垂。假设湍流黏度恒定，即\par}
\begin{sourceblock}
\includegraphics[page=176,pagebox=mediabox,trim={64.39bp 119.39bp 64.39bp 1846.77bp},clip,width=0.94\textwidth,max height=0.38\textheight]{Pope.pdf}
\end{sourceblock}
\begin{sourceblock}
\includegraphics[page=177,pagebox=mediabox,trim={64.39bp 1356.83bp 64.39bp 109.44bp},clip,width=\textwidth,max height=0.72\textheight]{Pope.pdf}
\par\vspace{0.45\baselineskip}
\small 图 5.26。自相似平面尾流中的归一化速度亏损轮廓。实线，来自 Wygnanski 等人的实验数据。 （1986）；虚线，恒定湍流黏度解决方案，方程。 （5.240）。
\end{sourceblock}
方程的解(5.238) 是

\begin{sourceblock}
\includegraphics[page=177,pagebox=mediabox,trim={64.39bp 1067.87bp 64.39bp 880.89bp},clip,width=0.94\textwidth,max height=0.38\textheight]{Pope.pdf}
\end{sourceblock}
其中 a In 2 0.693（参见练习 5.37）。湍流雷诺数
Wygnanski、Champagne 和 Marasli (1986) 报道了具有启发意义的实验。除了圆柱体之外，这些研究人员还使用对称翼型和薄矩形板来产生飞机尾流。在每种情况下，我们发现流动（令人信服地）是自相似的，平均速度分布如图 5.26 所示。可以看出，恒定湍流黏度解与数据吻合得很好，除了边缘处（如喷射流的情况）。

然而，来自不同发生器的尾流似乎并不趋向于完全相同的自相似状态。铺展率参数为 SPlate -- 0+073、S汽缸 -- 0.083 和 Sairfoi = l 0.103。与式一致。 （5.238），剪应力分布表现出相同程度的差异。对于板，峰值轴向有效值速度为 (u2 ),1,/,2/U s 0.32，而机翼的速度为 0.41。正如 George (1989) 所讨论的，观察到的不同状态与自相似性完全一致。然而，它们表明，随着湍流流体在下游对流。它保留有关的信息
从边界层方程出发，证明

\begin{sourceblock}
\includegraphics[page=178,pagebox=mediabox,trim={64.39bp 1639.79bp 64.39bp 258.97bp},clip,width=0.94\textwidth,max height=0.38\textheight]{Pope.pdf}
\end{sourceblock}
飞机尾迹；因此表明动量赤字流动
速率守恒（方程（5.233））。

\begin{sourceblock}
\includegraphics[page=178,pagebox=mediabox,trim={64.39bp 1354.83bp 64.39bp 456.94bp},clip,width=0.94\textwidth,max height=0.38\textheight]{Pope.pdf}
\end{sourceblock}
其中扩展参数 S 由等式定义。 （5.237）。来自
等式。 (5.232) 表明

\begin{sourceblock}
\includegraphics[page=178,pagebox=mediabox,trim={64.39bp 1141.86bp 64.39bp 780.90bp},clip,width=0.94\textwidth,max height=0.38\textheight]{Pope.pdf}
\end{sourceblock}
因此

\begin{sourceblock}
\includegraphics[page=178,pagebox=mediabox,trim={64.39bp 989.88bp 64.39bp 929.88bp},clip,width=0.94\textwidth,max height=0.38\textheight]{Pope.pdf}
\end{sourceblock}
缩放剪应力为

\begin{sourceblock}
\includegraphics[page=178,pagebox=mediabox,trim={64.39bp 846.89bp 64.39bp 1079.86bp},clip,width=0.94\textwidth,max height=0.38\textheight]{Pope.pdf}
\end{sourceblock}
表明（近似）边界层方程

\begin{sourceblock}
\includegraphics[page=178,pagebox=mediabox,trim={64.39bp 680.91bp 64.39bp 1217.85bp},clip,width=0.94\textwidth,max height=0.38\textheight]{Pope.pdf}
\end{sourceblock}
可以写成

\begin{sourceblock}
\includegraphics[page=178,pagebox=mediabox,trim={64.39bp 561.93bp 64.39bp 1383.83bp},clip,width=0.94\textwidth,max height=0.38\textheight]{Pope.pdf}
\end{sourceblock}
表明湍流黏度假设（均匀
相当于

\begin{sourceblock}
\includegraphics[page=178,pagebox=mediabox,trim={64.39bp 387.95bp 64.39bp 1557.80bp},clip,width=0.94\textwidth,max height=0.38\textheight]{Pope.pdf}
\end{sourceblock}
代入方程式(5.249) 代入方程。 (5.248), 求解

\begin{sourceblock}
\includegraphics[page=178,pagebox=mediabox,trim={64.39bp 248.97bp 64.39bp 1693.79bp},clip,width=0.94\textwidth,max height=0.38\textheight]{Pope.pdf}
\end{sourceblock}
\begin{sourceblock}
\includegraphics[page=178,pagebox=mediabox,trim={64.39bp 119.39bp 64.39bp 1830.77bp},clip,width=0.94\textwidth,max height=0.38\textheight]{Pope.pdf}
\end{sourceblock}
证明 cr 在 2 中，因此

\begin{sourceblock}
\includegraphics[page=179,pagebox=mediabox,trim={64.39bp 1694.79bp 64.39bp 208.97bp},clip,width=0.94\textwidth,max height=0.38\textheight]{Pope.pdf}
\end{sourceblock}
\subsection*{5.4.4 轴对称尾流}\phantomsection\addcontentsline{toc}{subsection}{5.4.4 轴对称尾流}
轴对称尾流的分析与平面尾流的分析非常相似。然而，实验数据显示出显着的差异。

轴对称尾流在圆形物体后面形成球体、椭球体或圆盘，例如保持在均匀流中，沿 x 方向以速度 U 流动。流动在统计上是轴对称的，统计数据取决于 x 和 r，但与 0 无关。中心线速度赤字 U,(x) 和流动半宽 rlI2(x)a 以明显的方式定义。

正如平面尾流一样，只有当 UJU 趋于零时，自相似性才有可能，此时扩展参数 S (Uc/U,)drllz/dx 为常数。然而，对于该流动，等于物体阻力的动量赤字流量与 pUcUsr:12 成正比。因此，Us 随 x 变化，并且随着 x 雷诺数减小为
-2/3 SO x-'I3。湍流黏度在整个流动中是均匀的假设导致与平面尾流相同的平均速度赤字分布（方程（5.239）-（5.24 1））。

Uberoi 和 Freymuth (1970) 报道了在球体（直径为 d）的尾迹中进行的测量，雷诺数 Red U,d/v 为 8,600。经过发展距离 (xld 50) 后，在所检查的 xld 范围内 (50 c x ld < 150) 观察到平均速度和雷诺应力的自相似性。将测得的平均速度损失曲线与图 5.27 中的恒定湍流黏度解进行比较，并得出均方根曲线。速度如图 5.28 所示。应该观察到峰值
2 1/2 /Us 约为 0.9，远高于我们检查的 (u ) 其他流程中的值。相应地，扩散参数为 S 0.51 至少是平面尾流中观察到的扩散参数的五倍。

湍动能的平衡（图 5.29）也与其他流动有很大不同。主要项是来自上游的对流 (e JklJx)，耗散和横向传输各约为一半。相比之下，在高峰期，产量仅为对流的 20\% 和 15\%。对流占主导地位，且产生量相对较小，表明湍流受到上游条件的强烈影响。

\begin{sourceblock}
\includegraphics[page=180,pagebox=mediabox,trim={64.39bp 1169.85bp 64.39bp 109.44bp},clip,width=\textwidth,max height=0.72\textheight]{Pope.pdf}
\par\vspace{0.45\baselineskip}
\small 图 5.27。自相似轴对称尾流中的平均速度亏损剖面。符号，Uberoi 和 Freymuth (1970) 的实验数据；线，恒定湍流黏度解 f(£) = exp(-¢? In 2)。
\end{sourceblock}
\begin{sourceblock}
\includegraphics[page=180,pagebox=mediabox,trim={64.39bp 526.93bp 64.39bp 922.88bp},clip,width=\textwidth,max height=0.72\textheight]{Pope.pdf}
\par\vspace{0.45\baselineskip}
\small 图 5.28。均方根自相似轴对称尾流中的速度分布。 Uberoi和Freymuth (1970)的实验数据：x, (u2)!/2/U,; (v?)'/?/U,; 0, (w?)!/U,.
\end{sourceblock}
\begin{sourceblock}
\includegraphics[page=181,pagebox=mediabox,trim={64.39bp 1058.87bp 64.39bp 109.44bp},clip,width=\textwidth,max height=0.72\textheight]{Pope.pdf}
\par\vspace{0.45\baselineskip}
\small 图 5.29。自相似轴对称尾流中的湍动能收支。 Uberoi 和 Freymuth (1970) 的实验数据。
\end{sourceblock}
\begin{center}\small 表 5.3。各种物体背后轴对称尾流的扩散参数和湍流强度\end{center}
\begin{sourceblock}
\includegraphics[page=181,pagebox=mediabox,trim={64.39bp 119.39bp 64.39bp 1141.86bp},clip,width=\textwidth,max height=0.72\textheight]{Pope.pdf}
\end{sourceblock}
\begin{sourceblock}
\includegraphics[page=182,pagebox=mediabox,trim={64.39bp 1611.80bp 64.39bp 202.97bp},clip,width=0.94\textwidth,max height=0.38\textheight]{Pope.pdf}
\end{sourceblock}
表明动量亏缺流速

\begin{sourceblock}
\includegraphics[page=182,pagebox=mediabox,trim={64.39bp 1407.82bp 64.39bp 436.94bp},clip,width=0.94\textwidth,max height=0.38\textheight]{Pope.pdf}
\end{sourceblock}
对于自相似尾流，用 Us 和 f = 重新表达 M
(U, - (U))/LJS,其中 r/r1/2.H 表明

\begin{sourceblock}
\includegraphics[page=182,pagebox=mediabox,trim={64.39bp 1209.85bp 64.39bp 689.91bp},clip,width=0.94\textwidth,max height=0.38\textheight]{Pope.pdf}
\end{sourceblock}
其中传播参数是

\begin{sourceblock}
\includegraphics[page=182,pagebox=mediabox,trim={64.39bp 1052.87bp 64.39bp 849.89bp},clip,width=0.94\textwidth,max height=0.38\textheight]{Pope.pdf}
\end{sourceblock}
从等式开始（5.253），表明

\begin{sourceblock}
\includegraphics[page=182,pagebox=mediabox,trim={64.39bp 930.88bp 64.39bp 1014.87bp},clip,width=0.94\textwidth,max height=0.38\textheight]{Pope.pdf}
\end{sourceblock}
对于自相似轴对称尾迹，其中 g(5) (uv) u:。因此

表明

\begin{sourceblock}
\includegraphics[page=182,pagebox=mediabox,trim={64.39bp 765.90bp 64.39bp 1181.85bp},clip,width=0.94\textwidth,max height=0.38\textheight]{Pope.pdf}
\end{sourceblock}
表明，如果调用均匀湍流黏度假设
( e . , g CTfr)，方程的解。 (5.258) 是

\begin{sourceblock}
\includegraphics[page=182,pagebox=mediabox,trim={64.39bp 590.93bp 64.39bp 1356.83bp},clip,width=0.94\textwidth,max height=0.38\textheight]{Pope.pdf}
\end{sourceblock}
还有那个

\begin{sourceblock}
\includegraphics[page=182,pagebox=mediabox,trim={64.39bp 469.94bp 64.39bp 1477.81bp},clip,width=0.94\textwidth,max height=0.38\textheight]{Pope.pdf}
\end{sourceblock}
\subsection*{5.4.5 均匀剪切流}\phantomsection\addcontentsline{toc}{subsection}{5.4.5 均匀剪切流}
正如视觉观察所表明的那样，我们检查的自由剪切流本质上在统计上是不均匀的。在时间混合层的中心，平均剪切速率 ( U ) / a y 似乎相当均匀，但是

\begin{sourceblock}
\includegraphics[page=183,pagebox=mediabox,trim={64.39bp 1214.85bp 64.39bp 109.44bp},clip,width=\textwidth,max height=0.72\textheight]{Pope.pdf}
\par\vspace{0.45\baselineskip}
\small 图 5.30。均匀剪切流中平均速度剖面的草图。
\end{sourceblock}
动能不随时间变化，但生成量超过耗散量 40\%。产生的多余能量被向外输送，所涉及的输送过程本质上取决于统计不均匀性。与这些流动相反，研究均匀湍流的信息非常丰富，因为均匀湍流中不存在这种类型的传输过程。根据均匀湍流的定义，速度 u(x,t) 和压力 p’(x,t) 的脉动分量在统计上是均匀的。由此可见，施加的平均速度梯度@(U;)/0x;必须是统一的，尽管它们可能随时间而变化（参见练习 5.40 和 5.41）。在这里，我们检查均匀剪切流，其中单个施加的平均速度梯度 S = 是恒定的。

在风洞实验中可以很好地近似均匀剪切流。通过控制上游的流阻，可以产生平均速度分布如图 5.30 所示的湍流。 （平均流动完全在 x 方向上，即 (V) = (W) = 0，并且 (U) 仅在 y 方向上变化。）在流动开始时 (x/h = 0)，雷诺应力在垂直于流动方向上是均匀的，并且这种均匀性在下游持续存在。图 5.31 显示了 Tavoularis 和 Corrsin (1981) 测量的雷诺应力的轴向变化。尽管存在这种轴向变化，但在以平均速度 U 移动的框架中，湍流几乎是均匀的。还进行了均匀剪切流的直接数值模拟（例如，Rogallo (1981) 和 Rogers 和 Moin (1987)），其结果与

\begin{sourceblock}
\includegraphics[page=184,pagebox=mediabox,trim={64.39bp 1296.84bp 64.39bp 109.44bp},clip,width=\textwidth,max height=0.72\textheight]{Pope.pdf}
\par\vspace{0.45\baselineskip}
\small 图 5.31。 Tavoularis 和 Corrsin (1981) 的均匀剪切流实验中的雷诺应力与轴向距离的关系：©, 0, (v?)； A，（w？）。
\end{sourceblock}
\begin{sourceblock}
\includegraphics[page=184,pagebox=mediabox,trim={64.39bp 519.93bp 64.39bp 1388.82bp},clip,width=0.94\textwidth,max height=0.38\textheight]{Pope.pdf}
\end{sourceblock}
\begin{sourceblock}
\includegraphics[page=184,pagebox=mediabox,trim={64.39bp 387.95bp 64.39bp 1536.81bp},clip,width=0.94\textwidth,max height=0.38\textheight]{Pope.pdf}
\end{sourceblock}
\begin{sourceblock}
\includegraphics[page=184,pagebox=mediabox,trim={64.39bp 204.97bp 64.39bp 1694.79bp},clip,width=0.94\textwidth,max height=0.38\textheight]{Pope.pdf}
\par\vspace{0.45\baselineskip}
\small 表 5.4。来自 Tavoularis 和 Corrsin (1981) 实验以及 Rogers 和 Moin (1987) DNS 的均匀湍流剪切流统计
\end{sourceblock}
\begin{sourceblock}
\includegraphics[page=185,pagebox=mediabox,trim={64.39bp 973.88bp 64.39bp 262.97bp},clip,width=\textwidth,max height=0.72\textheight]{Pope.pdf}
\end{sourceblock}
通过从纳维方程中减去雷诺方程（方程（4.12））

斯托克斯方程（方程（2.35））表明，脉动速度
演变为

\begin{sourceblock}
\includegraphics[page=185,pagebox=mediabox,trim={64.39bp 738.91bp 64.39bp 1179.85bp},clip,width=0.94\textwidth,max height=0.38\textheight]{Pope.pdf}
\end{sourceblock}
表明，对于均匀湍流，

\begin{sourceblock}
\includegraphics[page=185,pagebox=mediabox,trim={64.39bp 577.93bp 64.39bp 1319.83bp},clip,width=0.94\textwidth,max height=0.38\textheight]{Pope.pdf}
\end{sourceblock}
\begin{sourceblock}
\includegraphics[page=185,pagebox=mediabox,trim={64.39bp 384.95bp 64.39bp 1554.80bp},clip,width=0.94\textwidth,max height=0.38\textheight]{Pope.pdf}
\end{sourceblock}
因此动能的演变为

\begin{sourceblock}
\includegraphics[page=185,pagebox=mediabox,trim={64.39bp 119.39bp 64.39bp 1835.77bp},clip,width=0.94\textwidth,max height=0.38\textheight]{Pope.pdf}
\end{sourceblock}
（请注意，这意味着同质的必要条件
湍流的本质是平均速度梯度是均匀的。）

通过对 xk 求导雷诺方程（方程 (4.12)

写为 (Uj)，表明在均匀湍流中，速度
梯度演化为

因此表明平均应变率

和旋转

进化为

（请注意，平均压力场的形式为 (p(x, t)) = A(t)

Bi(t)xi Cij(t)xixj，并且可以选择 Cij(t) 来产生任何
期望的演化 dSjk/dt。另一方面，进化论的发展
平均旋转率完全由 Sij 和 Rij 决定。）

湍流 在没有平均速度梯度的情况下，均匀湍流会衰减，因为没有产生 (p 0)。在风洞实验中，通过使均匀流（x 方向速度为 Uo）通过类似于图 5.32 所示的网格（其特征为网格间距 M），可以很好地近似衰减均匀湍流。在实验室框架中，流在统计上是静止的，并且（在流的中心）统计数据仅在 x 方向上变化。在以平均速度 Uo 运动的框架中，湍流是（充分近似）均匀的，并且随时间 (t x/Uo) 演变。

\subsection*{5.4.6 网格湍流}\phantomsection\addcontentsline{toc}{subsection}{5.4.6 网格湍流}
\begin{sourceblock}
\includegraphics[page=187,pagebox=mediabox,trim={64.39bp 1281.84bp 64.39bp 109.44bp},clip,width=\textwidth,max height=0.72\textheight]{Pope.pdf}
\par\vspace{0.45\baselineskip}
\small 图 5.32。由直径为 d、网格间距为 M 的条组成的湍流生成网格的草图。
\end{sourceblock}
\begin{sourceblock}
\includegraphics[page=187,pagebox=mediabox,trim={64.39bp 554.93bp 64.39bp 772.90bp},clip,width=\textwidth,max height=0.72\textheight]{Pope.pdf}
\par\vspace{0.45\baselineskip}
\small 图 5.33。网格湍流中雷诺应力的衰减：平方，\{u？)/U2；圆(v?)/U2；三角形 k/U3；线，与 (x/M)\textasciitilde{}'3 成比例。 （摘自 Comte-Bellot 和 Corrsin (1966)。）
\end{sourceblock}
（理想）实验表明 (v2) 和 (w2) 相等，并且所有
是剪应力为零。可以看出，rms。轴向速度 (u 2 ) 1/2 比横向均方根大 10\% （五）。 （Comte-Bellot 和 Corrsin (1966) 展示了对实验的修改，产生相等的法向应力，从而更好地近似均匀各向同性湍流的理想状态。）

法律，在实验室框架中，可以写成

\begin{sourceblock}
\includegraphics[page=188,pagebox=mediabox,trim={64.39bp 1720.78bp 64.39bp 207.97bp},clip,width=0.94\textwidth,max height=0.38\textheight]{Pope.pdf}
\end{sourceblock}
其中 xo 是虚拟原点。文献中报道了 1.15 至 1.45 之间的衰减指数 n 的值；但 Mohamed 和 LaRue (1990) 表明几乎所有数据都与 n = 1.3（且 xo = 0）一致。 （A 的值根据网格的几何形状和雷诺数而有很大变化。）

在移动坐标系中，幂律（方程（5.273））可以写成

\begin{sourceblock}
\includegraphics[page=188,pagebox=mediabox,trim={64.39bp 1319.83bp 64.39bp 576.93bp},clip,width=0.94\textwidth,max height=0.38\textheight]{Pope.pdf}
\end{sourceblock}
其中to是任意参考时间，ko是当时k的值。微分后我们得到

\begin{sourceblock}
\includegraphics[page=188,pagebox=mediabox,trim={64.39bp 1112.86bp 64.39bp 784.90bp},clip,width=0.94\textwidth,max height=0.38\textheight]{Pope.pdf}
\end{sourceblock}
现在，对于衰减均匀湍流，k 演化的精确方程（方程（5.261））简化为

\begin{sourceblock}
\includegraphics[page=188,pagebox=mediabox,trim={64.39bp 922.88bp 64.39bp 984.88bp},clip,width=0.94\textwidth,max height=0.38\textheight]{Pope.pdf}
\end{sourceblock}
因此，最后两个方程的比较表明，它也衰减为

\begin{sourceblock}
\includegraphics[page=188,pagebox=mediabox,trim={64.39bp 750.91bp 64.39bp 1121.86bp},clip,width=0.94\textwidth,max height=0.38\textheight]{Pope.pdf}
\end{sourceblock}
与 nko/to.其他量的衰减指数如下所示
EO 练习 5.42。

随着湍流减弱，雷诺数降低，最终粘度的影响占主导地位。这导致最后一个衰减周期，如第 6.3 节中讨论的那样，其中衰减指数为 n =（参见第 205 页的练习 6.10）。

从某种意义上说，网格湍流（作为均匀各向同性湍流的近似）是最基本的湍流，因此在实验和理论上都得到了广泛的研究。然而，从另一种意义上来说，它是病态的：与湍流剪切流相反，不存在湍流产生机制（网格的下游）。

CA ctive' 网格由一系列移动襟翼组成，可产生明显更高的湍流水平，从而产生更高的雷诺数，

\section*{5.5 进一步的观测}\phantomsection\addcontentsline{toc}{section}{5.5 进一步的观测}
\subsection*{5.5.1 守恒标量}\phantomsection\addcontentsline{toc}{subsection}{5.5.1 守恒标量}
对于网格湍流，给定 k 和 (方程 (5.274) 的衰减定律
和 (5.277)，并取 n 1.3，验证以下行为：

（请注意，雷诺数 k ' / ' \textasciitilde{} / v 减小。

L不应该被误解。并不是那波涛汹涌
动作变得更大、更慢。相反，更小、更快的运动
衰减得更快，留下更大、更慢的运动。）

5.5 进一步观察我们对自由剪切流的检查主要集中在平均速度场和雷诺应力上。在本节中，我们首先通过考虑其均值 (+) 将这些考虑因素扩展到守恒无源标量，
(～4')。其方差 (+'2) 和标量通量 然后我们检查一阶矩和二阶矩未描述的数量和现象。

守恒标量 守恒无源标量 +(x, t) 的行为已在本章讨论的所有自由剪切流中进行了实验研究。在射流中，射流流体的温度可能会稍微过高，或者化学成分与环境略有不同。可以方便地对所得标量场进行归一化，以便在射流中统一，在环境中为零。这种归一化标量有时被称为（特别是在燃烧文献中）作为混合分数。类似地，在混合层中，两个流中的 可以归一化为 0 和 1。

对于平面流，均值的边界层方程为

\begin{sourceblock}
\includegraphics[page=189,pagebox=mediabox,trim={64.39bp 218.97bp 64.39bp 1726.78bp},clip,width=0.94\textwidth,max height=0.38\textheight]{Pope.pdf}
\end{sourceblock}
该方程与轴向速度 (U) 的方程非常相似（方程 (5.56)），

\begin{sourceblock}
\includegraphics[page=190,pagebox=mediabox,trim={64.39bp 1203.85bp 64.39bp 109.44bp},clip,width=\textwidth,max height=0.72\textheight]{Pope.pdf}
\par\vspace{0.45\baselineskip}
\small 图 5.34。自相似平面尾流中的归一化平均速度赤字 f(€) 和标量 g() = (¢)/(@)y=o。符号（实心f，空心@）Fabris（1979）的实验数据：实线，f（€）= exp（Ce虚线，= exp（-f207 In2）其中oy = 0.7。
\end{sourceblock}
举个例子，我们已经看到动量流量 p (\textasciitilde{} )d'y 是守恒的，由此我们推断出其比例为 x-'/'。类似地，由等式得出。 (5.278) 标量流量 p(U)(\$) dy 也是守恒的，因此必须以与 (U) 相同的方式缩放，即 x。

- 1/ 2 同样的结论适用于所有自相似的自由剪切流：以与 x 相同的方式缩放。

的横向轮廓再次与 (U) 的横向轮廓相似，但在所有情况下都发现它们稍微更宽。例如，图 5.34 显示了 Fabris (1979) 在自相似平面尾流中测量的 (U) 的轮廓。恒定湍流黏度的假设导致
f ( 5 ) 图 5.34 所示的归一化速度剖面 exp(-5' ln2)（参见方程（5.240））。类似地，假设湍流扩散率恒定

\begin{sourceblock}
\includegraphics[page=190,pagebox=mediabox,trim={64.39bp 427.95bp 64.39bp 1508.81bp},clip,width=0.94\textwidth,max height=0.38\textheight]{Pope.pdf}
\end{sourceblock}
导致标准化标量轮廓

\begin{sourceblock}
\includegraphics[page=190,pagebox=mediabox,trim={64.39bp 297.96bp 64.39bp 1644.79bp},clip,width=0.94\textwidth,max height=0.38\textheight]{Pope.pdf}
\end{sourceblock}
（参见练习 5.43）。该轮廓，湍流普朗特数设置为
0.7，与数据吻合较好。 （如练习 5.43 所示，标量和速度剖面宽度的 CT 比与 a, -112 成正比）。

ee或

图 5.35，归一化均方根，圆射流中的标量波动。从实验
Panchapakesan 和 Lumley (1993b) 的数据。

扩散和湍流黏度假设仅在有限的意义上是成功的，即近似横向通量 (vd') 和 (uv)。在飞机尾迹中，|0(@)/dx|与 U,/U, 趋于零相比，它小得微乎其微。然而实验数据（例如 Fabris (1979)）显示 |(ud’)|与 |(v@’)| 相当。在均匀剪切流中也清楚地证明了 V(@) 和 -(u@') 之间缺乏对齐。 Tavoularis 和 Corrsin (1981) 在 y 方向施加平均标量梯度，但发现 |(u@’)|大于两倍|(v@’)|：和-(ud’)之间的角度为65°。

由于湍动能 k = \$(u+u) 表征了脉动速度场中的能量，因此标量方差 (@?) \textasciitilde{} 或 r.m.s (67)！ \textasciitilde{} 表征标量波动的水平。图 5.35 显示了均方根值Panchapakesan 和 Lumley（1993b）在空气中的氦射流中测量的。轮廓的形状和波动水平（高达 25\%）与均方根轴向速度（u？）！/？ （见图 5.7）。

标量方差的演化方程可以写成

\begin{sourceblock}
\includegraphics[page=191,pagebox=mediabox,trim={64.39bp 294.96bp 64.39bp 1605.80bp},clip,width=0.94\textwidth,max height=0.38\textheight]{Pope.pdf}
\end{sourceblock}
（参见练习 5.44），其中标量方差产生式为

\begin{sourceblock}
\includegraphics[page=191,pagebox=mediabox,trim={64.39bp 176.98bp 64.39bp 1767.78bp},clip,width=0.94\textwidth,max height=0.38\textheight]{Pope.pdf}
\end{sourceblock}
{\small 2 在 x/d = 50-120 的测量范围内，平均密度变化高达 15\%，\par}
\begin{sourceblock}
\includegraphics[page=192,pagebox=mediabox,trim={64.39bp 1151.86bp 64.39bp 109.44bp},clip,width=\textwidth,max height=0.72\textheight]{Pope.pdf}
\par\vspace{0.45\baselineskip}
\small 图 5.36。圆射流中的标量方差预算：方程中的项。 (5.281) 通过 和 归一化（来自 Panchapakesan 和 Lumley (1993b) 的实验数据。）
\end{sourceblock}
\begin{sourceblock}
\includegraphics[page=192,pagebox=mediabox,trim={64.39bp 597.93bp 64.39bp 936.88bp},clip,width=\textwidth,max height=0.72\textheight]{Pope.pdf}
\par\vspace{0.45\baselineskip}
\small 图 5.37。圆射流的标量与速度时间尺度之比。 （来自 Panchapakesan 和 Lumley (1993b) 的实验数据。）
\end{sourceblock}
是

标量耗散

通量为

\begin{sourceblock}
\includegraphics[page=192,pagebox=mediabox,trim={64.39bp 253.97bp 64.39bp 1686.79bp},clip,width=0.94\textwidth,max height=0.38\textheight]{Pope.pdf}
\end{sourceblock}
（有时，方程被写为，而 2 的因子在 和 的定义中被省略

\begin{sourceblock}
\includegraphics[page=193,pagebox=mediabox,trim={64.39bp 1195.85bp 64.39bp 109.44bp},clip,width=\textwidth,max height=0.72\textheight]{Pope.pdf}
\par\vspace{0.45\baselineskip}
\small 图 5.38。不同雷诺数下自相似圆射流轴上的归一化标量方差。三角形、空气喷射（Dowling 和 Dimotakis 的实验（1990））；圆圈、水射流（Miller 的实验（1991））。 （来自米勒（1991）。）
\end{sourceblock}
与图 5.16 的比较表明，项平衡与动能平衡非常相似。当 P4 达到峰值时，比率 P\$/E@ 为 0 .85。

量 k / 定义了速度波动的特征时间尺度，并且类似地 ( ) 定义了类似的标量时间尺度。圆射流中时标比 z/z+ 的轮廓如图 5.37 所示。大部分轮廓都在 1.5 的 15\% 以内。在许多其他剪切流中，该时间尺度比在 1.5-2.5 范围内（Beguier 等人，1978）。

分子扩散率的影响是什么 这个问题可以用雷诺数 Uf,/v 和普朗特数或施密特数 a -= v / r 来重新表述。实验中会遇到从 0.3（对于空气中的氦）到 1,500（对于水中的染料）的 a 值。图 5.38 提供了解决该问题的有价值的信息。它显示了标量方差 n，由不同雷诺数的自相似圆射流中心线上的平均标量标准化。对于空气射流 (a 1)，在研究的十年中，Re 没有影响。对于水射流 (a lo3)，标量方差随着 Re 的增加而减小，并且似乎渐近趋于与空气射流中的值相同。因此，这些数据与以下观点一致：在足够高的雷诺数（此处为 Re 30,000）下，速度和标量的均值和方差不受 Re 和 a 的影响。

从其定义（方程（5.283））可以明显看出，标量耗散本质上是
随着雷诺数的增加，归一化标量梯度 CV4 也随之增加。因此，不是像其定义中模糊暗示的那样使用 r / C 2 进行缩放，而是按照 U/L 进行缩放，独立于（在足够高的雷诺数下）。导致此行为的过程将在下一章中描述。

当 y 是标量的指定值时，点 x 满足

\begin{sourceblock}
\includegraphics[page=194,pagebox=mediabox,trim={64.39bp 1479.81bp 64.39bp 415.95bp},clip,width=0.94\textwidth,max height=0.38\textheight]{Pope.pdf}
\end{sourceblock}
定义等标量曲面。在高雷诺数下，此类表面在中间尺度范围内表现出分形性质，分形维数约为 2.36（参见例如 Sreenivasan (199 1)）。然而，仔细检查表明，表面几何形状明显偏离完美的分形（Frederiksen 等人，1997）。

对于自相似平面尾流，忽略分子扩散率
平均标量方程（方程（5.278））是

\begin{sourceblock}
\includegraphics[page=194,pagebox=mediabox,trim={64.39bp 966.88bp 64.39bp 928.88bp},clip,width=0.94\textwidth,max height=0.38\textheight]{Pope.pdf}
\end{sourceblock}
令 40(x) 表示对称平面上的值；并且，
其中 y / \textasciitilde{} \textasciitilde{}le/t \textasciitilde{}th,e 自相似标量轮廓为

\begin{sourceblock}
\includegraphics[page=194,pagebox=mediabox,trim={64.39bp 799.90bp 64.39bp 1147.86bp},clip,width=0.94\textwidth,max height=0.38\textheight]{Pope.pdf}
\end{sourceblock}
通过积分方程。 (5.286) 对于所有 y，证明 Uc40(x)\textasciitilde{}1/2(xi)s
守恒，因此

\begin{sourceblock}
\includegraphics[page=194,pagebox=mediabox,trim={64.39bp 544.93bp 64.39bp 1335.83bp},clip,width=0.94\textwidth,max height=0.38\textheight]{Pope.pdf}
\end{sourceblock}
\begin{sourceblock}
\includegraphics[page=194,pagebox=mediabox,trim={64.39bp 449.94bp 64.39bp 1475.81bp},clip,width=0.94\textwidth,max height=0.38\textheight]{Pope.pdf}
\end{sourceblock}
恒定湍流扩散率假设可以写成

\begin{sourceblock}
\includegraphics[page=194,pagebox=mediabox,trim={64.39bp 288.96bp 64.39bp 1601.80bp},clip,width=0.94\textwidth,max height=0.38\textheight]{Pope.pdf}
\end{sourceblock}
其中 aT 是湍流普朗特数。证明，用这个
假设，标量方程（Eq. (5.286)）变为

\begin{sourceblock}
\includegraphics[page=194,pagebox=mediabox,trim={64.39bp 119.39bp 64.39bp 1826.77bp},clip,width=0.94\textwidth,max height=0.38\textheight]{Pope.pdf}
\end{sourceblock}
证明这个方程的解是

\begin{sourceblock}
\includegraphics[page=195,pagebox=mediabox,trim={64.39bp 1724.78bp 64.39bp 222.97bp},clip,width=0.94\textwidth,max height=0.38\textheight]{Pope.pdf}
\end{sourceblock}
\begin{sourceblock}
\includegraphics[page=195,pagebox=mediabox,trim={64.39bp 1582.80bp 64.39bp 287.96bp},clip,width=0.94\textwidth,max height=0.38\textheight]{Pope.pdf}
\end{sourceblock}
使用练习 5.36 的结果表明它可以重写

\begin{sourceblock}
\includegraphics[page=195,pagebox=mediabox,trim={64.39bp 1408.82bp 64.39bp 483.94bp},clip,width=0.94\textwidth,max height=0.38\textheight]{Pope.pdf}
\end{sourceblock}
令 为归一化标量轮廓的半宽 (ibe*,
1）。表明 和 相关的是

\begin{sourceblock}
\includegraphics[page=195,pagebox=mediabox,trim={64.39bp 1231.84bp 64.39bp 705.91bp},clip,width=0.94\textwidth,max height=0.38\textheight]{Pope.pdf}
\end{sourceblock}
根据雷诺分解 U ( U ) u 和

\begin{sourceblock}
\includegraphics[page=195,pagebox=mediabox,trim={64.39bp 1073.86bp 64.39bp 819.90bp},clip,width=0.94\textwidth,max height=0.38\textheight]{Pope.pdf}
\end{sourceblock}
因此，根据 \$(x, t) 的守恒方程（方程 (4.38)），显示
标量涨落的演化过程为

\begin{sourceblock}
\includegraphics[page=195,pagebox=mediabox,trim={64.39bp 759.90bp 64.39bp 1033.87bp},clip,width=0.94\textwidth,max height=0.38\textheight]{Pope.pdf}
\end{sourceblock}
通过乘以方程。 (5.297) 除以 24' 并取平均值，表明
标量方差根据方程演变。 （5.281）。

\subsection*{5.5.2 间歇性}\phantomsection\addcontentsline{toc}{subsection}{5.5.2 间歇性}
瞬间自由剪切流的目视观察。 （例如，第 4 页和第 5 页的图 1.1 和 1.2）表明湍流和环境流体之间存在尖锐（但高度不规则）的界面。从 Corrsin (1943) 的实验开始，许多实验都证实了这一图景：存在一个高度扭曲的移动表面 - 称为粘性超级层 - 将湍流和非湍流区域分开。湍流区域的特点是涡度大：均方根涡度 co' 与应变率一样，与柯尔莫哥洛夫尺度成比例，即 co' l/z、(\&/G)Re

\begin{sourceblock}
\includegraphics[page=196,pagebox=mediabox,trim={64.39bp 1414.82bp 64.39bp 109.44bp},clip,width=\textwidth,max height=0.72\textheight]{Pope.pdf}
\par\vspace{0.45\baselineskip}
\small 图 5.39。自由剪切流中间歇函数与时间的关系图（a）在无旋转非湍流环境中，（b）在间歇区域的外部，（c）在间歇区域的内部，以及（d）靠近流动中心。
\end{sourceblock}
朝向自由剪切流的边缘，运动有时是湍流，有时是非湍流，那里的流动是间歇性的。

为了避免误解，需要强调的是，整个粘性超级层不存在不连续性。从数学意义上讲，涡度和所有其他场在整个层中平滑变化：但与流宽度 6 相比，该层非常薄。

定量描述间歇性的起点是指示函数或间歇函数 I(x, t)。在湍流中定义为 I = 1，在非湍流中定义为 I 0。在操作上，它可以根据 Heaviside 函数获得：

\begin{sourceblock}
\includegraphics[page=196,pagebox=mediabox,trim={64.39bp 738.91bp 64.39bp 1208.85bp},clip,width=0.94\textwidth,max height=0.38\textheight]{Pope.pdf}
\end{sourceblock}
其中 是一个小的正阈值。图5.39显示了时间的草图
自由剪切流中不同位置的 I(t) 的 atresh 级数。

间歇因子 y (x, t) 是 (x, t) 处的流动为湍流的概率：

\begin{sourceblock}
\includegraphics[page=196,pagebox=mediabox,trim={64.39bp 474.94bp 64.39bp 1468.81bp},clip,width=0.94\textwidth,max height=0.38\textheight]{Pope.pdf}
\end{sourceblock}
对于所有自由剪切流，实验表明 y 的轮廓变得自相似。作为示例，图 5.40 显示了在平面尾流中测量的 y 轮廓。在其他流中，轮廓形状相似，但在射流和混合层中，间歇区域（例如 0.1 y 0.9）占据流宽度的较小部分。

间歇函数可用于获取条件统计数据。我们以标量 \$(x, t) 为例。作为样本空间变量，

\begin{sourceblock}
\includegraphics[page=197,pagebox=mediabox,trim={64.39bp 1245.84bp 64.39bp 109.44bp},clip,width=\textwidth,max height=0.72\textheight]{Pope.pdf}
\par\vspace{0.45\baselineskip}
\small 图 5.40。自相似平面尾流中间歇因子的分布。显示平均标量轮廓以供比较：yy 是半宽。 （来自 LaRue 和 Libby (1974, 1976) 的实验数据。）
\end{sourceblock}
湍流的贡献 y (x, t) f ;x , t)，以及
T(I/-J 非湍流；

\begin{sourceblock}
\includegraphics[page=197,pagebox=mediabox,trim={64.39bp 942.88bp 64.39bp 1000.87bp},clip,width=0.94\textwidth,max height=0.38\textheight]{Pope.pdf}
\end{sourceblock}
因此，f T 和 f N 分别是 \$ 以湍流 (I = I) 和非湍流 (I O) 为条件的 PDF。

在非湍流区域，正如涡度基本上为零一样，标量也基本上等于自由流值。事实上，该条件经常被用作替代指标
- \$thresh 湍流。因此，在近似范围内，非湍流 PDF 为

\begin{sourceblock}
\includegraphics[page=197,pagebox=mediabox,trim={64.39bp 573.93bp 64.39bp 1361.83bp},clip,width=0.94\textwidth,max height=0.38\textheight]{Pope.pdf}
\end{sourceblock}
\$ 的湍流均值和方差定义为

\begin{sourceblock}
\includegraphics[page=197,pagebox=mediabox,trim={64.39bp 383.95bp 64.39bp 1516.81bp},clip,width=0.94\textwidth,max height=0.38\textheight]{Pope.pdf}
\end{sourceblock}
\begin{sourceblock}
\includegraphics[page=197,pagebox=mediabox,trim={64.39bp 248.97bp 64.39bp 1651.79bp},clip,width=0.94\textwidth,max height=0.38\textheight]{Pope.pdf}
\end{sourceblock}
非湍流力矩 和 的定义类似，并且如果
(\$)N (\$',)2 由等式给出。 (5.302) 它们是 \$ 和 0。

\begin{sourceblock}
\includegraphics[page=198,pagebox=mediabox,trim={64.39bp 1398.82bp 64.39bp 109.44bp},clip,width=\textwidth,max height=0.72\textheight]{Pope.pdf}
\par\vspace{0.45\baselineskip}
\small 图 5.41。归一化无条件和湍流均值 (a) 和均方根的比较(b) 自相似平面尾流中的标量分布。来自 LaRue 和 Libby (1974) 的实验数据。
\end{sourceblock}
条件矩由

\begin{sourceblock}
\includegraphics[page=198,pagebox=mediabox,trim={64.39bp 1140.86bp 64.39bp 804.90bp},clip,width=0.94\textwidth,max height=0.38\textheight]{Pope.pdf}
\end{sourceblock}
\begin{sourceblock}
\includegraphics[page=198,pagebox=mediabox,trim={64.39bp 1056.87bp 64.39bp 885.89bp},clip,width=0.94\textwidth,max height=0.38\textheight]{Pope.pdf}
\end{sourceblock}
请注意，无条件方差包含条件均值之间差异的贡献。

对于湍流尾流， 、(\$)* 和 \$ 均为零。动荡的意思
和均方根\$ 与无条件对应项进行比较
为什么湍流/非=湍流界面如此尖锐？粘性超级层的性质和行为是什么？非湍流区速度脉动有何特点？

所有湍流剪切流的一个特征是它们夹带流体。例如，在圆射流中，质量流量随轴向距离 x 线性增加。由于间歇因子是自相似的，因此湍流流体的质量流量也随着x线性增加。因此，在 x l d 100 处，80\% 的湍流运动流体被夹带在 x l d 20 和 x l d 100 之间。夹带的流体源自环境，其中 w 和 \$ 为零。只有通过分子扩散，w 和 \$ 才能偏离零3，并且在湍流的高雷诺数下，

图 5.42。显示粘性超级层的湍流圆射流草图，以及
流体粒子从静态环境中的点 0 到超级层 E 的路径。

仅当梯度非常陡时，分子扩散才会产生显着影响。因此，在粘性超级层中，lo1 和 从零急剧上升，并且
IV6I)层的厚度(与 成反比)小。

粘性超级层的一些特性是从实验中已知的（例如，LaRue 和 Libby (1976)），但有些特性则不然。大规模的湍流运动对超级层进行对流，使得其横向位置的 PDF 可以非常准确地由高斯分布给出。相应地，y 的轮廓接近误差函数（见图 5.40）。湍流运动还会使超级层变形，使其随机波纹化，并且在中间的长度尺度范围内，它近似分形，分形维数为 2.36（Sreenivasan 等，1989）。由于质量流量
等人。湍流流体的速率随着下游距离的增加而增加，因此超级层相对于流体传播到无旋转环境中是超级层面积的乘积，并且该传播速度给出了体积夹带率，其比例为\textasciitilde{} \$(in6dep\textasciitilde{}end of Re)。

如果（如 Corrsin 和 Kistler (1954) 首次提出的那样）我们假设超级层相对于流体的传播速度 u, 与柯尔莫哥洛夫速度 u, U,Re (Eqw(5 .154)) 成正比，就会出现一致的情况。然后，为了使整体夹带率与 Re 无关，超级层的面积比例为 6 2 \textasciitilde{} e ' / 4T.h e 根据粗略模型，则表明超级层厚度与 Kolmogorov 比例 q 6Re 成比例（方程（5.152））。

图 5.42 是湍流圆射流的草图，显示了

\begin{sourceblock}
\includegraphics[page=200,pagebox=mediabox,trim={64.39bp 1375.83bp 64.39bp 109.44bp},clip,width=\textwidth,max height=0.72\textheight]{Pope.pdf}
\par\vspace{0.45\baselineskip}
\small 图 5.43。间歇因子 7、无条件平均轴向速度 (U) 和湍流 (U) 的剖面；和自相似混合层中的非湍流 (U)y 条件平均速度。来自 Wygnanski 和 Fiedler (1970) 的实验数据。
\end{sourceblock}
超级层。让 s 表示沿测量的流体粒子路径的弧长
r v2\$。从 E 向 0 方向。标量方程 D\$/Dt = 沿着这条路径通过常微分方程粗略地近似

\begin{sourceblock}
\includegraphics[page=200,pagebox=mediabox,trim={64.39bp 957.88bp 64.39bp 945.88bp},clip,width=0.94\textwidth,max height=0.38\textheight]{Pope.pdf}
\end{sourceblock}
适当的边界条件为 0（在环境中）和

\$(m) = \$(0) 的秩序统一，在超级层中。 Eqw( 5.307) 的解为

\begin{sourceblock}
\includegraphics[page=200,pagebox=mediabox,trim={64.39bp 777.90bp 64.39bp 1167.85bp},clip,width=0.94\textwidth,max height=0.38\textheight]{Pope.pdf}
\end{sourceblock}
其中长度尺度 A 与 v i \textasciitilde{} c o u ss 上层的估计厚度成正比，为 A T/ue。相对于局部流宽6，这

\begin{sourceblock}
\includegraphics[page=200,pagebox=mediabox,trim={64.39bp 513.93bp 64.39bp 1328.83bp},clip,width=0.94\textwidth,max height=0.38\textheight]{Pope.pdf}
\end{sourceblock}
其中 c 是常数，雷诺数为 U,6/v。 （回想一下，u,/U，尺度为 Re 1/ 4 ）因此，根据该模型，超级层厚度与 Kolmogorov 尺度成比例，为 6Re 实验证据（例如，LaRue 和 Libby (1976)）与此图一致，但不足以被视为令人信服的确认。

我们重点关注涡度 w 和标量 \$，因为在非湍流区域，它们的值只能通过分子效应来改变
rv2\$), (v v 2 w 和 这些可以忽略不计。另一方面，速度，

\begin{sourceblock}
\includegraphics[page=201,pagebox=mediabox,trim={64.39bp 1350.83bp 64.39bp 109.44bp},clip,width=\textwidth,max height=0.72\textheight]{Pope.pdf}
\end{sourceblock}
\subsection*{5.5.3 PDF 和更高矩}\phantomsection\addcontentsline{toc}{subsection}{5.5.3 PDF 和更高矩}
\begin{center}\small 图 5.44。无条件 (v*)、湍流和非湍流 (\#4) 的剖面？自相似混合层中轴向速度的方差。来自 Wygnanski 和 Fiedler (1970) 的实验数据。\end{center}
\begin{sourceblock}
\includegraphics[page=201,pagebox=mediabox,trim={64.39bp 454.94bp 64.39bp 1437.82bp},clip,width=0.94\textwidth,max height=0.38\textheight]{Pope.pdf}
\end{sourceblock}
\begin{sourceblock}
\includegraphics[page=202,pagebox=mediabox,trim={64.39bp 1161.85bp 64.39bp 109.44bp},clip,width=\textwidth,max height=0.72\textheight]{Pope.pdf}
\par\vspace{0.45\baselineskip}
\small 图 5.45。均匀剪切流中 (a) u、(b) v、(c) w 和 (d) @ 的标准化 PDF。虚线是标准化高斯曲线。 （来自 Tavoularis 和 Corrsin (1981)。）
\end{sourceblock}
速度和标量的联合 PDF 是联合法线；而在自由剪切流中，PDF 不是高斯分布的。图 5.45 显示了 Tavoularis 和 Corrsin (1981) 在均匀剪切流中测量的 u、v、w 的标准化边际 PDF。与标准化高斯分布没有明显的偏离。对于相同的流动，图 5.46 显示了速度-速度和速度-标量联合 PDF。同样，这些是通过联合正态分布准确描述的。这是一个重要且有价值的观察。

在自由剪切流中，情况则完全不同。图 5.47 显示了时间混合层中的标量 PDF f + ( y ; c )。中的标量值
1;并且，由于有界性，两个流都是属性（参见第 2.6 节），因此处处都位于零和单位之间。

+(x, t) 因此 f + ( y ; 对于 y 0 和 y 1 为零。

从图 5.47 中可以看出，在该层的中心有一个广泛的、大致钟形的分布，跨越了整个值范围。当测量位置向高速流移动时，PDF 会移动到更高的 y 值，并在上限 y 1 处形成幅度不断增加的尖峰。

在喷流（例如，Dahm 和 Dimetakis (1990)）和尾流（例如，LaRue 和 Libby (1974)）中发现了定性相似的 PDF 形状。 Dowling 和 Dimotakis (1990) 在圆射流中的测量特别有价值

\begin{sourceblock}
\includegraphics[page=203,pagebox=mediabox,trim={64.39bp 639.92bp 64.39bp 198.78bp},clip,width=\textwidth,max height=0.72\textheight]{Pope.pdf}
\par\vspace{0.45\baselineskip}
\small 图 5.46。在均匀剪切流中测量的标准化变量的联合 PDF 的等高线图：(a\} u 和 v，(b) @ 和 v，(c) u 和 @。等高线值为 0.15、0.10、0.05 和 0.01。虚线是具有相同相关系数的联合正态分布的相应等高线。（来自 Tavoularis 和 Corrsin (1981)。）
\end{sourceblock}
在自由剪切流的边缘，PDF 形状的变化强烈影响更高的力矩。图 5.48 显示了偏度
和峰度

\begin{sourceblock}
\includegraphics[page=204,pagebox=mediabox,trim={64.39bp 1374.83bp 64.39bp 109.44bp},clip,width=\textwidth,max height=0.72\textheight]{Pope.pdf}
\par\vspace{0.45\baselineskip}
\small 图 5.47。自相似时间混合层中不同横向位置的守恒无源标量的 PDF，来自 Rogers 和 Moser (1994) 的直接数值模拟。
\end{sourceblock}
\begin{sourceblock}
\includegraphics[page=204,pagebox=mediabox,trim={64.39bp 705.91bp 64.39bp 719.91bp},clip,width=\textwidth,max height=0.72\textheight]{Pope.pdf}
\par\vspace{0.45\baselineskip}
\small 图 5.48。无条件（Sy 和 Ky）和条件湍流（Sgr 和
\end{sourceblock}
KIT）自相似平面中守恒无源标量的偏度和峰度
醒来。来自 LaRue 和 Libby (1974) 的实验数据。

离高斯值不太远（S\# = O，K\# = 3），但它们在流的边缘增加到更大的值。图 5.48 中还显示了在湍流区域中通过条件平均获得的值 和。 \textasciitilde{} 在整个流程中，它主要更接近高斯值。

我们现在转向自由剪切流中轴向速度 U(x, t) 的 PDF f。图 5.49 显示了时间混合层中各个跨流位置的 f,(V)。在层的中心可以观察到熟悉的钟形曲线；但在边缘处，PDF 明显偏向于层内的速度。与标量不同，速度不受限制

\begin{sourceblock}
\includegraphics[page=205,pagebox=mediabox,trim={64.39bp 1254.84bp 64.39bp 109.44bp},clip,width=\textwidth,max height=0.72\textheight]{Pope.pdf}
\par\vspace{0.45\baselineskip}
\small 图 5.49。时间混合层中轴向速度的 PDF。距层中心的距离为 € = y/d。虚线对应于自由流速度U i。来自 Rogers 和 Moser (1994) 的 DNS 数据。
\end{sourceblock}
\begin{sourceblock}
\includegraphics[page=205,pagebox=mediabox,trim={64.39bp 606.92bp 64.39bp 830.90bp},clip,width=\textwidth,max height=0.72\textheight]{Pope.pdf}
\par\vspace{0.45\baselineskip}
\small 图 5.50。自相似圆射流中的偏斜度（实线\}和峰度（虚线）轮廓。来自 Wygnanski 和 Fiedler (1969) 的实验数据。
\end{sourceblock}
在层的边缘，高于自由流速度的速度很有可能出现。

与标量的情况一样，轴向速度的偏度 S 和峰度 K 朝着流动边缘显着偏离高斯值（0 和 3）。例如，图 5.50 显示了 Wygnanski 和 Fiedler (1969) 在自相似圆射流中测量的值。

总之：在均匀湍流中，速度标量联合 PDF 是联合法向的。在自由剪切流的中心，PDF 呈钟形，

\begin{sourceblock}
\includegraphics[page=206,pagebox=mediabox,trim={64.39bp 1495.81bp 64.39bp 109.44bp},clip,width=\textwidth,max height=0.72\textheight]{Pope.pdf}
\par\vspace{0.45\baselineskip}
\small 图 5.51。平面混合层流动的可视化。 8 atm 压力下氦（上）Up = 10.1 m s\textasciitilde{}' 和氮（下）U, = 3.8 m 之间混合层的火花阴影图。 （摘自 Brown 和 Roshko (1974)。）
\end{sourceblock}
在朝向流动边缘的间歇区域中明显。在间歇区域中，标量的 PDF 在边界处有一个尖峰，而速度的 PDF 是无界的。

\subsection*{5.5.4 大尺度湍流运动}\phantomsection\addcontentsline{toc}{subsection}{5.5.4 大尺度湍流运动}
尽管它们很重要，但我们研究过的单点统计数据（￠，，（U）等）仅提供了对湍流的非常有限的描述。两点速度相关性的测量（例如图 5.13）表明，速度在与流宽度相当的距离上具有明显的相关性。然而，即使这些两点统计数据也无法提供有关引起这些远程相关性的瞬时流场结构的信息。

流动可视化提供了有关大规模湍流运动的一些信息。混合层和轴对称尾流的示例如图 2 和 3 所示。 5.51 和 5.52。混合层中可见的主要特征是大型“滚筒”。发现这些区域是展向涡度集中的区域，在相当大的展向距离上是相干的。当滚筒向下游对流时，它们的尺寸和间距增大，因此数量减少。一个滚轮可以在“配对”过程中与相邻的滚轮合并，也可以被撕裂，其涡流被相邻的滚轮吸收。

已经在其他自由剪切流中研究了大尺度运动。由于它们的大小与流的宽度有关，因此不可避免地受到流的几何形状和边界条件的强烈影响，因此在不同的流中它们是不同的。稳定性理论在某种程度上

\begin{sourceblock}
\includegraphics[page=207,pagebox=mediabox,trim={64.39bp 798.90bp 64.39bp 109.44bp},clip,width=\textwidth,max height=0.72\textheight]{Pope.pdf}
\par\vspace{0.45\baselineskip}
\small 图 5.52。统计轴对称尾流的可视化：(a) 50\% 遮挡屏幕、(b) 60\%、(c\} 85\%、(d\} 100\%（即圆盘）。尾流的动量厚度为 0；烟线位于 x/@ = 10 和 x/@ = 85 处。（来自 Cannon 等人 (1993)，经 Springer-Verlag 许可。）
\end{sourceblock}
支持以下观点：这些动议是由各种流动所固有的不稳定性引起的。

有几种现象可以用大规模运动来合理地解释。作为一个明显的例子，Oster 和 Wygnanski（1982）在有和没有小振幅强迫的情况下对湍流混合层进行了实验。施力在于使分流板后缘处的襟翼摆动。皮瓣运动的幅度（通常
mm）非常小，因此在混合层的初始区域（x < 100 mm），对平均速度和雷诺应力分布的影响是不可察觉的。然而，如图 5.53 所示，更下游有

\begin{sourceblock}
\includegraphics[page=208,pagebox=mediabox,trim={64.39bp 1360.83bp 64.39bp 109.44bp},clip,width=\textwidth,max height=0.72\textheight]{Pope.pdf}
\par\vspace{0.45\baselineskip}
\small 图 5.53。混合层的厚度 8，相对于轴向距离 x，对于各种强迫频率 f：-，0 Hz，x，60 Hz； A，50赫兹； 0.40赫兹； a、30赫兹。 （来自 Oster 和 Wygnanski (1982)。）
\end{sourceblock}
在 60 Hz 下受力的层开始以几乎是非受力层的两倍的速度扩散。然而，随后（在 x \textasciitilde{} 900 毫米处），该层停止生长，并且实际上在恢复生长之前（在 x \textasciitilde{} 1,500 毫米处）略有收缩。在较低的强迫频率下，会发生相同的现象，但在下游距离较大时，第 6/U 层的特征时间尺度较大。在层停止生长的区域中，发现剪切应力 (uv) 改变符号，因此产生 P 为负：能量从湍流中提取并返回到平均流。

显然，湍流黏度假说无法解释这些观察结果，尤其是负产量 P。然而，出现了将强迫、流动中固有的不稳定性和大规模湍流运动联系起来的合理图景。预计频率 f 处的力会激发与 U,/f 成比例的长度模式 £。靠近分离板，层厚度6比@小。当6向£增加时，激发更大规模的运动，并且促进辊的配对或合并，从而增加运动的规模和层的宽度。实现了共振，从而大规模运动锁定了力的规模。随后导致更大尺度的配对或合并受到抑制，因此该层停止生长。在这个区域中，大规模运动的结构被发现与稳定性理论预测的模式非常相似。

对于非强制混合层，铺展参数 S 因实验而异，从 S \textasciitilde{} 0.06 到 S = 0.11 (Dimotakis 1991)。

混合层到一系列小的、不可控的流动扰动，这取决于特定的设备。

在钝体（例如圆柱体和球体）尾流的近场中，存在具有优选频率的大规模运动，最明显的例子是圆柱体的涡旋脱落。当然，这些运动取决于物体的几何形状，对于屏幕、球体和圆盘来说它们是不同的。从图 5.52 可以明显看出，大尺度运动的差异持续到远尾流中。同样，这些运动与基本不稳定性有关，并被认为是造成观察到的传播率之间巨大差异的原因（Cannon 和 Champagne (1991)，见表 5.3）。

从大型结构的研究中至少可以学到两个教训。首先，湍流和湍流可以表现出比湍流黏度假说所承认的更丰富的行为范围。第二个是湍流过程在空间和时间上是非局部的：湍流具有很长的记忆性，并且它在某一点的行为可能会受到远离该点的流动的强烈影响。

\chapter{湍流运动的尺度}
在检查自由剪切流时，我们观察到湍流运动的大小范围从流的宽度 6 到更小的尺度，随着雷诺数的增加，其逐渐变小（相对于 6）。我们还看到了雷诺应力中湍动能和各向异性的重要性。在本章中，我们考虑能量和各向异性如何在不同运动尺度之间分布；我们研究在这些尺度上发生的不同物理过程。

本章中两个重复的主题是能量级联和科尔诺戈洛夫假设。简而言之，能量级联的思想（由理查森（Richardson，1922）提出）是动能以最大的运动尺度进入湍流（通过产生机制）。然后，这种能量（通过无粘性过程）转移到越来越小的尺度，直到在最小的尺度上，能量通过粘性作用耗散。 Kolmogorov (1941b) 添加并量化了这张图片。他特别指出，最小的湍流尺度就是那些现在以他的名字命名的湍流。

\section*{6.1 能量级联与柯尔莫哥洛夫假设}\phantomsection\addcontentsline{toc}{section}{6.1 能量级联与柯尔莫哥洛夫假设}
更详细地描述。然后检查区分不同运动尺度的各种统计数据；即结构函数（第 6.2 节）、两点相关性（第 6.3 节）和谱（第 6.5 节）。作为频谱讨论的前奏，在第 6.4 节中，湍流速度场被表示为傅里叶模态之和，并根据 Navier-S tokes 方程推导了这些模态的演化。其余部分给出了能量级联的频谱视图（第 6.6 节），并讨论了柯尔莫哥洛夫假设的局限性（第 6.7 节）。

6.1 能量级联和柯尔莫哥洛夫假设 我们考虑具有特征速度 U 和长度尺度的高雷诺数下的完全湍流 需要强调的是，雷诺数

\subsection*{6.1.1 能量级联}\phantomsection\addcontentsline{toc}{subsection}{6.1.1 能量级联}
Re UL/v 数较大；事实上，如果考虑到非常高的雷诺数，这些概念是最容易理解的。

6.1.I 级联
能量理查森能量级联观点中的第一个概念是湍流可以被认为是由不同尺寸组成的。

涡流 大小的涡流具有特征速度 u([) 和时间尺度 [/u([)。 “涡流”无法精确定义，但它被认为是一种湍流运动，位于大小为 t 的区域内，在该区域上至少具有适度的相干性。被大涡流占据的区域也可能包含较小的涡流。

最大尺寸范围内的涡流的特征是长度尺度低高与流动尺度L相当，并且其特征速度

\#(lois) 以均方根为单位湍流强度 u' (fk)llZw hich uo -= 与 U 相当。因此，这些涡流的雷诺数 Reo uolo/v 很大（ e，与 Re 相当），因此粘度的直接影响很小，可以忽略不计。

理查森的观点是，大涡流不稳定并会破裂，将能量转移到较小的涡流上。这些较小的涡流经历类似的破碎过程，并将其能量转移到更小的涡流。这种能量级联，其中能量转移到越来越小的涡流，一直持续到雷诺数 Re([) u([)l/v 足够小，使涡流运动稳定，并且分子粘度有效地耗散动能。理查森（Richardson，1922）如此简洁地总结了这一问题：

大螺纹有小螺纹，
以速度为食

小螺纹有更小的螺纹，
依此类推粘度

（在分子意义上）。

这张图很重要的一个原因是它将耗散置于一系列过程的末尾。耗散的确定，
因此，通过序列中的第一个过程，即来自最大涡流的能量转移。这些涡流的能量为 u 级；和时间尺度
[o/uo，因此能量转移速率可以按比例缩放为 TO =
ui/[o. u;/zo 因此，与实验观察结果一致
ui/to，在自由剪切流中，这张级联图片表明缩放比例为

\subsection*{6.1.2 柯尔莫哥洛夫假设}\phantomsection\addcontentsline{toc}{subsection}{6.1.2 柯尔莫哥洛夫假设}
几个基本问题仍未得到解答。负责耗散能量的最小涡流的大小是多少？当 减少时，特征速度和时间尺度 u(￠) 会增加、减少还是保持不变？ （雷诺数 u(£)£/v 与 @ 的假设减小不足以确定这些趋势。）

科尔莫哥洛夫 (1941b) 提出的理论回答了这些问题以及更多问题！以三个假设的形式表述。该理论（柯尔莫哥洛夫用来激发假设）的一个结果是，速度和时间尺度 u(£) 都随着减小而减小。

第一个假设涉及小尺度运动的各向同性。一般来说，大涡流是各向异性的，并且受到流动边界条件的影响。科尔莫哥洛夫认为，大尺度的方向偏差在混沌尺度减小过程中消失，能量通过该过程依次转移到越来越小的涡流中。因此（近似表述）：

柯尔莫哥洛夫局部各向同性假说。在足够高的雷诺数下
数，小尺度湍流运动 (∫ < ) 在统计上是各向同性的。

（术语“局部各向同性”仅在小尺度上表示各向同性，在第 6.1.4 节中进行了更精确的定义。）引入长度尺度（例如 =）作为各向异性大涡流 (€ > 和各向同性小涡流 (€ < \&) 之间的分界线是有用的。（\&; 以及下面介绍的其他尺度的规范的理由在第 6.5 节中提供。）

48 当能量沿着级联传递时，大尺度的方向信息就会丢失，柯尔莫哥洛夫认为，有关大涡流几何形状的所有信息（由平均流场和边界条件决定）也会丢失。因此，小尺度运动的统计数据在某种意义上是普遍的-在每个高雷诺数湍流中都是相似的。

这种本质上普遍的状态取决于什么参数？在能量级联中（对于 £ <，两个主要过程是能量向连续较小尺度的转移和粘性耗散。那么，一个合理的假设是，重要的参数是小尺度从大尺度接收能量的速率（我们用
和运动粘度 v。正如我们将看到的，耗散率 ¢
！该论文的英文译本以 Kolmogorov (1991) 的形式转载于《Kolmogorov》的特刊中。

为纪念原版出版物五十周年而出版的《皇家学会学报》。

由能量转移速率 TE1s, o 决定，这两个速率几乎相等，即， 因此，小尺度的统计普遍状态由 v 决定的假设，大尺度的能量转移速率 TEcI a n 可表述为：

科尔莫哥洛夫的第一个相似性假设。在每一个足够的湍流中
高雷诺数，小尺度运动的统计（！

具有由 v 和唯一确定的通用形式
< tEI 尺寸范围 l 称为通用平衡范围。在
在t / u ( t )这个范围内，时间尺度比to/uo小，因此小涡流能够快速适应以与大涡流所施加的能量传递率TEIi保持动态平衡。

给定两个参数，就可以形成（在乘法常数内）唯一的长度、速度和时间尺度。这些是柯尔莫哥洛夫音阶：

\begin{sourceblock}
\includegraphics[page=213,pagebox=mediabox,trim={64.39bp 1145.86bp 64.39bp 801.90bp},clip,width=0.94\textwidth,max height=0.38\textheight]{Pope.pdf}
\end{sourceblock}
\begin{sourceblock}
\includegraphics[page=213,pagebox=mediabox,trim={64.39bp 1068.87bp 64.39bp 876.89bp},clip,width=0.94\textwidth,max height=0.38\textheight]{Pope.pdf}
\end{sourceblock}
\begin{sourceblock}
\includegraphics[page=213,pagebox=mediabox,trim={64.39bp 994.87bp 64.39bp 945.88bp},clip,width=0.94\textwidth,max height=0.38\textheight]{Pope.pdf}
\end{sourceblock}
源自这些定义的两个恒等式清楚地表明，柯尔莫哥洛夫尺度描述了最小的耗散涡流。首先，基于柯尔莫哥洛夫尺度的雷诺数是统一的，即 qu,/v 1，这与级联进行到越来越小的尺度直到雷诺数 u(l)t/v 小到足以使耗散有效的概念是一致的。其次，耗散率由下式给出

\begin{sourceblock}
\includegraphics[page=213,pagebox=mediabox,trim={64.39bp 609.92bp 64.39bp 1338.83bp},clip,width=0.94\textwidth,max height=0.38\textheight]{Pope.pdf}
\end{sourceblock}
表明 ( u ) 1 提供了耗散涡流速度梯度的一致表征。

确定了柯尔莫哥洛夫量表后，我们现在可以陈述假设的结果，证明其效力，并澄清短语“相似性假设”和“通用形式”的含义。考虑一次在高雷诺数湍流中的点 xo。就 (xo,t o) 处的柯尔莫哥洛夫尺度而言，无量纲坐标定义为

\begin{sourceblock}
\includegraphics[page=213,pagebox=mediabox,trim={64.39bp 218.97bp 64.39bp 1726.78bp},clip,width=0.94\textwidth,max height=0.38\textheight]{Pope.pdf}
\end{sourceblock}
无量纲速度差场定义为

\begin{sourceblock}
\includegraphics[page=213,pagebox=mediabox,trim={64.39bp 119.39bp 64.39bp 1848.77bp},clip,width=0.94\textwidth,max height=0.38\textheight]{Pope.pdf}
\end{sourceblock}
不可能从 和 v 形成无量纲参数；因此（在量纲的基础上）无量纲场 w(y) 的统计量的“通用形式”不能依赖于 和 v。因此，根据上述柯尔莫哥洛夫假设，当在不太大的尺度上检查无量纲速度场时（特别是 lyl <
在所有高雷诺数湍流中，它在所有点 (xo,to) 处都是统计各向同性且统计相同。在小尺度上，所有高雷诺数湍流速度场在统计上都是相似的；也就是说，当它们按柯尔莫哥洛夫尺度缩放时，它们在统计上是相同的（方程（6.5）和（6.6））。

最小与最大尺度的比率很容易根据柯尔莫哥洛夫尺度的定义和尺度 ui/lo 确定。的

\begin{sourceblock}
\includegraphics[page=214,pagebox=mediabox,trim={64.39bp 1219.85bp 64.39bp 675.91bp},clip,width=0.94\textwidth,max height=0.38\textheight]{Pope.pdf}
\end{sourceblock}
\begin{sourceblock}
\includegraphics[page=214,pagebox=mediabox,trim={64.39bp 1137.86bp 64.39bp 801.90bp},clip,width=0.94\textwidth,max height=0.38\textheight]{Pope.pdf}
\end{sourceblock}
\begin{sourceblock}
\includegraphics[page=214,pagebox=mediabox,trim={64.39bp 1056.87bp 64.39bp 885.89bp},clip,width=0.94\textwidth,max height=0.38\textheight]{Pope.pdf}
\end{sourceblock}
显然，在高雷诺数下，最小涡流 (u, 和 ) 的速度尺度和时间尺度与之前假设的最大涡流 (uo 和 TO) 相比要小。

不可避免的是，从流可视化（例如，第 5 页的图 1.2）可以明显看出，q / l o 的比率随着 Re 的增加而减小。因此，在足够高的雷诺数下，存在一个尺度范围 l 与 lo 相比非常小，但与 q 相比却非常大，即 lo l q。由于该范围内的涡流比耗散涡流大得多，因此可以认为它们的雷诺数lu(l)/v很大，因此它们的运动受粘度的影响很小。因此，根据这一点和第一个相似性假设，我们有（近似表述）：

K\textasciitilde{}lmogorov 的第二个相似性假设。在足够的每个湍流中
迄今为止的高雷诺数，C 尺度运动的统计
范围 Co >> C >> q 有一个通用形式，由下式唯一确定
独立于

引入长度标尺（比如 6011）很方便，所以
lD1 \textasciitilde{} D I 上述假设中的范围可以写成 lEI>> 这个长度尺度 lD1sp 将普遍平衡范围 (l < lEIi)n 分为两个子范围：惯性子区 (lE>I l > l ) 和耗散区 (l < l )。顾名思义，根据第二个相似性假设，

\begin{sourceblock}
\includegraphics[page=215,pagebox=mediabox,trim={64.39bp 1437.82bp 64.39bp 109.44bp},clip,width=\textwidth,max height=0.72\textheight]{Pope.pdf}
\par\vspace{0.45\baselineskip}
\small 图 6.1。雷诺数非常高时的涡流大小 é（以对数尺度），显示各种长度尺度和范围。
\end{sourceblock}
效应可以忽略不计，而只有耗散区内的运动才会经历显着的粘性效应，因此基本上是所有耗散的原因。不同的长度尺度和范围被勾画在
长度尺度、速度尺度和时间尺度不能单独形成。然而，给定涡流大小 l （在惯性子区内），涡流的特征速度尺度和时间尺度由 和 l 形成：

\begin{sourceblock}
\includegraphics[page=215,pagebox=mediabox,trim={64.39bp 773.90bp 64.39bp 1166.85bp},clip,width=0.94\textwidth,max height=0.38\textheight]{Pope.pdf}
\end{sourceblock}
\begin{sourceblock}
\includegraphics[page=215,pagebox=mediabox,trim={64.39bp 693.91bp 64.39bp 1252.84bp},clip,width=0.94\textwidth,max height=0.38\textheight]{Pope.pdf}
\end{sourceblock}
那么，第二个相似性假设的结果是（在惯性子区内）速度尺度和时间尺度 u(l) 和 \textasciitilde{} ( dle) 随着 l 减小而减小。

在能量级联的概念中，一个至关重要的量
l ( l ) 表示能量从大于 l 的涡流转移到小于 l 的涡流的速率。如果此转移过程主要通过大小与 l 相当的涡流完成，则 T(1) 预计为 \textasciitilde{} ( l ) \textasciitilde{} / \textasciitilde{}T(hle) i.d 实体

\begin{sourceblock}
\includegraphics[page=215,pagebox=mediabox,trim={64.39bp 250.97bp 64.39bp 1695.79bp},clip,width=0.94\textwidth,max height=0.38\textheight]{Pope.pdf}
\end{sourceblock}
源自方程式。因此，（6.10）和（6.11）特别具有启发性，因为它表明 T(l)i 独立于 l（对于惯性子区内的 l）。

\begin{sourceblock}
\includegraphics[page=216,pagebox=mediabox,trim={64.39bp 1437.82bp 64.39bp 109.44bp},clip,width=\textwidth,max height=0.72\textheight]{Pope.pdf}
\par\vspace{0.45\baselineskip}
\small 图 6.2。高雷诺数下的能量级联示意图。
\end{sourceblock}
我们有

\begin{sourceblock}
\includegraphics[page=216,pagebox=mediabox,trim={64.39bp 1249.84bp 64.39bp 695.91bp},clip,width=0.94\textwidth,max height=0.38\textheight]{Pope.pdf}
\end{sourceblock}
(对于 lEI> l > ) 即，来自大尺度的能量传递速率 TE1d 决定了通过惯性子区 7 ( l ) 的恒定能量传递速率；因此能量离开惯性子区并进入耗散区 7 的速率； d 因此是耗散率 E。该图如图 6.2 所示。

\subsection*{6.1.3 能谱}\phantomsection\addcontentsline{toc}{subsection}{6.1.3 能谱}
湍动能如何在不同大小的涡流之间分布仍有待确定。通过考虑第 3 章中引入的能谱函数（方程（3.166）），对于均匀湍流来说，这是最容易完成的。

回想一下 3.7 节，长度尺度的运动对应于波数 2n:/l，并且波数范围内的能量

\begin{sourceblock}
\includegraphics[page=216,pagebox=mediabox,trim={64.39bp 462.94bp 64.39bp 1437.82bp},clip,width=0.94\textwidth,max height=0.38\textheight]{Pope.pdf}
\end{sourceblock}
[第 6.5 节中，对 E(K) 进行了一些详细考虑，其中一个令人感兴趣的结果是，运动对耗散率的贡献

\begin{sourceblock}
\includegraphics[page=216,pagebox=mediabox,trim={64.39bp 213.97bp 64.39bp 1689.79bp},clip,width=0.94\textwidth,max height=0.38\textheight]{Pope.pdf}
\end{sourceblock}
从柯尔莫哥洛夫的第一个相似性假设出发，在宇宙中-
和 v。从第二个假设可以看出，在惯性范围内
< < 2n:/lDI)，频谱为 ( K E \textasciitilde{} K KDI --=

\begin{sourceblock}
\includegraphics[page=217,pagebox=mediabox,trim={64.39bp 1675.79bp 64.39bp 261.97bp},clip,width=0.94\textwidth,max height=0.38\textheight]{Pope.pdf}
\end{sourceblock}
其中 C 是通用常数。 （这些断言在第 6.5 节中得到了证明。）

为了理解柯尔莫哥洛夫谱的一些基本特征，我们考虑一般幂律谱

\begin{sourceblock}
\includegraphics[page=217,pagebox=mediabox,trim={64.39bp 1449.82bp 64.39bp 494.94bp},clip,width=0.94\textwidth,max height=0.38\textheight]{Pope.pdf}
\end{sourceblock}
其中 A 和 p 是常数。波数所包含的能量越大

\begin{sourceblock}
\includegraphics[page=217,pagebox=mediabox,trim={64.39bp 1247.84bp 64.39bp 609.92bp},clip,width=0.94\textwidth,max height=0.38\textheight]{Pope.pdf}
\end{sourceblock}
对于 p > 1，而积分对于 p 5 1 发散。类似地，波数的耗散小于

\begin{sourceblock}
\includegraphics[page=217,pagebox=mediabox,trim={64.39bp 1041.87bp 64.39bp 858.89bp},clip,width=0.94\textwidth,max height=0.38\textheight]{Pope.pdf}
\end{sourceblock}
对于 p < 3，而积分对于 p 2 3 发散。因此，p = 对应于柯尔莫哥洛夫谱，位于积分 k(,,) 和 收敛的范围 (1,3) 的中间附近。中的能量量

\&(OK) 高波数随着 k(,,) 的增加而减少，而低波数的耗散随着 k(,,) 的减少而减少

K 4/3 K 趋于零。

尽管柯尔莫哥洛夫谱仅适用于惯性范围，但所做的观察与大部分能量处于大尺度 (l > lEIo r < 2n:/lEI) 和大部分耗散处于小尺度 (l < lDoIr > 2n:/lDI) 的概念一致。

\subsection*{6.1.4 科尔莫哥洛夫假设的重述}\phantomsection\addcontentsline{toc}{subsection}{6.1.4 科尔莫哥洛夫假设的重述}
为了推断。因此，有必要在这里提供对 Kolmogorov (1941) 假设的更精确的陈述。柯尔莫哥洛夫用四维 x-t 空间中的 N 点分布来表示这些。然而，在这里，我们考虑固定时间 t 物理空间 (x) 中的 N 点分布，这对于大多数目的来说是足够通用的。

考虑湍流中的一个简单域，并令 x (0)，
速度差定义为

\begin{sourceblock}
\includegraphics[page=218,pagebox=mediabox,trim={64.39bp 1725.78bp 64.39bp 215.97bp},clip,width=0.94\textwidth,max height=0.38\textheight]{Pope.pdf}
\end{sourceblock}
\begin{sourceblock}
\includegraphics[page=218,pagebox=mediabox,trim={64.39bp 1645.79bp 64.39bp 299.96bp},clip,width=0.94\textwidth,max height=0.38\textheight]{Pope.pdf}
\end{sourceblock}
N 点 y (1), , Y 处的联合 PDF 表示为
局部同质性的定义。湍流是局部同质的
在域 G 中是慷慨的，如果对于每个固定的 N 和 Ycn)(n 1,2? . , N)，
做）（做），
N 点 PDF f 独立于 和 U t)。

局部各向同性的定义。湍流是局部各向同性的
域，如果它是局部齐次的，并且如果此外
PDF 对于旋转和反射而言是不变的
坐标轴。

局部各向同性假说。在任何具有足够的湍流
古代较大的雷诺数 (Re UC/v)，湍流为
良好的近似，如果域足够，则局部各向同性
小（ e mly (")\textasciitilde{} C，对于所有 n）并且不靠近边界
流或其其他奇点。

第一个相似性假设。对于局部各向同性湍流，
N 点 PDF 由粘度 v 和
耗散率 E .

第二个相似性假设。如果向量 y(") 和
它们的差异 y (4 y(") 与
柯尔莫哥洛夫尺度 q，则 N 点 PDF 唯一确定
由且不依赖于 v.

重要的是要观察到这些假设特别适用于速度差异。 Nlpoint PDF 的使用允许将假设应用于任何湍流，而波数谱的陈述仅适用于统计上均匀的流（至少在一个方向上）。

对于非均匀流，局部各向同性只能“达到良好的近似”（如假设中所述）。例如，取 y(') 和
- （其中 l 是指定长度，e 是指定单位向量），我们
- -埃尔有

\begin{sourceblock}
\includegraphics[page=218,pagebox=mediabox,trim={64.39bp 177.98bp 64.39bp 1698.79bp},clip,width=0.94\textwidth,max height=0.38\textheight]{Pope.pdf}
\end{sourceblock}
\section*{6.2 结构函数}\phantomsection\addcontentsline{toc}{section}{6.2 结构函数}
\begin{sourceblock}
\includegraphics[page=219,pagebox=mediabox,trim={64.39bp 1401.82bp 64.39bp 109.44bp},clip,width=\textwidth,max height=0.72\textheight]{Pope.pdf}
\par\vspace{0.45\baselineskip}
\small 图 6.3。以 x 和 y 形式显示点 x 和 x+r 的草图”。所有点都在域 G 内。
\end{sourceblock}
由大尺度不均匀性引起的 t/L 阶小各向异性分量。

6.2 结构函数 为了说明 Kolmogorov 假设的正确应用，我们像 Kolmogorov (1941b) 一样考虑二阶速度结构函数。推导假设的预测，然后与实验数据进行比较。

根据定义，二阶速度结构函数是两点 x r 和 x 之间速度差的协方差：

\begin{sourceblock}
\includegraphics[page=219,pagebox=mediabox,trim={64.39bp 729.91bp 64.39bp 1204.85bp},clip,width=0.94\textwidth,max height=0.38\textheight]{Pope.pdf}
\end{sourceblock}
很明显，柯尔莫哥洛夫假设适用于该统计量，但这可以通过用位置和速度差 y 重新表达并由等式定义来验证。 （6.20）和（6.21）：

\begin{sourceblock}
\includegraphics[page=219,pagebox=mediabox,trim={64.39bp 509.94bp 64.39bp 1428.82bp},clip,width=0.94\textwidth,max height=0.38\textheight]{Pope.pdf}
\end{sourceblock}
见图 6.3。我们假设满足所有其他条件（例如足够大的雷诺数）。

局部各向同性假设的第一个含义是（对于 r 1 \textasciitilde{} <1 < L ） Dij 独立于 do) E。那么，方程 (6.24) 表明 Dij 不依赖于它的第二个参数，即 Dij(r,x,t ) 独立于 x。从其定义（方程（6.23））可以明显看出，Dij(r,t ) 取决于 r = y -Y ( I，但不分别取决于 y(') 和 y(2)s。因此 Dij(r, t) 是 r 的各向同性函数。

\begin{sourceblock}
\includegraphics[page=220,pagebox=mediabox,trim={64.39bp 1436.82bp 64.39bp 109.44bp},clip,width=\textwidth,max height=0.72\textheight]{Pope.pdf}
\par\vspace{0.45\baselineskip}
\small 图 6.4。 r = eyr 的纵向和横向结构函数中涉及的速度分量草图。
\end{sourceblock}
由向量 r 形成的是 dij 和 rirj。因此 Dij 可以写成

\begin{sourceblock}
\includegraphics[page=220,pagebox=mediabox,trim={64.39bp 1187.85bp 64.39bp 729.91bp},clip,width=0.94\textwidth,max height=0.38\textheight]{Pope.pdf}
\end{sourceblock}
其中标量函数和 DNa 分别称为纵向和横向结构函数。如果选择坐标系使 r 位于 xl 方向 ( e r = elr)，那么我们得到
与图 6.4 一起显示了 DJALan d DNN 的重要性。

在 (U) 0 的均匀湍流中，连续性方程的结果是
参见练习 6.1。在当前上下文中，由于局部同质性，该方程也适用（良好的近似）。然后从等式得出。 （6.25）和（6.27）（见练习 6.2） DNNi 是由 DJALac 唯一确定的，根据

\begin{sourceblock}
\includegraphics[page=220,pagebox=mediabox,trim={64.39bp 377.95bp 64.39bp 1526.81bp},clip,width=0.94\textwidth,max height=0.38\textheight]{Pope.pdf}
\end{sourceblock}
因此，在局部各向同性湍流中，Dij(r,t) 由单个标量函数 DLL(rt,) 确定。

根据第一个相似性假设，给定 r (Irl << L)，Dij 由 和 v 唯一确定。 (\textasciitilde{}r)h\textasciitilde{}as/d\textasciitilde{}im 是速度平方的量，因此可用于使 Dij 无量纲。有
和 v，可以方便地视为 = r/y。因此，根据第一个相似性假设，存在一个通用的无量纲函数 DJdL(r/ys)u ch，（受假设的条件限制）

\begin{sourceblock}
\includegraphics[page=221,pagebox=mediabox,trim={64.39bp 1639.79bp 64.39bp 295.96bp},clip,width=0.94\textwidth,max height=0.38\textheight]{Pope.pdf}
\end{sourceblock}
根据第二个相似性假设，对于大的 r/y ( C: r y )，DJdLi 独立于 v。在这种情况下，不存在可以由 和 r 形成的无量纲群，因此 DLLi 由下式给出

\begin{sourceblock}
\includegraphics[page=221,pagebox=mediabox,trim={64.39bp 1381.83bp 64.39bp 565.93bp},clip,width=0.94\textwidth,max height=0.38\textheight]{Pope.pdf}
\end{sourceblock}
其中 C2 是通用常数（这意味着，对于较大的 r/y，DL 渐近趋于常数值 C2。）横向结构函数由方程 1 得出。 （6.28），

\begin{sourceblock}
\includegraphics[page=221,pagebox=mediabox,trim={64.39bp 1164.85bp 64.39bp 783.90bp},clip,width=0.94\textwidth,max height=0.38\textheight]{Pope.pdf}
\end{sourceblock}
因此，从方程式(6.25)，Dij 由下式给出

\begin{sourceblock}
\includegraphics[page=221,pagebox=mediabox,trim={64.39bp 992.87bp 64.39bp 903.89bp},clip,width=0.94\textwidth,max height=0.38\textheight]{Pope.pdf}
\end{sourceblock}
因此，在惯性子区 ( r y ) 中，柯尔莫哥洛夫假设足以确定用 E、r 和通用常数 C2 表示的二阶结构函数。

科尔莫哥洛夫假设的预测体现在方程中。 Saddoughi 和 Veeravalli (1994) 在湍流边界层中测试了 (6.32) 式
据称是实验室中达到的最高雷诺数边界层。主流动方向为 x x1，垂直于壁的距离为 y x2。测量结果在 y 400 mm 处报告，该位置边界层厚度为 6 × 1,090 mm，雷诺数（基于 6）为 3.6 x lo6，测量的柯尔莫哥洛夫标度为 q = 0.09 mm。以C=6为特征流长，比值UV=
12,000 表明存在非常大的尺度分离。

使用 elr 和 ry，根据方程进行预测。 (6.32) 可以写成

\begin{sourceblock}
\includegraphics[page=221,pagebox=mediabox,trim={64.39bp 313.96bp 64.39bp 1632.79bp},clip,width=0.94\textwidth,max height=0.38\textheight]{Pope.pdf}
\end{sourceblock}
\begin{sourceblock}
\includegraphics[page=221,pagebox=mediabox,trim={64.39bp 237.97bp 64.39bp 1704.79bp},clip,width=0.94\textwidth,max height=0.38\textheight]{Pope.pdf}
\end{sourceblock}
\begin{sourceblock}
\includegraphics[page=221,pagebox=mediabox,trim={64.39bp 161.98bp 64.39bp 1785.78bp},clip,width=0.94\textwidth,max height=0.38\textheight]{Pope.pdf}
\end{sourceblock}
\begin{sourceblock}
\includegraphics[page=222,pagebox=mediabox,trim={64.39bp 804.90bp 64.39bp 109.44bp},clip,width=\textwidth,max height=0.72\textheight]{Pope.pdf}
\par\vspace{0.45\baselineskip}
\small 图 6.5。在高雷诺数湍流边界层中测量的二阶速度结构函数。水平线显示了惯性子区中柯尔莫哥洛夫假设的预测，方程。 （6.33）和（6.34）。 （摘自 Saddoughi 和 Veeravalli (1994)。）
\end{sourceblock}
上述预测很容易得到检验。根据这些数据和其他数据建议的 C2 值，我们得出以下结论。

(i) 对于 7,000 q r 20q，Dll/(\&r) 在 C2 的 +15\% 范围内。

(ii) 和 之间没有明显差异
(iii) 对于 1,200 q '1 > >r 12q，D 2 2 / ( \textasciitilde{} r ) 在 15\% 以内
在上面给出的 r 范围内，D l l 和 D l l 分别变化 50 和 20 倍，因此相比之下，+15\% 的变化可以认为很小。

显然，这些实验观察为柯尔莫哥洛夫假设提供了实质性支持。然而，还需要考虑其他流量

\section*{6.3 两点相关}\phantomsection\addcontentsline{toc}{section}{6.3 两点相关}
可以检查。与实验数据的进一步比较在 6.5 节中给出。

对于 (U) 0 的均匀湍流，表明结构
ture函数Dij(r)（方程（6.23））和两点相关性R*
（方程（3.160））由下式相关：

\begin{sourceblock}
\includegraphics[page=223,pagebox=mediabox,trim={64.39bp 1410.82bp 64.39bp 534.93bp},clip,width=0.94\textwidth,max height=0.38\textheight]{Pope.pdf}
\end{sourceblock}
证明连续性方程的结果是

\begin{sourceblock}
\includegraphics[page=223,pagebox=mediabox,trim={64.39bp 1239.84bp 64.39bp 657.92bp},clip,width=0.94\textwidth,max height=0.38\textheight]{Pope.pdf}
\end{sourceblock}
（提示：请参阅第 78 页的练习 3.34 和 3.35。）

微分方程。 (6.25) 得到

\begin{sourceblock}
\includegraphics[page=223,pagebox=mediabox,trim={64.39bp 1033.87bp 64.39bp 864.89bp},clip,width=0.94\textwidth,max height=0.38\textheight]{Pope.pdf}
\end{sourceblock}
因此验证方程。 （6.28）。

表明，对于较小的 r (r q)，在各向同性湍流中，第二阶
阶次速度结构函数是

\begin{sourceblock}
\includegraphics[page=223,pagebox=mediabox,trim={64.39bp 754.90bp 64.39bp 1121.86bp},clip,width=0.94\textwidth,max height=0.38\textheight]{Pope.pdf}
\end{sourceblock}
\begin{sourceblock}
\includegraphics[page=223,pagebox=mediabox,trim={64.39bp 613.92bp 64.39bp 1262.84bp},clip,width=0.94\textwidth,max height=0.38\textheight]{Pope.pdf}
\end{sourceblock}
沃点

6.5 I 相关性 柯尔莫哥洛夫假设以及从中得出的推论，与纳维-斯托克斯方程没有直接联系（尽管，如前一节所述，通常会调用连续性方程）。尽管在能量级联的描述中，能量向连续较小尺度的转移已被认为是最重要的现象，但这种转移发生的精确机制尚未确定或量化。因此，尝试从纳维-斯托克斯方程中提取信息是很自然的。

（本节概述）是 Taylor (1935a) 以及 van \textasciitilde{} 6 r m a n 和 Howarth (1938) 的理论，它们基于两点相关性。接下来的两节给出了波数空间中两点相关性的傅立叶变换能谱的视图。

自相关函数 考虑均匀各向同性湍流，平均速度 , r.m,s 为零。速度 u'(t) 和耗散率 \textasciitilde{} ( t )B。由于同质性，两点

\begin{sourceblock}
\includegraphics[page=224,pagebox=mediabox,trim={64.39bp 1395.82bp 64.39bp 495.94bp},clip,width=0.94\textwidth,max height=0.38\textheight]{Pope.pdf}
\end{sourceblock}
独立于 x。在原点上它是
既没有产生也没有传输，因此湍动能 k(t) 2 \textasciitilde{} ' ( t )(\textasciitilde{}E q. (5.132)) 的演化方程简化为

\begin{sourceblock}
\includegraphics[page=224,pagebox=mediabox,trim={64.39bp 1068.87bp 64.39bp 843.89bp},clip,width=0.94\textwidth,max height=0.38\textheight]{Pope.pdf}
\end{sourceblock}
正如结构函数 Dij 一样，各向同性的结果是 Rij 可以用两个标量函数 f (r, t) 和 g(r, t) 表示：

\begin{sourceblock}
\includegraphics[page=224,pagebox=mediabox,trim={64.39bp 839.89bp 64.39bp 1076.86bp},clip,width=0.94\textwidth,max height=0.38\textheight]{Pope.pdf}
\end{sourceblock}
（参见等式（6.25））。有了 elr，这个方程就变成了

\begin{sourceblock}
\includegraphics[page=224,pagebox=mediabox,trim={64.39bp 666.92bp 64.39bp 1270.84bp},clip,width=0.94\textwidth,max height=0.38\textheight]{Pope.pdf}
\end{sourceblock}
R33 R22,, Rij 0, 对于 i j，从而将 f 和 g 分别确定为纵向和横向自相关函数。 （请注意，f 和 g 是无量纲的，其中 f (0, t) = g(0, t) 1。）再次与 Dij 的属性平行，连续性方程意味着 aRij/arj 0 （参见练习 3,35），与方程结合： [6.44), 导致

\begin{sourceblock}
\includegraphics[page=224,pagebox=mediabox,trim={64.39bp 282.96bp 64.39bp 1624.80bp},clip,width=0.94\textwidth,max height=0.38\textheight]{Pope.pdf}
\end{sourceblock}
因此，在各向同性湍流中，两点相关性 Rij(r,t ) 完全由纵向自相关函数 f (r,t) 确定。图 6.6 显示了在近各向同性网格生成的湍流中 f ( r , t ) 的测量结果

\begin{sourceblock}
\includegraphics[page=225,pagebox=mediabox,trim={64.39bp 1269.84bp 64.39bp 109.44bp},clip,width=\textwidth,max height=0.72\textheight]{Pope.pdf}
\par\vspace{0.45\baselineskip}
\small 图 6.6。网格湍流中纵向速度自相关函数 f(r, 2) 的测量：x;/M = A。（来自 Comte-Beilot 和 Corrsin (1971)。）
\end{sourceblock}
有两个不同的纵向长度尺度，Ll (t) 和 Af(t)，可以从 f 定义； a 然后有相应的横向长度尺度 L22(t)a 和 A, (t) 由 g 定义。

积分长度 hscales 从 f(r, t) 获得的第一个长度尺度是纵向积分

\begin{sourceblock}
\includegraphics[page=225,pagebox=mediabox,trim={64.39bp 747.91bp 64.39bp 1108.86bp},clip,width=0.94\textwidth,max height=0.38\textheight]{Pope.pdf}
\end{sourceblock}
我们已经遇到过（例如，第 5.1 节，第 110 页的图 5.13）。积分尺度 Lll(t) 只是 f (r, t) 曲线下的面积，因此检查图 6.6 立即发现 Lll 随时间增长（在网格湍流中）。正如之前所观察到的，Lll 是较大涡流的特征。在各向同性湍流中，横向积分尺度

\begin{sourceblock}
\includegraphics[page=225,pagebox=mediabox,trim={64.39bp 402.95bp 64.39bp 1499.81bp},clip,width=0.94\textwidth,max height=0.38\textheight]{Pope.pdf}
\end{sourceblock}
只是 Lll(t) 的一半（参见练习 6.4）。

显示方程。 (6.46) 可以重写

\begin{sourceblock}
\includegraphics[page=225,pagebox=mediabox,trim={64.39bp 119.39bp 64.39bp 1827.77bp},clip,width=0.94\textwidth,max height=0.38\textheight]{Pope.pdf}
\end{sourceblock}
因此，在各向同性湍流中，横向积分尺度

\begin{sourceblock}
\includegraphics[page=226,pagebox=mediabox,trim={64.39bp 1688.79bp 64.39bp 214.97bp},clip,width=0.94\textwidth,max height=0.38\textheight]{Pope.pdf}
\end{sourceblock}
是纵向尺度的一半，即

\begin{sourceblock}
\includegraphics[page=226,pagebox=mediabox,trim={64.39bp 1563.80bp 64.39bp 381.95bp},clip,width=0.94\textwidth,max height=0.38\textheight]{Pope.pdf}
\end{sourceblock}
由方程式可知（6.46）

\begin{sourceblock}
\includegraphics[page=226,pagebox=mediabox,trim={64.39bp 1394.82bp 64.39bp 507.94bp},clip,width=0.94\textwidth,max height=0.38\textheight]{Pope.pdf}
\end{sourceblock}
（假设对于较大的 r，f (r, t) 比 r-2 衰减得更快）。

泰勒微尺度 从 f(r,t) 获得的第二个长度尺度是纵向泰勒微尺度 if(t)。由于 f(r,t) 是 r 的偶函数且不大于 1，因此原点 f '(0, t) (8f /ar)r=Oi 处的一阶导数为零；而二阶导数 f "(0, t) f /ar2)r=0i 是非正的。正如我们将看到的，在湍流中 f"(0) 严格为负，因此 if( t) 定义为

\begin{sourceblock}
\includegraphics[page=226,pagebox=mediabox,trim={64.39bp 911.88bp 64.39bp 1035.87bp},clip,width=0.94\textwidth,max height=0.38\textheight]{Pope.pdf}
\end{sourceblock}
是实数、正数，并且具有长度尺寸。

几何结构使这个深奥的定义变得清晰。令 p(r) 为在 r 0 处密切 f ( r ) 的抛物线（ i。具有 p(0) f(O)、p'(0) f'(O) 和 pfl(0) ffl(0) 的抛物线）。显然 p(r) 是

\begin{sourceblock}
\includegraphics[page=226,pagebox=mediabox,trim={64.39bp 588.93bp 64.39bp 1295.84bp},clip,width=0.94\textwidth,max height=0.38\textheight]{Pope.pdf}
\end{sourceblock}
因此，如图 6.7 所示，密切抛物线与轴相交于
如以下操作所示，f"(0, t)（因此 if(t)）与速度导数相关：

-u f N( 0, t) = -ur2 lim f k ? )

\begin{sourceblock}
\includegraphics[page=227,pagebox=mediabox,trim={64.39bp 1245.84bp 64.39bp 109.44bp},clip,width=\textwidth,max height=0.72\textheight]{Pope.pdf}
\par\vspace{0.45\baselineskip}
\small 图 6.7。纵向速度自相关函数草图显示了泰勒微尺度 Ay 的定义。
\end{sourceblock}
\begin{sourceblock}
\includegraphics[page=227,pagebox=mediabox,trim={64.39bp 815.90bp 64.39bp 1059.87bp},clip,width=0.94\textwidth,max height=0.38\textheight]{Pope.pdf}
\end{sourceblock}
这样我们就得到

\begin{sourceblock}
\includegraphics[page=227,pagebox=mediabox,trim={64.39bp 638.92bp 64.39bp 1229.84bp},clip,width=0.94\textwidth,max height=0.38\textheight]{Pope.pdf}
\end{sourceblock}
横向泰勒微尺度 A (t)，定义为

\begin{sourceblock}
\includegraphics[page=227,pagebox=mediabox,trim={64.39bp 504.94bp 64.39bp 1439.82bp},clip,width=0.94\textwidth,max height=0.38\textheight]{Pope.pdf}
\end{sourceblock}
在各向同性湍流中，等于 if(t) /\$（参见练习 6.6）。然后根据这两个方程和关系 1 5 \textasciitilde{} ( ( 8 \textasciitilde{} \textasciitilde{} (/E8q*\textasciitilde{}( 5\textasciitilde{}.1)7\textasciitilde{}1))) 耗散由下式给出

\begin{sourceblock}
\includegraphics[page=227,pagebox=mediabox,trim={64.39bp 277.96bp 64.39bp 1668.79bp},clip,width=0.94\textwidth,max height=0.38\textheight]{Pope.pdf}
\end{sourceblock}
在一篇标志着各向同性湍流研究开始的经典论文中，Taylor (1935a) 定义了 A 并获得了上面的 He 方程，然后
负责能量耗散的最小涡流。从等式中扣除。 (6.58) 是不正确的，因为它错误地假设 u' 是耗散涡流的特征速度。相反，最小涡流的特征长度和速度尺度是柯尔莫哥洛夫尺度 q 和 u。

为了确定泰勒尺度和柯尔莫哥洛夫尺度之间的关系，我们将 L k 3/2 定义为表征大涡流的长度尺度，并将湍流雷诺数定义为

\begin{sourceblock}
\includegraphics[page=228,pagebox=mediabox,trim={64.39bp 1390.82bp 64.39bp 540.93bp},clip,width=0.94\textwidth,max height=0.38\textheight]{Pope.pdf}
\end{sourceblock}
然后微尺度由下式给出

\begin{sourceblock}
\includegraphics[page=228,pagebox=mediabox,trim={64.39bp 1164.85bp 64.39bp 675.91bp},clip,width=0.94\textwidth,max height=0.38\textheight]{Pope.pdf}
\end{sourceblock}
因此，在高雷诺数下，i l 的尺寸介于 q 和 L 之间。

泰勒尺度没有明确的物理解释。然而，它是一个经常使用的明确定义的量。特别是泰勒尺度雷诺数

\begin{sourceblock}
\includegraphics[page=228,pagebox=mediabox,trim={64.39bp 915.88bp 64.39bp 1030.87bp},clip,width=0.94\textwidth,max height=0.38\textheight]{Pope.pdf}
\end{sourceblock}
传统上用于表征网格湍流。观察，从方程。 (6.60)，Rn 随着积分尺度雷诺数的平方根而变化

\begin{sourceblock}
\includegraphics[page=228,pagebox=mediabox,trim={64.39bp 741.91bp 64.39bp 1205.85bp},clip,width=0.94\textwidth,max height=0.38\textheight]{Pope.pdf}
\end{sourceblock}
此外，还可以观察到该比率

\begin{sourceblock}
\includegraphics[page=228,pagebox=mediabox,trim={64.39bp 610.92bp 64.39bp 1312.83bp},clip,width=0.94\textwidth,max height=0.38\textheight]{Pope.pdf}
\end{sourceblock}
正确地描述了小涡流的时间尺度。

由方程式可知（6.46）

\begin{sourceblock}
\includegraphics[page=228,pagebox=mediabox,trim={64.39bp 354.95bp 64.39bp 1592.80bp},clip,width=0.94\textwidth,max height=0.38\textheight]{Pope.pdf}
\end{sourceblock}
因此横向泰勒微尺度

\begin{sourceblock}
\includegraphics[page=228,pagebox=mediabox,trim={64.39bp 224.97bp 64.39bp 1709.78bp},clip,width=0.94\textwidth,max height=0.38\textheight]{Pope.pdf}
\end{sourceblock}
与纵向尺度 if( t) 的关系为

\begin{sourceblock}
\includegraphics[page=228,pagebox=mediabox,trim={64.39bp 119.39bp 64.39bp 1836.77bp},clip,width=0.94\textwidth,max height=0.38\textheight]{Pope.pdf}
\end{sourceblock}
\begin{sourceblock}
\includegraphics[page=229,pagebox=mediabox,trim={64.39bp 1695.79bp 64.39bp 158.98bp},clip,width=0.94\textwidth,max height=0.38\textheight]{Pope.pdf}
\end{sourceblock}
Kurmun-Howarth 方程 von Karman 和 Howarth (1938) 从 Navier-Stokes 方程获得 f (r, t) 的演化方程。我们在这里概述了主要步骤、结果和一些含义：详细的推导可以在原始工作或标准参考文献中找到（例如，Hinze (1975)、Monin 和 Yaglom (1975)）。

Rij(r,x , t) 的时间导数可表示为

\begin{sourceblock}
\includegraphics[page=229,pagebox=mediabox,trim={64.39bp 1034.87bp 64.39bp 684.91bp},clip,width=0.94\textwidth,max height=0.38\textheight]{Pope.pdf}
\end{sourceblock}
然后是纳维-斯托克斯方程，即

\begin{sourceblock}
\includegraphics[page=229,pagebox=mediabox,trim={64.39bp 870.89bp 64.39bp 1049.87bp},clip,width=0.94\textwidth,max height=0.38\textheight]{Pope.pdf}
\end{sourceblock}
可用于消除等式右侧的时间导数。 （6.70）。出现了三种类型的项，分别对应于方程（1）中的对流项、压力梯度项和粘性项。 (6.7 1)。对于各向同性湍流，Rij(r,t ) 方程中的压力梯度项为零。

对流项涉及两点三重速度相关性，例如

\begin{sourceblock}
\includegraphics[page=229,pagebox=mediabox,trim={64.39bp 552.93bp 64.39bp 1356.83bp},clip,width=0.94\textwidth,max height=0.38\textheight]{Pope.pdf}
\end{sourceblock}
正如 Rij 由 f 唯一确定（方程（6.44））一样，在各向同性湍流 Sijk 中
由纵向相关性唯一确定

\begin{sourceblock}
\includegraphics[page=229,pagebox=mediabox,trim={64.39bp 328.96bp 64.39bp 1559.80bp},clip,width=0.94\textwidth,max height=0.38\textheight]{Pope.pdf}
\end{sourceblock}
可以证明K(r,t)是r的奇函数，并且连续性方程意味着0 t) = 0，因此它的级数展开式为

\begin{sourceblock}
\includegraphics[page=229,pagebox=mediabox,trim={64.39bp 158.98bp 64.39bp 1783.78bp},clip,width=0.94\textwidth,max height=0.38\textheight]{Pope.pdf}
\end{sourceblock}
通过此过程，可以从纳维-斯托克斯方程获得 f (r t) 的精确方程：它是 Kdrmun-Howarth 方程

\begin{sourceblock}
\includegraphics[page=230,pagebox=mediabox,trim={64.39bp 1641.79bp 64.39bp 254.97bp},clip,width=0.94\textwidth,max height=0.38\textheight]{Pope.pdf}
\end{sourceblock}
需要提出的主要意见如下。

(i) 存在闭合问题。这个单一方程涉及两个未知数
函数 f (r, t) 和 k(r, t)。

(ii) 和 v 中的术语代表惯性和粘性过程，分别是
(iii) 在 r 0 处，项 in 消失（根据方程（6.74））；同时，
由于 f 在 r 中是偶数，我们得到
因此，对于 r 0，Karman-Howarth 方程简化为 ( 2 次)

动能方程：

\begin{sourceblock}
\includegraphics[page=230,pagebox=mediabox,trim={64.39bp 924.88bp 64.39bp 966.88bp},clip,width=0.94\textwidth,max height=0.38\textheight]{Pope.pdf}
\end{sourceblock}
(iv) 理查森-柯尔莫哥洛夫关于高能级联的观点
雷诺数，能量从较大尺度到较小尺度的转移
是一个惯性过程（至少对于 r q 而言）。因此，本次转让
将能量转移到更小的尺度是通过中的术语来完成的

Kkrman-Howarth 方程。

(v) 如果 u(x,t) 是高斯场，则 k(r,t) 与所有三阶矩一样
将为零。因此，“能量级联取决于非高斯”

速度场的各个方面。

通过遵循类似于方程式的过程。 （6.55），证明

\begin{sourceblock}
\includegraphics[page=230,pagebox=mediabox,trim={64.39bp 221.97bp 64.39bp 1710.78bp},clip,width=0.94\textwidth,max height=0.38\textheight]{Pope.pdf}
\end{sourceblock}
其中 S 是速度导数的偏度（方程（6.85））。 （提示：

表明，在各向同性湍流中，纵向结构函数
和自相关函数的关系为

\begin{sourceblock}
\includegraphics[page=231,pagebox=mediabox,trim={64.39bp 1683.79bp 64.39bp 255.97bp},clip,width=0.94\textwidth,max height=0.38\textheight]{Pope.pdf}
\end{sourceblock}
三阶结构函数定义为

\begin{sourceblock}
\includegraphics[page=231,pagebox=mediabox,trim={64.39bp 1555.80bp 64.39bp 388.95bp},clip,width=0.94\textwidth,max height=0.38\textheight]{Pope.pdf}
\end{sourceblock}
表明

\begin{sourceblock}
\includegraphics[page=231,pagebox=mediabox,trim={64.39bp 1437.82bp 64.39bp 506.94bp},clip,width=0.94\textwidth,max height=0.38\textheight]{Pope.pdf}
\end{sourceblock}
表明Karman-Howarth方程(Eq.(6.75))重新表达
就结构函数而言是

对该方程进行积分可得

\begin{sourceblock}
\includegraphics[page=231,pagebox=mediabox,trim={64.39bp 1077.86bp 64.39bp 822.90bp},clip,width=0.94\textwidth,max height=0.38\textheight]{Pope.pdf}
\end{sourceblock}
进一步观察 卡门-豪沃斯方程已经被广泛研究，并且获得的结果比这里提到的要多得多。现在给出了一些更为人所知和信息最丰富的结果； Batchelor (1953)、Monin 和 Yaglom (1975) 以及 Hinze (1975) 提供了更全面的说明。

速度导数的偏度 量 k"'(0, t) 决定了 k(r, t) 的前导阶数（方程（6.74）），可以重新表示为
是速度导数偏度（被发现为负）参见练习 6.7。所以，这种偏态涡旋拉伸与不同尺度之间的能量传递之间存在联系。

柯尔莫哥洛夫定律：Kaman-Howarth 方程可以用结构函数 DLL(rt,) 和 DLLL(rt), , 重新表示

\begin{sourceblock}
\includegraphics[page=232,pagebox=mediabox,trim={64.39bp 1441.82bp 64.39bp 500.94bp},clip,width=0.94\textwidth,max height=0.38\textheight]{Pope.pdf}
\end{sourceblock}
结果（参见练习 6.9）是 Kolmogorov 方程（Kolmogorov 194 la）

\begin{sourceblock}
\includegraphics[page=232,pagebox=mediabox,trim={64.39bp 1262.84bp 64.39bp 643.92bp},clip,width=0.94\textwidth,max height=0.38\textheight]{Pope.pdf}
\end{sourceblock}
从中可以得出一些有用的结果。科尔莫哥洛夫认为，在局部各向同性湍流中，左侧的非定常项为零，并且粘性项在惯性子区中可以忽略不计。这导致了柯尔莫哥洛夫定律：我

\begin{sourceblock}
\includegraphics[page=232,pagebox=mediabox,trim={64.39bp 986.88bp 64.39bp 960.88bp},clip,width=0.94\textwidth,max height=0.38\textheight]{Pope.pdf}
\end{sourceblock}
Kol.mogorov 进一步论证了结构-功能偏度

\begin{sourceblock}
\includegraphics[page=232,pagebox=mediabox,trim={64.39bp 837.89bp 64.39bp 1104.86bp},clip,width=0.94\textwidth,max height=0.38\textheight]{Pope.pdf}
\end{sourceblock}
是常数，导致

\begin{sourceblock}
\includegraphics[page=232,pagebox=mediabox,trim={64.39bp 640.92bp 64.39bp 1256.84bp},clip,width=0.94\textwidth,max height=0.38\textheight]{Pope.pdf}
\end{sourceblock}
这与柯尔莫哥洛夫假设（方程（6.30））的预测相同。因此，对于惯性子区中的 D Lt)，这一发展表明了柯尔莫哥洛夫假设和纳维-斯托克斯方程之间的一致性。它还将常数 C2 与偏度联系起来
Loitsyans kii 积分。将 Karman-Howarth 方程（方程 (6.75)）乘以 r4 并在 0 和 R 之间积分，我们得到

\begin{sourceblock}
\includegraphics[page=232,pagebox=mediabox,trim={64.39bp 119.39bp 64.39bp 1797.77bp},clip,width=0.94\textwidth,max height=0.38\textheight]{Pope.pdf}
\end{sourceblock}
Loitsyanskii (1939) 考虑了极限 R m，并假设 f 和随 r 迅速减小，Loitsyanskii 积分

\begin{sourceblock}
\includegraphics[page=233,pagebox=mediabox,trim={64.39bp 1648.79bp 64.39bp 253.97bp},clip,width=0.94\textwidth,max height=0.38\textheight]{Pope.pdf}
\end{sourceblock}
收敛，使得 k ( \textasciitilde{} , t a)n d fr(R,t) 中的元组消失。有了这些假设，方程。 (6.91) 表明 B2 不随时间变化，s\textasciitilde{} B2 被称为 Loitsyanskii 不变量。然而，所做的假设是不正确的。根据各向同性湍流的创建方式，Loitsyanskii 积分可以是有限的，也可以是发散的 (Saffman 1967)。当它是有限的时，发现式中的项： (6.91) 不会随着 R 趋于无穷大而消失，事实上 B2 随着时间的推移而增加（例如，参见 Chasnov
衰变的最后阶段 随着各向同性湍流的衰变，雷诺数减小，因此相对于粘性过程，惯性效应减弱。最终，当雷诺数足够小时，惯性效应可以忽略不计。

对于最后的衰减期，Batchelor 和 Townsend (1948) 表明，忽略惯性项的 Kkrnhn-Howarth 方程承认自相似解

\begin{sourceblock}
\includegraphics[page=233,pagebox=mediabox,trim={64.39bp 872.89bp 64.39bp 1068.87bp},clip,width=0.94\textwidth,max height=0.38\textheight]{Pope.pdf}
\end{sourceblock}
这与实验数据非常吻合。需要强调的是，该解决方案适用于比通常感兴趣的低得多的非常低的雷诺数。

验证方程。 (6.93) 满足 Khrmin-Howarth 方程
涉及 I 的惯性项的最后一个周期；可以忽略不计。

证明最后湍动能以 k t-5/2 衰减
对于 h\textasciitilde{}moge\textasciitilde{}neouisso 热带湍流，考虑六阶十-

\begin{sourceblock}
\includegraphics[page=233,pagebox=mediabox,trim={64.39bp 280.96bp 64.39bp 1583.80bp},clip,width=0.94\textwidth,max height=0.38\textheight]{Pope.pdf}
\end{sourceblock}
这是一个各向同性张量，因此可以写成总和
标量系数乘以克罗内克德尔塔乘积，例如，
a16 6 6。有 15 种不同的 Kronecker Delta 产品（对应
ip 34 kr
鉴于对称性（例如，Xijkpq- - r ），一般表示
-- a16 6 6 a2(6iP6jk6qr 6jq\&k6pr 6krdijdpq)

IP 34 克朗
xi jkpqr
a3(6ip6jr6qk 6jq6ir6pk 6kr6ipbpj) a ( 6 6 6 airdpjdqk)

4 智商PK jr
a5(6 6 +6ij6gk6pr +6ikapjbqr
我\textasciitilde{} PK qr

\begin{sourceblock}
\includegraphics[page=234,pagebox=mediabox,trim={64.39bp 1539.81bp 64.39bp 402.95bp},clip,width=0.94\textwidth,max height=0.38\textheight]{Pope.pdf}
\end{sourceblock}
表明连续性方程意味着 -- 0，其中
导致关系

\begin{sourceblock}
\includegraphics[page=234,pagebox=mediabox,trim={64.39bp 1269.84bp 64.39bp 573.93bp},clip,width=0.94\textwidth,max height=0.38\textheight]{Pope.pdf}
\end{sourceblock}
通过考虑数量

（在均匀湍流中为零）表明为零，
这导致了以下关系

证明这四个关系式。 (6.96)-(6.97) 确定所有
以 a 表示的系数，如

\begin{sourceblock}
\includegraphics[page=234,pagebox=mediabox,trim={64.39bp 775.90bp 64.39bp 1167.85bp},clip,width=0.94\textwidth,max height=0.38\textheight]{Pope.pdf}
\end{sourceblock}
因此表明

\begin{sourceblock}
\includegraphics[page=234,pagebox=mediabox,trim={64.39bp 590.93bp 64.39bp 1285.84bp},clip,width=0.94\textwidth,max height=0.38\textheight]{Pope.pdf}
\end{sourceblock}
其中 S 是速度导数偏度（参见方程（6.84）和（6.85））。

（因此，在各向同性湍流中， 的 729 个分量是
完全由速度导数偏度 S 确定，并且
K01莫戈洛夫时间表

使用方程式。 (6.95) 和 (6.99) 得到结果

\begin{sourceblock}
\includegraphics[page=234,pagebox=mediabox,trim={64.39bp 256.97bp 64.39bp 1664.79bp},clip,width=0.94\textwidth,max height=0.38\textheight]{Pope.pdf}
\end{sourceblock}
\begin{sourceblock}
\includegraphics[page=234,pagebox=mediabox,trim={64.39bp 142.98bp 64.39bp 1756.78bp},clip,width=0.94\textwidth,max height=0.38\textheight]{Pope.pdf}
\end{sourceblock}
\section*{6.4 傅里叶模态}\phantomsection\addcontentsline{toc}{section}{6.4 傅里叶模态}
6.4 傅立叶模态 对于各向同性湍流，源自纳维-斯托克斯方程的 \textasciitilde{}hrrnan-Howarth 方程（方程（6.75））充分描述了两点速度相关性的动力学。然而，它并没有提供能量级联所涉及过程的非常清晰的图景。通过检查波数空间中的纳维-斯托克斯方程可以获得一些进一步的见解。在本节中，我们将研究平均速度为零的均匀湍流的纳维-斯托克斯方程所规定的离散傅里叶模式的行为。

第一小节提供了将速度场表示为三维傅立叶级数的数学背景

\begin{sourceblock}
\includegraphics[page=235,pagebox=mediabox,trim={64.39bp 1272.84bp 64.39bp 668.92bp},clip,width=0.94\textwidth,max height=0.38\textheight]{Pope.pdf}
\end{sourceblock}
（这将附录 E 的内容扩展到三维矢量场。）然后从纳维-斯托克斯方程推导出傅立叶模式 \& ( \textasciitilde{} , tis) 的演化方程。最后，波数 K 处的动能平衡方程，

\begin{sourceblock}
\includegraphics[page=235,pagebox=mediabox,trim={64.39bp 972.88bp 64.39bp 964.88bp},clip,width=0.94\textwidth,max height=0.38\textheight]{Pope.pdf}
\end{sourceblock}
被推导并讨论。

除了它提供的关于能量级联的见解之外，研究纳维-斯托克斯方程的傅立叶表示还有其他动机。正如第 9 章所讨论的，均匀湍流的直接数值模拟（DNS）通常在波数空间中进行（例如，通过求解 ( , t ) ），而在第 11 章中应用于受到非常大的平均速度梯度的均匀湍流的快速畸变理论（RDT）也设置在波数空间中。

\subsection*{6.4.1 傅里叶级数表示}\phantomsection\addcontentsline{toc}{subsection}{6.4.1 傅里叶级数表示}
傅里叶级数方程的含义。 （6.102）可知湍流速度场是周期性的。因此，我们考虑物理空间中的立方体 0 C，其中边长 C 与湍流积分尺度 LI1 相比较大。那么速度场应该是周期性的，即

\begin{sourceblock}
\includegraphics[page=235,pagebox=mediabox,trim={64.39bp 209.97bp 64.39bp 1735.78bp},clip,width=0.94\textwidth,max height=0.38\textheight]{Pope.pdf}
\end{sourceblock}
对于所有整数向量 N。这种人为强加的周期性的影响
在 xl 方向上，傅里叶模的形式为

\begin{sourceblock}
\includegraphics[page=236,pagebox=mediabox,trim={64.39bp 1722.78bp 64.39bp 222.97bp},clip,width=0.94\textwidth,max height=0.38\textheight]{Pope.pdf}
\end{sourceblock}
和 sin(lconlx l)，对于整数 nl，其中 是最低波数：

\begin{sourceblock}
\includegraphics[page=236,pagebox=mediabox,trim={64.39bp 1579.80bp 64.39bp 365.95bp},clip,width=0.94\textwidth,max height=0.38\textheight]{Pope.pdf}
\end{sourceblock}
或者，以复数形式，傅里叶模式是

\begin{sourceblock}
\includegraphics[page=236,pagebox=mediabox,trim={64.39bp 1436.82bp 64.39bp 508.94bp},clip,width=0.94\textwidth,max height=0.38\textheight]{Pope.pdf}
\end{sourceblock}
对于正整数和负整数同样，在其他两个方向
nl.傅里叶模是 ei \textasciitilde{}n2g x2 和 ei \textasciitilde{} g n \textasciitilde{} x , 3，一般的三维模就是一维模的乘积。通过定义波数向量

\begin{sourceblock}
\includegraphics[page=236,pagebox=mediabox,trim={64.39bp 1161.85bp 64.39bp 783.90bp},clip,width=0.94\textwidth,max height=0.38\textheight]{Pope.pdf}
\end{sourceblock}
我们可以将一般傅立叶模简写为

\begin{sourceblock}
\includegraphics[page=236,pagebox=mediabox,trim={64.39bp 1019.87bp 64.39bp 925.88bp},clip,width=0.94\textwidth,max height=0.38\textheight]{Pope.pdf}
\end{sourceblock}
方程给出的傅里叶模式。 (6.109) 可以解释为
= I K I = 波数矢量的大小和方向 e。让我们
K K / K 是物理空间中 e 方向的坐标，即 s e x。然后我们观察到

\begin{sourceblock}
\includegraphics[page=236,pagebox=mediabox,trim={64.39bp 748.91bp 64.39bp 1189.85bp},clip,width=0.94\textwidth,max height=0.38\textheight]{Pope.pdf}
\end{sourceblock}
因此 eircmisX 在垂直于（常数 s）的平面上保持不变，而在其方向上则作为波数的一维傅里叶模式而变化。 作为示例，图 6,8 显示了具有 (n1, n3) (4,2,0) 的傅里叶模式。

傅里叶模式是正交的。为了用方程简单地表述这个性质，我们引入两个定义。首先，给定两个波数向量
和 K'，我们定义

\begin{sourceblock}
\includegraphics[page=236,pagebox=mediabox,trim={64.39bp 357.95bp 64.39bp 1539.81bp},clip,width=0.94\textwidth,max height=0.38\textheight]{Pope.pdf}
\end{sourceblock}
其次，我们用 ( ) 表示立方体 0 xi C 上的体积平均值。正交性属性为

\begin{sourceblock}
\includegraphics[page=237,pagebox=mediabox,trim={64.39bp 1194.85bp 64.39bp 109.44bp},clip,width=\textwidth,max height=0.72\textheight]{Pope.pdf}
\par\vspace{0.45\baselineskip}
\small 图 6.8。对应于 « = Kkp(4, 2,0) 的傅立叶模式草图。斜线显示波峰，其中 H(e**) = cos x * x 是单位。
\end{sourceblock}
对于周期函数 g(x)（例如，给定时间的速度分量），其傅里叶级数为

\begin{sourceblock}
\includegraphics[page=237,pagebox=mediabox,trim={64.39bp 939.88bp 64.39bp 1002.87bp},clip,width=0.94\textwidth,max height=0.38\textheight]{Pope.pdf}
\end{sourceblock}
其中总和超过无限数量的离散波数

g ( \textasciitilde{} ) 是波数处的复傅里叶系数 因为 g(x) 是
g ( \textasciitilde{} ) 实数，满足共轭对称性，

\begin{sourceblock}
\includegraphics[page=237,pagebox=mediabox,trim={64.39bp 690.91bp 64.39bp 1254.84bp},clip,width=0.94\textwidth,max height=0.38\textheight]{Pope.pdf}
\end{sourceblock}
其中星号表示复共轭。

给定 g(x)，可以根据正交条件确定傅立叶系数（方程 (6.1 11)）：

\begin{sourceblock}
\includegraphics[page=237,pagebox=mediabox,trim={64.39bp 348.96bp 64.39bp 1595.80bp},clip,width=0.94\textwidth,max height=0.38\textheight]{Pope.pdf}
\end{sourceblock}
定义操作符很方便

\begin{sourceblock}
\includegraphics[page=237,pagebox=mediabox,trim={64.39bp 119.39bp 64.39bp 1755.78bp},clip,width=0.94\textwidth,max height=0.38\textheight]{Pope.pdf}
\end{sourceblock}
这样前面的方程就可以简单地写成

\begin{sourceblock}
\includegraphics[page=238,pagebox=mediabox,trim={64.39bp 1724.78bp 64.39bp 218.97bp},clip,width=0.94\textwidth,max height=0.38\textheight]{Pope.pdf}
\end{sourceblock}
因此算子 3 ) 确定了波数的傅立叶模的系数
调用傅里叶表示的主要原因之一是导数所采用的形式。在等式中。 ( 6 . 1 6 ) 如果将 g(x) 替换为 ag/axj，我们得到

\begin{sourceblock}
\includegraphics[page=238,pagebox=mediabox,trim={64.39bp 1178.85bp 64.39bp 765.90bp},clip,width=0.94\textwidth,max height=0.38\textheight]{Pope.pdf}
\end{sourceblock}
物理空间中对 x j 的微分对应于波数空间中的乘法。

湍流速度场的傅立叶级数为

\begin{sourceblock}
\includegraphics[page=238,pagebox=mediabox,trim={64.39bp 948.88bp 64.39bp 993.87bp},clip,width=0.94\textwidth,max height=0.38\textheight]{Pope.pdf}
\end{sourceblock}
其中速度的傅立叶系数为

\begin{sourceblock}
\includegraphics[page=238,pagebox=mediabox,trim={64.39bp 776.90bp 64.39bp 1165.85bp},clip,width=0.94\textwidth,max height=0.38\textheight]{Pope.pdf}
\end{sourceblock}
傅里叶模 e 是非随机的并且在时间上是固定的。因此，湍流速度场 u(x, t) 的时间依赖性、随机性意味着 l e 傅里叶系数 \&(Kt,) 是时间相关的且随机的。由于精益 (\textasciitilde{}(xt),) 为零，因此可以从等式得出： (6.120) 意味着 (\&(Kt,) ) re 也为零。请注意，对于每个 K，\&(Kt,) 是一个满足 \textasciitilde{}njugates ymmetry 的复向量，即

\begin{sourceblock}
\includegraphics[page=238,pagebox=mediabox,trim={64.39bp 426.95bp 64.39bp 1516.81bp},clip,width=0.94\textwidth,max height=0.38\textheight]{Pope.pdf}
\end{sourceblock}
证明，对于整数

\begin{sourceblock}
\includegraphics[page=238,pagebox=mediabox,trim={64.39bp 174.98bp 64.39bp 1719.78bp},clip,width=0.94\textwidth,max height=0.38\textheight]{Pope.pdf}
\end{sourceblock}
由方程式可知（6.118）

\begin{sourceblock}
\includegraphics[page=239,pagebox=mediabox,trim={64.39bp 1722.78bp 64.39bp 222.97bp},clip,width=0.94\textwidth,max height=0.38\textheight]{Pope.pdf}
\end{sourceblock}
假设体积平均速度 (u(x, t))L 为零，表明
第 0 阶傅立叶模态的系数为零：

\begin{sourceblock}
\includegraphics[page=239,pagebox=mediabox,trim={64.39bp 1548.80bp 64.39bp 396.95bp},clip,width=0.94\textwidth,max height=0.38\textheight]{Pope.pdf}
\end{sourceblock}
证明涡度 o = V x 的傅立叶系数 \&(K) 为

\begin{sourceblock}
\includegraphics[page=239,pagebox=mediabox,trim={64.39bp 1415.82bp 64.39bp 527.93bp},clip,width=0.94\textwidth,max height=0.38\textheight]{Pope.pdf}
\end{sourceblock}
证明 K,\& (K)a, nd 相互正交。

一般来说，不可压缩速度场可以写成
无旋分量 V+ 和旋转分量 V x B，
其中 +(x) 是速度势，B(x) 是矢量势，
无散度 (V B 0)：

\begin{sourceblock}
\includegraphics[page=239,pagebox=mediabox,trim={64.39bp 1110.86bp 64.39bp 837.89bp},clip,width=0.94\textwidth,max height=0.38\textheight]{Pope.pdf}
\end{sourceblock}
对于周期性速度场，证明 为零，并获得
3,\{B(x)) 与 F x \{ o ( x ) l 之间的关系

\subsection*{6.4.2 傅立叶模态的演变}\phantomsection\addcontentsline{toc}{subsection}{6.4.2 傅立叶模态的演变}
在波数空间中，速度散度为

\begin{sourceblock}
\includegraphics[page=239,pagebox=mediabox,trim={64.39bp 672.91bp 64.39bp 1214.85bp},clip,width=0.94\textwidth,max height=0.38\textheight]{Pope.pdf}
\end{sourceblock}
使得连续性方程 V u 0 表明 i 垂直于

\begin{sourceblock}
\includegraphics[page=239,pagebox=mediabox,trim={64.39bp 555.93bp 64.39bp 1392.82bp},clip,width=0.94\textwidth,max height=0.38\textheight]{Pope.pdf}
\end{sourceblock}
有理由更详细地检查矢量（例如 ii）r 相对于波数的方向任何矢量 G 都可以分解为
集成电路。与 K 平行的分量 GA 和与 K 垂直的分量 G A I，即 G G GA，见图 6.9。其中 e 为单位向量
==K / K K 的方向，我们有

\begin{sourceblock}
\includegraphics[page=239,pagebox=mediabox,trim={64.39bp 225.97bp 64.39bp 1706.78bp},clip,width=0.94\textwidth,max height=0.38\textheight]{Pope.pdf}
\end{sourceblock}
\begin{sourceblock}
\includegraphics[page=239,pagebox=mediabox,trim={64.39bp 119.39bp 64.39bp 1793.77bp},clip,width=0.94\textwidth,max height=0.38\textheight]{Pope.pdf}
\end{sourceblock}
\begin{sourceblock}
\includegraphics[page=240,pagebox=mediabox,trim={64.39bp 1386.83bp 64.39bp 109.44bp},clip,width=\textwidth,max height=0.72\textheight]{Pope.pdf}
\par\vspace{0.45\baselineskip}
\small 图 6.9。显示分解 A A al 的草图（在二维波数空间中）。将任意向量 G 分解为平行于 x 的分量和垂直于 K 的分量。
\end{sourceblock}
因此，从 G I G G”，我们得到

\begin{sourceblock}
\includegraphics[page=240,pagebox=mediabox,trim={64.39bp 1119.86bp 64.39bp 813.90bp},clip,width=0.94\textwidth,max height=0.38\textheight]{Pope.pdf}
\end{sourceblock}
\begin{sourceblock}
\includegraphics[page=240,pagebox=mediabox,trim={64.39bp 993.87bp 64.39bp 903.89bp},clip,width=0.94\textwidth,max height=0.38\textheight]{Pope.pdf}
\end{sourceblock}
其中投影张量是

\begin{sourceblock}
\includegraphics[page=240,pagebox=mediabox,trim={64.39bp 844.89bp 64.39bp 1076.86bp},clip,width=0.94\textwidth,max height=0.38\textheight]{Pope.pdf}
\end{sourceblock}
主要观察结果是投影张量决定

pjkekto 是 到垂直于 This 的平面上的投影
i = K 投影张量在下面用于在波数空间中编写纳维-斯托克斯方程。

波数空间 \&(Kt,) 中速度矢量的演化方程是通过将算子 \{ (Eq- (6.1 16)) 逐项应用于纳维-斯托克斯方程来获得：

时间导数很简单

\begin{sourceblock}
\includegraphics[page=240,pagebox=mediabox,trim={64.39bp 256.97bp 64.39bp 1639.79bp},clip,width=0.94\textwidth,max height=0.38\textheight]{Pope.pdf}
\end{sourceblock}
而对于粘性项，应用方程 (6.118) 两次会产生

\begin{sourceblock}
\includegraphics[page=240,pagebox=mediabox,trim={64.39bp 119.39bp 64.39bp 1798.77bp},clip,width=0.94\textwidth,max height=0.38\textheight]{Pope.pdf}
\end{sourceblock}
其中 \$ ( \textasciitilde{} , t ) 是动态的傅里叶系数
-T,\{p(x,t)/p) 压力 (plp)，压力梯度项为

\begin{sourceblock}
\includegraphics[page=241,pagebox=mediabox,trim={64.39bp 1636.79bp 64.39bp 259.97bp},clip,width=0.94\textwidth,max height=0.38\textheight]{Pope.pdf}
\end{sourceblock}
目前，非线性对流项写为
从而定义傅里叶系数 G。结合这些结果我们得到
当方程。 (6.139) 乘以 Kj，左侧消失（由于连续性方程 0，方程 (6.128)），留下

\begin{sourceblock}
\includegraphics[page=241,pagebox=mediabox,trim={64.39bp 1115.86bp 64.39bp 820.90bp},clip,width=0.94\textwidth,max height=0.38\textheight]{Pope.pdf}
\end{sourceblock}
这是在波数空间中从纳维-斯托克斯方程获得的压力泊松方程，即：

通过求解方程。 (6.140) 对于\$，我们得到方程中的压力项。 (6.139)

\begin{sourceblock}
\includegraphics[page=241,pagebox=mediabox,trim={64.39bp 737.91bp 64.39bp 1184.85bp},clip,width=0.94\textwidth,max height=0.38\textheight]{Pope.pdf}
\end{sourceblock}
也就是说，压力项精确地平衡了 G，即在方程右侧剩余的方向上的分量。 (6.139)

电脑。是 -G， -G 垂直于

\begin{sourceblock}
\includegraphics[page=241,pagebox=mediabox,trim={64.39bp 408.95bp 64.39bp 1440.82bp},clip,width=0.94\textwidth,max height=0.38\textheight]{Pope.pdf}
\end{sourceblock}
该方程中的粘性项具有简单的作用。例如，考虑各向同性湍流衰减的最后阶段，其中雷诺数非常低，以至于对流相对于粘度的影响可以忽略不计。然后，根据指定的初始条件 ( P C 0, , 方程 (6.143)，忽略 G 项，有解

\begin{sourceblock}
\includegraphics[page=241,pagebox=mediabox,trim={64.39bp 119.39bp 64.39bp 1835.77bp},clip,width=0.94\textwidth,max height=0.38\textheight]{Pope.pdf}
\end{sourceblock}
因此，在最后的衰减周期中，每个傅立叶系数的演化独立于所有其他模式，随时间以指数方式衰减，其速率为 V K H\textasciitilde{} i.g h 波数模式比低波数模式衰减得更快。

用 ii(rc) 表示（隐式依赖 t），非线性对流项为

\begin{sourceblock}
\includegraphics[page=242,pagebox=mediabox,trim={64.39bp 1115.86bp 64.39bp 688.91bp},clip,width=0.94\textwidth,max height=0.38\textheight]{Pope.pdf}
\end{sourceblock}
（此开发中的六个步骤分别调用方程 (6.13 8)、(6.1 18)、(6.11 9)、(6.11 6)、(6.112) 和 (6.1 1 1。) 将此结果代入方程 (6.143) 中，我们得到 \&(Kt,) 的最终形式的演化方程：

\begin{sourceblock}
\includegraphics[page=242,pagebox=mediabox,trim={64.39bp 815.90bp 64.39bp 1073.86bp},clip,width=0.94\textwidth,max height=0.38\textheight]{Pope.pdf}
\end{sourceblock}
左侧仅在 K 处涉及 i3。相反，右侧在 rcf 和 K \textasciitilde{}su,ch h 处涉及 i3，即 rcf rc" rc；事实上，rc' = rc 和 rc" rc 的贡献为零。因此，在波数空间中，对流项是非线性和非局部的，涉及波数三元组 rc 和 K" 的相互作用，使得 rc" rc.

令 rca、rcb 和 rcc 为三个波数向量，使得

\begin{sourceblock}
\includegraphics[page=242,pagebox=mediabox,trim={64.39bp 356.95bp 64.39bp 1586.80bp},clip,width=0.94\textwidth,max height=0.38\textheight]{Pope.pdf}
\end{sourceblock}
并定义 a(t) = ii(rca, t)、b(t) = h(rcb, t) 和 c(t) = i3(rcc, t)。考虑
周期性速度场，在时间 t = 0 时具有非零傅里叶
仅在六个波数 \_+rca、rcb 和 +K, 处的系数。考虑
由欧拉控制的速度场的初始演化

\begin{sourceblock}
\includegraphics[page=243,pagebox=mediabox,trim={64.39bp 1730.78bp 64.39bp 159.98bp},clip,width=0.94\textwidth,max height=0.38\textheight]{Pope.pdf}
\end{sourceblock}
\begin{sourceblock}
\includegraphics[page=243,pagebox=mediabox,trim={64.39bp 1473.81bp 64.39bp 285.96bp},clip,width=0.94\textwidth,max height=0.38\textheight]{Pope.pdf}
\end{sourceblock}
(d) 有 24 种变化率非零的模态。素描他们的
波数空间中的位置。

\subsection*{6.4.3 傅立叶模态的动能}\phantomsection\addcontentsline{toc}{subsection}{6.4.3 傅立叶模态的动能}
对于所考虑的周期性情况，方程。 (6.146) 是波数空间中的纳维-斯托克斯方程。它是傅立叶系数 ( ) 的一组确定性常微分方程。为了用统计方法描述湍流，我们现在考虑各种方法。

由于（根据假设）平均速度 (U(x, t)) 处处为零，因此 (U(x,t)) 的傅立叶系数，即 (ii(u,t)) 也为零。要考虑的下一个最简单的统计量是两个傅立叶系数的协方差，即

\begin{sourceblock}
\includegraphics[page=243,pagebox=mediabox,trim={64.39bp 811.90bp 64.39bp 1133.86bp},clip,width=0.94\textwidth,max height=0.38\textheight]{Pope.pdf}
\end{sourceblock}
练习 6.18 表明，这些系数是不相关的，除非波数之和为零，即 u' u 0，或等效的 u' -u。因此，所有协方差信息都包含在

\begin{sourceblock}
\includegraphics[page=243,pagebox=mediabox,trim={64.39bp 539.93bp 64.39bp 1345.83bp},clip,width=0.94\textwidth,max height=0.38\textheight]{Pope.pdf}
\end{sourceblock}
很容易证明 Rij(u,t ) 是两点的傅里叶系数

\begin{sourceblock}
\includegraphics[page=243,pagebox=mediabox,trim={64.39bp 372.95bp 64.39bp 1515.81bp},clip,width=0.94\textwidth,max height=0.38\textheight]{Pope.pdf}
\end{sourceblock}
（参见练习 6.18）使得 R m ( xt) 具有傅立叶级数

\begin{sourceblock}
\includegraphics[page=243,pagebox=mediabox,trim={64.39bp 179.98bp 64.39bp 1700.79bp},clip,width=0.94\textwidth,max height=0.38\textheight]{Pope.pdf}
\end{sourceblock}
速度谱张量定义为

\begin{sourceblock}
\includegraphics[page=244,pagebox=mediabox,trim={64.39bp 1719.78bp 64.39bp 212.97bp},clip,width=0.94\textwidth,max height=0.38\textheight]{Pope.pdf}
\end{sourceblock}
其中 是连续波数变量。显然是谱张量
ii： 是两点相关性的傅里叶变换，因为等式： (6.154) 和

\begin{sourceblock}
\includegraphics[page=244,pagebox=mediabox,trim={64.39bp 1439.82bp 64.39bp 397.95bp},clip,width=0.94\textwidth,max height=0.38\textheight]{Pope.pdf}
\end{sourceblock}
在等式中设置 r 0 (6.154) 和 (6.156)，然后我们得到

\begin{sourceblock}
\includegraphics[page=244,pagebox=mediabox,trim={64.39bp 1241.84bp 64.39bp 653.92bp},clip,width=0.94\textwidth,max height=0.38\textheight]{Pope.pdf}
\end{sourceblock}
因此，R j ( , t ) 被确定为波数傅里叶模式对雷诺应力的贡献，而
是波数空间指定区域 IC 中的模式对 (uiuj) 的贡献。 Rij 和 cDij 的其他性质在练习 6.20 中给出。

特别令人感兴趣的是傅里叶模式的动能，定义为

\begin{sourceblock}
\includegraphics[page=244,pagebox=mediabox,trim={64.39bp 804.90bp 64.39bp 1133.86bp},clip,width=0.94\textwidth,max height=0.38\textheight]{Pope.pdf}
\end{sourceblock}
湍动能是

\begin{sourceblock}
\includegraphics[page=244,pagebox=mediabox,trim={64.39bp 598.92bp 64.39bp 1266.84bp},clip,width=0.94\textwidth,max height=0.38\textheight]{Pope.pdf}
\end{sourceblock}
将 \&,t) 标识为波数 K 对 k 的贡献。

耗散率 \textasciitilde{} ( ti)s 也与 h ( \textasciitilde{} , t )b,y 有关

\begin{sourceblock}
\includegraphics[page=244,pagebox=mediabox,trim={64.39bp 170.98bp 64.39bp 1698.79bp},clip,width=0.94\textwidth,max height=0.38\textheight]{Pope.pdf}
\end{sourceblock}
分别来自傅里叶模式的动能和耗散率

t) 的演化方程可以从 i i ( \textasciitilde{} , t ) 的演化方程推导出来

\begin{sourceblock}
\includegraphics[page=245,pagebox=mediabox,trim={64.39bp 1575.80bp 64.39bp 285.96bp},clip,width=0.94\textwidth,max height=0.38\textheight]{Pope.pdf}
\end{sourceblock}
\begin{sourceblock}
\includegraphics[page=245,pagebox=mediabox,trim={64.39bp 1406.82bp 64.39bp 432.94bp},clip,width=0.94\textwidth,max height=0.38\textheight]{Pope.pdf}
\end{sourceblock}
\{ \} 表示实部。当对所有 K 求和时，方程的左侧为： (6.161) 为 dk/dt，而右侧最后一项的总和为 -E。对于各向同性湍流 dk/dt 等于 -E，因此（可以直接确认）o:f T 之和为零：

\begin{sourceblock}
\includegraphics[page=245,pagebox=mediabox,trim={64.39bp 1143.86bp 64.39bp 802.90bp},clip,width=0.94\textwidth,max height=0.38\textheight]{Pope.pdf}
\end{sourceblock}
因此，术语 T(K,t ) 表示模式之间的能量转移。

方程之间有直接对应关系。 (6.161) 对于 h ( \textasciitilde{}t)，以及对于 f t) 的 Kirmin-H\textasciitilde{}warthe 方程，方程： （6.75）。它们包含本质上相同的信息，但表达方式不同。傅里叶模式公式的一个优点是，它提供了不同运动尺度下能量的清晰量化，并且获得了能量传输速率 T(K,t ) 的明确表达式，该能量传输速率在能量级联中起着核心作用。事实上，使用各向同性湍流的直接数值模拟，可以测量 T(K,t ) 和相关量（参见 Domaradzki (19 92)）。

证明两个速度傅里叶系数的协方差可以

可以表示为

\begin{sourceblock}
\includegraphics[page=245,pagebox=mediabox,trim={64.39bp 207.97bp 64.39bp 1588.80bp},clip,width=0.94\textwidth,max height=0.38\textheight]{Pope.pdf}
\end{sourceblock}
通过替换 x = xf - r，并且从同质中的事实来看，
位置，表明最后的结果可以重新表示为

\begin{sourceblock}
\includegraphics[page=246,pagebox=mediabox,trim={64.39bp 1668.79bp 64.39bp 224.97bp},clip,width=0.94\textwidth,max height=0.38\textheight]{Pope.pdf}
\end{sourceblock}
（提示：参见方程（E.22）。）因此，通过设置 -K，验证方程。 （6.153）。

通过换人

\begin{sourceblock}
\includegraphics[page=246,pagebox=mediabox,trim={64.39bp 1443.82bp 64.39bp 461.94bp},clip,width=0.94\textwidth,max height=0.38\textheight]{Pope.pdf}
\end{sourceblock}
\begin{sourceblock}
\includegraphics[page=246,pagebox=mediabox,trim={64.39bp 1339.83bp 64.39bp 589.93bp},clip,width=0.94\textwidth,max height=0.38\textheight]{Pope.pdf}
\end{sourceblock}
表明，在均匀湍流中，

\begin{sourceblock}
\includegraphics[page=246,pagebox=mediabox,trim={64.39bp 1179.85bp 64.39bp 764.90bp},clip,width=0.94\textwidth,max height=0.38\textheight]{Pope.pdf}
\end{sourceblock}
因此建立方程。 （6.154）。

由(Eq.(6.152))的定义可知

\begin{sourceblock}
\includegraphics[page=246,pagebox=mediabox,trim={64.39bp 959.88bp 64.39bp 976.88bp},clip,width=0.94\textwidth,max height=0.38\textheight]{Pope.pdf}
\end{sourceblock}
\begin{sourceblock}
\includegraphics[page=246,pagebox=mediabox,trim={64.39bp 871.89bp 64.39bp 1066.87bp},clip,width=0.94\textwidth,max height=0.38\textheight]{Pope.pdf}
\end{sourceblock}
由共轭对称性可知

\begin{sourceblock}
\includegraphics[page=246,pagebox=mediabox,trim={64.39bp 721.91bp 64.39bp 1211.85bp},clip,width=0.94\textwidth,max height=0.38\textheight]{Pope.pdf}
\end{sourceblock}
由不可压缩条件 ii(rc) 0 可知

\begin{sourceblock}
\includegraphics[page=246,pagebox=mediabox,trim={64.39bp 583.93bp 64.39bp 1364.83bp},clip,width=0.94\textwidth,max height=0.38\textheight]{Pope.pdf}
\end{sourceblock}
请注意，所有这些属性也适用于速度谱
设 Y 为任意常数向量，并定义 Y 得到

\begin{sourceblock}
\includegraphics[page=246,pagebox=mediabox,trim={64.39bp 315.96bp 64.39bp 1632.79bp},clip,width=0.94\textwidth,max height=0.38\textheight]{Pope.pdf}
\end{sourceblock}
证明 和 ( ) 都是半正定的，即

\begin{sourceblock}
\includegraphics[page=246,pagebox=mediabox,trim={64.39bp 172.98bp 64.39bp 1773.78bp},clip,width=0.94\textwidth,max height=0.38\textheight]{Pope.pdf}
\end{sourceblock}
\section*{6.5 速度谱}\phantomsection\addcontentsline{toc}{section}{6.5 速度谱}
6.22 证明 (Eq. (6.158)) 是实数、非负数，其中

\begin{sourceblock}
\includegraphics[page=247,pagebox=mediabox,trim={64.39bp 1709.78bp 64.39bp 227.97bp},clip,width=0.94\textwidth,max height=0.38\textheight]{Pope.pdf}
\end{sourceblock}
6.23 从 u(x) 的谱表示（方程（6.119））开始，显示
的光谱表示为

\begin{sourceblock}
\includegraphics[page=247,pagebox=mediabox,trim={64.39bp 1466.82bp 64.39bp 431.95bp},clip,width=0.94\textwidth,max height=0.38\textheight]{Pope.pdf}
\end{sourceblock}
因此显示关系

(== KK eRi(K)

\begin{sourceblock}
\includegraphics[page=247,pagebox=mediabox,trim={64.39bp 1148.86bp 64.39bp 754.90bp},clip,width=0.94\textwidth,max height=0.38\textheight]{Pope.pdf}
\end{sourceblock}
\begin{sourceblock}
\includegraphics[page=247,pagebox=mediabox,trim={64.39bp 861.89bp 64.39bp 1024.87bp},clip,width=0.94\textwidth,max height=0.38\textheight]{Pope.pdf}
\end{sourceblock}
6.5 速度谱
在上一节中，速度谱张量 ®;,(x,t) 被定义（对于均匀湍流）为两点速度相关性 R,(r) 的傅立叶变换。 （我们现在使用 « 表示连续波数变量，代替上面使用的 \%。）在第 6.5.1 节中，回顾了 t) 的属性并介绍了相关量；主要是能谱函数 E(x,t) 和一维光谱。柯尔莫哥洛夫假设对这些光谱在高波数（即在普遍平衡范围内）的形式有影响。这些含义在第 6.5.2 节中介绍，实验测量的光谱作为对假设的进一步检验。 6.6节用能谱函数描述了波数空间中的能量级联

\subsection*{6.5.1 定义和属性}\phantomsection\addcontentsline{toc}{subsection}{6.5.1 定义和属性}
速度谱张量在均匀湍流中，两点速度相关性和速度谱张量形成傅里叶变换对：

\begin{sourceblock}
\includegraphics[page=248,pagebox=mediabox,trim={64.39bp 1524.81bp 64.39bp 360.95bp},clip,width=0.94\textwidth,max height=0.38\textheight]{Pope.pdf}
\end{sourceblock}
\begin{sourceblock}
\includegraphics[page=248,pagebox=mediabox,trim={64.39bp 1395.82bp 64.39bp 486.94bp},clip,width=0.94\textwidth,max height=0.38\textheight]{Pope.pdf}
\end{sourceblock}
这里 \{ K \textasciitilde{}\textasciitilde{} ,2 rc,3\} 是（连续）波数向量；并且，为了简化符号，没有明确显示 Rij 和时间的依赖性。速度谱张量是一个复量，具有

\begin{sourceblock}
\includegraphics[page=248,pagebox=mediabox,trim={64.39bp 1134.86bp 64.39bp 759.90bp},clip,width=0.94\textwidth,max height=0.38\textheight]{Pope.pdf}
\end{sourceblock}
方程 (6.181) 源于 Rij(r) 的对称性和
Kj(r) 事实 ithat 是真实的；而方程。 (6.182) 是不可压缩性的结果（见练习 6.20）。另外 @ , j ( \textasciitilde{}is) 是半正定的，即

\begin{sourceblock}
\includegraphics[page=248,pagebox=mediabox,trim={64.39bp 819.90bp 64.39bp 1124.86bp},clip,width=0.94\textwidth,max height=0.38\textheight]{Pope.pdf}
\end{sourceblock}
对于所有向量 Y（参见练习 6.21）。然后得出对角线
( . e mi j) 的 aij(PC) 分量是实数且非负的，并且 .-t 因此迹也是如此：

\begin{sourceblock}
\includegraphics[page=248,pagebox=mediabox,trim={64.39bp 627.92bp 64.39bp 1316.83bp},clip,width=0.94\textwidth,max height=0.38\textheight]{Pope.pdf}
\end{sourceblock}
速度谱张量是一个值得考虑的有用量

@-r\textasciitilde{}( K) 原因（如第 6.4.3 节所示）它表示波数空间中的雷诺应力密度：即 是贡献（每单位 v \textasciitilde{} h m ein
mij(K) 波数空间）来自傅里叶模式 eiKmtoX 雷诺应力 (uiui)。特别是，在等式中设置r 0。 (6.180) 我们得到

\begin{sourceblock}
\includegraphics[page=248,pagebox=mediabox,trim={64.39bp 228.97bp 64.39bp 1659.79bp},clip,width=0.94\textwidth,max height=0.38\textheight]{Pope.pdf}
\end{sourceblock}
（请注意，其维度为 ( v e l \textasciitilde{} c i t \textasciitilde{} ) \textasciitilde{} / ( w a v e n u m boerr )e\textasciitilde{}q,u ivalent
其中包含的信息可以分为三个部分。首先，下标 和 给出物理空间中速度的方向。例如，完全属于域 u2(x)。第二，
m22(rc) 波数方向 rc/lrcl 给出了傅里叶模式的物理空间方向。第三，波数的大小决定了模式的长度尺度，即l 2n/1rcI（见图6.8）。

速度导数信息也包含在 ( ) 中，特别是，

\begin{sourceblock}
\includegraphics[page=249,pagebox=mediabox,trim={64.39bp 1402.82bp 64.39bp 474.94bp},clip,width=0.94\textwidth,max height=0.38\textheight]{Pope.pdf}
\end{sourceblock}
使得耗散率为

\begin{sourceblock}
\includegraphics[page=249,pagebox=mediabox,trim={64.39bp 1201.85bp 64.39bp 701.91bp},clip,width=0.94\textwidth,max height=0.38\textheight]{Pope.pdf}
\end{sourceblock}
mij(")（见练习 6.23）。 和 积分长度尺度之间的关系如下（方程（6.210）和（6.213））。

能谱函数 @-(rc) 作为向量的二阶张量函数，包含大量信息。能谱函数 E(K) 提供了更简单但不太完整的描述，它是标量的标量函数。

T 能谱函数是通过去除所有
(达尔赫(rc)

CJ方向信息。关于速度方向的信息通过考虑（一半）轨迹来消除，即 ( ) 通过对所有幅度的波数 rc 进行积分来消除关于傅立叶模式方向的信息 为了用数学方式表达这一点，
llcl K. 我们用 S(K) 表示波数空间中的球体，以原点为中心，半径和该球体表面的积分表示为
)\textasciitilde{}S(KTh)u。能量谱函数定义为

\begin{sourceblock}
\includegraphics[page=249,pagebox=mediabox,trim={64.39bp 399.95bp 64.39bp 1503.81bp},clip,width=0.94\textwidth,max height=0.38\textheight]{Pope.pdf}
\end{sourceblock}
或者，考虑到狄拉克δ函数的筛选性质（参见方程（C.1）），等价表达式为

\begin{sourceblock}
\includegraphics[page=249,pagebox=mediabox,trim={64.39bp 175.98bp 64.39bp 1711.78bp},clip,width=0.94\textwidth,max height=0.38\textheight]{Pope.pdf}
\end{sourceblock}
E(K) 的性质直接源自
@(K)：E[K) 是实数，非负数，对于负 K，根据式 (1) 是未定义的。 (6.188)，或根据方程为零。 （6.189）。 E(K)over all 的积分与 over all 的积分相同 因此，从等式： (6.185) 和
:mii(K) K. (6.18 7)，我们得到湍动能

\begin{sourceblock}
\includegraphics[page=250,pagebox=mediabox,trim={64.39bp 1518.81bp 64.39bp 404.95bp},clip,width=0.94\textwidth,max height=0.38\textheight]{Pope.pdf}
\end{sourceblock}
和耗散

\begin{sourceblock}
\includegraphics[page=250,pagebox=mediabox,trim={64.39bp 1374.83bp 64.39bp 528.93bp},clip,width=0.94\textwidth,max height=0.38\textheight]{Pope.pdf}
\end{sourceblock}
显然，E d\textasciitilde{} 是波数空间中无穷小壳层 I d\textasciitilde{} 中所有波数对 k 的贡献。

一般来说，包含比 E(K) 更多的信息；但是，
各向同性湍流中的 mij(K) 完全由 E(K) 决定。如果
cDij(K) 湍流是各向同性的， 中的方向信息只能取决于
K 上的 cDij(K)，并且在标量倍数内，唯一可以
aij 由 are 和 iciKj 组成。因此，在各向同性湍流中，Oij(K) 由下式给出

\begin{sourceblock}
\includegraphics[page=250,pagebox=mediabox,trim={64.39bp 959.88bp 64.39bp 984.88bp},clip,width=0.94\textwidth,max height=0.38\textheight]{Pope.pdf}
\end{sourceblock}
其中 A(K) 和 B(K) 是 的标量函数 这些标量函数很容易确定（参见练习 6.25），得出的结果是，在各向同性湍流中，速度谱张量为

\begin{sourceblock}
\includegraphics[page=250,pagebox=mediabox,trim={64.39bp 653.92bp 64.39bp 1197.85bp},clip,width=0.94\textwidth,max height=0.38\textheight]{Pope.pdf}
\end{sourceblock}
其中 是投影张量（方程（6.13 3））。

如果假设在原点处是解析的，则 E(K)v 为
对于小（见练习 6.26）。然而，K\textasciitilde{}K mij(K) 也可能是非解析的，E(K) 的变化如下（Saffman 1967）。在直接数值模拟中，可以获得 rc2 和 rc4 行为（Chasnov 1995）。有人建议（例如，Reynolds (1987)）网格湍流会产生
行为，但证据并不确凿。

表明

\begin{sourceblock}
\includegraphics[page=250,pagebox=mediabox,trim={64.39bp 119.39bp 64.39bp 1784.78bp},clip,width=0.94\textwidth,max height=0.38\textheight]{Pope.pdf}
\end{sourceblock}
（提示：认为方程（6.195）中的积分是各向同性的，即
bij 的标量倍数。）

表明应用不可压缩条件 0
至方程式(6.192) 产量

\begin{sourceblock}
\includegraphics[page=251,pagebox=mediabox,trim={64.39bp 1594.80bp 64.39bp 345.96bp},clip,width=0.94\textwidth,max height=0.38\textheight]{Pope.pdf}
\end{sourceblock}
使用练习（6.24）的结果，表明能量谱
对应于方程的函数(6.192) 是

\begin{sourceblock}
\includegraphics[page=251,pagebox=mediabox,trim={64.39bp 1407.82bp 64.39bp 532.93bp},clip,width=0.94\textwidth,max height=0.38\textheight]{Pope.pdf}
\end{sourceblock}
因此推导出方程。 （6.193）。

如果在原点处是解析的，则它具有形式的扩展
mij（电脑：）

\begin{sourceblock}
\includegraphics[page=251,pagebox=mediabox,trim={64.39bp 1211.85bp 64.39bp 721.91bp},clip,width=0.94\textwidth,max height=0.38\textheight]{Pope.pdf}
\end{sourceblock}
其中 是常数张量。表明不可压缩性
状态为 0 0，并且 Qij(K) 的正半定性（参见
等式。 (6.183)) 指示 dli (1 jk ) = 0。然后证明 (对于
小 n) 能谱函数是

\begin{sourceblock}
\includegraphics[page=251,pagebox=mediabox,trim={64.39bp 942.88bp 64.39bp 1004.87bp},clip,width=0.94\textwidth,max height=0.38\textheight]{Pope.pdf}
\end{sourceblock}
泰勒假设 在湍流的直接数值模拟中，可以提取速度谱张量和能量谱函数 E(K).T O
mij(K) 通过实验确定这些量需要测量所有 r 的两点速度相关性 Rij(r)，这显然是不可行的。然而，使用单个探头（例如热线风速计），可以近似地沿一条线测量 Rij(r)。

一种技术是使用“飞热线”。探针以恒定速度 V 沿着我们认为平行于 xl 轴的直线快速移动穿过湍流（即，探针沿 el 基矢量的方向移动）。如果探针在时间 t = 0 时位于 xo，则在时间 t

\begin{sourceblock}
\includegraphics[page=251,pagebox=mediabox,trim={64.39bp 237.97bp 64.39bp 1708.78bp},clip,width=0.94\textwidth,max height=0.38\textheight]{Pope.pdf}
\end{sourceblock}
探头测得的速度为

\begin{sourceblock}
\includegraphics[page=251,pagebox=mediabox,trim={64.39bp 119.39bp 64.39bp 1844.77bp},clip,width=0.94\textwidth,max height=0.38\textheight]{Pope.pdf}
\end{sourceblock}
T\#het 从 U( m) 获得的时间自协方差

\begin{sourceblock}
\includegraphics[page=252,pagebox=mediabox,trim={64.39bp 1617.80bp 64.39bp 213.97bp},clip,width=0.94\textwidth,max height=0.38\textheight]{Pope.pdf}
\end{sourceblock}
其中 r l Vs 是探针在 s 时间内移动的距离。如果湍流在 xl 方向上在统计上是均匀的，那么，在探头速度 V 趋于无穷大的假设极限下，我们得到
因此，在这些情况下，探针测量的时间自协方差产生 (xo,O) 处的空间自协方差。对于有限 V 的实际情况，显然方程： (6.203) 是一个随着 V 增加而改善的近似值。

一种更简单且更常见的技术是使用单个固定探针。该方法适用于统计平稳流，其中（在测量位置）湍流强度 u' 与平均速度 (U) 相比较小，我们将平均速度 (U) 视为 xl 方向。在以平均速度移动的坐标系中，探头以速度 elV -(U) e l ) 移动，因此适用飞热线分析（使用 rl -(U1)s）。

通过时间相关性来近似空间相关性，例如等式： （6.203）被称为泰勒假设（Taylor 1938）或冻结湍流近似。其准确性取决于流量的特性和所测量的统计数据。在 u'/(U1) 1 的网格湍流中，它相当准确，并且可以进行高阶校正（Lumley 1965）。另一方面，在自由剪切流中，许多实验（例如 Tong 和 Warhaft (1995)）已表明泰勒的假设失败。

一维光谱几乎所有现有的湍流光谱实验数据都来自固定热线测量。从测量中推导出来的量（使用泰勒假设）的形式为

\begin{sourceblock}
\includegraphics[page=252,pagebox=mediabox,trim={64.39bp 183.98bp 64.39bp 1755.78bp},clip,width=0.94\textwidth,max height=0.38\textheight]{Pope.pdf}
\end{sourceblock}
\begin{sourceblock}
\includegraphics[page=252,pagebox=mediabox,trim={64.39bp 119.39bp 64.39bp 1846.77bp},clip,width=0.94\textwidth,max height=0.38\textheight]{Pope.pdf}
\end{sourceblock}
其中f和g是纵向和横向自相关函数（方程（6.45））。这里我们定义并推导一维的性质
从 R j ( e l 1t ) 获得 Eij(rcl,t ) 并显示它们与光谱的关系
- 各向同性湍流中的 (PCt) 和 E(K,t )。下一节将进行实验
Eij(Kvl,t ) 上的 iJ 数据用于评估柯尔莫哥洛夫假设。

一维谱 Eij(rcl) 定义为 Rij(elrl) 的一维傅里叶变换的两倍：

\begin{sourceblock}
\includegraphics[page=253,pagebox=mediabox,trim={64.39bp 1419.82bp 64.39bp 475.94bp},clip,width=0.94\textwidth,max height=0.38\textheight]{Pope.pdf}
\end{sourceblock}
（此后，对 t 的依赖性不再明确显示。）对于对角分量 - 以 i == 2 为例 - R22(elr l ) 是实数，并且是 rl 的偶函数。因此 E22(rc1i)s 也是实数且偶数，因此方程： (6.206) 可以重写

\begin{sourceblock}
\includegraphics[page=253,pagebox=mediabox,trim={64.39bp 1128.86bp 64.39bp 767.90bp},clip,width=0.94\textwidth,max height=0.38\textheight]{Pope.pdf}
\end{sourceblock}
用反演公式

\begin{sourceblock}
\includegraphics[page=253,pagebox=mediabox,trim={64.39bp 965.88bp 64.39bp 936.88bp},clip,width=0.94\textwidth,max height=0.38\textheight]{Pope.pdf}
\end{sourceblock}
（参见方程（D.8）和（D.9））。添加 EV-(lcli)s 定义中的因子 2，以便（在方程（6.208）中设置 r1 0）我们得到

\begin{sourceblock}
\includegraphics[page=253,pagebox=mediabox,trim={64.39bp 750.91bp 64.39bp 1150.85bp},clip,width=0.94\textwidth,max height=0.38\textheight]{Pope.pdf}
\end{sourceblock}
一维谱与速度谱张量相关

\begin{sourceblock}
\includegraphics[page=253,pagebox=mediabox,trim={64.39bp 501.94bp 64.39bp 1401.82bp},clip,width=0.94\textwidth,max height=0.38\textheight]{Pope.pdf}
\end{sourceblock}
（参见练习 6.27）。应该理解的是，E 2 2 (\textasciitilde{} 1ha)s 来自平面 el rcl 中所有波数的贡献，因此波数
PC 傅里叶模式对 E22(rc1)c 的影响的 PC 幅度是相当大的
JPCJ比K大\textasciitilde{}。

一维谱 EI1(rc1) 与纵向自相关函数的关系为

\begin{sourceblock}
\includegraphics[page=253,pagebox=mediabox,trim={64.39bp 119.39bp 64.39bp 1810.77bp},clip,width=0.94\textwidth,max height=0.38\textheight]{Pope.pdf}
\end{sourceblock}
\begin{sourceblock}
\includegraphics[page=254,pagebox=mediabox,trim={64.39bp 1418.82bp 64.39bp 109.44bp},clip,width=\textwidth,max height=0.72\textheight]{Pope.pdf}
\par\vspace{0.45\baselineskip}
\small 图 6.10。显示径向坐标定义的波数空间草图
\end{sourceblock}
和纵向结构函数

\begin{sourceblock}
\includegraphics[page=254,pagebox=mediabox,trim={64.39bp 1153.85bp 64.39bp 768.90bp},clip,width=0.94\textwidth,max height=0.38\textheight]{Pope.pdf}
\end{sourceblock}
当 rcl 0 时，等式(6.211) 纵向积分尺度的产量

\begin{sourceblock}
\includegraphics[page=254,pagebox=mediabox,trim={64.39bp 984.88bp 64.39bp 914.88bp},clip,width=0.94\textwidth,max height=0.38\textheight]{Pope.pdf}
\end{sourceblock}
横向相关性也有类似的结果。

在各向同性湍流中，一维谱由能谱函数 E(rc) 确定。写方程。 (6.210) 为 Ell(rcl) 并代入方程。 (6.193) 因为我们得到

\begin{sourceblock}
\includegraphics[page=254,pagebox=mediabox,trim={64.39bp 643.92bp 64.39bp 1221.85bp},clip,width=0.94\textwidth,max height=0.38\textheight]{Pope.pdf}
\end{sourceblock}
积分是在固定 rcl 的平面上进行的，并且被积函数关于 rcl 轴径向对称。因此，引入径向坐标 rc，（参见

\begin{sourceblock}
\includegraphics[page=254,pagebox=mediabox,trim={64.39bp 427.95bp 64.39bp 1452.82bp},clip,width=0.94\textwidth,max height=0.38\textheight]{Pope.pdf}
\end{sourceblock}
并注意到 2n:rc, drc, 2nrc drc （对于固定 rcl），等式： (6.214) 可以重写

\begin{sourceblock}
\includegraphics[page=254,pagebox=mediabox,trim={64.39bp 248.97bp 64.39bp 1644.79bp},clip,width=0.94\textwidth,max height=0.38\textheight]{Pope.pdf}
\end{sourceblock}
这个公式强调了之前的观察结果，即 El (rcl) 包含的波数 rc 的贡献大于称为
rcl 的单调递减函数，因此无论 E (rc) 的形状如何，El1 在零波数处最大。

可以将上面的公式倒置（参见练习 6.28）以获得 Ell(ul) 的 E(K)i（对于各向同性湍流）：

\begin{sourceblock}
\includegraphics[page=255,pagebox=mediabox,trim={64.39bp 1557.80bp 64.39bp 339.96bp},clip,width=0.94\textwidth,max height=0.38\textheight]{Pope.pdf}
\end{sourceblock}
正如在各向同性湍流中，横向速度自相关函数 g(r) 由其纵向对应函数 g(r) 确定（方程（6.46））一样，EZ2(rc1i)s 也由 Ell(rcl) 确定。可以通过对方程（6.46）进行余弦傅立叶变换来获得该关系。 （6.46）（另见方程（6.211））：

\begin{sourceblock}
\includegraphics[page=255,pagebox=mediabox,trim={64.39bp 1256.84bp 64.39bp 690.91bp},clip,width=0.94\textwidth,max height=0.38\textheight]{Pope.pdf}
\end{sourceblock}
从方程。 (6.180) 证明

\begin{sourceblock}
\includegraphics[page=255,pagebox=mediabox,trim={64.39bp 992.87bp 64.39bp 910.88bp},clip,width=0.94\textwidth,max height=0.38\textheight]{Pope.pdf}
\end{sourceblock}
并从方程式(6.208) 证明

\begin{sourceblock}
\includegraphics[page=255,pagebox=mediabox,trim={64.39bp 807.90bp 64.39bp 1095.86bp},clip,width=0.94\textwidth,max height=0.38\textheight]{Pope.pdf}
\end{sourceblock}
因此验证方程。 (6.210)。

微分方程。 (6.216) 得到

\begin{sourceblock}
\includegraphics[page=255,pagebox=mediabox,trim={64.39bp 594.92bp 64.39bp 1302.84bp},clip,width=0.94\textwidth,max height=0.38\textheight]{Pope.pdf}
\end{sourceblock}
\begin{sourceblock}
\includegraphics[page=255,pagebox=mediabox,trim={64.39bp 469.94bp 64.39bp 1425.82bp},clip,width=0.94\textwidth,max height=0.38\textheight]{Pope.pdf}
\end{sourceblock}
因此验证方程。 (6.217)。使用方程式。 (6.221) 证明 Ell(rcl) 是
rcl 的单调递减函数，因此最大值为
证明在各向同性湍流中，

\begin{sourceblock}
\includegraphics[page=255,pagebox=mediabox,trim={64.39bp 119.39bp 64.39bp 1719.78bp},clip,width=0.94\textwidth,max height=0.38\textheight]{Pope.pdf}
\end{sourceblock}
表明，在各向同性湍流中，纵向积分尺度为

\begin{sourceblock}
\includegraphics[page=256,pagebox=mediabox,trim={64.39bp 1689.79bp 64.39bp 211.97bp},clip,width=0.94\textwidth,max height=0.38\textheight]{Pope.pdf}
\end{sourceblock}
表明，在各向同性湍流中，EZ2(rc1a) 和 E (rc) 的关系为

\begin{sourceblock}
\includegraphics[page=256,pagebox=mediabox,trim={64.39bp 1518.81bp 64.39bp 371.95bp},clip,width=0.94\textwidth,max height=0.38\textheight]{Pope.pdf}
\end{sourceblock}
\begin{sourceblock}
\includegraphics[page=256,pagebox=mediabox,trim={64.39bp 1400.82bp 64.39bp 496.94bp},clip,width=0.94\textwidth,max height=0.38\textheight]{Pope.pdf}
\end{sourceblock}
幂律谱 在检验柯尔莫哥洛夫假设（第 6.5.2 节）时，我们感兴趣的是：n 形式的幂律谱

\begin{sourceblock}
\includegraphics[page=256,pagebox=mediabox,trim={64.39bp 1116.86bp 64.39bp 830.89bp},clip,width=0.94\textwidth,max height=0.38\textheight]{Pope.pdf}
\end{sourceblock}
（在 rcl 的某个范围内），其中 C1 是常数，A 是归一化因子（例如，A (u2, ) L-1-P ）此类光谱在附录 G 中详细检查。 (6.228) 那么从方程可以得出： (6.217) 则 E(K) 是

\begin{sourceblock}
\includegraphics[page=256,pagebox=mediabox,trim={64.39bp 894.89bp 64.39bp 1051.87bp},clip,width=0.94\textwidth,max height=0.38\textheight]{Pope.pdf}
\end{sourceblock}
\begin{sourceblock}
\includegraphics[page=256,pagebox=mediabox,trim={64.39bp 771.90bp 64.39bp 1174.85bp},clip,width=0.94\textwidth,max height=0.38\textheight]{Pope.pdf}
\end{sourceblock}
并从方程式(6.218) EZ2(rcli)s

\begin{sourceblock}
\includegraphics[page=256,pagebox=mediabox,trim={64.39bp 637.92bp 64.39bp 1308.83bp},clip,width=0.94\textwidth,max height=0.38\textheight]{Pope.pdf}
\end{sourceblock}
\begin{sourceblock}
\includegraphics[page=256,pagebox=mediabox,trim={64.39bp 512.93bp 64.39bp 1431.82bp},clip,width=0.94\textwidth,max height=0.38\textheight]{Pope.pdf}
\end{sourceblock}
因此，三个光谱的幂律指数 p 是相同的，并且常数 C、C 和 Ci 是相关的。

光谱比较 在第 6.5.3 节中，我们引入了模型光谱（方程（6.246）），它提供了测量的湍流光谱的相当准确的表示。图 6.11 显示了该模型针对雷诺数 Ri 500 的各向同性湍流给出的各种光谱 E(K)、E 11(\textasciitilde{}1)a 和 nd EZ2(rc1)。

\begin{sourceblock}
\includegraphics[page=257,pagebox=mediabox,trim={64.39bp 1065.87bp 64.39bp 109.44bp},clip,width=\textwidth,max height=0.72\textheight]{Pope.pdf}
\end{sourceblock}
\subsection*{6.5.2 柯尔莫哥洛夫谱}\phantomsection\addcontentsline{toc}{subsection}{6.5.2 柯尔莫哥洛夫谱}
\begin{center}\small 图 6.11。 R 处各向同性湍流光谱的比较； = 500：实线，E(x)；虚线，点划线，E22(<1)。从模型谱来看，方程。 （6.246）。 （任意单位，）\end{center}
(i) 在波数范围的中心，所有光谱都表现出功率-
规律行为 p = 在此范围内，与方程一致。 (6.230)

和 (6.232)，Ell、E22 和 E 的值的比率为 1- - -
(ii) 在高波数下，光谱衰减比功率衰减得更快
K，与基础速度场无限一致
(iii) 在低波数下，当 K \textasciitilde{}In 时，E(K)t 结束为零。相比之下，一
维度光谱在零波数处最大。又是这个
说明了一个事实，即一维光谱包含贡献
的波数大于（见方程（6.216））。

(iv) 在低波数时，一维光谱 E l l 和 E22 为
比例2:l与整体长度比例一致
Lll 和 L22（见方程（6.51）和（6.213））。

柯尔莫哥罗谱 根据柯尔莫哥洛夫假设，在足够的任意湍流中
光谱采用特定的通用形式。这个结论以及柯尔莫哥洛夫谱的形式可以通过两种不同的途径获得。柯尔莫哥洛夫假设对二阶速度结构函数的影响在第 6.2 节中给出（例如方程（61.29）和（6.30））。第一条途径是获得柯尔莫哥洛夫谱作为结构函数的适当傅里叶变换。然而，我们遵循第二条路线，虽然不太严格，但更简洁：即将柯尔莫哥洛夫假设直接应用于光谱。

回想一下（对于任何足够高雷诺数的湍流），柯尔莫哥洛夫假设适用于小长度尺度上的速度场，特别是在由 ! 定义的通用平衡范围内。 lEI.I n 波数空间对应的范围为
根据局部各向同性假设，普适平衡范围内的速度统计是各向同性的。因此，对于 ICEI，速度谱张量 cDij(\textasciitilde{}) 由方程 1 的能谱函数 E(K) 给出。 (6.193)；以及 Ell(lcl)、E 2 2 ( \textasciitilde{} 1a)n,d E 之间的各向同性关系适用（参见方程（6.214）-（6.218））。

根据第一个相似性假设，属于普适平衡范围的速度统计 I 具有由 和 v 唯一确定的普适形式。因此，对于 KEI，E(K) 是 K、E 和 v 的普适函数。利用 和 v 对和 E(K) 进行无维化，简单的量纲分析表明，这个普适关系可以写成

\begin{sourceblock}
\includegraphics[page=258,pagebox=mediabox,trim={64.39bp 753.90bp 64.39bp 1139.86bp},clip,width=0.94\textwidth,max height=0.38\textheight]{Pope.pdf}
\end{sourceblock}
其中 是通用无量纲函数 Kolmogorov

\textasciitilde{} ( K V ) 谱函数。或者，如果 和 用于对 E 无量纲化，则关系为

\begin{sourceblock}
\includegraphics[page=258,pagebox=mediabox,trim={64.39bp 531.93bp 64.39bp 1408.82bp},clip,width=0.94\textwidth,max height=0.38\textheight]{Pope.pdf}
\end{sourceblock}
其中 是补偿后的柯尔莫哥洛夫谱函数。这些通用函数通过以下方式相关联

\begin{sourceblock}
\includegraphics[page=258,pagebox=mediabox,trim={64.39bp 341.96bp 64.39bp 1599.80bp},clip,width=0.94\textwidth,max height=0.38\textheight]{Pope.pdf}
\end{sourceblock}
和等式。 (6.233)和(6.234)适用于KE\textasciitilde{}w，hch对应
通用平衡范围

惯性子区 耗散区

图 6.12。雷诺数非常高时的波数（对数尺度）

显示各种范围。

即，q << i？ lo，或更准确地说 lD < 1 i？ <我？波数空间中对应的范围为KDI，见图6.12；或者，就 I C \textasciitilde{} 而言，
集成电路

\begin{sourceblock}
\includegraphics[page=259,pagebox=mediabox,trim={64.39bp 1107.86bp 64.39bp 839.89bp},clip,width=0.94\textwidth,max height=0.38\textheight]{Pope.pdf}
\end{sourceblock}
\begin{sourceblock}
\includegraphics[page=259,pagebox=mediabox,trim={64.39bp 960.88bp 64.39bp 912.88bp},clip,width=0.94\textwidth,max height=0.38\textheight]{Pope.pdf}
\end{sourceblock}
在惯性子域中，根据第二相似性假设，E(K)具有由独立于v唯一确定的通用形式。 (6.234) 对于 E(Ic)，v 仅通过 q 进入。因此，该假设意味着，当其参数 lcq 趋于零（ i , \textasciitilde{} q 1, 参见方程 (6.237)）时，函数 Y 变得独立于其参数，即，它趋向于常数 C。因此，第二个相似性假设预测，在惯性子区内，能谱函数为

\begin{sourceblock}
\includegraphics[page=259,pagebox=mediabox,trim={64.39bp 549.93bp 64.39bp 1398.82bp},clip,width=0.94\textwidth,max height=0.38\textheight]{Pope.pdf}
\end{sourceblock}
( e mE, q. (6.234) with Y C.) 这就是著名的柯尔莫哥洛夫谱，C 是通用柯尔莫哥洛夫常数。实验数据支持 C 1.5 值（参见下图 6.1 7 和 Sreenivasan (1995)）。

根据假设，在惯性子域中，Oij(IC) 是各向同性张量函数，E(K) 是幂律谱 i f Eq。 (6.228) 与 p ) 因此，如第 6.5.1 节所示，给出一维谱

\begin{sourceblock}
\includegraphics[page=259,pagebox=mediabox,trim={64.39bp 169.98bp 64.39bp 1777.78bp},clip,width=0.94\textwidth,max height=0.38\textheight]{Pope.pdf}
\end{sourceblock}
\begin{sourceblock}
\includegraphics[page=259,pagebox=mediabox,trim={64.39bp 119.39bp 64.39bp 1843.77bp},clip,width=0.94\textwidth,max height=0.38\textheight]{Pope.pdf}
\end{sourceblock}
\begin{sourceblock}
\includegraphics[page=260,pagebox=mediabox,trim={64.39bp 1726.78bp 64.39bp 162.98bp},clip,width=0.94\textwidth,max height=0.38\textheight]{Pope.pdf}
\end{sourceblock}
\begin{sourceblock}
\includegraphics[page=260,pagebox=mediabox,trim={64.39bp 1639.79bp 64.39bp 296.96bp},clip,width=0.94\textwidth,max height=0.38\textheight]{Pope.pdf}
\end{sourceblock}
（参见方程（6.228）-（6.232））。

幂律谱的一些性质在附录 G 中给出。E l l ( q ) (Eq. (6.240)) 的形式与二阶速度结构函数的形式之间有直接的对应关系

\begin{sourceblock}
\includegraphics[page=260,pagebox=mediabox,trim={64.39bp 1375.83bp 64.39bp 571.93bp},clip,width=0.94\textwidth,max height=0.38\textheight]{Pope.pdf}
\end{sourceblock}
（方程（6.30））在惯性范围内。 p 和 q 的幂 (E(K) K-P, DLL(r)

rq) 与 p 相关
= 1 q = 1 2 m 和常数（到一个极好的

\begin{sourceblock}
\includegraphics[page=260,pagebox=mediabox,trim={64.39bp 1156.85bp 64.39bp 733.91bp},clip,width=0.94\textwidth,max height=0.38\textheight]{Pope.pdf}
\end{sourceblock}
（参见方程（(3.25)）。

\subsection*{6.5.3 模型谱}\phantomsection\addcontentsline{toc}{subsection}{6.5.3 模型谱}
在检查实验数据以进一步检验柯尔莫哥洛夫假设之前，我们介绍一个用于比较的简单模型谱。能谱函数的模型是

\begin{sourceblock}
\includegraphics[page=260,pagebox=mediabox,trim={64.39bp 736.91bp 64.39bp 1187.85bp},clip,width=0.94\textwidth,max height=0.38\textheight]{Pope.pdf}
\end{sourceblock}
其中 和 f 是指定的无量纲函数。功能
f L f L 决定了能量包含范围的形状，并且对于大KL 趋于统一。类似地，f，决定了耗散区的形状，对于较小的KV，它趋于统一。在惯性子区内，f，和f，基本上是统一的，因此恢复了具有常数C的柯尔莫哥洛夫谱。

的规格是

\begin{sourceblock}
\includegraphics[page=260,pagebox=mediabox,trim={64.39bp 289.96bp 64.39bp 1599.80bp},clip,width=0.94\textwidth,max height=0.38\textheight]{Pope.pdf}
\end{sourceblock}
其中po取2，是一个正常数。显然倾向于
对于大 KL，CL f L 统一，而指数 +po 导致 E(K)v 变化为
KPO K 2 表示小 KL。（使用替代选择 po 4，方程（6.247）被称为

\begin{sourceblock}
\includegraphics[page=261,pagebox=mediabox,trim={64.39bp 1313.84bp 64.39bp 109.44bp},clip,width=\textwidth,max height=0.72\textheight]{Pope.pdf}
\par\vspace{0.45\baselineskip}
\small 图 6.13。 R, = 500 的模型谱（方程（6.246））通过柯尔莫哥洛夫尺度归一化。
\end{sourceblock}
f 的规范是

\begin{sourceblock}
\includegraphics[page=261,pagebox=mediabox,trim={64.39bp 1069.86bp 64.39bp 868.89bp},clip,width=0.94\textwidth,max height=0.38\textheight]{Pope.pdf}
\end{sourceblock}
其中 和 c 是正常数。请注意，对于 c, 0，这减少为

\begin{sourceblock}
\includegraphics[page=261,pagebox=mediabox,trim={64.39bp 914.88bp 64.39bp 1032.87bp},clip,width=0.94\textwidth,max height=0.38\textheight]{Pope.pdf}
\end{sourceblock}
由于速度场是无限可微的，因此对于较大的 K，能谱函数的衰减速度比 的任何幂都要快（参见附录 G）。因此出现指数衰减（如 Kraichnan (1959) 所建议的）。一些实验支持 5.2 的指数形式（参见 Saddoughi 和 Veeravalli (1994)）。然而，对于小 icq 来说，简单指数（方程（6.249））偏离单位太快，并且 的值被限制为 2.1（见练习 6.33）。这些缺陷可以通过等式来弥补。 （6.248）。

对于 k、E 和 v 的指定值，模型谱由方程 1 确定。 (6.246)-(6.248)，C 1.5 和 5.2。或者，无量纲模型谱由 Rid 的指定值唯一确定。 常数 和 分别由 E(K) 和 积分到 k 和 的要求确定： 在高雷诺数下

\textasciitilde{} V K \textasciitilde{} E ( K ) E，它们的值为 6.78 和 c 0.40（见练习 6.32）。对于各向同性

CL tl 湍流，一维谱 E 和 E2\&c2)a 的相应模型可从方程 (1) 中获得。 (6.2 16)-(6.2 18)。

图 6.13 是模型谱的双对数图（采用 Kolmogorov
惯性子区中的 E[K) 很明显，大 K 处的衰减也是如此。

表明，在非常高的雷诺数下，模型的积分
所有产量的频谱（方程（6.246））

\begin{sourceblock}
\includegraphics[page=262,pagebox=mediabox,trim={64.39bp 1478.81bp 64.39bp 423.95bp},clip,width=0.94\textwidth,max height=0.38\textheight]{Pope.pdf}
\end{sourceblock}
表明，由方程给出。 （6.247）（与 2），积分

\begin{sourceblock}
\includegraphics[page=262,pagebox=mediabox,trim={64.39bp 1293.84bp 64.39bp 568.93bp},clip,width=0.94\textwidth,max height=0.38\textheight]{Pope.pdf}
\end{sourceblock}
因此表明，对于 C = 1.5，高雷诺数渐近线

\begin{sourceblock}
\includegraphics[page=262,pagebox=mediabox,trim={64.39bp 1114.86bp 64.39bp 826.90bp},clip,width=0.94\textwidth,max height=0.38\textheight]{Pope.pdf}
\end{sourceblock}
\subsection*{6.5.4 耗散谱}\phantomsection\addcontentsline{toc}{subsection}{6.5.4 耗散谱}
在本节和接下来的四小节中，实验数据、柯尔莫哥洛夫假设和模型谱用于检查湍流中的速度谱。在大多数相关实验中，都会调用泰勒假设以获得一维光谱的测量结果

图 6.14 是用 Kolmogorov 标度标注的 E l l ( \textasciitilde{} l )p,l 测量值的汇编。与 E(K)( 方程 (6.233)) 的情况一样，柯尔莫哥洛夫假设意味着缩放谱 E \textasciitilde{} \textasciitilde{} ( K \textasciitilde{} ) / ( EisV a \textasciitilde{}un) i'v/e\textasciitilde{}rs al
K I \textasciitilde{}at 的 cpll 函数，足够高的雷诺数，并且对于 >

图 6.14 中所示的 \textasciitilde{}1 KEI 数据来自许多不同的流，泰勒尺度雷诺数从 23 到 3,180。可以看出，对于 > 0.1，所有数据都位于一条曲线上。高雷诺数数据在 rclq < 0.1 时表现出幂律行为，幂律区域的范围通常随着 RA 的增加而增加。因此，数据与 E I I ( \textasciitilde{} l ) )/ ( \textasciitilde{}b\textasciitilde{}ein g lclq 的通用函数 > 一致，但偏离通用
K\textasciitilde{} KEI，图 6.14 中由含能区 < The 引起的行为
凯凯。模型光谱（也显示在图 6.14 中的各种 RA）似乎相当准确地表示了数据。

\begin{sourceblock}
\includegraphics[page=263,pagebox=mediabox,trim={64.39bp 539.93bp 64.39bp 298.76bp},clip,width=\textwidth,max height=0.72\textheight]{Pope.pdf}
\par\vspace{0.45\baselineskip}
\small 图 6.14。 R i = 30、70、130、300、600 和 1,500（线）的一维纵向速度谱（符号）和模型谱（方程（6.246））的测量。实验数据取自 Saddoughi 和 Veeravalli (1994)，其中给出了对各种实验的参考。对于每个实验，密钥中的最终数字就是 R 的值。
\end{sourceblock}
缩放比例如图 6.15 中的线性对数图所示，它强调了耗散区。对于 lclq > 0.1，网格湍流 (Ri 60) 和湍流边界层中的测量结果非常一致
ss (RIM 600)，再次支持高波数规格的通用性

\begin{sourceblock}
\includegraphics[page=264,pagebox=mediabox,trim={64.39bp 1087.86bp 64.39bp 109.44bp},clip,width=\textwidth,max height=0.72\textheight]{Pope.pdf}
\par\vspace{0.45\baselineskip}
\small 图 6.15。补偿一维速度谱。 Comte-Bellot 和 Corrsin (1971) 在 R, \textasciitilde{} 60（三角形）处的网格湍流中进行的测量，以及 Saddoughi 和 Veeravalli (1994) 在 R, * 600（圆圈）处的湍流边界层中进行的测量。实线，模型谱方程。 (6.246) 对于 R，； = 600；虚线，指数谱(6.253)；点划线，Pao 谱，方程。 （6.254）。
\end{sourceblock}
最高波数处光谱的指数衰减。同样，模型谱准确地代表了数据。

图 6.15 还显示了从 f (lcq) 的两个替代模型推导出来的一维谱。这些是指数

\begin{sourceblock}
\includegraphics[page=264,pagebox=mediabox,trim={64.39bp 619.92bp 64.39bp 1311.83bp},clip,width=0.94\textwidth,max height=0.38\textheight]{Pope.pdf}
\end{sourceblock}
其中由方程式给出。 (6.258) 和 Pao 谱

\begin{sourceblock}
\includegraphics[page=264,pagebox=mediabox,trim={64.39bp 493.94bp 64.39bp 1453.82bp},clip,width=0.94\textwidth,max height=0.38\textheight]{Pope.pdf}
\end{sourceblock}
（参见 Pao (1965) 和第 6.6 节）。从图 6.15 可以明显看出，这些替代方案不能像模型谱那样代表数据。

确定模型谱准确地描述了耗散区后，我们现在用它来量化耗散运动的规模。图 6.16 显示了根据以下公式计算的耗散谱 D(K)：

\textasciitilde{} V I C \textasciitilde{} E ( K ) Riv 600 的模型，也是累积耗散

\begin{sourceblock}
\includegraphics[page=264,pagebox=mediabox,trim={64.39bp 119.39bp 64.39bp 1801.77bp},clip,width=0.94\textwidth,max height=0.38\textheight]{Pope.pdf}
\end{sourceblock}
\begin{sourceblock}
\includegraphics[page=265,pagebox=mediabox,trim={64.39bp 1137.86bp 64.39bp 109.44bp},clip,width=\textwidth,max height=0.72\textheight]{Pope.pdf}
\par\vspace{0.45\baselineskip}
\small 图 6.16。模型谱方程对应的耗散谱（实线）和累积耗散（虚线）。 (6.246) 对于 R； = 600。 ¢ = 27/x 是波数 xk 对应的波长。
\end{sourceblock}
横坐标显示波数和相应的波长 l 27r/rc，两者均通过柯尔莫哥洛夫尺度 q 归一化。从这些曲线获得的特征波数和波长在

\begin{center}\small 表 6.1。可以看出，耗散谱的峰值出现在 xy = 0.26 处，对应于 cent/y = 24，而质心（其中 £0）= 48）出现在 Ky \textasciitilde{} 0.34 处，对应于 £/y \textasciitilde{} 18。因此，造成大部分耗散的运动（0.1 < Ky < 0.75，或 60 >> 8）远大于 Kolmogorov 尺度。 （有\end{center}
\begin{sourceblock}
\includegraphics[page=265,pagebox=mediabox,trim={64.39bp 238.97bp 64.39bp 1345.83bp},clip,width=\textwidth,max height=0.72\textheight]{Pope.pdf}
\end{sourceblock}
从模型谱积分得到（方程（6.246））为

\begin{sourceblock}
\includegraphics[page=265,pagebox=mediabox,trim={64.39bp 119.39bp 64.39bp 1810.77bp},clip,width=0.94\textwidth,max height=0.38\textheight]{Pope.pdf}
\par\vspace{0.45\baselineskip}
\small 表 6.1。耗散谱的特征波数和长度尺度（基于 Ry = 600 时的模型谱方程（6.246））
\end{sourceblock}
\begin{sourceblock}
\includegraphics[page=266,pagebox=mediabox,trim={64.39bp 1425.82bp 64.39bp 261.97bp},clip,width=\textwidth,max height=0.72\textheight]{Pope.pdf}
\end{sourceblock}
证明，如果 f，由指数方程给出。 (6.253)，那么
方程中的积分(6.256) 是

\begin{sourceblock}
\includegraphics[page=266,pagebox=mediabox,trim={64.39bp 1231.84bp 64.39bp 671.91bp},clip,width=0.94\textwidth,max height=0.38\textheight]{Pope.pdf}
\end{sourceblock}
因此证明 的高雷诺数渐近线是

\begin{sourceblock}
\includegraphics[page=266,pagebox=mediabox,trim={64.39bp 1098.86bp 64.39bp 847.89bp},clip,width=0.94\textwidth,max height=0.38\textheight]{Pope.pdf}
\end{sourceblock}
对于 C 1.5。确认 Pao 谱（Eq. (6.254)）满足

\subsection*{6.5.5 惯性子区}\phantomsection\addcontentsline{toc}{subsection}{6.5.5 惯性子区}
第二个柯尔莫哥洛夫假设预测惯性子区内的谱。图 6.14 中明显的幂律行为最好通过绘制柯尔莫哥洛夫方程的补偿谱 E \textasciitilde{} \textasciitilde{} ( I cF\textasciitilde{}or) 来检验
c 2 I 3 'I3 假设预测该量在惯性范围内采用常数值C1（见式（6.240））。高雷诺数边界层的补偿光谱与图 6.17 中的 Kolmogorov 预测进行了比较。可以看出，在二十个波数的范围内，数据与预测值的误差在 20\% 以内
IC 增加了 2,000 倍以上。

当然，模型谱的构建是为了产生惯性子区内的柯尔莫哥洛夫行为，这种行为在图 6.17 中很明显。图 6.17 还提供了这种明显各向异性湍流中局部各向同性的一些证据：对于 > 2 x Ep 和 非常相似，
I C \textasciitilde{} \textasciitilde{} E33 和（根据局部各向同性预测）其补偿光谱的平台值是 Ell 的平台值的 4（方程（6.243））。更直接的测试（包括

\begin{sourceblock}
\includegraphics[page=267,pagebox=mediabox,trim={64.39bp 570.93bp 64.39bp 267.77bp},clip,width=\textwidth,max height=0.72\textheight]{Pope.pdf}
\par\vspace{0.45\baselineskip}
\small 图 6.17。在 R 为 1,450 的湍流边界层中测量的补偿一维光谱。实线，实验数据 Saddoughi 和 Veeravalli (1994)；虚线，来自方程式的模型光谱。 (6.246)；长虚线C和C；对应于柯尔莫哥洛夫惯性范围谱。 （对于 E,;、E22 和 E33，模型光谱分别为 Rz, = 1,450、690 和 910，对应于测量值
\end{sourceblock}
假设各向同性，根据 El1(rcl) 计算得出，方程： (6.218)。 Saddoughi 和 Veeravalli (1994) 进行了这项测试，发现测量值和计算值在整个过程中相差不超过 10\%。

\begin{center}\small 表 6.2。能谱的特征波数和长度尺度（基于 R; = 600 处的模型谱方程（6.246））\end{center}
\begin{sourceblock}
\includegraphics[page=268,pagebox=mediabox,trim={64.39bp 884.89bp 64.39bp 262.97bp},clip,width=\textwidth,max height=0.72\textheight]{Pope.pdf}
\end{sourceblock}
\subsection*{6.5.6 含能区}\phantomsection\addcontentsline{toc}{subsection}{6.5.6 含能区}
适合归一化的尺度是湍动能 k 和纵向积分长度尺度 L1l。对于各向同性湍流，E(K)h 作为积分属性

\begin{sourceblock}
\includegraphics[page=268,pagebox=mediabox,trim={64.39bp 633.92bp 64.39bp 1268.84bp},clip,width=0.94\textwidth,max height=0.38\textheight]{Pope.pdf}
\end{sourceblock}
\begin{sourceblock}
\includegraphics[page=268,pagebox=mediabox,trim={64.39bp 479.94bp 64.39bp 1421.82bp},clip,width=0.94\textwidth,max height=0.38\textheight]{Pope.pdf}
\end{sourceblock}
通过这些缩放，图 6.18 显示了肋 60 处网格湍流中 E(K)i 的测量结果，以及 R 的模型谱； 60 和 1,000。可以看出（通过这种缩放）频谱的形状不会随着雷诺数的变化而剧烈变化，并且该模型提供了数据的合理表示。图 6.18 中还显示了模型频谱 Eq. (6.246) 与 po 4 （对于小 KL 给出 E(K) rc4）。与po 2相比，形状差别不大，相差10\%

\begin{sourceblock}
\includegraphics[page=269,pagebox=mediabox,trim={64.39bp 1289.84bp 64.39bp 109.44bp},clip,width=\textwidth,max height=0.72\textheight]{Pope.pdf}
\par\vspace{0.45\baselineskip}
\small 图 6.18。由 k 和 L,, 归一化的各向同性湍流中的能谱函数。 Comte-Bellot 和 Corrsin (1971) 的符号、网格湍流实验：O,Ra = 71;0,R, = 65;A,R,; = 61. 线、模型谱、方程。 (6.246)：固体，pp = 2，R = 60；虚线，po = 2，Ry = 1,000；点划线 py = 4, R, = 60。
\end{sourceblock}
\begin{sourceblock}
\includegraphics[page=269,pagebox=mediabox,trim={64.39bp 500.94bp 64.39bp 826.90bp},clip,width=\textwidth,max height=0.72\textheight]{Pope.pdf}
\par\vspace{0.45\baselineskip}
\small 图 6.19。模型光谱中相对于波数 x 和波长 ¢ = 2x/x 的累积湍动能。
\end{sourceblock}
\begin{sourceblock}
\includegraphics[page=269,pagebox=mediabox,trim={64.39bp 240.97bp 64.39bp 1661.79bp},clip,width=0.94\textwidth,max height=0.38\textheight]{Pope.pdf}
\end{sourceblock}
\begin{sourceblock}
\includegraphics[page=270,pagebox=mediabox,trim={64.39bp 1436.82bp 64.39bp 109.44bp},clip,width=\textwidth,max height=0.72\textheight]{Pope.pdf}
\par\vspace{0.45\baselineskip}
\small 图 6.20。各种雷诺数的模型谱，按 (a) k 和 Li 以及 (b) Kolmogorov 尺度进行缩放。
\end{sourceblock}
< < 6Lll。在此基础上，我们采用表征 LL 的长度尺度，
eEr含能运动为Lll和L\textasciitilde{}\textasciitilde{}。

\subsection*{6.5.7 雷诺数的影响}\phantomsection\addcontentsline{toc}{subsection}{6.5.7 雷诺数的影响}
图 6.20(a) 显示了一系列雷诺数下通过 k 和 Lll 归一化的模型谱。可以看出，光谱的能量范围（例如 0.1 < rcLll < 10）非常相似，而随着 RA 的增加，-5 区域的范围增加，并且耗散区（光谱滚降的地方）移动到更高的 rcLll 值。

图 6+20(b)N 显示了相同的光谱，但通过 Kolmogorov 尺度进行了归一化。耗散区（例如 uq 0.1）非常相似，而随着 Rl 的增加，含能区移动到较低的 icy 值。

图 6.21 对比了高雷诺数 (Riv 1,000) 和低雷诺数 (Rlv 30) 能量和耗散谱。与此类对数线性图的通常做法一样，将光谱乘以 u，以便曲线 KE(Kr)e 下方的面积表示能量。即波数范围（K,, ）内的能量

\begin{sourceblock}
\includegraphics[page=270,pagebox=mediabox,trim={64.39bp 428.95bp 64.39bp 1473.81bp},clip,width=0.94\textwidth,max height=0.38\textheight]{Pope.pdf}
\end{sourceblock}
通过采用柯尔莫哥洛夫标度，高雷诺数谱包含更多能量（即，更大的k/u2值）。因此，能谱按不同的数值因子进行缩放，以便可以在同一图上对它们进行比较。图 6.21 的重要观察结果是，在低雷诺数时，能量谱和耗散谱显着重叠：没有明显的尺度分离。

\begin{sourceblock}
\includegraphics[page=271,pagebox=mediabox,trim={64.39bp 1249.84bp 64.39bp 109.44bp},clip,width=\textwidth,max height=0.72\textheight]{Pope.pdf}
\par\vspace{0.45\baselineskip}
\small 图 6.21。通过 Kolmogorov 尺度在 R, = 1,000（实线）和 = 30（虚线）处标准化的模型能量和耗散谱。 （注意 E(x) 的缩放比例。）
\end{sourceblock}
\begin{sourceblock}
\includegraphics[page=271,pagebox=mediabox,trim={64.39bp 677.92bp 64.39bp 805.90bp},clip,width=\textwidth,max height=0.72\textheight]{Pope.pdf}
\par\vspace{0.45\baselineskip}
\small 图 6.22。波数大于 « 的能量分数和分数。 R 处模型光谱的波数小于 k 时的耗散； = 1,000（实线）且位于 R； = 30（虚线）。对于两个雷诺数，水平条标识能量谱和耗散谱之间的“最重叠的十年波数”。
\end{sourceblock}
\begin{sourceblock}
\includegraphics[page=272,pagebox=mediabox,trim={64.39bp 1410.82bp 64.39bp 109.44bp},clip,width=\textwidth,max height=0.72\textheight]{Pope.pdf}
\par\vspace{0.45\baselineskip}
\small 图 6.23。最大重叠波数十年贡献的能量和耗散的分数 fo 作为 R 的函数；为模型谱。
\end{sourceblock}
( 1 0 \textasciitilde{}an\textasciitilde{}d )“重叠分数” f o 可以通过以下方式识别和定义

\begin{sourceblock}
\includegraphics[page=272,pagebox=mediabox,trim={64.39bp 1183.85bp 64.39bp 760.90bp},clip,width=0.94\textwidth,max height=0.38\textheight]{Pope.pdf}
\end{sourceblock}
因此，波数十进制 ( 1 0 \textasciitilde{}contr) 对能量和耗散都有一个略小于 f o 的分数。对于 RA 30 和 1,000，f o 的值为 0.75 和 0.11。图 6.23 显示了 fo 与 RAf 或模型频谱的函数关系。显然，需要非常大的雷诺数才能有十年的波数，其中能量和耗散都可以忽略不计。

能量级联图中的一个重要原则是（在高雷诺数时）能量耗散速率为 ui/lo，其中 uo 和 lo 是含能涡流的特征速度尺度和长度尺度。取 uo = k 1/ 2 和 lo = L l l，这个原则是 \textasciitilde{}k 3/2
LL.现在，根据定义 L k 3/2 我们有
使得 k 3 / 2 / \textasciitilde{} 1is1 的缩放相当于 L l l / L 的恒定性。图 6.24 显示了长度比例 L l l / L 作为 Rn 的函数。显然，根据模型谱，在高雷诺数时，L l l/ L 渐近趋于值 0.43。然而，随着 RAd 的减小，该比率显着增加，例如，它超过了渐近值 50\%
最后，图 6.25 显示了不同湍流雷诺数之间的关系。根据 ReL、Rn 和 Eq 的定义： (6.64) 我们有

\begin{sourceblock}
\includegraphics[page=272,pagebox=mediabox,trim={64.39bp 119.39bp 64.39bp 1806.77bp},clip,width=0.94\textwidth,max height=0.38\textheight]{Pope.pdf}
\end{sourceblock}
\begin{sourceblock}
\includegraphics[page=273,pagebox=mediabox,trim={64.39bp 1474.82bp 64.39bp 109.44bp},clip,width=\textwidth,max height=0.72\textheight]{Pope.pdf}
\par\vspace{0.45\baselineskip}
\small 图 6.24。纵向积分长度尺度与 L = 的比率作为模型光谱雷诺数的函数。
\end{sourceblock}
\begin{sourceblock}
\includegraphics[page=273,pagebox=mediabox,trim={64.39bp 998.87bp 64.39bp 580.93bp},clip,width=\textwidth,max height=0.72\textheight]{Pope.pdf}
\par\vspace{0.45\baselineskip}
\small 图 6.25。湍流雷诺数Re，（实线）和Re； （虚线）作为 R 的函数；为模型谱。
\end{sourceblock}
然而，基于 u' 和 L l l，我们有

\begin{sourceblock}
\includegraphics[page=273,pagebox=mediabox,trim={64.39bp 731.91bp 64.39bp 1156.85bp},clip,width=0.94\textwidth,max height=0.38\textheight]{Pope.pdf}
\end{sourceblock}
在湍流中，流动雷诺数 Re UL/V 通常比 ReT 大一个数量级（例如，u'/U 0.2、L l l / L 0.5），导致
DM。粗略估计Ri
对于高雷诺数和惯性波数

子范围，使用柯尔莫哥洛夫谱（方程（6.239））来估计
波数运动产生的能量比例更大

\begin{sourceblock}
\includegraphics[page=273,pagebox=mediabox,trim={64.39bp 202.97bp 64.39bp 1617.80bp},clip,width=0.94\textwidth,max height=0.38\textheight]{Pope.pdf}
\end{sourceblock}
这个估计值与表 6.2 中给出的值相比如何？

\subsection*{6.5.8 剪应力谱}\phantomsection\addcontentsline{toc}{subsection}{6.5.8 剪应力谱}
到目前为止，我们仅检查了各向同性湍流的速度谱，或局部各向同性湍流的速度谱的各向同性部分。在这些情况下，剪切应力谱，例如 < D 1 2 ( \textasciitilde{} ) 和 E l 2 ( \textasciitilde{} la)r,e 为零。在简单剪切流中，S i7(U1)/i7x2 > 0 是唯一显着的平均速度梯度，平均剪切速率 S 导致湍流具有各向异性。这种各向同性在雷诺应力中很明显（例如，( u l u 2 ) / k -0.3），并且考虑到以下关系

\begin{sourceblock}
\includegraphics[page=274,pagebox=mediabox,trim={64.39bp 1398.82bp 64.39bp 503.94bp},clip,width=0.94\textwidth,max height=0.38\textheight]{Pope.pdf}
\end{sourceblock}
因此，光谱至少在部分波数范围内也必须是各向异性的。鉴于剪切应力在动量传递和湍流能量产生中发挥的重要作用，确定不同运动尺度对 (ulu2) 的贡献非常重要。出现的简单、一致的图景是（不可避免地）对 (ulu2) 的主要贡献来自含能区内的波数，并且在较高波数 E 1 2 (\textasciitilde{}d1e)c 上比 El 更快（与局域各向同性一致）。

1.f是波数运动的特征时标则

Z(K) K，平均剪切 S 的影响由无量纲参数 ST(KI)t 来表征。可以合理地假设，如果 SZ(K) 较小，则也如此。 是平均剪切产生的各向异性水平。

在耗散区内，适当的时间尺度是 z。因此，正如 Corrsin (1958) 首先提出的，最小尺度各向同性的标准是

\begin{sourceblock}
\includegraphics[page=274,pagebox=mediabox,trim={64.39bp 711.91bp 64.39bp 1232.84bp},clip,width=0.94\textwidth,max height=0.38\textheight]{Pope.pdf}
\end{sourceblock}
参数 ST 随着 R' 的变化而变化（参见练习 6.35），因此方程： (6.270) 相当于高雷诺数要求。

在惯性子区内，适当的时间尺度是由 和 形成的，即 ( K \textasciitilde{} E ) - ' /H\textasciitilde{}e.n ce 波数各向同性的准则

\begin{sourceblock}
\includegraphics[page=274,pagebox=mediabox,trim={64.39bp 420.95bp 64.39bp 1521.81bp},clip,width=0.94\textwidth,max height=0.38\textheight]{Pope.pdf}
\end{sourceblock}
长度尺度 Ls 定义为

\begin{sourceblock}
\includegraphics[page=274,pagebox=mediabox,trim={64.39bp 291.96bp 64.39bp 1655.79bp},clip,width=0.94\textwidth,max height=0.38\textheight]{Pope.pdf}
\end{sourceblock}
该标准可以重新表示为

\begin{sourceblock}
\includegraphics[page=274,pagebox=mediabox,trim={64.39bp 165.98bp 64.39bp 1781.78bp},clip,width=0.94\textwidth,max height=0.38\textheight]{Pope.pdf}
\end{sourceblock}
在雷诺数非常高时，存在一个波数范围

\begin{sourceblock}
\includegraphics[page=275,pagebox=mediabox,trim={64.39bp 1714.78bp 64.39bp 225.97bp},clip,width=0.94\textwidth,max height=0.38\textheight]{Pope.pdf}
\end{sourceblock}
在惯性子区内。在此范围内的小水平各向异性可以假设为背景各向同性状态的小扰动（由 S 引起），其特征为 从该假设得出，剪切应力谱 E12(rc1i)s 由 K 和 S 确定，此外，作为小扰动，它随 S 线性变化。尺寸分析得出

\begin{sourceblock}
\includegraphics[page=275,pagebox=mediabox,trim={64.39bp 1334.83bp 64.39bp 600.92bp},clip,width=0.94\textwidth,max height=0.38\textheight]{Pope.pdf}
\end{sourceblock}
其中 E12i 是无量纲函数，速度尺度 us 是

\begin{sourceblock}
\includegraphics[page=275,pagebox=mediabox,trim={64.39bp 1161.85bp 64.39bp 779.90bp},clip,width=0.94\textwidth,max height=0.38\textheight]{Pope.pdf}
\end{sourceblock}
（练习 6.35 表明我们通常是 E12 与 S 的线性关系然后决定 E12，产生

\begin{sourceblock}
\includegraphics[page=275,pagebox=mediabox,trim={64.39bp 925.88bp 64.39bp 972.88bp},clip,width=0.94\textwidth,max height=0.38\textheight]{Pope.pdf}
\end{sourceblock}
\begin{sourceblock}
\includegraphics[page=275,pagebox=mediabox,trim={64.39bp 802.90bp 64.39bp 1143.86bp},clip,width=0.94\textwidth,max height=0.38\textheight]{Pope.pdf}
\end{sourceblock}
其中 C12i 是常数。这个结果归功于 Lumley (1967a)。

图 6.26 显示了在四个不同位置测量的 E 1 2 （\textasciitilde{}(1sc)a，由 Ls 和 us 主导）以及湍流边界层中的雷诺数。显然，对于 LS，数据与方程（1）合理一致。 (6.277) 与
当然，重要的是 E 1 2 (\textasciitilde{}Id,e)c 比
与此相比，各向异性随着
这可以直接从光谱相干性 H 1 2 (\textasciitilde{} ,w) 中看出，h ich 是傅里叶模式的相关系数。图 6.27 显示
湍流边界层中 H12(rcli) 的 1x1-u2 测量。在此基础上以及其他数据 Saddoughi 和 Veeravalli (1994) 提出了标准

\begin{sourceblock}
\includegraphics[page=275,pagebox=mediabox,trim={64.39bp 213.97bp 64.39bp 1732.78bp},clip,width=0.94\textwidth,max height=0.38\textheight]{Pope.pdf}
\end{sourceblock}
为光谱的局部各向同性区域。这与

\begin{sourceblock}
\includegraphics[page=276,pagebox=mediabox,trim={64.39bp 1214.85bp 64.39bp 109.44bp},clip,width=\textwidth,max height=0.72\textheight]{Pope.pdf}
\par\vspace{0.45\baselineskip}
\small 图 6.26。由 us 和 Lg 缩放的剪切应力谱：线，方程。 (6.277) 与 Cy, = 0.15;符号，Saddoughi 和 Veeravalli (1994) 的实验数据，来自 R + 500 至 1,450 的湍流边界层。
\end{sourceblock}
\begin{sourceblock}
\includegraphics[page=276,pagebox=mediabox,trim={64.39bp 609.92bp 64.39bp 871.89bp},clip,width=\textwidth,max height=0.72\textheight]{Pope.pdf}
\par\vspace{0.45\baselineskip}
\small 图 6.27。在 R 处的湍流边界层中测量的光谱相干性； = 1,400（Saddoughi 和 Veeravalli 1994）。
\end{sourceblock}
假设）Saddoughi 和 Veeravalli (19 94) 的数据表明 (2n/!e,,)Ls 6。

对于具有 S d(U1)/i?x2P, /E 11a, nd 的简单湍流剪切流
a -(ulu2)/k 0.3，得到如下结果：

\begin{sourceblock}
\includegraphics[page=276,pagebox=mediabox,trim={64.39bp 119.39bp 64.39bp 1808.77bp},clip,width=0.94\textwidth,max height=0.38\textheight]{Pope.pdf}
\end{sourceblock}
\section*{6.6 能量级联的谱观点}\phantomsection\addcontentsline{toc}{section}{6.6 能量级联的谱观点}
\begin{sourceblock}
\includegraphics[page=277,pagebox=mediabox,trim={64.39bp 1671.79bp 64.39bp 151.98bp},clip,width=0.94\textwidth,max height=0.38\textheight]{Pope.pdf}
\end{sourceblock}
\begin{sourceblock}
\includegraphics[page=277,pagebox=mediabox,trim={64.39bp 1566.80bp 64.39bp 342.96bp},clip,width=0.94\textwidth,max height=0.38\textheight]{Pope.pdf}
\end{sourceblock}
\begin{sourceblock}
\includegraphics[page=277,pagebox=mediabox,trim={64.39bp 1423.82bp 64.39bp 465.94bp},clip,width=0.94\textwidth,max height=0.38\textheight]{Pope.pdf}
\end{sourceblock}
6.6 能量级联的谱视图 在第 6.2-6.5 节中，我们介绍了用于量化不同尺度上的湍流运动的几种统计数据，并且我们通过实验数据、柯尔莫哥洛夫假设和简单的模型谱检查了这些统计数据。我们现在能够提供比第 6.1 节中给出的更全面的能量级联说明。因此，本节旨在总结和巩固之前的发展。

含能运动我们再次考虑极高雷诺数流，因此含能运动和耗散运动尺度之间存在明显的分离
长度尺度 l 的运动，与积分长度尺度 L l l (6Lll l lEI,s ay) 相当，其特征速度为 k112.S 量级，因为它们
l6 \textasciitilde{} 1 1 尺寸与流动尺寸 L 相当，这些大规模运动会受到流动几何形状的强烈影响。此外，它们的时间尺度 Lll/k112 与平均流量时间尺度相比较大（参见第 13I 页的表 5.2），因此它们受到流量历史记录的显着影响。换句话说，与普遍平衡范围相反，含能运动不具有统计平衡带来的普遍形式。

所有的各向异性都局限于含能量的运动，因此所有的湍流的产生也是如此。另一方面，粘性耗散可以忽略不计。相反，在级联的初始步骤中，能量被无粘过程移除，并以 TEl 速率转移到较小的尺度 (l < lEIa)，其尺度为 k 3 / 2 / \textasciitilde{} 1T1h。 i 的传递过程取决于非万能含能运动，因此
能谱平衡

对于均匀湍流（施加平均速度梯度），该图通过能谱函数 E(K,t ) 的平衡方程进行量化。这个方程（在 Hinze (1975) 和 Monin 和 Ya,glom (1975) 中详细推导）可以写成

\begin{sourceblock}
\includegraphics[page=278,pagebox=mediabox,trim={64.39bp 1513.81bp 64.39bp 393.95bp},clip,width=0.94\textwidth,max height=0.38\textheight]{Pope.pdf}
\end{sourceblock}
右侧的三个术语代表产生、光谱传输和耗散。

生产谱 P 由平均速度梯度 d(Ui)/dxj 和谱张量的各向异性部分的乘积给出。波数范围对生产的贡献表示为

\begin{sourceblock}
\includegraphics[page=278,pagebox=mediabox,trim={64.39bp 1150.85bp 64.39bp 708.91bp},clip,width=0.94\textwidth,max height=0.38\textheight]{Pope.pdf}
\end{sourceblock}
并且，就所有各向异性都包含在能量包含范围内而言，我们因此有

\begin{sourceblock}
\includegraphics[page=278,pagebox=mediabox,trim={64.39bp 980.88bp 64.39bp 966.88bp},clip,width=0.94\textwidth,max height=0.38\textheight]{Pope.pdf}
\end{sourceblock}
\begin{sourceblock}
\includegraphics[page=278,pagebox=mediabox,trim={64.39bp 900.89bp 64.39bp 1045.87bp},clip,width=0.94\textwidth,max height=0.38\textheight]{Pope.pdf}
\end{sourceblock}
在等式右侧的第二项中。 (6.284), I,(Kis) 光谱能量转移率：它是能量从波数较低的模式转移到波数高于 K 的模式的净速率。这与 l ( l ) 能量从大于 l 的涡流到小于 l 的涡流的转移率简单相关

\begin{sourceblock}
\includegraphics[page=278,pagebox=mediabox,trim={64.39bp 592.92bp 64.39bp 1354.83bp},clip,width=0.94\textwidth,max height=0.38\textheight]{Pope.pdf}
\end{sourceblock}
由此导致的波数范围内的能量增益率

( K , , K \textasciitilde{} ) 谱传递为

\begin{sourceblock}
\includegraphics[page=278,pagebox=mediabox,trim={64.39bp 400.95bp 64.39bp 1521.81bp},clip,width=0.94\textwidth,max height=0.38\textheight]{Pope.pdf}
\end{sourceblock}
由于 ?;I 在零和无限波数处消失，因此该传递项对湍动能 k 的平衡没有贡献。

的精确表达式可以从纳维-斯托克斯方程获得（参见，例如，Hinze (19 75)）。有两个贡献
一个由波数模式三波相互作用产生的结果，类似于方程。 (6.162)；的

\begin{sourceblock}
\includegraphics[page=279,pagebox=mediabox,trim={64.39bp 1084.86bp 64.39bp 109.44bp},clip,width=\textwidth,max height=0.72\textheight]{Pope.pdf}
\par\vspace{0.45\baselineskip}
\small 图 6.28。对于非常高雷诺数的均匀湍流，绘制了 (a) 能量和耗散谱、(b) 对 E(x, t) 平衡方程的贡献（方程 (6.284)）和 (c) 谱能量转移率的草图。
\end{sourceblock}
平均速度梯度对频谱的影响。方程中的最后一项。 (6.284) 是耗散谱 D(K,t ) t)。

图 6.28 是 E ( K , \textasciitilde{} )I 平衡方程中出现的量的草图。在含能区内，除耗散外，所有项都很重要。具有近似值 k(O,KEI) k、0 和
P，当它在含能区 (O,rcEI) 上积分时，P(o,KEI) 方程。 (6.284) 产量

\begin{sourceblock}
\includegraphics[page=279,pagebox=mediabox,trim={64.39bp 552.93bp 64.39bp 1392.82bp},clip,width=0.94\textwidth,max height=0.38\textheight]{Pope.pdf}
\end{sourceblock}
GI 其中 T\&cEI)I. 在惯性子区内，谱传递是唯一重要的过程，因此（当它从方程（6.284）积分时
KEI \textasciitilde{}\textasciitilde{} 1 ) 产量

\begin{sourceblock}
\includegraphics[page=279,pagebox=mediabox,trim={64.39bp 319.96bp 64.39bp 1626.79bp},clip,width=0.94\textwidth,max height=0.38\textheight]{Pope.pdf}
\end{sourceblock}
其中 TD1 7 \& c D I )W。在耗散区内，光谱传输平衡耗散，以便（当它从无穷大积分时）

KDI 方程(6.284) 产量

\begin{sourceblock}
\includegraphics[page=279,pagebox=mediabox,trim={64.39bp 119.39bp 64.39bp 1850.77bp},clip,width=0.94\textwidth,max height=0.38\textheight]{Pope.pdf}
\end{sourceblock}
将它们相加后，最后三个方程给出（未经近似）湍动能方程 dk/dt p
上述方程再次强调了能量级联的基本特征。从含能区 TEId 的能量转移速率以非通用方式取决于多个因素，包括平均速度梯度和光谱含能区的细节。然而，此传输速率随后建立了具有 TE;I 的通用特征的惯性子区，最后光谱的高波数部分以与以下速率相同的速率耗散能量：

埃尔。它收到它。因此 7 和 都由 和 确定，并且等于， 通常，当考虑“耗散”时，例如，将惯性距离谱表征为 E(K) CE 时，它在概念上是优越的
= 2/3 K -5/3 考虑 TEIi 代替 E。

级联时间尺度 An。有效性存疑的类比是，惯性子区内的能量流动就像不可压缩流体通过变面积管道的流动。恒定流量为 TEI（单位为每时间能量），而级联容量（类似于管道面积）为 E(K)（单位为每波数能量）。因此能量通过级联的速度（以每次波数为单位）为

\begin{sourceblock}
\includegraphics[page=280,pagebox=mediabox,trim={64.39bp 827.90bp 64.39bp 1104.86bp},clip,width=0.94\textwidth,max height=0.38\textheight]{Pope.pdf}
\end{sourceblock}
后一个表达式是从 Kolmogorov 谱和替换 TEI E 中获得的。请注意，该速度随着增加而迅速增加
从方程 d \textasciitilde{} / d = t k 的解可以得出，根据这个类比，能量从波数流出所需的时间 t(
到更高的波数是

\begin{sourceblock}
\includegraphics[page=280,pagebox=mediabox,trim={64.39bp 393.95bp 64.39bp 1470.81bp},clip,width=0.94\textwidth,max height=0.38\textheight]{Pope.pdf}
\end{sourceblock}
与关系 27r/eE1,eEI 与 L I \textasciitilde{} / L
KEI == L L \textasciitilde{}\textasciitilde{} , 0.4，该公式得出

\begin{sourceblock}
\includegraphics[page=280,pagebox=mediabox,trim={64.39bp 212.97bp 64.39bp 1726.78bp},clip,width=0.94\textwidth,max height=0.38\textheight]{Pope.pdf}
\end{sourceblock}
给出能量进入惯性系后的寿命估计

光谱能量传递模型 在普适平衡范围内（k > 光谱能量方程（方程（6.284））中的平衡点位于能量传递和耗散之间，参见图 6.28（b）。因此（在任何时刻 ￠）方程（6.284）简化为

\begin{sourceblock}
\includegraphics[page=281,pagebox=mediabox,trim={64.39bp 1566.80bp 64.39bp 368.95bp},clip,width=0.94\textwidth,max height=0.38\textheight]{Pope.pdf}
\end{sourceblock}
在 1940 年至 1970 年间，提出了许多光谱能量转移率 J 的模型，这些模型允许从方程 1 中推导出光谱 E(x) 的形式。 （6.296）。 Panchev (1971) 回顾了奥布霍夫 (Obukhov) (1941)、海森堡 (1948) 和许多其他人的建议。适合级联的物理原理，大多数模型都是非局部的，因为
假设取决于 E(x’)，对于 x’ \# x。然而，为了说明该方法，我们考虑 Pao (1965) 提出的简单局部模型。类似于等式。 (6.293)，能量传输速度 x(x) 定义为

\begin{sourceblock}
\includegraphics[page=281,pagebox=mediabox,trim={64.39bp 1116.86bp 64.39bp 830.89bp},clip,width=0.94\textwidth,max height=0.38\textheight]{Pope.pdf}
\end{sourceblock}
Pao 模型中的单一（尽管很强）假设是 k 仅取决于 ¢ 和 x。尺寸分析然后确定

\begin{sourceblock}
\includegraphics[page=281,pagebox=mediabox,trim={64.39bp 932.88bp 64.39bp 1014.87bp},clip,width=0.94\textwidth,max height=0.38\textheight]{Pope.pdf}
\end{sourceblock}
其中 « 是常数。用这个表达式 7,, 等式： (6.296) 可以积分（见练习 6.36）来产生 Pao 谱

\begin{sourceblock}
\includegraphics[page=281,pagebox=mediabox,trim={64.39bp 740.91bp 64.39bp 1196.85bp},clip,width=0.94\textwidth,max height=0.38\textheight]{Pope.pdf}
\end{sourceblock}
参见等式。 （6.254）。与图 6.15 中的实验数据进行比较。

练习 6.36 代入方程(6.298) 代入等式。 (6.296) 获得

\begin{sourceblock}
\includegraphics[page=281,pagebox=mediabox,trim={64.39bp 435.94bp 64.39bp 1470.81bp},clip,width=0.94\textwidth,max height=0.38\textheight]{Pope.pdf}
\end{sourceblock}
然后积分得到

\begin{sourceblock}
\includegraphics[page=281,pagebox=mediabox,trim={64.39bp 251.97bp 64.39bp 1632.79bp},clip,width=0.94\textwidth,max height=0.38\textheight]{Pope.pdf}
\end{sourceblock}
其中 f 是积分的（维）常数。论证说，对于
与柯尔莫哥洛夫谱（对于小xn）的一致性，是重新

\begin{sourceblock}
\includegraphics[page=282,pagebox=mediabox,trim={64.39bp 1709.78bp 64.39bp 158.98bp},clip,width=0.94\textwidth,max height=0.38\textheight]{Pope.pdf}
\end{sourceblock}
常数C.确认

因此这与柯尔莫哥洛夫相同

即，用 Ce 2/3 和 C，方程。 (6.301) 产生 Pao 谱，

\section*{6.7 局限、不足与改进}\phantomsection\addcontentsline{toc}{section}{6.7 局限、不足与改进}
考虑到不同尺度的湍流运动，能量级联、涡旋拉伸和柯尔莫哥洛夫假设的概念提供了一个宝贵的概念框架。然而，无论是概念上还是实践上，都存在一些缺陷。事实上，自 1960 年左右以来，主要的研究方向（理论、实验和计算）一直是检验这些缺点并尝试改进柯尔莫哥洛夫假设。虽然在这里对这些问题进行一些讨论是适当的，但应该意识到它们对湍流的研究和建模影响较小。这仅仅是因为小尺度（L < LEI）包含很少的能量（并且各向异性较小），因此对流动几乎没有直接影响。

\subsection*{6.7.1 雷诺数}\phantomsection\addcontentsline{toc}{subsection}{6.7.1 雷诺数}
柯尔莫哥洛夫假设的局限性在于它们仅适用于高雷诺数流，并且没有提供“足够高的雷诺数”的标准。许多实验室和实际流动具有相当高的雷诺数（例如，Re m 10,000、Ri m 150），但甚至发现耗散尺度上的运动也是各向异性的（例如，参见 George 和 Hussein (1991)）。

对惯性范围谱的仔细检查表明，柯尔莫哥洛夫
随着雷诺数的增加，光谱慢慢接近。通过在相当高雷诺数 (Ri 50-500) 下的网格湍流实验，Mydlarski 和 Warhaft (1998) 得出结论，惯性范围谱确实是幂律 E(K)，但指数 p 取决于 Ri（见图 6.29）。如图 6.29 中的曲线所示，很可能
Ri p 在 R 非常大时接近，但在 200（这是许多实验室流程的典型情况）时，p 约为 1.5。

假设能量级联肯定过于简单化

\begin{sourceblock}
\includegraphics[page=283,pagebox=mediabox,trim={64.39bp 1364.83bp 64.39bp 109.44bp},clip,width=\textwidth,max height=0.72\textheight]{Pope.pdf}
\par\vspace{0.45\baselineskip}
\small 图 6.29。谱幂律指数 p (E(k) \textasciitilde{} x7?) 作为网格湍流中雷诺数的函数：符号、Mydlarski 和 Warhaft (1998) 的实验数据；虚线，p = 3；实线，经验曲线 p = 3 \_ 4
\end{sourceblock}
尺寸稍小（例如 LC），并且这种能量传递仅取决于尺寸 C 的运动。光谱能量传递几乎不可能通过实验测量（但请参见 Kellogg 和 Corrsin (19 80)），而它可以从直接数值模拟（仅限于中或低雷诺数）中提取。 Domaradzki 和 Rogallo（1990）的 DNS 研究得出的结论是，能量向较小尺度和较大尺度转移，净转移向较小尺度。在波数空间中，能量传递是通过以下方式完成的
也就是说，具有波数三重相互作用 ua、ub 和 uC 的三种模式之间的相互作用，使得 ua ub uC 0（参见方程（6.162））。 DNS 结果表明，传输主要是局部的（例如，在 IuaI 1ubl 的模式 a 和 b 之间），但它受到与波数明显较小的第三模式的相互作用的影响（ e Iu "<<uaI()F。 Domaradzki (1992) 和 Zhou (1993) 进行了进一步的研究。）

\subsection*{6.7.2 高阶统计量}\phantomsection\addcontentsline{toc}{subsection}{6.7.2 高阶统计量}
本章到目前为止考虑的所有实验数据都属于二阶速度统计（速度的二次统计）。这些是最重要的量，因为它们决定动能和雷诺应力。

高阶统计量最简单的例子是归一化速度导数矩

图 6.30。 Van Atta 和 Antonia (1980) 编制的测量值（符号）

作为雷诺数函数的速度导数峰度。实线是
对于 n 3 和 n 4，它们是速度导数偏度 S 和峰度 K。（回想一下，对于高斯随机变量，S 为零，K 为 3。）根据柯尔莫哥洛夫假设，对于每个 n，M 是通用常数。然而，我们发现S和K并不是恒定的，而是随着雷诺数的增加而增加。例如，图 6.30 显示峰度测量值从低雷诺数网格湍流中的 K 4 增加到测量的最高雷诺数时的 K 40。与之前讨论的雷诺数效应相反，这里 K 似乎没有达到渐近线，但数据与无限增加一致，可能为 K R,3 /8。

速度导数矩 M（例如，偏度 S 和峰度 K ）属于耗散区。与惯性子区相关的最简单的高阶统计量是纵向速度结构函数，

\begin{sourceblock}
\includegraphics[page=284,pagebox=mediabox,trim={64.39bp 525.93bp 64.39bp 1419.82bp},clip,width=0.94\textwidth,max height=0.38\textheight]{Pope.pdf}
\end{sourceblock}
即，速度差的矩（在 x，t 处）定义为

\begin{sourceblock}
\includegraphics[page=284,pagebox=mediabox,trim={64.39bp 358.95bp 64.39bp 1587.80bp},clip,width=0.94\textwidth,max height=0.38\textheight]{Pope.pdf}
\end{sourceblock}
回想一下，第 6.2 节和第 6.3 节中考虑了二阶和三阶结构函数 D2(r) 和 D3(r)，其中它们用 DLL 表示，
根据柯尔莫哥洛夫的第二个假设，对于惯性范围分离

图 6.31。 Anselmet 等人编制的测量值（符号）。 （1984）的长期
惯性子区内的潮流速度结构函数指数，D,(r) \$"。

实线是 Kolmogorov (1941) 的预测，= 'n;虚线是
精确的相似性假设的预测，方程。 (6.323) 为 0.25。

\begin{sourceblock}
\includegraphics[page=285,pagebox=mediabox,trim={64.39bp 1142.86bp 64.39bp 744.91bp},clip,width=0.94\textwidth,max height=0.38\textheight]{Pope.pdf}
\end{sourceblock}
其中 C2、C3 是常数。测量证实了 n 2 的这一预测（使用 C2 2.0，参见第 194 页的图 6.5）；对于 n 3，柯尔莫哥洛夫定律（方程（6.88））产生 C3 高阶结构函数，Anselmet 等人测量了 n 高达 18 的情况。 （1984）。在惯性子区内，观察到对 r 的幂律依赖性，

\begin{sourceblock}
\includegraphics[page=285,pagebox=mediabox,trim={64.39bp 828.90bp 64.39bp 1119.86bp},clip,width=0.94\textwidth,max height=0.38\textheight]{Pope.pdf}
\end{sourceblock}
但测量的指数（如图 6.3 1 所示）与 Kolmogorov 不同
在预测方程中(6.306)，即 4 3。

检查构成这些高阶矩的 PDF 是很有启发性的。图 6.32 显示了在大气边界层中测量的 aul,/axl 的标准化 PDF（极高雷诺数流）。该 PDF 表示为 f z (z)，其中 Z 是标准化导数
观察分布的尾部（例如，超过四个标准导数）在该图上接近直线，对应于指数尾部：图 6.32 中所示的虚线是近似值

\begin{sourceblock}
\includegraphics[page=285,pagebox=mediabox,trim={64.39bp 180.98bp 64.39bp 1764.78bp},clip,width=0.94\textwidth,max height=0.38\textheight]{Pope.pdf}
\end{sourceblock}
\begin{sourceblock}
\includegraphics[page=285,pagebox=mediabox,trim={64.39bp 119.39bp 64.39bp 1848.77bp},clip,width=0.94\textwidth,max height=0.38\textheight]{Pope.pdf}
\end{sourceblock}
图 6.32。归一化速度导数 Z 的 PDF fi(z)

Van Atta 和 Chen (1970) 在大气中测量
边界层（高Re）。实线是高斯分布；虚线对应
到指数尾（方程（6.309）和（6.3 10））。

（请注意，负值的较慢衰减与观察到的负偏度 S 一致。）这种指数衰减当然比标准化高斯的衰减慢得多，如图 6.32 所示。

这些尾巴有什么意义呢？首先，它们对应于罕见事件：采用等式： (6.310) 作为大的近似值，可以得出以下结论
lz 121 超过 5 的概率小于 0.3\%。然而，这些低概率的尾部可以对更高的矩做出巨大贡献。表 6.3 显示了尾部对矩的贡献 (12 > 5)

f z (z) 由方程式给出。 （6.3 10）。例如，观察到，与高斯值 15 相比，对超偏度 M6 的贡献为 220。Belin 等人描述了 aul/axl 的 PDF 及其在雷诺数范围内的矩的一些实验室测量。 （1997）。

\subsection*{6.7.3 内部间歇性}\phantomsection\addcontentsline{toc}{subsection}{6.7.3 内部间歇性}
柯尔莫哥洛夫预测与高阶矩 M 和 D,(r) 的实验值之间的差异归因于内部间歇性，并且很大程度上在

\begin{center}\small 表 6.3。根据方程 1，Z 的 PDF 的指数尾部 (|Z| > 5) 对矩 M 的贡献 M©。 (6.310) 和 (6.311)\end{center}
\clearpage
\begin{sourceblock}
\includegraphics[page=287,pagebox=mediabox,trim={64.39bp 460.94bp 64.39bp 264.97bp},clip,width=\textwidth,max height=0.72\textheight]{Pope.pdf}
\end{sourceblock}
早在 1949 年（例如，Batchelor 和 Townsend (1949)），实验就重新
Zo 发现瞬时耗散间歇性达到很大
Zo/\& 值。观测到的峰值随着雷诺数的增加而增加：在实验室实验（中等RA）中，Meneveau和Sreenivasan（199 1）观察到峰值为ZO/\&15，而在大气表层（高RA），相应的观测值为50。Kolmogorov（1962）推测均方耗散涨落尺度为

\begin{sourceblock}
\includegraphics[page=287,pagebox=mediabox,trim={64.39bp 119.39bp 64.39bp 1848.77bp},clip,width=0.94\textwidth,max height=0.38\textheight]{Pope.pdf}
\end{sourceblock}
\begin{sourceblock}
\includegraphics[page=288,pagebox=mediabox,trim={64.39bp 1729.78bp 64.39bp 159.98bp},clip,width=0.94\textwidth,max height=0.38\textheight]{Pope.pdf}
\end{sourceblock}
其中 是间歇指数的正常数。实验证实了方程。 (6.31 7) 对于替代项 8，并确定 0.25 f 0.05 (Sreenivasan 和 Kailasnath 1993)。从定义上可以看出
（方程（6.314））这正是速度导数峰度 K。

( E , ) / E \textasciitilde{} 因此，采用并回顾关系 L/q R，方程。 （6.316）对于代理导致

\begin{sourceblock}
\includegraphics[page=288,pagebox=mediabox,trim={64.39bp 1403.82bp 64.39bp 533.93bp},clip,width=0.94\textwidth,max height=0.38\textheight]{Pope.pdf}
\end{sourceblock}
与实验数据一致（见图6.30）。

6.7.4 精炼的相似性假说考虑到速度增量 A,u（方程（6.305）），第一个（原始）柯尔莫哥洛夫假设指出 A,u（对于 r << L）的统计量是通用的，由平均耗散和 v 决定。

\subsection*{6.7.4 精炼的相似性假设}\phantomsection\addcontentsline{toc}{subsection}{6.7.4 精炼的相似性假设}
不受平均耗散影响，而是受局部值（在距离 r 上取平均值）影响，即，因此第一个拒绝的相似性假设是（对于 r << L）A,u 条件的统计量是通用的，由 和 v 确定。第二个精炼的相似性假设是，对于 q << r << L，这些条件统计量仅依赖于独立于 v。

将第二个改进的相似性假设应用于矩

\begin{sourceblock}
\includegraphics[page=288,pagebox=mediabox,trim={64.39bp 687.91bp 64.39bp 1202.85bp},clip,width=0.94\textwidth,max height=0.38\textheight]{Pope.pdf}
\end{sourceblock}
其中 C 是通用常数（参见方程（6.306）），并且是样本空间变量。结构函数 ( r ) 是无条件均值，然后获得为

\begin{sourceblock}
\includegraphics[page=288,pagebox=mediabox,trim={64.39bp 421.95bp 64.39bp 1471.81bp},clip,width=0.94\textwidth,max height=0.38\textheight]{Pope.pdf}
\end{sourceblock}
对于 n 3，因为等于原始值，并且经过改进的假设使
= (E,)E，同样的预测；其中，C3 就是柯尔莫哥洛夫定律（方程（6.88））。对于 n 6，并使用等式。 (6.31 7) 的预测是

\begin{sourceblock}
\includegraphics[page=288,pagebox=mediabox,trim={64.39bp 207.97bp 64.39bp 1737.78bp},clip,width=0.94\textwidth,max height=0.38\textheight]{Pope.pdf}
\end{sourceblock}
即 r ( \textasciitilde{} q (.6 .307)) 的幂律，指数为 2 1.75
对于 n 的其他值， ) 可以根据 E, 的 PDF 确定。 Obukhov (1962) 和 Kolmogorov (1962) 推测 是对数正态分布，即 in(\#\&, 具有高斯分布。根据对数正态性的这一假设以及 (方程 (6.317)) 的缩放，可以得出（参见练习 6.37）尺度矩为

\begin{sourceblock}
\includegraphics[page=289,pagebox=mediabox,trim={64.39bp 1553.80bp 64.39bp 387.95bp},clip,width=0.94\textwidth,max height=0.38\textheight]{Pope.pdf}
\end{sourceblock}
因此，结构函数（方程（6.320））预计将缩放为

\begin{sourceblock}
\includegraphics[page=289,pagebox=mediabox,trim={64.39bp 1393.82bp 64.39bp 497.94bp},clip,width=0.94\textwidth,max height=0.38\textheight]{Pope.pdf}
\end{sourceblock}
对于 n 不太大（例如 n <\_ 10），该预测与图 6.31 中所示的数据相当一致。对于较大的 n，差异归因于对数正态假设的缺陷，该假设受到 Mandelbrot (1974) 等人的严厉批评。

对于二阶结构函数 D2(r)，式： (6.323) 产量

\begin{sourceblock}
\includegraphics[page=289,pagebox=mediabox,trim={64.39bp 1093.86bp 64.39bp 854.89bp},clip,width=0.94\textwidth,max height=0.38\textheight]{Pope.pdf}
\end{sourceblock}
\begin{sourceblock}
\includegraphics[page=289,pagebox=mediabox,trim={64.39bp 922.88bp 64.39bp 959.88bp},clip,width=0.94\textwidth,max height=0.38\textheight]{Pope.pdf}
\end{sourceblock}
因此，预计对频谱的修改非常小。

对于速度导数矩，改进的假设产生

\begin{sourceblock}
\includegraphics[page=289,pagebox=mediabox,trim={64.39bp 720.91bp 64.39bp 1174.85bp},clip,width=0.94\textwidth,max height=0.38\textheight]{Pope.pdf}
\end{sourceblock}
其中 C, 是常数，因此

\begin{sourceblock}
\includegraphics[page=289,pagebox=mediabox,trim={64.39bp 501.94bp 64.39bp 1385.83bp},clip,width=0.94\textwidth,max height=0.38\textheight]{Pope.pdf}
\end{sourceblock}
使用对数正态假设来评估（方程（6.322））的矩，对于偏度和峰度，我们获得

\begin{sourceblock}
\includegraphics[page=289,pagebox=mediabox,trim={64.39bp 222.97bp 64.39bp 1671.79bp},clip,width=0.94\textwidth,max height=0.38\textheight]{Pope.pdf}
\end{sourceblock}
因此，无论 p 的值如何，预测都是

\begin{sourceblock}
\includegraphics[page=289,pagebox=mediabox,trim={64.39bp 119.39bp 64.39bp 1844.77bp},clip,width=0.94\textwidth,max height=0.38\textheight]{Pope.pdf}
\end{sourceblock}
\begin{sourceblock}
\includegraphics[page=290,pagebox=mediabox,trim={64.39bp 1367.83bp 64.39bp 109.44bp},clip,width=\textwidth,max height=0.72\textheight]{Pope.pdf}
\par\vspace{0.45\baselineskip}
\small 图 6.33。 Van Atta 和 Antonia (1980) 编制的速度导数偏度 S 和峰度 K 的测量结果。该线为-S \textasciitilde{} K*/8。
\end{sourceblock}
图 6.33 显示数据确实与此预测一致，尽管存在相当大的分散。

该领域的研究文献非常丰富。 Nelkin (1994) 和 Stolovitzky 等人提供了有用的评论。 （1995）。

正随机变量服从对数正态分布，如果

\begin{sourceblock}
\includegraphics[page=290,pagebox=mediabox,trim={64.39bp 886.89bp 64.39bp 1059.87bp},clip,width=0.94\textwidth,max height=0.38\textheight]{Pope.pdf}
\end{sourceblock}
是高斯分布，其中 是正常数。令其为高斯分布
(a) 证明 的矩是

\begin{sourceblock}
\includegraphics[page=290,pagebox=mediabox,trim={64.39bp 655.92bp 64.39bp 1287.84bp},clip,width=0.94\textwidth,max height=0.38\textheight]{Pope.pdf}
\end{sourceblock}
(b) 证明，如果 eRf 取为 ( E )，则

\begin{sourceblock}
\includegraphics[page=290,pagebox=mediabox,trim={64.39bp 519.93bp 64.39bp 1417.82bp},clip,width=0.94\textwidth,max height=0.38\textheight]{Pope.pdf}
\end{sourceblock}
(c) 如果 的均方取决于参数 r，根据

\begin{sourceblock}
\includegraphics[page=290,pagebox=mediabox,trim={64.39bp 359.95bp 64.39bp 1588.80bp},clip,width=0.94\textwidth,max height=0.38\textheight]{Pope.pdf}
\end{sourceblock}
其中 A、L 和 是正常数，表明

\begin{sourceblock}
\includegraphics[page=290,pagebox=mediabox,trim={64.39bp 228.97bp 64.39bp 1711.78bp},clip,width=0.94\textwidth,max height=0.38\textheight]{Pope.pdf}
\end{sourceblock}
\begin{sourceblock}
\includegraphics[page=290,pagebox=mediabox,trim={64.39bp 152.98bp 64.39bp 1786.77bp},clip,width=0.94\textwidth,max height=0.38\textheight]{Pope.pdf}
\end{sourceblock}
\subsection*{6.7.5 结束语}\phantomsection\addcontentsline{toc}{subsection}{6.7.5 结束语}
大量的研究工作集中在原始柯尔莫哥洛夫假设的缺点以及改进的相似性假设的细节上。然而，需要再次强调的是，在湍流中的平均速度场和雷诺应力的背景下，这些问题并不重要。湍流能量和各向异性主要包含在大尺度运动中。另一方面，内部间歇性涉及仅在小规模量的高阶统计中显现的罕见事件。在能量级联和湍流的背景下，比内部间歇性更重要的问题是决定能量从含能尺度转移的速率的问题。这是一个困难的问题，因为大尺度并不普遍。这个问题将在第 10 章中讨论。

\section*{7.1 槽道流}\phantomsection\addcontentsline{toc}{section}{7.1 槽道流}
\subsection*{7.1.1 流程描述}\phantomsection\addcontentsline{toc}{subsection}{7.1.1 流程描述}
与第 5 章中考虑的自由剪切流相反，大多数湍流（至少部分地）受到一个或多个固体表面的限制。示例包括内部流动，例如通过管道和管道的流动；外部流动，例如飞机和船体周围的流动；以及在大气边界层、河流等环境中的流动。

我们考虑其中三个最简单的流程（如图 7 4 所示）

\chapter{壁面流}
平板边界层。在每种流动中，平均速度矢量都（或几乎）平行于壁面，并且正如我们将看到的，每种情况下的近壁行为都非常相似。这些简单的流动具有实际重要性，并且在湍流研究的历史发展中发挥了突出的作用。

核心问题是平均速度分布的形式和摩擦定律，它们描述了流体在壁上施加的剪切应力。另外，第7.1.7节介绍了混合长度；雷诺应力的平衡方程在第 7.3.5 节中推导和检验；适当的正交分解（POD）在 7.4 节中描述。

7.1 槽道流量
7.1.1 的描述 如图 7.1 所示，我们考虑通过高度为 h 26 的矩形管道的流动。该管道很长 (L/6 >> 1)，并且具有较大的长宽比 (b/6 >> 1)。平均流量主要在轴向（x 方向，
= X I ），平均速度主要在横流方向（y 方向）变化。底壁和顶壁分别位于 y 0 和 y = 26 处，中平面为 y = 6。与 6 相比，通道在展向 (z = x3) 方向上的范围较大，因此（远离端壁）

\begin{sourceblock}
\includegraphics[page=293,pagebox=mediabox,trim={64.39bp 985.88bp 64.39bp 109.44bp},clip,width=\textwidth,max height=0.72\textheight]{Pope.pdf}
\par\vspace{0.45\baselineskip}
\small 图 7.1。 (a) 槽道流、(b) 管道流和 (c) 平板边界层的草图。
\end{sourceblock}
流量在统计上与 z 无关。中心线由 y 6, z 0 定义。

( u l，三个坐标方向上的速度为( U, V
U2, U?) 波动 (u,u, ( u l , u ? u3)+T 平均横流速度 ( W ) 为零。

靠近管道入口 (x 0) 存在流动发展区域。然而，我们将注意力集中在完全发达的区域（大x），其中速度统计数据不再随x变化。因此，完全发展的槽道流被认为是统计平稳和统计一维的，速度统计仅取决于 y。实验证实了自然期望，即流动关于中平面 y 6 统计对称：y 处的 (U, V, W ) 统计量与 26 y 处的 (U, -V, W ) 统计量相同（参见方程（4.31）-（4.34））。

用于表征流动的雷诺数为

\begin{sourceblock}
\includegraphics[page=293,pagebox=mediabox,trim={64.39bp 210.97bp 64.39bp 1733.78bp},clip,width=0.94\textwidth,max height=0.38\textheight]{Pope.pdf}
\end{sourceblock}
\begin{sourceblock}
\includegraphics[page=293,pagebox=mediabox,trim={64.39bp 119.39bp 64.39bp 1849.77bp},clip,width=0.94\textwidth,max height=0.38\textheight]{Pope.pdf}
\end{sourceblock}
(u),=a 其中 Uo -- 是中心线速度， 是体积速度

\begin{sourceblock}
\includegraphics[page=294,pagebox=mediabox,trim={64.39bp 1679.79bp 64.39bp 212.97bp},clip,width=0.94\textwidth,max height=0.38\textheight]{Pope.pdf}
\end{sourceblock}
Re < 1,350 时流动为层流，Re > 1,800 时流动为完全湍流，尽管在 Re 3,000 之前过渡效应很明显（参见 Patel 和 Head (1969)）。

\subsection*{7.1.2 平均力的平衡}\phantomsection\addcontentsline{toc}{subsection}{7.1.2 平均力的平衡}
平均连续性方程简化为

\begin{sourceblock}
\includegraphics[page=294,pagebox=mediabox,trim={64.39bp 1269.84bp 64.39bp 629.92bp},clip,width=0.94\textwidth,max height=0.38\textheight]{Pope.pdf}
\end{sourceblock}
因为 ( W ) 为零，并且 ( U ) 与 x 无关。对于边界条件 (V),=o，这表明 (V) 对于所有 y 都为零，因此顶壁 0 处的边界条件也得到满足。

横向平均动量方程简化为

\begin{sourceblock}
\includegraphics[page=294,pagebox=mediabox,trim={64.39bp 976.88bp 64.39bp 921.88bp},clip,width=0.94\textwidth,max height=0.38\textheight]{Pope.pdf}
\end{sourceblock}
边界条件为 0 时，积分为

\begin{sourceblock}
\includegraphics[page=294,pagebox=mediabox,trim={64.39bp 850.89bp 64.39bp 1089.86bp},clip,width=0.94\textwidth,max height=0.38\textheight]{Pope.pdf}
\end{sourceblock}
其中 pw = (p(x,0 ,O)) 是 bott0.m 壁上的平均压力。从该方程中得出的一个重要推论是，平均轴向压力梯度在整个流动中是均匀的：

\begin{sourceblock}
\includegraphics[page=294,pagebox=mediabox,trim={64.39bp 618.92bp 64.39bp 1286.84bp},clip,width=0.94\textwidth,max height=0.38\textheight]{Pope.pdf}
\end{sourceblock}
轴向平均动量方程，

\begin{sourceblock}
\includegraphics[page=294,pagebox=mediabox,trim={64.39bp 456.94bp 64.39bp 1443.82bp},clip,width=0.94\textwidth,max height=0.38\textheight]{Pope.pdf}
\end{sourceblock}
可以重写

\begin{sourceblock}
\includegraphics[page=294,pagebox=mediabox,trim={64.39bp 311.96bp 64.39bp 1584.80bp},clip,width=0.94\textwidth,max height=0.38\textheight]{Pope.pdf}
\end{sourceblock}
其中总剪应力 z(y) 为

\begin{sourceblock}
\includegraphics[page=294,pagebox=mediabox,trim={64.39bp 161.98bp 64.39bp 1740.78bp},clip,width=0.94\textwidth,max height=0.38\textheight]{Pope.pdf}
\end{sourceblock}
（方程（7.9））相当于力的平衡：轴向法向应力梯度由横流剪切应力梯度平衡。

由于 是仅 y 的函数，并且 p 是仅 x 的函数，因此从等式可以明显看出： (7.9) d.r/dy 和 dpw/dx 都是常数。 7(y) 和 dp,/dx 的解可以根据壁剪应力明确地写出

\begin{sourceblock}
\includegraphics[page=295,pagebox=mediabox,trim={64.39bp 1548.80bp 64.39bp 397.95bp},clip,width=0.94\textwidth,max height=0.38\textheight]{Pope.pdf}
\end{sourceblock}
因为 7(y) 关于中平面反对称，所以 \textasciitilde{} ( 6 i) 为零；顶壁处的应力为 \textasciitilde{} ( 2 6 ) 因此，解为

\begin{sourceblock}
\includegraphics[page=295,pagebox=mediabox,trim={64.39bp 1320.83bp 64.39bp 555.93bp},clip,width=0.94\textwidth,max height=0.38\textheight]{Pope.pdf}
\end{sourceblock}
\begin{sourceblock}
\includegraphics[page=295,pagebox=mediabox,trim={64.39bp 1165.85bp 64.39bp 703.91bp},clip,width=0.94\textwidth,max height=0.38\textheight]{Pope.pdf}
\end{sourceblock}
通过参考速度归一化的壁剪切应力称为表面摩擦系数。在 Uo 的基础上，我们定义

\begin{sourceblock}
\includegraphics[page=295,pagebox=mediabox,trim={64.39bp 986.88bp 64.39bp 941.88bp},clip,width=0.94\textwidth,max height=0.38\textheight]{Pope.pdf}
\end{sourceblock}
\begin{sourceblock}
\includegraphics[page=295,pagebox=mediabox,trim={64.39bp 902.89bp 64.39bp 1026.87bp},clip,width=0.94\textwidth,max height=0.38\textheight]{Pope.pdf}
\end{sourceblock}
总结一下：流动是由通道入口和出口之间的压降驱动的。在完全发育的区域中，存在恒定（负）平均压力梯度 d(p)/i?x dp,/dx，它由剪应力梯度 d.r/dy 6 平衡。对于给定的压力梯度 dp,/dx 和通道半宽 6，线性剪应力分布由方程 1 给出。 （7.12）和（7.13）独立于流体属性（例如， p 和 v ），并且独立于流体运动状态（即层流或湍流）。请注意，如果流量由 p、v、6 和 dp、/dx 定义，则 Uo 和 U 是先验未知的。或者，在实验中可以施加 U，然后压力梯度是未知的。在这两种情况下，表面摩擦系数都是未知的。当然，正如以下练习所示，层流的所有这些量都很容易确定。

对于层流，从方程。 (7.10) 和 (7.13) 表明均值
速度剖面为

\begin{sourceblock}
\includegraphics[page=295,pagebox=mediabox,trim={64.39bp 119.39bp 64.39bp 1812.77bp},clip,width=0.94\textwidth,max height=0.38\textheight]{Pope.pdf}
\end{sourceblock}
\begin{sourceblock}
\includegraphics[page=296,pagebox=mediabox,trim={64.39bp 1382.83bp 64.39bp 109.44bp},clip,width=\textwidth,max height=0.72\textheight]{Pope.pdf}
\par\vspace{0.45\baselineskip}
\small 图 7.2。来自 Kim 等人的 DNS 的完全发展的湍流槽道流中的平均速度剖面。 (1987)：虚线，Re = 5,600；实线，Re = 13,750。
\end{sourceblock}
从而得到如下结果：

\begin{sourceblock}
\includegraphics[page=296,pagebox=mediabox,trim={64.39bp 1146.86bp 64.39bp 776.90bp},clip,width=0.94\textwidth,max height=0.38\textheight]{Pope.pdf}
\end{sourceblock}
\begin{sourceblock}
\includegraphics[page=296,pagebox=mediabox,trim={64.39bp 1050.87bp 64.39bp 887.89bp},clip,width=0.94\textwidth,max height=0.38\textheight]{Pope.pdf}
\end{sourceblock}
\begin{sourceblock}
\includegraphics[page=296,pagebox=mediabox,trim={64.39bp 831.89bp 64.39bp 967.88bp},clip,width=0.94\textwidth,max height=0.38\textheight]{Pope.pdf}
\end{sourceblock}
的定义为

摩擦速度

\begin{sourceblock}
\includegraphics[page=296,pagebox=mediabox,trim={64.39bp 709.91bp 64.39bp 1237.84bp},clip,width=0.94\textwidth,max height=0.38\textheight]{Pope.pdf}
\end{sourceblock}
表明，一般而言，

\begin{sourceblock}
\includegraphics[page=296,pagebox=mediabox,trim={64.39bp 586.93bp 64.39bp 1354.83bp},clip,width=0.94\textwidth,max height=0.38\textheight]{Pope.pdf}
\end{sourceblock}
以及拉米哈尔流

\begin{sourceblock}
\includegraphics[page=296,pagebox=mediabox,trim={64.39bp 414.95bp 64.39bp 1483.81bp},clip,width=0.94\textwidth,max height=0.38\textheight]{Pope.pdf}
\end{sourceblock}
计算层流上限的 u,/Uo，即 Re 1,350。

\subsection*{7.1.3 近壁剪应力}\phantomsection\addcontentsline{toc}{subsection}{7.1.3 近壁剪应力}
图 7.2 显示了 Kim (1987) 获得的平均速度剖面
等人。

\begin{sourceblock}
\includegraphics[page=297,pagebox=mediabox,trim={64.39bp 1428.82bp 64.39bp 109.44bp},clip,width=\textwidth,max height=0.72\textheight]{Pope.pdf}
\par\vspace{0.45\baselineskip}
\small 图 7.3。湍流槽道流中的粘性剪切应力和雷诺剪切应力分布：Kim 等人的 DNS 数据。 (1987)：虚线，Re = 5,600；实线，Re = 13,750。
\end{sourceblock}
为 Re 5,600 和 Re 13,750。本节和下一小节的目的是解释和量化这些概况。

总剪切应力 7(y)（方程（7.10））是粘性应力之和
和雷诺应力 -p(uv)。在壁处，边界 pv d(U)/dv 条件 U(x, t) = 0 指示所有雷诺应力为零。因此，壁剪切应力完全是由于粘性的贡献，

\begin{sourceblock}
\includegraphics[page=297,pagebox=mediabox,trim={64.39bp 846.89bp 64.39bp 1044.87bp},clip,width=0.94\textwidth,max height=0.38\textheight]{Pope.pdf}
\end{sourceblock}
粘性和雷诺剪切应力的分布如图 7.3 所示。

重要的观察结果是，壁上的粘性应力占主导地位，这与自由剪切流中的情况形成鲜明对比。在那里，雷诺数很高，与雷诺应力相比，粘性应力到处都小得可以忽略不计。此外，在壁附近，由于粘度是一个有影响的参数，因此与自由剪切流相比，速度分布再次取决于雷诺数（如图 7.2 所示）。

显然，靠近壁面时，粘度和壁面剪应力是重要的参数。根据这些量（和 p），我们定义粘性尺度，即近壁区域中适当的速度尺度和长度尺度。这些是摩擦速度

\begin{sourceblock}
\includegraphics[page=297,pagebox=mediabox,trim={64.39bp 204.97bp 64.39bp 1717.78bp},clip,width=0.94\textwidth,max height=0.38\textheight]{Pope.pdf}
\end{sourceblock}
Moser、Kim 和 Mansour (1999) 简要介绍了较高雷诺数数据。一个

\begin{sourceblock}
\includegraphics[page=298,pagebox=mediabox,trim={64.39bp 1385.83bp 64.39bp 109.44bp},clip,width=\textwidth,max height=0.72\textheight]{Pope.pdf}
\par\vspace{0.45\baselineskip}
\small 图 7.4。粘性应力和雷诺应力对总应力的分数贡献曲线。 Kim 等人的 DNS 数据。 (1987)：虚线，Re = 5,600；实线，Re = 13,750。
\end{sourceblock}
和粘性长度尺度

\begin{sourceblock}
\includegraphics[page=298,pagebox=mediabox,trim={64.39bp 1113.86bp 64.39bp 820.90bp},clip,width=0.94\textwidth,max height=0.38\textheight]{Pope.pdf}
\end{sourceblock}
基于粘性尺度 u,6,/v 的雷诺数完全相同，而摩擦雷诺数定义为

\begin{sourceblock}
\includegraphics[page=298,pagebox=mediabox,trim={64.39bp 875.89bp 64.39bp 1026.87bp},clip,width=0.94\textwidth,max height=0.38\textheight]{Pope.pdf}
\end{sourceblock}
（在 Kim 等人 1987 年的 DNS 中，摩擦雷诺数在 Re 5,600 处为 Re, -- 180，在 Re 13,750 处为 Re, 395。）

以粘性长度或壁单位测量的距壁的距离表示为

\begin{sourceblock}
\includegraphics[page=298,pagebox=mediabox,trim={64.39bp 619.92bp 64.39bp 1326.83bp},clip,width=0.94\textwidth,max height=0.38\textheight]{Pope.pdf}
\end{sourceblock}
请注意，y+ 类似于局部雷诺数，因此可以预期其大小可以确定粘性过程和湍流过程的相对重要性。为了支持这一假设，图 7.4 显示了槽道流近壁区域中的粘性应力和雷诺应力对总应力的贡献分数。当针对 y f 绘制时，两个雷诺数的分布几乎崩溃。粘性贡献从壁 (y+ 0) 处的 100\% 下降到 y+ 12 处的 50\%，并且 y + = 50 处小于 10\%。

近壁流中的不同区域或层是根据 y+ 定义的。在粘性壁区域 y+ < 50，存在分子的直接影响
粘度的直接影响可以忽略不计。在粘性壁区域内，在粘性子层 y+ < 5 中，雷诺剪应力与粘性应力相比可以忽略不计。随着流动的雷诺数增加，粘性壁区域占据的通道部分减少，因为 6,/6 随着 Re--' 变化（方程（7.27））。

在完全发展的湍流通道上进行了实验

Re lo5 处的流量。流体是水 (v 1.14 x m2 s ) 并且
通道半高为6×2厘米。求出皮肤摩擦系数
为 Cf 4.4 x 确定：U、u、/U、Re、和 6、/6。什么是
粘性壁区域和粘性底层的厚度，
既是 6 的分数又是毫米？

\subsection*{7.1.4 平均速度剖面}\phantomsection\addcontentsline{toc}{subsection}{7.1.4 平均速度剖面}
完全展开的槽道流量完全由 p、v、6 和 dp、/dx 指定；或者，等效地，通过 p、v、6 和 u，，因为我们有
只有两个独立的无量纲群可以由 p、v、6、u 和 y 组成（例如 y/6 和 Re, = u,6/v），因此平均速度分布可以写成

\begin{sourceblock}
\includegraphics[page=299,pagebox=mediabox,trim={64.39bp 679.91bp 64.39bp 1264.84bp},clip,width=0.94\textwidth,max height=0.38\textheight]{Pope.pdf}
\end{sourceblock}
其中Fo是待确定的通用无量纲函数。

虽然这种确定平均速度分布的方法看起来很自然，但是最好以稍微不同的方式进行。我们考虑的是速度梯度 d(U)/dy，而不是 (U)，它是动态重要量。例如，粘性应力和湍流产生均由 d(U)/dy 决定。再次基于量纲，d(U)/dy 仅取决于两个无量纲参数，因此（无需任何假设）我们可以写出
其中 是通用无量纲函数。选择背后的想法
壁区域 (y+ 50)，而 6 是外层 (y+ 50) 中的适当比例。关系

\begin{sourceblock}
\includegraphics[page=300,pagebox=mediabox,trim={64.39bp 1680.79bp 64.39bp 267.97bp},clip,width=0.94\textwidth,max height=0.38\textheight]{Pope.pdf}
\end{sourceblock}
不可避免地表明，这两个参数包含与 y / 6 和 Re 相同的信息（方程（7.30））。

壁定律 Prandtl (1925) 假设，在高雷诺数时，靠近壁 (y/3 << 1) 存在一个内层，其中平均速度剖面由粘性尺度决定，与 6 和 Uo 无关。从数学上讲，这意味着方程中的函数 @(y/6,,,y / 6 ) (7.3 1) 仅当 y / 6 趋于零时才渐近趋于 y/6 的函数，因此等式： (7.31) 变为

\begin{sourceblock}
\includegraphics[page=300,pagebox=mediabox,trim={64.39bp 1000.87bp 64.39bp 745.91bp},clip,width=0.94\textwidth,max height=0.38\textheight]{Pope.pdf}
\end{sourceblock}
其中 y+ y/6 和 u (y ) 定义为

\begin{sourceblock}
\includegraphics[page=300,pagebox=mediabox,trim={64.39bp 849.89bp 64.39bp 1050.87bp},clip,width=0.94\textwidth,max height=0.38\textheight]{Pope.pdf}
\end{sourceblock}
等式。 (7.33) 也可以写成

\begin{sourceblock}
\includegraphics[page=300,pagebox=mediabox,trim={64.39bp 688.91bp 64.39bp 1206.85bp},clip,width=0.94\textwidth,max height=0.38\textheight]{Pope.pdf}
\end{sourceblock}
方程的积分。 (7.36) 是墙定律：

\begin{sourceblock}
\includegraphics[page=300,pagebox=mediabox,trim={64.39bp 573.93bp 64.39bp 1371.83bp},clip,width=0.94\textwidth,max height=0.38\textheight]{Pope.pdf}
\end{sourceblock}
\begin{sourceblock}
\includegraphics[page=300,pagebox=mediabox,trim={64.39bp 414.95bp 64.39bp 1432.82bp},clip,width=0.94\textwidth,max height=0.38\textheight]{Pope.pdf}
\end{sourceblock}
重要的一点不是等式。 (7.38)，但事实是（根据普朗特的假设）u+ 仅取决于 y+，因为 y / 6 << 1。

对于不太接近过渡的雷诺数，有大量的实验验证函数 f 是通用的，不仅对于槽道流，而且对于管流和b 边界层。如现在所示，函数 fw(y+)c 的形式可以针对小值和大值确定

\begin{sourceblock}
\includegraphics[page=301,pagebox=mediabox,trim={64.39bp 1387.83bp 64.39bp 109.44bp},clip,width=\textwidth,max height=0.72\textheight]{Pope.pdf}
\par\vspace{0.45\baselineskip}
\small 图 7.5。来自 Kim 等人的 DNS 数据的平均速度近壁剖面。 (1987)：虚线，Re = 5,600；实线，Re = 13,750；点划线，ut = y*。
\end{sourceblock}
粘性子层 无滑移条件 (U),=, 0 对应于 fw(0) 0，而壁处的粘性应力定律（方程 (7.24)）得出导数

\begin{sourceblock}
\includegraphics[page=301,pagebox=mediabox,trim={64.39bp 1047.87bp 64.39bp 897.89bp},clip,width=0.94\textwidth,max height=0.38\textheight]{Pope.pdf}
\end{sourceblock}
（这只是粘性尺度标准化的结果。）因此，对于小 y+，f w ( y f ) 的泰勒级数展开式为

\begin{sourceblock}
\includegraphics[page=301,pagebox=mediabox,trim={64.39bp 849.89bp 64.39bp 1092.86bp},clip,width=0.94\textwidth,max height=0.38\textheight]{Pope.pdf}
\end{sourceblock}
（事实上，仔细检查就会发现，在线性项之后，下一个非零项的阶数是 y，请参见练习 7.9。）

图 7.5 显示了通过直接数值模拟获得的近壁区域 u+ 的分布。与线性关系 uf = y+ 的偏离在粘性子层 (y+ 5) 中可以忽略不计，但对于 yf 12 来说非常显着（大于 25\%）。

对数定律 内层通常定义为 y / 6 < 0.1。在高雷诺数下，内层的外部部分对应于大的 y f，即 y+ 0.16 /6，！奥。 l \textasciitilde{} e >, > 1。正如已经讨论过的，对于较大的 y+，可以认为粘度几乎没有影响。因此，在等式中。 (7.33) 中，(Dl(y/6,) 对 v（通过 6,）的依赖性消失，因此采用由下式表示的常数值

\begin{sourceblock}
\includegraphics[page=301,pagebox=mediabox,trim={64.39bp 119.39bp 64.39bp 1762.78bp},clip,width=0.94\textwidth,max height=0.38\textheight]{Pope.pdf}
\end{sourceblock}
\begin{sourceblock}
\includegraphics[page=302,pagebox=mediabox,trim={64.39bp 1414.82bp 64.39bp 109.44bp},clip,width=\textwidth,max height=0.72\textheight]{Pope.pdf}
\par\vspace{0.45\baselineskip}
\small 图 7.6。平均速度的近壁剖面：实线，Kim 等人的 DNS 数据。 (1987): 回复 = 13, 750;点-虚线，u* = y*；虚线，对数定律，方程。 (7.43)-(7.44)。
\end{sourceblock}
因此，在该区域中，平均速度梯度为

\begin{sourceblock}
\includegraphics[page=302,pagebox=mediabox,trim={64.39bp 1146.86bp 64.39bp 753.90bp},clip,width=0.94\textwidth,max height=0.38\textheight]{Pope.pdf}
\end{sourceblock}
它集成到

\begin{sourceblock}
\includegraphics[page=302,pagebox=mediabox,trim={64.39bp 999.87bp 64.39bp 927.88bp},clip,width=0.94\textwidth,max height=0.38\textheight]{Pope.pdf}
\end{sourceblock}
其中 B 是常数。这是 von Karmin (1930) 提出的墙的对数定律，或者简称为对数定律，并且是 von \textasciitilde{} i r m h nconstant。

- - rc 在文献中，对数定律常数的值存在一些变化，但通常在 5\% 以内

\begin{sourceblock}
\includegraphics[page=302,pagebox=mediabox,trim={64.39bp 741.91bp 64.39bp 1205.85bp},clip,width=0.94\textwidth,max height=0.38\textheight]{Pope.pdf}
\end{sourceblock}
图 7.6 显示了对数定律与通道内部 DNS 数据（y/6 0.25）的比较。显然 y+ 30 具有极好的一致性。

对数定律在半对数图中更清楚地揭示。图 7.7 显示了雷诺数在 Reo 3,000 和 Reo 40,000 之间时测量的湍流槽道流的 u+(y +) 分布。可以看出，数据折叠成一条曲线，证实了壁面定律，并且对于 y+ 30，数据符合对数定律，除了通道中平面附近（每个雷诺数的最后几个数据点）。

5) 和对数律区域
粘性子层 (y+ (yf > 30) 之间的区域称为缓冲层。它是流动的粘性主导部分和湍流主导部分之间的过渡区域。用于描述近壁流的各种区域和层是

\begin{sourceblock}
\includegraphics[page=303,pagebox=mediabox,trim={64.39bp 1156.86bp 64.39bp 109.44bp},clip,width=\textwidth,max height=0.72\textheight]{Pope.pdf}
\par\vspace{0.45\baselineskip}
\small 图 7.7。 Wei 和 Willmarth (1989) 测量的完全发展的湍流槽道流中的平均速度剖面：©, Reo = 2,970； 0，雷奥 = 14,914； A，Reo = 22,776； V, 雷奥 = 39, 582;线，对数定律，方程。 (7.43)-(7.44)。
\end{sourceblock}
\begin{center}\small 表 7.1。墙区域和层及其定义属性\end{center}
\begin{sourceblock}
\includegraphics[page=303,pagebox=mediabox,trim={64.39bp 119.39bp 64.39bp 1050.87bp},clip,width=\textwidth,max height=0.72\textheight]{Pope.pdf}
\end{sourceblock}
\begin{sourceblock}
\includegraphics[page=304,pagebox=mediabox,trim={64.39bp 1281.84bp 64.39bp 109.44bp},clip,width=\textwidth,max height=0.72\textheight]{Pope.pdf}
\par\vspace{0.45\baselineskip}
\small 图 7.8。草图显示了以 yt = y/d 和 y/6 定义的各种壁区域和层，用于高雷诺数 (Re, = 104) 的湍流槽道流。
\end{sourceblock}
\begin{sourceblock}
\includegraphics[page=304,pagebox=mediabox,trim={64.39bp 991.87bp 64.39bp 853.89bp},clip,width=0.94\textwidth,max height=0.38\textheight]{Pope.pdf}
\end{sourceblock}
用 cDo 代替方程式中的 Q。 (7.31) 并对 y 和 6 进行积分，然后根据 'von K i r min (1930) 得出速度亏损定律：

\begin{sourceblock}
\includegraphics[page=304,pagebox=mediabox,trim={64.39bp 780.90bp 64.39bp 1119.86bp},clip,width=0.94\textwidth,max height=0.38\textheight]{Pope.pdf}
\end{sourceblock}
\begin{sourceblock}
\includegraphics[page=304,pagebox=mediabox,trim={64.39bp 618.92bp 64.39bp 1227.85bp},clip,width=0.94\textwidth,max height=0.38\textheight]{Pope.pdf}
\end{sourceblock}
根据定义，速度亏损是平均速度 (U) 与中心线值 Uo 之间的差值。速度亏损定律指出，由 u 标准化的速度亏损仅取决于 y/6。与墙定律函数 f,(y+) 不同，这里并不表明 FD(y/6) 是通用的：它在不同的流中是不同的。

当雷诺数足够高（大约 Re 20,000）时，内层 (y/6 0.1) 和外层 (\textasciitilde{} / 6 , 50) 之间存在重叠区域（见图 7.8）。在这个区域中，两个方程。 (7.33) 和 (7.45) 是有效的，产生（根据方程（7.31））

\begin{sourceblock}
\includegraphics[page=304,pagebox=mediabox,trim={64.39bp 119.39bp 64.39bp 1799.77bp},clip,width=0.94\textwidth,max height=0.38\textheight]{Pope.pdf}
\end{sourceblock}
\begin{sourceblock}
\includegraphics[page=305,pagebox=mediabox,trim={64.39bp 1387.83bp 64.39bp 109.44bp},clip,width=\textwidth,max height=0.72\textheight]{Pope.pdf}
\par\vspace{0.45\baselineskip}
\small 图 7.9。湍流槽道流中的平均速度亏损。实线，Kim 等人的 DNS。 (1987)，Re = 13,750；虚线，对数定律，方程。 \{7.43)-\{7.44)。
\end{sourceblock}
该方程只能在重叠区域中满足并且是

(DI (Do 常数，这导致

\begin{sourceblock}
\includegraphics[page=305,pagebox=mediabox,trim={64.39bp 1074.86bp 64.39bp 824.90bp},clip,width=0.94\textwidth,max height=0.38\textheight]{Pope.pdf}
\end{sourceblock}
密立根 (Millikan, 1938) 提出的这一论点提供了对数定律的另一种推导。它还建立了小速度亏损定律的形式

\begin{sourceblock}
\includegraphics[page=305,pagebox=mediabox,trim={64.39bp 833.89bp 64.39bp 1090.86bp},clip,width=0.94\textwidth,max height=0.38\textheight]{Pope.pdf}
\end{sourceblock}
其中 B1 是与流量相关的常数。 （重叠区域和导致对数定律的参数将在第 7.3.4 节中进一步考虑。）

图 7.9 显示了湍流槽道流 DNS 中的速度亏损。可以看出，在 y/6 = 0.08 (y+ 30) 和 y/6 = 0.3 之间非常严格地遵循对数定律。即使在通道的中心部分 (0.3 y/6 1.0)，与对数定律的偏差也相当小；但应该认识到，导致对数定律的论点不适用于该地区。

令Uo,log 表示通过对数定律外推获得的中心线上的( U ) 值。对于 y/6 1，方程(7.50) 然后得出

\begin{sourceblock}
\includegraphics[page=305,pagebox=mediabox,trim={64.39bp 244.97bp 64.39bp 1658.79bp},clip,width=0.94\textwidth,max height=0.38\textheight]{Pope.pdf}
\end{sourceblock}
这提供了一种确定 B1 的便捷方法。从图 7.9 可以看出，Uo - Uo,l,, 之差非常小，约为 Uo 的 1\%
根据对许多测量的调查，Dean (1978) 建议 B\textasciitilde{} 0.7。 B1 中的不确定性影响不大：关键是它很小。

在边界层的外层，与对数定律的偏差更大。因此，速度亏损定律将在该上下文中进一步讨论（第 7.3 节）。

\subsection*{7.1.5 摩擦定律和雷诺数}\phantomsection\addcontentsline{toc}{subsection}{7.1.5 摩擦定律和雷诺数}
描述了平均速度分布的特征后，我们现在可以确定 slkin 摩擦系数和其他量的雷诺数依赖性。主要任务是建立速度 Uo 和 u, 之间的关系。

通过使用对数定律（方程（7.50））在整个通道上近似 (U) 可以获得对体积速度 U 的良好估计。 （为了在 Y = 6 时保持一致，需要取 B1 = 0。）正如我们所见，在通道中心，与对数定律的偏离非常小（图 7.9）：靠近壁 (y+ 30) 的近似性很差（图 7.6），但该区域对积分的贡献可以忽略不计（雷诺数非常低时除外）。通过这种近似得到的结果是

\begin{sourceblock}
\includegraphics[page=306,pagebox=mediabox,trim={64.39bp 819.90bp 64.39bp 975.88bp},clip,width=0.94\textwidth,max height=0.38\textheight]{Pope.pdf}
\end{sourceblock}
这个估计与分散在 2 和 3 之间的实验数据（Dean 1978）以及 DNS 值为 2.6 非常吻合。

内层的对数定律（式（7.43））可以写成

\begin{sourceblock}
\includegraphics[page=306,pagebox=mediabox,trim={64.39bp 573.93bp 64.39bp 1325.83bp},clip,width=0.94\textwidth,max height=0.38\textheight]{Pope.pdf}
\end{sourceblock}
v 而在外层则是（方程（7.50））

\begin{sourceblock}
\includegraphics[page=306,pagebox=mediabox,trim={64.39bp 414.95bp 64.39bp 1485.81bp},clip,width=0.94\textwidth,max height=0.38\textheight]{Pope.pdf}
\end{sourceblock}
当这两个方程相加时，y 相关性消失为

\begin{sourceblock}
\includegraphics[page=306,pagebox=mediabox,trim={64.39bp 119.39bp 64.39bp 1777.78bp},clip,width=0.94\textwidth,max height=0.38\textheight]{Pope.pdf}
\end{sourceblock}
\begin{sourceblock}
\includegraphics[page=307,pagebox=mediabox,trim={64.39bp 1226.85bp 64.39bp 109.44bp},clip,width=\textwidth,max height=0.72\textheight]{Pope.pdf}
\par\vspace{0.45\baselineskip}
\small 图 7.10。槽道流的表面摩擦系数 cr = ty A5eU2) 相对于雷诺数 (Re = 2U6/v)：符号，由 Dean (1978) 编制的实验数据；实线，来自方程。 （7.55）；虚线，层流摩擦定律，cr = 16/(3Re)。
\end{sourceblock}
对于给定的 Reo，该方程可以解出 Uo/u,，从而确定皮肤摩擦系数 cf = T , / ( ; \textasciitilde{} U \textasciitilde{} ) = 2(u,/ Uo)2。 (7.52)、Re 2U6/v 和 Cf z,/(\$d,J 2 ) 然后也可以确定。

图 7.10 显示了从方程 7.10 中获得的表面摩擦系数 cf。 (7.55) 作为 Re 的函数。还显示了层流关系和 Dean (1978) 编制的实验数据。对于 Re 3,000，方程式(7.55) 很好地表示了表面摩擦系数。有趣的是，Pate1 和 Head (1969) 发现 Re 3,000 是观察到具有通用常数的对数定律的最低雷诺数。

平均流量与粘性比例的比率如图 1 和 2 所示。 7.11 和 7.12。长度比例 6/6 Re 几乎与 Re 线性增加，Re 是一个很好的近似值，Re 为 0.09Re。因此，在高雷诺数时，粘性长度尺度可能非常小。例如，对于 6 2 ' cm 的通道，在 Re lo5 下，粘度标度为 6, 10- 5 m，因此位置 y+ 100 距壁仅 1 mm。不用说，在高雷诺数实验室流的粘性壁区域进行测量存在相当大的困难。

相反，速度比随 Re 增加非常缓慢（图 7.12）

一个简单的近似是 Uo/u, 5 logloRe。因此，壁和壁之间的平均速度增加的很大一部分

\begin{sourceblock}
\includegraphics[page=308,pagebox=mediabox,trim={64.39bp 1280.84bp 64.39bp 109.44bp},clip,width=\textwidth,max height=0.72\textheight]{Pope.pdf}
\par\vspace{0.45\baselineskip}
\small 图 7.11。对于湍流槽道流，外到内长度比例为 6/6 = Re，作为雷诺数的函数（从方程 \{7.55）获得）。
\end{sourceblock}
\begin{sourceblock}
\includegraphics[page=308,pagebox=mediabox,trim={64.39bp 611.92bp 64.39bp 780.90bp},clip,width=\textwidth,max height=0.72\textheight]{Pope.pdf}
\par\vspace{0.45\baselineskip}
\small 图 7.12。湍流槽道流的外到内速度尺度比作为雷诺数的函数（从方程（7.55）获得）：实线，U/u，；虚线Uo/u；。
\end{sourceblock}
根据以上 (6 = 2 -cm, Re = lo5)i 可知，在 yf = 10 ( e . , y 0.1 mm) 处，平均速度超过中心线值 Uo 的 30\%。

图 7.13 显示了 y 位置的雷诺数依赖性；它描绘了各个区域和层。根据该图，Re 3,000 存在 loga w 区域 (306, y 0.36)，与 Patel 和 Head (1969) 的实验观察结果一致。另一方面，雷诺数需要超过 20,000 才能出现重叠

\begin{sourceblock}
\includegraphics[page=309,pagebox=mediabox,trim={64.39bp 1101.86bp 64.39bp 109.44bp},clip,width=\textwidth,max height=0.72\textheight]{Pope.pdf}
\par\vspace{0.45\baselineskip}
\small 图 7.13。湍流通道中的区域和层作为雷诺数的函数。 -
\end{sourceblock}
观察到，对数定律在重叠论证所暗示的区域之外仍然存在。

\subsection*{7.1.6 雷诺应力}\phantomsection\addcontentsline{toc}{subsection}{7.1.6 雷诺应力}
图 7.14-7.16 显示了从 Re 13,750 处槽道流的 DNS 获得的雷诺应力和一些相关统计数据。为了讨论这些统计数据，将流动分为三个区域是有用的：粘性区域
< < < < < 壁区域 (y+ 50) ;对数律区域（506、y 0.36 或在此雷诺数下 50 y+ 120）；和核心（y 0.36）。

在对数律区域中存在近似自相似性。归一化雷诺应力 (uiuj)/ k 基本上是均匀的，生产时也是如此。

PIE、耗散比和归一化平均剪切速率 S k / \&（其中 Sa(U)/ay）。其值在表 7.2 中给出。有趣的是，(uiuj)/k 的值与 Tavoularis 和 Corrsin (1981) 在均匀剪切流中测量的值相差几个百分点（参见第 157 页的表 5.4）。产生 F 和耗散几乎平衡，相比之下 k 的粘性和湍流传输非常小。

在中心线上，平均速度梯度和剪应力

\begin{sourceblock}
\includegraphics[page=310,pagebox=mediabox,trim={64.39bp 1386.83bp 64.39bp 109.44bp},clip,width=\textwidth,max height=0.72\textheight]{Pope.pdf}
\par\vspace{0.45\baselineskip}
\small 图 7.14。雷诺应力和动能根据 Re = 13,750 处槽道流 DNS 的摩擦速度对 y\textasciitilde{} 进行归一化（Kim 等人，1987 年）。
\end{sourceblock}
\begin{sourceblock}
\includegraphics[page=310,pagebox=mediabox,trim={64.39bp 816.90bp 64.39bp 680.91bp},clip,width=\textwidth,max height=0.72\textheight]{Pope.pdf}
\par\vspace{0.45\baselineskip}
\small 图 7.15。雷诺应力分布由 Re = 13,750 处槽道流 DNS 的湍动能标准化（Kim 等人，1987 年）。
\end{sourceblock}
\begin{sourceblock}
\includegraphics[page=310,pagebox=mediabox,trim={64.39bp 257.97bp 64.39bp 1239.84bp},clip,width=\textwidth,max height=0.72\textheight]{Pope.pdf}
\par\vspace{0.45\baselineskip}
\small 图 7.16。 Re = 13,750 时河道流量 DNS 的生产与耗散比 (P/e)、归一化平均剪切速率 (Sk/e) 和剪切应力相关系数 (p,.) 剖面图（Kim 等人，1987 年）。
\end{sourceblock}
\begin{center}\small 表 7.2。湍流河道流的统计数据，从 Kim 等人的 DNS 数据中获得。 (1987)，回复 = 13,750\end{center}
\begin{sourceblock}
\includegraphics[page=311,pagebox=mediabox,trim={64.39bp 1059.87bp 64.39bp 259.97bp},clip,width=\textwidth,max height=0.72\textheight]{Pope.pdf}
\end{sourceblock}
粘性壁区域 (y+50) 包含最剧烈的湍流活动。产生、耗散、湍动能和各向异性都在 y+ 小于 20 时达到峰值。我们将更详细地研究该区域的行为。

壁面处的边界条件 U 0 决定了对于较小的 y，雷诺应力偏离零的方式。对于固定的 x、z 和 t，并且对于较小的 y，脉动速度分量可以写为以下形式的泰勒级数

\begin{sourceblock}
\includegraphics[page=311,pagebox=mediabox,trim={64.39bp 655.92bp 64.39bp 1291.84bp},clip,width=0.94\textwidth,max height=0.38\textheight]{Pope.pdf}
\end{sourceblock}
\begin{sourceblock}
\includegraphics[page=311,pagebox=mediabox,trim={64.39bp 573.93bp 64.39bp 1373.83bp},clip,width=0.94\textwidth,max height=0.38\textheight]{Pope.pdf}
\end{sourceblock}
\begin{sourceblock}
\includegraphics[page=311,pagebox=mediabox,trim={64.39bp 485.94bp 64.39bp 1448.82bp},clip,width=0.94\textwidth,max height=0.38\textheight]{Pope.pdf}
\end{sourceblock}
这些系数是零均值随机变量，对于完全发展的河道流量，它们在统计上与 x、z 和 t 无关。对于 y 0，无滑移条件产生 u al 0 和 w a3 0；和类似地
= = = = 不渗透条件得出 v = a2 = 0。在壁处，由于 u 和 w 对于所有 x 和 z 均为零，因此导数 (d u / d \textasciitilde{} ) , , \textasciitilde{}a 和 ( d w / d \textasciitilde{} ) , =a\textasciitilde{}re 也为零。因此连续性方程得出

\begin{sourceblock}
\includegraphics[page=311,pagebox=mediabox,trim={64.39bp 122.98bp 64.39bp 1803.77bp},clip,width=0.94\textwidth,max height=0.38\textheight]{Pope.pdf}
\end{sourceblock}
\begin{sourceblock}
\includegraphics[page=312,pagebox=mediabox,trim={64.39bp 1373.83bp 64.39bp 109.44bp},clip,width=\textwidth,max height=0.72\textheight]{Pope.pdf}
\par\vspace{0.45\baselineskip}
\small 图 7.17。由湍流槽道流粘性壁区域中的摩擦速度归一化的雷诺应力和动能分布：Kim 等人的 DNS 数据。 （1987）。回复 = 13,750。
\end{sourceblock}
系数 b2 为零的意义在于，非常接近 \#all，存在双组分流。也就是说，为了订购 y，v 为零，而 u 和
均非零。由此产生的运动对应于平行平面中的流动
墙。 （这称为二分量流，而不是二维牵引，因为 u 和 w 在 y 方向上变化。）

只需采用该级数的乘积即可从展开式（方程（7.56）y7.58））获得雷诺应力。取 3 个为零的系数（ i . a - ? 、 a3 和 b2 ），以求导
al, x der in y :雷诺应力为

\begin{sourceblock}
\includegraphics[page=312,pagebox=mediabox,trim={64.39bp 772.90bp 64.39bp 1168.85bp},clip,width=0.94\textwidth,max height=0.38\textheight]{Pope.pdf}
\end{sourceblock}
\begin{sourceblock}
\includegraphics[page=312,pagebox=mediabox,trim={64.39bp 694.91bp 64.39bp 1247.84bp},clip,width=0.94\textwidth,max height=0.38\textheight]{Pope.pdf}
\end{sourceblock}
\begin{sourceblock}
\includegraphics[page=312,pagebox=mediabox,trim={64.39bp 614.92bp 64.39bp 1326.83bp},clip,width=0.94\textwidth,max height=0.38\textheight]{Pope.pdf}
\end{sourceblock}
\begin{sourceblock}
\includegraphics[page=312,pagebox=mediabox,trim={64.39bp 536.93bp 64.39bp 1407.82bp},clip,width=0.94\textwidth,max height=0.38\textheight]{Pope.pdf}
\end{sourceblock}
I l \textasciitilde{} u sw、hile (u2)、(w2) 和 k 从零开始增加，而 -(uv) 和 (v2) 分别随着 y3 和 y4 的增加而缓慢增加。这些行为可以清楚地看到 (uiuj) 与 y 的 n 对数对数分区（未显示），并且它们在 'ig 中也很明显。 7.17，显示了粘性壁区域中 (uiuj) 和 k 的分布。

对于充分发展的槽道流，湍流动力学的平衡方程

\begin{sourceblock}
\includegraphics[page=312,pagebox=mediabox,trim={64.39bp 167.98bp 64.39bp 1695.79bp},clip,width=0.94\textwidth,max height=0.38\textheight]{Pope.pdf}
\end{sourceblock}
\begin{sourceblock}
\includegraphics[page=313,pagebox=mediabox,trim={64.39bp 1352.83bp 64.39bp 109.44bp},clip,width=\textwidth,max height=0.72\textheight]{Pope.pdf}
\par\vspace{0.45\baselineskip}
\small 图 7.18。槽道流粘性壁区域中的湍动能收支：方程中的项。 （7.64）通过粘性尺度标准化。来自 Kim 等人的 DNS 数据。 （1987）。回复 = 13,750。
\end{sourceblock}
墙区域。这些术语依次为产生、赝耗散、粘性扩散、湍流对流和压力传递。

与 -(uv) 类似，产量 P 从零开始增加为 y3。它在缓冲层内的 y+ 12 处达到峰值。事实上，可以证明（练习 7.6），峰值产量恰好发生在粘性应力和雷诺剪切应力相等的地方。在这座山峰周围，
（PIE= 产量超过耗散 1.8，产生的多余能量被输送走。压力输送很小，而湍流对流将能量输送到壁面并进入对数定律区域。粘性输送 v d 2k dy2 将动能一直输送到壁面。

也许令人惊讶的是，峰值耗散发生在壁处，那里的动能为零。尽管脉动速度在 y 0 处消失，但脉动应变率和耗散却没有消失（练习 7.7）。

sij 壁面处的耗散通过粘性传输来平衡，

\begin{sourceblock}
\includegraphics[page=313,pagebox=mediabox,trim={64.39bp 498.94bp 64.39bp 1398.82bp},clip,width=0.94\textwidth,max height=0.38\textheight]{Pope.pdf}
\end{sourceblock}
等式中的其他项。 (7.64) 为零。 （另请参见练习 7.5 和 7.7。）

对于完全湍流，此处考虑的统计数据（通过粘性尺度归一化）仅对内层中的雷诺数 (y/6 0.1) 具有微弱的依赖性。图 7.19 显示了均方根的曲线。 u 和 v 在不同雷诺数下测量。 u'/u 的峰值与 Re 无关；但在 yf 50 处（除最低雷诺数外，均位于内层），ur/u 的值在 Reo 14,914 和 Reo 39,582 之间增加 20\%。这些和其他雷诺数效应

\begin{sourceblock}
\includegraphics[page=314,pagebox=mediabox,trim={64.39bp 1280.84bp 64.39bp 109.44bp},clip,width=\textwidth,max height=0.72\textheight]{Pope.pdf}
\par\vspace{0.45\baselineskip}
\small 图 7.19。 rm.s 简介Wei 和 Willmarth (1989) 在不同雷诺数下测量了渠道流中的速度。开放符号：u//u, = (u’)’/?/u,; O，Reg = 2,970； O，Reo = 14,914； A、Reo = 22,776； V，Reo = 39,582。实心符号：v’/u, = (v?)!/? /u，在相同的雷诺数处。
\end{sourceblock}
\begin{sourceblock}
\includegraphics[page=314,pagebox=mediabox,trim={64.39bp 851.89bp 64.39bp 1044.87bp},clip,width=0.94\textwidth,max height=0.38\textheight]{Pope.pdf}
\end{sourceblock}
\begin{sourceblock}
\includegraphics[page=314,pagebox=mediabox,trim={64.39bp 691.91bp 64.39bp 1158.85bp},clip,width=0.94\textwidth,max height=0.38\textheight]{Pope.pdf}
\end{sourceblock}
\begin{sourceblock}
\includegraphics[page=314,pagebox=mediabox,trim={64.39bp 291.96bp 64.39bp 1605.80bp},clip,width=0.94\textwidth,max height=0.38\textheight]{Pope.pdf}
\end{sourceblock}
剪切 stresbsc 可写

\begin{sourceblock}
\includegraphics[page=315,pagebox=mediabox,trim={64.39bp 1703.79bp 64.39bp 214.97bp},clip,width=0.94\textwidth,max height=0.38\textheight]{Pope.pdf}
\end{sourceblock}
因此生产率是

\begin{sourceblock}
\includegraphics[page=315,pagebox=mediabox,trim={64.39bp 1543.81bp 64.39bp 373.95bp},clip,width=0.94\textwidth,max height=0.38\textheight]{Pope.pdf}
\end{sourceblock}
从 P 的这个表达式，表明峰值产量发生
在位置 jl 处，粘性应力和雷诺应力相等。

证明这个峰值是

\begin{sourceblock}
\includegraphics[page=315,pagebox=mediabox,trim={64.39bp 1270.84bp 64.39bp 638.92bp},clip,width=0.94\textwidth,max height=0.38\textheight]{Pope.pdf}
\end{sourceblock}
显示静止固体壁处的应变波动率

\begin{sourceblock}
\includegraphics[page=315,pagebox=mediabox,trim={64.39bp 855.89bp 64.39bp 1038.87bp},clip,width=0.94\textwidth,max height=0.38\textheight]{Pope.pdf}
\end{sourceblock}
其中 bl 和 b3 是方程中的系数。 （7.56）和（7.58）。因此，
根据 Hanjalik 和 Launder (1976) 得到以下结果：

\begin{sourceblock}
\includegraphics[page=315,pagebox=mediabox,trim={64.39bp 543.93bp 64.39bp 1292.84bp},clip,width=0.94\textwidth,max height=0.38\textheight]{Pope.pdf}
\end{sourceblock}
证明 和 “在 y 0 处相等。

\begin{sourceblock}
\includegraphics[page=315,pagebox=mediabox,trim={64.39bp 291.96bp 64.39bp 1512.81bp},clip,width=0.94\textwidth,max height=0.38\textheight]{Pope.pdf}
\end{sourceblock}
使用图 7.18 来估计 Kolmogorov 的值
墙上的比例是

\begin{sourceblock}
\includegraphics[page=315,pagebox=mediabox,trim={64.39bp 119.39bp 64.39bp 1805.77bp},clip,width=0.94\textwidth,max height=0.38\textheight]{Pope.pdf}
\end{sourceblock}
壁面雷诺剪应力的展开式（方程（7.63））

可以写成

\begin{sourceblock}
\includegraphics[page=316,pagebox=mediabox,trim={64.39bp 1694.79bp 64.39bp 250.97bp},clip,width=0.94\textwidth,max height=0.38\textheight]{Pope.pdf}
\end{sourceblock}
其中无量纲系数可以假设为inde-
雷诺数的下垂。从动量方程可知
化（方程（7.10）和（7.13））这意味着以下展开

\begin{sourceblock}
\includegraphics[page=316,pagebox=mediabox,trim={64.39bp 1402.82bp 64.39bp 442.94bp},clip,width=0.94\textwidth,max height=0.38\textheight]{Pope.pdf}
\end{sourceblock}
为什么墙定律的展开式是

\begin{sourceblock}
\includegraphics[page=316,pagebox=mediabox,trim={64.39bp 1278.84bp 64.39bp 643.92bp},clip,width=0.94\textwidth,max height=0.38\textheight]{Pope.pdf}
\end{sourceblock}
\subsection*{7.1.7 长度尺度和混合长度}\phantomsection\addcontentsline{toc}{subsection}{7.1.7 长度尺度和混合长度}
对数律域的三个基本属性是 .themean 的形式

\begin{sourceblock}
\includegraphics[page=316,pagebox=mediabox,trim={64.39bp 918.88bp 64.39bp 926.88bp},clip,width=0.94\textwidth,max height=0.38\textheight]{Pope.pdf}
\end{sourceblock}
事实上生产和消耗几乎是平衡的，

\begin{sourceblock}
\includegraphics[page=316,pagebox=mediabox,trim={64.39bp 802.90bp 64.39bp 1143.86bp},clip,width=0.94\textwidth,max height=0.38\textheight]{Pope.pdf}
\end{sourceblock}
以及归一化雷诺剪应力接近恒定，

\begin{sourceblock}
\includegraphics[page=316,pagebox=mediabox,trim={64.39bp 671.91bp 64.39bp 1265.84bp},clip,width=0.94\textwidth,max height=0.38\textheight]{Pope.pdf}
\end{sourceblock}
从这三个属性中得出的第四个属性是湍流与平均剪切时间尺度比率的接近恒定性

\begin{sourceblock}
\includegraphics[page=316,pagebox=mediabox,trim={64.39bp 471.94bp 64.39bp 1427.82bp},clip,width=0.94\textwidth,max height=0.38\textheight]{Pope.pdf}
\end{sourceblock}
根据这些关系，可以通过代数推导出湍流长度尺度*L k 3 / 2 /v\textasciitilde{}a 如下

\begin{sourceblock}
\includegraphics[page=316,pagebox=mediabox,trim={64.39bp 260.97bp 64.39bp 1635.79bp},clip,width=0.94\textwidth,max height=0.38\textheight]{Pope.pdf}
\end{sourceblock}
在高雷诺数时，在重叠区域 (506, 0.16)，雷诺应力基本恒定，因此 L 随 y 线性变化：

\begin{sourceblock}
\includegraphics[page=316,pagebox=mediabox,trim={64.39bp 119.39bp 64.39bp 1844.77bp},clip,width=0.94\textwidth,max height=0.38\textheight]{Pope.pdf}
\end{sourceblock}
\begin{sourceblock}
\includegraphics[page=317,pagebox=mediabox,trim={64.39bp 1702.79bp 64.39bp 163.98bp},clip,width=0.94\textwidth,max height=0.38\textheight]{Pope.pdf}
\end{sourceblock}
请注意，S 和 与 y 成反比变化，而 L 和 z = k/\& 与 y 成线性变化。

（\textasciitilde{}DNS 中可获取中等雷诺数，没有重叠区域，并且剪应力在对数律区域内明显变化。这与近似方程 (7.79)-(7.8 1) 中的缺陷一起导致方程 (7.84) 为从 DNS 获得的 L 提供了较差的近似值。）

定义湍流黏度 vT(y)，以便雷诺剪切应力由下式给出

\begin{sourceblock}
\includegraphics[page=317,pagebox=mediabox,trim={64.39bp 1245.84bp 64.39bp 651.92bp},clip,width=0.94\textwidth,max height=0.38\textheight]{Pope.pdf}
\end{sourceblock}
它可以表示为速度尺度 u" 和长度尺度的乘积

\begin{sourceblock}
\includegraphics[page=317,pagebox=mediabox,trim={64.39bp 1095.86bp 64.39bp 802.90bp},clip,width=0.94\textwidth,max height=0.38\textheight]{Pope.pdf}
\end{sourceblock}
可以随意指定其中一个尺度，然后再确定另一个尺度
一个有利的（隐含的）规范是 VT。

\begin{sourceblock}
\includegraphics[page=317,pagebox=mediabox,trim={64.39bp 926.88bp 64.39bp 1016.87bp},clip,width=0.94\textwidth,max height=0.38\textheight]{Pope.pdf}
\end{sourceblock}
通过代入等式。 (7.87) 和 (7.88) 代入方程。 (7.86) 并取绝对值，我们得到显式关系

\begin{sourceblock}
\includegraphics[page=317,pagebox=mediabox,trim={64.39bp 716.91bp 64.39bp 1204.85bp},clip,width=0.94\textwidth,max height=0.38\textheight]{Pope.pdf}
\end{sourceblock}
（在通道的上半部分 (6 y 26)，速度梯度 d(U)/dy 为负，雷诺应力 (uv) 为正。方程 (7.88) 和 (7.89) 中的绝对值确保 u" 对于所有 y 均为非负值。）

在高雷诺数下出现的重叠区域 (506, y 0.16) 中，剪切应力 -(uv) 与 u2 相差不大，平均速度梯度为 u,/(rcy )。因此，urne 等于 u,，然后方程。 (7.89) 确定

\begin{sourceblock}
\includegraphics[page=317,pagebox=mediabox,trim={64.39bp 342.96bp 64.39bp 1549.80bp},clip,width=0.94\textwidth,max height=0.38\textheight]{Pope.pdf}
\end{sourceblock}
em 与 L k3 / 2 一样，长度尺度随 y 线性变化。

上述关系构成了普朗特的混合长度假设（Prandtl 1925）。总之，湍流黏度由下式给出

\begin{sourceblock}
\includegraphics[page=317,pagebox=mediabox,trim={64.39bp 119.39bp 64.39bp 1806.77bp},clip,width=0.94\textwidth,max height=0.38\textheight]{Pope.pdf}
\end{sourceblock}
ern 其中是混合长度。在重叠区域中，随 y 线性变化，比例常数为卡门常数，
为了使用混合长度假设作为湍流模型，它
ern 需要指定重叠区域之外，即粘性壁区域和核心区域。该主题的讨论推迟到第 7.3 节。

从方程式。 (7.79)-(7.80)，获得 Kol- 的以下估计值
对数律区域中的莫戈洛夫尺度：

\begin{sourceblock}
\includegraphics[page=318,pagebox=mediabox,trim={64.39bp 1293.84bp 64.39bp 598.92bp},clip,width=0.94\textwidth,max height=0.38\textheight]{Pope.pdf}
\end{sourceblock}
\section*{7.2 管流}\phantomsection\addcontentsline{toc}{section}{7.2 管流}
自 1883 年雷诺兹实验以来，管道流在我们对湍流理解的发展中发挥了重要作用。特别是，测量完全发展的湍流管流长度上的压降非常简单，因此可以确定表面摩擦系数 Cf。在 20 世纪 30 年代进行的一组有影响力的实验中，\textasciitilde{}ikuradse\textasciitilde{}测量了光滑管道和壁粗糙度不同的管道的 Cf 与雷诺数的函数关系。

本节的主要目的是描述壁粗糙度的影响，这在管道、通道和边界层流中是相似的。首先简要概述光滑壁情况。

\subsection*{7.2.1 光滑管道的摩擦定律}\phantomsection\addcontentsline{toc}{subsection}{7.2.1 光滑管道的摩擦定律}
我们考虑在圆形横截面的长直管中充分发展的湍流，内径见图 7.1。在极柱坐标 ( , ) 中，速度统计仅取决于径向坐标 r。平均中心线速度用 Uo 表示，

\begin{sourceblock}
\includegraphics[page=318,pagebox=mediabox,trim={64.39bp 325.96bp 64.39bp 1620.80bp},clip,width=0.94\textwidth,max height=0.38\textheight]{Pope.pdf}
\end{sourceblock}
体积速度为

\begin{sourceblock}
\includegraphics[page=318,pagebox=mediabox,trim={64.39bp 154.98bp 64.39bp 1738.78bp},clip,width=0.94\textwidth,max height=0.38\textheight]{Pope.pdf}
\end{sourceblock}
\begin{sourceblock}
\includegraphics[page=319,pagebox=mediabox,trim={64.39bp 1275.84bp 64.39bp 109.44bp},clip,width=\textwidth,max height=0.72\textheight]{Pope.pdf}
\par\vspace{0.45\baselineskip}
\small 图 7.20。完全发展的湍流管流中的平均速度剖面。符号，Zagarola 和 Smits (1997) 在六个雷诺数处的实验数据（Re \textasciitilde{} 32 x 10、99 x 10°、409 x 10°、1.79 x 108、7.71 x 10°、29.9 x 10°）。实线，对数定律，« = 0.436 且 B = 6.13；虚线，对数定律，x = 0.41，B = 5.2。
\end{sourceblock}
kLl 其中 R 是管道的半径。取6R为特征流宽，则常规定义的雷诺数为

\begin{sourceblock}
\includegraphics[page=319,pagebox=mediabox,trim={64.39bp 877.89bp 64.39bp 1025.87bp},clip,width=0.94\textwidth,max height=0.38\textheight]{Pope.pdf}
\end{sourceblock}
对于渠道流量，我们定义 y 为距壁面的距离，即4，

\begin{sourceblock}
\includegraphics[page=319,pagebox=mediabox,trim={64.39bp 709.91bp 64.39bp 1237.84bp},clip,width=0.94\textwidth,max height=0.38\textheight]{Pope.pdf}
\end{sourceblock}
大量的实验数据表明，内部区域的平均速度剖面(y/6 0.1)符合壁面普遍定律u+ f,(yf )。图 7.20 显示了 Zagarola 和 Smits (1997) 在雷诺数从 Re 30 x lo3 到 Re 30 x lo6 的完全发展的湍流管流中测量的平均速度剖面。为了进行比较，对数定律用标准常数 0.41、B 5.2) 和
（rc == 最适合数据 0.436，B 6.13）。可以看出，当 y+ > 30 时，轮廓遵循随 Re 增加的 y+ 范围内的对数定律；正如预期的那样，当接近管道中心线时，轮廓会偏离对数定律。限制为 y/R 0.1 的相同数据如图 7.21 所示。显然，对于所有雷诺数，y+ > 30 和 y/R < 0.1 时测得的速度与对数定律差别不大。

\begin{sourceblock}
\includegraphics[page=320,pagebox=mediabox,trim={64.39bp 1266.84bp 64.39bp 109.44bp},clip,width=\textwidth,max height=0.72\textheight]{Pope.pdf}
\par\vspace{0.45\baselineskip}
\small 图 7.21。完全发展的湍流管流中的平均速度剖面。符号，Zagarola 和 Smits (1997) 的实验数据，y/R < 0.1，Re 值与图 7.20 相同。直线，对数定律，« = 0.436 且 B = 6.13。
\end{sourceblock}
摩擦系数

\begin{sourceblock}
\includegraphics[page=320,pagebox=mediabox,trim={64.39bp 998.87bp 64.39bp 948.88bp},clip,width=0.94\textwidth,max height=0.38\textheight]{Pope.pdf}
\end{sourceblock}
其中 是轴向距离 L 上的压降。这只是四
Ap 乘以表面摩擦系数 Cf（见练习 7.12）。正如渠道流一样，可以通过使用对数定律来近似整个流的速度剖面来获得摩擦定律（参见练习 7.14 和 7.15）。对常数进行小幅调整后，结果（方程（7.110））为
光滑管道的普朗特摩擦定律

\begin{sourceblock}
\includegraphics[page=320,pagebox=mediabox,trim={64.39bp 562.93bp 64.39bp 1332.83bp},clip,width=0.94\textwidth,max height=0.38\textheight]{Pope.pdf}
\end{sourceblock}
这隐含地产生了作为 Re 的函数的 f。从图 7.22 中可以看出，该摩擦定律与整个湍流雷诺数范围内的实验数据非常吻合。

从极柱坐标中的雷诺方程开始

（方程（5.45）-（5.47））表明，对于充分发展的湍流管流，
剪应力

\begin{sourceblock}
\includegraphics[page=320,pagebox=mediabox,trim={64.39bp 119.39bp 64.39bp 1806.77bp},clip,width=0.94\textwidth,max height=0.38\textheight]{Pope.pdf}
\end{sourceblock}
\begin{sourceblock}
\includegraphics[page=321,pagebox=mediabox,trim={64.39bp 1251.84bp 64.39bp 109.44bp},clip,width=\textwidth,max height=0.72\textheight]{Pope.pdf}
\par\vspace{0.45\baselineskip}
\small 图 7.22。光滑管道中充分流动的摩擦系数 f 与雷诺数的关系。虚线，层流的哈根-泊肃叶摩擦定律；实线，
\end{sourceblock}
湍流的普朗特摩擦定律，方程： (7.98);符号、测量值编译
施利希廷（1979）。 （经麦格劳-希尔公司许可转载。）

被给予

\begin{sourceblock}
\includegraphics[page=321,pagebox=mediabox,trim={64.39bp 965.88bp 64.39bp 981.88bp},clip,width=0.94\textwidth,max height=0.38\textheight]{Pope.pdf}
\end{sourceblock}
其中 pw(x)i 是壁面处的平均压力（参见：通道的方程 (7.9)

流）。由此得到关系式
其中壁剪应力 zw 为正值，定义为

\begin{sourceblock}
\includegraphics[page=321,pagebox=mediabox,trim={64.39bp 630.92bp 64.39bp 1315.83bp},clip,width=0.94\textwidth,max height=0.38\textheight]{Pope.pdf}
\end{sourceblock}
（注意，在此处使用的坐标系中，(uv) 为正，
且速度梯度 d(U)/dr 为负。）

摩擦系数 f 由方程式定义。 （7.97），和皮肤-
摩擦系数 Cf 的计算公式： (7.15), 得到关系式

\begin{sourceblock}
\includegraphics[page=321,pagebox=mediabox,trim={64.39bp 366.95bp 64.39bp 1578.80bp},clip,width=0.94\textwidth,max height=0.38\textheight]{Pope.pdf}
\end{sourceblock}
\begin{sourceblock}
\includegraphics[page=321,pagebox=mediabox,trim={64.39bp 259.97bp 64.39bp 1659.79bp},clip,width=0.94\textwidth,max height=0.38\textheight]{Pope.pdf}
\end{sourceblock}
其中 u 是摩擦速度。

对于层流，求解方程。 (7.99) 表明速度剖面
速度。获得充分发展的
哈根-泊肃叶摩擦定律
层流管流，

\begin{sourceblock}
\includegraphics[page=322,pagebox=mediabox,trim={64.39bp 1675.79bp 64.39bp 248.97bp},clip,width=0.94\textwidth,max height=0.38\textheight]{Pope.pdf}
\end{sourceblock}
通过对数缺陷近似平均速度分布

定律（方程（7.50））与 BL 0，即

\begin{sourceblock}
\includegraphics[page=322,pagebox=mediabox,trim={64.39bp 1444.82bp 64.39bp 456.94bp},clip,width=0.94\textwidth,max height=0.38\textheight]{Pope.pdf}
\end{sourceblock}
获得估计

\begin{sourceblock}
\includegraphics[page=322,pagebox=mediabox,trim={64.39bp 1299.84bp 64.39bp 620.92bp},clip,width=0.94\textwidth,max height=0.38\textheight]{Pope.pdf}
\end{sourceblock}
（参见槽道流量的方程（7.52））。 （根据施利希廷（Schlichting，1979）的说法，
值4.07与实验数据更加吻合。）

\begin{sourceblock}
\includegraphics[page=322,pagebox=mediabox,trim={64.39bp 984.88bp 64.39bp 823.90bp},clip,width=0.94\textwidth,max height=0.38\textheight]{Pope.pdf}
\end{sourceblock}
假设这在轴上成立，并使用等式： （7.104），
获得摩擦定律

\begin{sourceblock}
\includegraphics[page=322,pagebox=mediabox,trim={64.39bp 730.91bp 64.39bp 1122.86bp},clip,width=0.94\textwidth,max height=0.38\textheight]{Pope.pdf}
\end{sourceblock}
\begin{sourceblock}
\includegraphics[page=322,pagebox=mediabox,trim={64.39bp 652.92bp 64.39bp 1292.84bp},clip,width=0.94\textwidth,max height=0.38\textheight]{Pope.pdf}
\end{sourceblock}
0.41 和 B 5.2。

让 jj 表示平均速度梯度的 y 位置
d(U)/dy 等于 U/6。假设位于对数律区域，

表明

\begin{sourceblock}
\includegraphics[page=322,pagebox=mediabox,trim={64.39bp 354.95bp 64.39bp 1593.80bp},clip,width=0.94\textwidth,max height=0.38\textheight]{Pope.pdf}
\end{sourceblock}
\begin{sourceblock}
\includegraphics[page=322,pagebox=mediabox,trim={64.39bp 214.97bp 64.39bp 1689.79bp},clip,width=0.94\textwidth,max height=0.38\textheight]{Pope.pdf}
\end{sourceblock}
估计 Re lo4 和 lo6 的 y/6 和 y+。

\begin{sourceblock}
\includegraphics[page=323,pagebox=mediabox,trim={64.39bp 1182.85bp 64.39bp 109.44bp},clip,width=\textwidth,max height=0.72\textheight]{Pope.pdf}
\end{sourceblock}
\subsection*{7.2.2 壁面粗糙度}\phantomsection\addcontentsline{toc}{subsection}{7.2.2 壁面粗糙度}
\begin{center}\small 图 7.23。摩擦系数 f 与雷诺数的关系，用于在各种粗糙度的管道中充分发展流动。虚线，层流摩擦定律；实线，光滑管道中湍流的普朗特摩擦定律，方程。 （7.98）； Nikuradse 的符号、测量值。 （经 McGraw-Hill 许可改编自 Schlichting (1979)。）\end{center}
粗糙度 到目前为止，我们假设通道和管道的壁是完全光滑的。当然，在实践中，每个表面都在某种程度上偏离理想状态，并且对于第一近似，这种偏离的特征在于突出或凹进的长度尺度s。对于给定的流量 ( e，给定 R，需要解决的主要问题如下：是否存在 s 值，(s*,s ay)，低于该值流量与 s 无关，因此壁是有效光滑的？对于 s s*，粗糙度如何影响流量？

Nikuradse 在管道上进行了实验，沙子尽可能紧密地粘在管壁上，颗粒尺寸 s 从 s/R = 1/15 到 s/R = 1/500。摩擦系数 f 的测量结果如图 7.23 所示。可以看出，粗糙度对层流状态的影响很小，并且对转变的影响显然也很小。然后，每个粗糙度的曲线遵循同一条线，即光滑管道的普朗特定律，直到达到某个雷诺数，然后向上转，并达到渐近线。在最高雷诺数下，摩擦因数与 Re 无关，渐近 vzilue 随 s/R 增加。

观察到的行为可以通过将壁定律扩展到包含粗糙度来解释。对于给定的表面几何形状（以便

\begin{sourceblock}
\includegraphics[page=324,pagebox=mediabox,trim={64.39bp 1703.79bp 64.39bp 161.98bp},clip,width=0.94\textwidth,max height=0.38\textheight]{Pope.pdf}
\end{sourceblock}
其中 是通用无量纲函数（参见方程（7.31））。正如之前一样，假设内层中的 Cb 与 y/6 不相关（y/6 0.1）。

在高雷诺数时，可以考虑两种极端情况。如果 s/6 非常小，则有充分的理由假设流动不受粗糙度的影响，然后恢复壁的标准定律：

\begin{sourceblock}
\includegraphics[page=324,pagebox=mediabox,trim={64.39bp 1314.83bp 64.39bp 578.93bp},clip,width=0.94\textwidth,max height=0.38\textheight]{Pope.pdf}
\end{sourceblock}
\begin{sourceblock}
\includegraphics[page=324,pagebox=mediabox,trim={64.39bp 1145.86bp 64.39bp 694.91bp},clip,width=0.94\textwidth,max height=0.38\textheight]{Pope.pdf}
\end{sourceblock}
对于大 y/6,，对粘度的依赖性消失的假设意味着 @ 渐近趋向于常数 1 / \textasciitilde{} a, n d 那么方程： (7.1 14) 积分到对数定律，即

\begin{sourceblock}
\includegraphics[page=324,pagebox=mediabox,trim={64.39bp 879.89bp 64.39bp 1008.87bp},clip,width=0.94\textwidth,max height=0.38\textheight]{Pope.pdf}
\end{sourceblock}
\begin{sourceblock}
\includegraphics[page=324,pagebox=mediabox,trim={64.39bp 713.91bp 64.39bp 1120.86bp},clip,width=0.94\textwidth,max height=0.38\textheight]{Pope.pdf}
\end{sourceblock}
是一个普适常数。

在第二种极端情况下，粗糙度尺度 s 与粘性尺度 6 相比较大。那么粗糙度元素上的流动的局部雷诺数就很大 (u7s/v s/6, 1)。动量从流体到壁的传递是通过粗糙度元素上的阻力来完成的，在高雷诺数下，阻力主要是通过压力，而不是粘性应力。那么可以假设 v 和 6 不是相关参数，因此等式： (7.1 13) 可以重写

\begin{sourceblock}
\includegraphics[page=324,pagebox=mediabox,trim={64.39bp 252.97bp 64.39bp 1653.79bp},clip,width=0.94\textwidth,max height=0.38\textheight]{Pope.pdf}
\end{sourceblock}
其中 是通用无量纲函数（对于给定的粗糙度
QR 几何）。

\begin{sourceblock}
\includegraphics[page=325,pagebox=mediabox,trim={64.39bp 1243.84bp 64.39bp 109.44bp},clip,width=\textwidth,max height=0.72\textheight]{Pope.pdf}
\par\vspace{0.45\baselineskip}
\small 图 7.24。对数定律 B（方程（7.121））中的加性常数是由粘性长度 6 归一化的粗糙度尺度 s 的函数。虚线，完全粗糙 B = 8.5；实线，平滑（方程（7.122））；符号，来自 Nikuradse 的数据。 （经 McGraw-Hill 许可改编自 Schlichting (1979)。）
\end{sourceblock}
\begin{sourceblock}
\includegraphics[page=325,pagebox=mediabox,trim={64.39bp 656.92bp 64.39bp 1074.86bp},clip,width=0.94\textwidth,max height=0.38\textheight]{Pope.pdf}
\end{sourceblock}
\begin{sourceblock}
\includegraphics[page=325,pagebox=mediabox,trim={64.39bp 326.96bp 64.39bp 1518.81bp},clip,width=0.94\textwidth,max height=0.38\textheight]{Pope.pdf}
\end{sourceblock}
\begin{sourceblock}
\includegraphics[page=325,pagebox=mediabox,trim={64.39bp 155.98bp 64.39bp 1741.78bp},clip,width=0.94\textwidth,max height=0.38\textheight]{Pope.pdf}
\end{sourceblock}
\begin{sourceblock}
\includegraphics[page=326,pagebox=mediabox,trim={64.39bp 1692.79bp 64.39bp 156.98bp},clip,width=0.94\textwidth,max height=0.38\textheight]{Pope.pdf}
\end{sourceblock}
实验确实证实了粗糙壁的对数定律（方程（7.12 1）），并且根据 Nikuradse 的数据，加性常数已确定为 s/6 的函数，见图 7.24。显然，对于 s/6., 70，比如说，墙完全是
Bz p(m) 大致为 8.5。在另一个极端，测量值与方程一致。 (7.122) 高达 s/6, 5，例如，将其作为允许粗糙度的极限，即极限 s*，低于该极限，壁实际上是光滑的。

完全粗糙情况的对数定律得出精确的摩擦定律（练习 7.17 的方程（7.124）），给出摩擦系数 f 作为粗糙度 s/R 的函数（与雷诺数无关）。

证明完全粗糙墙的对数定律(7.119) 连同
等式。 (7.107) 得出摩擦定律

\begin{sourceblock}
\includegraphics[page=326,pagebox=mediabox,trim={64.39bp 878.89bp 64.39bp 1008.87bp},clip,width=0.94\textwidth,max height=0.38\textheight]{Pope.pdf}
\end{sourceblock}
将此定律与图 7.23 中的实验数据进行比较。

（Schlichting (1979) 建议稍微修改 2.0 和
1.74 代替 1.99 和 1.71。）

\section*{7.3 边界层}\phantomsection\addcontentsline{toc}{section}{7.3 边界层}
要考虑的最简单的边界层是匀速非湍流流过光滑平板时形成的边界层（见图 7.1）。与给定平均压力梯度的完全发育的槽道流相比，主要区别是：

(i) 边界层沿流动方向连续发展，
边界层厚度 6 (x) 随着 x 的增加而增加；

(ii) 壁剪应力 \textasciitilde{} , ( x ) 事先未知；和
(iii) 流动的外部部分由间歇性湍流/非-
尽管存在这些差异，内层 (y/d(x) 0.1) 的行为本质上与槽道流中的行为相同。对此进行了演示，并更详细地检查了缓冲层中的行为。在缺陷层 (y/6(x) 0 1 ) 中，与对数定律的偏离更为显着，这需要对速度-缺陷定律进行更仔细的检查。

\subsection*{7.3.1 流程描述}\phantomsection\addcontentsline{toc}{subsection}{7.3.1 流程描述}
如图 7.1 所示，坐标系与渠道流所用的坐标系相同。板的表面（i，壁）位于 x2 y 0 处，XI = x 2 0，前缘为 x = 0, y -- 0。平均流量主要在 x 方向，自由流速度（边界层外侧）由 Uo(x) 表示。统计数据主要在 y 方向上变化，并且与 z 方向无关。然而，与槽道流不同的是，边界层不断发展，因此统计数据既取决于 x 也取决于 y。速度分量为 U、V 和 W，其中 ( W ) 为零。

自由流压力 po(x) 通过伯努利方程 (Eq. (2.67)) 将 po(x) + \textasciitilde{} u \textasciitilde{} ( x c) 与速度 Uo(x) 联系起来，因此压力

\begin{sourceblock}
\includegraphics[page=327,pagebox=mediabox,trim={64.39bp 918.88bp 64.39bp 940.88bp},clip,width=0.94\textwidth,max height=0.38\textheight]{Pope.pdf}
\end{sourceblock}
加速流动 (dUo/dx 0) 对应于负面或有利
压力梯度。相反，减速流动会产生正的、逆的压力梯度，之所以这样称呼是因为它会导致边界层与表面分离。在航空应用中，通常需要附接边界层。我们的大部分注意力都集中在零压力梯度情况，对应于 Uo(x) 恒定。

边界层厚度 d(x) 通常定义为 (U(x,y )) 等于自由流速度 Uo(x) 的 99\% 时的 y 值。这是一个条件较差的量，因为它取决于微小速度差的测量。更可靠的是整体测量，例如位移厚度

\begin{sourceblock}
\includegraphics[page=327,pagebox=mediabox,trim={64.39bp 312.96bp 64.39bp 1584.80bp},clip,width=0.94\textwidth,max height=0.38\textheight]{Pope.pdf}
\end{sourceblock}
和动量厚度

\begin{sourceblock}
\includegraphics[page=327,pagebox=mediabox,trim={64.39bp 151.98bp 64.39bp 1744.78bp},clip,width=0.94\textwidth,max height=0.38\textheight]{Pope.pdf}
\end{sourceblock}
\begin{sourceblock}
\includegraphics[page=328,pagebox=mediabox,trim={64.39bp 1703.79bp 64.39bp 159.98bp},clip,width=0.94\textwidth,max height=0.38\textheight]{Pope.pdf}
\end{sourceblock}
零压力梯度边界层，从前缘 (x 0) 到 Re 达到临界点的位置存在层流
lo6 标志着过渡的开始。 （Recritv aries 的值相当大，取决于自由流中扰动的性质和水平，参见例如 Schlichting (1979)。）过渡发生在一定距离内（可能是距前缘距离的 30\%），之后边界层完全湍流。在一些实验中，将金属丝或其他装置放置在流动上，以触发层流边界层，即促进向湍流的转变。 （有关转型主题的更多信息，读者可以参阅 Arnal 和 Michel (1990)、Kachanov (1994) 及其参考文献。）

\subsection*{7.3.2 平均动量方程}\phantomsection\addcontentsline{toc}{subsection}{7.3.2 平均动量方程}
当然，边界层方程适用于流动在 x 方向上缓慢发展的情况，与横流梯度相比，轴向应力梯度较小。横向平均动量方程（Eq.（5.52））

\begin{sourceblock}
\includegraphics[page=328,pagebox=mediabox,trim={64.39bp 891.89bp 64.39bp 1012.87bp},clip,width=0.94\textwidth,max height=0.38\textheight]{Pope.pdf}
\end{sourceblock}
请注意，由于 ( v 2 ) 在壁面处为零，因此壁面压力 p,(x) 等于自由流压力 po(x)。

在边界层近似中，平均轴向动量方程

\begin{sourceblock}
\includegraphics[page=328,pagebox=mediabox,trim={64.39bp 521.93bp 64.39bp 1402.82bp},clip,width=0.94\textwidth,max height=0.38\textheight]{Pope.pdf}
\end{sourceblock}
其中 z(x,y) 是总剪应力，

\begin{sourceblock}
\includegraphics[page=328,pagebox=mediabox,trim={64.39bp 383.95bp 64.39bp 1553.80bp},clip,width=0.94\textwidth,max height=0.38\textheight]{Pope.pdf}
\end{sourceblock}
（见方程（5.55））。在壁面处，对流项为零，因此剪切应力和压力梯度平衡。如果压力梯度为零，则

\begin{sourceblock}
\includegraphics[page=328,pagebox=mediabox,trim={64.39bp 165.98bp 64.39bp 1731.78bp},clip,width=0.94\textwidth,max height=0.38\textheight]{Pope.pdf}
\end{sourceblock}
\begin{sourceblock}
\includegraphics[page=329,pagebox=mediabox,trim={64.39bp 1407.82bp 64.39bp 109.44bp},clip,width=\textwidth,max height=0.72\textheight]{Pope.pdf}
\par\vspace{0.45\baselineskip}
\small 图 7.25。平板上零压力梯度层流边界层 Blasius 解的归一化速度和剪切应力分布：y 归一化为 5，= x/Rel？ =
\end{sourceblock}
对边界层动量方程（Eq.（7.130））积分即可得到（见练习7.18）。对于
von Kcirmun 的积分动量方程在零压力梯度情况下结果为

\begin{sourceblock}
\includegraphics[page=329,pagebox=mediabox,trim={64.39bp 1013.87bp 64.39bp 888.89bp},clip,width=0.94\textwidth,max height=0.38\textheight]{Pope.pdf}
\end{sourceblock}
或者，对于皮肤摩擦系数，

\begin{sourceblock}
\includegraphics[page=329,pagebox=mediabox,trim={64.39bp 831.89bp 64.39bp 1062.87bp},clip,width=0.94\textwidth,max height=0.38\textheight]{Pope.pdf}
\end{sourceblock}
方程（7.133）量化了流的动量流速的减少或由壁剪切应力引起的动量赤字的增加。

对于层流零压力梯度边界层，方程（1）有一个相似的解。 （7.130）归因于 Blasius（1908），由 Schlichting（1979）详细描述。缩放速度 U(x, y)/ Uo 仅取决于缩放横流坐标 y/dX，其中长度尺度 6, 为 6, (xv/Uo) x/Re。该解决方案如图 7.25 所示。从溶液中获得的各种厚度为

\begin{sourceblock}
\includegraphics[page=329,pagebox=mediabox,trim={64.39bp 313.96bp 64.39bp 1593.80bp},clip,width=0.94\textwidth,max height=0.38\textheight]{Pope.pdf}
\end{sourceblock}
l he Skin-1摩擦系数是

\begin{sourceblock}
\includegraphics[page=329,pagebox=mediabox,trim={64.39bp 171.98bp 64.39bp 1775.78bp},clip,width=0.94\textwidth,max height=0.38\textheight]{Pope.pdf}
\end{sourceblock}
\begin{sourceblock}
\includegraphics[page=330,pagebox=mediabox,trim={64.39bp 1383.83bp 64.39bp 109.44bp},clip,width=\textwidth,max height=0.72\textheight]{Pope.pdf}
\par\vspace{0.45\baselineskip}
\small 图 7.26。零压力梯度湍流边界层中的平均速度、剪切应力和间歇因子的剖面，Reg = 8,000。 （来自 Klebanoff (1954) 的实验数据。）
\end{sourceblock}
根据平均连续性（方程（5.49））和动量（方程（5.55））

方程组，得到结果

\begin{sourceblock}
\includegraphics[page=330,pagebox=mediabox,trim={64.39bp 949.88bp 64.39bp 998.87bp},clip,width=0.94\textwidth,max height=0.38\textheight]{Pope.pdf}
\end{sourceblock}
从 y 0 到 y co 积分得到
从而得到

uon Kdrmdn 积分动量方程

\begin{sourceblock}
\includegraphics[page=330,pagebox=mediabox,trim={64.39bp 616.92bp 64.39bp 1277.84bp},clip,width=0.94\textwidth,max height=0.38\textheight]{Pope.pdf}
\end{sourceblock}
（回想一下 6' 和 0 是位移和动量厚 -
内斯，等式。 （7.126）和（7.127）。）

\subsection*{7.3.3 平均速度剖面}\phantomsection\addcontentsline{toc}{subsection}{7.3.3 平均速度剖面}
图 7.26 显示了 Klebanoff (1954) 在动量厚度雷诺数 Reo 8,000 的零压力梯度湍流边界层中测量的平均速度、剪切应力和间歇因子 y（在第 5.5.2 节中定义）的分布图。请注意，平均速度剖面上升

\begin{sourceblock}
\includegraphics[page=331,pagebox=mediabox,trim={64.39bp 1323.83bp 64.39bp 109.44bp},clip,width=\textwidth,max height=0.72\textheight]{Pope.pdf}
\par\vspace{0.45\baselineskip}
\small 图 7.27。以墙单位表示的平均速度分布。 Circles，Klebanoff (1954) 的边界层实验，Reg = 8,000；虚线，Spalart 的边界层 DNS (1988)，Reg = 1,410；点划线，Kim 等人的槽道流 DNS。 (1987)，Re = 13,750；实线，范德里斯特墙定律，方程。 (7.144)-(7.145)。
\end{sourceblock}
然后从墙上看就更平坦了。平均速度分布的“平坦度”通过形状因子 H 进行量化，形状因子 H 定义为位移和动量厚度之比：

\begin{sourceblock}
\includegraphics[page=331,pagebox=mediabox,trim={64.39bp 951.88bp 64.39bp 993.87bp},clip,width=0.94\textwidth,max height=0.38\textheight]{Pope.pdf}
\end{sourceblock}
对于 Blasius 剖面，平坦度因子为 H 2.6：对于 Klebanoff 边界层，平坦度因子为 H 1.3。平均速度剖面与河道流中的速度剖面非常相似（图 7.2）；并且，正如槽道流一样，随着雷诺数的增加，(U)/Uo 相对 y/d 绘制的轮廓在壁面处变陡，远离壁面变得更平坦。相应地，形状因子H随着Re的减小而减小。

剪切应力分布与相应的层流分布相似（图 7.29，尽管剪切应力的来源完全不同）。现在更详细地检查各层中的速度分布。

的

对于三种流动，由粘性尺度 u 和 d 归一化的壁速度剖面定律如图 7.27 中所示： Klebanoff 边界层（Res - 8,000）；由 Spalart (1988) 执行的边界层 DNS (Ree 1,410)；和 \textasciitilde{} ha \textasciitilde{} nfelolw (Re 13,750)。剖面之间的一致性说明了墙体定律的普遍性，不仅在对数定律区域，而且在缓冲层中。在粘性子层 (y+ 5) 中，壁定律为 f,(y+) y+，而对数定律适用于 y+ 大于 30 或 50（且
Y < 50）？ van Driest (19 56) 在混合长度假设的背景下提供了一个纯粹的经验答案，但仍然是一个受启发的答案。

对于边界层（d ( U ) / d y 0）；根据混合长度假设，总剪应力为

\begin{sourceblock}
\includegraphics[page=332,pagebox=mediabox,trim={64.39bp 1462.82bp 64.39bp 334.96bp},clip,width=0.94\textwidth,max height=0.38\textheight]{Pope.pdf}
\end{sourceblock}
随着定义

\begin{sourceblock}
\includegraphics[page=332,pagebox=mediabox,trim={64.39bp 1358.83bp 64.39bp 589.93bp},clip,width=0.94\textwidth,max height=0.38\textheight]{Pope.pdf}
\end{sourceblock}
当它通过粘性尺度标准化时，方程： (7.140) 变为

\begin{sourceblock}
\includegraphics[page=332,pagebox=mediabox,trim={64.39bp 1181.85bp 64.39bp 708.91bp},clip,width=0.94\textwidth,max height=0.38\textheight]{Pope.pdf}
\end{sourceblock}
这是 du+/dyf 的二次方程，其解为

\begin{sourceblock}
\includegraphics[page=332,pagebox=mediabox,trim={64.39bp 1020.87bp 64.39bp 879.89bp},clip,width=0.94\textwidth,max height=0.38\textheight]{Pope.pdf}
\end{sourceblock}
在内层中，z/zW之比基本为1，因此壁面定律是根据混合长度对式（1）的积分得到的。 （7.143）：

\begin{sourceblock}
\includegraphics[page=332,pagebox=mediabox,trim={64.39bp 807.90bp 64.39bp 1083.86bp},clip,width=0.94\textwidth,max height=0.38\textheight]{Pope.pdf}
\end{sourceblock}
在对数律区域中，混合长度的适当规格为 em
或等效的l+（方程（7.90））。如果相同规格
ICY rn ICY+用于粘性底层，隐含湍流应力 d ( U ) / d y 将
VT随着y2增加，而-(uv)+缓慢增加矿石，作为y3。那么显然，
靠近墙壁的规格需要减小或阻尼。据此，van Driest (1956) 提出了规范

\begin{sourceblock}
\includegraphics[page=332,pagebox=mediabox,trim={64.39bp 457.94bp 64.39bp 1457.82bp},clip,width=0.94\textwidth,max height=0.38\textheight]{Pope.pdf}
\end{sourceblock}
其中 A+ - 是一个常数，其值为 A+ 26。 平方项
( [ I ) 括号.s 称为 van Driest 阻尼函数。

方程（7.144）和（7.143）给出的壁面定律如图 7.27 所示。显然，它提供了极好的数据表示。

对于较大的yf，阻尼函数趋于统一，恢复对数定律。请注意（对于给定的情况，存在一对一的对应关系
IC) 在对数中 van Driest 常数 A+ 和加性常数 B 之间
唤醒贡献

图 7.28。湍流边界层中的平均速度剖面显示了该定律
的觉醒。符号，Klebanoff (1954) 的实验数据；虚线，对数定律
(x = 0.41，B = 5.2)；点-虚线，尾流贡献（II = 0.5）；坚硬的
线，对数定律和尾流贡献之和（方程（7.148））。

墙定律的普遍性在文献中已被广泛考虑（参见 Bradshaw 和 Huang（1995））。在具有各种压力梯度 dpo/dx（以及剪切应力梯度 (0t/dy)）-9）的边界层和管道流中，观察到对数定律 « 接近 = 0.41。然而，当壁面剪应力梯度变化时，A* 显着增加
超过 2 倍（Huffman 和 Bradshaw 1972）。

练习 7.19 证明，根据混合长度假设，非常接近
壁 (yt < 1) 雷诺剪应力为

\begin{sourceblock}
\includegraphics[page=333,pagebox=mediabox,trim={64.39bp 657.92bp 64.39bp 1239.84bp},clip,width=0.94\textwidth,max height=0.38\textheight]{Pope.pdf}
\end{sourceblock}
表明 van Driest 规范（方程 (7.145)）得出

\begin{sourceblock}
\includegraphics[page=333,pagebox=mediabox,trim={64.39bp 497.94bp 64.39bp 1399.82bp},clip,width=0.94\textwidth,max height=0.38\textheight]{Pope.pdf}
\end{sourceblock}
将此结果与 (uv) 对 y 的正确依赖性进行对比（对于非常
小 y)，等式。 （7.63）。

速度-缺陷定律 在缺陷层（例如 y/é > 0.2）中，平均速度偏离对数定律。这可以在图 7.27 中看到，并且在图 7.28 中看得更清楚。

根据对边界层数据的广泛检查，Coles (1956)

可以很好地用两个函数的和来表示。第一个函数是墙定律 f,(y+)，它取决于 y/6：第二个函数称为尾流定律，取决于 y/6。这个表示被写成

\begin{sourceblock}
\includegraphics[page=334,pagebox=mediabox,trim={64.39bp 1591.80bp 64.39bp 305.96bp},clip,width=0.94\textwidth,max height=0.38\textheight]{Pope.pdf}
\end{sourceblock}
尾流函数 w(y/6) 被假定为通用的（对于所有边界层都相同），并被定义为满足归一化条件 \textasciitilde{} ( 0 ) 0 和 w(1) 2。Coles (1956) 将 w(y/6) 制成表格（基于实验数据），但更方便的近似是

\begin{sourceblock}
\includegraphics[page=334,pagebox=mediabox,trim={64.39bp 1320.83bp 64.39bp 626.92bp},clip,width=0.94\textwidth,max height=0.38\textheight]{Pope.pdf}
\end{sourceblock}
无量纲量称为尾流强度参数，其值与流量相关。

对于克莱巴诺夫边界层，图 7.28 说明了 (U)/u 的表示形式，即这两个“定律”的总和。虚线是 log 1（即，f 的极好近似，对于 y+ > 50，y/6 < 0.2），而点划线是尾流定律。它们的总和，即实线，与实验数据非常吻合。顾名思义，函数 w(y/6) 的形状类似于平面尾流中的速度剖面，对称平面位于 y = 0。但是，这两个流之间并不存在详细的相似性。

方程（7.148）也可以写成速度-缺陷定律的形式。通过对数定律近似 f，并施加条件 (U)y=n；哦，

\begin{sourceblock}
\includegraphics[page=334,pagebox=mediabox,trim={64.39bp 628.92bp 64.39bp 1225.85bp},clip,width=0.94\textwidth,max height=0.38\textheight]{Pope.pdf}
\end{sourceblock}
该定律与图 7.29 中 Klebanoff 的数据进行了比较。

使用相同的近似值，方程。 (7.148) 在 y 6 处评估导致

\begin{sourceblock}
\includegraphics[page=334,pagebox=mediabox,trim={64.39bp 386.95bp 64.39bp 1464.82bp},clip,width=0.94\textwidth,max height=0.38\textheight]{Pope.pdf}
\end{sourceblock}
\begin{sourceblock}
\includegraphics[page=334,pagebox=mediabox,trim={64.39bp 248.97bp 64.39bp 1639.79bp},clip,width=0.94\textwidth,max height=0.38\textheight]{Pope.pdf}
\end{sourceblock}
对于给定的 Red，该方程可以解出 u7/ Uo，从而确定皮肤摩擦系数 \textasciitilde{} ( U , / U \textasciitilde{}M)o\textasciitilde{}re。方便的显式形式

\begin{sourceblock}
\includegraphics[page=335,pagebox=mediabox,trim={64.39bp 1346.83bp 64.39bp 109.44bp},clip,width=\textwidth,max height=0.72\textheight]{Pope.pdf}
\par\vspace{0.45\baselineskip}
\small 图 7.29。速度亏损定律。符号，Klebanoff (1954)的实验数据：虚线，对数定律；实线，对数定律和尾流贡献之和
\end{sourceblock}
\begin{sourceblock}
\includegraphics[page=335,pagebox=mediabox,trim={64.39bp 743.91bp 64.39bp 707.91bp},clip,width=\textwidth,max height=0.72\textheight]{Pope.pdf}
\par\vspace{0.45\baselineskip}
\small 图 7.30。从湍流边界层的直接数值模拟推导出湍流黏度和混合长度（Spalart 1988）。实线，来自 DNS 的 VR；点划线，4，，来自 DNS；虚线和 vy 根据 van Driest 的规范（方程（7.145））。
\end{sourceblock}
舒尔茨-格鲁诺公式，

\begin{sourceblock}
\includegraphics[page=335,pagebox=mediabox,trim={64.39bp 465.94bp 64.39bp 1480.81bp},clip,width=0.94\textwidth,max height=0.38\textheight]{Pope.pdf}
\end{sourceblock}
在缺陷层中，剪应力r(y)小于r，且速度梯度i ? (\textasciitilde{})/我？ yi 大于由对数定律给出的值 u , / ( \textasciitilde{} y )。显然，湍流黏度 z/(i?(U)/dy) 的值较小
VT 比对数律公式 u , \textasciitilde{} y 给出的值大；因此，
VT混合长度小于冰。图 7.30 证实了这一点，该图显示了从 Spalart 的 DNS 边界层数据推导出来的 vT(y) 和 l,(y)

因此要修改缺陷层中的公式￣,=xy。 Escudier (1966) 提出了一个简单的修改，即设置为 xy 和 0.096 的最小值。混合长度模型的其他变体（例如，Smith 和 Cebeci (1967) 以及 Baldwin 和 Lomax (1978)）通过不同的方式实现了类似的效果（参见 Wilcox (1993)）。

\subsection*{7.3.4 重新考虑重叠区域}\phantomsection\addcontentsline{toc}{subsection}{7.3.4 重新考虑重叠区域}
考虑到近壁湍流运动及其所涉及过程的复杂性，管流、槽道流和边界层中的平均速度剖面如何用简单的公式（尤其是通用对数定律）很好地表示是值得注意的。然而，理论的实证成功并不一定意味着它所基于的假设的有效性：不同的假设可能会导致预测的准确性相当。

多年来，Barenblatt 和 Monin (1979)、Long 和 Chen (1981)、George、Castillo 和 Knecht (1996) 等人提出了反对对数定律的论点。一个中心问题是重叠区域雷诺数的影响。在导致对数定律的论证中做出的强有力的假设（参见方程 (7.41)）是（对于 Re > 1、yt > 1 和 y/5 < 1）yéut/éy 独立于 Uo、6 和 v，因此也独立于雷诺数。从实验中可以清楚地看出（例如，参见图 7.19 以及 Gad-el-Hak 和 Bandyopadhyay (1994)），重叠区域中的雷诺应力分布取决于雷诺数，因此所涉及的湍流过程并不完全独立于 Re。

现在，我们在高雷诺数下充分发展的湍流管流的背景下考虑较弱的替代假设。内层速度剖面为

\begin{sourceblock}
\includegraphics[page=336,pagebox=mediabox,trim={64.39bp 550.93bp 64.39bp 1393.82bp},clip,width=0.94\textwidth,max height=0.38\textheight]{Pope.pdf}
\end{sourceblock}
（参见方程（7.37）），其中函数 f；可能取决于雷诺数。在外层，缺陷定律被重写

\begin{sourceblock}
\includegraphics[page=336,pagebox=mediabox,trim={64.39bp 354.95bp 64.39bp 1566.80bp},clip,width=0.94\textwidth,max height=0.38\textheight]{Pope.pdf}
\end{sourceblock}
其中 9 = y/6（参见方程（7.46）），其中 u 是外层的速度尺度，可能与 u 不同，并且函数 F 可能取决于雷诺数。

在重叠区域 (6, < y < 6) 中，f,(yt) 的渐近线（对于

\begin{sourceblock}
\includegraphics[page=337,pagebox=mediabox,trim={64.39bp 1014.87bp 64.39bp 109.44bp},clip,width=\textwidth,max height=0.72\textheight]{Pope.pdf}
\par\vspace{0.45\baselineskip}
\small 图 7.31。六个雷诺数（从左到右：Re 32 x 10°、99 x 103、409 x 10°、1.79 x 109、7,71 x 10° 和 29.9 x 10°）的湍流管流中平均速度剖面的双对数图。 u* 的标度属于最低雷诺数：后续轮廓连续向下移动 1.1 倍。显示的范围是重叠区域，506，< y < 0.1 R。符号，Zagarola 的实验数据和
\end{sourceblock}
斯密茨（1997）；虚线，对数定律，0.436，B = 6.13；实线，力量
定律（方程（7.157）），其功效由数据的最佳拟合确定。

符合此匹配要求的配置文件：滞后定律

\begin{sourceblock}
\includegraphics[page=337,pagebox=mediabox,trim={64.39bp 598.92bp 64.39bp 1308.83bp},clip,width=0.94\textwidth,max height=0.38\textheight]{Pope.pdf}
\end{sourceblock}
和幂律

\begin{sourceblock}
\includegraphics[page=337,pagebox=mediabox,trim={64.39bp 472.94bp 64.39bp 1470.81bp},clip,width=0.94\textwidth,max height=0.38\textheight]{Pope.pdf}
\end{sourceblock}
（参见练习 7.20 和 Barenblatt (1993)）。此开发中所做的假设允许正系数 B 、 a 和 C 取决于雷诺数。相反，如果系数与雷诺数无关，则称这些定律是通用的。

图 7.31 显示了在不同雷诺数下湍流管流重叠区域中测量的平均速度剖面。将数据与通用对数定律（0.436 和 B 6.13）和幂律进行比较

\begin{sourceblock}
\includegraphics[page=338,pagebox=mediabox,trim={64.39bp 1155.86bp 64.39bp 109.44bp},clip,width=\textwidth,max height=0.72\textheight]{Pope.pdf}
\par\vspace{0.45\baselineskip}
\small 图 7.32。管道流量的幂律关系 u* = C(y*)t = 中的指数 « = I/n (方程 (7.158)) 作为雷诺数的函数。
\end{sourceblock}
数量。很明显，a 随 Re 显着下降：Zagarola、Perry 和 Smits (1997) 表明，经验确定的值可以很好地近似为以下公式

\begin{sourceblock}
\includegraphics[page=338,pagebox=mediabox,trim={64.39bp 757.90bp 64.39bp 1164.85bp},clip,width=0.94\textwidth,max height=0.38\textheight]{Pope.pdf}
\end{sourceblock}
如图 7.32 所示。

从图 7.31 可以看出，对数定律和幂律都提供了测量速度剖面的相当准确的表示。尽管数据已经过仔细审查（例如，Zagarola al.

et (1997))，由于对数律和幂律预测之间的微小差异，得出的结论可能仍然存在争议（参见，例如，Barenblatt 和 Chorin (19 98)）。无论基本假设有什么优点，对数定律都具有普适性的实际优势。

出现的另一个问题是外层速度尺度 u 的适当选择。通用对数定律意味着 u, u,，而通用幂律（如图 7.31 所示，它不存在）意味着 u, Uo（参见练习 7.20）。人们提出了各种建议：例如 u, -- Uo- U 表示管道流量（Zagarola 和 Smits 1997），以及 u, = Uo
考虑在高雷诺数下充分发展的湍流管流

伯尔。在内层和外层中，给出平均速度分布
由方程式分别为（7.154）和（7.155）。

(a) 显示 ( U ) 和 d(U)/dy 的匹配由下式给出
等式。 (7.154) 和 (7.155) 收益率

\begin{sourceblock}
\includegraphics[page=339,pagebox=mediabox,trim={64.39bp 1468.81bp 64.39bp 439.94bp},clip,width=0.94\textwidth,max height=0.38\textheight]{Pope.pdf}
\end{sourceblock}
\begin{sourceblock}
\includegraphics[page=339,pagebox=mediabox,trim={64.39bp 1344.83bp 64.39bp 553.93bp},clip,width=0.94\textwidth,max height=0.38\textheight]{Pope.pdf}
\end{sourceblock}
在重叠区域 (6, << y << 6)，不同的渐近形式
（对于大 y+）和 F（对于小 q ）是从替代方案中获得的

\begin{sourceblock}
\includegraphics[page=339,pagebox=mediabox,trim={64.39bp 1112.86bp 64.39bp 747.91bp},clip,width=0.94\textwidth,max height=0.38\textheight]{Pope.pdf}
\end{sourceblock}
\begin{sourceblock}
\includegraphics[page=339,pagebox=mediabox,trim={64.39bp 983.88bp 64.39bp 899.89bp},clip,width=0.94\textwidth,max height=0.38\textheight]{Pope.pdf}
\end{sourceblock}
其中 K、C 和 a 是正常数（即 e i 独立于 y，但是
不一定是Re）。

(b) 证明方程的通解。 (7.160) 和 (7.161) 是

\begin{sourceblock}
\includegraphics[page=339,pagebox=mediabox,trim={64.39bp 706.91bp 64.39bp 1200.85bp},clip,width=0.94\textwidth,max height=0.38\textheight]{Pope.pdf}
\end{sourceblock}
（参见方程（7.43））和

\begin{sourceblock}
\includegraphics[page=339,pagebox=mediabox,trim={64.39bp 544.93bp 64.39bp 1354.83bp},clip,width=0.94\textwidth,max height=0.38\textheight]{Pope.pdf}
\end{sourceblock}
（参见方程（7.50））其中 B 和 B1 是积分常数。

Show，w 方程。 (7.159) 然后得出摩擦定律方程。 （7.55）。

(c) 证明方程的一般解。 (7.160) 和 (7.162) 是

\begin{sourceblock}
\includegraphics[page=339,pagebox=mediabox,trim={64.39bp 327.96bp 64.39bp 1616.80bp},clip,width=0.94\textwidth,max height=0.38\textheight]{Pope.pdf}
\end{sourceblock}
\begin{sourceblock}
\includegraphics[page=339,pagebox=mediabox,trim={64.39bp 247.97bp 64.39bp 1698.79bp},clip,width=0.94\textwidth,max height=0.38\textheight]{Pope.pdf}
\end{sourceblock}
\begin{sourceblock}
\includegraphics[page=339,pagebox=mediabox,trim={64.39bp 119.39bp 64.39bp 1764.78bp},clip,width=0.94\textwidth,max height=0.38\textheight]{Pope.pdf}
\end{sourceblock}
b 和 bl 是积分常数。 (7.159)

\begin{sourceblock}
\includegraphics[page=340,pagebox=mediabox,trim={64.39bp 1661.79bp 64.39bp 198.97bp},clip,width=0.94\textwidth,max height=0.38\textheight]{Pope.pdf}
\end{sourceblock}
(d) 表明方程。 (7.163) 和 (7.164) 是通用对数定律（即
系数独立于 Re) if, 和 onTy if, B, 和
B1 与 Re、u 无关，与 u 成比例，并且摩擦力
定律方程(7.55) 满足。

(e) 表明方程。 (7.165) 和 (7.166) 是普适幂律

\begin{sourceblock}
\includegraphics[page=340,pagebox=mediabox,trim={64.39bp 1092.86bp 64.39bp 577.93bp},clip,width=\textwidth,max height=0.72\textheight]{Pope.pdf}
\par\vspace{0.45\baselineskip}
\small 图 7.32.) (f) 取 « = + 并忽略 随 Re 的变化。然后证明方程。 (7.169) 摩擦系数的产量
\end{sourceblock}
\begin{sourceblock}
\includegraphics[page=340,pagebox=mediabox,trim={64.39bp 885.89bp 64.39bp 1061.87bp},clip,width=0.94\textwidth,max height=0.38\textheight]{Pope.pdf}
\end{sourceblock}
其中 是一个常数。

（注：当 b3 0.3 164 时，方程（7.170）是经验 Blasius 阻力
公式，与普朗特公式（方程（7.98））一致
lo4 < Re < lo5 为 3\%，但误差为 30 0/o。在再
洛7。它
强调这个推导错误地假设存在

普遍幂律。）

考虑幂律

\begin{sourceblock}
\includegraphics[page=340,pagebox=mediabox,trim={64.39bp 422.95bp 64.39bp 1491.81bp},clip,width=0.94\textwidth,max height=0.38\textheight]{Pope.pdf}
\end{sourceblock}
作为速度剖面的近似值（对于 0 y R）

管内形成湍流。证明面积平均速度

\begin{sourceblock}
\includegraphics[page=340,pagebox=mediabox,trim={64.39bp 193.97bp 64.39bp 1676.79bp},clip,width=0.94\textwidth,max height=0.38\textheight]{Pope.pdf}
\end{sourceblock}
图 7.33。由摩擦归一化的雷诺应力和动能分布
Reo = 1,410 时湍流边界层中的速度：(a) 穿过边界层
(b) 在粘性近壁区域。来自 Spalart (1988) 的 DNS 数据。

层，考虑幂律分布
其中 n 为正数。对于此配置文件表明

\begin{sourceblock}
\includegraphics[page=341,pagebox=mediabox,trim={64.39bp 813.90bp 64.39bp 1092.86bp},clip,width=0.94\textwidth,max height=0.38\textheight]{Pope.pdf}
\end{sourceblock}
因此形状因子是

\begin{sourceblock}
\includegraphics[page=341,pagebox=mediabox,trim={64.39bp 538.93bp 64.39bp 1369.83bp},clip,width=0.94\textwidth,max height=0.38\textheight]{Pope.pdf}
\end{sourceblock}
\subsection*{7.3.5 雷诺应力平衡}\phantomsection\addcontentsline{toc}{subsection}{7.3.5 雷诺应力平衡}
对于自由剪切流和槽道流，我们都检查了湍动能的平衡。在这里，我们使用 Spalart (1988) 的边界层 DNS 数据，进一步检查各个雷诺应力的平衡。首先，简要描述雷诺应力分布和动能平衡。

\begin{sourceblock}
\includegraphics[page=342,pagebox=mediabox,trim={64.39bp 1373.83bp 64.39bp 109.44bp},clip,width=\textwidth,max height=0.72\textheight]{Pope.pdf}
\par\vspace{0.45\baselineskip}
\small 图 7.34。 Reg = 1,410 时湍流边界层中的湍动能收支：方程式中的项。 (7.177) (a) 标准化为 y 的函数，以便各项的平方和为 1；(b) 通过粘性尺度标准化。来自 Spalart (1988) 的 DNS 数据。
\end{sourceblock}
与槽道流中的流量几乎没有什么不同（参见图 7.14 和 7.17）。当接近边界层的边缘时，所有雷诺应力平滑地趋向于零（对应于非湍流自由流）。正如在槽道流中一样，雷诺应力分布表现出弱的雷诺数依赖性（参见例如 Marusit、Uddin 和 Perry (19 97)）。

动能平衡 在边界层近似中，湍动能方程（方程（5.164））为
该平衡方程中的项依次为平均流对流、产生、赝耗散、粘性扩散、湍流对流和压力传输。该方程本质上与槽道流方程（方程（7.64））相同，但增加了平均流对流项。

对方程式的各种贡献的概况。 (7.177) 示于

图 7.34(a) 贡献已被局部归一化，因此它们的平方和为 1。从 y+ 40 到 y]b 0.4，主要平衡是生产和 Ld 消耗之间。在边界层之外，产量变小；，，-并且平衡是在耗散和各种传输项之间实现的。

雷诺应力方程 从纳维-斯托克斯方程推导出的雷诺应力的传输方程（参见练习 7.23）为

\begin{sourceblock}
\includegraphics[page=343,pagebox=mediabox,trim={64.39bp 1348.83bp 64.39bp 553.93bp},clip,width=0.94\textwidth,max height=0.38\textheight]{Pope.pdf}
\end{sourceblock}
皮吉在哪里

生产张量

\begin{sourceblock}
\includegraphics[page=343,pagebox=mediabox,trim={64.39bp 1190.85bp 64.39bp 735.91bp},clip,width=0.94\textwidth,max height=0.38\textheight]{Pope.pdf}
\end{sourceblock}
是

速度-压力-梯度张量
并且是

耗散张量

\begin{sourceblock}
\includegraphics[page=343,pagebox=mediabox,trim={64.39bp 885.89bp 64.39bp 1031.87bp},clip,width=0.94\textwidth,max height=0.38\textheight]{Pope.pdf}
\end{sourceblock}
为了将这些对称二阶张量与更熟悉的量联系起来，我们观察到雷诺应力方程的一半迹是动能方程，特别是，我们有

\begin{sourceblock}
\includegraphics[page=343,pagebox=mediabox,trim={64.39bp 655.92bp 64.39bp 1285.84bp},clip,width=0.94\textwidth,max height=0.38\textheight]{Pope.pdf}
\end{sourceblock}
\begin{sourceblock}
\includegraphics[page=343,pagebox=mediabox,trim={64.39bp 554.93bp 64.39bp 1373.83bp},clip,width=0.94\textwidth,max height=0.38\textheight]{Pope.pdf}
\end{sourceblock}
图 7.35-7.38 显示了方程式中各项的分布。 (7.178) 对于每个非零雷诺应力。这些术语依次称为平均对流、湍流对流、粘性扩散、产生、压力和耗散。首先检查法向应力 和。

正应力平衡 在简单剪切流中，a(U)/ay 是唯一重要的平均速度梯度，正应力产生式为

\begin{sourceblock}
\includegraphics[page=343,pagebox=mediabox,trim={64.39bp 119.39bp 64.39bp 1808.77bp},clip,width=0.94\textwidth,max height=0.38\textheight]{Pope.pdf}
\end{sourceblock}
\begin{sourceblock}
\includegraphics[page=344,pagebox=mediabox,trim={64.39bp 1398.82bp 64.39bp 109.44bp},clip,width=\textwidth,max height=0.72\textheight]{Pope.pdf}
\par\vspace{0.45\baselineskip}
\small 图 7.35。湍流边界层中 (uv) 的预算：条件和归一化与图 7.34 中的相同。
\end{sourceblock}
\begin{sourceblock}
\includegraphics[page=344,pagebox=mediabox,trim={64.39bp 827.90bp 64.39bp 658.92bp},clip,width=\textwidth,max height=0.72\textheight]{Pope.pdf}
\par\vspace{0.45\baselineskip}
\small 图 7.36。湍流边界层中(v?)的预算：条件和归一化与图7.34相同。
\end{sourceblock}
\begin{sourceblock}
\includegraphics[page=344,pagebox=mediabox,trim={64.39bp 649.92bp 64.39bp 1296.84bp},clip,width=0.94\textwidth,max height=0.38\textheight]{Pope.pdf}
\end{sourceblock}
即，所有动能产生都在 (u2) 中。因此，正如预期的那样，在大部分边界层上，plli 是 (u2) 的主要来源（见图 7.35）。

在湍动能平衡中，p'仅作为输运出现
inii 项（即 - -V (up'lp)），它在大部分边界层上相对较小（见图 7.34）。相反，在雷诺应力方程中，压力项起着核心作用。在大部分边界层上
是 (u2) 平衡中的主要汇，而 和 是
n22 n33 (v2) 和 (w2) 余额中的主要来源。因此脉动压力的主要作用是在各成分之间重新分配能量
从 (u2) 提取能量并将其转移到 (v2) 和*(w2)。

\begin{sourceblock}
\includegraphics[page=345,pagebox=mediabox,trim={64.39bp 1397.82bp 64.39bp 109.44bp},clip,width=\textwidth,max height=0.72\textheight]{Pope.pdf}
\par\vspace{0.45\baselineskip}
\small 图 7.37。湍流边界层中(w?)的预算：条件和归一化与图7.34中的相同。
\end{sourceblock}
\begin{sourceblock}
\includegraphics[page=345,pagebox=mediabox,trim={64.39bp 833.90bp 64.39bp 658.92bp},clip,width=\textwidth,max height=0.72\textheight]{Pope.pdf}
\par\vspace{0.45\baselineskip}
\small 图 7.38。湍流边界层中-(uv)的预算：条件和归一化与图7.34相同。
\end{sourceblock}
\begin{sourceblock}
\includegraphics[page=345,pagebox=mediabox,trim={64.39bp 553.93bp 64.39bp 1344.83bp},clip,width=0.94\textwidth,max height=0.38\textheight]{Pope.pdf}
\end{sourceblock}
其中压力-速率-ofis 列车张量 Rij 是
T (P) 是压力传递（参见练习 7.24）。考虑到连续性
kij 方程 V u 0，压力-应变率率张量收缩到零，因此从湍动能方程中消失。在边界层中，存在能量转移，速率为 -Rll R22 R33f，从 (u2) 到 (v2)

\begin{sourceblock}
\includegraphics[page=346,pagebox=mediabox,trim={64.39bp 1176.85bp 64.39bp 109.44bp},clip,width=\textwidth,max height=0.72\textheight]{Pope.pdf}
\par\vspace{0.45\baselineskip}
\small 图 7.39。 Reg = 1,410 处湍流边界层中的归一化耗散分量：来自 Spalart (1988) 的 DNS 数据，其中 6 =
\end{sourceblock}
剪应力平衡 图 7.38 显示了 - (uv) 的预算。由于 (uv) 为负，因此 -(uv) 中的“增益”对应于剪切应力大小的增加。可以看出，在边界层的大部分区域（y+ 40 到 y / 8 0.5），有
a ( u ) / a \textasciitilde{} ) 是产生量 ( e , -P12 (v2) 和压力项 -H12 之间的近似平衡。与法向应力平衡相反，除了靠近壁面之外，耗散相对较小。

在局部各向同性湍流中，耗散张量是各向同性的，即
2- 靠近墙壁，图 7.39 显示了由 -E 标准化的轮廓。

Eij 3 I' 的各向异性明显较大；但是，对于 y / b > 0.2 (y+ > 130)，存在近似各向同性 ( e。 /(: 的对角线分量接近
Eij E) 统一，非对角线分量接近于零)。小而独特
（对于 y / 8 > 0.1）可归因于模拟的雷诺数中各向异性水平相对较低：Saddoughi 和 Veeravalli (1994) 的边界层实验清楚地表明，在高雷诺数下存在局部各向同性。

从雷诺应力预算可以清楚地看出，速度-压力梯度相关性是一个重要的量，当然，速度-压力梯度相关性也是一个重要的量。

nij 生产 P bamnd 耗散 E\textasciitilde{}T,.h e 雷诺应力方程 和 是
根据雷诺分解 U - ( U ) u 和 def-
初始化 D/Dt a / a t U V 和 a / a t ( U ) V 显示

\begin{sourceblock}
\includegraphics[page=347,pagebox=mediabox,trim={64.39bp 1550.80bp 64.39bp 342.96bp},clip,width=0.94\textwidth,max height=0.38\textheight]{Pope.pdf}
\end{sourceblock}
因此，根据脉动速度 u(x, t) 的输运方程
(Eq. (5.138))，得到雷诺应力方程

\begin{sourceblock}
\includegraphics[page=347,pagebox=mediabox,trim={64.39bp 1349.83bp 64.39bp 569.93bp},clip,width=0.94\textwidth,max height=0.38\textheight]{Pope.pdf}
\end{sourceblock}
其中 P +a+n d II
由等式定义。 （7.179）和（7.180）。表明
方程中的粘性项。 (7.190) 可以重新表示为

\begin{sourceblock}
\includegraphics[page=347,pagebox=mediabox,trim={64.39bp 1145.86bp 64.39bp 794.90bp},clip,width=0.94\textwidth,max height=0.38\textheight]{Pope.pdf}
\end{sourceblock}
其中由方程式定义。 （7.181）。因此验证方程。 （7.178）。

证明 I I i j (Eq. (7.180)) 可以分解为

\begin{sourceblock}
\includegraphics[page=347,pagebox=mediabox,trim={64.39bp 904.89bp 64.39bp 1020.87bp},clip,width=0.94\textwidth,max height=0.38\textheight]{Pope.pdf}
\end{sourceblock}
是由方程定义的压力-应变率率张量。 (7.187),

\begin{sourceblock}
\includegraphics[page=347,pagebox=mediabox,trim={64.39bp 698.91bp 64.39bp 1150.85bp},clip,width=0.94\textwidth,max height=0.38\textheight]{Pope.pdf}
\end{sourceblock}
证明 Rij 是偏态的，并且 IIIij 和 Rij 在以下方面相等
均匀湍流。

由以上结果可知，雷诺应力方程可

\begin{sourceblock}
\includegraphics[page=347,pagebox=mediabox,trim={64.39bp 406.95bp 64.39bp 1443.82bp},clip,width=0.94\textwidth,max height=0.38\textheight]{Pope.pdf}
\end{sourceblock}
其中雷诺应力通量 Tkiji s

\begin{sourceblock}
\includegraphics[page=347,pagebox=mediabox,trim={64.39bp 261.97bp 64.39bp 1671.79bp},clip,width=0.94\textwidth,max height=0.38\textheight]{Pope.pdf}
\end{sourceblock}
\begin{sourceblock}
\includegraphics[page=347,pagebox=mediabox,trim={64.39bp 119.39bp 64.39bp 1750.78bp},clip,width=0.94\textwidth,max height=0.38\textheight]{Pope.pdf}
\end{sourceblock}
\subsection*{7.3.6 附加效果}\phantomsection\addcontentsline{toc}{subsection}{7.3.6 附加效果}
\begin{sourceblock}
\includegraphics[page=348,pagebox=mediabox,trim={64.39bp 1731.78bp 64.39bp 157.98bp},clip,width=0.94\textwidth,max height=0.38\textheight]{Pope.pdf}
\end{sourceblock}
因此表明雷诺应力方程的一半迹

（方程（7.194））与动能方程（方程（5.164））相同。

对于简单剪切流，i?(Ul)/dx2b 是唯一的非零均值
速度梯度，表明产生张量 P 为

\begin{sourceblock}
\includegraphics[page=348,pagebox=mediabox,trim={64.39bp 1389.82bp 64.39bp 465.94bp},clip,width=0.94\textwidth,max height=0.38\textheight]{Pope.pdf}
\end{sourceblock}
根据雷诺应力各向异性 Zk6.m 可知
生产张量是

\begin{sourceblock}
\includegraphics[page=348,pagebox=mediabox,trim={64.39bp 1172.85bp 64.39bp 723.91bp},clip,width=0.94\textwidth,max height=0.38\textheight]{Pope.pdf}
\end{sourceblock}
给出一个流程示例，其中正应力产生 Pll
是负数。

讨论压力-应变率率张量对雷诺兹的影响
剪应力

(a) 在 Rij 的主轴上，并且
(b) 在 (uiuj) 的主轴上。

根据 (Eq. 7 18 1 ) 的定义可知，在墙处，

\begin{sourceblock}
\includegraphics[page=348,pagebox=mediabox,trim={64.39bp 720.91bp 64.39bp 1177.85bp},clip,width=0.94\textwidth,max height=0.38\textheight]{Pope.pdf}
\end{sourceblock}
墙上哪些分量不为零？

7.3.6 附加效应 到目前为止，在本节中我们已经考虑了平板上的零压力梯度湍流边界层。现在简要描述非零压力梯度和表面曲率的影响。

平均压力梯度 提到平均压力梯度对边界层的影响，首先是因为在应用中它们经常出现，其次是因为影响可能很大。有利的压力梯度 (dpo/dx < 0)

例如，在机翼的前部。相反，逆压梯度 (dpoldx > 0) 对应于减速的自由蒸汽 (dUo/dx < 0)，并发生在扩压器或机翼的后部。

有利的平均压力梯度会导致平均速度剖面变陡，从而使形状因子 H 减小并且表面摩擦系数
增加。间歇区域的宽度增加（Fiedler 和 Head cf 1966），非湍流流体可以一直渗透到壁。事实上，在低雷诺数下，足够强的有利压力梯度可以导致边界层重新层化（Narasimha 和 Sreenivasan
正如所预料的，温和的逆压梯度会产生相反的效果：平均速度剖面变平，形状因子 H 增加，表面摩擦系数减小（Bradshaw 1967，Spalart 和
参见沃特莫夫 1993）。然而，强大的、长期的逆压梯度会导致边界层分离或脱离表面（Simpson 1989）。分离伴随着大规模的不稳定
( ( u ) 的流动，并且反向流动< 0) 分离的下游。

表面曲率 曲面上的边界层在许多应用中都很重要，例如涡轮机械中压缩机和涡轮叶片上的流动。在这些应用中，曲率处于主导流动方向。在机翼的上（吸力）表面上，曲率是凸的；而高弧度翼型的下（压力）表面的一部分有凹曲率。

在翼型上的流动中，边界层同时受到曲率和平均压力梯度的影响。这些影响可以在实验室实验中单独研究。 Muck、Hoffmann 和 Bradshaw (1985) 在恒定自由流速度的边界层上进行了实验，该边界层在具有平面上游部分的表面上发展，随后是凸曲率部分（具有恒定的曲率半径）。Hoffmann、Muck 和 Bradshaw (1985) 描述了使用凹面进行的类似实验。

在这些实验中，边界层厚度6与曲率半径的比率约为0.01。 Gillis 和 Johnston (1983) 报道了凸曲率高达 d / R , 0.1 的实验。

瑞利准则 (Rayleigh 191 6) 正确预测了曲率的稳定或不稳定效应。对于凸曲率，曲率中心位于表面下方，并且在边界层中角动量

\section*{7.4 湍流结构}\phantomsection\addcontentsline{toc}{section}{7.4 湍流结构}
根据瑞利准则，这种增加的角动量是稳定的，并且实验数据确实表明，与平面“边界层”相比，雷诺应力和皮肤摩擦系数有所降低。

在雷诺应力方程（以适当的曲线坐标系编写）中，存在相对大小的附加产生项（由于曲率）。然而，正如 Bradshaw (1973) 所指出的，曲率对湍流的影响比这种直接机制所能解释的要大一个数量级。

对于凹曲率表面上的流动，曲率中心位于表面上方的流体内。在边界层中，流体的角动量随半径减小（e，it 随着接近表面而减小）。根据瑞利准则，这是不稳定的。实验a1观察到形成了纵向泰勒-戈特勒涡，并且雷诺应力和皮肤摩擦系数增加（与平面边界层相比）。同样，影响的幅度比简单的缩放参数所暗示的要大得多。

对于缓和弯曲表面上的边界层，非零平均应变率为

\begin{sourceblock}
\includegraphics[page=350,pagebox=mediabox,trim={64.39bp 919.88bp 64.39bp 976.88bp},clip,width=0.94\textwidth,max height=0.38\textheight]{Pope.pdf}
\end{sourceblock}
（这是在局部笛卡尔坐标系中，其中基向量 e2 垂直于表面，el 和 e3 与表面相切，平均流动沿 el 方向。）可以看出，曲率产生额外应变率 -L ( U 1 )，与 (Ul) /ax2 相比，该额外应变率较小
通常发现湍流剪切流对这种额外的应变率表现出不成比例的大响应（布拉德肖
/,4 湍流结构 与自由剪切流的情况一样，自 1960 年以来，大部分关于壁限流的实验工作都针对湍流结构或准相干结构。这些结构通过流动可视化、条件采样技术或其他推导方法（如下所述）来识别；但它们很难准确定义。这个想法是，它们是空间和时间的区域（明显大于最小的流动或湍流尺度），其中流场具有
不同的位置和时间及其流场当然在细节上有所不同：但它们具有共同的特征连贯模式。

湍流结构实验研究的动机之一是希望

(i) 在明显的混乱中寻求秩序，
(ii)“解释”流动可视化中看到的模式，
(iii) 用基本原理“解释”流程中的重要“机制”

结构，以及
(iv) 识别“重要”结构，以便对其进行修改
为了实现减少阻力等工程目标

增强传热。

毫无疑问，这些研究取得了宝贵的成果，并实现了一些目标。然而，其他目标尚未实现，也不会实现。充满牛顿力学的头脑寻求对现象的简单确定性解释。只有在非常有限的意义上，相干结构才能简单地“解释”近壁湍流的行为。随机背景中有许多结构，它们之间的确定性和随机性相互作用远非清晰，而且并不简单。基于少量结构的动力学相互作用发展近壁湍流定量理论的目标尚未实现，而且很可能无法实现。

壁流结构的描述 Kline 和 Robinson (1990) 以及 Robinson (1991) 提供了槽道流和边界层中准相干结构的有用分类。确定的八个类别如下：

1. 0 < yf 1 10 区域内的低速条纹。

2. 低速流体从壁向外喷射。

3. 高速流体扫向墙壁。

4. 几种拟议形式的垂直结构。

5. 壁区有强的内部剪切层(y+ I8 0)。

6. 近壁袋，观察到在某些区域中没有明显液体的区域
流可视化的类型。

7. 背面：流向速度穿过的表面（比例为 6）

突然改变。

8. 外层的大规模运动（包括边界

\begin{sourceblock}
\includegraphics[page=352,pagebox=mediabox,trim={64.39bp 1549.81bp 64.39bp 109.44bp},clip,width=\textwidth,max height=0.72\textheight]{Pope.pdf}
\par\vspace{0.45\baselineskip}
\small 图 7.40。湍流边界层中的染料条纹显示低速近壁流体的喷射。 （来自 Kline 等人的实验（1967）。）
\end{sourceblock}
Kline 和 Robinson (1990)、Robinson (19 91)、Sreenivasan (1989)、Cantwell (198I)、Gad-el-Hak 和 Bandyopadhyay (19 94) 等人对这些结构的实验工作和讨论进行了回顾。这里我们仅提及一些主要发现。

许多流动可视化实验揭示了近壁区域的条纹，yf < 40。在一项关于水道边界层的实验中，Kline 等人。 (1967) 使用一根细线放置在流动上（沿 z 方向）作为电极来产生微小的氢气泡。将导线放置在壁和 yf 10 之间时，在壁平面中清晰可见沿流向 (x) 方向的长条纹，对应于氢气泡的积累。这些和随后的实验（例如 Kim、Kline 和 Reynolds (1971) 以及 Smith 和 Metzler (1983)）已经确定了这些结构的许多特征。在壁附近 (yf < 7)，条纹之间的间距随机分布在大约 806 和 1206, 之间，与雷诺数无关，并且它们的长度（在 x 方向）可以超过 1,0006,。条纹对应于相对缓慢移动的流体，其流向速度约为局部平均值的一半，而条纹之间的流体（不可避免地）移动相对较快。

条纹具有一种特征行为，称为破裂。随着下游距离的增加，条纹慢慢地远离墙壁；但随后，在某个点（通常在 yf lo 附近），它会转向并更快地远离墙壁，这一过程称为条纹提升或弹出。当它被提升时，条纹表现出快速振荡，然后分解为更精细的运动。图 7.40 显示了典型喷射中的染料条纹。

随着流体在喷射中远离壁，连续性要求在其他一些区域流向壁。 Corino 和 Brodkey (1969) 确定了高速流体 ( e . , u 0 ) 向壁移动的区域

\begin{sourceblock}
\includegraphics[page=353,pagebox=mediabox,trim={64.39bp 1355.83bp 64.39bp 109.44bp},clip,width=\textwidth,max height=0.72\textheight]{Pope.pdf}
\par\vspace{0.45\baselineskip}
\small 图 7.41。 u-v 样本空间显示四个象限的编号，以及对应于喷射和扫描的象限。
\end{sourceblock}
一个重要的问题是喷射和扫掠在湍流产生中的重要性。脉动速度的 u-v 样本空间被轴划分为四个象限，如图 7.41 所示。在象限 2 和象限 4 中，乘积为负，因此这些区域中的事件对应
uv 为正值（回想一下 P = -(uv) ( U )l ay）。因此，喷射（9u adrant 2）和扫掠（象限 4）都会产生湍流能量。测量各个象限的贡献（例如华莱士、
(uv) Eckelmann 和 Brodkey (1972) 以及 Willmarth 和 Lu (1972)) 被引用作为喷射和扫掠对生产重要性的证据。然而，应该认识到，象限 2 事件不一定是喷射；相关系数约为 -0.5 的简单事实表明，无论湍流结构如何，象限 2 和 -4 事件的可能性是象限 1 和 -3 事件的两倍（参见练习 7.3 1）。

在近壁区域 (y < loo)，成对的反向旋转流向涡流或如图 7.42 所示已被确定为主要涡流。

卷 - -“旋涡结构”（Bakewell 和 Lumley 1967，Blackwelder 和 Haritonidis 1983）。靠近壁，在辊之间，壁平面内存在流动收敛（a W l a z < 0），这解释了观察到的条纹。在图 7.42 的简化图中，离开辊之间壁的流体具有相对减小的轴向速度，这导致了图 7.42 中所示的速度分布。这些剖面包含拐点，因此无粘性不稳定，并被推测与破裂有关（Holmes、Lumley 和 Berkooz（1996））。

Head 和 Bandyopadhyay (1981) 提出，边界层中距壁较远的主要涡旋结构是或

\begin{sourceblock}
\includegraphics[page=354,pagebox=mediabox,trim={64.39bp 1296.84bp 64.39bp 109.44bp},clip,width=\textwidth,max height=0.72\textheight]{Pope.pdf}
\par\vspace{0.45\baselineskip}
\small 图 7.42。近壁区域反向旋转辊的草图。 （摘自 Holmes 等人 (1996)。）
\end{sourceblock}
\begin{sourceblock}
\includegraphics[page=354,pagebox=mediabox,trim={64.39bp 797.90bp 64.39bp 762.90bp},clip,width=\textwidth,max height=0.72\textheight]{Pope.pdf}
\par\vspace{0.45\baselineskip}
\small 图 7.43。 Head 和 Bandyopadhyay (1981) 提出的发夹涡。
\end{sourceblock}
为6，但总长度可以为6阶，因此它们在高雷诺数下被极度拉长。正如 Theodorsen (1952) 先前预测的那样，它们与墙倾斜约 45 英寸。

Head 和 Bandyopadhyay (1981) 提出，较大的结构可以由发夹涡流集合组成；事实上，Perry 和他的同事已经证明，这种基本涡流结构的适当分布可以解释边界层中的测量统计数据（Perry 和 Chong 1982，Perry、Henbest 和 Chong 1986，Perry 和 Marusit 1995）。

在边界层的最外层，流动是间歇性的。如第 5.5.2 节所述，前方有一个薄湍流
- 粘性超强-
将湍流边界层流体与无旋层自由流流体分离。粘性超层的大尺度特征

\begin{sourceblock}
\includegraphics[page=355,pagebox=mediabox,trim={64.39bp 1506.81bp 64.39bp 109.44bp},clip,width=\textwidth,max height=0.72\textheight]{Pope.pdf}
\par\vspace{0.45\baselineskip}
\small 图 7.44。 Reg \textasciitilde{} 4,000 处湍流边界层的大尺度特征。不规则的线-近似于粘性上层-是充满烟雾的湍流流体和清澈的自由流流体之间的边界。 （来自 Falco 1977 的实验。）
\end{sourceblock}
冒烟，而自由流动的液体是清澈的。作为一个例子，图 7.44 显示了 Falco (19 77) 实验中充满烟雾的区域的轮廓（解释为粘性超级层的位置）。

图 7.44 中识别的特征的性质和行为已在基于热线测量条件采样的大量实验中进行了量化（例如 Corrsin 和 Kistler (1954)、Kovasznay、Kibens 和 Blackwelder (19 70)、Blackwelder 和 Kovasznay (19 72) 以及 Murlis、Tsai 和 Bradshaw (1982)）。对于给定的超级层的 y 位置（非常近似）服从正态分布，平均值为 0.86，标准差为 0.156。相应地，间歇性因子 y 具有误差函数曲线（见图 7.26）。有一些非湍流流体的山谷深入边界层。这些山谷将大涡流或凸起分开，这些涡流或凸起以 2CL25" 的特征角倾斜。凸起的长度通常为 6 至 36（x 方向），宽度（z 方向）大约为 6 至 36。当它们以 0.9Uo 至 0.97Uo 的速度向下游对流时，它们会缓慢演化，并以与平均旋转相同的方式缓慢旋转（图 7.44 中的顺时针方向）。凸起的大尺度行为由外部尺度 Uo 和 6 来表征。

大涡流和上层还包含更精细尺度的结构：Falco (1974) 称这些为典型涡流，而如前所述，Head 和 Bandyopadhyay (1981) 认为它们是倾斜 45" 的发夹涡流。Murlis 等人 (1982) 的测量支持这些作者的论点，即精细尺度结构的长度为 6 阶，但宽度为 6 阶。

假设 u 和 v 是具有相关系数的联合正态分布

\begin{sourceblock}
\includegraphics[page=356,pagebox=mediabox,trim={64.39bp 1708.78bp 64.39bp 156.98bp},clip,width=0.94\textwidth,max height=0.38\textheight]{Pope.pdf}
\end{sourceblock}
对于联合法线 u 和 v，以及 -0.5（典型的边界
层），表明象限 1 和 -2 事件的概率为
（提示：考虑联合正态随机变量 u 和
au Dv，其中 a 和 的选择使得 6-5 协方差矩阵
是各向同性的。）

导出技术上述准相干结构主要通过低雷诺数湍流的可视化（来自实验或来自 DNS）来识别。这种主观的方法不可避免地会导致对结构的性质和意义的争议。因此，多年来，针对客观教育技术、识别结构并量化其重要性提出了许多建议。

基于条件单点测量的推导技术的例子有象限分析（Willmarth 和 Lu 1972）和可变间隔时间平均（VITA）（Blackwelder 和 Kaplan 1976）。 DNS 提供的对流场信息的访问允许基于压力的推导技术（Robinson 199I）和速度梯度张量的各种不变量（例如，Blackburn、Mansour 和 Cantwell（1996）以及 Jeong 等人。

本征正交分解 (POD) POD 是一种基于两点速度相关性的推导技术。它识别出平均而言包含最多能量的运动。 POD，也称为 Karhunen-Lohe 分解，由 Lumley (1967b) 将其引入湍流研究中，Berkooz、Holmes 和 Lumley (1993) 以及 Holmes 等人给出了完整的描述。 （1996）。

对于湍流，POD 提供脉动速度场 t) 的表示。然而，为了便于解释，这里描述的是区间 0 i x 5 L 中的随机标量函数 u(x)。

u(x) 的正交分解为

\begin{sourceblock}
\includegraphics[page=356,pagebox=mediabox,trim={64.39bp 119.39bp 64.39bp 1812.77bp},clip,width=0.94\textwidth,max height=0.38\textheight]{Pope.pdf}
\end{sourceblock}
其中 \{qn(x), n 1 2, .。 .\} 是一组实数基函数，满足

\begin{sourceblock}
\includegraphics[page=357,pagebox=mediabox,trim={64.39bp 1638.79bp 64.39bp 202.97bp},clip,width=0.94\textwidth,max height=0.38\textheight]{Pope.pdf}
\end{sourceblock}
和 \{a,, n 1, 2 .。是基函数系数。这些由下式给出

\begin{sourceblock}
\includegraphics[page=357,pagebox=mediabox,trim={64.39bp 1470.81bp 64.39bp 424.95bp},clip,width=0.94\textwidth,max height=0.38\textheight]{Pope.pdf}
\end{sourceblock}
通过乘以方程可以很容易地推导出来。 (7.203) 通过 qm(x) 和积分。基函数 qn(x)a 是非随机的，而系数 a 是随机的。基函数有无限多种选择，最熟悉的是傅里叶模式，（为了符合目前的符号）可以写成
为了点头，

\begin{sourceblock}
\includegraphics[page=357,pagebox=mediabox,trim={64.39bp 1105.86bp 64.39bp 798.90bp},clip,width=0.94\textwidth,max height=0.38\textheight]{Pope.pdf}
\end{sourceblock}
u(x) 在域上的平均能量定义为

\begin{sourceblock}
\includegraphics[page=357,pagebox=mediabox,trim={64.39bp 930.88bp 64.39bp 963.88bp},clip,width=0.94\textwidth,max height=0.38\textheight]{Pope.pdf}
\end{sourceblock}
代入方程式(7.203) 对于 u(x) 并调用方程。 (7.204)，我们得到

\begin{sourceblock}
\includegraphics[page=357,pagebox=mediabox,trim={64.39bp 641.92bp 64.39bp 1265.84bp},clip,width=0.94\textwidth,max height=0.38\textheight]{Pope.pdf}
\end{sourceblock}
因此 (a2)i 是第 n 个模式 anqn(x)a 对能量的贡献，以及部分和

\begin{sourceblock}
\includegraphics[page=357,pagebox=mediabox,trim={64.39bp 416.95bp 64.39bp 1465.82bp},clip,width=0.94\textwidth,max height=0.38\textheight]{Pope.pdf}
\end{sourceblock}
是前 N 个模式中包含的能量。这一发展适用于任何正交分解；也就是说，q,(x) 的任何选择都与方程一致。 （7.204）。

正交分解的一个定义属性是选择基函数来最大化部分和 EN（对于 N 1 2 , . ）：POD 的前 N 个模式包含的能量比
因此函数被确定为满足以下条件的函数

q \textasciitilde{} , ( \textasciitilde{} ) 正交条件（n 的方程 (7.204)），并且最大化

\begin{sourceblock}
\includegraphics[page=358,pagebox=mediabox,trim={64.39bp 1643.79bp 64.39bp 252.97bp},clip,width=0.94\textwidth,max height=0.38\textheight]{Pope.pdf}
\end{sourceblock}
显然，了解两点相关性就足以确定 ( x )；变分法的直接应用（参见 Holmes 等人 (1996)）表明 go,(x) 是积分的本征函数

\begin{sourceblock}
\includegraphics[page=358,pagebox=mediabox,trim={64.39bp 1362.83bp 64.39bp 491.94bp},clip,width=0.94\textwidth,max height=0.38\textheight]{Pope.pdf}
\end{sourceblock}
特征值A等于(a2)（见练习7.32）；因此特征值必须按降序排列 ( i2 A(2) 2 A(3) . . . ) 以便部分和最大。

POD 在湍流中的应用需要将这些想法扩展到三维速度场 u(x, t) - 参见 Holmes 等人。 （1996）。通常分解仅在空间中进行（而不是在时间和空间中），因此 POD 的形式为

\begin{sourceblock}
\includegraphics[page=358,pagebox=mediabox,trim={64.39bp 930.88bp 64.39bp 983.88bp},clip,width=0.94\textwidth,max height=0.38\textheight]{Pope.pdf}
\end{sourceblock}
基函数 cp,(x) 是不可压缩矢量场，在某种意义上是“特征涡流”（Moin 和 Moser 1989）。系数是时间的随机函数。

Bakewell 和 Lumle y (19 67) 从测量的两点相关性中推断出，完全发展的湍流管流中的第一个 POD 基函数对应于一对反向旋转的辊，类似于
伯科兹等人。 （1993）和霍姆斯等人。 （1996）。

通过代入方程。 （7.203），表明方程的左侧。 (7.211)

可以重新表示为

\begin{sourceblock}
\includegraphics[page=358,pagebox=mediabox,trim={64.39bp 119.39bp 64.39bp 1783.78bp},clip,width=0.94\textwidth,max height=0.38\textheight]{Pope.pdf}
\end{sourceblock}
乘以等式。 (7.211) 通过 vt(x) 积分得到

\begin{sourceblock}
\includegraphics[page=359,pagebox=mediabox,trim={64.39bp 1722.78bp 64.39bp 224.97bp},clip,width=0.94\textwidth,max height=0.38\textheight]{Pope.pdf}
\end{sourceblock}
显示 t：不同的 POD 系数是不相关的，并且
POD 特征值是

\begin{sourceblock}
\includegraphics[page=359,pagebox=mediabox,trim={64.39bp 1546.81bp 64.39bp 398.95bp},clip,width=0.94\textwidth,max height=0.38\textheight]{Pope.pdf}
\end{sourceblock}
证明两点相关性以 POD 形式给出
特征值和特征函数为

\begin{sourceblock}
\includegraphics[page=359,pagebox=mediabox,trim={64.39bp 1346.83bp 64.39bp 597.92bp},clip,width=0.94\textwidth,max height=0.38\textheight]{Pope.pdf}
\end{sourceblock}
动力学模型 从第 6.4.2 节中可以看出，波数空间中的纳维-斯托克斯方程由傅里叶系数的无限常微分方程组组成（方程（6.146））。这些方程是非线性的，涉及模式之间的三元相互作用。同理，对于正交分解方程： (7.2 12) 中，纳维-斯托克斯方程意味着一组无限的非线性常微分方程，用于描述系数 a,(t) 的演化。尝试用正交分解的 N 个系数和 N 个常微分方程的截断组来描述主要结构之间的相互作用是很自然的。

奥布里等人。 (1988) 在前五种 POD 模式的基础上开发了近壁区域的动力学模型。该模型在定性描述近壁行为的一些特征（包括爆裂现象）方面取得了一些成功。这种方法的进一步发展是 Holmes、Lumley 和 Berkooz (1996) 所著书籍的主题。

Waleffe (1997) 开发了一个稍微不同的模型，但同样基于正交模式。该模型展示了一个自我维持的过程，其中流向滚动和平均剪切之间的相互作用导致条纹的发展。然后条纹变得不稳定并分解以重新形成卷。

结论 实验和 DNS（主要是在低雷诺数下）揭示了几种准相干结构；在壁限流中（至少八个ac-
7 壁流
导出技术，特别是适当的正交分解，已被开发用于识别和量化结构的特征。获得的见解很有用，例如在减阻装置的开发中（参见 Bushnell 和 McGinley (1989)、Choi (1991)、Schoppa 和 Hussain (1998)）。然而，应该意识到，近壁湍流行为的完整定量图像还远未获得。特别是大多数结构的动力学以及它们之间的相互作用在很大程度上是未知的，并且可能很复杂。

\part{建模与模拟}
建模简介和

\chapter{建模与模拟概论}
\section*{8.1 挑战}\phantomsection\addcontentsline{toc}{section}{8.1 挑战}
与其他科学探究领域一样，在湍流研究中，最终目标是获得易于处理的定量理论或模型，可用于计算感兴趣的数量和实际相关性。一个世纪的经验表明，“湍流问题”是出了名的困难，并且没有简单的分析理论的前景。相反，希望利用数字计算机不断增强的能力来实现计算湍流相关属性的目标。在随后的章节中，描述和研究了五种主要的湍流计算方法。

从一开始就值得反思湍流的特殊性质，这些性质使得开发精确的、易于处理的理论或模型变得困难。速度场 U(x, t) 是三维的、与时间相关的且随机的。最大的湍流运动几乎与流动的特征宽度一样大，因此直接受到边界几何形状的影响（因此不具有普遍性）。时间尺度和长度尺度的范围很大。相对于最大尺度，柯尔莫哥洛夫时间尺度随着 Re 的减小而减小
柯尔莫哥洛夫长度尺度为 Re 在壁界流中，最具能量的运动（负责产生峰值湍流）与粘性长度尺度 6 成比例，与外部尺度 6 相比较小，并且大约减小（相对于 6）为 Re
纳维-斯托克斯方程中的非线性对流项会产生困难，而压力梯度项则更是如此。当用速度表示时（通过泊松方程的解，第 19 页的方程（2.49）），压力梯度既是非线性的又是非局部的。

\section*{8.2 方法概览}\phantomsection\addcontentsline{toc}{section}{8.2 方法概览}
后续章节中描述的方法采用偏微分方程组的形式，在某些情况下辅以代数方程。对于给定的流动，在指定适当的初始条件和边界条件的情况下，可以对这些方程进行数值求解。

在湍流模拟中，方程求解的是与时间相关的速度场，该速度场在某种程度上代表了湍流的一种实现的速度场 U(x)。相反，在湍流模型中，方程求解的是一些平均量，例如 (U)、(uju;) 和 ¢。（当不需要区分时，“模型”一词既指模拟又指湍流模型。）

所描述的两种模拟方法是直接数值模拟（DNS，第 9 章）和大涡模拟（LES，第 13 章）。在 DNS 中，求解纳维-斯托克斯方程以确定流的一种实现的 U(x,t)。由于必须解析所有长度尺度和时间尺度，因此 DNS 的计算成本很高；并且，由于计算成本随着 Re' 的增加而增加，因此该方法仅限于具有低到中等雷诺数的流。在 LES 中，求解“过滤”速度场 U(x,0) 的方程，该速度场代表更大规模的湍流运动。求解的方程包括未直接表示的较小尺度运动影响的模型。

第 10 章和第 11 章中描述的方法称为（雷诺平均纳维斯托克斯）(RANS)，因为它们涉及求解雷诺方程以确定平均速度场 (U)。在第一种方法中，雷诺应力是从湍流黏度模型获得的。湍流黏度可以从代数关系（例如在混合长度模型中）获得，也可以从湍流量（例如求解模拟输运方程的’k 和 «）获得。在雷诺应力模型（第 11 章）中，模型输运方程针对雷诺应力进行求解，从而消除了对湍流黏度的需要。

平均速度 (U) 和雷诺应力是速度 f(V ; x,t) 的欧拉 PDF 的一阶矩和二阶矩。在 PDF 方法（第 12 章）中，求解模型传输方程以获得 PDF，例如

\section*{8.3 模型评价准则}\phantomsection\addcontentsline{toc}{section}{8.3 模型评价准则}
本节的目的是概述中使用的标准

目前正在使用中。重要的是要认识到，存在着广泛的湍流，也存在着广泛的需要解决的问题。因此，拥有各种复杂性、准确性和其他属性各异的模型是有用且适当的。

可用于评估不同模型的主要标准是

(i) 描述水平，
(ii) 完整性，
(iii) 成本和易用性，
(iv) 适用范围，以及
作为阐述这些标准的例子，我们考虑两种模型：混合长度模型和 DNS，它们是方法范围的极端。

回想一下，DNS（直接数值模拟）在于求解纳维-斯托克斯方程以确定一种流动实现的瞬时速度场 U(x, t)。混合长度模型（应用于统计平稳的二维边界层流）由 (U(x,y )) 和 (V y)) 的边界层方程组成，雷诺剪切应力和湍流黏度从模型中获得

\begin{sourceblock}
\includegraphics[page=365,pagebox=mediabox,trim={64.39bp 809.90bp 64.39bp 1039.87bp},clip,width=0.94\textwidth,max height=0.38\textheight]{Pope.pdf}
\end{sourceblock}
\begin{sourceblock}
\includegraphics[page=365,pagebox=mediabox,trim={64.39bp 686.91bp 64.39bp 1232.84bp},clip,width=0.94\textwidth,max height=0.38\textheight]{Pope.pdf}
\end{sourceblock}
混合长度 tm(x,y ) 被指定为位置的函数。

描述级别 在 DNS 中，流是通过瞬时速度 U(x, t) 来描述的，由此可以确定所有其他信息。例如，可以执行流动可视化来检查湍流结构，并且可以提取多时间、多点的统计数据。另一方面，在混合长度模型中，描述是在平均流量级别：除了指定的混合长度之外，唯一直接表示的量是（U）和（V）。例如，没有提供有关速度 PDF、两点相关性或湍流结构的信息。平均流闭合（例如 thz 混合长度模型）提供的有限描述在许多方面是足够的

\begin{center}\small 表 8.1。不同计算级别的湍流示例\end{center}
\clearpage
\begin{sourceblock}
\includegraphics[page=366,pagebox=mediabox,trim={64.39bp 580.93bp 64.39bp 216.97bp},clip,width=\textwidth,max height=0.72\textheight]{Pope.pdf}
\end{sourceblock}
不完整的模型对于狭窄类别内的流动（例如，翼型上的附加边界层）很有用，对于这些流动，存在关于适当的流动相关规范的半经验知识。然而，一般来说，完整性显然是可取的。

成本和易用性 对于除最简单的流动之外的所有流动，都需要数值方法来求解模型方程。执行湍流模型计算的难度取决于流动和模型。

表 8.1 根据湍流的性质对湍流进行了分类
流动的特性：如果流动在统计上是平稳的，则它会减小，如果可以使用边界层方程，则它会进一步减小。在某些方法（例如 DNS）中，计算成本是流的雷诺数的快速增加函数；而在其他情况下（例如混合长度），雷诺数导致的成本增加微不足道或不存在。

对特定流执行湍流计算的任务可以分为两部分。首先，必须获得或开发求解模型方程的计算机程序，并针对当前的流动进行设置（例如，通过指定适当的边界条件）。其次，执行计算机程序进行计算，并提取所需的结果。第一步的成本和难度取决于可用的软件和算法，以及模型的复杂性。为特定类别的流程和模型开发计算机程序所需的工作量可能非常巨大，因此对评估和使用需要新程序的新模型来说是一个重大障碍。然而，这是“一次性成本”。

第二部分执行计算的成本和难度

取决于所需计算机的规模（例如工作站或超级计算机）、执行计算所需的人力时间和技能以及消耗的计算机资源。这些都是“经常性成本”。

就消耗的计算机时间而言，什么计算成本是可以接受的？根据具体情况，答案可能会相差一百万倍。彼得森等人。 (1989) 表明，最强大的超级计算机上大约 200 小时的 CPU 时间是“大规模研究”计算的合理上限。 （Kim 等人（1987）的渠道流 D N S 需要 250 小时。）很少进行这种规模的计算。对于“应用程序”彼得森等人。 (1989) 建议更合理的时间是最强大的超级计算机上 15 分钟的 CPU 时间，这相当于速度百分之一的工作站上 25 小时的 CPU 时间。为了在需要“重复”湍流计算的工作站上执行工程设计研究，每次计算的 CPU 时间最好为一分钟或更少。这些“重复”、“应用”和“大规模研究”计算的大小比率为 1: 1 .5 x lo3: 1 .2 x lo6。

在给定时间内可以执行的以浮点运算（浮点运算）为单位的计算量由计算机的速度决定，以兆浮点运算、千兆浮点运算或万亿浮点运算为单位，即 lo6、l o9 或
每秒失败次数。图8.1显示了最大的峰值速度
请注意，Jlops 是操作数，而 rnegamps 是速率，即操作数

\begin{sourceblock}
\includegraphics[page=368,pagebox=mediabox,trim={64.39bp 1262.84bp 64.39bp 109.44bp},clip,width=\textwidth,max height=0.72\textheight]{Pope.pdf}
\par\vspace{0.45\baselineskip}
\small 图 8.1。最快超级计算机与其推出年份的速度（每秒浮点数）。该线显示每十年增长率为 30 倍。 （经 Addison Wesley Longman 许可，改编自 Foster (1995)。）
\end{sourceblock}
30年的超级计算机。可以看出，“速度呈指数级增长，每十年增长 30 倍”。这是一个惊人的增长率，20 年内增长了 1000 倍，40 年内增长了 100 万倍。尽管没有可靠的基础可以推断几年后的情况，但通常认为这种趋势将持续下去（Foster 1995）。因此，今天的“研究”方法可能在20年后“应用”，在40年后“重复”计算。另一方面，“大规模研究”和“重复性”计算之间的 40 年跨度再次说明需要一系列不同计算要求的模型。

（需要谨慎看待图8.1中所示的绝对速度。实际实现的速度可能比峰值速度低一个甚至两个数量级。通常只使用并行计算机的一小部分（例如八分之一）处理器，并且每个处理器上仅实现峰值速度的2\%和50\%。还应该意识到，虽然这里讨论的重点是CPU时间，但内存也可能是一个限制因素。）

适用范围 并非所有模型都适用于所有流量。例如，有许多基于速度谱或两点相关性的模型，它们仅适用于均匀湍流。 （此类模型不考虑
作为第二个例子，特定的混合长度模型通常在混合长度的规范中对流动几何形状做出假设，因此它们的适用性仅限于该几何形状的流动。计算要求对某些模型的适用性提出了另一个尽管确实存在的限制。特别是，对于 DNS，计算要求随着雷诺数的增加而急剧上升，该方法适用于
因此仅适用于低或中等雷诺数的流动。第 9 章更详细地讨论了这一限制。

在本书中，注意力集中在恒定密度流中的速度场。然而，应当理解，在应用湍流模型的许多流中，存在附加现象，例如传热和传质、化学反应、浮力、压缩性和多相流。因此，一个重要的考虑因素是这里考虑的方法适用于或可以扩展到这些更复杂的流程的程度。

需要强调的是，在这些考虑中，我们将适用性与准确性分开。如果模型方程适定并且可以求解，则模型适用于流动，无论解是否准确。

准确性 不言而喻，准确性是任何模型的理想属性。在应用于特定流时，可以通过将模型计算与实验测量进行比较来确定模型的准确性。模具测试过程至关重要，值得仔细考虑。如图 8.2 所示，该过程由多个步骤组成，其中几个步骤会引入错误。

由于多种原因，计算中的边界条件不需要与测量流量的边界条件完全对应。流可以是近似二维的，但是在计算中可以假设完全是二维的。某些属性的边界条件不一定是已知的，因此必须进行估计；或者，它们可能取自包含一些测量误差的实验数据。

模型方程的数值解不可避免地存在数值误差。这可能有多种来源，但通常主要是空间截断误差。例如，在有限差分或有限元方法中，此误差按网格间距 Ax 的正幂进行缩放，而计算成本可能随着 A 可能作为正幂而增加。许多已发表的湍流模型计算包含重要的数值

\begin{sourceblock}
\includegraphics[page=370,pagebox=mediabox,trim={64.39bp 835.90bp 64.39bp 109.44bp},clip,width=\textwidth,max height=0.72\textheight]{Pope.pdf}
\par\vspace{0.45\baselineskip}
\small 图 8.2。该图表说明测量和计算的流量特性之间的差异 e 源于模型 €moue 的不准确；数字错误，€pum；测量误差 €meas;以及边界条件的差异，ep...（给出的 € 等式仅是暗示性的：误差不会线性增加。）
\end{sourceblock}
网格间距足够细，或者，坦率地说，是因为计算时没有足够仔细或考虑到数值精度。

总之，如图 8.2 所示，测量的流动特性与计算的流动特性之间的差异源于
(i) 模型的不准确性，
(ii) 数值误差，
(iii) 测量误差，以及
(iv) 边界条件的差异。

重要的结论是，只有当 (ii\}\{iv) 产生的误差相对较小时，测量的流动特性和计算的流动特性之间的比较才能确定模型的准确性。特别地，有一个
包含较大或未量化数值误差的计算。 Coleman 和 Stern (1997) 进一步讨论了这些问题。

最后的评论 特定模型对于特定湍流问题的适用性取决于上面讨论的标准的加权组合，并且各种标准的重要性的相对权重很大程度上取决于问题。因此，正如一开始所提到的，现在和未来，不存在一个“最佳”模型，而是存在一系列可以有效地应用于广泛的湍流问题的模型。

\chapter{直接数值模拟}
直接数值模拟 (DNS) 包括求解 Navier-S tokes 方程，解析所有运动尺度，并考虑适合流动的初始条件和边界条件。每个模拟都会产生流程的单个实现。直到 20 世纪 70 年代，有了足够强大的计算机，DNS 方法才变得可行。尽管它是建模方法中的后来者，但首先讨论 DNS 是合乎逻辑的。从概念上讲，它是最简单的方法，当它可以应用时，它的准确性和提供的描述水平是无与伦比的。然而，重要的是要认识到成本非常高；并且计算机要求随着雷诺数的增加而迅速增加，以至于该方法的适用性仅限于低或中等雷诺数的流动。

\section*{9.1 均匀湍流}\phantomsection\addcontentsline{toc}{section}{9.1 均匀湍流}
并详细检查计算要求。然后我们考虑非均匀湍流的 DNS，为此需要相当不同的数值方法。

9.1 均匀湍流 对于均匀湍流，方法（由 Orszag 首创）

伪谱法、Patterson (19 72) 和 Rogallo (19 8 1)) 是首选的数值方法，因为它们具有较高的精度。这些方法的基本原理在第 9.1.1 节中进行了描述，它允许估计 DNS 的计算成本（在第 9.1.2 节中）。

\subsection*{9.1.1 伪谱方法}\phantomsection\addcontentsline{toc}{subsection}{9.1.1 伪谱方法}
在均匀各向同性湍流的 DNS 中，解域是边长为 L: 的立方体，速度场 u(x, t) 表示为有限傅里叶

\begin{sourceblock}
\includegraphics[page=373,pagebox=mediabox,trim={64.39bp 1726.78bp 64.39bp 216.97bp},clip,width=0.94\textwidth,max height=0.38\textheight]{Pope.pdf}
\end{sourceblock}
参见第 6.4.1 节。总共表示 N3 个波数，其中偶数 N 决定模拟的大小，从而决定可以获得的雷诺数。通常，N 被选择为丰富的 2 次方（例如，N = 128 或 N = 192）。在震级上，最低的非零波数为 2n:/C，表示的 N 3 波数为

\begin{sourceblock}
\includegraphics[page=373,pagebox=mediabox,trim={64.39bp 1379.83bp 64.39bp 566.93bp},clip,width=0.94\textwidth,max height=0.38\textheight]{Pope.pdf}
\end{sourceblock}
对于 1 和 'N 之间的 ni 整数值。在每个方向上，表示的最大波数为

\begin{sourceblock}
\includegraphics[page=373,pagebox=mediabox,trim={64.39bp 1160.85bp 64.39bp 764.90bp},clip,width=0.94\textwidth,max height=0.38\textheight]{Pope.pdf}
\end{sourceblock}
这种谱表示相当于在均匀间距的 N3 网格上表示物理空间中的 u(x, t)

\begin{sourceblock}
\includegraphics[page=373,pagebox=mediabox,trim={64.39bp 945.88bp 64.39bp 980.88bp},clip,width=0.94\textwidth,max height=0.38\textheight]{Pope.pdf}
\end{sourceblock}
离散傅里叶变换（DFT，参见附录F）给出了N3个网格节点处的傅里叶系数t）和速度u（x，t）之间的一对一映射；快速傅立叶变换（FFT）可用于在波数空间（ i e ii(rc, t)）和物理空间（即+，u(x, t)）之间进行 N3l og N 次运算的变换。

谱方法涉及根据波数空间中的纳维-斯托克斯方程（方程（6.146））以小时间步长 A t 推进傅立叶模式 G(K, t)。对该方程中的三重相互作用求和需要 N6 级运算。为了避免这种巨大的成本，在伪谱方法中，纳维-斯托克斯方程中的非线性项的计算方式不同：速度场转换为物理空间；形成非线性项 (imem, uiuj)；然后将它们变换到波数空间。此过程需要按 N3 log N 次操作的顺序。还引入了必须消除或控制的混叠错误（参见附录 F）。

伪谱方法中的主要数值和计算问题是时间步进策略、混叠误差的控制以及分布式内存并行计算机上的实现。 Rogallo (198 I) 开发的算法对这些问题中的每一个都进行了优雅的处理，该算法形成了
除了各向同性湍流之外，DNS 还应用于具有各种施加的平均速度梯度的均匀湍流（参见第 158 页的练习 5.41）。这包括平均旋转的各向同性湍流（Bardina、Ferziger 和 Rogallo 1985）、均匀剪切流（Rogers 和 Moin 1987、Lee 等 1990）以及无旋转平均应变
等人。 （李和雷诺兹，1985）。对于这些情况，周期解域和波数向量会因平均变形而扭曲，如第 11.4 节所述。

计算成本

\subsection*{9.1.2 计算成本}\phantomsection\addcontentsline{toc}{subsection}{9.1.2 计算成本}
解决方案要求。盒子尺寸 LC 必须足够大才能代表含能运动；并且网格间距 Ax 必须足够小才能解决耗散尺度。此外，出于数值精度的考虑，用于推进求解的时间步长 At 受到限制。现在对这些因素进行更详细的研究。

或 r 是具有给定光谱的反向性湍流，L 的合理下限是八个积分长度尺度 (LC 8L11)，就最低波数而言，这意味着

\begin{sourceblock}
\includegraphics[page=374,pagebox=mediabox,trim={64.39bp 895.89bp 64.39bp 1020.87bp},clip,width=0.94\textwidth,max height=0.38\textheight]{Pope.pdf}
\end{sourceblock}
从图 6.18（第 241 页）和表 6.2（第 240 页）可以观察到，能谱的峰值为 t lcLll 1.3，并且 10\% 的能量位于 lcLll 1.0 以下的波数处：当 lcoLll 0.8 时，光谱低波数端的约 5\% 能量未解析。

速度场的傅里叶表示意味着必须在解域上施加周期性边界条件。图 9.1 显示了这些人为条件如何在自相关函数 f (r) 中体现。对于无限低波数分辨率 (LC/ Lll m)，自相关函数单调趋于零，而对于 L/LI1 8，它是周期性的，并且 ( r ( / L l l 3) 存在明显差异。有限框尺寸的影响尚未得到系统研究；但是，有代价的是，可以通过增加 L / L l l 来减少人为影响。

最小耗散运动（以柯尔莫哥洛夫标度 q 为特征）的分辨率需要足够小的网格间距 Axlq，或者相应地，足够大的最大波数。从图 6.16（第 237 页）可以观察到，耗散谱在超过 lcq - 1.5 时非常小，而且实际上经验表明，在

\begin{sourceblock}
\includegraphics[page=375,pagebox=mediabox,trim={64.39bp 888.89bp 64.39bp 1059.87bp},clip,width=0.94\textwidth,max height=0.38\textheight]{Pope.pdf}
\end{sourceblock}
\begin{sourceblock}
\includegraphics[page=375,pagebox=mediabox,trim={64.39bp 219.97bp 64.39bp 1677.79bp},clip,width=0.94\textwidth,max height=0.38\textheight]{Pope.pdf}
\end{sourceblock}
\begin{sourceblock}
\includegraphics[page=375,pagebox=mediabox,trim={64.39bp 119.39bp 64.39bp 1841.77bp},clip,width=0.94\textwidth,max height=0.38\textheight]{Pope.pdf}
\end{sourceblock}
\begin{sourceblock}
\includegraphics[page=376,pagebox=mediabox,trim={64.39bp 1224.85bp 64.39bp 109.44bp},clip,width=\textwidth,max height=0.72\textheight]{Pope.pdf}
\par\vspace{0.45\baselineskip}
\small 图 9.2。充分解析各向同性湍流所需的每个方向的傅里叶模式（或网格节点）数量 N。实线，等式。 （9.7）；虚线，渐近线，方程。 （9.8）。右轴显示所需的模式总数 Ne。
\end{sourceblock}
由（这与 N\textasciitilde{} 0.1 RA Reynolds (19 901.) 的估计值合理一致。）

或者为了使解在时间上前进得准确，流体粒子必须在时间步At中仅移动网格间距Ax的一小部分。在实践中，发现由此施加的库朗数大约为

\begin{sourceblock}
\includegraphics[page=376,pagebox=mediabox,trim={64.39bp 731.91bp 64.39bp 1168.85bp},clip,width=0.94\textwidth,max height=0.38\textheight]{Pope.pdf}
\end{sourceblock}
模拟的持续时间通常约为湍流时间尺度的四倍，k / \textasciitilde{} .H，因此所需的时间步数为

\begin{sourceblock}
\includegraphics[page=376,pagebox=mediabox,trim={64.39bp 526.93bp 64.39bp 1380.83bp},clip,width=0.94\textwidth,max height=0.38\textheight]{Pope.pdf}
\end{sourceblock}
（练习 9.1 给出了 At 的替代估计，因此也给出了 Me 的替代估计）

近似而言，执行模拟所需的浮点运算数量与模式数量和步数 M（模式步数）的乘积成正比。前面的结果产生

\begin{sourceblock}
\includegraphics[page=376,pagebox=mediabox,trim={64.39bp 258.97bp 64.39bp 1681.79bp},clip,width=0.94\textwidth,max height=0.38\textheight]{Pope.pdf}
\end{sourceblock}
显示雷诺数急剧上升。

为了完成这个图景，并有理由遵循，我们假设 1,000

\begin{center}\small 表 9.1。对不同雷诺数下各向同性湍流的 DNS 的估计，每个方向所需的模态 N（方程（9.7））；模式总数，N？；时间步数，M（方程（9.11））；模式步数，N?M；以及以 1 gigaflop 执行仿真的时间（假设每个模式每步执行 1,000 次操作）\end{center}
\begin{sourceblock}
\includegraphics[page=377,pagebox=mediabox,trim={64.39bp 876.89bp 64.39bp 392.95bp},clip,width=\textwidth,max height=0.72\textheight]{Pope.pdf}
\end{sourceblock}
从这些估计中得出的明显结论是，计算成本随着雷诺数（如 R！或 \textasciitilde{} e i th）急剧增加，对于千兆浮点计算机而言，远高于 RA 100 是不切实际的。图 9.3 中的纵坐标 (TG) 有另一种解释。::,i：t 是在一天内执行 DNS 所需的改进因子（在千兆浮点计算机上，以及在数值方法中）。因此，要实现 RIw -- 1,000，需要进行百万倍的改进。

检查计算量在不同尺度的湍流运动中的分布具有启发性。在三维波数空间中，所表示的模式位于边长为 2\textasciitilde{},,, 的立方体内。在分辨率良好的模拟中（例如 1.5），只有半径球内的模态
是动态显着的。如图 9.4 所示，耗散区 Kmax 对应于波数在 和 之间的球壳
KDI K,,,, 其中 0.1lq 是最大耗散运动的波数（参见
KDI 第 6.5.4 节）。在半径范围内存在着包含能量和
KDI（在足够高的雷诺数下）惯性范围运动。图9.4中半径的内球对应于能谱中的峰值

\begin{sourceblock}
\includegraphics[page=378,pagebox=mediabox,trim={64.39bp 1181.85bp 64.39bp 109.44bp},clip,width=\textwidth,max height=0.72\textheight]{Pope.pdf}
\par\vspace{0.45\baselineskip}
\small 图 9.3。在千兆浮点计算机上执行均匀各向同性湍流的 DNS 所需的时间（以天为单位）Tg 作为雷诺数的函数。实线，根据方程估计。 (9.7)、(9.11)和(9.13)；虚线，渐近线 (R;/70)°；符号，基于 40 节点 IBM SP2 的 DNS 计时。
\end{sourceblock}
较小。这是一个简单的算术问题（练习 9.2）来表明 99.98\% 的模式的波数大于
KDI；少于 0.02\% 的模态表示含能区或惯性子区内的运动。

由库朗数（方程（9.10））决定的时间步At为
目前数值方法所施加的限制

受雇。湍流对 At 的内在限制
是A t l z，应该很小。如果库朗数限制（方程
(9.10)) 被替换为

\begin{sourceblock}
\includegraphics[page=378,pagebox=mediabox,trim={64.39bp 431.95bp 64.39bp 1466.82bp},clip,width=0.94\textwidth,max height=0.38\textheight]{Pope.pdf}
\end{sourceblock}
获得修订后的估计

\begin{sourceblock}
\includegraphics[page=378,pagebox=mediabox,trim={64.39bp 307.96bp 64.39bp 1639.79bp},clip,width=0.94\textwidth,max height=0.38\textheight]{Pope.pdf}
\end{sourceblock}
\begin{sourceblock}
\includegraphics[page=378,pagebox=mediabox,trim={64.39bp 228.97bp 64.39bp 1714.78bp},clip,width=0.94\textwidth,max height=0.38\textheight]{Pope.pdf}
\end{sourceblock}
\begin{sourceblock}
\includegraphics[page=378,pagebox=mediabox,trim={64.39bp 119.39bp 64.39bp 1799.77bp},clip,width=0.94\textwidth,max height=0.38\textheight]{Pope.pdf}
\end{sourceblock}
\begin{sourceblock}
\includegraphics[page=379,pagebox=mediabox,trim={64.39bp 919.88bp 64.39bp 109.44bp},clip,width=\textwidth,max height=0.72\textheight]{Pope.pdf}
\par\vspace{0.45\baselineskip}
\small 图 9.4。各向同性湍流伪谱 DNS 的波数空间解域。所表示的模式位于边长立方体（虚线）内。显示的三个球体为： 半径为所有方向上解析的最大波数 (Kmax4 = 1.5)；半径为 Kp 的最大耗散运动的波数 (kp;7 = 0.1)；半径为 xg 的波数对应于
\end{sourceblock}
能谱 = 70 (KL \textasciitilde{} I I = 1.3)。仅代表 0.016\% 的模式
位于半径对应于包含能量的运动的球体内

范围和惯性子区。

（对于 Riv 1,000，与
由方程式给出的值 TG 8.5 x lo6 (9+13)。)

考虑均匀各向同性湍流的 D N S，其中
小尺度分辨率的特点是icrnaxq 1.5，如图所示
如图 9.4 所示。耗散区定义为（参见第 6.5.4 节）：

IKI > KDI，此处为 0.1。显示 lV3 的分数

表示不在耗散区内的波数模式

为 15 3 n/6 0.00016；相应地，99.98\% 的模式
均在耗散区内。估计在 Ri 70 时，峰值
能谱出现在 r c \textasciitilde{} O.Olrc,,,。其中的几分之几
模式是微不足道的，i.:.， IKI > Kma7x。

\subsection*{9.1.3 人为修改和不完全解决}\phantomsection\addcontentsline{toc}{subsection}{9.1.3 人为修改和不完全解决}
9.1.3 人工调制和不完全分辨率 D N S 的高成本激发了几种不同的方法，通过不完全解析低波数或高波数模式来获得更高的雷诺数（具有给定的模式数）。

低波数强迫 通过人为地强迫低波数模态，为其提供能量，可以在 D N S 中获得统计平稳的均匀湍流。因此，含能量的运动是不自然的，并且不受纳维-斯托克斯方程的控制。然而，只要小尺度运动是普遍的，就可以提取有关它们的有用信息。例如，在使用 2563 模式和强迫的模拟中，可以实现 RIu 180，其光谱超过十年的波数（Overholt 和 Pope 1996）：在没有强迫的情况下，可以获得的雷诺数为 RA = 60（图 9.2）。陈等人。 (1993) 报道了雷诺数 Ri 202 的强制 512' 模拟。

大涡模拟在n波数空间中，小于1x1的离散波数模式rc0n的数量随着lc3的增加而增加。因此，在完全解析的 DNS 中，绝大多数模式都处于耗散区内（见图 9.4 和练习 9.2）。这一观察结果为降低耗散区内的分辨率要求的方法提供了强有力的动力，特别是如果含能尺度是主要关注点的话。

在大涡模拟 (LES) 中，仅解析含能量的运动，并对未解析模式的影响进行建模。这种 LES 方法是第 13 章的主题。

超粘性另一种方法是用 vH(-1) v2mU 代替粘性项 v V2u，对于整数 m > 1，其中 是超粘性系数（参见，例如 Borue
VH 和 Orszag (1996))。在波数空间中，这对应于改变
+to V \textasciitilde{} K \textasciitilde{} " \textasciitilde{} w(Khic)h，导致耗散区更窄，V K U(K) 因此需要更少的模式来解析。

稀疏模式方法

( I K I 在这种方法中，含能量运动 < 被完全解析，
K E \textasciitilde{} ) 而在较高波数的壳中 (\textasciitilde{}"-'KEI - < 1x1 < 2"KEr, m 1,2，仅表示傅立叶模式的一小部分 2 (参见，例如 Meneguzzi

\section*{9.2 非均匀流动}\phantomsection\addcontentsline{toc}{section}{9.2 非均匀流动}
应该注意的是，上述所有方法都强制使用稀疏模式，这相当于与纳维-斯托克斯方程的重大背离。因此，它们不是湍流的直接数值模拟，而是每个模型的准确性事先未知。

9.2 非均匀流
与均匀湍流相比，将 DNS 应用于非均匀流的主要区别在于：

(i) 傅立叶表示不能用于不均匀的方向，
(ii) 物理边界条件（相对于周期性条件）是
必需的，并且
(iii) 近壁运动，以粘性长度尺度 6 为特征，施加
额外的分辨率要求。

现在针对特定流程说明其中一些问题。

\subsection*{9.2.1 槽道流量}\phantomsection\addcontentsline{toc}{subsection}{9.2.1 槽道流量}
域是一个尺寸为 La,x h x C, 的矩形框。流量在平均流量 (x) 和展向方向上在统计上是均匀的，允许分别使用 N 、 和 N z 模式进行傅立叶表示。相应地，在物理空间中，在x和z方向上存在均匀间距Ax C , / N 和Az C , / N 的网格节点。与均匀湍流的情况一样，傅里叶表示所隐含的周期性边界条件是人为的，但只要周期 C1 和 C z 与流动尺度相比足够大，它们的影响就很小。

在横流方向上，解域从底部 y = 0 延伸到顶壁 y = h。在两个边界处均适用无滑移条件 U = 0。

x 和 z 方向上的 N x 和 Nz 傅里叶模式的基本特征是它们提供一组正交基函数。在非周期 y 方向上，通过使用 N 切比雪夫多项式（在区间 1 1 1 1 上定义）可以实现相同的效果。因此
形式的

\begin{sourceblock}
\includegraphics[page=382,pagebox=mediabox,trim={64.39bp 1690.79bp 64.39bp 207.97bp},clip,width=0.94\textwidth,max height=0.38\textheight]{Pope.pdf}
\end{sourceblock}
快速变换（涉及 FFT）用于在 N, x N, x N 网格上的物理空间中的基函数系数 U(n,,n,,n,,t) 和速度 U(x, y,z,t) 之间进行变换。在 y 方向上，网格间距不均匀，方便的是，边界附近的网格间距更细。快速变换允许使用伪谱方法根据纳维-斯托克斯方程推进速度场。

从指定的初始条件开始，求解随时间推进，直到达到统计稳态。指定的初始条件可能会影响达到静止状态所需的步数 Mo，但不会影响静止状态下的统计数据。然后继续求解（另一个 M；时间步长），以便可以对统计数据进行时间平均。 （一点统计数据也在两个同质方向上进行平均。）表 9.2 显示了两个不同雷诺数下槽道流 DNS 中使用的数值参数。请注意，壁上 y 方向的网格间距大约为 34，其中 6 = v/u，在粘性长度尺度中。 （另请注意，较大雷诺数模拟在大小为四分之一的域上进行，持续时间为较低雷诺数模拟的四分之一。）根据基于 6 的分辨率要求，Reynolds (1990) 估计所需的模态总数随着 Re’ 的增加而增加（参见均匀湍流的 Re？”）。

\subsection*{9.2.2 自由剪切流}\phantomsection\addcontentsline{toc}{subsection}{9.2.2 自由剪切流}
随时间演变的混合层和平面尾流是使用 DNS 研究的统计一维自由剪切流的示例。对于混合层，Rogers 和 Moser (1994) 使用了伪谱方法，在 x 和 z 方向上使用傅里叶模式，在 y 方向上使用雅可比多项式。由于周期性的强加，与随时间增加的混合层厚度相比，域£ 1 和£ 的尺寸必须很大。雅可比多项式扩展到 y = +0，便于指定自由流中的边界条件，即平行于 x 轴的均匀流。

对于随时间变化的流动，指定适当的初始条件至关重要。在混合层模拟期间，5(t) 增加了三倍，找到了初始条件

\begin{center}\small 表 9.2。 Re, = 180（Kim 等人，1987 年）和 Re, = 595（N. N. Mansour 博士，个人通讯）时槽道流 DNS 的数值参数（文本中未定义的变量是模态总数 Nyy；y 中的最大和最小网格间距、Aymax 和 AYmin；以及模拟持续时间 T）\end{center}
\begin{sourceblock}
\includegraphics[page=383,pagebox=mediabox,trim={64.39bp 787.90bp 64.39bp 384.95bp},clip,width=\textwidth,max height=0.72\textheight]{Pope.pdf}
\end{sourceblock}
空间发展的混合层在统计上是静止的和二维的。在DNS中，流入边界条件（例如0处）和流出条件（at处）是必需的。粗略地说，这里的流入条件与时间混合层中的初始条件起着相同的作用。毫不奇怪，我们发现流对这些流入条件的细节很敏感（Buell and Mansour 1989）。

\subsection*{9.2.3 流过向后台阶}\phantomsection\addcontentsline{toc}{subsection}{9.2.3 流过向后台阶}
该流程的 DNS 概述如图 9.5 所示，由 Le、Moin 和 Kim (1997) 执行。湍流边界层（厚度为 6·1.2h，自由流速度为 Uo）从左侧边界进入，在台阶处分离
0），然后重新连接下游（7h）。该流动是统计上的（ x = x 静止且二维。考虑为 Re Uoh/v 5,100 的雷诺数相当低，例如与 Re 500,000 相比

\begin{sourceblock}
\includegraphics[page=384,pagebox=mediabox,trim={64.39bp 1547.81bp 64.39bp 109.44bp},clip,width=\textwidth,max height=0.72\textheight]{Pope.pdf}
\par\vspace{0.45\baselineskip}
\small 图 9.5。 Le 等人使用的解决方案域的草图。 (1997) 用于向后步骤的 DNS 流。尺寸以台阶高度 h 为单位。
\end{sourceblock}
\begin{center}\small 表 9.3。向后步骤中流的 DNS 数值参数（Le et al. 1997）\end{center}
\begin{sourceblock}
\includegraphics[page=384,pagebox=mediabox,trim={64.39bp 936.88bp 64.39bp 654.92bp},clip,width=\textwidth,max height=0.72\textheight]{Pope.pdf}
\end{sourceblock}
表 9.3 给出了 DNS 的一些数字特征。尽管雷诺数较低且消耗大量计算机时间（约 54 天），但分辨率仍然很低。根据作者的说法，最好将入口长度（台阶上游）加倍，并将翼展方向的节点数量 N 加倍。此外，周期性展向方向上相对较短的长度 Lz 4h 导致某些位置的两点速度自相关对于所有展向间隔超过 0.7（而不是可以忽略不计）

L）。模拟结果与 Jovic 和 Driver (1994) 的实验非常一致，并且包括雷诺应力预算。

\section*{9.3 讨论}\phantomsection\addcontentsline{toc}{section}{9.3 讨论}
在可以应用的地方，DNS 提供的描述和准确性水平是其他方法无法比拟的。人们已经开发出了高精度的数值方法来求解纳维-斯托克斯方程，并且 DNS 的实践者通常对数值精度有很高的标准。 w
湍流和湍流实验的边缘。例如，DNS已被用来提取拉格朗日统计数据（例如，Yeung和Pope（1989））和压力波动的统计数据（例如，压力应变张量，Spalart（1988）），这些统计数据几乎不可能通过实验获得。近壁流的细节，以及受到各种变形的均匀湍流，在 DNS 中也比在实验中更容易研究+。一些 DNS 结果在第 5 章和第 7 章中介绍；但读者可以参考 Moin 和 Mahesh (1998)，以更全面地回顾 DNS 研究对我们理解湍流和湍流所做的贡献。

应当理解，并非所有基于纳维-斯托克斯方程的模拟都是“直接的”。由于人工修改、不完全分辨率或非物理边界条件，模拟不需要直接对应于可实现的湍流。然而，可以明智地使用非物理模拟来分离和研究特定现象（参见 Jimknez 和 Moin (1991) 以及 Perot 和 Moin (1995)）。

DNS 的缺点当然是它的计算成本非常大，而且这个成本随着雷诺数（大约为 Re3）而迅速增加。在超级计算机上，计算机时间通常为 200 小时，因此只能模拟低或中等雷诺数的流动。

对 DNS 计算成本的观察不仅仅表明了当前计算机的局限性。它们还表明 DNS 与确定湍流中的平均速度和含能量运动的目标之间存在不匹配。在 DNS 中，超过 99\% 的工作都致力于耗散区（参见图 9.4 和练习 9.2），并且这种工作随着雷诺数的增加而大幅增加。相比之下，平均流量 aI d。 d 含能量运动的统计数据仅显示出较弱的 \textasciitilde{} e \textasciitilde{} nold s - n um bdeerp 端点。

\chapter{湍流黏度模型}
在本章和下一章中，我们考虑 RANS 模型，其中求解平均速度场的雷诺方程。在雷诺方程中显示为未知数的雷诺应力由湍流模型确定，可以通过湍流黏度假设或更直接地根据建模的雷诺应力传递方程（第 11 章）。

湍流黏度模型基于湍流黏度假设，该假设在第 4 章中介绍并在后续章节中使用。根据假设，雷诺应力由下式给出

\begin{sourceblock}
\includegraphics[page=386,pagebox=mediabox,trim={64.39bp 940.88bp 64.39bp 950.88bp},clip,width=0.94\textwidth,max height=0.38\textheight]{Pope.pdf}
\end{sourceblock}
或者，在简单剪切流中，剪切应力由下式给出

\begin{sourceblock}
\includegraphics[page=386,pagebox=mediabox,trim={64.39bp 807.90bp 64.39bp 1137.86bp},clip,width=0.94\textwidth,max height=0.38\textheight]{Pope.pdf}
\end{sourceblock}
给定湍流黏度场 v \textasciitilde{} ( xt),, 方程： (10.1) 为方程提供了最方便的闭合，其形式与纳维-斯托克斯方程相同（第 93 页的方程（4.46））。因此，不幸的是，对于许多流量来说，假设的准确性很差。 10.1 节回顾了上面提到的湍流黏度假说的许多缺陷。

如果湍流黏度假设被认为是充分的近似，那么剩下的就是确定湍流黏度 v \textasciitilde{} ( xt), 的适当规格。这可以写成速度的乘积

\begin{sourceblock}
\includegraphics[page=386,pagebox=mediabox,trim={64.39bp 285.96bp 64.39bp 1607.80bp},clip,width=0.94\textwidth,max height=0.38\textheight]{Pope.pdf}
\end{sourceblock}
指定 vr 的任务通常通过 u* 和 l* 的规范来完成。在代数模型（第 10.2 节）中，混合长度模型，

\section*{10.1 湍流黏度假设}\phantomsection\addcontentsline{toc}{section}{10.1 湍流黏度假设}
例如，根据流动的几何形状来指定。在二元方程模型（第 10.4 节）中，k-Em 模型是最好的例子 - u' 和 l' 与 k 相关，并且求解了建模的输运方程。

10. I 湍流黏度假说 湍流黏度假说可以分为两部分。首先，有

\subsection*{10.1.1 内在假设}\phantomsection\addcontentsline{toc}{subsection}{10.1.1 内在假设}
各向异性 (uiuj) 2kgm
由平均速度梯度决定

aij - 3 a(Ui)/axj。其次，有一个具体的假设，即 a ( U i ) / a x j 之间的关系为

\begin{sourceblock}
\includegraphics[page=387,pagebox=mediabox,trim={64.39bp 1148.86bp 64.39bp 685.91bp},clip,width=0.94\textwidth,max height=0.38\textheight]{Pope.pdf}
\end{sourceblock}
Sij 其中是平均应变率张量。当然，这直接类似于牛顿流体中粘性应力的关系：

\begin{sourceblock}
\includegraphics[page=387,pagebox=mediabox,trim={64.39bp 915.88bp 64.39bp 1028.87bp},clip,width=0.94\textwidth,max height=0.38\textheight]{Pope.pdf}
\end{sourceblock}
内在假设 为了讨论内在假设，我们首先描述一个简单的流程，但它是完全错误的。然后表明，在一个关键方面，湍流的物理学与导致粘性应力定律（方程（10.6））的分子过程的物理学有很大不同。然而，最后观察到，对于简单的剪切流，湍流黏度假设仍然是相当合理的。

轴对称收缩 图 10.1 是风洞实验的草图，最初由 Uberoi (1956) 进行，旨在研究轴对称收缩对湍流的影响。空气流经湍流产生网格进入第一个直线部分，其中平均速度 (U1)（理想情况下）是均匀的。在这一部分有
（Sij 没有平均应变 O），并且湍流（几乎是各向同性的）开始衰减。

在第一个直线部分之后有一个轴对称收缩

\begin{sourceblock}
\includegraphics[page=388,pagebox=mediabox,trim={64.39bp 1415.82bp 64.39bp 109.44bp},clip,width=\textwidth,max height=0.72\textheight]{Pope.pdf}
\par\vspace{0.45\baselineskip}
\small 图 10.1。与 Uberoi (1956) 和 Tucker (1970) 使用的设备类似的设备草图，用于研究轴对称平均应变对网格湍流的影响。
\end{sourceblock}
= a ( U , ) / a x l = S,，因此横向压缩应变率均匀，
S,,-S。量 S L k / \&（在收缩开始时评估）测量相对于湍流时间尺度的平均应变率。图 10.2 显示了标准化各向异性的测量结果
1 来自 Tucker (1970) 的实验 (bij ( u \textasciitilde{} u \textasciitilde{} ) / ( u \textasciitilde{}-u \textasciitilde{} ) lav/k) 与 S n k / e 2.1。图 10.2 中还显示了 Lee 和 Reynolds (1985) 获得的 S\&/E 55.7 的 DNS 结果。对于 Sik/\& 的这个大值，快速畸变理论（RDT，参见第 11.4 节）准确地描述了雷诺应力的演变。根据 RDT，雷诺应力不是由应变率决定，而是由湍流所经历的平均应变总量决定。在这些情况下，湍流的行为不像粘性流体，而更像弹性固体（Crow 1968）：湍流黏度假设在性质上是不正确的。

在图 10.1 所示的实验中，收缩之后出现第二个直线部分。由于该部分不存在平均应变，因此湍流黏度假说不可避免地预测雷诺应力各向异性为零。然而，Warhaft (19 80) 的实验数据却表明，在湍流时间尺度 k / \& 上，收缩中产生的各向异性衰减得相当缓慢（见图 10.2）。这些持续的各向异性的存在并不是因为局部平均应变率（为零），而是因为湍流所承受的先前的应变历史。

显然，对于该流动，无论是在收缩段还是在下游直段，湍流黏度假说的内在假设都是不正确的：雷诺应力各向异性 str 未确定

\begin{sourceblock}
\includegraphics[page=389,pagebox=mediabox,trim={64.39bp 1262.84bp 64.39bp 109.44bp},clip,width=\textwidth,max height=0.72\textheight]{Pope.pdf}
\par\vspace{0.45\baselineskip}
\small 图 10.2。轴对称应变期间和之后的雷诺应力各向异性。收缩：Tucker(1970)的实验数据，S,k/e = 2.1； Lee 和 Reynolds (1985) 的 DNS 数据，S,k/e = 55.7；从收缩开始的飞行时间 t 通过平均应变率归一化 直线部分：Warhaft 的实验数据（1980）；从直线段开始的飞行时间被归一化
\end{sourceblock}
根据那里的湍流时间尺度。

与动力学理论的比较 理想气体的简单动力学理论（参见 Vincenti 和 Kruger (1965) 以及 Chapman 和 Cowling (19 70)）得出牛顿粘性应力定律（方程 ( 1 0 ) ），其运动粘度为

\begin{sourceblock}
\includegraphics[page=389,pagebox=mediabox,trim={64.39bp 769.90bp 64.39bp 1178.85bp},clip,width=0.94\textwidth,max height=0.38\textheight]{Pope.pdf}
\end{sourceblock}
其中 是平均分子速度，A 是平均自由程。通过与动力学理论的类比来证明湍流黏度假说是很自然的，从而赋予 u* 和 l*b 与 A 的类比物理意义。然而，对所涉及的各种时间尺度的简单检查表明，这种类比不具有普遍有效性。

在简单层流剪切流（剪切速率为 8Ul/8xz -- S UIL）中，
AlC S-' 分子时间尺度和剪切时间尺度的比率为

\begin{sourceblock}
\includegraphics[page=389,pagebox=mediabox,trim={64.39bp 333.96bp 64.39bp 1572.80bp},clip,width=0.94\textwidth,max height=0.38\textheight]{Pope.pdf}
\end{sourceblock}
这通常非常小（例如 10，参见练习 10.1）。意义
-分子时间尺度相对较小是因为分子运动的统计状态快速调整以适应施加的应变。相比之下，对于湍流剪切流，湍流时间尺度 k / \& 与
约 3（第 131 页表 5.2）；在均匀湍流剪切流实验中，通常为 6（第 157 页的表 5.4）；在受到快速扭曲的湍流中，它可能会大几个数量级。因此，正如已经观察到的，湍流不会快速调整以适应施加的平均应变，因此（与分子运动的情况相反）应力和应变率之间的局部关系没有一般基础。

简单剪切流上面给出的快速轴对称变形的例子和时间尺度的考虑表明，一般来说，湍流黏度假设是不正确的。尽管存在这些普遍的反对意见，但对于一些重要的特定流动，该假设更为合理。在简单的湍流剪切流（例如，圆射流、混合层、槽道流和边界层）中，湍流特性和平均速度梯度变化相对较慢（遵循平均流）。因此，局部平均速度梯度表征了湍流所遭受的平均畸变的历史；雷诺应力平衡由局部过程产生、耗散、压力-应变率率张量主导，非局部传输过程相比较小（例如，参见第 316-3 17 页上的图 7.35-7.38）。那么，在这些情况下，假设雷诺应力和局部平均速度梯度之间存在关系是更合理的。

一个重要的观察是，在这些特定的流动中（其中湍流特性跟随平均流动缓慢变化），湍动能的产生和耗散近似平衡，即 P I E s s 1。相比之下，在轴对称收缩实验（图 1.1）中，收缩部分的 P I Ei 远大于 1，而在下游的直线部分 P I E 为零：在这两种情况下，湍流黏度假说是不正确的。

梯度扩散假说 与湍流黏度假说相关的是梯度扩散假说

\begin{sourceblock}
\includegraphics[page=390,pagebox=mediabox,trim={64.39bp 327.96bp 64.39bp 1604.80bp},clip,width=0.94\textwidth,max height=0.38\textheight]{Pope.pdf}
\end{sourceblock}
据此，标量通量与平均标量梯度一致（参见第 4.4 节）。上面的大多数观察结果同样适用于梯度扩散假设。在均匀剪切流中，发现标量通量的方向与均值通量的方向显着不同。

湍流热流（在通常的坐标系中）标量方程

\begin{sourceblock}
\includegraphics[page=391,pagebox=mediabox,trim={64.39bp 1711.78bp 64.39bp 234.97bp},clip,width=0.94\textwidth,max height=0.38\textheight]{Pope.pdf}
\end{sourceblock}
rT 可用于定义，因此不涉及任何假设（对于该组件）。湍流普朗特数可用于关联 vT 和 rT、rnT
- vT/uTman d 对于简单剪切流，具有阶单位（参见
即 - 、OT，例如第 162 页上的图 5.34）。

v 和 都可以写成速度尺度和长度尺度的乘积（方程（10.3））。它们也可以表示为速度尺度和时间尺度的平方的乘积：

\begin{sourceblock}
\includegraphics[page=391,pagebox=mediabox,trim={64.39bp 1314.83bp 64.39bp 624.92bp},clip,width=0.94\textwidth,max height=0.38\textheight]{Pope.pdf}
\end{sourceblock}
rT 如12.4节所示，在理想情况下，可以与湍流的统计相关：u'是均方根。速度 u'，T' 是拉格朗日积分时标 TL（参见第 500 页的方程（12.15 8））。

根据简单的动力学理论（参见 Vincenti 和 Kruger
(1965)) 理想气体的运动粘度为

\begin{sourceblock}
\includegraphics[page=391,pagebox=mediabox,trim={64.39bp 924.88bp 64.39bp 1005.87bp},clip,width=0.94\textwidth,max height=0.38\textheight]{Pope.pdf}
\end{sourceblock}
平均分子速度是声速的1.35倍

a.显示由分子归一化的剪切速率 S U I L

\begin{sourceblock}
\includegraphics[page=391,pagebox=mediabox,trim={64.39bp 707.91bp 64.39bp 1164.85bp},clip,width=0.94\textwidth,max height=0.38\textheight]{Pope.pdf}
\end{sourceblock}
其中马赫数和克努森数由 M a 定义
U/a 和 K n AIL。

使用关系式 a2 yplp（y 为 1.4）来表明
粘性剪切应力 q2t o 法向应力（压力）为

\begin{sourceblock}
\includegraphics[page=391,pagebox=mediabox,trim={64.39bp 432.94bp 64.39bp 1477.81bp},clip,width=0.94\textwidth,max height=0.38\textheight]{Pope.pdf}
\end{sourceblock}
使用值 a 332 m s-I 和 v 1.33 x m2 s-I
（对应于大气条件下的空气）且S = 1 s-'，
获得以下估计值：

\begin{sourceblock}
\includegraphics[page=391,pagebox=mediabox,trim={64.39bp 135.98bp 64.39bp 1717.78bp},clip,width=0.94\textwidth,max height=0.38\textheight]{Pope.pdf}
\end{sourceblock}
10.1.2 spec\#c 假设

\subsection*{10.1.2 具体假设}\phantomsection\addcontentsline{toc}{subsection}{10.1.2 具体假设}
雷诺应力和平均速度梯度由 \textasciitilde{} q(1.0.1)（或 e：等价地，方程（10.4）或（10.5））给出。

对于简单剪切流，单个感兴趣的雷诺应力 (uv) 与单个显着平均速度梯度 a ( U ) / a y 通过方程 1 相关。 （10.2）。本质上，不涉及任何假设，而是方程定义了 v \textasciitilde{}。对于圆射流，图 5.10（第 108 页）和图 7.30（第 307 页）给出了边界层的 v\textasciitilde{} 剖面的示例。

一般来说，湍流黏度假说中的具体假设是雷诺应力各向异性张量 aij 与均值线性相关
Sij 应变率张量通过标量湍流黏度，方程： （10.5）。即使对于最简单的流程，这显然也是不正确的。在湍流剪切流中
（Sll Sz2 S33 法向应变率为零 0），但法向应变率为零
-- = = 雷诺应力彼此显着不同（参见第 157 页的表 5.4）。对同一观察结果的另一种观点是， 的主轴线（很大程度上）与 的主轴线不一致
aij Sij（参见第 90 页的练习 4.5）。

简单线性应力定律适用于粘性应力（方程（10.6））但不适用于雷诺应力的原因可以再次根据时间尺度比和各向异性水平来理解。与分子尺度相比，应变非常弱 ( S I I C I)，因此它与各向同性的偏差非常小：在简单层流剪切流中，各向异性和各向同性分子应力之比为

\begin{sourceblock}
\includegraphics[page=392,pagebox=mediabox,trim={64.39bp 677.91bp 64.39bp 1227.85bp},clip,width=0.94\textwidth,max height=0.38\textheight]{Pope.pdf}
\end{sourceblock}
（参见练习 1 0 1 ） isi 通常非常小。因此，有充分的理由预期各向异性应力与速度梯度线性相关。牛顿粘性应力定律（方程（10.6））是与应力张量的数学性质一致的最一般可能的线性关系。相比之下，在湍流剪切流中，各向异性-各向同性应力比 -(uv)/(?k) 接近 0.5。那么，在湍流情况下，应变率相对较大（S k l \textasciitilde{}> 1），并且导致相对较大的各向异性。因此，没有理由假设该关系是线性的。

有几类流动，其中平均速度梯度张量比简单剪切流中的更复杂，并且已知湍流黏度假设显着失败。例子是

\section*{10.2 代数模型}\phantomsection\addcontentsline{toc}{section}{10.2 代数模型}
巨大的流线曲率（Bradshaw 1973，Patel 和 Sotiropoulos 1997），以及非圆形横截面管道中充分发展的流动（Melling 和 Whitelaw 1976，Bradshaw 1987）。

代替方程式。 (10.5)，一个可能的非线性湍流黏度假设

\begin{sourceblock}
\includegraphics[page=393,pagebox=mediabox,trim={64.39bp 1553.80bp 64.39bp 386.95bp},clip,width=0.94\textwidth,max height=0.38\textheight]{Pope.pdf}
\end{sourceblock}
其中系数 和 可能取决于平均速度-
V T \textasciitilde{} 、VT2 VT3 梯度不变量，例如 S（以及湍流量）。请注意，对平均旋转率的依赖性（例如，通过
vT2)i 中的 Qij 项要求 的主轴不与
所以那些 Sij 的人。已经开发出获得非线性湍流黏度定律的合理方法，并在第 11 .9 节中进行了描述。

总之：由 d(Ui)/dxj 局部确定的湍流黏度假说的内在假设不具有普遍有效性。

aij 然而，对于简单的剪切流，其中平均速度梯度和湍流特性演化缓慢（跟随平均流），该假设更为合理。在这种流动中，P/Ei 接近统一，这表明雷诺应力方程中平均剪切产生的产物与其他局部过程的重新分布和耗散之间存在近似局部平衡。

1 0.2 代数模型 前面章节已经介绍过的代数模型有

\subsection*{10.2.1 均匀湍流黏度}\phantomsection\addcontentsline{toc}{subsection}{10.2.1 均匀湍流黏度}
现在根据第 8 章中描述的标准进行评估。

10.2.1 均匀湍流黏度 在平面二维自由剪切流的应用中，均匀湍流黏度模型可以写成

\begin{sourceblock}
\includegraphics[page=393,pagebox=mediabox,trim={64.39bp 365.95bp 64.39bp 1557.80bp},clip,width=0.94\textwidth,max height=0.38\textheight]{Pope.pdf}
\end{sourceblock}
其中 Uo(x) 和 6(x) 是平均流量的特征速度尺度和长度尺度，解释湍流雷诺数的 RT 是与流量相关的常数。因此，湍流黏度在整个流动（y 方向）上被视为恒定，但它在平均流动方向上变化。

\begin{center}\small 表 10.1。测得的撒播率 S 和相应的值\end{center}
\begin{sourceblock}
\includegraphics[page=394,pagebox=mediabox,trim={64.39bp 1227.85bp 64.39bp 216.97bp},clip,width=\textwidth,max height=0.72\textheight]{Pope.pdf}
\end{sourceblock}
对于适用的简单自由剪切流，该模型是不完整的，因为必须指定 RT。适当的值取决于流量的性质以及为 6(x) 和 Uo(x) 选择的定义。第五章表明，对于每个自相似自由剪切流，扩散速率S与湍流雷诺数RT之间存在反比关系。表 10.1 总结了测得的铺展率和相应的 RT 值。 （对于每个流，S、6 和 Uo 的定义在第 5 章中给出。）

在自相似的自由剪切流中，凭经验确定的湍流黏度在整个流动中相当均匀，但当接近自由流时，它会减小到零（例如，参见第 108 页的图 5.10）。相应地，除流动边缘外，均匀粘度模型预测的平均速度剖面与实验数据非常吻合（例如，第 119 页的图 5.15）。

原则上，均匀湍流黏度模型可以应用于简单的壁面流动。然而，由于湍流黏度实际上在整个流动中变化很大（参见第 307 页上的图 7.30），因此对于大多数目的而言，所得的预测平均速度剖面将是毫无用处的不准确。

总之，均匀湍流黏度模型提供了自相似自由剪切流中平均速度分布的有用基本描述。然而，它是一个不完整的模型，适用范围非常有限。

\subsection*{10.2.2 混合长度模型}\phantomsection\addcontentsline{toc}{subsection}{10.2.2 混合长度模型}
在二维边界层流的应用中，[混合长度
获得为

\begin{sourceblock}
\includegraphics[page=395,pagebox=mediabox,trim={64.39bp 1707.78bp 64.39bp 239.97bp},clip,width=0.94\textwidth,max height=0.38\textheight]{Pope.pdf}
\end{sourceblock}
如7.1.7节所示，在对数律区域，混合长度的适当规格AS为lm ICYa,n d，则湍流黏度为\textasciitilde{} , K J J。

等式的几个概括提出方程（10.19）是为了允许将混合长度假设应用于所有流。基于
Sij 平均应变率 Smagorinsky (1963) 提出

\begin{sourceblock}
\includegraphics[page=395,pagebox=mediabox,trim={64.39bp 1347.83bp 64.39bp 591.93bp},clip,width=0.94\textwidth,max height=0.38\textheight]{Pope.pdf}
\end{sourceblock}
Rij，而根据 Baldwin 和 Lomax (19 78) 提出的平均旋转率

\begin{sourceblock}
\includegraphics[page=395,pagebox=mediabox,trim={64.39bp 1143.86bp 64.39bp 796.90bp},clip,width=0.94\textwidth,max height=0.38\textheight]{Pope.pdf}
\end{sourceblock}
（在 i3(Ul)/i3x2 是唯一非零平均速度梯度的情况下，这两个公式都简化为方程 (10.19)。）

混合长度模型的广义形式适用于所有湍流，并且可以说是最简单的湍流模型。然而，其主要缺点是其不完整性：需要指定混合长度 lm(x)h，并且适当的规格不可避免地取决于流动的几何形状。对于尚未研究的复杂流
em@) 之前， 的规范需要大量猜测，因此人们对计算得出的平均速度场的准确性几乎没有信心。另一方面，一些技术上重要的流已经被广泛研究，因此 lm(x) 的适当规范已经得到很好的建立。主要的例子是航空应用中的边界层流。 Cebeci-Smith 模型（Smith 和 Cebeci 1967）和 Baldwin-Lomax 模型（Baldwin 和 Lomax 1978）提供了混合长度规范，可以对附着的边界层进行相当准确的计算。 Wilcox (1993) 提供了这些模型及其性能的详细信息。

正如下面的练习所示，混合长度模型也可以应用于自由剪切流。预测的平均速度剖面与实验数据非常吻合（参见 Schlichting (19 79)）。该解决方案的一个有趣的（尽管非物理的）特征是混合层有一个明确的边缘，在该边缘平均速度为零
考虑自相似时间混合层，其中均值
横向速度 (V) 为零，轴向速度 (u) 取决于
仅 y 和 t。速度差为 Us，因此边界
条件为 ( U ) = ++Us at y -- - 层的厚度
6(t) 的定义（如 \textasciitilde{} i \textasciitilde{}\textasciitilde{} 5o.n 2pa1ge 140）使得 ( U ) +ZU，在
混合长度模型应用于该流动，其中混合
整个流量的长度是均匀的，并且与流量成正比
宽度，即 ad，其中 a 是指定常数。

从雷诺方程出发

\begin{sourceblock}
\includegraphics[page=396,pagebox=mediabox,trim={64.39bp 1253.84bp 64.39bp 669.92bp},clip,width=0.94\textwidth,max height=0.38\textheight]{Pope.pdf}
\end{sourceblock}
表明混合长度假设意味着

\begin{sourceblock}
\includegraphics[page=396,pagebox=mediabox,trim={64.39bp 1089.86bp 64.39bp 807.90bp},clip,width=0.94\textwidth,max height=0.38\textheight]{Pope.pdf}
\end{sourceblock}
证明该方程有以下形式的自相似解

(U) U,f(()，其中 y/6；并且 f ( ( ) 满足普通

\begin{sourceblock}
\includegraphics[page=396,pagebox=mediabox,trim={64.39bp 890.89bp 64.39bp 999.87bp},clip,width=0.94\textwidth,max height=0.38\textheight]{Pope.pdf}
\end{sourceblock}
其中 S d6/dt 是铺展率。

显示方程。 (10.24) 承认两个不同的解（表示为

\begin{sourceblock}
\includegraphics[page=396,pagebox=mediabox,trim={64.39bp 666.92bp 64.39bp 1257.84bp},clip,width=0.94\textwidth,max height=0.38\textheight]{Pope.pdf}
\end{sourceblock}
\begin{sourceblock}
\includegraphics[page=396,pagebox=mediabox,trim={64.39bp 579.93bp 64.39bp 1366.83bp},clip,width=0.94\textwidth,max height=0.38\textheight]{Pope.pdf}
\end{sourceblock}
其中 A、B 和 C 是任意常数。

f 的适当解由三部分组成。对于
大于特定值 f 是常数 ( i . :

\begin{sourceblock}
\includegraphics[page=396,pagebox=mediabox,trim={64.39bp 299.96bp 64.39bp 1581.80bp},clip,width=0.94\textwidth,max height=0.38\textheight]{Pope.pdf}
\end{sourceblock}
显示 < < 满足的适当解决方案
证明铺展率与混合长度常数有关

\begin{sourceblock}
\includegraphics[page=397,pagebox=mediabox,trim={64.39bp 1689.79bp 64.39bp 200.97bp},clip,width=0.94\textwidth,max height=0.38\textheight]{Pope.pdf}
\end{sourceblock}
并利用 6 ( e . f (l) 2) 的定义得到

\begin{sourceblock}
\includegraphics[page=397,pagebox=mediabox,trim={64.39bp 1554.80bp 64.39bp 387.95bp},clip,width=0.94\textwidth,max height=0.38\textheight]{Pope.pdf}
\end{sourceblock}
v\textasciitilde{} 在整个流程中如何变化？

\section*{10.3 湍动能模型}\phantomsection\addcontentsline{toc}{section}{10.3 湍动能模型}
将湍流黏度写为

\begin{sourceblock}
\includegraphics[page=397,pagebox=mediabox,trim={64.39bp 1198.85bp 64.39bp 744.91bp},clip,width=0.94\textwidth,max height=0.38\textheight]{Pope.pdf}
\end{sourceblock}
在混合长度模型中，长度尺度为，速度尺度为
Qm 是（简单剪切流）

\begin{sourceblock}
\includegraphics[page=397,pagebox=mediabox,trim={64.39bp 1016.87bp 64.39bp 904.89bp},clip,width=0.94\textwidth,max height=0.38\textheight]{Pope.pdf}
\end{sourceblock}
这意味着速度尺度是由平均速度梯度局部确定的；特别是，当 a ( U ) / a y 为零时，u* 为零。事实上，与此含义相反，存在多种情况，其中速度梯度为零，但湍流速度尺度不为零。一个例子是衰减的网格湍流；另一个是圆射流的中心线，直接测量显示 vT 远离零（参见第 108 页的图 5.10）。

Kolmogorov (1942) 和 Prandtl (1945) 独立地建议，最好将速度尺度建立在湍动能的基础上，即

\begin{sourceblock}
\includegraphics[page=397,pagebox=mediabox,trim={64.39bp 517.93bp 64.39bp 1425.82bp},clip,width=0.94\textwidth,max height=0.38\textheight]{Pope.pdf}
\end{sourceblock}
其中 c 是常数。如果再次将长度尺度视为混合长度，则湍流黏度变为

\begin{sourceblock}
\includegraphics[page=397,pagebox=mediabox,trim={64.39bp 337.96bp 64.39bp 1609.80bp},clip,width=0.94\textwidth,max height=0.38\textheight]{Pope.pdf}
\end{sourceblock}
如练习 10.3 所示，常数 c 0.55 的值在对数律区域中产生正确的行为。

为了方程。要使用 (10.34)，k(x,t) 的值必须已知或估计。 \textasciitilde{}ol mogorova 和 Prandtl 建议通过解决以下问题来实现这一目标
模型输运方程仅针对一个湍流量求解，即
在讨论 k 的模型传输方程之前，逐项列出模型的所有组成部分会很有帮助：

(i) 指定混合长度 Qm(x t)；

(ii) 求解模型输运方程 k(x, t)；

(iii) 湍流黏度由 vT 定义
(iiv) 雷诺应力由湍流黏度 hy- 获得
假设，方程。 (10.1) 和
(v) 求解雷诺方程 (U(x,t )) 和 (p(x,t ))。

因此，从规范和解决方案到精确和
Qm 模型方程，确定以下字段。d： (U)、(p)、em,k 和
VT，（uiuj）。这些被称为“已知”。

我们现在考虑 k 的模型传输方程。精确方程

\begin{sourceblock}
\includegraphics[page=398,pagebox=mediabox,trim={64.39bp 1036.87bp 64.39bp 759.90bp},clip,width=0.94\textwidth,max height=0.38\textheight]{Pope.pdf}
\end{sourceblock}
其中通量 T'（方程（5.140））为

\begin{sourceblock}
\includegraphics[page=398,pagebox=mediabox,trim={64.39bp 902.89bp 64.39bp 1027.87bp},clip,width=0.94\textwidth,max height=0.38\textheight]{Pope.pdf}
\end{sourceblock}
在等式中。 （10.35），任何完全由“已知”决定的项是
Dic/Dt P 据说是“封闭形式”。具体来说， 和 都是封闭形式。相反，其余项和 V T') 是“未知的”，并且为了获得一组封闭的模型方程，必须对这些项进行建模。也就是说，需要根据已知来对未知进行建模的“闭合近似”。

正如第 6 章中广泛讨论的，在高雷诺数下，耗散率尺度为 u i / l o，其中 uo 和 loa 是含能运动的速度尺度和长度尺度。因此，建模为合理的

\begin{sourceblock}
\includegraphics[page=398,pagebox=mediabox,trim={64.39bp 399.95bp 64.39bp 1542.81bp},clip,width=0.94\textwidth,max height=0.38\textheight]{Pope.pdf}
\end{sourceblock}
其中 CD 是模型常数。事实上，对对数律区域（练习 10.3）的检查得出了与 CD c 的关系
建模假设，例如方程。 （10.37）值得仔细审查。平等-

表达式（10.34）和（10.37）可以结合起来消除，得到

\begin{sourceblock}
\includegraphics[page=398,pagebox=mediabox,trim={64.39bp 141.98bp 64.39bp 1797.77bp},clip,width=0.94\textwidth,max height=0.38\textheight]{Pope.pdf}
\end{sourceblock}
\begin{sourceblock}
\includegraphics[page=399,pagebox=mediabox,trim={64.54bp 1266.68bp 64.54bp 109.48bp},clip,width=\textwidth,max height=0.72\textheight]{Pope.pdf}
\par\vspace{0.45\baselineskip}
\small 图 10.3。 vre/k的简介？ （参见例（10.39））来自 Re = 13,750 处槽道流的 DNS（Kim 等人，1987）。
\end{sourceblock}
或者，等价地，

\begin{sourceblock}
\includegraphics[page=399,pagebox=mediabox,trim={64.54bp 1048.74bp 64.54bp 898.77bp},clip,width=0.94\textwidth,max height=0.38\textheight]{Pope.pdf}
\end{sourceblock}
对于简单的剪切流，k 和 是可以测量的，所以
E、VT表明该建模假设可以直接进行检验。图 10.3 显示了方程的左侧。 （10.39）从完全发展的湍流河道流的 DNS 数据中提取。可以看出（除了靠近墙壁的地方，y+ < 50）这个量确实近似恒定，值在0.09左右。图 10.4 显示了时间混合层的相同数量：除了边缘附近之外，该值到处都接近 0.08。

湍动能方程中剩下的未知数是能量通量 T'（方程（10.36））。这是用梯度扩散建模的

\begin{sourceblock}
\includegraphics[page=399,pagebox=mediabox,trim={64.54bp 499.87bp 64.54bp 1364.66bp},clip,width=0.94\textwidth,max height=0.38\textheight]{Pope.pdf}
\end{sourceblock}
其中动能 1 i 的“湍流普朗特数”通常取 1.0。物理上，方程。 (10.40) 断言（由于速度和压力波动）存在 k 的通量沿着 k 的梯度。从数学上讲，该项确保所得到的 k 模型输运方程产生平滑解，并且可以在解域边界上的任意位置对 k 施加边界条件。

符号 ok 是标准符号。但请注意，ak 是标量，并且“k”不是 中的后缀
笛卡尔张量后缀表示法的意义。

\begin{sourceblock}
\includegraphics[page=400,pagebox=mediabox,trim={64.39bp 1108.86bp 64.39bp 109.44bp},clip,width=\textwidth,max height=0.72\textheight]{Pope.pdf}
\par\vspace{0.45\baselineskip}
\small 图 10.4。 vre/k的简介？ （参见方程（10.39））来自时间混合层的 DNS，（来自 Rogers 和 Moser (1994) 的数据）。）
\end{sourceblock}
综上所述，基于k的一方程模型由以下模型组成

\begin{sourceblock}
\includegraphics[page=400,pagebox=mediabox,trim={64.39bp 806.90bp 64.39bp 1041.87bp},clip,width=0.94\textwidth,max height=0.38\textheight]{Pope.pdf}
\end{sourceblock}
其中 vT = 和 = CDk Q 以及湍流黏度

\textasciitilde{} k l / \textasciitilde{} Q , E 3/2 m，假设（方程（10.1））和 Q, 的规范。

模型预测与实验数据的比较（Wilcox 1993）表明，这一单方程模型在精度方面比混合长度模型具有一定的优势。然而，不完整性的主要缺点仍然存在：必须指定长度尺度 Qm(x)m。

考虑壁面流动的对数律区域。使用对数定律
a n d 规范 = rcy 以表明适当的值
常数 c（在关系 vT ck1j2Q 中）m 是

\begin{sourceblock}
\includegraphics[page=400,pagebox=mediabox,trim={64.39bp 237.97bp 64.39bp 1707.78bp},clip,width=0.94\textwidth,max height=0.38\textheight]{Pope.pdf}
\end{sourceblock}
使用关系式表明

\begin{sourceblock}
\includegraphics[page=400,pagebox=mediabox,trim={64.39bp 119.39bp 64.39bp 1840.77bp},clip,width=0.94\textwidth,max height=0.38\textheight]{Pope.pdf}
\end{sourceblock}
\section*{10.4 k-ε 模型}\phantomsection\addcontentsline{toc}{section}{10.4 k-ε 模型}
\begin{sourceblock}
\includegraphics[page=401,pagebox=mediabox,trim={64.39bp 1725.78bp 64.39bp 156.98bp},clip,width=0.94\textwidth,max height=0.38\textheight]{Pope.pdf}
\end{sourceblock}
对于应用于简单剪切流的一方程模型，表示为
以 k 和 a ( U ) y 表示的产量。因此（取 CD = c-7
在等式中(10.37)) 表明一元方程中的速度尺度
模型 和 在混合长度模型中的关系为

\begin{sourceblock}
\includegraphics[page=401,pagebox=mediabox,trim={64.39bp 1376.83bp 64.39bp 513.93bp},clip,width=0.94\textwidth,max height=0.38\textheight]{Pope.pdf}
\end{sourceblock}
证明一般流的对应关系为

\begin{sourceblock}
\includegraphics[page=401,pagebox=mediabox,trim={64.39bp 1223.85bp 64.39bp 715.91bp},clip,width=0.94\textwidth,max height=0.38\textheight]{Pope.pdf}
\end{sourceblock}
（参见等式（10.20））。

10.4 k-E模型

\subsection*{10.4.1 概述}\phantomsection\addcontentsline{toc}{subsection}{10.4.1 概述}
k-E模型属于哪个模型类
在两个方程模型中，输运方程针对两个湍流量 i-emk 和 k-E 模型进行求解。从这两个量可以形成长度尺度（L k 3/2 a 时间尺度 k / \textasciitilde{} ）a、维度量（k2/e）等。

= E ) = 因此，二方程模型可以依赖于流量
不需要完整的规范，例如 Qrn(xa)r e。

k-E 模型是使用最广泛的完整湍流模型，并且被纳入大多数商业 CFD 代码中。与所有湍流模型一样，概念和细节都随着时间的推移而演变。但 Jones 和 Launder (1972) 因开发“标准”k-e 模型而受到适当的赞誉，Launder 和 Sharma (1974) 提供了模型常数的改进值。 Davidov (1961)、Harlow 和 Nakayama (1968)、Hanjalik (1970) 以及 Launder 和 Spalding (1972) 引用的其他人做出了重要的早期贡献。

除了湍流黏度假设之外，k-E 模型还包括
(i) k 的模型传输方程（与
一方程模型，Eq. (10.4 一));

(iii) 湍流黏度的规格为

\begin{sourceblock}
\includegraphics[page=402,pagebox=mediabox,trim={64.39bp 1719.78bp 64.39bp 219.97bp},clip,width=0.94\textwidth,max height=0.38\textheight]{Pope.pdf}
\end{sourceblock}
其中 C, 0.09 是五个模型常数之一。

如果假设 vT 仅取决于湍流量 k 和
（独立于 a(Ui)/dxj 等），则方程(1.0.47) 是不可避免的。 oneequation 模型隐含着相似的关系 vT = ck （参见练习 10.3） sol 模型常数的关系为
在简单的湍流剪切流中，k-E 模型得出

\begin{sourceblock}
\includegraphics[page=402,pagebox=mediabox,trim={64.39bp 1322.83bp 64.39bp 614.92bp},clip,width=0.94\textwidth,max height=0.38\textheight]{Pope.pdf}
\end{sourceblock}
（参见练习 10.5），因此规范 C, 0.09 (0.3)\textasciitilde{} 源自 P/Ei 接近统一的区域中的经验观察 ( ( u v ) ( / k 0.3。

vT\&/k2p 的数量如图 1 和 2 所示。 10.3 和 10.4 是槽道流和时间混合层的 C 的“测量”。可以看出，除了流边界附近之外，vT\&/k2 在任何地方都接近值 0.09。

考虑应用于简单湍流剪切流的 k-E 模型：

S ( U ) y 是唯一的非零平均速度梯度。获取

\begin{sourceblock}
\includegraphics[page=402,pagebox=mediabox,trim={64.39bp 766.90bp 64.39bp 1094.86bp},clip,width=0.94\textwidth,max height=0.38\textheight]{Pope.pdf}
\end{sourceblock}
\begin{sourceblock}
\includegraphics[page=402,pagebox=mediabox,trim={64.39bp 657.92bp 64.39bp 1257.84bp},clip,width=0.94\textwidth,max height=0.38\textheight]{Pope.pdf}
\end{sourceblock}
从而验证方程。 （10.48）。

证明 (uv) 满足 Cauchy-Schwartz 不等式（Eq. (3.100)

当且仅当 C 满足

\begin{sourceblock}
\includegraphics[page=402,pagebox=mediabox,trim={64.39bp 370.95bp 64.39bp 1529.81bp},clip,width=0.94\textwidth,max height=0.38\textheight]{Pope.pdf}
\end{sourceblock}
或者，等价地，
显示方程。 (10.50) 也适用于一般流程。

\subsection*{10.4.2 e 的模型方程}\phantomsection\addcontentsline{toc}{subsection}{10.4.2 e 的模型方程}
在建立 k 和 k 的模型输运方程时，采用了完全不同的方法。 k 方程相当于精确方程（方程（10.35）），其中湍流通量 T' 被建模为梯度扩散（方程（10.40））。其他三个项 Dk/Dt 和 是闭合形式（给定湍流黏度假设）。

也可以导出 的精确方程，但它不是模型方程的有用起点。这是因为（如第 6 章中所讨论的）c 最好被视为级联中的能量流速，并且它由大规模运动决定，与粘度（在高雷诺数下）无关。相比之下， 的精确方程适用于耗散区内的过程。因此，c 的标准模型方程最好被视为完全是经验性的，而不是基于精确的方程：它
Launder 和 Sharma (1974) 提出的所有模型常数的标准值为

\begin{sourceblock}
\includegraphics[page=403,pagebox=mediabox,trim={64.39bp 934.88bp 64.39bp 1009.87bp},clip,width=0.94\textwidth,max height=0.38\textheight]{Pope.pdf}
\end{sourceblock}
通过研究方程在各种流中的行为可以了解该方程。我们首先检查均匀湍流，其中 k 和方程变为

\begin{sourceblock}
\includegraphics[page=403,pagebox=mediabox,trim={64.39bp 683.91bp 64.39bp 1223.85bp},clip,width=0.94\textwidth,max height=0.38\textheight]{Pope.pdf}
\end{sourceblock}
\begin{sourceblock}
\includegraphics[page=403,pagebox=mediabox,trim={64.39bp 558.93bp 64.39bp 1348.83bp},clip,width=0.94\textwidth,max height=0.38\textheight]{Pope.pdf}
\end{sourceblock}
衰减湍流 在没有平均速度梯度的情况下，产生量为零，并且湍流衰减。方程就有解

\begin{sourceblock}
\includegraphics[page=403,pagebox=mediabox,trim={64.39bp 254.97bp 64.39bp 1634.79bp},clip,width=0.94\textwidth,max height=0.38\textheight]{Pope.pdf}
\end{sourceblock}
其中 k 和 在参考时间的值为 ko 和

\begin{sourceblock}
\includegraphics[page=403,pagebox=mediabox,trim={64.39bp 119.39bp 64.39bp 1833.77bp},clip,width=0.94\textwidth,max height=0.38\textheight]{Pope.pdf}
\end{sourceblock}
衰减指数 n 为

\begin{sourceblock}
\includegraphics[page=404,pagebox=mediabox,trim={64.39bp 1690.79bp 64.39bp 213.97bp},clip,width=0.94\textwidth,max height=0.38\textheight]{Pope.pdf}
\end{sourceblock}
这种幂律衰减正是在网格湍流中观察到的（参见选择 5.4.6，方程（5.274）和（5.277）），因此 c 方程的行为对于该流是正确的4
报告的衰减指数 n 的实验值通常在 1.15-1.45 范围内，Mohamed 和 LaRue (1990) 表明大多数数据与 n 1.3 一致。方程 (10.59) 可以重新排列以给出 CE2i n 的 n 项：

\begin{sourceblock}
\includegraphics[page=404,pagebox=mediabox,trim={64.39bp 1284.84bp 64.39bp 645.92bp},clip,width=0.94\textwidth,max height=0.38\textheight]{Pope.pdf}
\end{sourceblock}
n 1.15、1.3、1.45对应的CE2c值为1.87、1.77、1.69。那么可以看出，标准值（CE2 1492）在一定程度上超出了实验观察到的范围。下面讨论其原因。

考虑应用于衰减湍流的 k-c 模型。令 s(t) 为
归一化时间定义为

(a) 获得显式表达式
(b) 推导并求解 s 中 k 和 的演化方程（即
(c) 检查 x / M 40 和 x / M 之间的网格湍流
200. 这对应于标准化时间的什么间隔？

证明 k-E 模型给出了最终结果的正确行为
衰减期（参见第 205 页的练习 6.11），如果 CE2i 修改为
均匀剪切 如第 5.4.5 节中所观察到的，在均匀湍流剪切流中，主要的实验观察结果是雷诺应力变得自相似，并且无量纲参数 S k l 和 P/Eb 变得自相似。

常数，S \textasciitilde{} /imEpl 的恒定性意味着湍流时间尺度 z k / \textasciitilde{} 也是固定的。从 k 和方程（方程（10.55）和（10.56））我们得到

\begin{sourceblock}
\includegraphics[page=405,pagebox=mediabox,trim={64.39bp 1642.79bp 64.39bp 257.97bp},clip,width=0.94\textwidth,max height=0.38\textheight]{Pope.pdf}
\end{sourceblock}
显然，该模型预测不会随时间变化
PIE，特定值

\begin{sourceblock}
\includegraphics[page=405,pagebox=mediabox,trim={64.39bp 1434.82bp 64.39bp 468.94bp},clip,width=0.94\textwidth,max height=0.38\textheight]{Pope.pdf}
\end{sourceblock}
比实验和 DNS 中观察到的值要高得多。

方程的解释 现在我们根据以下关系提供方程的解释
= \textasciitilde{} l k 湍流频率 w 和特征平均应变率 S 之间的关系。

一个过于简单的模型是

\begin{sourceblock}
\includegraphics[page=405,pagebox=mediabox,trim={64.39bp 1027.87bp 64.39bp 920.88bp},clip,width=0.94\textwidth,max height=0.38\textheight]{Pope.pdf}
\end{sourceblock}
fl 其中 是常数；等于 3，比如说。然后该模型预测
fl S k l \textasciitilde{} S/o。等于所有流中的常数值3。在几种剪切流中，确实测量了该 S k l c 值（参见第 131、157 和 283 页上的表 5.2、5.4 和 7.2）。然而，对于其他流动，该模型是错误的：在网格湍流 (S 0) 中，w 不为零；在均匀剪切流中，S k l \textasciitilde{} 约为 6。

不要将 w 设置为 SIP，而是考虑一个模型，该模型使得
SIP。 w 松弛到 Or （作为产生所需结果的设计）考虑使 w2 松弛到的模型 当应用于均匀湍流时，该模型由以下方程描述

\begin{sourceblock}
\includegraphics[page=405,pagebox=mediabox,trim={64.39bp 441.94bp 64.39bp 1454.82bp},clip,width=0.94\textwidth,max height=0.38\textheight]{Pope.pdf}
\end{sourceblock}
其中 a 是常数，aw 是弛豫率。这是一个代数问题（参见练习 10.8），如果常数指定为，则该方程完全等价于方程（方程（10.56））

\begin{sourceblock}
\includegraphics[page=405,pagebox=mediabox,trim={64.39bp 224.97bp 64.39bp 1721.78bp},clip,width=0.94\textwidth,max height=0.38\textheight]{Pope.pdf}
\end{sourceblock}
\begin{sourceblock}
\includegraphics[page=405,pagebox=mediabox,trim={64.39bp 119.39bp 64.39bp 1797.77bp},clip,width=0.94\textwidth,max height=0.38\textheight]{Pope.pdf}
\end{sourceblock}
因此，c 方程可以通过方程 2 来解释。 (10.64)：在速率 am 时，湍流频率（平方）向 SIP（平方）松弛。

对于均匀湍流，根据 k 的模型方程
（方程（10.55））和 m2（方程（10.64）），表明相应的
c 的模型方程为

\begin{sourceblock}
\includegraphics[page=406,pagebox=mediabox,trim={64.39bp 1410.82bp 64.39bp 482.94bp},clip,width=0.94\textwidth,max height=0.38\textheight]{Pope.pdf}
\end{sourceblock}
通过使用等式。 (10.50) 消去S，将方程重新表示为

\begin{sourceblock}
\includegraphics[page=406,pagebox=mediabox,trim={64.39bp 1229.84bp 64.39bp 663.92bp},clip,width=0.94\textwidth,max height=0.38\textheight]{Pope.pdf}
\end{sourceblock}
通过将此结果与标准 c 方程（方程（10.56））进行比较，
验证 a 和 /?,a 和 CE1a 和 Cg2( 等式 (10.65)

对数律区域中的行为 对于非均匀流，c 方程 ( e , V [ ( v \textasciitilde{} / o \textasciitilde{} ) \textasciitilde{}haEs] t)h e 中的扩散项与 k 方程中的类似项具有相同的优点：它确保我们获得平滑解，并且允许在边界上的任何位置指定 c 上的边界条件。正如现在所说明的，扩散项在近壁流动中起着重要作用。

考虑高雷诺数、充分发展的槽道流。感兴趣的量（ ( U ) 、 k 和 E）仅取决于 y，因此 k-c 方程

\begin{sourceblock}
\includegraphics[page=406,pagebox=mediabox,trim={64.39bp 478.94bp 64.39bp 1370.83bp},clip,width=0.94\textwidth,max height=0.38\textheight]{Pope.pdf}
\end{sourceblock}
\begin{sourceblock}
\includegraphics[page=406,pagebox=mediabox,trim={64.39bp 339.96bp 64.39bp 1555.80bp},clip,width=0.94\textwidth,max height=0.38\textheight]{Pope.pdf}
\end{sourceblock}
我们现在关注对数律区域。生产和耗散平衡，均等于u3/(rcy)。因此，在 k 方程中，扩散项为零，这意味着 k 是均匀的，这大致是正确的。在 c 方程中， 和 c 相等导致净汇（等于 -(CE2 c c l ) c 2 / k ）

练习 10.9 表明方程 (Eq. (10-70)) 满足

\begin{sourceblock}
\includegraphics[page=407,pagebox=mediabox,trim={64.39bp 1683.79bp 64.39bp 239.97bp},clip,width=0.94\textwidth,max height=0.38\textheight]{Pope.pdf}
\end{sourceblock}
并且常数之间的关系为

\begin{sourceblock}
\includegraphics[page=407,pagebox=mediabox,trim={64.39bp 1565.80bp 64.39bp 373.95bp},clip,width=0.94\textwidth,max height=0.38\textheight]{Pope.pdf}
\end{sourceblock}
方程 (10.72) 得出 k-E 模型 0.43) 隐含的冯卡门常数的值：或者，它可以用于调整模型常数（例如，o）以产生特定的值
考虑壁面湍流的对数律区域。显示
湍流黏度假设 v\textasciitilde{} C,,k2/\&a 和
对数定律意味着

\begin{sourceblock}
\includegraphics[page=407,pagebox=mediabox,trim={64.39bp 1129.86bp 64.39bp 789.90bp},clip,width=0.94\textwidth,max height=0.38\textheight]{Pope.pdf}
\end{sourceblock}
从表达式给出以获得结果

\begin{sourceblock}
\includegraphics[page=407,pagebox=mediabox,trim={64.39bp 944.88bp 64.39bp 949.88bp},clip,width=0.94\textwidth,max height=0.38\textheight]{Pope.pdf}
\end{sourceblock}
因此验证方程。 （10.71）并将其与式（10.71）进行比较。 （10.43）。替补
等式。 (10.71) 代入方程 (Eq. (10.70), 其中 k 个独立
y) 以获得方程。 （10.72）。验证这些方程得出
长度尺度的变化

\begin{sourceblock}
\includegraphics[page=407,pagebox=mediabox,trim={64.39bp 660.92bp 64.39bp 1243.84bp},clip,width=0.94\textwidth,max height=0.38\textheight]{Pope.pdf}
\end{sourceblock}
自由流边缘的行为 如第 5.5.2 节所述，在湍流和非旋转、非湍流自由流或静止环境之间存在一个间歇区域。在非湍流流体 (y m) 中，k 和 均为零。

k-E 模型没有考虑间歇性，并且它产生的解在湍流和非湍流区域之间存在尖锐的边缘（Cazalbou、Spalart 和 Bradshaw 1994）。对于统计二维边界层或自由剪切流，用 ye@) 表示边缘的位置。那么，对于 y ye@)，k 和 为零，并且平均流是无旋的。就在湍流区域内部，k 和 a ( U ) / a y 随着以下各项的正幂而变化
由于夹带，存在从非湍流侧通过边缘的平均流量。因此，该边缘对应于通过湍流扩散传播（相对于流体）的湍流前沿。这种传播前沿的 k-E 模型方程的解在练习 10.10 中得到。

考虑统计上平稳的一维流，没有转动
bulence 产生，其中 k-E 方程简化为

\begin{sourceblock}
\includegraphics[page=408,pagebox=mediabox,trim={64.39bp 1188.85bp 64.39bp 708.91bp},clip,width=0.94\textwidth,max height=0.38\textheight]{Pope.pdf}
\end{sourceblock}
其中 Uo 是均匀平均速度，为正值。这些
方程有一个弱解（k O 和 O 代表 x 0
和 k 0 和 0 对于 x 0) 对应于前面之间
湍流 (x 0) 和非湍流 (x I0)。对于小
正 x 的解是

\begin{sourceblock}
\includegraphics[page=408,pagebox=mediabox,trim={64.39bp 834.89bp 64.39bp 1060.87bp},clip,width=0.94\textwidth,max height=0.38\textheight]{Pope.pdf}
\end{sourceblock}
其中p、q、ko和+6 0 是正常数，其中ko3/ 2 /co。

通过代入等式。 （10.78）和（10.79）代入等式。 （10.76）和（10.77），
获得结果

\begin{sourceblock}
\includegraphics[page=408,pagebox=mediabox,trim={64.39bp 473.94bp 64.39bp 1474.81bp},clip,width=0.94\textwidth,max height=0.38\textheight]{Pope.pdf}
\end{sourceblock}
\begin{sourceblock}
\includegraphics[page=408,pagebox=mediabox,trim={64.39bp 335.96bp 64.39bp 1557.80bp},clip,width=0.94\textwidth,max height=0.38\textheight]{Pope.pdf}
\end{sourceblock}
并表明（对于小 x）对流和扩散平衡，其中
相比之下，耗散项可以忽略不计。从方程。 (10.8 1)

\begin{sourceblock}
\includegraphics[page=408,pagebox=mediabox,trim={64.39bp 119.39bp 64.39bp 1758.78bp},clip,width=0.94\textwidth,max height=0.38\textheight]{Pope.pdf}
\end{sourceblock}
其中 oE/akS。如何（对于小 X） k, \&, L k 3/2 /E, o, -= \textasciitilde{} / , k
并随 x 的以下幂变化：

\begin{sourceblock}
\includegraphics[page=409,pagebox=mediabox,trim={64.39bp 1647.79bp 64.39bp 254.97bp},clip,width=0.94\textwidth,max height=0.38\textheight]{Pope.pdf}
\end{sourceblock}
分别，因此所有这些量在
当 x 接近 0 时的极限，前提是 o 介于 1 和

表明湍流前沿的传播速度（相对
到平均流量）是
证明适用于自相似温度的 k-E 模型方程
多孔混合层可以写成

\begin{sourceblock}
\includegraphics[page=409,pagebox=mediabox,trim={64.39bp 1146.86bp 64.39bp 749.91bp},clip,width=0.94\textwidth,max height=0.38\textheight]{Pope.pdf}
\end{sourceblock}
\begin{sourceblock}
\includegraphics[page=409,pagebox=mediabox,trim={64.39bp 1029.87bp 64.39bp 866.89bp},clip,width=0.94\textwidth,max height=0.38\textheight]{Pope.pdf}
\end{sourceblock}
6(t) 为层的宽度，相似度变量为

\begin{sourceblock}
\includegraphics[page=409,pagebox=mediabox,trim={64.39bp 836.89bp 64.39bp 1017.87bp},clip,width=0.94\textwidth,max height=0.38\textheight]{Pope.pdf}
\end{sourceblock}
其中 U s 是（恒定）速度差。变换方程。 ( 10.85)

和 (10.86) 获得

\begin{sourceblock}
\includegraphics[page=409,pagebox=mediabox,trim={64.39bp 611.92bp 64.39bp 1264.84bp},clip,width=0.94\textwidth,max height=0.38\textheight]{Pope.pdf}
\end{sourceblock}
其中 S 是撒播率 S U-' d6ldt。

让 表示湍流区域的边缘（因此 k 和 是
2 为零，x 为距边缘的距离，

\begin{sourceblock}
\includegraphics[page=409,pagebox=mediabox,trim={64.39bp 264.97bp 64.39bp 1681.79bp},clip,width=0.94\textwidth,max height=0.38\textheight]{Pope.pdf}
\end{sourceblock}
显示方程的左侧。 (10.88) 可以写成

\begin{sourceblock}
\includegraphics[page=409,pagebox=mediabox,trim={64.39bp 119.39bp 64.39bp 1812.77bp},clip,width=0.94\textwidth,max height=0.38\textheight]{Pope.pdf}
\end{sourceblock}
因此表明（对于小 x）方程具有幂律
练习 10.10 中给出的解（k \textasciitilde{} x? 和 ¢ \textasciitilde{} x‘），

\subsection*{10.4.3 讨论}\phantomsection\addcontentsline{toc}{subsection}{10.4.3 讨论}
k-s模型可以说是最简单的完全湍流模型，因此它具有最广泛的适用范围。它被纳入大多数商业 CFD 代码中，并已应用于各种问题，包括传热、燃烧和多相流。

关于其准确性的讨论推迟到下一章（第 11.10 节），其中将其性能与其他湍流模型的性能进行了比较。简要地;尽管对于简单流它通常是可以接受的准确度，但对于复杂流它可能非常不准确，以至于计算出的平均流模式在质量上可能是不正确的。误差源于湍流黏度假设和 Δ 方程。

对近壁处理的讨论也推迟到下一章（第 11.7 节）。为了将其应用于粘性近壁区域，需要对标准 k--g 模型进行修改。例如，从图 10.3 中可以看出，随着 y+ 降低到 50 以下，C 的适当值也随之降低。

标准 k-e 模型常数（方程 10.54）的值代表了一种折衷。对于任何特定的流量，可以通过调整常数来提高模型计算的准确性。对于衰减湍流，C,. =1.77（对应衰减指数n=1.3的值）比C,更合适。 = 1.92（对应于衰减指数 n = 1.09 的值）。 k-e 模型的一个众所周知的缺陷是它明显高估了圆射流的扩散速率。这个问题可以通过调整 Cy 或 Cy 的值来解决（对于圆射流）。然而，这种依赖于流量的临时调整的价值有限。对于一个完整的、普遍适用的模型，需要单一的常数规范；标准值代表（通过主观判断）选择的折衷方案，以便为一系列流量提供“最佳”性能。

多年来，人们提出了对标准 k-e 模型的许多“修改”，通常的动机是纠正特定类别流的不良性能。例如，Pope (1978) 提出了在形式的 Δ 方程中添加一个附加源项，以便修改后的模型能够产生圆射流的正确扩散率。 Hanjali 和 Launder (1980) 以及 Bardina、Ferziger 和 Reynolds (1983) 等人提出了对 Δ 方程的其他修改（基于 Q，;）。当这些
修改后的模型应用于一系列流量，一般经验（Launder 11990，H anjalib 1994）是它们的整体性能不如标准模型。

此处提出的模型方程完全是经验性的；但这只是一种观点。重正化群方法 (RNG) 已用于从 Navier-Stokes 方程获得 k-E 方程（参见，例如，Yakhot 和 Orszag (1986)、Smith 和 Reynolds (1992)、Smith 和 Woodruff (1998) 以及其中的参考文献）。来自 R N G 分析的常数值为

\begin{sourceblock}
\includegraphics[page=411,pagebox=mediabox,trim={64.39bp 1380.83bp 64.39bp 563.93bp},clip,width=0.94\textwidth,max height=0.38\textheight]{Pope.pdf}
\end{sourceblock}
（Orszag 等 1996），参见。等式。 （10.54）。在 R N G k-E 模型中还有一个
等人。方程中的附加项，这是一个临时模型，不是从 R N G 理论导出的。正是这个术语造成了标准模型和 R N G 模型的性能差异。

\section*{10.5 其他湍流黏度模型}\phantomsection\addcontentsline{toc}{section}{10.5 其他湍流黏度模型}
\subsection*{10.5.1 k-ω 模型}\phantomsection\addcontentsline{toc}{subsection}{10.5.1 k-ω 模型}
历史上，已经提出了许多二方程模型。在大多数情况下，k被视为变量之一，但第二个变量有多种选择。例如量纲为 kL (Rotta 1951)、co (Kolmogorov 1942)、(Saffman 1970) 和 (Speziale、Abid 和 Anderson 1992) 的量。

对于均匀湍流，第二个变量的选择并不重要，因为各个方程之间存在精确的对应关系，并且它们的形式本质上是相同的（参见练习10.12）。对于非均匀流，差异在于扩散项。例如，考虑以下 co \textasciitilde{} / k 的模型方程：

大众) Pco
Dt 基于该方程的 k - o 模型与 k-E 模型有何不同？回答这个问题的一种方法是推导 k-E 模型隐含的协方程（参见练习 10.13）。为简单起见，取 o, o，结果
Dco Pco
V * ( 2 V o ) 1) k (CE2 1)co2

\begin{sourceblock}
\includegraphics[page=411,pagebox=mediabox,trim={64.39bp 119.39bp 64.39bp 1808.77bp},clip,width=0.94\textwidth,max height=0.38\textheight]{Pope.pdf}
\par\vspace{0.45\baselineskip}
\small 表 10.2。 Cz 的值，以及各种 Z 规格的 Cz =（参见练习 10.12）
\end{sourceblock}
\begin{sourceblock}
\includegraphics[page=412,pagebox=mediabox,trim={64.39bp 1122.86bp 64.39bp 255.97bp},clip,width=\textwidth,max height=0.72\textheight]{Pope.pdf}
\end{sourceblock}
第二个最广泛使用的双方程模型是 k-wb 模型，它由 Wilcox 等人开发了 20 多年（参见 Wilcox (1993)）。该模型中v\textasciitilde{}的表达式和k方程与k-E模型中的相同。区别在于方程式的使用。 (10.93) 对于 w 而不是方程。 （10.53）（或隐含的 w 方程，方程（10.94））。正如 Wilcox (1993) 详细描述的，对于边界层流，k-o, 模型在处理粘性近壁区域和解释流向压力梯度的影响方面都更胜一筹。然而，非湍流自由流边界的处理是有问题的：需要 w 上的非零（非物理）边界条件，并且计算的流量对指定的值敏感。

Menter (1994) 提出了一个二方程模型，旨在产生 k-r 和 k-o 模型的最佳行为。它被写成一个（非标准）k-w 模型，方程的形式为： （10.94），但是最后一项乘以“混合函数”。靠近墙壁，混合函数为零（导致标准协方程），而远离墙壁，混合函数为单位（对应于标准方程）。 （Cazalbou 等人 (1994) 已经分析了自由流边界处的 k-c 模型的行为；另请参见练习 10.11。）

考虑由下式定义的数量 Z

\begin{sourceblock}
\includegraphics[page=412,pagebox=mediabox,trim={64.39bp 119.39bp 64.39bp 1839.77bp},clip,width=0.94\textwidth,max height=0.38\textheight]{Pope.pdf}
\end{sourceblock}
对于给定的 C1、p 和 q。根据标准 k-E 模型方程，显示
在均匀湍流中，Z 的隐含模型方程为

\begin{sourceblock}
\includegraphics[page=413,pagebox=mediabox,trim={64.39bp 1602.80bp 64.39bp 253.97bp},clip,width=0.94\textwidth,max height=0.38\textheight]{Pope.pdf}
\end{sourceblock}
\begin{sourceblock}
\includegraphics[page=413,pagebox=mediabox,trim={64.39bp 1539.81bp 64.39bp 405.95bp},clip,width=0.94\textwidth,max height=0.38\textheight]{Pope.pdf}
\end{sourceblock}
\begin{sourceblock}
\includegraphics[page=413,pagebox=mediabox,trim={64.39bp 1461.82bp 64.39bp 485.94bp},clip,width=0.94\textwidth,max height=0.38\textheight]{Pope.pdf}
\end{sourceblock}
因此验证表 10.2 中 CZ1a 和 CZ2i 的值。

给定标准 k-E 模型方程，表明隐含模型
equal1o:n 表示 is

\begin{sourceblock}
\includegraphics[page=413,pagebox=mediabox,trim={64.39bp 1010.87bp 64.39bp 795.90bp},clip,width=0.94\textwidth,max height=0.38\textheight]{Pope.pdf}
\end{sourceblock}
\subsection*{10.5.2 Spalart-Allmaras模型}\phantomsection\addcontentsline{toc}{subsection}{10.5.2 Spalart-Allmaras模型}
Spalart 和 Allmaras (1994) 描述了一种为空气动力学应用而开发的单方程模型，其中求解湍流黏度 vT 的单一模型传输方程。 Nee 和 Kovasznay (1969) 以及 Baldwin 和 Barth (1990) 描述了这种模型的早期建议。

从一开始就了解 Spalart-Allmaras 模型的发展背景是很有用的。，上述代数模型、一方程、二方程自然进展到下一章中描述的雷诺应力模型。每个连续的级别都提供了对湍流的更全面的描述，从而消除了前一个级别的质量缺陷。如果准确性是选择开发和应用模型的唯一标准，那么选择自然会倾向于描述水平较高的模型。然而，正如第 8.3 节中所讨论的，成本和易用性也是有利于更简单模型的重要标准。因此，对于模型开发人员来说，在每个描述级别上努力构建最佳模型是很有用的。

可以说，一元方程 m \textasciitilde{} delfor vT 是满足以下条件的最低水平：

消除了基于 k 的代数模型和一方程模型的不完整性，并且具有比二方程模型计算更简单的模式 1。该模型专为空气动力学流动而设计，例如翼型上的跨音速流动，包括边界层分离。

模型方程的形式为

\begin{sourceblock}
\includegraphics[page=414,pagebox=mediabox,trim={64.39bp 1514.81bp 64.39bp 382.95bp},clip,width=0.94\textwidth,max height=0.38\textheight]{Pope.pdf}
\end{sourceblock}
Sv，其中源项取决于层流和湍流黏度 v 和平均涡度（或旋转率）Q；湍流黏度
VT；梯度IVvTI；以及到最近墙壁的距离 L,。该模型的细节非常复杂：读者可以参考原始论文，该论文提供了关于模型构建以实现特定期望行为的启发性说明。

在空气动力学流动的应用中，该模型已被证明相当成功（参见 Godin、Zingg 和 Nelson (1997)）。然而，作为通用模型，它有明显的局限性。例如;它无法解释各向同性湍流的衰减，这意味着，
均匀湍流中的 VT 不受无旋平均应变的影响；

VT 将飞机喷射的扩散率高估了近 40\%。

将 Spalart-allmar 视为应用于高雷诺数的模型
数均匀湍流。认为层流粘度
v 和到最近墙壁的距离 l 不是相关的量
关系。因此表明尺寸和其他考虑因素减少了

\begin{sourceblock}
\includegraphics[page=414,pagebox=mediabox,trim={64.39bp 620.92bp 64.39bp 1249.84bp},clip,width=0.94\textwidth,max height=0.38\textheight]{Pope.pdf}
\end{sourceblock}
其中 是一个常数。评论这个方程的形式
无旋平均应变。

对于自相似均匀湍流剪切流（其中
和 S 相等），根据关系 a ( U ) / a y，表明
'由 Eqw(10.101) 演变而来

\begin{sourceblock}
\includegraphics[page=414,pagebox=mediabox,trim={64.39bp 282.96bp 64.39bp 1608.80bp},clip,width=0.94\textwidth,max height=0.38\textheight]{Pope.pdf}
\end{sourceblock}
使用实验数据（第 157 页表 5.4）来估计
根据方程式(10.102) 并将其与 Spalart-Allmaras 进行比较
值为 0.135。

\chapter{雷诺应力及相关模型}
\section*{11.1 引言}\phantomsection\addcontentsline{toc}{section}{11.1 引言}
雷诺应力闭合 在雷诺应力模型中，模型传输方程针对各个雷诺应力 (uiui) 和耗散（或其他量，例如提供湍流的长度或时间尺度）求解。因此，不需要湍流黏度假设；这样就消除了上一章中描述的模型的一个主要缺陷。

雷诺应力的精确输运方程可从第 3 19-3 19 页练习 7.23-7.25 中的纳维-斯托克斯方程获得：

\begin{sourceblock}
\includegraphics[page=415,pagebox=mediabox,trim={64.39bp 875.89bp 64.39bp 1021.87bp},clip,width=0.94\textwidth,max height=0.38\textheight]{Pope.pdf}
\end{sourceblock}
在雷诺应力模型中，“已知”为 (U)、(p)、(uiuj)，因此方程 1 中的“已知”为 (U)、(p)、(uiuj)。 (11.1) 平均流对流 D ( u + um ) / Dat 和 n d 生产张量 Pij （方程（7.179））均为闭合形式。然而，耗散张量模型（方程 (7 . 1 1 ) ）压力-应变率率张量 Rij
需要 Eij（方程（7.187））和雷诺应力通量 Tkij（方程（7.195））。

本章概述 到目前为止，要建模的最重要的量是压力-应变率率张量 Rij。在接下来的四节中，该术语将在均匀湍流的背景下得到广泛考虑。这包括（第 11.4 节）快速扭曲理论（RDT）的描述，该理论适用于极限情况，并提供了有用的见解。第 11.6 节描述了非均匀流的扩展，第 11.7 节描述了 k-E 和雷诺应力模型的特殊近壁处理。关于雷诺应力模型有大量文献，有许多不同的建议和变体：这里的重点是基本概念和方法。

在雷诺应力模型中，Rij 被建模为 (u-u-) 的局部函数
和 8 (Ui) / a x j。椭圆松弛模型（第 11.8 节）提供了更高级别的闭合，从而允许 Rij 的模型是非局部的。

代数应力模型和非线性湍流黏度模型在第 11.9 节中描述。这些是可以从雷诺应力闭合导出的更简单的模型。本章最后对所描述的一系列模型的相对优点进行了评估。

这里对耗散进行了概括性的处理。对于高雷诺数流动，局部各向同性的结果是1
这被视为对于中等雷诺数流动的模型，
然而，这种各向同性关系并不完全准确（例如，参见第 318 页的图 7.39）。然而，在某种程度上，这并不重要，因为各向异性分量 ( e . ?Edij) 具有与 Rij,and 相同的数学属性
E i j - 3 因此可以被吸收到第 11.7 节中讨论的 Rij.A S 模型中，靠近壁的耗散是各向异性的，并且不同的模型是合适的。

雷诺数

大多数雷诺应力模型不包含雷诺数依赖性（近壁处理除外），因此它们隐含地假设建模项与雷诺数无关。为了简化说明，我们遵循这个权宜的假设。然而，最好记住，在适度的雷诺数实验中，特别是在 DNS 中，可能存在（通常是适度的）雷诺数效应。

\section*{11.2 压力-应变率张量}\phantomsection\addcontentsline{toc}{section}{11.2 压力-应变率张量}
脉动压力最直接地作为速度-压力-梯度张量出现在雷诺应力方程（方程（7.178））中

\begin{sourceblock}
\includegraphics[page=416,pagebox=mediabox,trim={64.39bp 204.97bp 64.39bp 1695.79bp},clip,width=0.94\textwidth,max height=0.38\textheight]{Pope.pdf}
\end{sourceblock}
在这种情况下，没有必要区分耗散和伪耗散
这可以分解为压力（参见第 3 19 页的练习 7.24）：

/ 斧头。。输运项-aT，方程 (7 192) 和压力-应变率率张量
Rij 的迹为零 (Rii 2(p'V
O），因此
u / p ) = t e r n 没有出现在动能方程中：它用于在雷诺应力之间重新分配能量。

根据 Lumley (1975) 的观察，H i 分解为重分布项和传输项并不是唯一的。例如，另一种选择

\begin{sourceblock}
\includegraphics[page=417,pagebox=mediabox,trim={64.39bp 1155.85bp 64.39bp 585.93bp},clip,width=0.94\textwidth,max height=0.38\textheight]{Pope.pdf}
\end{sourceblock}
\begin{sourceblock}
\includegraphics[page=417,pagebox=mediabox,trim={64.39bp 1084.86bp 64.39bp 857.89bp},clip,width=0.94\textwidth,max height=0.38\textheight]{Pope.pdf}
\end{sourceblock}
其意义在于，由于压力传输而产生的动能的来源是

\begin{sourceblock}
\includegraphics[page=417,pagebox=mediabox,trim={64.39bp 907.89bp 64.39bp 1034.87bp},clip,width=0.94\textwidth,max height=0.38\textheight]{Pope.pdf}
\end{sourceblock}
在均匀湍流中，压力传输为零，并且所有再分配项都是等效的（例如，IIij = Rij = R (a . ) ）在检查此类流动（在本节和接下来的三个部分中）时，我们重点关注压力-应变率率张量 Rij。对于非均匀流，使用哪种分解只是方便的问题；在这种情况下，正如第 11.6 节中所讨论的，有理由支持用 R:;) 进行分解，方程(11 .5)。

再分布的重要性 值得回顾一下湍流边界层中的行为。在 (u2) 的预算中（第 316 页图 7.35），1111r 以大约两倍的速率转移能量。这两个汇（大约）由生产 Pll 平衡。在 (v2) 和 (w2) 的预算中（图 7.36 和 7.37）没有生产，但 1122a 和 1133a 资源大约平衡，因此能量从最大正应力（即

\textasciitilde{} 3 3。具有所有的能量产生）到较小的法向应力（没有能量产生）。在剪应力预算中（第 317 页图 7.38），产生 P12 近似由 -1112 平衡，耗散为

显然，除了产生和耗散之外，再分配也是雷诺应力平衡的一个决定性过程。因此，其建模至关重要，也是广泛研究的主题。

pf 的泊松方程 通过检查压力的泊松方程可以深入了解压力-应变率率张量（参见第 2.5 节）。该方程的雷诺分解（在练习 1 1.1 中执行）得出 p' 的泊松方程，其中包含两个源项：

\begin{sourceblock}
\includegraphics[page=418,pagebox=mediabox,trim={64.39bp 1323.83bp 64.39bp 594.92bp},clip,width=0.94\textwidth,max height=0.38\textheight]{Pope.pdf}
\end{sourceblock}
根据该方程，脉动压力场可以分解为三个贡献：

\begin{sourceblock}
\includegraphics[page=418,pagebox=mediabox,trim={64.39bp 1147.86bp 64.39bp 793.90bp},clip,width=0.94\textwidth,max height=0.38\textheight]{Pope.pdf}
\end{sourceblock}
快速压力p(')满足
慢压力 p@)s 满足

\begin{sourceblock}
\includegraphics[page=418,pagebox=mediabox,trim={64.39bp 785.90bp 64.39bp 1108.86bp},clip,width=0.94\textwidth,max height=0.38\textheight]{Pope.pdf}
\end{sourceblock}
谐波贡献 p(h) 满足拉普拉斯方程 V 2 p (h) 边界条件在 pch) 上指定，具体取决于 per) 和 p@) 上的边界条件，以便 pf 满足所需的边界条件。

快速压力之所以如此命名，是因为（与 p(") 不同，它立即响应平均速度梯度的变化。此外，在快速变形极限 ( e . Sk/\& co) 中，快速压力场 per) 具有前序效应，而 p(Qi 可以忽略不计（第 11.4 节）。对应于 p (0, p(S)a, nd p(h,)，压力-应变率率张量可以也可以分解为三个贡献，R!!) R (s) 和 R (h . )，具有明显的定义，例如，

\begin{sourceblock}
\includegraphics[page=418,pagebox=mediabox,trim={64.39bp 252.97bp 64.39bp 1640.79bp},clip,width=0.94\textwidth,max height=0.38\textheight]{Pope.pdf}
\end{sourceblock}
正如 Chou (1945) 的开创性著作所示，泊松方程（方程（2.48））的格林函数解可用于表达压力-速率-
R(irj) 均匀湍流，对

\begin{sourceblock}
\includegraphics[page=419,pagebox=mediabox,trim={64.39bp 1672.79bp 64.39bp 224.97bp},clip,width=0.94\textwidth,max height=0.38\textheight]{Pope.pdf}
\end{sourceblock}
其中四阶张量 M 由两点速度相关性 Rij(r) (ui(x)uj(x r)) 的 n 积分给出：

\begin{sourceblock}
\includegraphics[page=419,pagebox=mediabox,trim={64.39bp 1457.82bp 64.39bp 435.94bp},clip,width=0.94\textwidth,max height=0.38\textheight]{Pope.pdf}
\end{sourceblock}
（参见练习 11.2）。

从这些考虑中得出的有价值的结论是，Rij 具有三种性质不同的贡献。快速压力涉及平均速度梯度，并且在均匀湍流中 R!?与 (Wk) / a x e 成正比（方程（11.14））。在大多数情况下，慢压力-应变率率张量 R ( s ) 预计会很显着（快压力除外）

ij 失真）；事实上，在衰减的均匀各向异性湍流 中，R(-s ) 是三个贡献中唯一一个非零的。谐波分量 R(h)i 在均匀湍流中为零，并且仅在壁附近很重要：第 11.7.5 节对此进行了讨论。

表明压力的泊松方程（方程（2.42））可以交替
积极地写

因此表明平均压力满足

且波动压力满足式（1）： （11.9）。

考虑均匀湍流（其中平均速度梯度
dient a ( U k ) / a x e 是均匀的）。从方程式。 (2.48) 和 ( ) 表明
x 处快速压力与随机场之间的相关性
+(x) 由下式给出

\begin{sourceblock}
\includegraphics[page=419,pagebox=mediabox,trim={64.39bp 155.98bp 64.39bp 1703.79bp},clip,width=0.94\textwidth,max height=0.38\textheight]{Pope.pdf}
\end{sourceblock}
积分收敛所需的两点相关性。显示
对 R 的贡献是

\begin{sourceblock}
\includegraphics[page=420,pagebox=mediabox,trim={64.39bp 1613.80bp 64.39bp 282.96bp},clip,width=0.94\textwidth,max height=0.38\textheight]{Pope.pdf}
\end{sourceblock}
使用由 r y x 定义的分离向量，表明该等式
化可以根据两点速度相关性重写

Rij(r) (ui(x)uj(x r)) 为
因此验证方程。 （11.14）。

\section*{11.3 向各向同性回归模型}\phantomsection\addcontentsline{toc}{section}{11.3 向各向同性回归模型}
11.3.1 检查慢压力-应变率率张量的最简单情况是衰减均匀各向异性湍流。在这种情况下，没有生产或运输，并且 R( r . ) 和 R!:) 为零，因此精确的雷诺应力方程为

\subsection*{11.3.1 罗塔模型}\phantomsection\addcontentsline{toc}{subsection}{11.3.1 罗塔模型}
\begin{sourceblock}
\includegraphics[page=420,pagebox=mediabox,trim={64.39bp 739.91bp 64.39bp 1167.85bp},clip,width=0.94\textwidth,max height=0.38\textheight]{Pope.pdf}
\end{sourceblock}
由于 R(isj) 的迹为零，因此该项对湍动能没有影响。它的影响在于雷诺应力之间的能量分布，这可以通过归一化各向异性张量来检查
取 - 为各向同性 (Eq. l l )，bij 的演化方程为

\begin{sourceblock}
\includegraphics[page=420,pagebox=mediabox,trim={64.39bp 316.96bp 64.39bp 1599.80bp},clip,width=0.94\textwidth,max height=0.38\textheight]{Pope.pdf}
\end{sourceblock}
（参见练习 11 .3）。

很自然地假设湍流在衰减时具有各向异性减弱的趋势，并且实际上这种向各向同性回归的趋势是
提出模型

\begin{sourceblock}
\includegraphics[page=421,pagebox=mediabox,trim={64.39bp 1660.79bp 64.39bp 215.97bp},clip,width=0.94\textwidth,max height=0.38\textheight]{Pope.pdf}
\end{sourceblock}
其中 CR 是“罗塔常数”。将其代入方程。 (11 .23) 产量

\begin{sourceblock}
\includegraphics[page=421,pagebox=mediabox,trim={64.39bp 1502.81bp 64.39bp 398.95bp},clip,width=0.94\textwidth,max height=0.38\textheight]{Pope.pdf}
\end{sourceblock}
表明罗塔模型对应于各向同性的回归。

线性 显然需要大于 1 的 CR 值。

根据 b i j 的定义（方程（11.22））和雷诺应力
演化方程 (Eq. ( ) )，表明 b i j 的精确方程
在衰减湍流中是

\begin{sourceblock}
\includegraphics[page=421,pagebox=mediabox,trim={64.39bp 1074.86bp 64.39bp 813.90bp},clip,width=0.94\textwidth,max height=0.38\textheight]{Pope.pdf}
\end{sourceblock}
因此表明方程。 (11.23) 由假设的各向同性得出

表明，如果相反，被认为与 ( u i u j )t,h en 成正比
bij 的结果方程为

\begin{sourceblock}
\includegraphics[page=421,pagebox=mediabox,trim={64.39bp 772.90bp 64.39bp 1125.86bp},clip,width=0.94\textwidth,max height=0.38\textheight]{Pope.pdf}
\end{sourceblock}
表明，如果在该方程中使用 Rotta 模型，则结果为
与等式相同(11.25) 但用 CR 代替 CR 1。

11.3.2 为了检查 Rotta 模型（和其他向各向同性回归模型）的行为，能够比各向异性张量 b 的六个分量 bij 更简单地表征雷诺应力各向异性是有用的。事实上，正如现在所描述的，各向异性的状态可以仅由两个变量（由 q 和 q 表示）来表征，这允许简单的图形表示。

\subsection*{11.3.2 雷诺应力各向异性的表征}\phantomsection\addcontentsline{toc}{subsection}{11.3.2 雷诺应力各向异性的表征}
各向异性张量的迹为零 (bii 0)，因此它只有两个独立的不变量（参见附录 B）。把这两个独立的不变量看作是方便吗？数量和 q 定义为

\begin{sourceblock}
\includegraphics[page=421,pagebox=mediabox,trim={64.39bp 119.39bp 64.39bp 1844.77bp},clip,width=0.94\textwidth,max height=0.38\textheight]{Pope.pdf}
\par\vspace{0.45\baselineskip}
\small 表 11.1。雷诺应力张量的特殊状态，以不变量 € 和 n、b 的特征值以及练习 11.8 中定义的雷诺应力椭球的形状表示（另请参见图 11.1 和练习 11.4）
\end{sourceblock}
\clearpage
\begin{sourceblock}
\includegraphics[page=422,pagebox=mediabox,trim={64.39bp 489.94bp 64.39bp 306.96bp},clip,width=\textwidth,max height=0.72\textheight]{Pope.pdf}
\end{sourceblock}
在任何湍流中的任何点和时间，[ 和 y 都可以根据雷诺应力确定，并将结果绘制为 [-y 平面上的一个点。雷诺应力张量的一些特殊状态对应于该平面中的特定点和曲线。这些状态在练习 11.4 中确定，总结在表 11.1 中，并在图 11.1 中绘制，被称为（即使一侧不是 t a i g h t ）。 ～
拉姆利三角

湍流中可能出现的每个雷诺应力（ e，即
事实上，Lumley (1978) 使用 和 -IIb 作为坐标（而不是 ( 和 )，在这种情况下，没有

\begin{sourceblock}
\includegraphics[page=423,pagebox=mediabox,trim={64.39bp 1259.84bp 64.39bp 109.44bp},clip,width=\textwidth,max height=0.72\textheight]{Pope.pdf}
\par\vspace{0.45\baselineskip}
\small 图 11.1。雷诺应力各向异性张量的不变量 € 和 y 平面上的 Lumley 三角形。线和顶点对应于特殊状态（见表11.1）。圆圈：来自槽道流的 DNS（Kim 等人，1987）。正方形：来自湍流混合层的实验（Bell 和 Mehta 1990）。 1C，单组分； 2C，双组分。
\end{sourceblock}
可实现）对应于 Lumley 三角形中的一个点。外部点对应于具有负或复特征值的不可实现的雷诺应力。

图 11.1 还显示了在湍流槽道流和湍流混合层中获得的 (5, \textasciitilde{}o) 值。在槽道流中非常靠近壁面（例如 yf < 5），湍流本质上是双分量的，(u2) 远小于 (u2) 和 (w2)。各向异性（由 y 测量）在 yf 7 处达到峰值。在通道的其余部分，雷诺应力接近正轴对称。在湍流混合层的中心，雷诺应力接近轴对称，
正值，各向异性略低于槽道流的对数区域。在混合层边缘，雷诺应力接近负轴对称。

表明 和 y（方程（11.28）和（11.29））与
b 的特征值（方程 (11 .30)）

\begin{sourceblock}
\includegraphics[page=423,pagebox=mediabox,trim={64.39bp 233.97bp 64.39bp 1694.79bp},clip,width=0.94\textwidth,max height=0.38\textheight]{Pope.pdf}
\end{sourceblock}
\begin{sourceblock}
\includegraphics[page=423,pagebox=mediabox,trim={64.39bp 159.98bp 64.39bp 1779.78bp},clip,width=0.94\textwidth,max height=0.38\textheight]{Pope.pdf}
\end{sourceblock}
张量（主轴上的法向应力）+验证条目

表 11.1 列出了以下雷诺应力状态：

(b) 双组分，轴对称； (u,) (f i ) (u3) 0 ;和
(c) 单组分； 0。

显示轴对称雷诺应力 ( e , ( 6 ( 6 11, 12)

从而验证表 11 .1 中给出的轴对称状态。

数量 F 被定义为归一化的行列式
雷诺应力张量

\begin{sourceblock}
\includegraphics[page=424,pagebox=mediabox,trim={64.39bp 1099.86bp 64.39bp 792.90bp},clip,width=0.94\textwidth,max height=0.38\textheight]{Pope.pdf}
\end{sourceblock}
因此它在各向同性湍流中是统一的。推导以下等价
F 的 alent 表达式

\begin{sourceblock}
\includegraphics[page=424,pagebox=mediabox,trim={64.39bp 715.91bp 64.39bp 1225.85bp},clip,width=0.94\textwidth,max height=0.38\textheight]{Pope.pdf}
\end{sourceblock}
观察到 F 在二分量湍流中为零（例如，(6:) 0），
从而验证二元曲线的方程

图11.1是

\begin{sourceblock}
\includegraphics[page=424,pagebox=mediabox,trim={64.39bp 479.94bp 64.39bp 1462.82bp},clip,width=0.94\textwidth,max height=0.38\textheight]{Pope.pdf}
\end{sourceblock}
对于均匀剪切流 (i)，来自实验
数据（参见第 157 页表 5.4），以及 (ii) 根据 k-E 模型。

找到图 11.1 上的相应点。

证明，根据 Rotta 模型（方程 (11 . 5 ) ），不变量

\begin{sourceblock}
\includegraphics[page=424,pagebox=mediabox,trim={64.39bp 119.39bp 64.39bp 1829.77bp},clip,width=0.94\textwidth,max height=0.38\textheight]{Pope.pdf}
\end{sourceblock}
因此表明 4 和 q（方程 (11 .28) 和 (11 .29)）通过以下方式演化

\begin{sourceblock}
\includegraphics[page=425,pagebox=mediabox,trim={64.39bp 1689.79bp 64.39bp 220.97bp},clip,width=0.94\textwidth,max height=0.38\textheight]{Pope.pdf}
\end{sourceblock}
\begin{sourceblock}
\includegraphics[page=425,pagebox=mediabox,trim={64.39bp 1575.80bp 64.39bp 332.96bp},clip,width=0.94\textwidth,max height=0.38\textheight]{Pope.pdf}
\end{sourceblock}
比率 q 是如何演变的？显示 5-q 中的轨迹
由方程生成的平面。 (11 .38) 和 (11 .39) 是有向的直线
朝向原点（CR > 1）。

表明

\begin{sourceblock}
\includegraphics[page=425,pagebox=mediabox,trim={64.39bp 1266.84bp 64.39bp 624.92bp},clip,width=0.94\textwidth,max height=0.38\textheight]{Pope.pdf}
\end{sourceblock}
分别是雷诺数的必要条件和充分条件

k-E 模型给出的应力可实现。

在湍流中的给定点，对于给定的 0，考虑
数量

\begin{sourceblock}
\includegraphics[page=425,pagebox=mediabox,trim={64.39bp 982.88bp 64.39bp 959.88bp},clip,width=0.94\textwidth,max height=0.38\textheight]{Pope.pdf}
\end{sourceblock}
其中 Cyl 是雷诺应力倒数的 i-j 分量
张量，Cij(u\textasciitilde{}u,)。验证方程

\begin{sourceblock}
\includegraphics[page=425,pagebox=mediabox,trim={64.39bp 789.90bp 64.39bp 1158.85bp},clip,width=0.94\textwidth,max height=0.38\textheight]{Pope.pdf}
\end{sourceblock}
在速度空间中定义一个椭球，中心位于 (U)，主函数
轴与 (uiuj) 的轴以及与 (uiuj) 的半长对齐
主轴为 crfiii)，其中 u r2 (1,2,3) 是特征值
(uiuj).验证雷诺应力椭球的形状
如果 U 是联合正态分布（参见方程（3.1 1 6 ）），则
由 E(V,a ) 1 给出的椭球表面是一个表面
恒定的概率密度。证明U说谎的概率
椭球内是

\begin{sourceblock}
\includegraphics[page=425,pagebox=mediabox,trim={64.39bp 246.97bp 64.39bp 1655.79bp},clip,width=0.94\textwidth,max height=0.38\textheight]{Pope.pdf}
\end{sourceblock}
其中 Q 是 y -E q 定义的不完全伽玛函数。 (12.202)。

（这是具有三个自由度的分布。）

\begin{sourceblock}
\includegraphics[page=426,pagebox=mediabox,trim={64.39bp 1090.86bp 64.39bp 109.44bp},clip,width=\textwidth,max height=0.72\textheight]{Pope.pdf}
\par\vspace{0.45\baselineskip}
\small 图 11.2。由 Sarkar 和 Speziale (1990) 模型给出的 €-9 平面上的轨迹（方程（11.51）和（11.57））。
\end{sourceblock}
\subsection*{11.3.3 非线性向各向同性回归模型}\phantomsection\addcontentsline{toc}{subsection}{11.3.3 非线性向各向同性回归模型}
当均匀各向异性湍流衰减时，其演化状态对应于平面上的轨迹。根据罗塔模型，这些轨迹是指向原点的直线（参见练习 11.6）；和归一化各向同性返回率

\begin{sourceblock}
\includegraphics[page=426,pagebox=mediabox,trim={64.39bp 657.92bp 64.39bp 1249.84bp},clip,width=0.94\textwidth,max height=0.38\textheight]{Pope.pdf}
\end{sourceblock}
独立于 和 y。对现有经验的仔细检查 \textasciitilde{} \textasciitilde{} \textasciitilde{}dat taa (lse e Chung and Kim (1995)) 揭示了不同的行为。主要观察结果如下。

(i) 向各向同性回归的轨迹不是直线
更像是图 11.2 所示：有一种趋势
轴对称状态为正。

(ii) 归一化各向同性返回率不独立于和
y。特别是，回报率通常较大（例如，通过一个因子
两个）对于负数的影响比对于正数的影响要大。

(iii) 在网格湍流实验的雷诺数范围内以及
DNS，回报率很大程度上取决于雷诺数
Rotta 模型的缺点促使人们考虑更通用的慢压力-应变率率张量 Rij 模型（例如 Shih 和 Lumley (1985)、Sarkar 和 Speziale (1990) 以及 Chung 和 Kim (1995)）。这些模型的发展在一个相对简单的环境中说明了雷诺应力建模中使用的几个一般原理。因此，详细说明所涉及的步骤很有用。

在应用于衰减各向异性湍流的雷诺应力模型中，湍流的特征为 (uiuj) 并且唯一相关的材料属性是 v。因此，一般模型可以写为（根据未指定的张量函数 '1) ）

\begin{sourceblock}
\includegraphics[page=427,pagebox=mediabox,trim={64.39bp 1325.83bp 64.39bp 601.92bp},clip,width=0.94\textwidth,max height=0.38\textheight]{Pope.pdf}
\end{sourceblock}
在衰减各向异性湍流的风洞实验中，平均速度 (U) 也以雷诺应力闭合表示，因此也可以考虑 (U) 的依赖性。然而，原则
伽利略不变性的 IJ 排除了这种依赖性：在所有惯性系中都是相同的，而 (U) 显然不是（参见第 2.9 节）。尽管在符号中没有明确表示，但也应该理解的是，方程。 (11.45) 包含局部性假设：假设在位置 x 和时间 t
lj 依赖 (uiuj) 并且在相同的位置和时间。

尺寸分析的简单应用允许方程。 (11 .45) 以无量纲形式重新表达为

\begin{sourceblock}
\includegraphics[page=427,pagebox=mediabox,trim={64.39bp 762.90bp 64.39bp 1131.86bp},clip,width=0.94\textwidth,max height=0.38\textheight]{Pope.pdf}
\end{sourceblock}
其中 是无量纲张量函数，雷诺数为 ReL k2/(\&v)。可以注意到，Rotta 模型（方程（11 .24））对应于规范

\begin{sourceblock}
\includegraphics[page=427,pagebox=mediabox,trim={64.39bp 549.93bp 64.39bp 1383.83bp},clip,width=0.94\textwidth,max height=0.38\textheight]{Pope.pdf}
\end{sourceblock}
量 R ( i s i ) 是二阶张量，并且（正如符号已经暗示的那样）只有模型也是二阶张量才是合理的
tj 张量。这似乎太明显了，无需提及。然而，重要的是要认识到，坚持张量正确的模型可以保证模型隐含的行为独立于坐标系的选择。未以张量形式编写的模型（例如，为特定组件编写的模型，或者涉及伪向量或交替符号的模型）

Eijk）不保证坐标系不变性，应该用
菲；'怀疑。此外，为了重现 R (s' /E) 的属性，模型
下一个问题是张量函数的表示。超越方程式。 ( l )，用张量 b 和标量 ReL 表示的张量最通用的表示是什么？形式上，我们可以写

\begin{sourceblock}
\includegraphics[page=428,pagebox=mediabox,trim={64.39bp 1608.80bp 64.39bp 299.96bp},clip,width=0.94\textwidth,max height=0.38\textheight]{Pope.pdf}
\end{sourceblock}
其中 , n 1,2,3,... 是所有对称的偏张量，它们可以： 由 b 形成，并且标量系数取决于 ReL 和 b 的不变量。可以由 b 形成的张量的唯一形式是幂
?;jn) b.因此，偏张量可以被视为

\begin{sourceblock}
\includegraphics[page=428,pagebox=mediabox,trim={64.39bp 1335.83bp 64.39bp 597.92bp},clip,width=0.94\textwidth,max height=0.38\textheight]{Pope.pdf}
\end{sourceblock}
然而，鉴于凯莱-汉密尔顿定理（参见附录 B），方程 1 中仅包含 I('a)nd 7( 2) 就足够了。 (11.48)，因为所有其他张量 ( e 7 , n > 2) 都可以表示为这两者的线性组合。因此方程。 (11.48) 可以简化为

\begin{sourceblock}
\includegraphics[page=428,pagebox=mediabox,trim={64.39bp 1065.87bp 64.39bp 858.89bp},clip,width=0.94\textwidth,max height=0.38\textheight]{Pope.pdf}
\end{sourceblock}
该方程表明 (2)（张量的张量函数）是
ij 由两个已知张量乘以标量系数给出。因此，建模任务已简化为确定这两个标量系数（f 和 f）。显然，使用凯莱-汉密尔顿定理获得此类张量表示的过程（由 Robertson（1940）提出）在湍流建模中具有重要价值。

在本例中，程序很简单。在湍流研究中出现的其他情况远非简单。对于此类情况，力学文献中有大量工作描述了一系列张量表示定理（例如 Spencer (197 I)、Pennisi 和 Trovato (1987)、Pennisi (1992) 以及其中的参考文献）。

到目前为止所采取的步骤已将一般模型（方程（11.45））简化为更具体的形式（方程（1 1 . 0）），这（与方程（11.23））导致bij演化的建模方程

\begin{sourceblock}
\includegraphics[page=428,pagebox=mediabox,trim={64.39bp 338.96bp 64.39bp 1569.80bp},clip,width=0.94\textwidth,max height=0.38\textheight]{Pope.pdf}
\end{sourceblock}
特定型号对应于特定规格和 f
( ) (2) 作为 5 和 ReL 的函数。 Rotta 模型由 y f --2CR, f 0 给出。

对可实现性的考虑对规范施加了限制

FC2）。 f c ' )a 和 对于给定的规范和给定的初始条件

\begin{sourceblock}
\includegraphics[page=429,pagebox=mediabox,trim={64.39bp 1095.86bp 64.39bp 109.44bp},clip,width=\textwidth,max height=0.72\textheight]{Pope.pdf}
\par\vspace{0.45\baselineskip}
\small 图 11.3。拉姆利三角形显示了三种类型的轨迹：(a) 违反了可实现性； (b) 满足弱可实现性； (c) 满足较强的可实现性。 （请注意，其他类型的轨迹也是可能的。）
\end{sourceblock}
\begin{sourceblock}
\includegraphics[page=429,pagebox=mediabox,trim={64.39bp 590.93bp 64.39bp 1305.84bp},clip,width=0.94\textwidth,max height=0.38\textheight]{Pope.pdf}
\end{sourceblock}
\begin{sourceblock}
\includegraphics[page=429,pagebox=mediabox,trim={64.39bp 279.96bp 64.39bp 1588.80bp},clip,width=0.94\textwidth,max height=0.38\textheight]{Pope.pdf}
\end{sourceblock}
图 11.3 中所示的三个轨迹通过二分量线 (F 0) 上 d F / d t 的行为来区分。轨迹 (a) 有

\begin{sourceblock}
\includegraphics[page=430,pagebox=mediabox,trim={64.39bp 1610.80bp 64.39bp 275.96bp},clip,width=0.94\textwidth,max height=0.38\textheight]{Pope.pdf}
\end{sourceblock}
导致负F和不可实现性。轨迹 (b) 的特点是

\begin{sourceblock}
\includegraphics[page=430,pagebox=mediabox,trim={64.39bp 1380.83bp 64.39bp 536.93bp},clip,width=0.94\textwidth,max height=0.38\textheight]{Pope.pdf}
\end{sourceblock}
显然，如果沿二分量线始终满足此条件，则任何轨迹都不会穿过该线，并且可保证可实现性。4 轨迹 (c) 与二分量线密切相关，其特征为
Schumann (1977) 和 Lumley (1978) 将可实现性的概念引入雷诺应力模型中，并认为模型应该产生 (c) 类型的轨迹。许多模型开发人员都遵循这个想法，并相应地构建了满足式（1）的模型。 1 1 5 6 ) 据说这样的模型具有很强的可实现性。另一方面，Pope（1985）认为（b）类型的轨迹不违反可实现性，因此（一般来说）强可实现性条件太强了。具有类型 (b) 或 (c) 行为的模型，满足方程(11 .55) 或等式。 1 1.56，据说满足弱可实现性。

对于非线性向各向同性回归模型（方程 1 1 . 0 ) ），Sarkar 描述了弱可实现性的约束和强加的约束
f(') f(2) Speziale (1990)，以及 Chung 和 Kim (1995) 的强可实现性。此处针对由向各向同性回归模型给出的雷诺应力的特定情况描述了可实现性。然而，应该意识到，可实现性是一个通用概念，适用于所有模型和所有统计数据。

向各向同性回归模型开发的最后一步是指定模型系数 f q、ReL) 和 f ReL)。对于
(I)(<, (2)(<，这一步，实验数据在建议适当的规范和测试结果模型方面都起着至关重要的作用。在 Sarkar 和

\begin{sourceblock}
\includegraphics[page=431,pagebox=mediabox,trim={64.39bp 1193.85bp 64.39bp 109.44bp},clip,width=\textwidth,max height=0.72\textheight]{Pope.pdf}
\par\vspace{0.45\baselineskip}
\small 图 11.4。两个实验（或 DNS）平面上的轨迹（A 和 B）草图，其中初始光谱不同，但 b 的初始值相同。雷诺应力模型从初始点 O 产生独特的轨迹。
\end{sourceblock}
Speziale (1990)，规格很简单

\begin{sourceblock}
\includegraphics[page=431,pagebox=mediabox,trim={64.39bp 894.89bp 64.39bp 1047.87bp},clip,width=0.94\textwidth,max height=0.38\textheight]{Pope.pdf}
\end{sourceblock}
而在 Chung 和 Kim (1995) 的模型中，规范是 \&、7 和 Re; 的复杂函数。前一个模型生成的轨迹如第 398 页的图 11.2 所示。

我们以警告的方式结束本节。开发通用模型的理性过程中看似完美的逻辑可能会导致对模型描述当前现象的能力的高估。考虑衰减均匀各向异性湍流的两个直接数值模拟。选择初始条件以使 (uju;)、\& 和 v 在两次模拟中相同，但光谱却截然不同。不可避免的是，两个模拟中湍流的后续演变是不同的，导致飞机上的轨迹不同，如图 11.4 所示。然而，无论如何指定 和 f®，一般模型（方程（11.50））不可避免地会从初始条件产生独特的轨迹，因此无法描述两个模拟的不同行为。当然，原因是湍流状态仅部分地由以下表示的变量来表征：

证明 b 的三个主要不变量（见方程（J3.3 1）-（B.33））

\begin{sourceblock}
\includegraphics[page=432,pagebox=mediabox,trim={64.39bp 1653.79bp 64.39bp 252.97bp},clip,width=0.94\textwidth,max height=0.38\textheight]{Pope.pdf}
\end{sourceblock}
\begin{sourceblock}
\includegraphics[page=432,pagebox=mediabox,trim={64.39bp 1572.80bp 64.39bp 373.95bp},clip,width=0.94\textwidth,max height=0.38\textheight]{Pope.pdf}
\end{sourceblock}
\begin{sourceblock}
\includegraphics[page=432,pagebox=mediabox,trim={64.39bp 1493.81bp 64.39bp 449.94bp},clip,width=0.94\textwidth,max height=0.38\textheight]{Pope.pdf}
\end{sourceblock}
使用凯莱-汉密尔顿定理（方程（B.39））表明

\begin{sourceblock}
\includegraphics[page=432,pagebox=mediabox,trim={64.39bp 1355.83bp 64.39bp 583.93bp},clip,width=0.94\textwidth,max height=0.38\textheight]{Pope.pdf}
\end{sourceblock}
显示方程。 (11 .51) 意味着以下演化方程

\begin{sourceblock}
\includegraphics[page=432,pagebox=mediabox,trim={64.39bp 1171.85bp 64.39bp 764.90bp},clip,width=0.94\textwidth,max height=0.38\textheight]{Pope.pdf}
\end{sourceblock}
\begin{sourceblock}
\includegraphics[page=432,pagebox=mediabox,trim={64.39bp 1048.87bp 64.39bp 858.89bp},clip,width=0.94\textwidth,max height=0.38\textheight]{Pope.pdf}
\end{sourceblock}
因此表明 F（方程（11 .35））演化为

\begin{sourceblock}
\includegraphics[page=432,pagebox=mediabox,trim={64.39bp 894.89bp 64.39bp 1013.87bp},clip,width=0.94\textwidth,max height=0.38\textheight]{Pope.pdf}
\end{sourceblock}
\section*{11.4 快速畸变理论}\phantomsection\addcontentsline{toc}{section}{11.4 快速畸变理论}
均匀湍流可以受到时间相关的均匀平均速度梯度的影响，其大小可以通过以下方式表征

\begin{sourceblock}
\includegraphics[page=432,pagebox=mediabox,trim={64.39bp 552.93bp 64.39bp 1386.83bp},clip,width=0.94\textwidth,max height=0.38\textheight]{Pope.pdf}
\end{sourceblock}
（参见第 158 页的练习 5.41）？如上所述，在湍流剪切流中，湍流与平均剪切时间尺度比 zS = S \textasciitilde{} /E 通常在 3-6 范围内。相反，在本节中，我们考虑快速畸变极限，其中 S \textasciitilde{} / isE 任意大。在这种极限情况下，湍流的演变可以通过快速畸变理论（RDT）准确地描述。 RDT 提供了一些有用的见解，特别是关于快速压力-应变率率张量，下一节将讨论其模型。

显然，固体旋转需要不同的表征，例如 (

\subsection*{11.4.1 快速畸变方程}\phantomsection\addcontentsline{toc}{subsection}{11.4.1 快速畸变方程}
在均匀湍流中，脉动速度演变为（方程（5.138））

\begin{sourceblock}
\includegraphics[page=433,pagebox=mediabox,trim={64.39bp 1630.79bp 64.39bp 292.96bp},clip,width=0.94\textwidth,max height=0.38\textheight]{Pope.pdf}
\end{sourceblock}
p' p(r) 的泊松方程为 (Eq. (11.9))

\begin{sourceblock}
\includegraphics[page=433,pagebox=mediabox,trim={64.39bp 1464.82bp 64.39bp 450.94bp},clip,width=0.94\textwidth,max height=0.38\textheight]{Pope.pdf}
\end{sourceblock}
在这两个方程的右侧，第一项表示湍流场 u 和平均速度梯度 之间的相互作用；而第二项代表湍流-湍流相互作用。给定时间 t 处的湍流场 u(x, t)，湍流-湍流项已确定，并且与 a(Ui)/dxj 无关。另一方面，平均速度梯度项与 S 成线性比例。因此，显然，在快速变形极限 (e.Sm) 中，与 S 成比例的项占主导地位，相比之下，所有其他项都可以忽略不计。因此，在这个限制下，方程。 (11.66) 和 (11.67) 简化为快速畸变方程

\begin{sourceblock}
\includegraphics[page=433,pagebox=mediabox,trim={64.39bp 951.88bp 64.39bp 940.88bp},clip,width=0.94\textwidth,max height=0.38\textheight]{Pope.pdf}
\end{sourceblock}
\begin{sourceblock}
\includegraphics[page=433,pagebox=mediabox,trim={64.39bp 834.89bp 64.39bp 1057.87bp},clip,width=0.94\textwidth,max height=0.38\textheight]{Pope.pdf}
\end{sourceblock}
由平均速度梯度引起的变形可以用速率 S(t),t 量（从时间 0 到 t）来考虑

\begin{sourceblock}
\includegraphics[page=433,pagebox=mediabox,trim={64.39bp 629.92bp 64.39bp 1269.84bp},clip,width=0.94\textwidth,max height=0.38\textheight]{Pope.pdf}
\end{sourceblock}
和变形的几何形状

\begin{sourceblock}
\includegraphics[page=433,pagebox=mediabox,trim={64.39bp 466.94bp 64.39bp 1428.82bp},clip,width=0.94\textwidth,max height=0.38\textheight]{Pope.pdf}
\end{sourceblock}
请注意，s p 和 Gij 都是无量纲量。快速畸变理论的一个有趣特征是，湍流场取决于几何结构和畸变量，但与速率 S(t) 无关

表明湍流黏度假说对于快速扭曲来说在性质上是不正确的（Crow 1968）。为了显示快速畸变方程的这一性质，我们使用 s 代替 t 作为自变量，并定义

\begin{sourceblock}
\includegraphics[page=433,pagebox=mediabox,trim={64.39bp 119.39bp 64.39bp 1830.77bp},clip,width=0.94\textwidth,max height=0.38\textheight]{Pope.pdf}
\end{sourceblock}
然后（除以 5 时）快速畸变方程（方程（11.68）和

\begin{sourceblock}
\includegraphics[page=434,pagebox=mediabox,trim={64.39bp 1660.79bp 64.39bp 199.97bp},clip,width=0.94\textwidth,max height=0.38\textheight]{Pope.pdf}
\end{sourceblock}
\begin{sourceblock}
\includegraphics[page=434,pagebox=mediabox,trim={64.39bp 1547.80bp 64.39bp 355.95bp},clip,width=0.94\textwidth,max height=0.38\textheight]{Pope.pdf}
\end{sourceblock}
给定初始湍流场 u(x,0) 和畸变几何形状，这些方程可以在 s 中向前积分，以确定随后的湍流场，作为畸变量 s 的函数（独立于 S(t) 和 t）。 （做出这一观察后，我们回到更熟悉的方程（11.68）和（11.69）变量。）

为了在快速畸变方程的分析上取得进展，有必要绕过或求解 p 的泊松方程。在 RDT 的第一个著作中，Prandtl (1933) 和 Taylor (1935b) 考虑了湍流涡量方程（方程 (11.68) 的旋度），从而消除了 pTM。对于无旋平均畸变 (Q,; = 0)，该涡度方程为

\begin{sourceblock}
\includegraphics[page=434,pagebox=mediabox,trim={64.39bp 992.87bp 64.39bp 933.88bp},clip,width=0.94\textwidth,max height=0.38\textheight]{Pope.pdf}
\end{sourceblock}
因此，（脉动涡量场的）涡线随平均速度场移动并被平均速度场拉伸；和（如无粘流）涡度 |@|与拉伸量成比例增加。

在轴对称收缩中，S,, > 9 且 = \$3 = -48,； （参见第 360 和 361 页上的图 10.1 和 10.2），涡流线向 x 轴倾斜，并在 x 方向上拉伸，从而导致 (u3) 和 (u3) 相对于图 10.2 中观察到的增加，

\subsection*{11.4.2 傅立叶模的演变}\phantomsection\addcontentsline{toc}{subsection}{11.4.2 傅立叶模的演变}
Taylor 和 Batchelor (1949) 引入的 RDT 替代方法（并由 Batchelor 和 Proudman (1954) 和 Craya (1958) 进一步发展）取代了涡度方程，即考虑傅立叶模态。初始湍流速度场 u(x,0) 可以表示为傅里叶模态之和（参见第 6.4 节）。快速畸变方程（方程（11.68）和（11.69））在 x 和 p" 中是线性的，因此每个傅立叶模式独立演化。因此，我们首先考虑单一模式，最初指定为

\begin{sourceblock}
\includegraphics[page=434,pagebox=mediabox,trim={64.39bp 119.39bp 64.39bp 1833.77bp},clip,width=0.94\textwidth,max height=0.38\textheight]{Pope.pdf}
\end{sourceblock}
\begin{sourceblock}
\includegraphics[page=435,pagebox=mediabox,trim={64.39bp 1140.86bp 64.39bp 109.44bp},clip,width=\textwidth,max height=0.72\textheight]{Pope.pdf}
\par\vspace{0.45\baselineskip}
\small 图 11.5。场 t) 的波峰由 = 0 演化： (a) 初始条件 geTM*, = K3 > 0, x3 = 0； (b) 平面应变后 = > 0)； (c) 剪切后 > 0。
\end{sourceblock}
ii(0) 0 以满足连续性方程， （6.128）。 （应该理解，还有共轭模 ii'(0)e -ix0*x
，以便这两种模式的总和产生真实的速度场。）

在处理快速畸变方程之前，为了发展必要的概念，我们首先考虑一个更简单的方程。具体来说，令 \$(x,t) 为固定于平均流的标量场，即

\begin{sourceblock}
\includegraphics[page=435,pagebox=mediabox,trim={64.39bp 632.92bp 64.39bp 1268.84bp},clip,width=0.94\textwidth,max height=0.38\textheight]{Pope.pdf}
\end{sourceblock}
最初是傅里叶模式

\begin{sourceblock}
\includegraphics[page=435,pagebox=mediabox,trim={64.39bp 510.94bp 64.39bp 1426.82bp},clip,width=0.94\textwidth,max height=0.38\textheight]{Pope.pdf}
\end{sourceblock}
对于某些指定的 和 由于不随平均流变化，因此 的最大值始终等于。

图 11.5(a) 显示 4(x,O) 的波峰（其中 ）代表 0，
0.这些是平行平面（在图11S(a)中显示为线）。图 11.5 还显示了 (b) 平面应变和 (c) 剪切后场 \$(x, t) 的波峰。波峰保持平行，但它们的间距和方向因变形而改变。

根据这些观察结果，可以假设该场

\begin{sourceblock}
\includegraphics[page=436,pagebox=mediabox,trim={64.39bp 1739.78bp 64.39bp 159.98bp},clip,width=0.94\textwidth,max height=0.38\textheight]{Pope.pdf}
\end{sourceblock}
对于一些 R(t)。很容易证明（练习 11.11）这确实是一个
b+/bt 解为 0，前提是波数按照以下规律演化

\begin{sourceblock}
\includegraphics[page=436,pagebox=mediabox,trim={64.39bp 1531.81bp 64.39bp 370.95bp},clip,width=0.94\textwidth,max height=0.38\textheight]{Pope.pdf}
\end{sourceblock}
with:h 初始条件 2(O)

对于快速畸变方程，获得了相同类型的解，

\begin{sourceblock}
\includegraphics[page=436,pagebox=mediabox,trim={64.39bp 1335.83bp 64.39bp 561.93bp},clip,width=0.94\textwidth,max height=0.38\textheight]{Pope.pdf}
\end{sourceblock}
\begin{sourceblock}
\includegraphics[page=436,pagebox=mediabox,trim={64.39bp 1261.84bp 64.39bp 685.91bp},clip,width=0.94\textwidth,max height=0.38\textheight]{Pope.pdf}
\end{sourceblock}
傅里叶系数 k(t) 演变为

\begin{sourceblock}
\includegraphics[page=436,pagebox=mediabox,trim={64.39bp 1098.86bp 64.39bp 807.90bp},clip,width=0.94\textwidth,max height=0.38\textheight]{Pope.pdf}
\end{sourceblock}
（参见练习 11.13）。因此，在快速畸变极限下，湍流速度场（方程（11 .81））的每个傅里叶模式根据 R(t) 和 ii(t) 的一对常微分方程独立演化， （11.80）和（11.83）。

波数向量 R(t) 可以分解为其幅度 2(t) (iZ-(t)l 及其方向

\begin{sourceblock}
\includegraphics[page=436,pagebox=mediabox,trim={64.39bp 721.91bp 64.39bp 1174.85bp},clip,width=0.94\textwidth,max height=0.38\textheight]{Pope.pdf}
\end{sourceblock}
称为单位波矢。从方程式中可以立即看出。 (11.83) 傅立叶系数 ii(t) 的演化与波数的大小无关：该方程可以用单位波矢重写为
此外，它还可以从等式得出。 (11.80) 单位波数演化为

\begin{sourceblock}
\includegraphics[page=436,pagebox=mediabox,trim={64.39bp 288.96bp 64.39bp 1611.80bp},clip,width=0.94\textwidth,max height=0.38\textheight]{Pope.pdf}
\end{sourceblock}
这对方程完全确定了 k(t) 的演化，显然，它不依赖于波数的大小 k(t)。

单位向量可以表示为单位球面上的一个点，并且

\begin{sourceblock}
\includegraphics[page=437,pagebox=mediabox,trim={64.39bp 885.89bp 64.39bp 109.44bp},clip,width=\textwidth,max height=0.72\textheight]{Pope.pdf}
\par\vspace{0.45\baselineskip}
\small 图 11.6。根据 (a) 轴对称收缩、(b) 轴对称膨胀、(c) 平面应变和 (d) 剪切的随机初始条件，单位球面上单位波矢 @(t) 的轨迹。 @，方向为水平，@方向为垂直，@；方向是进入页面。这些符号标记了变形后轨迹的末端。
\end{sourceblock}
作为单位球面上的轨迹。表 11.2（第 415 页）中定义的各种平均速度梯度的轨迹（来自不同的初始条件）如图 11 .6 所示。

傅立叶系数 k(t) 与 k t ) 正交，因此与 \$(t) 正交（因为连续性）。因此，如图 11.7 所示，k(t) 可以表示为单位球体在 Yt) 处的切平面中的向量。当点 2(t) 在单位球面上移动时，系数 i ( t ) 根据方程 (1) 演化。 (11 .85)，并且保持在局部切平面内。

练习 11 .11 考虑具有均匀平均速度梯度的均匀湍流

\begin{sourceblock}
\includegraphics[page=438,pagebox=mediabox,trim={64.39bp 1339.83bp 64.39bp 109.44bp},clip,width=\textwidth,max height=0.72\textheight]{Pope.pdf}
\par\vspace{0.45\baselineskip}
\small 图 11.7。显示单位波矢 @t) 的单位球面草图。速度 \#(t) 的傅里叶分量与 @t) 正交，因此它位于 é(t) 处单位球面的切平面内。
\end{sourceblock}
速度场为

证明 \$(x, t) 由方程给出。 (11 .79) 演变为

\begin{sourceblock}
\includegraphics[page=438,pagebox=mediabox,trim={64.39bp 786.90bp 64.39bp 1110.86bp},clip,width=0.94\textwidth,max height=0.38\textheight]{Pope.pdf}
\end{sourceblock}
因此表明，如果 k(t) 按等式演化，则 D \$ / D \textasciitilde{} 为零。 (11480)4

表明，如果 \$ 演化为 = 0，则其梯度 g 演化为
评论一下这个结果和方程之间的联系。 (11 .80)*
由方程给出 (r) 和 k(t)，(11 . 8 ) (11 8 2 和 (18 ) 表明
p(') 的泊松方程（方程 (1 1.69)）可以写成
傅里叶系数为

\begin{sourceblock}
\includegraphics[page=438,pagebox=mediabox,trim={64.39bp 270.97bp 64.39bp 1631.79bp},clip,width=0.94\textwidth,max height=0.38\textheight]{Pope.pdf}
\end{sourceblock}
因此表明快速压力梯度为

\begin{sourceblock}
\includegraphics[page=438,pagebox=mediabox,trim={64.39bp 119.39bp 64.39bp 1794.77bp},clip,width=0.94\textwidth,max height=0.38\textheight]{Pope.pdf}
\end{sourceblock}
\begin{sourceblock}
\includegraphics[page=439,pagebox=mediabox,trim={64.39bp 1704.79bp 64.39bp 158.98bp},clip,width=0.94\textwidth,max height=0.38\textheight]{Pope.pdf}
\end{sourceblock}
因此，如果 ii(t) 演化，则满足快速畸变方程

\begin{sourceblock}
\includegraphics[page=439,pagebox=mediabox,trim={64.39bp 1515.81bp 64.39bp 340.96bp},clip,width=0.94\textwidth,max height=0.38\textheight]{Pope.pdf}
\end{sourceblock}
u(x,t) 是傅里叶模式方程。 ( 8 ) , 是如何
连续性方程要求 R(t) 和 f(t) 正交（即，
ii(t) 0)。假设这个条件最初满足（i+e
ii(0) = 0)，表明 k(t) 和 f ( t ) 的演化方程
（方程（11 .80）和（11 .83））保持这种正交性。

\subsection*{11.4.3 频谱的演变}\phantomsection\addcontentsline{toc}{subsection}{11.4.3 频谱的演变}
我们现在考虑由多种模式组成的随机脉动速度场，而不是单一的傅里叶模式。速度谱张量 @-*(K,
IJ t) 提供该领域的统计描述；值得注意的是，频谱的演变完全由 RDT 决定。

脉动速度场的初始条件（t = 0）可以表示为傅立叶模态之和

\begin{sourceblock}
\includegraphics[page=439,pagebox=mediabox,trim={64.39bp 783.90bp 64.39bp 1161.85bp},clip,width=0.94\textwidth,max height=0.38\textheight]{Pope.pdf}
\end{sourceblock}
其中 d' 表示一组波数中的一个（可以是随机的），f(rc0,0 ) 是相应的（随机）傅里叶系数。这些初始条件的要求是 和 f 是正交的（以满足
KO 连续性），并且模态为复共轭对（因此 u(x, 0) 为实数）。

每种模式根据快速畸变方程独立演化，如上一小节所述。因此，随后时间的速度场可以简单地表示为

\begin{sourceblock}
\includegraphics[page=439,pagebox=mediabox,trim={64.39bp 318.96bp 64.39bp 1622.80bp},clip,width=0.94\textwidth,max height=0.38\textheight]{Pope.pdf}
\end{sourceblock}
回想一下，波数 R(t) 从初始条件 R(0) KO 根据 Eq+( 11.80) 演化，傅里叶系数 f(lc0,t) 由 Eq+( 11.80) 演化。 (11 *83)+
湍流速度场（参见第 6.5 节）。根据速度场的傅里叶表示（方程（11.95）），它由下式明确给出

\begin{sourceblock}
\includegraphics[page=440,pagebox=mediabox,trim={64.39bp 1626.79bp 64.39bp 252.97bp},clip,width=0.94\textwidth,max height=0.38\textheight]{Pope.pdf}
\end{sourceblock}
这是方程。 （6.155）适应当前情况。在计算中，考虑由初始参数化的频谱也很方便

\begin{sourceblock}
\includegraphics[page=440,pagebox=mediabox,trim={64.39bp 1366.83bp 64.39bp 465.94bp},clip,width=0.94\textwidth,max height=0.38\textheight]{Pope.pdf}
\end{sourceblock}
速度谱张量的重要属性之一是它代表波数空间中的雷诺应力密度：对方程进行积分。我们得到的总和为 (11.96) 或 (11.97)

\begin{sourceblock}
\includegraphics[page=440,pagebox=mediabox,trim={64.39bp 968.88bp 64.39bp 908.89bp},clip,width=0.94\textwidth,max height=0.38\textheight]{Pope.pdf}
\end{sourceblock}
速度谱张量的演化方程可以从其定义（方程（11.96））以及 2 和 li 的演化方程中导出（参见练习 11.15）。结果是

\begin{sourceblock}
\includegraphics[page=440,pagebox=mediabox,trim={64.39bp 634.92bp 64.39bp 1275.84bp},clip,width=0.94\textwidth,max height=0.38\textheight]{Pope.pdf}
\end{sourceblock}
这是一个封闭的统计方程，准确地描述了快速畸变极限下湍流谱的演变。除了指定的平均速度梯度和自变量 t) 之外，方程中涉及的唯一量是频谱本身。

等式右侧的第一项。 (11 .99) 表示波数空间中频谱的传输，这是平均速度梯度导致波数矢量 2(t) 失真的结果。 (K,t ) (Eq. (11 .97)) 的演化方程与 @ij(., t) (Eq. (11 -99)) 相同，
ZJ 除了省略第一个术语之外。

方程中的接下来两项。 (11 .99) 代表产量，其余项源于快速压力。请注意快速压力项
当 的方程在波数空间上积分时，输运项消失，使得 Rey 诺兹应力方程的形式为

\begin{sourceblock}
\includegraphics[page=441,pagebox=mediabox,trim={64.39bp 1557.80bp 64.39bp 363.95bp},clip,width=0.94\textwidth,max height=0.38\textheight]{Pope.pdf}
\end{sourceblock}
其中四阶张量 M 是

\begin{sourceblock}
\includegraphics[page=441,pagebox=mediabox,trim={64.39bp 1245.84bp 64.39bp 630.92bp},clip,width=0.94\textwidth,max height=0.38\textheight]{Pope.pdf}
\end{sourceblock}
（可以验证（参见练习 11 .16）MijkiCs 的定义等价于上面介绍的方程 1 1 1 5 ）雷诺应力方程为
为封闭形式，因为张量 M 涉及有关傅立叶模态方向的信息，该信息不包含在雷诺应力本身中。

对方程求微分。 (11.96) 得到结果

\begin{sourceblock}
\includegraphics[page=441,pagebox=mediabox,trim={64.39bp 640.92bp 64.39bp 1131.86bp},clip,width=0.94\textwidth,max height=0.38\textheight]{Pope.pdf}
\end{sourceblock}
替代方程。 (11.80) 对于 dfcldt 以便重新表达第一项
在右侧为

\begin{sourceblock}
\includegraphics[page=441,pagebox=mediabox,trim={64.39bp 429.95bp 64.39bp 1470.81bp},clip,width=0.94\textwidth,max height=0.38\textheight]{Pope.pdf}
\end{sourceblock}
（提示：利用 2 独立于 K 的事实，以及筛选属性
的 delta 函数。）

替代方程。 (11 .83) 为 diij/dt 以获得结果

\begin{sourceblock}
\includegraphics[page=441,pagebox=mediabox,trim={64.39bp 172.98bp 64.39bp 1734.78bp},clip,width=0.94\textwidth,max height=0.38\textheight]{Pope.pdf}
\end{sourceblock}
考虑函数 mijke(r)，它被定义为
泊松方程

\begin{sourceblock}
\includegraphics[page=442,pagebox=mediabox,trim={64.39bp 1680.79bp 64.39bp 265.97bp},clip,width=0.94\textwidth,max height=0.38\textheight]{Pope.pdf}
\end{sourceblock}
使用方程式（2.48），写出格林函数的解
方程，因此表明原点处的解为

\begin{sourceblock}
\includegraphics[page=442,pagebox=mediabox,trim={64.39bp 1442.82bp 64.39bp 442.94bp},clip,width=0.94\textwidth,max height=0.38\textheight]{Pope.pdf}
\end{sourceblock}
证明方程的傅里叶变换。 (11 .105) 是

\begin{sourceblock}
\includegraphics[page=442,pagebox=mediabox,trim={64.39bp 1305.84bp 64.39bp 641.92bp},clip,width=0.94\textwidth,max height=0.38\textheight]{Pope.pdf}
\end{sourceblock}
其中 hijk(Ke) 表示 mijke( r) 的傅立叶变换。由此得到

\begin{sourceblock}
\includegraphics[page=442,pagebox=mediabox,trim={64.39bp 1128.86bp 64.39bp 787.90bp},clip,width=0.94\textwidth,max height=0.38\textheight]{Pope.pdf}
\end{sourceblock}
表明 Mijke 的两个定义（方程 11.15 和 11.101）是
包含 en t。

\begin{sourceblock}
\includegraphics[page=442,pagebox=mediabox,trim={64.39bp 762.90bp 64.39bp 988.88bp},clip,width=0.94\textwidth,max height=0.38\textheight]{Pope.pdf}
\end{sourceblock}
Mijkei 的 sion

\begin{sourceblock}
\includegraphics[page=442,pagebox=mediabox,trim={64.39bp 649.92bp 64.39bp 1281.84bp},clip,width=0.94\textwidth,max height=0.38\textheight]{Pope.pdf}
\end{sourceblock}
其中a和y是常数。显示方程。 (11.109) 确定
这些常数是

\begin{sourceblock}
\includegraphics[page=442,pagebox=mediabox,trim={64.39bp 438.94bp 64.39bp 1450.82bp},clip,width=0.94\textwidth,max height=0.38\textheight]{Pope.pdf}
\end{sourceblock}
\begin{sourceblock}
\includegraphics[page=442,pagebox=mediabox,trim={64.39bp 325.96bp 64.39bp 1611.80bp},clip,width=0.94\textwidth,max height=0.38\textheight]{Pope.pdf}
\end{sourceblock}
使用练习 11.18 的结果表明快速压力-
应变率张量（对于各向同性湍流）是

\begin{sourceblock}
\includegraphics[page=442,pagebox=mediabox,trim={64.39bp 160.98bp 64.39bp 1774.78bp},clip,width=0.94\textwidth,max height=0.38\textheight]{Pope.pdf}
\par\vspace{0.45\baselineskip}
\small 表 11.2。简单变形的非零平均速度梯度：S，是 Si 的最大特征值；
\end{sourceblock}
\clearpage
\begin{sourceblock}
\includegraphics[page=443,pagebox=mediabox,trim={64.39bp 119.39bp 64.39bp 262.97bp},clip,width=\textwidth,max height=0.72\textheight]{Pope.pdf}
\end{sourceblock}
\subsection*{11.4.4 初始各向同性湍流的快速畸变}\phantomsection\addcontentsline{toc}{subsection}{11.4.4 初始各向同性湍流的快速畸变}
\begin{sourceblock}
\includegraphics[page=444,pagebox=mediabox,trim={64.39bp 1190.85bp 64.39bp 109.44bp},clip,width=\textwidth,max height=0.72\textheight]{Pope.pdf}
\par\vspace{0.45\baselineskip}
\small 图 11.8。雷诺应力的演变为轴对称收缩快速变形。虚线(u?)/k(0)，实线(u3)/k(0)=(u3)/k(0)。
\end{sourceblock}
轴对称收缩 对于轴对称收缩，波数演化（方程（1 1.80））：

显然，iZ2 和 i13 呈指数增长，而 ill 则下降。因此，如图16(a)所示，单位波矢2(t)的轨迹趋于与0对应的本初子午线。

(u:) 的产生式为负 (Pll - ( ) )，而 (ui) 的产生式为
(\textasciitilde{}22)s\textasciitilde{})：正 ( P (u:) 在统计上与 ( u ) 相同。因此，雷诺应力演化（图 11.8）符合预期。从图 11.9 中可以看出，在大部分时间 (ui) 随着 exp(Sjbt) 增加。这是产生的结果：相比之下，快速压力-应变率率张量趋于零。在早期，(u;) 增加得比随着压力-应变率率张量将能量重新分配给 ( u )，波矢在很大程度上趋于本初子午线 (El O)，并且傅里叶分量 ii 与子午线 (dl 0) 对齐。

轴对称膨胀虽然可以在实验中近似轴对称收缩中的快速变形，但由于流动分离，扩散器中的快速轴对称膨胀是不可能的。然而，可以通过以下方式研究流量

\begin{sourceblock}
\includegraphics[page=445,pagebox=mediabox,trim={64.39bp 1191.85bp 64.39bp 109.44bp},clip,width=\textwidth,max height=0.72\textheight]{Pope.pdf}
\par\vspace{0.45\baselineskip}
\small 图 11.9。 (v3)（在对数尺度上）轴对称收缩快速变形的演变（实线）。虚线为\}exp(Sjt)，表示渐近增长率。
\end{sourceblock}
在某些方面，这种流动与轴对称收缩相反。波数 fcl 呈指数增长，而 22 和 23 则减少，因此 2(t) 在单位球面上的轨迹（图 11.6（b））趋于极点
(u:) 的产生为正 (Pll 4(u:)Sj*)，而 (ui) 和 (ui) 的产生为负 (P22 -2(u;)SA)。然而，对于该流动，快速压力强烈抑制 (u) A 的增长，当波矢趋于 Z1 +-1 时，连续性方程（由快速压力强制执行）要求 G1 趋于零。

图 11.10 显示了雷诺应力的演变。在大多数情况下，这些比例是

\begin{sourceblock}
\includegraphics[page=445,pagebox=mediabox,trim={64.39bp 532.93bp 64.39bp 1407.82bp},clip,width=0.94\textwidth,max height=0.38\textheight]{Pope.pdf}
\end{sourceblock}
并且每个都以 exp(S?"t) 的形式增长。因此，可以推断出雷诺应力平衡中的项

\begin{sourceblock}
\includegraphics[page=445,pagebox=mediabox,trim={64.39bp 309.96bp 64.39bp 1592.80bp},clip,width=0.94\textwidth,max height=0.38\textheight]{Pope.pdf}
\end{sourceblock}
当它们除以 Sibk 时，a re

\begin{sourceblock}
\includegraphics[page=445,pagebox=mediabox,trim={64.39bp 119.39bp 64.39bp 1761.78bp},clip,width=0.94\textwidth,max height=0.38\textheight]{Pope.pdf}
\end{sourceblock}
\begin{sourceblock}
\includegraphics[page=446,pagebox=mediabox,trim={64.39bp 1176.85bp 64.39bp 109.44bp},clip,width=\textwidth,max height=0.72\textheight]{Pope.pdf}
\par\vspace{0.45\baselineskip}
\small 图 11.10。轴对称膨胀快速变形的雷诺应力的演变。虚线显示渐进增长为
\end{sourceblock}
快速压力消耗的能量是生产所增加能量的四分之三。

该生产不仅增加了能量，而且往往会增加各向异性。对于这个流（在大部分时间），pij 和 R D gbiv en 由方程 1 表示。 (11.120) 相关联

\begin{sourceblock}
\includegraphics[page=446,pagebox=mediabox,trim={64.39bp 786.90bp 64.39bp 1146.86bp},clip,width=0.94\textwidth,max height=0.38\textheight]{Pope.pdf}
\end{sourceblock}
表明快速压力抵消了 90\% 的各向异性产生（参见方程 (11 .1 16)）。

平面应变 对于平面应变，波数演化为
导致如图 11.6(\textasciitilde{})A 所示的轨迹。在很大程度上，波矢趋向于+I极点。

雷诺应力产生率为

\begin{sourceblock}
\includegraphics[page=446,pagebox=mediabox,trim={64.39bp 252.97bp 64.39bp 1694.79bp},clip,width=0.94\textwidth,max height=0.38\textheight]{Pope.pdf}
\end{sourceblock}
因此，图 11.11 中观察到的 的衰减和增长是
(u:) (u:) 是预料之中的。然而，由于波矢趋向于极点 +l(a t

\begin{sourceblock}
\includegraphics[page=447,pagebox=mediabox,trim={64.39bp 715.91bp 64.39bp 1212.85bp},clip,width=0.94\textwidth,max height=0.38\textheight]{Pope.pdf}
\end{sourceblock}
\begin{sourceblock}
\includegraphics[page=448,pagebox=mediabox,trim={64.39bp 1191.85bp 64.39bp 109.44bp},clip,width=\textwidth,max height=0.72\textheight]{Pope.pdf}
\par\vspace{0.45\baselineskip}
\small 图 11.12。剪切快速变形的雷诺应力各向异性的演变。
\end{sourceblock}
\begin{sourceblock}
\includegraphics[page=448,pagebox=mediabox,trim={64.39bp 461.94bp 64.39bp 835.90bp},clip,width=\textwidth,max height=0.72\textheight]{Pope.pdf}
\par\vspace{0.45\baselineskip}
\small 图 11.13。剪切快速变形的雷诺应力不变量的演化。从原点（对应于各向同性）开始，每个符号给出剪切量 St = 0.5 后的状态。
\end{sourceblock}
快速加压具有抑制最大雷诺应力增长速度的效果。事实上，渐近增长似乎是线性的而不是指数的（图）。然而，对于剪切，抑制是
11..1 4)。间接的。在极点 +1 处，连续性方程不排除

\begin{sourceblock}
\includegraphics[page=449,pagebox=mediabox,trim={64.39bp 1200.85bp 64.39bp 109.44bp},clip,width=\textwidth,max height=0.72\textheight]{Pope.pdf}
\par\vspace{0.45\baselineskip}
\small 图 11.14。快速变形剪切的湍动能的演变。
\end{sourceblock}
然而，在这些极点处被排除，因此剪切应力 (u1u2)（以及
ti2 生产 Pll)a 受到抑制。在大多数情况下，R22i 为负值，不断从已经很小的法向应力 (u;) 中移除能量。

\subsection*{11.4.5 最后评论}\phantomsection\addcontentsline{toc}{subsection}{11.4.5 最后评论}
RDT应用于均匀湍流的主要结论如下。

(i) RDT 在极限范围内适用于平均速度的大小
与湍流率相比，梯度变得非常大。的
需要条件 ST Sk/r 1 才能应用 RDT
含能量运动：将其应用于所有运动尺度，
需要更严格的条件 ST, 1。 RDT 效果如何
代表不满足这些条件时的湍流进行讨论

作者：萨维尔 (Savill) (1987) 以及亨特 (Hunt) 和卡拉瑟斯 (Carruthers) (1990)。

(ii) 湍流的演变取决于几何形状和
变形量，但与速率无关。

(iii) 在快速畸变方程（方程（11.68）和（1169））中，作用
快速压力的计算是为了强制执行连续性方程 V u 0。

(iv) 快速畸变方程是线性的，因此傅立叶方程
系数 a(t) 通过常微分方程演化(11 .80)

(v) 傅里叶系数的演变取决于单位波向量-
tor = 但不取决于波数的大小
(vi) 在快速畸变极限下，速度谱张量 cP,,(l\textasciitilde{}, t)

由精确的闭合方程演化（方程（11 -99））

(vii) 雷诺应力方程（Eq. (11.100)）不是封闭形式
因为四阶张量 Mijke（决定了快速
压力-应变率率张量 R(I rr j ) 涉及有关以下信息
傅立叶模式的方向（方程（11 .101）） 两个同质场
新湍流可以具有相同的雷诺应力张量

\begin{sourceblock}
\includegraphics[page=450,pagebox=mediabox,trim={64.39bp 1287.84bp 64.39bp 654.92bp},clip,width=0.94\textwidth,max height=0.38\textheight]{Pope.pdf}
\end{sourceblock}
但不同的值

因此，雷诺应力演化并不是唯一阻止
由雷诺应力开采。 (viii) 四个简单变形的快速变形方程的解
根据几何形状的不同，化学反应揭示了截然不同的行为。

对于大倍轴对称膨胀，快速压力
抵消了 90\% 的各向异性产生（方程（11 121））

各向同性湍流受到任何扭曲时的初始瞬间，
相应的数字是60\%。然而，对于轴对称构造
牵引力（大部分时间）为零：快速压力-应变率率
张量相对于生产消失。在剪切力下，组件
没有生产，但快速的压力却将能量从
该组件。

Kassinos 和 Reynolds (1994) 以及 Van Slooten 和 Pope (1997) 开发了一些不同的均匀湍流 RDT 方法。正如 Hunt (1973) 以及 Hunt 和 Carruthers (1990) 所描述的，RDT 也可以应用于非均匀流。

\section*{11.5 压力-应变率模型}\phantomsection\addcontentsline{toc}{section}{11.5 压力-应变率模型}
前两节考虑零平均速度梯度 (Sk/\& 0) 和非常大的平均速度的极限情况下的均匀湍流

\subsection*{11.5.1 基本模型LRR-IP}\phantomsection\addcontentsline{toc}{subsection}{11.5.1 基本模型LRR-IP}
完全由于缓慢的压力（Rij）；在第二种情况下，它是由于
= il 完全与快速压力 ( e . R ij R (f) ) 这里我们考虑与湍流剪切流更相关的情况，其中 S k / \& 是阶数统一的，并且慢压力场和快速压力场都很重要。

对于均匀湍流，耗散被视为各向同性，雷诺应力方程（方程（11 .1））为
压力-应变率率张量 Rij 是唯一不采用封闭形式的项。

11.5.1 (LRR-IP)

基本型号 Riji s 的基本型号

\begin{sourceblock}
\includegraphics[page=451,pagebox=mediabox,trim={64.39bp 1127.86bp 64.39bp 809.90bp},clip,width=0.94\textwidth,max height=0.38\textheight]{Pope.pdf}
\end{sourceblock}
第一项是 R(=s) 的 Rotta 模型，第二项是 Naot、Shavit 和 Wolfshtein (1970) 提出的 R(mf) pmro 的各向同性生产 (IP) 模型。 IP 模型假设快速压力部分抵消了生产增加雷诺应力各向异性的影响。这确实是在快速畸变轴对称膨胀中观察到的效应（参见方程（11.121））。此外，如果模型常数取为 C2，则 IP 模型会产生各向同性湍流对所有快速畸变的正确初始响应（参见方程（11.116））。

Rotta 模型和 IP 模型的这种组合是 Launder、Reece 和 Rodi (1975) 提出的两个模型中的第一个，也称为 LRR-IP。 Launder (1996) 建议 Rotta 常数的 CR 值为 1.8； IP常数取C2，与RDT一致。

图 11.15 显示了 LRR-IP 模型的计算结果与均匀剪切流的 DNS 数据的比较。从雷诺应力方程（Eq. (11 .128)）和 Rij 模型（Eq. (1. 2 9 ) ）以及建模方程（Eq. (10.56)）获得一组封闭的常微分方程。唯一的非零平均速度梯度 d(Ul)/dx2 S 是常数；雷诺应力最初是各向同性的；因此表征方程解的单个无量纲参数是S k l \textasciitilde{}w 的初始值，取2.36 以匹配DNS。

从图 11.15 可以看出，LRR-IP 模型合理地体现了 b12 和 bll 的演化，差异通常较小。

\begin{sourceblock}
\includegraphics[page=452,pagebox=mediabox,trim={64.39bp 394.95bp 64.39bp 1524.81bp},clip,width=0.94\textwidth,max height=0.38\textheight]{Pope.pdf}
\end{sourceblock}
\begin{sourceblock}
\includegraphics[page=452,pagebox=mediabox,trim={64.39bp 119.39bp 64.39bp 1802.77bp},clip,width=0.94\textwidth,max height=0.38\textheight]{Pope.pdf}
\end{sourceblock}
\begin{sourceblock}
\includegraphics[page=453,pagebox=mediabox,trim={64.39bp 1725.78bp 64.39bp 158.98bp},clip,width=0.94\textwidth,max height=0.38\textheight]{Pope.pdf}
\end{sourceblock}
由此得到

\begin{sourceblock}
\includegraphics[page=453,pagebox=mediabox,trim={64.39bp 1403.82bp 64.39bp 540.93bp},clip,width=0.94\textwidth,max height=0.38\textheight]{Pope.pdf}
\end{sourceblock}
11 S。 2 人们还提出了许多其他压力-应变率率模型，包括： Hanjalic 和 Launder (1972) 的 HL 模型； Launder (1975) 的准各向同性模型 (LRR-QI)； Shih 和 Lumley 的 SL 模型

\subsection*{11.5.2 其他压力-应变率率模型}\phantomsection\addcontentsline{toc}{subsection}{11.5.2 其他压力-应变率率模型}
等人。 （1985）； Jones 和 Musonge (1988) 的 JM 模型； Speziale、Sarkar 和 Gatski (1991) 的 SSG 模型；以及 Fu、Launder 和 Tselepidakis (1987) 的 FLT 模型。其中前五个（以及 LRR-IP 模型）可以写成
其中无量纲、对称、偏张量 7 在表 11.3 中给出，并根据归一化平均应变率定义
和标准化平均旋转率

表 11.4 给出了定义某些模型的系数 f。 （FLT 模型涉及额外的张量。）

7( 1.) 中的术语对应于 Rotta 模型贡献，即 Rotta
存在-f(1)。 LRR 和 JM 模型将 Rotta 系数设为常数，而 SSG 和 SL 模型则包含依赖性
馅饼。的系数；模型如表11.4所示，仅

\begin{center}\small 表 11.3。方程中无量纲、对称、偏张量的定义。 (11.135)\end{center}
\begin{sourceblock}
\includegraphics[page=454,pagebox=mediabox,trim={64.39bp 1143.86bp 64.39bp 256.97bp},clip,width=\textwidth,max height=0.72\textheight]{Pope.pdf}
\end{sourceblock}
LRR-IP 模型是第一个被广泛使用的雷诺应力模型，但应该理解的是，所有这些模型都建立在早期贡献的基础上。例如，LRR、JM和SSG模型都是Hanjalik和Launder（1972）提出的一般模型的特例；并且 HL 模型的特定形式包括相同的依赖性
PIE Roltta 系数出现在 SSG 模型中。

理想的压力-列车速率模型在所有情况下都是准确的，因此符合 Rij 的所有已知属性。这些属性如下。

(i) 在快速畸变极限（S k l \textasciitilde{} a）下，R ij 均值呈线性
速度梯度。

( 由张量决定 (方程 (11.14)-(11.15)) 其中
满足方程中给出的精确关系。 (11.109)。

(iii) R(. r.) 在雷诺应力中是线性的，因为 Milk[ 在
频谱（方程（11.101））。

(iv) Rij 使得雷诺应力张量仍然可实现。

(v) 在二维湍流极限下，Rijs 满足材料框架
冷漠（参见第（2.9）节和Speziale（198 1））。

在很大程度上，这一领域的研究是出于开发符合这些特性的模型的愿望。基本模型（LRR-IP）不满足（ii），LRR-QI就是为此而开发的。 Shih-Lumley模型是为了满足强可实现性而开发的；但是，在这样做的过程中，燕鸥
LRR-IP LRR-QI SSG
朗德·琼斯 (Launder Jones) 和斯佩齐亚莱·施 (Speziale Shih) 以及
等人。等人。穆松格·拉姆利
fl 值 - 参见 Shih
和 Lumley 常数
其他模型在除病理情况外的所有情况下都满足弱可实现性，并且可以修改它们的系数以满足所有情况（Durbin 和 Speziale 1994）。提到的模型均不满足 (v)，但 Haworth 和 Pope (1986) 以及 Ristorcelli、Lumley 和 Abid (1995) 提出的模型满足 (v)。

需要理解的一个微妙但重要的一点是，上面列举的属性与压力-应变率率模型有关。然而，没有
理想的这种模型是存在的。正如已经观察到的（见图 11.4），各向异性湍流的衰减并不是由雷诺应力中的“已知”唯一决定的
米克封闭。同样，张量也不是唯一确定的（参见
( u i u j ) 等式。 (11.126) 和 ( l l l 2 ) ) 在缺乏理想模型的情况下，自然会寻求“最佳”模型“最佳”，因为该模型在应用于感兴趣的流量范围时尽可能准确。上面列举的许多属性都属于极端情况，例如快速畸变极限、二维湍流、二分量湍流
模型形式和系数的规范 当应用于感兴趣的流量时，模型的定量性能比满足极端限制中的精确约束更重要。

与LRR-IP相比，LRR-QI模型具有符合性质(ii)的理论优势。然而，经验表明，在实践中，LRR-IP 明显优越（Launder 1996）。 SSG 模型符合的属性比 LRR-IP 少，它是非线性的
(a) 表 11.4 给出的模型中哪个与
性质 (i) 和 (iii)？

(b) 哪个模型对于初始快速畸变是正确的
各向同性湍流？

(c) 哪个模型对于任意快速扭曲是正确的？

11.22 证明 LRR-QI 是在 bij 中呈线性的最通用模型
和 8(Ui)/8xj 这与 Mijke 的性质一致

\section*{11.6 向非均匀流动的扩展}\phantomsection\addcontentsline{toc}{section}{11.6 向非均匀流动的扩展}
前面几节的发展重点是均匀湍流，而雷诺应力模型的主要应用是非均匀湍流。本节描述的是获得完整的非均匀流雷诺应力模型模型的其余步骤。下一节将介绍对近壁区域的其他处理。

精确雷诺应力方程（方程（11 .1））中要建模的项是耗散张量 eij、压力-应变率率张量 Rij 以及第 11.1 节中讨论的输运 TkijAS，耗散可取为
'\&aij 是各向同性的。（靠近墙壁的除外）。本节稍后将描述所使用的模型耗散方程的形式。

\subsection*{11.6.1 重新分配}\phantomsection\addcontentsline{toc}{subsection}{11.6.1 重新分配}
正如 11.2 节中所讨论的，有不同的分解
nij 将速度-压力-梯度张量转换为再分配和传输项；

Mansour、Kim 和 Moin（1988）对河道流 DNS 数据的分析强烈表明，对于这种近壁流，最好使用 R！” 的分解，而不是 Rij 的分解。

inii 发现动能的压力传递相当小（参见
尼吉无花果。第 285 和 314 页上的 7.18 和 7.34），实际上这本身就差不多
nij 再分配。 （第 3 16-3 17 页上的图 7.35-7.38）的轮廓表现出简单的行为，其中 Iij 在壁处为零（因为 u 在那里为零）。相比之下，aT、/dy 表现出更复杂的行为，包括几个符号变化。尽管这两项在壁处均为零，但在粘性子层内它们很大，但几乎抵消。鉴于这些考虑，我们今后使用分解方程。 (11 .5) 就 R (a) 和 T (PI .

对于非均匀流，重新分布根据局部量进行建模。也就是说，R (a)( x, t) 根据 (uiuj)、a ( u i ) / a x j 进行建模，并在 (x, t) 处进行评估，就像在均匀湍流中一样（方程 (11 135)）。就所得模型方程的求解容易性而言，与包含非局部量的替代方案相比，这无疑是一个权宜之计的假设。然而，应该认识到，p'(x, t) 受泊松方程 (Eq. ( 1 9 ) ) 控制，因此它受到距 x 一定距离的 a(Ui)/ axj 等量的影响。因此，与均匀湍流相比，Rij 模型对于非均匀流的安全性较差。 （相比之下，在第 11.8 节的椭圆松弛模型中，R(aa)m 的建模是非局部的。）

\subsection*{11.6.2 雷诺应力输运}\phantomsection\addcontentsline{toc}{subsection}{11.6.2 雷诺应力输运}
nij 速度-压力-梯度张量根据方程分解。 (11 .5，雷诺应力演化的精确方程为

\begin{sourceblock}
\includegraphics[page=457,pagebox=mediabox,trim={64.39bp 501.94bp 64.39bp 1396.82bp},clip,width=0.94\textwidth,max height=0.38\textheight]{Pope.pdf}
\end{sourceblock}
其中三个通量是粘性扩散

\begin{sourceblock}
\includegraphics[page=457,pagebox=mediabox,trim={64.39bp 341.96bp 64.39bp 1559.80bp},clip,width=0.94\textwidth,max height=0.38\textheight]{Pope.pdf}
\end{sourceblock}
\begin{sourceblock}
\includegraphics[page=457,pagebox=mediabox,trim={64.39bp 219.97bp 64.39bp 1664.79bp},clip,width=0.94\textwidth,max height=0.38\textheight]{Pope.pdf}
\end{sourceblock}
和湍流对流

\begin{sourceblock}
\includegraphics[page=457,pagebox=mediabox,trim={64.39bp 119.39bp 64.39bp 1842.77bp},clip,width=0.94\textwidth,max height=0.38\textheight]{Pope.pdf}
\end{sourceblock}
除粘性壁区域外，粘性扩散可以忽略不计；而且，由于该术语是封闭形式，因此不需要进一步讨论。

压力传输 在 DNS 出现之前，关于压力相关性的可靠信息很少。 Lumley (1978) 在对近均匀湍流分析的基础上提出了模型

\begin{sourceblock}
\includegraphics[page=458,pagebox=mediabox,trim={64.39bp 1428.82bp 64.39bp 471.94bp},clip,width=0.94\textwidth,max height=0.38\textheight]{Pope.pdf}
\end{sourceblock}
由于压力传输由方程给出。 (11.140) 是各向同性的，可以通过动能方程中的相应项来检验

\begin{sourceblock}
\includegraphics[page=458,pagebox=mediabox,trim={64.39bp 1257.84bp 64.39bp 673.91bp},clip,width=0.94\textwidth,max height=0.38\textheight]{Pope.pdf}
\end{sourceblock}
Lumley 的模型是

\begin{sourceblock}
\includegraphics[page=458,pagebox=mediabox,trim={64.39bp 1124.86bp 64.39bp 807.90bp},clip,width=0.94\textwidth,max height=0.38\textheight]{Pope.pdf}
\end{sourceblock}
这是对流通量的 -2

\begin{sourceblock}
\includegraphics[page=458,pagebox=mediabox,trim={64.39bp 986.88bp 64.39bp 943.88bp},clip,width=0.94\textwidth,max height=0.38\textheight]{Pope.pdf}
\end{sourceblock}
槽道流（第 285 页上的图 7.18）和边界层（第 314 页上的图 7.34）的动能预算表明，靠近壁面的压力传递不是很显着，并且 Lumley 模型在质量上是不正确的。然而，在边界层边缘，压力传输更为重要，Lumley 的模型至少在质量上是正确的。

为了检查自由剪切流的压力传输，自相似时间混合层的动能预算如图 11.16 所示。从图11.16(a)'可以看出，大部分层位的压力传输都比较小，Lumley的模型是相当合理的。

图 11.16(b) 更详细地检查了层的边缘。层内的旋转湍流涨落会引起非湍流区域的无旋转涨落（例如 y/é > 1）。这种能量传递受到波动压力场的影响，因此它在动能收支中表现为压力传递。从图 11.16(b) 可以看出，在层的边缘，这种压力传输占主导地位。

德穆伦等人。 (1996) 使用 DNS 数据分别检查由于慢速和快速压力引起的压力传输，并为每种贡献提出了模型。在大多数雷诺应力模型中，压力传递要么被忽略，要么（隐式或显式）与压力传递一起建模

\begin{sourceblock}
\includegraphics[page=459,pagebox=mediabox,trim={64.39bp 1361.83bp 64.39bp 109.44bp},clip,width=\textwidth,max height=0.72\textheight]{Pope.pdf}
\par\vspace{0.45\baselineskip}
\small 图 11.16。来自 Rogers 和 Moser (1994) 的 DNS 数据的时间混合层中的动能预算： (a) 整个流； (b) 层边缘的放大视图。对预算的贡献是生产 P、耗散 \textasciitilde{}e、变化率 -dk/dt、湍流传输和压力传输（虚线）。所有量均通过速度差和层厚度 5 进行归一化（参见
\end{sourceblock}
第 140 页图 5.21）。

梯度扩散模型

\begin{sourceblock}
\includegraphics[page=459,pagebox=mediabox,trim={64.39bp 888.89bp 64.39bp 936.88bp},clip,width=0.94\textwidth,max height=0.38\textheight]{Pope.pdf}
\end{sourceblock}
其中 C 是模型常数。更普遍使用的是 Daly 和 Harlow 的模型（1970 年），该模型使用雷诺应力张量来定义各向异性

\begin{sourceblock}
\includegraphics[page=459,pagebox=mediabox,trim={64.39bp 653.92bp 64.39bp 1203.85bp},clip,width=0.94\textwidth,max height=0.38\textheight]{Pope.pdf}
\end{sourceblock}
对于这个模型，Launder (1990) 建议常数 C 的值为 0.22。

分别建模，那么一致的模型是
IJ 要求对于所有三个指数都是对称的。 Mellor 和 Herring (19731,
以及 Hanjalik 和 Launder (1972)，

\begin{sourceblock}
\includegraphics[page=459,pagebox=mediabox,trim={64.39bp 119.39bp 64.39bp 1847.77bp},clip,width=0.94\textwidth,max height=0.38\textheight]{Pope.pdf}
\end{sourceblock}
然而，T i 的一致模型需要关于 a 和 j 对称，这只是上述所有四个模型都满足的要求。

检查三重相关的输运方程

( U i U j U k ) 可以。激发更复杂的模型，包括平均速度梯度。但从业者普遍的经验是，T的造型，'。并不是整个模型中的关键因素，相对简单的 Daly-Harlow 模型就足够了（Launder 1990）。这种观点受到 Parneix、Laurence 和 Durbin (1998b) 的质疑，他们认为传输模型的缺陷是导致向后台阶流量计算不准确的原因。 (u) 模型与实验数据的直接测试可以在 Schwarz 和 Bradshaw (1994) 及其参考文献中找到。

耗散方程 雷诺应力模型中使用的标准模型方程是由 Hanjalik 和 Launder (1972) 提出的，

\begin{sourceblock}
\includegraphics[page=460,pagebox=mediabox,trim={64.39bp 954.88bp 64.39bp 934.88bp},clip,width=0.94\textwidth,max height=0.38\textheight]{Pope.pdf}
\end{sourceblock}
其中 C、0.15、C、1.44 和 CC2 1.92（Launder 1990）。该方程与 k-E 模型（方程（10.53））中使用的方程有两个不同之处。首先，直接根据雷诺应力评估产量
2'vTSijSij。其次，扩散项涉及各向异性扩散率。

正如第 10.4.3 节中提到的，已经提出了对耗散方程的一些修改。在雷诺应力模型中，各向异性张量 bij 的不变量可以用于此类修改。例如，lhunder (1996) 建议将系数修改为

\begin{sourceblock}
\includegraphics[page=460,pagebox=mediabox,trim={64.39bp 486.94bp 64.39bp 1459.82bp},clip,width=0.94\textwidth,max height=0.38\textheight]{Pope.pdf}
\end{sourceblock}
其中 --variants q 和 F 由等式定义。 （11 .28）和（11 .34）。 Hanjalik、Jakirlik 和 Hadiik (1997) 建议进行其他修改。

\subsection*{11.6.3 耗散方程}\phantomsection\addcontentsline{toc}{subsection}{11.6.3 耗散方程}
对数定律，卡门常数由下式给出

\subsection*{11.7.1 近壁效应}\phantomsection\addcontentsline{toc}{subsection}{11.7.1 近壁效应}
评估 LRR-IP 模型给出的 bij（方程（11 .133））

\section*{11.7 近壁处理}\phantomsection\addcontentsline{toc}{section}{11.7 近壁处理}
在湍流中，壁的存在会导致许多不同的影响。例如，对于通常坐标系中的边界层，其中一些影响是：

(i) 低雷诺数湍流雷诺数ReL
当接近墙壁时，k2/(\&v)t 结束为零；

(ii) 高剪切速率 最高平均剪切速率 a ( U ) / a y 出现在
墙（见练习 11.24）；

(iii) 小 y 的双分量湍流，(v2) 随 v4 变化，而 (u2)

和 (w2) 变化为（方程（7.60）-（7.63）），因此，当墙是 ap-
继续下去，湍流趋于二分量极限（图 11 1）；

(iv) 通过压力场 (pch) 的壁阻挡，方程（11.10）即时
渗透率条件 V 0（在 y 0 处）影响流量达到（的
阶）从墙上的积分比例（Hunt 和 Graham 1978）。

这些效应需要对基本 k-E 和雷诺应力湍流模型进行修改。 （研究不存在 (i) 和 (ii) 之一或两者的湍流，例如自由表面处的湍流 (Brumley 1984) 和壁附近的无剪切湍流 (Thomas 和 Hancock 1977，Perot 和 Moin 1995，Aronson、Johansson 和 Lofdahl 1997)，也能提供丰富的信息。）

自 20 世纪 70 年代以来，基本 k-E 和雷诺应力模型的形式没有改变。对于近壁处理来说情况并非如此。 1月9日8月0日末，粘性近壁区域的DNS数据变得可用，该数据很难通过实验进行研究。这些数据使现有模型在质量上是错误的，并导致了新的、持续的发展。

练习 11.24 考虑 \textasciitilde{} , a ( U ) / a ya 是剪切速率的无量纲测量。

证明它在墙上的价值是

\begin{sourceblock}
\includegraphics[page=461,pagebox=mediabox,trim={64.39bp 119.39bp 64.39bp 1824.77bp},clip,width=0.94\textwidth,max height=0.38\textheight]{Pope.pdf}
\end{sourceblock}
（参见第 287 页的练习 7.8）；而在对数律区域中它是

\begin{sourceblock}
\includegraphics[page=462,pagebox=mediabox,trim={64.39bp 1688.79bp 64.39bp 211.97bp},clip,width=0.94\textwidth,max height=0.38\textheight]{Pope.pdf}
\end{sourceblock}
小于 0.3（对于 y+ 30）。另一种无量纲
剪切速率的测量是 S k l \textasciitilde{} H。它的变化接近于
11.7.2 管状粘度 可以回忆一下，混合长度规格 l, \textasciitilde{}y 在近壁区域（例如 y+ 30）太大，并且从 van Driest 阻尼函数获得了一个大大改进的规格，方程： （7.145）。同样，

\subsection*{11.7.2 湍流黏度}\phantomsection\addcontentsline{toc}{subsection}{11.7.2 湍流黏度}
VT P 在近壁区域，如图 10.3 所示。

Jones 和 Launder (1972) 的原始 k-E 模型包括各种阻尼函数，允许该模型在粘性近壁区域内使用。湍流黏度由下式给出

\begin{sourceblock}
\includegraphics[page=462,pagebox=mediabox,trim={64.39bp 941.88bp 64.39bp 979.88bp},clip,width=0.94\textwidth,max height=0.38\textheight]{Pope.pdf}
\end{sourceblock}
其中阻尼函数 f 取决于湍流雷诺数，根据：

f p = e x P ( 1 ReL/50 ) 在壁面 f, 很小，在大 ReL,s o 处趋于统一，恢复了标准 k-E 公式。

图 10.3（第 371 页）中标注的量 vT\&/k2p 是渠道流量 f,C 的“测量值”。 Rodi 和 Mansour (1993) 基于类似的 DNS 数据评估了一些关于 f 的提案：最早期的提案，例如等式 1。 ( 1 . 1 6 )，非常不准确。 Rodi 和 Mansour (1993) 建议采用经验关系

\begin{sourceblock}
\includegraphics[page=462,pagebox=mediabox,trim={64.39bp 339.96bp 64.39bp 1600.80bp},clip,width=0.94\textwidth,max height=0.38\textheight]{Pope.pdf}
\end{sourceblock}
这种方法的一个站得住脚的观点是，阻尼函数 f 缺乏物理合理性，它被用来补偿基本模型中不正确的物理现象。没有理由假设方程。 (11 .157) 在应用于分离点或重新附着点附近的其他流时是准确的，对于

\begin{sourceblock}
\includegraphics[page=463,pagebox=mediabox,trim={64.39bp 1197.85bp 64.39bp 109.44bp},clip,width=\textwidth,max height=0.72\textheight]{Pope.pdf}
\end{sourceblock}
\subsection*{11.7.3 k 和 e 的模型方程。}\phantomsection\addcontentsline{toc}{subsection}{11.7.3 k 和 e 的模型方程。}
\begin{center}\small 图 11.17。 Re = 13,750 时槽道流的湍流黏度相对于 y+。 Kim 等人的符号、DNS 数据。 （1987）；实线，0,09k?/e；虚线，0.22(v?)k/e。\end{center}
Durbin (1991) 认为，交叉流速度 v 是边界层流中湍流传输的原因，并且湍流黏度的适当表达式为

\begin{sourceblock}
\includegraphics[page=463,pagebox=mediabox,trim={64.39bp 858.89bp 64.39bp 1068.87bp},clip,width=0.94\textwidth,max height=0.38\textheight]{Pope.pdf}
\end{sourceblock}
其中 C' 是常数。 （这也是Daly-Harlow模型中的跨流扩散系数，方程（11.47））由图11.17可知（11.158）确实提供了整个近壁区域 vT 的出色表现。正如 Launder (1986) 所指出的，这意味着 f 和 ( v 2 )/ k 的轮廓形状相似。

近壁区域中 ( v 2 ) （相对于 k ）的强烈抑制（参见第 282 页和第 313 页的图 7.14 和 7.33）是由于无粘性壁阻挡效应，而不是粘性效应。因此，将 f 称为“低雷诺数校正”是不合适的。

和 的模型方程 大多数提出的 k 和 的模型方程与 Jones 和 Launder (1972) 的原始建议具有相同的形式。 k 的方程（包括分子粘度的传输）为

\begin{sourceblock}
\includegraphics[page=463,pagebox=mediabox,trim={64.39bp 119.39bp 64.39bp 1797.77bp},clip,width=0.94\textwidth,max height=0.38\textheight]{Pope.pdf}
\end{sourceblock}
在壁上，耗散值为

\begin{sourceblock}
\includegraphics[page=464,pagebox=mediabox,trim={64.39bp 1675.79bp 64.39bp 208.97bp},clip,width=0.94\textwidth,max height=0.38\textheight]{Pope.pdf}
\end{sourceblock}
（见方程（7.73））。 Jones 和 Launder (1972) 选择作为第二个变量

\begin{sourceblock}
\includegraphics[page=464,pagebox=mediabox,trim={64.39bp 1558.80bp 64.39bp 388.95bp},clip,width=0.94\textwidth,max height=0.38\textheight]{Pope.pdf}
\end{sourceblock}
从而可以应用更简单的边界条件 EyZo 0。 E 的模型方程为

\begin{sourceblock}
\includegraphics[page=464,pagebox=mediabox,trim={64.39bp 1341.83bp 64.39bp 552.93bp},clip,width=0.94\textwidth,max height=0.38\textheight]{Pope.pdf}
\end{sourceblock}
不同规格的阻尼功能和

Patel、Rodi 和 Scheuerer (1985) 以及 Rodi 和 Mansour (1993) 中给出了 fz 附加项 E。 DNS 数据揭示了所有提出的模型之间存在重大差异（Rodi 和 Mansour 1993）。

湍流时间尺度 z k / 在壁面处趋于零，当然所有相关的物理时间尺度都是严格正的。 Durbin (1991) 认为（常数乘以）柯尔莫哥洛夫时间尺度 z，（v 是最小的相关尺度，从而定义了修正的时间尺度

\begin{sourceblock}
\includegraphics[page=464,pagebox=mediabox,trim={64.39bp 917.88bp 64.39bp 1026.87bp},clip,width=0.94\textwidth,max height=0.38\textheight]{Pope.pdf}
\end{sourceblock}
取 CT 6。（T 由 CT7 确定，直到 y 5，更大的由 z 确定。）用 T 替换编写的标准耗散方程为

\begin{sourceblock}
\includegraphics[page=464,pagebox=mediabox,trim={64.39bp 699.91bp 64.39bp 1189.85bp},clip,width=0.94\textwidth,max height=0.38\textheight]{Pope.pdf}
\end{sourceblock}
在没有阻尼函数或附加项的情况下对该方程的解产生的轮廓与 DNS 数据非常一致 (Durbin 199 1)。

\subsection*{11.7.4 耗散张量}\phantomsection\addcontentsline{toc}{subsection}{11.7.4 耗散张量}
在雷诺应力模型中，耗散张量和压力速率
Eij of-strain-t ensor Rij 是主要受壁面影响的量。在高雷诺数且远离墙壁时，耗散被认为是各向同性的，

\begin{sourceblock}
\includegraphics[page=464,pagebox=mediabox,trim={64.39bp 208.97bp 64.39bp 1735.78bp},clip,width=0.94\textwidth,max height=0.38\textheight]{Pope.pdf}
\end{sourceblock}
然而，正如第 7 章中所观察到的（例如，第 318 页的图 7.39），随着墙
根据 ( \textasciitilde{} q (.7 .181)) 的定义和 ui 的级数展开式（方程 (7.56)-(7.58)），可以直接证明，当接近墙壁 (J O) 时，由下式给出

\begin{sourceblock}
\includegraphics[page=465,pagebox=mediabox,trim={64.39bp 1491.81bp 64.39bp 423.95bp},clip,width=0.94\textwidth,max height=0.38\textheight]{Pope.pdf}
\end{sourceblock}
Rotta (1951) 提出了简单模型

\begin{sourceblock}
\includegraphics[page=465,pagebox=mediabox,trim={64.39bp 1213.85bp 64.39bp 684.91bp},clip,width=0.94\textwidth,max height=0.38\textheight]{Pope.pdf}
\end{sourceblock}
为低雷诺数湍流的极限。在墙上，这与等式一致。 ( 1 1 . 1 6 ) 因为 (v2)/k 和 (uv)/k 在 y 0 处为零。但是，对于较小的 y，Elq. (11 .167) 分别低估 c12 和 2 和 4 倍。

为小 y 提供更准确的近似，与方程一致。 ( 1 1 6 )，Launder 和 Reynolds (1983) 以及 Kebede、Launder 和 Younis (1985) 引入了这个量
单位垂直于墙的位置在哪里。分母在壁处是统一的，并且被指定为使得 的迹是 2 \textasciitilde{}I.t 很容易验证，对于
IJ 墙在 y = 0 e nj = 6j2)，方程。 (11 .168) 恢复为方程。 (11.166)。

与等式。 (11.168) 给出墙处的值，方程(11.165)

Eij 远离墙壁，很自然地将模型建模为这两个函数的混合：

\begin{sourceblock}
\includegraphics[page=465,pagebox=mediabox,trim={64.39bp 444.94bp 64.39bp 1500.81bp},clip,width=0.94\textwidth,max height=0.38\textheight]{Pope.pdf}
\end{sourceblock}
其中混合函数 f 从壁处的 1 减小到远离壁的零。 Lai 和 So (1990) 提出

\begin{sourceblock}
\includegraphics[page=465,pagebox=mediabox,trim={64.39bp 222.97bp 64.39bp 1677.79bp},clip,width=0.94\textwidth,max height=0.38\textheight]{Pope.pdf}
\end{sourceblock}
图 11 .18 显示了边界层数据的剖面

图 11.18，湍流边界层中的归一化耗散分量
Re, = 1,410：符号，Spalart 的 DNS 数据（1988）；虚线，罗塔模型，
等式，(11.167)；实线，等式。 (11.169)。

可以看出，两个模型对于粘性子层都相当准确，更复杂的模型（方程（11.169））没有明显的优势。在流动中更远的地方 (yt > 200)，耗散的各向异性比 Rotta 模型预测的要小。另一方面，对于这种低雷诺数流，很明显，例如中的混合函数。随着 y 的增加，(11.169) 导致建模的 ¢ 比数据更快地变得各向同性。 （由于 Hanjalic 和 Jakirlic (1993) 以及 Launder (1996) 对 f 提出了更复杂的建议，它们与数据的一致性更好。）

上述一些近壁处理涉及到壁的距离 y（例如方程（11.157））或壁的法线\#（例如方程（11.168））。对于几何简单的流动，例如平板边界层，定义和评估这些量是很简单的。对于一般情况，£,(x)（定义为从点 x 到墙上最近点的距离）可以用来代替 y；和向量

\begin{sourceblock}
\includegraphics[page=466,pagebox=mediabox,trim={64.39bp 345.96bp 64.39bp 1599.80bp},clip,width=0.94\textwidth,max height=0.38\textheight]{Pope.pdf}
\end{sourceblock}
可以用来代替a。 a” 的性质是它在光滑表面上的每个点都等于 m，其大小 |n*| 小于或等于 1，并且可以随 x 不连续变化。

对于使用和“在”中的可取性存在不同的意见。

类似的行为。例如，在墙壁上 vk1l2 是 n 方向上的矢量，并且等式： (1.1.160)c 可用于给出

\begin{sourceblock}
\includegraphics[page=467,pagebox=mediabox,trim={64.39bp 1633.79bp 64.39bp 267.97bp},clip,width=0.94\textwidth,max height=0.38\textheight]{Pope.pdf}
\end{sourceblock}
这种关系在壁上是精确的，并且通过粘性底层提供了良好的近似。 Craft and Launder (19 96) 使用这种近似值来计算 nrn，在方程 1 中。 (11.168) 等式(1 1.169) 成为 c, 的局部模型。我

\subsection*{11.7.5 脉动压力}\phantomsection\addcontentsline{toc}{subsection}{11.7.5 脉动压力}
谐波压力 为了将脉动压力 p' 精确分解为快速 p(r)、慢速 p 和谐波 p(h) 贡献（方程 1 1 1 0 ）1 1 1 2 ），有必要指定边界条件。 p' 上的边界条件由脉动速度的演化方程 Eq. 获得。 （5.138）。在通常的坐标系中，在壁 (y 0) 处计算的该方程的法向分量产生诺依曼条件

\begin{sourceblock}
\includegraphics[page=467,pagebox=mediabox,trim={64.39bp 936.88bp 64.39bp 1011.87bp},clip,width=0.94\textwidth,max height=0.38\textheight]{Pope.pdf}
\end{sourceblock}
显然，壁上的法向压力梯度完全是由于粘性效应造成的：对于无粘性流动，dp'ldy 将为零。

p' 泊松方程（方程 (11 .9)）中的源项与粘度无关。因此，很自然地在 p(r) 和 p(" 上指定零正态梯度（无粘性）边界条件，并在谐波贡献 P ( \textasciitilde{} ) 上指定粘性条件（方程 (11 .173)）（M ansour 等人 (1988) 指的是 p(h)a 斯托克斯压力。）

对河道流 DNS 数据的检查（Mansour 等人，1988）表明，谐波压力在雷诺应力预算中起着次要作用，并且在 yf 15 之外不显着。这证实了这样一个事实，即壁对脉动压力场的显着影响是由不渗透条件引起的无粘性阻塞效应，而不是由无滑移条件引起的粘性效应。

练习 11.25 考虑无限平面上方的半无限湍流体
y 0 处的壁。湍流在 x 方向上统计上是均匀的
在墙上是

\begin{sourceblock}
\includegraphics[page=468,pagebox=mediabox,trim={64.39bp 1692.79bp 64.39bp 204.97bp},clip,width=0.94\textwidth,max height=0.38\textheight]{Pope.pdf}
\end{sourceblock}
其中 rcl 和 rc3 是归一化波数， 是归一化
傅里叶系数。确定相应的谐波压力
场，因此表明它随 y 呈指数衰减。

墙壁反射 “惯性”压力定义为

\begin{sourceblock}
\includegraphics[page=468,pagebox=mediabox,trim={64.39bp 1312.83bp 64.39bp 627.92bp},clip,width=0.94\textwidth,max height=0.38\textheight]{Pope.pdf}
\end{sourceblock}
p (i) 的格林函数解提供了对墙的阻挡效应的一些见解。考虑无限平面壁上方的湍流
0. 惯性压力由与 p' (Eq. 1 1.9)) 相同的泊松方程控制，可写为

\begin{sourceblock}
\includegraphics[page=468,pagebox=mediabox,trim={64.39bp 1054.87bp 64.39bp 892.89bp},clip,width=0.94\textwidth,max height=0.38\textheight]{Pope.pdf}
\end{sourceblock}
源 S(x, t) 由等式右侧给出。 (11.9) （乘以 为简单起见，假设远离感兴趣区域 (1x1 m)，流动是非湍流。因此 p li）在那里为零。壁面处的边界条件为

\begin{sourceblock}
\includegraphics[page=468,pagebox=mediabox,trim={64.39bp 771.90bp 64.39bp 1125.86bp},clip,width=0.94\textwidth,max height=0.38\textheight]{Pope.pdf}
\end{sourceblock}
源被定义在上半空间 y 2 0 中。这个定义通过关于墙反射 S 来扩展：

\begin{sourceblock}
\includegraphics[page=468,pagebox=mediabox,trim={64.39bp 596.92bp 64.39bp 1341.83bp},clip,width=0.94\textwidth,max height=0.38\textheight]{Pope.pdf}
\end{sourceblock}
由于 S 关于 y 0 对称，因此整个域中泊松方程的解 pci) 也关于 y 0 对称，因此满足边界条件 Eq. (11.177)。这个解决方案是

\begin{sourceblock}
\includegraphics[page=468,pagebox=mediabox,trim={64.39bp 328.96bp 64.39bp 1556.80bp},clip,width=0.94\textwidth,max height=0.38\textheight]{Pope.pdf}
\end{sourceblock}
（参见等式（2.49））。

如图 11.19 所示，对于上半空间中的每个点 x' (x', y', )，图像点 (xu,y ", 定义为

\begin{sourceblock}
\includegraphics[page=468,pagebox=mediabox,trim={64.39bp 119.39bp 64.39bp 1840.77bp},clip,width=0.94\textwidth,max height=0.38\textheight]{Pope.pdf}
\end{sourceblock}
\begin{sourceblock}
\includegraphics[page=469,pagebox=mediabox,trim={64.39bp 1540.81bp 64.39bp 109.44bp},clip,width=\textwidth,max height=0.72\textheight]{Pope.pdf}
\par\vspace{0.45\baselineskip}
\small 图 11.19。点 x' 及其图像 x" 的草图，显示出现在格林函数解（方程式 1）中的向量 r' 和 r"。 （11.181）和（11.182）。
\end{sourceblock}
然后利用定义 v' r x' x 和 v" x" x，可以重写解

\begin{sourceblock}
\includegraphics[page=469,pagebox=mediabox,trim={64.39bp 1231.84bp 64.39bp 666.92bp},clip,width=0.94\textwidth,max height=0.38\textheight]{Pope.pdf}
\end{sourceblock}
因此 p(') 和随机场 \$(x, t) 在 x 处的相关性为

\begin{sourceblock}
\includegraphics[page=469,pagebox=mediabox,trim={64.39bp 1007.87bp 64.39bp 847.89bp},clip,width=0.94\textwidth,max height=0.38\textheight]{Pope.pdf}
\end{sourceblock}
将应变波动率替换为压力-
sij 应变率张量 Rij，对于速度压力梯度张量 Uij 可以得到类似的公式。

方程 (11 .18 2) 显示压力相关性可以被视为有两个贡献：一是由于自由空间格林函数 v；二是由于自由空间格林函数 v。第二个称为反射贡献（或
r”（-l，墙墙回波）。估计这些 twd 贡献的相对幅度与 L-I 的比率是一件简单的事情，其中 L 是特征
,: - 4. << S 和 的相关长度尺度 因此，远离墙壁 (L,/y 1)，墙壁反射可以忽略不计，但当湍流长度尺度与距墙壁的距离相当时，它可能很重要，因为它在整个对数律区域中。

在一些雷诺应力闭合中，使用附加的重新分布项来考虑壁反射。通过 LRR-IP 模型，Gibson 和 Launder (1978) 提出了附加的慢速 Rij 和快速 Rij 项，前者是
其中湍流长度尺度为 L k 3/2 因子 L/y 一致
对数律区域，L/y 2.5（参见方程（7.84））。该项的作用是减少 (v2) 和 (uv) 并增加 (u2) 和 ( w Z )（参见练习 11 .26）。

SSG 模型（没有壁反射项）在 log-1,aw 区域表现良好；和 Speziale (1996) 认为，其他模型中需要壁面反射项是由于其压力-应变率率的缺陷
R,,(0.modmels(8) 和 Rij
练习 11.26 对于通常坐标系中的近壁流，写成矩阵
形成由方程给出的 Rij 的分量。 (11.183)。

对于边界层的对数律区域，估计 L / y 并
比较 R2(s), , R(S.W) 2 项的相对大小。

和-38
LRR-IP 模型。

函数 近壁区域增加了执行湍流的湍流模型计算任务的复杂性和费用。如前面小节所述，近壁效应需要对基本模型进行添加或修改；并且需要解析陡峭的剖面（例如，(U) 和 y+ 30）可能导致大部分计算工作投入到近壁区域。同时，如果平均流近似平行于壁面，则适用对数律关系，给出对数律区域中的湍流模型变量的简单代数关系。

“壁面函数”方法（Launder 和 Spalding 1972）的思想是在距壁一定距离处应用边界条件（基于对数定律关系），这样湍流模型方程就不会在靠近壁的地方求解（例如，在壁和应用边界条件的位置之间）。

\subsection*{11.7.6 墙功能}\phantomsection\addcontentsline{toc}{subsection}{11.7.6 墙功能}
现在描述通常坐标系中的流动。壁面函数边界条件应用于对数律区域中的位置 y y（例如，yf 约为 50）。下标 ' 表示在 y 处评估的数量，例如 9 ( U p k，和
对于高雷诺数零压力梯度边界层，对数定律（方程（7.43））为

\begin{sourceblock}
\includegraphics[page=470,pagebox=mediabox,trim={64.39bp 119.39bp 64.39bp 1791.77bp},clip,width=0.94\textwidth,max height=0.38\textheight]{Pope.pdf}
\end{sourceblock}
生产与消耗产量的平衡

\begin{sourceblock}
\includegraphics[page=471,pagebox=mediabox,trim={64.39bp 1708.78bp 64.39bp 215.97bp},clip,width=0.94\textwidth,max height=0.38\textheight]{Pope.pdf}
\end{sourceblock}
以及湍流黏度的 k-E 表达式（方程 (10 .48)）

\begin{sourceblock}
\includegraphics[page=471,pagebox=mediabox,trim={64.39bp 1564.80bp 64.39bp 375.95bp},clip,width=0.94\textwidth,max height=0.38\textheight]{Pope.pdf}
\end{sourceblock}
壁面函数不是直接使用这些方程，而是被设计为在所有情况下提供“鲁棒”边界条件，并在理想条件下恢复到上述关系。

方程 (11 .186) 用于将标称摩擦速度定义为

\begin{sourceblock}
\includegraphics[page=471,pagebox=mediabox,trim={64.39bp 1302.84bp 64.39bp 642.92bp},clip,width=0.94\textwidth,max height=0.38\textheight]{Pope.pdf}
\end{sourceblock}
y+ 的相应估计为

\begin{sourceblock}
\includegraphics[page=471,pagebox=mediabox,trim={64.39bp 1135.86bp 64.39bp 795.90bp},clip,width=0.94\textwidth,max height=0.38\textheight]{Pope.pdf}
\end{sourceblock}
然后从对数定律获得标称平均速度：

\begin{sourceblock}
\includegraphics[page=471,pagebox=mediabox,trim={64.39bp 968.88bp 64.39bp 920.88bp},clip,width=0.94\textwidth,max height=0.38\textheight]{Pope.pdf}
\end{sourceblock}
在平均动量方程中，不是通过指定 U p 而是通过指定剪切应力来应用边界条件：

\begin{sourceblock}
\includegraphics[page=471,pagebox=mediabox,trim={64.39bp 756.90bp 64.39bp 1134.86bp},clip,width=0.94\textwidth,max height=0.38\textheight]{Pope.pdf}
\end{sourceblock}
这是一个稳健的条件，因为剪切应力与分离点或重新附着点 ( ( U ) , 0 ) 处的速度 U p 具有相反的符号，一切都明确定义（即使 u, 为零）；如果(U)超过标称值(U)'，则-(uv)超过，从而提供“恢复力”。

根据方程。 (11 18 5，边界条件为

\begin{sourceblock}
\includegraphics[page=471,pagebox=mediabox,trim={64.39bp 439.94bp 64.39bp 1477.81bp},clip,width=0.94\textwidth,max height=0.38\textheight]{Pope.pdf}
\end{sourceblock}
同时将零法向梯度条件应用于 k 和法向应力。

壁面函数通常以比上述更复杂的方式实现，该处理被纳入数值求解过程中使用的有限体积方程中（Jones 1994）。位置 yp 被视为远离墙壁的第一个网格节点。

它们很有吸引力，并且广泛应用于商业 CFD 代码对复杂湍流的应用。然而，在许多流动条件下
.，强压力梯度，分离流和冲击流，它们的物理基础是不确定的，并且它们的准确性很差（参见，例如，Wilcox（1993））。

壁面函数引入 y 作为人工参数。对于对数律关系准确的边界层流，整体解对 y（在对数律区域内）的选择不敏感。然而，在其他flow\$中发现整体解决方案对这个选择很敏感。因此，可能无法获得数值精确的网格独立解（因为细化网格通常意味着减少 y ）

练习 11.27 k-E 或雷诺应力模型方程可通过
完全发展的槽道流的有限差分法，无需
墙功能的使用。为了提供足够的空间分辨率，网格
间距 Ay 指定为

\begin{sourceblock}
\includegraphics[page=472,pagebox=mediabox,trim={64.39bp 1029.87bp 64.39bp 895.89bp},clip,width=0.94\textwidth,max height=0.38\textheight]{Pope.pdf}
\end{sourceblock}
其中 6 为通道半高，a、b、c 为正数
常数（b > c）。表明网格点数N，

\begin{sourceblock}
\includegraphics[page=472,pagebox=mediabox,trim={64.39bp 776.90bp 64.39bp 1061.87bp},clip,width=0.94\textwidth,max height=0.38\textheight]{Pope.pdf}
\end{sourceblock}
使用关系 Re, 6/6, 0.09Reo。 t \textasciitilde{}( 请参阅第 7.1.5 页
[278] 估计高雷诺数渐近线

\begin{sourceblock}
\includegraphics[page=472,pagebox=mediabox,trim={64.39bp 531.93bp 64.39bp 1328.83bp},clip,width=0.94\textwidth,max height=0.38\textheight]{Pope.pdf}
\end{sourceblock}
充分发展的槽道流，在 y y, 处应用壁面函数。

从 y yp 延伸到 y 6 的有限差分网格有

\begin{sourceblock}
\includegraphics[page=472,pagebox=mediabox,trim={64.39bp 333.96bp 64.39bp 1552.80bp},clip,width=0.94\textwidth,max height=0.38\textheight]{Pope.pdf}
\end{sourceblock}
与 0 c b．证明所需网格点数N,,为

\begin{sourceblock}
\includegraphics[page=472,pagebox=mediabox,trim={64.39bp 159.98bp 64.39bp 1738.78bp},clip,width=0.94\textwidth,max height=0.38\textheight]{Pope.pdf}
\end{sourceblock}
\section*{11.8 椭圆松弛模型}\phantomsection\addcontentsline{toc}{section}{11.8 椭圆松弛模型}
雷诺数，* 而如果 y 被指定为
= 50) 那么 N y 随着 ln(Re) 的增加而增加。

1 1.8 椭圆松弛模型 从理论角度来看，重分布 R 的建模
(iaj ) 对于极不均匀的流动，局部量项是可疑的。来自
从实际角度来看，对于近壁流，各种阻尼函数和壁面反射项都不是很令人满意。阻尼函数本质上是修正，通常发现它不准确
其开发中使用的特定流程（例如，Rodi 和 Scheuerer (1986)）。使用距墙壁的距离 t 和法线方向是将非局部信息合并到其他局部模型中的不完美方法。在模型方程中，最好仅通过施加在墙壁上的边界条件来感知墙壁的存在。

Durbin (1993) 提出了更高级别的封闭来解决这些问题。本质的创新在于，压力-应变率率模型不是基于局部的，而是基于椭圆方程的解（类似于泊松方程6）。这和相关的（杜宾
椭圆松弛模型 1991、Dreeben 和 Pope 1997a）既不使用阻尼函数也不使用壁面反射项，在计算各种近壁流方面取得了相当大的成功。

现在描述 Durbin (1993) 的模型。起点是精确的雷诺应力方程

\begin{sourceblock}
\includegraphics[page=473,pagebox=mediabox,trim={64.39bp 637.92bp 64.39bp 1272.84bp},clip,width=0.94\textwidth,max height=0.38\textheight]{Pope.pdf}
\end{sourceblock}
（参见方程（11.138））

\begin{sourceblock}
\includegraphics[page=473,pagebox=mediabox,trim={64.39bp 467.94bp 64.39bp 1421.82bp},clip,width=0.94\textwidth,max height=0.38\textheight]{Pope.pdf}
\end{sourceblock}
要建模的主要量是重新分布项，它由两个部分组成：速度压力梯度张量 R(ea)m 的偏量部分（在壁处为零）；和之间的区别
以及简单的各向异性 (Rotta) 模型 E ( u \textasciitilde{} u \textasciitilde{} )T/h\textasciitilde{}e。耗散以这种方式分解，因此，在墙壁上，方程式中的耗散。 (11.197) 是精确的，相应地 R(iej) 为零。

事实上，求解的椭圆方程 (Eq, (11.203)) 更接近修正的亥姆霍兹方程，即
湍流传输项建模为

\begin{sourceblock}
\includegraphics[page=474,pagebox=mediabox,trim={64.39bp 1680.79bp 64.39bp 215.97bp},clip,width=0.94\textwidth,max height=0.38\textheight]{Pope.pdf}
\end{sourceblock}
其中时间尺度为 T = max(k/e,6t,)，如先前定义（参见方程（11.163））。除了粘性子层外，该传输模型与 Daly-Harlow 模型相同（方程 (11.147)）。

数量 R 定义为
这与 Rj 的基本 LRR-IP 模型相同； （方程（11.129）），除了使用 1/T 代替 ∠/k，并且使用 Cp-1 代替 Cz 来解释方程（11.129）中的各向异性耗散。 （11.197）（见练习11.3）。

在粘性子层之外（使得 T = k/s），如果局部模型

\begin{sourceblock}
\includegraphics[page=474,pagebox=mediabox,trim={64.39bp 1114.86bp 64.39bp 815.90bp},clip,width=0.94\textwidth,max height=0.38\textheight]{Pope.pdf}
\end{sourceblock}
如果使用，则所得雷诺应力模型将与基本 LRR-IP 模型相同。然而，代替方程式。 (11.201), RY 由以下方程确定

\begin{sourceblock}
\includegraphics[page=474,pagebox=mediabox,trim={64.39bp 900.89bp 64.39bp 1033.87bp},clip,width=0.94\textwidth,max height=0.38\textheight]{Pope.pdf}
\end{sourceblock}
\begin{sourceblock}
\includegraphics[page=474,pagebox=mediabox,trim={64.39bp 794.90bp 64.39bp 1152.85bp},clip,width=0.94\textwidth,max height=0.38\textheight]{Pope.pdf}
\end{sourceblock}
其中长度尺度 Lp 是从 L = k*/?/¢ 和柯尔莫哥洛夫尺度 y 获得的
其中 C, = 0.2 且 C, = 80。方程 (11.203) 是量 f;; 的椭圆松弛方程。椭圆算子由恒等算子 I 和按 Lp 缩放的拉普拉斯算子组成。在均匀湍流中（其中 V?f；；为零），或者如果 Lp 为零，则方程： (11.203) 简单地产生 f;; = R,,/k，因此局部模型方程： (11.201) 被恢复。然而，一般来说，由于方程中的拉普拉斯算子。 (11.203), fi;(x,t) 由 R;,(x’,t) 和 f;j 上施加的边界条件非局部确定。

对于非均匀流，无论是从纳维-斯托克斯方程还是从实验观察，都很难为模型提供令人信服的理由；当然没有这样的理由
压力的泊松方程 (Durbin 1991, 1993)。如果泊松方程写成十

\begin{sourceblock}
\includegraphics[page=475,pagebox=mediabox,trim={64.39bp 1683.79bp 64.39bp 256.97bp},clip,width=0.94\textwidth,max height=0.38\textheight]{Pope.pdf}
\end{sourceblock}
\$(x) 是一个随机场，那么相关性的解 (pf(x)\$(x))i s

\begin{sourceblock}
\includegraphics[page=475,pagebox=mediabox,trim={64.39bp 1503.81bp 64.39bp 396.95bp},clip,width=0.94\textwidth,max height=0.38\textheight]{Pope.pdf}
\end{sourceblock}
参见等式。 （11 .179）和（11 .182）。 （积分是在流域上进行的，并且添加谐波分量 (P(\textasciitilde{})(x)\$(xc)a) n 以满足适当的边界条件。）现在假设两点相关性可以为

\begin{sourceblock}
\includegraphics[page=475,pagebox=mediabox,trim={64.39bp 1236.84bp 64.39bp 711.91bp},clip,width=0.94\textwidth,max height=0.38\textheight]{Pope.pdf}
\end{sourceblock}
然后，等式。 (11.206) 变为

\begin{sourceblock}
\includegraphics[page=475,pagebox=mediabox,trim={64.39bp 1046.87bp 64.39bp 846.89bp},clip,width=0.94\textwidth,max height=0.38\textheight]{Pope.pdf}
\end{sourceblock}
当它乘以大括号 (\{ \}) 中的项时，就是方程中运算符的格林函数。 (11.203)。因此，压力泊松方程和椭圆松弛方程之间建立了松散的联系。

也许比椭圆松弛方程的具体形式更重要的是它允许应用额外的边界条件。在局部雷诺应力模型中，每个雷诺应力只能在壁上应用一个边界条件 ( e (t iiuj) O)，然后雷诺应力随距离的渐近变化由模型确定。另一方面，在椭圆松弛模型中，可以将两个边界条件应用于每个 (uiuj) 和 f 对。附加边界条件可用于确保正确的渐近雷诺应力变化。

使用椭圆松弛的雷诺应力模型已成功用于计算具有不利压力梯度和凸曲率的边界层（Durbin 1993）。对于渠道流（DNS 数据可用），雷诺应力预算的再现相当准确（Dreeben 和 Pope 1997a）。

具有椭圆松弛的 k-E-v2 模型 在椭圆松弛雷诺应力模型之前，Durbin (1991) 介绍了
考虑到，v2是(v2)的模型，并且使用湍流黏度模型

\begin{sourceblock}
\includegraphics[page=476,pagebox=mediabox,trim={64.39bp 1682.79bp 64.39bp 201.97bp},clip,width=0.94\textwidth,max height=0.38\textheight]{Pope.pdf}
\end{sourceblock}
（参见方程（11.158）和图11.17）。

该模型也已成功应用于许多复杂的流（例如，Durbin (1995) 和 Parneix、Durbin 和 Behnia (1998a)）。那里
被解释为一个标量，它与墙壁相邻，是法向速度方差的模型，即 v2 (tdiu,)ninj。

EX E R C I S E 11.29 作为椭圆松弛方程的推广（方程（11.203）），
考虑方程

\begin{sourceblock}
\includegraphics[page=476,pagebox=mediabox,trim={64.39bp 1140.86bp 64.39bp 736.91bp},clip,width=0.94\textwidth,max height=0.38\textheight]{Pope.pdf}
\end{sourceblock}
Durbin 模型的指数为 p = q - 0，而 Wizman 等人。

(1996) 建议使用 p 0 和 q 2 的“中性”模型，以及另一个
具有 p -2 和 q 2 的模型。考虑中的对数律区域
通常的坐标-，和a-ss ume，湍流统计是自
相似（与 LD y、k yo、Rij） 4' 和 Y ）表明
方程的解(11 .210) 是-

\begin{sourceblock}
\includegraphics[page=476,pagebox=mediabox,trim={64.39bp 756.90bp 64.39bp 1135.86bp},clip,width=0.94\textwidth,max height=0.38\textheight]{Pope.pdf}
\end{sourceblock}
评论一下上面提到的三个模型的效果。

\section*{11.9 代数应力与非线性黏度模型}\phantomsection\addcontentsline{toc}{section}{11.9 代数应力与非线性黏度模型}
11.9.1 通过引入输运项的近似值，雷诺应力模型可以简化为一组代数方程。这些方程形成了代数应力模型 (ASM)，该模型隐式地将雷诺应力（局部）确定为 k 和平均速度梯度的函数。由于涉及近似，代数应力模型本质上不如雷诺应力模型通用且准确。然而，由于它们相对简单，它们已被用作湍流模型（与 k 和 的模型方程结合）此外，

\subsection*{11.9.1 代数应力模型}\phantomsection\addcontentsline{toc}{subsection}{11.9.1 代数应力模型}
它也可用于获得非线性湍流黏度模型。

标准模型雷诺应力传递方程为

这是一组六个偏微分方程的耦合组。右侧的项是 (Ui) /axj, (uiuj) 的局部代数函数，并且它们不涉及雷诺应力的导数。在代数应力模型中，输运项 V（位于方程 (11 .212) 的左侧）由代数表达式近似，因此整个方程变为代数方程。具体来说，等式。 (11.212) 成为一组六个代数方程，隐式地将雷诺应力确定为 k 和平均速度梯度的函数。

在某些情况下（例如，高雷诺数完全发展的槽道流的对数律区域），方程中的传输项（11.212）可以忽略不计，因此（在某种意义上）雷诺应力与施加的速度梯度处于局部平衡。然而，完全忽略运输条件是不一致的，除非 P/F: 是统一的，因为一半的痕迹

\begin{sourceblock}
\includegraphics[page=477,pagebox=mediabox,trim={64.39bp 930.88bp 64.39bp 970.88bp},clip,width=0.94\textwidth,max height=0.38\textheight]{Pope.pdf}
\end{sourceblock}
Rodi (1972) 引入了更一般的弱均衡假设。雷诺应力可分解为

\begin{sourceblock}
\includegraphics[page=477,pagebox=mediabox,trim={64.39bp 728.91bp 64.39bp 1176.85bp},clip,width=0.94\textwidth,max height=0.38\textheight]{Pope.pdf}
\end{sourceblock}
因此 (uiuj) 的空间和时间变化可以被认为是由于 k 和 bij 的变化造成的。在弱平衡假设中，(uiuj)/k（或等效的 bij）的变化被忽略，但由于 k 的变化而导致的 (uiuj) 的变化被保留。对于平均对流项，这导致近似

\begin{sourceblock}
\includegraphics[page=477,pagebox=mediabox,trim={64.39bp 321.96bp 64.39bp 1598.80bp},clip,width=0.94\textwidth,max height=0.38\textheight]{Pope.pdf}
\end{sourceblock}
同样的近似适用于整个运输期限收益率

\begin{sourceblock}
\includegraphics[page=477,pagebox=mediabox,trim={64.39bp 152.98bp 64.39bp 1754.78bp},clip,width=0.94\textwidth,max height=0.38\textheight]{Pope.pdf}
\end{sourceblock}
在模型雷诺应力方程（方程（11.212））中使用弱平衡假设（方程（11.216））得出代数应力模型
这包括五个独立的代数方程（因为轨迹不包含信息），可用于确定 k 的 (uiuj)/k（或等效的 bpe）i 和 a(Ui)/axj。泽杰
作为 ASM 可以提供的见解的一个示例，如果 Rij 由 LRR-IP 模型给出，则等式： (11.217) 可以被操纵以产生

\begin{sourceblock}
\includegraphics[page=478,pagebox=mediabox,trim={64.39bp 1336.83bp 64.39bp 582.93bp},clip,width=0.94\textwidth,max height=0.38\textheight]{Pope.pdf}
\end{sourceblock}
参见练习 11.20。因此，可以看出该模型的含义是雷诺应力各向异性与生产各向异性成正比。

对于简单的剪切流，方程： (11.218) 很容易求解得到
PIE 各向异性 bij 作为 的函数（参见练习 11.20）：这些被绘制在

\begin{sourceblock}
\includegraphics[page=478,pagebox=mediabox,trim={64.39bp 789.90bp 64.39bp 1109.86bp},clip,width=0.94\textwidth,max height=0.38\textheight]{Pope.pdf}
\end{sourceblock}
用于定义 C，则可以推断（练习 11.31）ASM

\begin{sourceblock}
\includegraphics[page=478,pagebox=mediabox,trim={64.39bp 581.93bp 64.39bp 1255.84bp},clip,width=0.94\textwidth,max height=0.38\textheight]{Pope.pdf}
\end{sourceblock}
因此，如图 11.21 所示，LRR-IP 模型暗示的 C 值随 P I E 减小，对应于“剪切稀化”行为 C，随着剪切力 S k 的增加而减小
对于平均流对流，弱平衡假设（方程（11 .215））相当于

\begin{sourceblock}
\includegraphics[page=478,pagebox=mediabox,trim={64.39bp 246.97bp 64.39bp 1662.79bp},clip,width=0.94\textwidth,max height=0.38\textheight]{Pope.pdf}
\end{sourceblock}
其中 bij 是雷诺应力各向异性张量 b 的分量，称为惯性系。方程中体现的部分假设。 (11.221)

\begin{sourceblock}
\includegraphics[page=479,pagebox=mediabox,trim={64.39bp 1190.85bp 64.39bp 109.44bp},clip,width=\textwidth,max height=0.72\textheight]{Pope.pdf}
\par\vspace{0.45\baselineskip}
\small 图 11.20。根据 LRR-IP 代数应力模型，雷诺应力各向异性是 P/e 的函数。虚线根据 k-e 模型显示。
\end{sourceblock}
\begin{sourceblock}
\includegraphics[page=479,pagebox=mediabox,trim={64.39bp 415.95bp 64.39bp 864.89bp},clip,width=\textwidth,max height=0.72\textheight]{Pope.pdf}
\par\vspace{0.45\baselineskip}
\small 图 11.21。 C 值作为 P/e 的函数，由 LRR-IP 代数应力模型（方程（11.220））给出。
\end{sourceblock}
框架，遵循平均流）。对于某些流动，例如那些具有显着平均流线曲率的流动，Girimaji (1997) 认为更好的假设是 b 的分量相对于特定的值是固定的。

\begin{center}\small 表 11.5。方程中的系数。 (11.223) 适用于各种型号\end{center}
\begin{sourceblock}
\includegraphics[page=480,pagebox=mediabox,trim={64.39bp 1345.83bp 64.39bp 216.97bp},clip,width=\textwidth,max height=0.72\textheight]{Pope.pdf}
\end{sourceblock}
考虑基于 LRR-IP 的代数应力模型应用于
简单的剪切流（方程 1 1 130-(I 1 1 3 4 ) 并获得表达式
对于极限 P / E a.A 中的 bij，这些值可以实现吗？他们是吗
与 RDT 一致吗？

操纵方程。 (11.219) 获得

\begin{sourceblock}
\includegraphics[page=480,pagebox=mediabox,trim={64.39bp 1028.87bp 64.39bp 895.89bp},clip,width=0.94\textwidth,max height=0.38\textheight]{Pope.pdf}
\end{sourceblock}
因此使用方程。 (11 .13 4) 验证 C 方程的表达式。 (11.220)。

考虑压力-应变率率张量的通用模型：

在 bij 和平均速度梯度中是线性的，即方程： (11.135)

其中 f 为 0。表明相应的代数应力
模型（方程（11.217））可以写成

\begin{sourceblock}
\includegraphics[page=480,pagebox=mediabox,trim={64.39bp 676.91bp 64.39bp 1202.85bp},clip,width=0.94\textwidth,max height=0.38\textheight]{Pope.pdf}
\end{sourceblock}
其中系数（不太一般，对于 LRR 模型）是

表 11.5 给出。

\subsection*{11.9.2 非线性湍流黏度}\phantomsection\addcontentsline{toc}{subsection}{11.9.2 非线性湍流黏度}
代数应力模型方程(11.217) 是一个方程
隐式 (uiuj)/k，或等效于各向异性 bij。显然，获得以下形式的关系是有好处的

\begin{sourceblock}
\includegraphics[page=480,pagebox=mediabox,trim={64.39bp 207.97bp 64.39bp 1728.78bp},clip,width=0.94\textwidth,max height=0.38\textheight]{Pope.pdf}
\end{sourceblock}
其中 S 和 是归一化平均应变率和旋转张量

\begin{center}\small 表 11.6。完整的独立、对称、偏函数集）。。。 “\textasciitilde{} . . \textasciitilde{} 偏对称张量 S 和反对称张量 Q 的 IT，以矩阵表示法表示：大括号表示迹，例如 = SiSji\end{center}
\begin{sourceblock}
\includegraphics[page=481,pagebox=mediabox,trim={64.39bp 1292.84bp 64.39bp 309.96bp},clip,width=\textwidth,max height=0.72\textheight]{Pope.pdf}
\end{sourceblock}
形式的最一般可能的表达 Eq. (11a22r4) 可以是

\begin{sourceblock}
\includegraphics[page=481,pagebox=mediabox,trim={64.39bp 1100.86bp 64.39bp 741.91bp},clip,width=0.94\textwidth,max height=0.38\textheight]{Pope.pdf}
\end{sourceblock}
其中张量 7A 在表 11 .6 中给出，系数可以取决于
一个aA;,-. n 基于五个不变量 S 2 -S3 + -C J；?iSjia，nd C22mS（TP ope 1975）。像bij一样，每个
张量 7 是无量纲、对称且偏量的。作为一个集合，它们形成了一个含义，即每个对称的偏差秒-
完整性基础，由 S 形成的阶张量，可以表示为这十个的线性组合。 （这一点的证明基于 Cayley-Hamilton 定理，Pope (1975)。）

指定 G = -C P，对于 n 1，G(") = O，方程： (11 .2-25)r 转换为线性 k-E 湍流黏度公式

\begin{sourceblock}
\includegraphics[page=481,pagebox=mediabox,trim={64.39bp 422.95bp 64.39bp 1414.82bp},clip,width=0.94\textwidth,max height=0.38\textheight]{Pope.pdf}
\end{sourceblock}
当 n > 1 时，G(") 的一个重要规范会产生一个非线性湍流黏度模型，即，一个显式的非线性公式
( u i u j ) 平均速度梯度。

对于统计上二维的流，情况是考虑-
--(I)\textasciitilde{} ( 2 ) (3) 简单得多。张量 7 、 7 和 7 形成完整性基础，并且
gzk 只有两个独立的不变量和\#--。 （教皇 1975 年，加茨基和
7 可以吸收到修改后的压力（参见练习 11.33）项中，以便 G(3) 的值对平均速度场没有影响。与
-- 0，统计二维流动的非线性粘度模型 \textasciitilde{} ( j : )

\begin{sourceblock}
\includegraphics[page=482,pagebox=mediabox,trim={64.39bp 1440.82bp 64.39bp 294.96bp},clip,width=0.94\textwidth,max height=0.38\textheight]{Pope.pdf}
\end{sourceblock}
获得系数 G(") 的合适规格的一种方法是从代数应力模型中获得。由于非线性湍流黏度公式 Eq.. (11.225) 提供了关于平均速度梯度的 bij 的完全通用表达式，因此每个代数应力模型都对应一个非线性粘度模型。确定相应系数 G(") 是一个代数问题。例如，对于统计二维流，LRR-IP代数应力模型对应的系数G(")为

\begin{sourceblock}
\includegraphics[page=482,pagebox=mediabox,trim={64.39bp 1008.87bp 64.39bp 930.88bp},clip,width=0.94\textwidth,max height=0.38\textheight]{Pope.pdf}
\end{sourceblock}
\begin{sourceblock}
\includegraphics[page=482,pagebox=mediabox,trim={64.39bp 726.91bp 64.39bp 1155.85bp},clip,width=0.94\textwidth,max height=0.38\textheight]{Pope.pdf}
\end{sourceblock}
参见练习 11.36。图 11.22 显示了 -G C 和 -G i C 作为 S k / \& 和 SZk/\& 的函数（其中 SZ (2QijQij)

由方程定义的非线性粘度模型。 (11.230)-(11.232) 并不完全明确，因为 ic 的定义包含 P I E -2bijS A ij。 Girimaji (1996) 给出了完全明确的公式，通过求解 A 的三次方程获得，请参见练习 11.36。 Taulbee (1992) 以及 Gatski 和 Speziale (1993) 将这种方法扩展到三维流，其中，一般来说，所有十个系数 G 都不为零。

非线性粘度模型不是基于代数应力模型，由 Yoshizawa (1984)、Speziale (1987)、Rubinstein 和 Barton (1990)、Craft、Launder 和 Suga (1996) 等人提出。前三个提到的平均速度梯度是二次的，因此 G(') G(4) 不为零。在 Craft 等人的模型中。 G (5) 也不为零。除了平均速度
(a) C, = G 和 (b) -G(\textasciitilde{}),f 或 LRR-IP 非线性的等高线图
粘度模型（1
考虑统计上的二维湍流

平面（因此 0 和 0）。通过评估每个
= (Ui)/ax3 =
组件，表明

\begin{sourceblock}
\includegraphics[page=483,pagebox=mediabox,trim={64.39bp 1005.87bp 64.39bp 927.88bp},clip,width=0.94\textwidth,max height=0.38\textheight]{Pope.pdf}
\end{sourceblock}
\begin{sourceblock}
\includegraphics[page=483,pagebox=mediabox,trim={64.39bp 845.89bp 64.39bp 1059.87bp},clip,width=0.94\textwidth,max height=0.38\textheight]{Pope.pdf}
\end{sourceblock}
因此表明

\begin{sourceblock}
\includegraphics[page=483,pagebox=mediabox,trim={64.39bp 666.92bp 64.39bp 1265.84bp},clip,width=0.94\textwidth,max height=0.38\textheight]{Pope.pdf}
\end{sourceblock}
\begin{sourceblock}
\includegraphics[page=483,pagebox=mediabox,trim={64.39bp 533.93bp 64.39bp 1400.82bp},clip,width=0.94\textwidth,max height=0.38\textheight]{Pope.pdf}
\end{sourceblock}
I - n 非线性粘度模型（方程（11.225）），贡献来自
平均动量方程中的 I 为

\begin{sourceblock}
\includegraphics[page=483,pagebox=mediabox,trim={64.39bp 312.96bp 64.39bp 1595.80bp},clip,width=0.94\textwidth,max height=0.38\textheight]{Pope.pdf}
\end{sourceblock}
对于正在考虑的二维流，表明这可以
重新表达为

\begin{sourceblock}
\includegraphics[page=483,pagebox=mediabox,trim={64.39bp 119.39bp 64.39bp 1805.77bp},clip,width=0.94\textwidth,max height=0.38\textheight]{Pope.pdf}
\end{sourceblock}
（这项可以被吸收到修改后的压力中，以便 G
对平均速度场没有影响。）

计算 7 的简单剪切流，其中 a ( \textasciitilde{} \textasciitilde{} ) / a xS \textasciitilde{} 是
仅非零平均速度梯度。因此表明（对于这个
flow) 由非线性粘度模型给出的剪切应力 (u1u2)

（方程（11 .229））取决于 G 但独立于 G (2)。

考虑一个压力-应变率率模型，该模型在
平均速度梯度和各向异性。相应的
代数应力模型为 Eq. (11.223)。对于湍流
统计二维，相应的非线性粘度
模型由 (Pope 1975) 给出

\begin{sourceblock}
\includegraphics[page=484,pagebox=mediabox,trim={64.39bp 1176.85bp 64.39bp 676.91bp},clip,width=0.94\textwidth,max height=0.38\textheight]{Pope.pdf}
\end{sourceblock}
其中 g、yl 和 的定义见表 11.5。使用此结果
验证方程(11.230)。

通过使用关系式 P/E -2bVs 可知 il (Eq. (11.231))

满足三次方程。

Rubinstein 和 Barton (1990) 的模型可以写成

\begin{sourceblock}
\includegraphics[page=484,pagebox=mediabox,trim={64.39bp 606.92bp 64.39bp 1290.84bp},clip,width=0.94\textwidth,max height=0.38\textheight]{Pope.pdf}
\end{sourceblock}
其中 Q 的定义使得右侧的迹消失，
常数为CT1 0.034、CT2 0.104、CT3 -0.014。

以方程的形式重新表达该模型。 (11.225) 并确定
工艺等人。 (1996) 提出了一个非线性粘度模型，涉及
张量 7 以及

\begin{sourceblock}
\includegraphics[page=484,pagebox=mediabox,trim={64.39bp 224.97bp 64.39bp 1699.79bp},clip,width=0.94\textwidth,max height=0.38\textheight]{Pope.pdf}
\end{sourceblock}
用 ?;, 重新表达 7 A'。

\section*{11.10 讨论}\phantomsection\addcontentsline{toc}{section}{11.10 讨论}
\begin{center}\small 表 11.7。各种RANS湍流模型的属性（前四个是湍流黏度模型）\end{center}
\clearpage
\begin{sourceblock}
\includegraphics[page=485,pagebox=mediabox,trim={64.39bp 119.39bp 64.39bp 265.97bp},clip,width=\textwidth,max height=0.72\textheight]{Pope.pdf}
\end{sourceblock}
到目前为止提到的模型（表 11.7 中的前四个）是湍流黏度模型。最重要的雷诺应力张量不是直接表示的，而是建模为 k 和 (U+)/dx 的局部函数。然而，这是通过各向同性或非线性粘度模型或代数应力模型来完成的，描述级别是相同的；并且 (uIu . ) 由 Ui)/ axj 局部确定的假设是不可避免的。

在雷诺应力模型中，雷诺应力的直接表示消除了湍流黏度模型的这种内在假设，即 (upu -) 由 d (Ui) xi 局部确定。在雷诺应力方程中，平均对流 D ( u i u j )/ D t 由四个过程产生、耗散、再分布和湍流传输来平衡，其中前三个过程通常占主导地位（且可比较）。在这个闭合水平上，生产和耗散（至少其各向同性部分）处于闭合形式，因此要建模的主要量是由于波动压力场 Rij 引起的重新分配。

在强不均匀性区域，特别是靠近墙壁的区域，Rij 局部确定（由 (uiuj) 和 d(Ui)/axj）的假设是值得怀疑的。脉动压力场 p' 由泊松方程控制，因此可能会产生显着的非局部效应，例如壁阻塞。椭圆松弛模型中的更高级别的描述允许将非局部效应纳入重新分布的建模中（如 R！” kf ij，方程（11 .202））。

EJ 特别是在墙壁附近，R( e) 受到施加在墙壁上的边界条件的强烈影响。

列出的最后三个模型，建模中使用的长度尺度和时间尺度是从（通过 L = k 3/2 和 z = k ）获得的。在某些情况下，例如，当湍流近似自相似时
单一尺度可能足以对湍流状态进行参数化。然而，一般来说，特别是对于快速变化的流动，可能需要多个参数来表征不同尺度的湍流运动的状态和响应。 Hanjalik、Launder 和 Schiestel (1980)、Wilcox (1988) 和 Lumley (1992) 等人提出了引入附加尺度的模型。这些模型尚未得到广泛应用。

在什么情况下，或者以什么假设或选择
参数，下面列表中的每个模型是否都简化为它的成功
(b) 雷诺应力模型。

(c) 代数应力模型（使用 k-E）。

(d) 非线性粘度模型（使用 k-E）。

(e) 标准（各向同性粘度）k-E 模型。

(f) 一方程 k 模型。

(g) 混合长度模型。

混合长度和一方程 k-ern 模型并不完整，因为它们需要 Irn(xt,) 的规范。其余模型均已完成。

Spalart-Allmaras 模型涉及到墙壁的距离，其他模型中的一些近墙处理涉及 I，而正常情况
由于 I 和 可以针对一般流程明确定义（例如，方程（11171）），因此它们的使用不会使模型不完整。

成本和易用性 在本次讨论中，我们以 k-E 模型作为参考：该模型被纳入大多数商业 CFD 代码中，并且当与壁面函数结合使用时，通常认为它易于使用且计算成本低廉。

如果不采用壁面函数，则对粘性近壁区域执行 k-E 计算的任务将明显更加困难和昂贵。这是因为需要求解 k 和 （在近壁区域变化很大）；而且这些方程中的源项在靠近墙的地方变得非常大。 （在对数律区域中，术语 CEZ2\&k 变化为 u4/y2。）

对于近壁空气动力学流动，Spalart-Allmaras 模型的设计要简单得多且成本更低。这是因为 与 k 相比，湍流黏度在近壁区域表现良好，并且更容易求解。

与 k-E 模型相比，雷诺应力模型有些困难且成本较高，因为
(i) 一般而言，需要求解七个湍流方程（对于 (uiuj)

而不是两个（对于 k 和 E），.

(ii) 模型雷诺斯特勒方程要复杂得多
在平均动量方程中，术语

\begin{sourceblock}
\includegraphics[page=488,pagebox=mediabox,trim={64.39bp 1708.78bp 64.39bp 238.97bp},clip,width=0.94\textwidth,max height=0.38\textheight]{Pope.pdf}
\end{sourceblock}
导致流量和流量之间的数值耦合不太有利

湍流方程比相应项

在 k-E 模型中。

通常，雷诺应力模型计算所需的 CPU 时间可能比 k-E 计算多出两倍。

使用代数应力模型的主要动机是避免求解雷诺应力模型方程的成本和难度。然而，普遍的经验是这些好处并没有实现。代数应力模型方程是耦合非线性方程，通常具有多个根，从经济角度求解这些方程并不简单。另外，对于上述第(iii)项，代数应力模型具有与雷诺应力模型相同的缺点。正如上一节所讨论的，代数应力模型可以改写为非线性粘度模型。这些几乎不会增加 k-c-模型计算的成本和难度。

椭圆松弛模型引入的另外六个字段 fij(x,t) 显然会导致额外的计算成本。椭圆方程（Eq. (11 .203)）的求解并不是特别困难或昂贵，但它的结构与 CFD 代码通常用来求解的对流扩散方程不同。

适用范围 基本的 k-E 和雷诺应力模型可应用于任何湍流。它们还提供可用于附加过程建模的长度尺度和时间尺度信息。因此，它们为湍流反应流、多相流等的建模提供了基础。标量通量的模型输运方程可以与雷诺应力模型结合求解，以提供平均标量方程的闭合。这种所谓的二阶矩闭合已成功扩展到浮力效应显着的大气流动（例如，Zeman 和 Lumley（1976））。尽管原则上它可以应用于任何湍流（在所考虑的类别中），但 Spalart-Allmaras 模型仅适用于
准确性 评估其准确性的理想方法是比较各种湍流模型在各种测试流中的性能。外部误差（图 8.2）必须很小；测试流程应与模型开发中使用的流程不同；并且测试不应由模型开发人员执行。尽管它们在所有方面都不符合这一理想，但以下著作对于评估各种湍流模型的准确性很有用：Bradshaw、Launder 和 Lumley (1996)、Luo 和 Lakshminarayana (1997)、Kral、Mani 和 Ladd (1996)、Behnia、Parneix 和 Durbin (1998)、Godin 等人。 (1997)、Menter (1994) 以及 Kline、Cantwell 和 Lilley (1982)。 Hanjalik (19 94) 和 Wilcox (19 93) 的评论和评估也很有价值。

引用的作品支持以下结论。

(i) k-E 模型对于二维薄层表现相当好
she a \textasciitilde{}流动，其中平均流线曲率和平均压力
梯度较小。

(ii) 对于具有强压力梯度的边界层，k-E 模型

表现不佳。然而，k-co 模型的表现令人满意，
事实上，它的性能对于许多流程来说都是优越的。

(iii) 对于远离简单剪切的流动（例如，撞击射流和
三维流），k-E 模型可能会严重失败。

(iv) 使用非线性粘度模型是有益的并且允许
二次流量的计算（不能使用
各向同性粘度假设）。

(v) 雷诺应力模型可以成功（而湍流黏度模型
模型不是）在计算具有显着平均流线的流量时
曲率、强旋流或平均旋转流、二次流
在管道中，并且平均流量快速变化。

(vi) 雷诺应力模型计算对以下细节敏感
压力-应变率率张量的建模，包括壁面
反射术语。

(vii) 椭圆松弛模型（雷诺应力和 k-E-v2）有
在应用于许多具有挑战性的两方面方面取得了相当成功

维度流，包括撞击射流和分离的束缚
芳层。 (viii) 耗散方程经常被指责为性能不佳的原因
一个模型的。对于许多 Eow 来说，可以大大提高性能
修正条款。不对耗散方程进行修正
已发现在所有流程中均有效。

总之，特别是对于复杂流动，雷诺应力模型已被证明优于二方程模型。直到 1990 年左右，由于所使用的模型存在缺陷，而且对复杂流程的测试很少，这种优势还没有得到很好的确立。自从 Bradshaw（1987）写道“最好的现代方法使得几乎所有流量的计算精度都高于最明智的猜测”以来，RANS 模型已经有了相当大的改进，这意味着即使这些方法无法取代实验，它们也是真正有用的。

\chapter{PDF 方法}
平均速度 ( U(x, t)) 和雷诺应力 (uiuj) 是速度 f ( V ;x , t ) 的欧拉 PDF 的一阶矩和二阶矩（方程 (3.153)）。在 PDF 方法中，求解模型传输方程以获得 PDF，例如
f ( V ;x , t) 的精确 t.transport 方程源自附录 H 中的纳维-斯托克斯方程，并在 12.1 节中进行了讨论。在这个方程中，所有对流传输都是封闭形式，与平均动量方程中的 (uiuj) xi 项和 / a x i 相比
( U \textasciitilde{} Uj \textasciitilde{} k ) 雷诺应力方程。第 12.2 节给出了基于广义 Langevin 模型 (GLM) 的 PDF 闭合模型方程，并显示了它与压力-应变率率张量 Rij 模型的密切关系。

PDF 方法的核心是随机拉格朗日模型，它涉及新概念并需要额外的数学工具。附录 J 中给出了扩散过程和随机微分方程的必要背景。最简单的随机拉格朗日模型是朗之万方程，它提供了跟随流体粒子的速度模型。该模型在 12.3 节中进行了介绍和研究。

闭合不能仅基于速度的 PDF，因为该 PDF 不包含有关湍流时间尺度的信息。获得闭合的一种方法是用模型耗散方程补充 PDF 方程。第 12.5 节中描述的一种更好的方法是考虑速度和湍流频率的联合 PDF。

在实践中，模型 PDF 方程是通过拉格朗日粒子方法（第 12.6 节）求解的，拉格朗日粒子方法本身提供了有价值的见解。更多模型（GLM 之外）在第 12.7 节中描述。

\section*{12.1 速度的欧拉 PDF}\phantomsection\addcontentsline{toc}{section}{12.1 速度的欧拉 PDF}
\subsection*{12.1.1 定义和属性}\phantomsection\addcontentsline{toc}{subsection}{12.1.1 定义和属性}
如方程式所定义。 (3.153)，f ( V , x, t) 是速度 U(x,t) 的单点、一次性欧拉 PDF，其中 V \{ V , V ', V3) 是样本空间速度空间中的自变量。整个速度的积分
JP ) dV1 dV2d V3 缩写为 ) d V，因此空格 -.= (
Jm J m 归一化条件写十
对于任何函数 Q( U(x, t))，其均值（或期望）由下式定义

\begin{sourceblock}
\includegraphics[page=492,pagebox=mediabox,trim={64.39bp 1251.84bp 64.39bp 646.92bp},clip,width=0.94\textwidth,max height=0.38\textheight]{Pope.pdf}
\end{sourceblock}
因此平均速度和雷诺应力由下式给出

\begin{sourceblock}
\includegraphics[page=492,pagebox=mediabox,trim={64.39bp 1086.86bp 64.39bp 816.90bp},clip,width=0.94\textwidth,max height=0.38\textheight]{Pope.pdf}
\end{sourceblock}
和（带有缩写符号）

\begin{sourceblock}
\includegraphics[page=492,pagebox=mediabox,trim={64.39bp 922.88bp 64.39bp 977.88bp},clip,width=0.94\textwidth,max height=0.38\textheight]{Pope.pdf}
\end{sourceblock}
PDF 的可实现性问题比雷诺应力的可实现性问题简单得多。为了使 f ( V ;x , t) 成为可实现的 PDF，它对于所有 V 都必须是非负的，并且满足归一化条件（式 1）。 （12.1）。此外，由于动量和能量是有界的，因此 f ( V ; x, t) 必须使得 ( U ) 和 *) （方程 (12.3) 和 (12.4)) 是有界的。

在精确的 PDF 输运方程中，“未知数”表现为条件期望，现在对其定义和属性进行回顾。令 \#(x, 为随机场，并令 f u+(V , y, ;x , t) 为 U(x, t ) 和 \#(x, t) 的（单点、一次性）欧拉联合 PDF，其中 y 是对应于 的样本空间变量。以 U(x, t) = V 为条件的 \#(x, t) 的 PDF 定义为

\begin{sourceblock}
\includegraphics[page=492,pagebox=mediabox,trim={64.39bp 352.96bp 64.39bp 1590.80bp},clip,width=0.94\textwidth,max height=0.38\textheight]{Pope.pdf}
\end{sourceblock}
（参见等式（3.95））；的条件均值定义为

\begin{sourceblock}
\includegraphics[page=492,pagebox=mediabox,trim={64.39bp 169.98bp 64.39bp 1734.78bp},clip,width=0.94\textwidth,max height=0.38\textheight]{Pope.pdf}
\end{sourceblock}
获得为

\begin{sourceblock}
\includegraphics[page=493,pagebox=mediabox,trim={64.39bp 1530.81bp 64.39bp 219.97bp},clip,width=0.94\textwidth,max height=0.38\textheight]{Pope.pdf}
\end{sourceblock}
\subsection*{12.1.2 PDF传输方程}\phantomsection\addcontentsline{toc}{subsection}{12.1.2 PDF传输方程}
f ( V ; x ,t ) 的精确传输方程在附录 H 中导出。该方程可以用多种信息方式表示，其中两种是

\begin{sourceblock}
\includegraphics[page=493,pagebox=mediabox,trim={64.39bp 1170.85bp 64.39bp 725.91bp},clip,width=0.94\textwidth,max height=0.38\textheight]{Pope.pdf}
\end{sourceblock}
\begin{sourceblock}
\includegraphics[page=493,pagebox=mediabox,trim={64.39bp 1019.87bp 64.39bp 927.88bp},clip,width=0.94\textwidth,max height=0.38\textheight]{Pope.pdf}
\end{sourceblock}
第三个是方程式。 （H.25）第 706 页。

左侧的项代表变化率和对流率。它们仅涉及方程的主题 f 和自变量 x、V 和 t：因此这些项呈封闭形式。

第一个方程，方程。 (12.8) 是一个数学恒等式，因此不包含任何物理学（除了假定不可压缩性之外）。然而，它表明决定 PDF 演化的基本量是条件加速度 (DUIDtl V)。当使用纳维-斯托克斯方程将加速度 (DU IDt) 替换为引起加速度的特定力 ( v V2U Vplp) 时，物理学就进入了。这与压力的雷诺分解一起得出第二个方程，Eq. （12.9）。

平均压力梯度项也是封闭形式，即方程的第一项或右侧“侧”。 （12.9）。 （这是因为；给定 f ( V ; x , t )，平均速度和雷诺应力场是已知的（从方程 (12.3) 和 (2. 4) ）；因此 (p) 的泊松方程中的源（方程 (4.13)）是已知的；因此 d(p)/axi 由泊松方程的解确定。）

练习 12.3 和 12.4 说明了 f 矩的演化方程（例如，afid 可以从 PDF 中获得）

形式的数量

\begin{sourceblock}
\includegraphics[page=494,pagebox=mediabox,trim={64.39bp 1718.78bp 64.39bp 229.97bp},clip,width=0.94\textwidth,max height=0.38\textheight]{Pope.pdf}
\end{sourceblock}
在此处遇到的所有情况下，这种形式的量都为零：如练习 12.1 所示，Qij 为零的充分条件是均值 Al) 存在。

Qij 由方程定义。 ( 1 2 1 0 ) 用高斯定理证明

\begin{sourceblock}
\includegraphics[page=494,pagebox=mediabox,trim={64.39bp 1328.83bp 64.39bp 574.93bp},clip,width=0.94\textwidth,max height=0.38\textheight]{Pope.pdf}
\end{sourceblock}
其中 ( )dS(V)d 表示在球体表面上的积分
速度空间中的半径 V。由此得到

\begin{sourceblock}
\includegraphics[page=494,pagebox=mediabox,trim={64.39bp 1149.85bp 64.39bp 783.90bp},clip,width=0.94\textwidth,max height=0.38\textheight]{Pope.pdf}
\end{sourceblock}
其中 \}s(v)d 表示半径 V 球体上的面积平均值。

表明 (JAl) 由下式给出

\begin{sourceblock}
\includegraphics[page=494,pagebox=mediabox,trim={64.39bp 943.88bp 64.39bp 955.88bp},clip,width=0.94\textwidth,max height=0.38\textheight]{Pope.pdf}
\end{sourceblock}
因此 (JAl) 的存在是 Qij 的充分条件
为零。

获得总体结果

\begin{sourceblock}
\includegraphics[page=494,pagebox=mediabox,trim={64.39bp 725.91bp 64.39bp 1221.85bp},clip,width=0.94\textwidth,max height=0.38\textheight]{Pope.pdf}
\end{sourceblock}
其中 Bk 表示导数

\begin{sourceblock}
\includegraphics[page=494,pagebox=mediabox,trim={64.39bp 500.94bp 64.39bp 1356.83bp},clip,width=0.94\textwidth,max height=0.38\textheight]{Pope.pdf}
\end{sourceblock}
获得平均动量方程。表明这与
等式； （4.12）（第 85 页）。同样的过程会得到什么结果？

应用于方程。 （12.8）？

精确的雷诺应力方程（对于 ,(ujuk)）可以通过以下公式导出
乘法方程(12.9) 通过并在速度空间上积分，
其中定义为

\begin{sourceblock}
\includegraphics[page=494,pagebox=mediabox,trim={64.39bp 169.98bp 64.39bp 1775.78bp},clip,width=0.94\textwidth,max height=0.38\textheight]{Pope.pdf}
\end{sourceblock}
(a) 项 a f / a t 导致 a ( u j u \& ) / a t；

(b) 对流项可以重新表示为

\begin{sourceblock}
\includegraphics[page=495,pagebox=mediabox,trim={64.39bp 1606.80bp 64.39bp 260.97bp},clip,width=0.94\textwidth,max height=0.38\textheight]{Pope.pdf}
\end{sourceblock}
因此该术语解释了平均对流 , tur-
bu1e:nt运输和生产，以及
(c) 平均压力梯度项消失。

\subsection*{12.1.3 脉动速度的PDF}\phantomsection\addcontentsline{toc}{subsection}{12.1.3 脉动速度的PDF}
脉动速度 u(x, t) 的 PDF 用 g(v; x, t) 表示，其中 v 是样本空间变量。 g 的一阶矩为零，

\begin{sourceblock}
\includegraphics[page=495,pagebox=mediabox,trim={64.39bp 1105.86bp 64.39bp 792.90bp},clip,width=0.94\textwidth,max height=0.38\textheight]{Pope.pdf}
\end{sourceblock}
g 不包含有关平均速度的信息。然而，( U(x, t)) 和 g(v ;x , t) 中包含的信息与 f ( V ; x, t) 中包含的信息相同。

附录 H 中导出的 g(v; x, t) 传输方程为

\begin{sourceblock}
\includegraphics[page=495,pagebox=mediabox,trim={64.39bp 614.92bp 64.39bp 1275.84bp},clip,width=0.94\textwidth,max height=0.38\textheight]{Pope.pdf}
\end{sourceblock}
其中条件压力应变张量为

\begin{sourceblock}
\includegraphics[page=495,pagebox=mediabox,trim={64.39bp 478.94bp 64.39bp 1467.81bp},clip,width=0.94\textwidth,max height=0.38\textheight]{Pope.pdf}
\end{sourceblock}
条件耗散张量为

\begin{sourceblock}
\includegraphics[page=495,pagebox=mediabox,trim={64.39bp 285.96bp 64.39bp 1611.80bp},clip,width=0.94\textwidth,max height=0.38\textheight]{Pope.pdf}
\end{sourceblock}
可以看出，g 方程的左侧（方程（12.19））比 f 方程（方程（12.9））的左侧复杂得多。最后一项（左侧）对应于雷诺应力下的生产
f 的方程。第二种形式显示了雷诺应力方程中与粘性扩散、压力传递、重新分布和耗散相对应的项的分解。

在均匀湍流中，方程式(12.19) 大大简化为

\begin{sourceblock}
\includegraphics[page=496,pagebox=mediabox,trim={64.39bp 1554.80bp 64.39bp 340.96bp},clip,width=0.94\textwidth,max height=0.38\textheight]{Pope.pdf}
\end{sourceblock}
对于非对称张量函数 Hij(v)，分部积分为

\begin{sourceblock}
\includegraphics[page=496,pagebox=mediabox,trim={64.39bp 1302.84bp 64.39bp 554.93bp},clip,width=0.94\textwidth,max height=0.38\textheight]{Pope.pdf}
\end{sourceblock}
\begin{sourceblock}
\includegraphics[page=496,pagebox=mediabox,trim={64.39bp 1125.86bp 64.39bp 822.90bp},clip,width=0.94\textwidth,max height=0.38\textheight]{Pope.pdf}
\end{sourceblock}
表明当方程（1）时获得一致的结果。 (12.19) 是在-
在速度空间上积分，当它乘以 和
在速度空间上积分。

\begin{sourceblock}
\includegraphics[page=496,pagebox=mediabox,trim={64.39bp 854.89bp 64.39bp 1022.87bp},clip,width=0.94\textwidth,max height=0.38\textheight]{Pope.pdf}
\end{sourceblock}
类似地对于

乘以等式。 (12.19) 通过 并在速度空间上积分得
获得雷诺应力张量的精确输运方程

与等式进行逐项比较。 (7.194) 页

\section*{12.2 速度 PDF 模型方程}\phantomsection\addcontentsline{toc}{section}{12.2 速度 PDF 模型方程}
在从纳维-斯托克斯方程导出的 PDF 输运方程（方程（12.9））中，对流和平均压力以封闭形式出现，而粘性项和脉动压力的条件期望则需要模型。 Lundgren (1969) 和 Pope (1981b) 提出了不同类型的模型；但自推出以来（Pope 1983b），流行的模型类型一直是广义朗之万模型（GLM）。在本节中，使用 GLM 的 PDF 方程的行为和属性
22.2.1 广义朗厄恩模型

\subsection*{12.2.1 广义朗之万模型}\phantomsection\addcontentsline{toc}{subsection}{12.2.1 广义朗之万模型}
其中 Co(x, t ) 和 Gij(x,t ) 是定义特定模型的系数：Co 是无量纲的，而 Gij 具有反时间的量纲。对于均匀湍流，脉动速度的 PDF 的相应方程为

\begin{sourceblock}
\includegraphics[page=497,pagebox=mediabox,trim={64.39bp 1191.85bp 64.39bp 729.91bp},clip,width=0.94\textwidth,max height=0.38\textheight]{Pope.pdf}
\end{sourceblock}
在讨论系数的说明和这些方程的行为之前，我们进行以下一般性观察。

(i) 当平均动量方程由方程组成时： (12.26)

（通过乘以 V 并积分），建模项消失，
正确地留下雷诺方程。 （粘性项 v V2(U)

h.为简单起见被省略，但可以很容易地包括在内。）

(ii) 由于耗散的出现 \textasciitilde{} ( xt),, 等式。 (12.26) 通过
本身不是封闭的，但是 f 和 e 的模型方程一起为
(iii) 一般来说，系数 Co(x,t) 和 Gij(x,t ) 取决于局部
(uiuj) 和 t3(Ui)/t3xj 的值，但它们与 V 无关。

(iv) G h a 和 Co 中的两项联合建模 和 中的精确项
p'。每个模型项并不对应于确切的项之一。

(v) 如本节后面所示，系数取决于
这确保了动能在均质中正确演化

新湍流。

(vi) 一个非常重要的观察（Pope 1985，Durbin 和 Speziale 1994）

GLM 方程（方程（12.26））确保了可实现性，前提是
只是 Co 是非负的并且 Co 和 Gij 有界。 （这个

\begin{sourceblock}
\includegraphics[page=498,pagebox=mediabox,trim={64.39bp 1265.84bp 64.39bp 109.44bp},clip,width=\textwidth,max height=0.72\textheight]{Pope.pdf}
\end{sourceblock}
\subsection*{12.2.2 PDF的演变}\phantomsection\addcontentsline{toc}{subsection}{12.2.2 PDF的演变}
\begin{center}\small 图 12.1。扩散对 PDF 形状的影响：方程的解。 (12.29)，Dt = 0、0.02、0.2 和 1。虚线是高斯分布，其均值 (0) 和方差 (3) 与 Dt = 1 时 PDF 的均值 (0) 和方差 (3) 相同。\end{center}
\begin{sourceblock}
\includegraphics[page=498,pagebox=mediabox,trim={64.39bp 585.93bp 64.39bp 1259.84bp},clip,width=0.94\textwidth,max height=0.38\textheight]{Pope.pdf}
\end{sourceblock}
对于几个 PDF 方程，例如 Eq.根据式（12.29），可以得到相应特征函数的解析解。在附录 I 中，回顾了特征函数的属性，第 710 页的表 1.1 提供了有用结果的摘要。

漂移系数最简单的说明是它是各向同性的，即

\begin{sourceblock}
\includegraphics[page=499,pagebox=mediabox,trim={64.39bp 1532.81bp 64.39bp 404.95bp},clip,width=0.94\textwidth,max height=0.38\textheight]{Pope.pdf}
\end{sourceblock}
其中约束方程。 (12.28) 确定时间尺度 TL 为

\begin{sourceblock}
\includegraphics[page=499,pagebox=mediabox,trim={64.39bp 1368.83bp 64.39bp 547.93bp},clip,width=0.94\textwidth,max height=0.38\textheight]{Pope.pdf}
\end{sourceblock}
该规范（Co 为常数）称为简化朗之万模型 (SLM)。

为了说明漂移项的影响，我们考虑以下方程

\begin{sourceblock}
\includegraphics[page=499,pagebox=mediabox,trim={64.39bp 1118.86bp 64.39bp 783.90bp},clip,width=0.94\textwidth,max height=0.38\textheight]{Pope.pdf}
\end{sourceblock}
对于零均值随机过程 u(t) 的 PDF，将 TL 设为常数。很容易证明零均值被保留，并且方差呈指数下降：

\begin{sourceblock}
\includegraphics[page=499,pagebox=mediabox,trim={64.39bp 903.89bp 64.39bp 1039.87bp},clip,width=0.94\textwidth,max height=0.38\textheight]{Pope.pdf}
\end{sourceblock}
因此标准差为 o(t) = o(0)e
方程的独特行为。 (12.32) 的优点是它保留了 PDFa 的形状 正如可以很容易验证的那样，时间 t 的解根据初始条件（t = 0）给出：

\begin{sourceblock}
\includegraphics[page=499,pagebox=mediabox,trim={64.39bp 592.92bp 64.39bp 1303.84bp},clip,width=0.94\textwidth,max height=0.38\textheight]{Pope.pdf}
\end{sourceblock}
对于之前考虑的非物理初始条件，该解如图 12.2 中 t/TL 0 和 1 所示。

对于衰减均匀湍流（使用 (Ui) / a x O），简化的朗之万模型 (SLM) 简化为

\begin{sourceblock}
\includegraphics[page=499,pagebox=mediabox,trim={64.39bp 294.96bp 64.39bp 1604.80bp},clip,width=0.94\textwidth,max height=0.38\textheight]{Pope.pdf}
\end{sourceblock}
可以对 v2 和 v3 进行积分，得到由 g,(v; t) 表示的边际 PDF 的方程：

\begin{sourceblock}
\includegraphics[page=499,pagebox=mediabox,trim={64.39bp 119.39bp 64.39bp 1812.77bp},clip,width=0.94\textwidth,max height=0.38\textheight]{Pope.pdf}
\end{sourceblock}
\begin{sourceblock}
\includegraphics[page=500,pagebox=mediabox,trim={64.39bp 1256.84bp 64.39bp 109.44bp},clip,width=\textwidth,max height=0.72\textheight]{Pope.pdf}
\par\vspace{0.45\baselineskip}
\small 图 12.2。方程（12.34）的解(12.32) 对于 t/T，= 0, 1 1。
\end{sourceblock}
\begin{sourceblock}
\includegraphics[page=500,pagebox=mediabox,trim={64.39bp 605.92bp 64.39bp 775.90bp},clip,width=\textwidth,max height=0.72\textheight]{Pope.pdf}
\par\vspace{0.45\baselineskip}
\small 图 12.3。 PDF g,(v;t£) 根据简化的 Langevin 模型演化，方程： （12.36）。有时会显示标准差为 1、0.99、0.9、0.75 和 0.5 的 PDF。 （常数 Cp 取为 2.1。）
\end{sourceblock}
两项均保留平均值 ( u l ) 0。扩散项导致方差 (u:) 以 Coe 速率增加，而漂移项导致方差以 (Co 假设 (u;) \% 的速率减少。净效应是
?3 )E - 3 产生正确的方差衰减率 (ut),n amely

图 12.3 显示了方程 12.3 的解。 （12.36）来自非物理初始条件。可以看出，随着时间的推移，PDF 趋于高斯分布并变得更窄。方程的解析解。 (12.36) 是在
对于一般情况，广义 Langevin 模型的定性行为与上面示例中的相同。漂移项使 PDF g(v ;t ) 变形，但不会对其形状产生定性影响；而扩散项使其趋向于各向同性的关节法线。对于均匀湍流，如练习 12.12 所示，GLM 允许联合法线解，并且该解从任何初始条件都趋向于联合法线。这是正确的物理行为，因为均匀湍流中的测量确实显示速度的单点 PDF 是联合法向的（参见第 5.5.3 节和第 175 页的图 5.46）。

通过积分方程。 (12.29) 表明 (a) 满足归一化
条件，(b) 保留 u(t) 的均值，(c) 导致方差
以2V的速率增加。

表明，对应于方程。 (12.29)，特征函数

\begin{sourceblock}
\includegraphics[page=501,pagebox=mediabox,trim={64.39bp 1109.86bp 64.39bp 831.89bp},clip,width=0.94\textwidth,max height=0.38\textheight]{Pope.pdf}
\end{sourceblock}
\begin{sourceblock}
\includegraphics[page=501,pagebox=mediabox,trim={64.39bp 965.88bp 64.39bp 904.89bp},clip,width=0.94\textwidth,max height=0.38\textheight]{Pope.pdf}
\end{sourceblock}
（提示：请参阅第 7 10 页上的表 I.1。）通过对该方程求微分
对于 s，验证属性 (a)、(b) 和 (c)。表明
方程的解(12.38) 是

\begin{sourceblock}
\includegraphics[page=501,pagebox=mediabox,trim={64.39bp 739.91bp 64.39bp 1208.85bp},clip,width=0.94\textwidth,max height=0.38\textheight]{Pope.pdf}
\end{sourceblock}
显示（对于 t 0）随机数的特征函数

\begin{sourceblock}
\includegraphics[page=501,pagebox=mediabox,trim={64.39bp 569.93bp 64.39bp 1320.83bp},clip,width=0.94\textwidth,max height=0.38\textheight]{Pope.pdf}
\end{sourceblock}
使用此结果表明，从任何初始 PDF 中，解决方案
扩散方程(12.29) 在很大程度上趋于高斯。

证明，如果 PDF g,(v; t) 按方程式演化： （12.32），然后校正
相应地，特征函数（方程（12.37））演化为
从而得到解

\begin{sourceblock}
\includegraphics[page=502,pagebox=mediabox,trim={64.39bp 1724.78bp 64.39bp 218.97bp},clip,width=0.94\textwidth,max height=0.38\textheight]{Pope.pdf}
\end{sourceblock}
g(a;?) 12..11 令 为 u(t) 的标准化 PDF，即

\begin{sourceblock}
\includegraphics[page=502,pagebox=mediabox,trim={64.39bp 1587.80bp 64.39bp 360.95bp},clip,width=0.94\textwidth,max height=0.38\textheight]{Pope.pdf}
\end{sourceblock}
\begin{sourceblock}
\includegraphics[page=502,pagebox=mediabox,trim={64.39bp 1439.82bp 64.39bp 497.94bp},clip,width=0.94\textwidth,max height=0.38\textheight]{Pope.pdf}
\end{sourceblock}
表明，如果 g 根据 SLM 演化（方程（12.36）），则

\begin{sourceblock}
\includegraphics[page=502,pagebox=mediabox,trim={64.39bp 1231.84bp 64.39bp 622.92bp},clip,width=0.94\textwidth,max height=0.38\textheight]{Pope.pdf}
\end{sourceblock}
以 ( 3 , ) 为 u(t)/o(t) 的特征函数，证明
它的演变是

\begin{sourceblock}
\includegraphics[page=502,pagebox=mediabox,trim={64.39bp 1023.87bp 64.39bp 873.89bp},clip,width=0.94\textwidth,max height=0.38\textheight]{Pope.pdf}
\end{sourceblock}
利用特征法求得解

\begin{sourceblock}
\includegraphics[page=502,pagebox=mediabox,trim={64.39bp 895.89bp 64.39bp 1042.87bp},clip,width=0.94\textwidth,max height=0.38\textheight]{Pope.pdf}
\end{sourceblock}
对长期行为进行评论。将 o 确定为以下函数
t，从而得到 Yu(s?t) 的解。 12.12 定义

\begin{sourceblock}
\includegraphics[page=502,pagebox=mediabox,trim={64.39bp 645.92bp 64.39bp 1273.84bp},clip,width=0.94\textwidth,max height=0.38\textheight]{Pope.pdf}
\end{sourceblock}
表明 g(v,t) 的 GLM（方程（12.27））可以重写

\begin{sourceblock}
\includegraphics[page=502,pagebox=mediabox,trim={64.39bp 480.94bp 64.39bp 1419.82bp},clip,width=0.94\textwidth,max height=0.38\textheight]{Pope.pdf}
\end{sourceblock}
以及特征函数的相应方程

表明，沿着由下式给出的轨迹 3(s,,t)

\begin{sourceblock}
\includegraphics[page=502,pagebox=mediabox,trim={64.39bp 119.39bp 64.39bp 1839.77bp},clip,width=0.94\textwidth,max height=0.38\textheight]{Pope.pdf}
\end{sourceblock}
\begin{sourceblock}
\includegraphics[page=503,pagebox=mediabox,trim={64.39bp 1576.80bp 64.39bp 159.98bp},clip,width=0.94\textwidth,max height=0.38\textheight]{Pope.pdf}
\end{sourceblock}
因此表明方程的解。 (12.51) 是

\begin{sourceblock}
\includegraphics[page=503,pagebox=mediabox,trim={64.39bp 1439.82bp 64.39bp 488.94bp},clip,width=0.94\textwidth,max height=0.38\textheight]{Pope.pdf}
\end{sourceblock}
\begin{sourceblock}
\includegraphics[page=503,pagebox=mediabox,trim={64.39bp 1289.84bp 64.39bp 570.93bp},clip,width=0.94\textwidth,max height=0.38\textheight]{Pope.pdf}
\end{sourceblock}
证明，如果 g(v, t ) 最初是关节法线，那么它仍然是关节
正常或一直。还表明，在任何初始条件下，
分布趋于联合正态（协方差 C），前提是
趋于零。

\subsection*{12.2.3 相应的雷诺应力模型}\phantomsection\addcontentsline{toc}{subsection}{12.2.3 相应的雷诺应力模型}
从 PDF 的 GLM 方程（方程（12.26）），可以直接推导出一阶矩和二阶矩 (U(x, t)) 和 (uiuj\} 的相应方程。正如已经观察到的，一阶矩方程是雷诺方程（方程 (4.12)，忽略粘性项）。二阶矩方程是部分建模的雷诺应力方程
前三个术语表示对流、湍流传输和产生
是准确的。

上述方程可以与精确的雷诺应力方程进行比较(7.194) 第 3 页 19. 如果忽略粘性和压力传输，则模型项（方程 (12.57) 中的 Gij 和 Co 中的项）对应于重新分布和耗散，即

\begin{sourceblock}
\includegraphics[page=503,pagebox=mediabox,trim={64.39bp 265.97bp 64.39bp 1674.79bp},clip,width=0.94\textwidth,max height=0.38\textheight]{Pope.pdf}
\end{sourceblock}
因此，假设耗散是各向同性的，GLM 隐含的压力-应变率率模型为

\begin{sourceblock}
\includegraphics[page=503,pagebox=mediabox,trim={64.39bp 119.39bp 64.39bp 1845.77bp},clip,width=0.94\textwidth,max height=0.38\textheight]{Pope.pdf}
\end{sourceblock}
Rjk 是可再分配的要求 ( e , Rii 0) 导致约束方程： (12.28) GLM 系数。

简化的 Langevin 模型 (SLM) 该模型由 Gij 的最简单可能选择定义，即

\begin{sourceblock}
\includegraphics[page=504,pagebox=mediabox,trim={64.39bp 1545.81bp 64.39bp 393.95bp},clip,width=0.94\textwidth,max height=0.38\textheight]{Pope.pdf}
\end{sourceblock}
将此规格代入方程式。 (12.59) 收益率

\begin{sourceblock}
\includegraphics[page=504,pagebox=mediabox,trim={64.39bp 1372.83bp 64.39bp 524.93bp},clip,width=0.94\textwidth,max height=0.38\textheight]{Pope.pdf}
\end{sourceblock}
这是 Rotta 模型（第 393 页上的方程（11.24）），其系数为

\begin{sourceblock}
\includegraphics[page=504,pagebox=mediabox,trim={64.39bp 1250.84bp 64.39bp 696.91bp},clip,width=0.94\textwidth,max height=0.38\textheight]{Pope.pdf}
\end{sourceblock}
因此，简化的 Langevin 模型（在雷诺应力水平上）与 CR (1 2Co) 的 Rotta 模型相对应。当将其应用于均匀湍流时，SLM 会产生与 Rotta 模型相同的雷诺应力演化。

型号

生产各向同性 Gij 和 Co 的什么规格对应于 LRR-IP 模型？均匀湍流中 g(v, t) 的 GLM 方程（方程（12.27））可以是

\begin{sourceblock}
\includegraphics[page=504,pagebox=mediabox,trim={64.39bp 748.91bp 64.39bp 1107.86bp},clip,width=0.94\textwidth,max height=0.38\textheight]{Pope.pdf}
\end{sourceblock}
平均速度梯度中的项导致雷诺应力方程中的产生 Pij。因此，显然，对 Gij 的贡献为

\begin{sourceblock}
\includegraphics[page=504,pagebox=mediabox,trim={64.39bp 538.93bp 64.39bp 1382.83bp},clip,width=0.94\textwidth,max height=0.38\textheight]{Pope.pdf}
\end{sourceblock}
导致雷诺应力方程中的项 -CzPij，即根据 IP 模型的生产反作用。因此，对应于 LRR-IP 模型的 Gij 规范是该 IP 贡献（方程 (12 .64)）和对应于 Rotta 模型的各向同性贡献之和：

\begin{sourceblock}
\includegraphics[page=504,pagebox=mediabox,trim={64.39bp 246.97bp 64.39bp 1649.79bp},clip,width=0.94\textwidth,max height=0.38\textheight]{Pope.pdf}
\end{sourceblock}
约束方程。 (12.28) 则得出

\begin{sourceblock}
\includegraphics[page=504,pagebox=mediabox,trim={64.39bp 119.39bp 64.39bp 1805.77bp},clip,width=0.94\textwidth,max height=0.38\textheight]{Pope.pdf}
\end{sourceblock}
\begin{sourceblock}
\includegraphics[page=505,pagebox=mediabox,trim={64.39bp 1717.78bp 64.39bp 159.98bp},clip,width=0.94\textwidth,max height=0.38\textheight]{Pope.pdf}
\end{sourceblock}
总结一下：对于给定的 CR 和 C2 值，Gij 和 co 的规范
根据方程式(12.65) 和(12.67) 产生对应于LRR-IP 的广义Langevin 模型。在均匀湍流中，两个模型给出的雷诺应力的演变是相同的。

Haworth-Pope 模型 Haworth 和 Pope (1986) 提出了模型1
其中四阶张量 H 是

\begin{sourceblock}
\includegraphics[page=505,pagebox=mediabox,trim={64.39bp 1030.87bp 64.39bp 797.90bp},clip,width=0.94\textwidth,max height=0.38\textheight]{Pope.pdf}
\end{sourceblock}
a(n 中的项是仅涉及雷诺应力各向异性 bij 的最一般可能项，而 H 中的项在平均速度梯度和各向异性中都是线性的。系数值（视为常数）由 Haworth 和 Pope (1986) 以及 Haworth 和 Pope 给出
压力-应变率率张量的相应模型很容易从方程中推导出来。 （12.59）。它的一般形式为 Eq. ( 1 1 . 1 5 )，系数由 和 确定（参见 Pope (1994b) 和
f(") ( ) 习题 12.13，方程 (12.72))。

拉格朗日 IP 模型 (LIPM) 此类中一个有用的模型是由 Pope (1994b) 提出的 LIPM。指定了三个常数：Co 2.1、C2 0.6 和 a2 3.5。剩余的
-- - 系数由 = - 3 \textasciitilde{} 2 , - '(1
和--(I - Cz)？ = 0, y j - = -

\begin{sourceblock}
\includegraphics[page=505,pagebox=mediabox,trim={64.39bp 255.97bp 64.39bp 1653.79bp},clip,width=0.94\textwidth,max height=0.38\textheight]{Pope.pdf}
\end{sourceblock}
该模型的性能与 LRR-IP 模型（方程（12.65）-（12-67））差别不大，但具有恒定值的优点
Rotta 常数的隐含值为

\begin{sourceblock}
\includegraphics[page=506,pagebox=mediabox,trim={64.39bp 1707.78bp 64.39bp 239.97bp},clip,width=0.94\textwidth,max height=0.38\textheight]{Pope.pdf}
\end{sourceblock}
（参见方程（12.66）），其中 F 是归一化 Reynoldsst.ress 张量的行列式，方程(11 -34)。

一般情况

式(12.59)表明，对于每个G L M，对于再分布Rij，都有一个唯一的模型。而且，该模型是可实现的(仅满足系数有界、Co非负的条件)。

反之则更为微妙。我们已经确定了与 Rotta 模型和 LRR-IP 模型相对应的 Co 和 G m 的规格，但是
ij 这些规范并不是唯一的：请注意，Rij 只有五个独立组件，而 Gij 有九个。 Pope (1994b) 更深入地讨论了这些问题。

12.13 练习（12.59）表明压力-应变率率模型相关
响应 Haworth-Pope 模型方程。 (12.68) 为等式。 (11.135)

系数由下式给出

\begin{sourceblock}
\includegraphics[page=506,pagebox=mediabox,trim={64.39bp 426.95bp 64.39bp 1143.86bp},clip,width=0.94\textwidth,max height=0.38\textheight]{Pope.pdf}
\end{sourceblock}
12.14 证明对于任意 y5，剩余系数 yi) 中
可以指定 Haworth-Pope 模型，以便相应的
压力-应变率率模型是LRR-IF吗？ 12.15 证明对应于的系数（方程（12.72））

Haworth-Pope 模型满足

\begin{sourceblock}
\includegraphics[page=506,pagebox=mediabox,trim={64.39bp 119.39bp 64.39bp 1836.77bp},clip,width=0.94\textwidth,max height=0.38\textheight]{Pope.pdf}
\end{sourceblock}
\begin{sourceblock}
\includegraphics[page=507,pagebox=mediabox,trim={64.39bp 815.90bp 64.39bp 109.44bp},clip,width=\textwidth,max height=0.72\textheight]{Pope.pdf}
\par\vspace{0.45\baselineskip}
\small 图 12.4。欧拉建模方法。
\end{sourceblock}
讨论系数的这种线性相关性的重要性

\subsection*{12.2.4 欧拉和拉格朗日建模方法}\phantomsection\addcontentsline{toc}{subsection}{12.2.4 欧拉和拉格朗日建模方法}
到目前为止，在考虑雷诺应力和 PDF 闭合时，我们遵循欧拉建模方法，总结如下：

\subsection*{12.2.5 拉格朗日PDF和欧拉PDF之间的关系}\phantomsection\addcontentsline{toc}{subsection}{12.2.5 拉格朗日PDF和欧拉PDF之间的关系}
纳维-斯托克斯方程。例如，考虑以下流体粒子速度演化的模型方程：

\begin{sourceblock}
\includegraphics[page=508,pagebox=mediabox,trim={64.39bp 1643.79bp 64.39bp 247.97bp},clip,width=0.94\textwidth,max height=0.38\textheight]{Pope.pdf}
\end{sourceblock}
其中系数 Gij(x,t ) 是 (uiuj)、a(Ui)/axj 和 e 的局部函数。 Gij 中的术语将比力（由于波动的压力梯度和粘度）建模为速度波动的线性函数。

从该模型方程导出的欧拉 PDF f ( V ; x , t) 的演化方程为

\begin{sourceblock}
\includegraphics[page=508,pagebox=mediabox,trim={64.39bp 1329.83bp 64.39bp 597.92bp},clip,width=0.94\textwidth,max height=0.38\textheight]{Pope.pdf}
\end{sourceblock}
这直接来自等式。 （12.8）。该方程和模型耗散方程构成一个闭集。事实上，等式。 (12.75) 与扩散系数设置为零 ( e . Co 0) 的广义 Langevin 模型方程（方程 (12.26)）相同。因此，闭合模型 PDF 方程也可以从流体粒子速度模型（例如方程 1）获得。 (12 .74)。

等式的右侧。 （12.74）是流体粒子上的特定力的确定性模型。就目前情况而言，这不是一个令人满意的模型，因为方程中没有扩散项。 (12.75) 导致 PDF 的形状松弛至关节法线。下一节将介绍朗之万方程。这些是流体粒子上特定力的随机模型，可产生令人满意的 PDF 行为。图 12.5 总结了这种替代拉格朗日建模方法。

练习 12.16 平均速度的必要条件是什么？

等式。 (12.75) 索非尼？写下论证的步骤

表明 f 和 (方程 (12.75) 和 (10.53)) 的方程
形成一个闭集。 12 .1 7 推导出平均动量和雷诺应力方程
等式。 （12.75）。证明 (uiuj) 和 的方程是封闭的
均匀湍流，但不适用于一般情况。

拉格万量和欧拉量之间的关系 尽管后续章节中描述的模型是拉格朗日量，但理解它们对欧拉量的影响当然是有必要的。

\begin{sourceblock}
\includegraphics[page=509,pagebox=mediabox,trim={64.39bp 1080.86bp 64.39bp 109.44bp},clip,width=\textwidth,max height=0.72\textheight]{Pope.pdf}
\par\vspace{0.45\baselineskip}
\small 图 12.5。拉格朗日建模方法。
\end{sourceblock}
此处建立的拉格朗日和欧拉 PDF，方程。 （12.80）和（12.87）。

xf(t, uf(t,
可以回忆一下第 2.2 节，Y) 和 Y ) 表示在参考时间从位置 Y 出发的流体粒子的位置和速度。这些是拉格朗日量，它们的联合 PDF
到。因此称为拉格朗日 PDF，并用 f (, x; tl Y) 表示。这可以用细粒度拉格朗日 PDF f 来表示；作为

\begin{sourceblock}
\includegraphics[page=509,pagebox=mediabox,trim={64.39bp 627.92bp 64.39bp 1320.83bp},clip,width=0.94\textwidth,max height=0.38\textheight]{Pope.pdf}
\end{sourceblock}
\begin{sourceblock}
\includegraphics[page=509,pagebox=mediabox,trim={64.39bp 499.94bp 64.39bp 1443.82bp},clip,width=0.94\textwidth,max height=0.38\textheight]{Pope.pdf}
\end{sourceblock}
（细粒度 PDF 的属性，例如 f ; 在附录 H 中进行了描述。）

现在考虑所有初始点的积分，即 dY。自从
xf(t, 点 Y 和 Y 之间存在一对一映射)，该积分可以重新表示为

\begin{sourceblock}
\includegraphics[page=509,pagebox=mediabox,trim={64.39bp 157.98bp 64.39bp 1747.78bp},clip,width=0.94\textwidth,max height=0.38\textheight]{Pope.pdf}
\end{sourceblock}
考虑不可压缩流时，连续性方程的结果是 J 是单位（参见第 16 页的练习 2.5）。因此方程。 (12.78)

\begin{sourceblock}
\includegraphics[page=510,pagebox=mediabox,trim={64.39bp 1507.81bp 64.39bp 246.97bp},clip,width=0.94\textwidth,max height=0.38\textheight]{Pope.pdf}
\end{sourceblock}
δ函数的筛选特性挑选出流体粒子

xf(t, 位于 Y) x，其速度为 U(x, t)。可以认识到 6(U(x, t) - V) 是细粒度欧拉 PDF (Eq. ( H ) )s o 即 Eq. ( H ) 的期望。 (12.79) 得出所需的结果：

\begin{sourceblock}
\includegraphics[page=510,pagebox=mediabox,trim={64.39bp 1200.85bp 64.39bp 695.91bp},clip,width=0.94\textwidth,max height=0.38\textheight]{Pope.pdf}
\end{sourceblock}
通常将 Y 视为固定的初始位置，因此（在湍流中）X+(t, Y) 是随机的。然而，我们可以认为 X' 是固定的 ( i . X f(t, Y ) x)，在这种情况下，Y 是在时间 t 处 x 处的流体粒子的随机初始位置。从这个角度可以看出，f L ( V , x ; t l Y ) 有第二种解释：即，对于固定的 x 和 t，为
uf(t, 时间 t 处 X+ ( t , ' \textasciitilde{}=) x 处的流体粒子的速度 Y ) U(x,t) 和初始位置 Y 的联合 PDF。因此方程的左侧。 (12.80)代表边际PDT；流体粒子在 x 处的速度（无论其初始位置如何），与欧拉 PDF 相同。

结果方程。 (12.80)同样适用于脉动速度的PDF。欧拉速度波动 u(x, t) 的 PDF g(v; x, t) 与 f 的关系如下：

\begin{sourceblock}
\includegraphics[page=510,pagebox=mediabox,trim={64.39bp 584.93bp 64.39bp 1344.83bp},clip,width=0.94\textwidth,max height=0.38\textheight]{Pope.pdf}
\end{sourceblock}
xf(t, Y 的联合 P D F) 和波动

\begin{sourceblock}
\includegraphics[page=510,pagebox=mediabox,trim={64.39bp 447.94bp 64.39bp 1489.81bp},clip,width=0.94\textwidth,max height=0.38\textheight]{Pope.pdf}
\end{sourceblock}
由 gf(v,x;tlY) 表示，并且与以下相关：

\begin{sourceblock}
\includegraphics[page=510,pagebox=mediabox,trim={64.39bp 306.96bp 64.39bp 1630.79bp},clip,width=0.94\textwidth,max height=0.38\textheight]{Pope.pdf}
\end{sourceblock}
它直接从方程式得出。 (12.80)-(12.83) gf 和 g 是

\begin{sourceblock}
\includegraphics[page=510,pagebox=mediabox,trim={64.39bp 119.39bp 64.39bp 1742.78bp},clip,width=0.94\textwidth,max height=0.38\textheight]{Pope.pdf}
\end{sourceblock}
\section*{12.3 朗之万方程}\phantomsection\addcontentsline{toc}{section}{12.3 朗之万方程}
均匀湍流 根据定义，在均匀湍流中，脉动欧拉速度场 u(x, t) 在统计上是均匀的：其 PDF g(v; t) 独立于
要考虑的适当拉格朗日量是速度波动 uf (t, Y ) 和位移

\begin{sourceblock}
\includegraphics[page=511,pagebox=mediabox,trim={64.39bp 1545.81bp 64.39bp 396.95bp},clip,width=0.94\textwidth,max height=0.38\textheight]{Pope.pdf}
\end{sourceblock}
由于统计同质性，用 gi(v, v; t ) 表示的 uf (t, Y) 和 Rf (t, Y ) 的联合 PDF 与 Y 无关。 uf(t , Y ) 的边际 PDF

\begin{sourceblock}
\includegraphics[page=511,pagebox=mediabox,trim={64.39bp 1338.83bp 64.39bp 563.93bp},clip,width=0.94\textwidth,max height=0.38\textheight]{Pope.pdf}
\end{sourceblock}
也独立于 Y。

根据这些定义，对于均匀湍流，方程： (12.84) 可以重新表示为

\begin{sourceblock}
\includegraphics[page=511,pagebox=mediabox,trim={64.39bp 965.88bp 64.39bp 800.90bp},clip,width=0.94\textwidth,max height=0.38\textheight]{Pope.pdf}
\end{sourceblock}
因此，在均匀湍流中，u(x, t) 和 uf (t, Y ) 的 PDF 是相同的，与 x 和 Y 无关。就单点一次性统计而言，所有点，无论是固定 (x) 还是移动，在统计上都是等效的。在平均速度为零的均匀湍流的更简单情况下，U(x, t)、u(x, t)、U+( t, Y) 和 uf (t, Y) 的单点一次性 PDF 均相等：

\begin{sourceblock}
\includegraphics[page=511,pagebox=mediabox,trim={64.39bp 518.93bp 64.39bp 1427.82bp},clip,width=0.94\textwidth,max height=0.38\textheight]{Pope.pdf}
\end{sourceblock}
12.3 Langevin 方程 Langevin 方程最初被提出（Langevin 1908）作为经历布朗运动的微观粒子的速度的随机模型。本节描述并显示了该方程，为湍流中流体粒子的速度提供了一个良好的模型。

Ornstein-Uhlenbeck (OU) 过程，其 PDF 由 Fokker-Planck 方程演化。在随机过程的术语中，U'(t) 是扩散过程，朗之万方程是随机微分方程（SDE）。附录 J 总结了与扩散过程相关的必要数学，其中描述了 OU 过程的属性。

12.3.1 我们首先考虑均匀各向同性湍流的最简单情况，通过人工强迫使其达到统计平稳。平均速度为零，k 和 的值恒定。在这些情况下，所有流体粒子在统计上是相同的，并且它们的速度的三个分量在统计上也是相同的。因此，考虑流体质点速度的一个分量就足够了，用 U+(t) 表示。由 Langevin 方程给出的 Uf (t) 模型用 U*(t) 表示。

\subsection*{12.3.1 静止各向同性湍流}\phantomsection\addcontentsline{toc}{subsection}{12.3.1 静止各向同性湍流}
朗之万方程是随机微分方程

\begin{sourceblock}
\includegraphics[page=512,pagebox=mediabox,trim={64.39bp 964.88bp 64.39bp 924.88bp},clip,width=0.94\textwidth,max height=0.38\textheight]{Pope.pdf}
\end{sourceblock}
其中 TL 和 a2 是正常数。对于不熟悉随机微分方程的读者，可以通过有限差分方程来理解该方程

\begin{sourceblock}
\includegraphics[page=512,pagebox=mediabox,trim={64.39bp 686.91bp 64.39bp 1202.85bp},clip,width=0.94\textwidth,max height=0.38\textheight]{Pope.pdf}
\end{sourceblock}
当 A t 趋于零时，处于极限。这里 c(t) 是标准化高斯随机
((t(t)) 变量 = 0, (<(t)" = 1) 在不同时间独立于自身 ((5(t)5(tf)) = 0, 对于 t t ) 并且 nd 在过去时间独立于 U*(t) (例如, ( ( t ) ( t ) ) = 0 对于 t' 5 t)。在方程 (12.89) 中，确定性漂移项 (-U' dt/TL) 导致速度在时间尺度 TL 上向零松弛，而扩散项则添加标准差 a t / 2 d t l TL 的零均值随机增量。

图 12.6 显示了由 Langevin 方程生成的 OU 过程 U*(t) 的实现或示例路径。因为等式。 (12.89) 是随机的（ e . i t 包含随机性），来自相同初始条件的两个样本路径是不同的。

这里的上标表示 \textasciitilde{}agrangianq 数量，不要与归一化混淆

图 12.6。 Langevin 生成的 Ornstein-Uhlenbeck 过程的示例路径
方程，方程。 （12.89）。

U* (t) 的 PDF f ) 由 Fokker-Planck 方程演化而来

\begin{sourceblock}
\includegraphics[page=513,pagebox=mediabox,trim={64.39bp 925.88bp 64.39bp 975.88bp},clip,width=0.94\textwidth,max height=0.38\textheight]{Pope.pdf}
\end{sourceblock}
OU 过程的属性在附录 J 中描述。简而言之，U'(t) 是一个统计平稳的高斯马尔可夫过程，具有不可微分的连续样本路径。作为平稳高斯过程，它完全由均值（零）、方差 a2, 和自相关函数来表征，即

\begin{sourceblock}
\includegraphics[page=513,pagebox=mediabox,trim={64.39bp 627.92bp 64.39bp 1314.83bp},clip,width=0.94\textwidth,max height=0.38\textheight]{Pope.pdf}
\end{sourceblock}
如果 p(s) 是拉格朗日速度自相关函数，则拉格朗日积分时间尺度定义为

\begin{sourceblock}
\includegraphics[page=513,pagebox=mediabox,trim={64.39bp 411.95bp 64.39bp 1487.81bp},clip,width=0.94\textwidth,max height=0.38\textheight]{Pope.pdf}
\end{sourceblock}
显然，等式。式(12.92)与这个定义是一致的，所以朗之万方程(Eq.(12.89))中的系数TL确实是过程的积分时间尺度。

与观测结果的比较 OU 过程 U'(t) 在多大程度上模拟了速度的行为
产生速度的高斯 PDF。在各向同性湍流中，拉格朗日速度 Uf (t) 的一次性 PDF 与单点一次性欧拉 PDF 相同（方程 (12.88)）。来自实验和 D N S（例如，Yeung 和 Pope (1989)）的明确证据是这些 PDF 非常接近高斯。

\begin{sourceblock}
\includegraphics[page=514,pagebox=mediabox,trim={64.39bp 1559.80bp 64.39bp 328.96bp},clip,width=0.94\textwidth,max height=0.38\textheight]{Pope.pdf}
\end{sourceblock}
产生正确的速度方差。

Langevin 方程的一个限制是 U+(t) 可微，而 U'(t) 则不可微。因此，如果在无穷小的时间尺度上检查 U*(t)，则该模型在质量上是不正确的。然而，考虑高雷诺数湍流，其中积分时间尺度 TL 和柯尔莫哥洛夫时间尺度之间存在较大分离，让我们检查惯性范围时间尺度 s、TL s 上的 U+(t)。这最好通过拉格朗日结构函数来完成

\begin{sourceblock}
\includegraphics[page=514,pagebox=mediabox,trim={64.39bp 1114.86bp 64.39bp 828.90bp},clip,width=0.94\textwidth,max height=0.38\textheight]{Pope.pdf}
\end{sourceblock}
柯尔莫哥洛夫假设（最初的，来自 194 1，并于 1962 年改进）预测

\begin{sourceblock}
\includegraphics[page=514,pagebox=mediabox,trim={64.39bp 939.88bp 64.39bp 1004.87bp},clip,width=0.94\textwidth,max height=0.38\textheight]{Pope.pdf}
\end{sourceblock}
其中 Co 是通用常数；而朗之万方程得出

\begin{sourceblock}
\includegraphics[page=514,pagebox=mediabox,trim={64.39bp 727.91bp 64.39bp 1144.86bp},clip,width=0.94\textwidth,max height=0.38\textheight]{Pope.pdf}
\end{sourceblock}
（参见方程式（J.46））。因此，朗之万方程在惯性范围内产生 DL 对 s 的线性依赖性方面与柯尔莫哥洛夫假设一致。

Langevin 方程中的两个参数可以代替 a2 和 TL，用 k 和 E 表示。方程 (12.94) 给出用 k 表示的 a2，我们通过以下关系式引入模型系数 3 Co

\begin{sourceblock}
\includegraphics[page=514,pagebox=mediabox,trim={64.39bp 319.96bp 64.39bp 1580.80bp},clip,width=0.94\textwidth,max height=0.38\textheight]{Pope.pdf}
\end{sourceblock}
因此时间尺度由下式给出

\begin{sourceblock}
\includegraphics[page=514,pagebox=mediabox,trim={64.39bp 164.98bp 64.39bp 1746.78bp},clip,width=0.94\textwidth,max height=0.38\textheight]{Pope.pdf}
\end{sourceblock}
\begin{sourceblock}
\includegraphics[page=515,pagebox=mediabox,trim={64.39bp 1175.85bp 64.39bp 109.44bp},clip,width=\textwidth,max height=0.72\textheight]{Pope.pdf}
\par\vspace{0.45\baselineskip}
\small 图 12.7。拉格朗日速度自相关函数。线，Langevin 模型 p(s) = 实心符号，Sato 和 Yamamoto (1987) 的实验数据 tb，Ry = 46 且 <，R； = 66；开放符号，Yeung 和 Pope (1989) 的 DNS 数据，R； = 90。
\end{sourceblock}
通过以这种方式重新表达系数，朗之万方程变为

\begin{sourceblock}
\includegraphics[page=515,pagebox=mediabox,trim={64.39bp 808.90bp 64.39bp 1116.86bp},clip,width=0.94\textwidth,max height=0.38\textheight]{Pope.pdf}
\end{sourceblock}
如果模型系数 Co 被取为 Kolmogorov 常数 Co，则 Langevin 方程在数量上与 Kolmogorov 假设一致。此时 DL(s)（由方程（12.96）给出）和 DL(s)（由方程（12.97）给出）是相同的（对于 << s << TL）。Co 的值和 Co 的选择将在 12.4 节中讨论。

高雷诺数流中的拉格朗日统计数据已证明可以用于实验和 DNS。然而，在低或中等雷诺数时，这两种技术都已用于测量拉格朗日自相关函数 p ( \textasciitilde{} )F。图 12.7 显示了各向同性湍流中 p(s) 的测量值与朗之万方程给出的指数（方程 (12 .92)）的比较。在非常小的时间 (s/TL 1)，朗之万方程产生的指数在性质上是不正确的，因为它具有负斜率，这是 U*(t) 不可微分的结果。然而，对于更大的时间，指数提供了一个非常合理的
练习 1-2.18 在朗之万方程（方程（12.90））的有限差分形式中，
随机变量被定义为独立于 U*(tr)f 或
t' 5 吨。证明 c(t) 和 U*(t At) 是相关的。

通过取等式的平均值。 (12.90)，证明均值 (U*(t))e 旋转
证明方程的平方的平均值。 (12.90) 是

\begin{sourceblock}
\includegraphics[page=516,pagebox=mediabox,trim={64.39bp 1324.83bp 64.39bp 569.93bp},clip,width=0.94\textwidth,max height=0.38\textheight]{Pope.pdf}
\end{sourceblock}
因此表明方差 (u ) r (U ) - 演化为

\begin{sourceblock}
\includegraphics[page=516,pagebox=mediabox,trim={64.39bp 1161.85bp 64.39bp 742.91bp},clip,width=0.94\textwidth,max height=0.38\textheight]{Pope.pdf}
\end{sourceblock}
从初始条件 (U*) = V1 和 ( u * ) = 0 at t = tl 得到
方程式的解决方案。 （12.101）和（12.102）。表明这些是一致的
与方程。 (5.41) 第 720 页。

\begin{sourceblock}
\includegraphics[page=516,pagebox=mediabox,trim={64.39bp 866.89bp 64.39bp 977.88bp},clip,width=0.94\textwidth,max height=0.38\textheight]{Pope.pdf}
\end{sourceblock}
因此表明，对于平稳 U"(t)，自协方差为

\begin{sourceblock}
\includegraphics[page=516,pagebox=mediabox,trim={64.39bp 748.91bp 64.39bp 1198.85bp},clip,width=0.94\textwidth,max height=0.38\textheight]{Pope.pdf}
\end{sourceblock}
如何确定 s O 的 R(s)？ 12.21 以 Y(s, t) 为 U*(t) 的特征函数，证明
等式。 (12.90) 对应于

\begin{sourceblock}
\includegraphics[page=516,pagebox=mediabox,trim={64.39bp 491.94bp 64.39bp 1410.82bp},clip,width=0.94\textwidth,max height=0.38\textheight]{Pope.pdf}
\end{sourceblock}
通过将右侧展开到 At 的一阶，证明

\begin{sourceblock}
\includegraphics[page=516,pagebox=mediabox,trim={64.39bp 285.96bp 64.39bp 1555.80bp},clip,width=0.94\textwidth,max height=0.38\textheight]{Pope.pdf}
\end{sourceblock}
证明该方程的傅里叶变换是 Fokker-
普朗克方程，Eq. (12.91)-（参见第 725 页的练习 5.4 了解
该方程的解。）

广义 Lmgevin 模型 现在描述 Langevin 方程（方程（12.100））对非均匀湍流的基本扩展，结果是简化的

\subsection*{12.3.2 广义朗之万模型}\phantomsection\addcontentsline{toc}{subsection}{12.3.2 广义朗之万模型}
被视为 SLM 的延伸。

这些朗之万方程的主题是位置为 X'(t) 的粒子的速度 U*(t)。该粒子模拟流体粒子的行为，因此它以自己的速度移动，即

\begin{sourceblock}
\includegraphics[page=517,pagebox=mediabox,trim={64.39bp 1380.83bp 64.39bp 547.93bp},clip,width=0.94\textwidth,max height=0.38\textheight]{Pope.pdf}
\end{sourceblock}
简化的 Langevin 模型被写为随机微分方程

\begin{sourceblock}
\includegraphics[page=517,pagebox=mediabox,trim={64.39bp 1126.86bp 64.39bp 672.91bp},clip,width=0.94\textwidth,max height=0.38\textheight]{Pope.pdf}
\end{sourceblock}
wj) 其中 W(t) 是向量值维纳过程（具有属性 (d Wi d dt dij，以及系数 (d(p)/axi、k、E 和 (Ui)），在粒子位置 X*( t ) 处进行评估。与稳态各向同性湍流的标量朗之万方程 (Eq. (12.100)) 相比，SLM 模型的差异如下。

(i) SLM 适用于粒子矢量速度 \textasciitilde{} ' ( t )。

(ii) 添加平均压力梯度中的漂移项：这是精确的
纳维-斯托克斯方程中的项（用 (p) p' 表示）。

(iii) 在第二个漂移项中，粒子速度 U'(t) 向
局部欧拉平均值 (U)。 （在静止各向同性湍流中，
哪个方程。 (12.100) 成立，平均速度为零。）

(iv) 系数的附加导致正确的
能量耗散率如下所示。 （或者，相反，这个 1
从等式中省略。 (12.100) 因为人工强迫恰好
这些修改（相对于方程（12.100））是与平均动量和动能方程保持一致所需的最小限度。

广义 Langevin 模型 (Pope 198 3b) 是

\begin{sourceblock}
\includegraphics[page=517,pagebox=mediabox,trim={64.39bp 177.98bp 64.39bp 1741.78bp},clip,width=0.94\textwidth,max height=0.38\textheight]{Pope.pdf}
\end{sourceblock}
(uiuj) 和 a ( \textasciitilde{} i ) / a xA\textasciitilde{}S 的局部值。上面提到的GLM是一类pm odels；特定模型对应于 G 和 Co 的特定规范。特定模型的示例是给出 Gij 的 SLM

\begin{sourceblock}
\includegraphics[page=518,pagebox=mediabox,trim={64.39bp 1587.80bp 64.39bp 285.96bp},clip,width=0.94\textwidth,max height=0.38\textheight]{Pope.pdf}
\end{sourceblock}
以及 IP 模型，其中 Gij 和 Co 由等式给出。 （12.65）和（12.67）。

GLM 的特点是漂移项在 U' 中呈线性，扩散系数为各向同性，即 (C \textasciitilde{} E ) " \textasciitilde{}C\textasciitilde{}on.sider 代替一般扩散系数 表示，因此扩散项为 B, dW-。此时，扩散过程给出的二阶拉格朗日结构函数为

\begin{sourceblock}
\includegraphics[page=518,pagebox=mediabox,trim={64.39bp 1171.85bp 64.39bp 704.91bp},clip,width=0.94\textwidth,max height=0.38\textheight]{Pope.pdf}
\end{sourceblock}
（参见方程（12 .97））；而柯尔莫哥洛夫假设产生

\begin{sourceblock}
\includegraphics[page=518,pagebox=mediabox,trim={64.39bp 969.88bp 64.39bp 917.88bp},clip,width=0.94\textwidth,max height=0.38\textheight]{Pope.pdf}
\end{sourceblock}
（参见等式（12.96））。因此，为了与局部各向同性保持一致，GLM 中各向同性扩散系数的选择是必要的；并且，从这些 D 方程的比较中可以明显看出；-(s) 与
和 DL，
。 11 ij Kolmogorov 假设由规范实现

\begin{sourceblock}
\includegraphics[page=518,pagebox=mediabox,trim={64.39bp 700.91bp 64.39bp 1245.84bp},clip,width=0.94\textwidth,max height=0.38\textheight]{Pope.pdf}
\end{sourceblock}
与 Co 公司。

给定扩散项的形式，U* 中 Gij 的线性确保 GLM 在均匀湍流中产生联合法向速度分布，与观测结果一致。 （Durbin 和 Speziale (1994) 进一步考虑在 Langevin 模型中使用各向异性扩散。）

广义 Langevin 模型中的拉格朗日 P D F 方程，X*(t) 和 U*(t) 模式表示流体-粒子属性，因此表示它们的联合 PDF f；是由于流体的模型
佛罗里达州。粒子被定义为在时间 to ( e . X f (to, Y) Y) 处源自 Y，这是 f 的适当定义； V , x ; tl Y) 是 X*(t) 和 U*(t) 的联合 PDF，以 X*(to) Y 为条件。

（方程（5.56））由 X*(t) 和 U*(t) 方程获得（方程（12.108）和（12.110））。它是

\begin{sourceblock}
\includegraphics[page=519,pagebox=mediabox,trim={64.39bp 1596.80bp 64.39bp 348.96bp},clip,width=0.94\textwidth,max height=0.38\textheight]{Pope.pdf}
\end{sourceblock}
其中所有系数均在 (x, t) 处计算（参见练习 12.7.2）。

欧拉 PDF 方程 根据拉格朗日 PDF 和欧拉 PDF 之间的关系（方程（12.80）），广义 Langevin 模型对应的模型欧拉 PDF 为

\begin{sourceblock}
\includegraphics[page=519,pagebox=mediabox,trim={64.39bp 1228.84bp 64.39bp 715.91bp},clip,width=0.94\textwidth,max height=0.38\textheight]{Pope.pdf}
\end{sourceblock}
f 的演化方程很容易通过对 f 的方程进行积分来获得： (Eq. (12.115)) 在所有 Y 上。由于该方程不包含对 Y 的依赖（除了 f ），因此结果很简单

\begin{sourceblock}
\includegraphics[page=519,pagebox=mediabox,trim={64.39bp 908.89bp 64.39bp 1038.87bp},clip,width=0.94\textwidth,max height=0.38\textheight]{Pope.pdf}
\end{sourceblock}
这个方程正是 12.2 节中讨论的模型欧拉 PDF 方程，即方程 1。 (12.26) （尽管这里 PDF 用 f 表示，以强调它源自 X*(t) 和 U*(t) 的随机拉格朗日模型方程，方程（12.108）和（12.110））因此，第 12.2 节中所做的所有推论都源于方程（12.26）。 （12.26）反映了随机拉格朗日模型的属性。

图 12.8 总结了随机拉格朗日模型 (GLM) 和各种统计演化方程之间的联系。特别是，U*(t)（方程 (12.1 lo)）的 G L M 方程意味着雷诺应力 (uiuj) 根据方程 2 演变。 （12.57）。

可实现性 随机拉格朗日模型 X*(t) 和 U'(t) 是可实现的，仅受 G L M 方程（方程（12.1.10））中的系数为实数且有界的最低条件影响。鉴于 X*(t) 和 U'(t) 是可实现的，因此它们的所有统计量也是可实现的，包括拉格朗日和 l3ulerian PDF f；，和 f'，以及雷诺应力

\begin{sourceblock}
\includegraphics[page=520,pagebox=mediabox,trim={64.39bp 1224.85bp 64.39bp 109.44bp},clip,width=\textwidth,max height=0.72\textheight]{Pope.pdf}
\par\vspace{0.45\baselineskip}
\small 图 12.8。由随机拉格朗日模型得出的统计演化方程。
\end{sourceblock}
因此，PDF 方法不仅可以轻松实现可实现性，而且图 12.8 中所示的从随机模型到雷诺应力方程的路线提供了一种简单的方法来解决雷诺应力闭合中的可实现性问题。特别是，压力-应变率率模型可以写成等式的右侧。如果 Gij 有界且 Co 为正且有界，则 (12.59) 是可实现的。 Pope (1994b) 以及 Durbin 和 Speziale (1994) 进一步考虑了这些问题。

考虑六维扩散过程 Z * ( t)，其中
Z * 的分量为 \{ Z , Z ;, Z;, 2 , Z ;, Z \} \{ X X ;, Xi, U;, U;,
U\}。 X*(t) 和 U*(t) 根据方程演变。 (12.108) 和
(12 . 11 0 s Z*(t)e 如何通过一般扩散方程旋转

\begin{sourceblock}
\includegraphics[page=520,pagebox=mediabox,trim={64.39bp 482.94bp 64.39bp 1461.82bp},clip,width=0.94\textwidth,max height=0.38\textheight]{Pope.pdf}
\end{sourceblock}
并写下漂移向量 ai 的显式表达式，即 diffu-
sion 矩阵 Bv 和扩散系数 Bij BikBjkH。有效使用
Fokker-Planck方程的一般形式（方程（5.56））来验证
代替纳维-斯托克斯方程，并作为
等式。 (12.74)，考虑方程

\begin{sourceblock}
\includegraphics[page=520,pagebox=mediabox,trim={64.39bp 119.39bp 64.39bp 1795.77bp},clip,width=0.94\textwidth,max height=0.38\textheight]{Pope.pdf}
\end{sourceblock}
其中 是可微函数。表明对应的
U(X,t ), f ( V ;x , t) 的欧拉 PDF 演化为
证明欧拉 PDF 方程 (Eq. (12.1 17)) 源于
G L M 可以写成相同的形式

\begin{sourceblock}
\includegraphics[page=521,pagebox=mediabox,trim={64.39bp 1438.82bp 64.39bp 468.94bp},clip,width=0.94\textwidth,max height=0.38\textheight]{Pope.pdf}
\end{sourceblock}
表明，如果 f 是关节法线，则 A 由下式给出

\begin{sourceblock}
\includegraphics[page=521,pagebox=mediabox,trim={64.39bp 1292.84bp 64.39bp 648.92bp},clip,width=0.94\textwidth,max height=0.38\textheight]{Pope.pdf}
\end{sourceblock}
其中 C-' 是雷诺应力倒数的 i-j 分量
张量 Cij (uiuj)。因此认为，对于均匀湍流，
其中 f * 是关节法线，等式： (12.119) 和 (12.122) 提供
具有相同欧拉 PDF 演化的确定性模型
GLM 的。对于一般流程来说，这种方式有什么缺点
确定性模型？ 12.24 通过比较纳维-斯托克斯方程（方程（2.35））与
G L M 方程 (Eq. ( 2 0 ) )，表示广义朗之万
模型相当于对脉动压力梯度进行建模，
粘性项为

\begin{sourceblock}
\includegraphics[page=521,pagebox=mediabox,trim={64.39bp 719.91bp 64.39bp 1178.85bp},clip,width=0.94\textwidth,max height=0.38\textheight]{Pope.pdf}
\end{sourceblock}
其中 是白噪声（见方程（5.29））。 12.25 设 u*(t) 为速度随
粒子，其定义为

\begin{sourceblock}
\includegraphics[page=521,pagebox=mediabox,trim={64.39bp 488.94bp 64.39bp 1455.82bp},clip,width=0.94\textwidth,max height=0.38\textheight]{Pope.pdf}
\end{sourceblock}
对于 X*(t)e 通过方程旋转。 (12.108) 和 U*(t)e 通过扩散旋转
过程，表明 u*(t) 演化为

\begin{sourceblock}
\includegraphics[page=521,pagebox=mediabox,trim={64.39bp 160.98bp 64.39bp 1686.79bp},clip,width=0.94\textwidth,max height=0.38\textheight]{Pope.pdf}
\end{sourceblock}
(Eq. 2 1 0 ) ) 表示 u*(t) 如何演化

\begin{sourceblock}
\includegraphics[page=522,pagebox=mediabox,trim={64.39bp 1675.79bp 64.39bp 243.97bp},clip,width=0.94\textwidth,max height=0.38\textheight]{Pope.pdf}
\end{sourceblock}
验证雷诺应力的演化方程

从方程获得的张量（ujut） (12.126)与式+(12.57)一致。

特别是，表明右侧第一项
等式。 (12.126) 得出生产 pij*
从等式获得(12.126) g*( v ;t ) 的演化方程，
u*(t) 的 PDF。验证该方程式与方程式相同。 （12.27）。

\section*{12.4 湍流扩散}\phantomsection\addcontentsline{toc}{section}{12.4 湍流扩散}
在前面的章节中，随机拉格朗日模型（例如 GLM）已被用作湍流建模的基础，专门用于获得速度欧拉 PDF 的闭合模型方程。从历史上看，随机拉格朗日模型在湍流弥散的背景下使用得更早、更广泛。在本节中，考虑网格湍流中线源背后的色散的最简单情况，这既是因为它的历史意义，也是因为它揭示了朗之万方程的性能。

所考虑的问题 湍流建模应用中解决的常见问题的简单表述是：给定流动的几何形状以及初始和边界条件，确定平均速度场 (U(x, t)) 和一些湍流特征（例如，(uiui) 和湍流弥散研究解决不同的问题，即：给定湍流流场（以 (U(Xt,) )、(u\textasciitilde{}u,) 等形式表示），确定源自特定源的守恒无源标量的平均场 (\$(x, t)) 此类研究的主要应用涉及污染物在大气中以及河流、湖泊和海洋中的扩散（参见 e+g+F、ischer 等人+ (1979) 和 Hunt (1985)）。在这些应用中，源可能在时间上是连续的（例如，从管道排放到海洋中的废水）。仅限于短时间内（例如，有毒气体的意外释放） 与域的大小相比，源的大小通常非常小。

欧拉方法 确定平均标量场 (\$(x, t)) 的一种方法是求解
(4) (～4～)。带有标量通量模型的 (Eq. ( 4 4 1 ) ) 输运方程 例如，如果使用湍流扩散假设（Eq. (4.42)），则得到的模型方程为

\begin{sourceblock}
\includegraphics[page=523,pagebox=mediabox,trim={64.39bp 1507.81bp 64.39bp 394.95bp},clip,width=0.94\textwidth,max height=0.38\textheight]{Pope.pdf}
\end{sourceblock}
r rT 其中 和 分别是分子扩散率和湍流扩散率。

拉格朗赞方法

1 fie 湍流扩散的另一种拉格朗日方法起源于 G. I. Taylor 的经典 (192 1) 论文“连续运动的扩散”。泰勒认为（在高雷诺数时）空间传输
与平均流动和湍流运动造成的对流传输相比，由于分子扩散造成的影响可以忽略不计。在完全忽略分子扩散的情况下，在流体粒子（D+/Dt O）之后是守恒的，因此平均场的演变可以根据流体粒子运动的统计来确定。现在考虑相关统计数据。

粒子位置的 PDF 回想一下，x + ( t , Y ) 表示源自时间 t 时位置 Y 的流体粒子在时间 t 时的位置。 X+(t, Y ) 的 PDF 称为质点位置的正向 P D F，可表示为

\begin{sourceblock}
\includegraphics[page=523,pagebox=mediabox,trim={64.39bp 680.91bp 64.39bp 1267.84bp},clip,width=0.94\textwidth,max height=0.38\textheight]{Pope.pdf}
\end{sourceblock}
细粒度PDF在哪里

\begin{sourceblock}
\includegraphics[page=523,pagebox=mediabox,trim={64.39bp 543.93bp 64.39bp 1402.82bp},clip,width=0.94\textwidth,max height=0.38\textheight]{Pope.pdf}
\end{sourceblock}
它与拉格朗日 PDF（方程（12.76））的关系为

\begin{sourceblock}
\includegraphics[page=523,pagebox=mediabox,trim={64.39bp 371.95bp 64.39bp 1531.81bp},clip,width=0.94\textwidth,max height=0.38\textheight]{Pope.pdf}
\end{sourceblock}
在这些定义中，Y 被认为是一个（非随机）自变量，然后 X+(t, Y ) 给出了流体粒子随时间向前移动的随机轨迹，对于 t > to。

考虑时间上向后的流体粒子轨迹也很有用。为此，我们将 yf (t, x) 定义为流体粒子在时间 to 时的位置，

\begin{sourceblock}
\includegraphics[page=524,pagebox=mediabox,trim={64.39bp 1238.84bp 64.39bp 109.44bp},clip,width=\textwidth,max height=0.72\textheight]{Pope.pdf}
\par\vspace{0.45\baselineskip}
\small 图 12.9。向前 (a) 和向后 (b) 流体粒子轨迹的草图（基于湍流的不同实现）。 (a) 前向轨迹-在时间 t0 处起始于 Y 的流体粒子路径。 (b) 向后轨迹-在时间 f 到达 x 的流体粒子路径。
\end{sourceblock}
\textasciitilde{} ' ( t ,x ) 是随机初始位置。向前和向后的流体粒子轨迹如图 12.9 所示。

函数 \textasciitilde{} ' ( t ,Y ) 提供初始点 Y 和最终点 X+ 之间的一对一映射；类似地，\textasciitilde{}' ( t ,x ) 提供最终点和初始点 Y f 之间的一对一映射。当然，这些函数是互逆的；那就是

\begin{sourceblock}
\includegraphics[page=524,pagebox=mediabox,trim={64.39bp 738.91bp 64.39bp 1204.85bp},clip,width=0.94\textwidth,max height=0.38\textheight]{Pope.pdf}
\end{sourceblock}
\begin{sourceblock}
\includegraphics[page=524,pagebox=mediabox,trim={64.39bp 654.92bp 64.39bp 1288.84bp},clip,width=0.94\textwidth,max height=0.38\textheight]{Pope.pdf}
\end{sourceblock}
对于正在考虑的恒定密度流，映射的雅可比行列式

\begin{sourceblock}
\includegraphics[page=524,pagebox=mediabox,trim={64.39bp 437.94bp 64.39bp 1399.82bp},clip,width=0.94\textwidth,max height=0.38\textheight]{Pope.pdf}
\end{sourceblock}
（参见练习 2.5）。因此，逆矩阵的雅可比行列式也是统一的：

\begin{sourceblock}
\includegraphics[page=524,pagebox=mediabox,trim={64.39bp 287.96bp 64.39bp 1639.79bp},clip,width=0.94\textwidth,max height=0.38\textheight]{Pope.pdf}
\end{sourceblock}
\textasciitilde{}' ( t ,x ) 的 PDF 称为流体质点位置的后向 P D F，可写为

\begin{sourceblock}
\includegraphics[page=524,pagebox=mediabox,trim={64.39bp 119.39bp 64.39bp 1840.77bp},clip,width=0.94\textwidth,max height=0.38\textheight]{Pope.pdf}
\end{sourceblock}
细粒度PDF在哪里

\begin{sourceblock}
\includegraphics[page=525,pagebox=mediabox,trim={64.39bp 1719.78bp 64.39bp 221.97bp},clip,width=0.94\textwidth,max height=0.38\textheight]{Pope.pdf}
\end{sourceblock}
一个简单而有价值但绝不明显的关系是后向 PDF 和前向 PDF 是相等的：

\begin{sourceblock}
\includegraphics[page=525,pagebox=mediabox,trim={64.39bp 1538.81bp 64.39bp 392.95bp},clip,width=0.94\textwidth,max height=0.38\textheight]{Pope.pdf}
\end{sourceblock}
这从发展中可以看出

\begin{sourceblock}
\includegraphics[page=525,pagebox=mediabox,trim={64.39bp 1188.85bp 64.39bp 528.93bp},clip,width=0.94\textwidth,max height=0.38\textheight]{Pope.pdf}
\end{sourceblock}
在第三行中，所有 y 的积分被重新表示为所有 X+ 的积分。雅可比行列式是统一的，因此下一行是根据 delta 函数的筛选性质得出的；其中 X+(tYye)q uals x 和 y 等于 Y+(t, x)。

拉格朗日公式回到湍流弥散问题，我们考虑无界湍流，其中源在当时将标量场确定性地初始化为某个值

\begin{sourceblock}
\includegraphics[page=525,pagebox=mediabox,trim={64.39bp 747.91bp 64.39bp 1199.85bp},clip,width=0.94\textwidth,max height=0.38\textheight]{Pope.pdf}
\end{sourceblock}
如果完全忽略分子扩散，则\$在流体粒子之后守恒。因此，\$ 在流体粒子轨迹的初始点和终点具有相同的值：

\begin{sourceblock}
\includegraphics[page=525,pagebox=mediabox,trim={64.39bp 385.95bp 64.39bp 1415.82bp},clip,width=0.94\textwidth,max height=0.38\textheight]{Pope.pdf}
\end{sourceblock}
最后一个方程的期望产生所需的结果

\begin{sourceblock}
\includegraphics[page=525,pagebox=mediabox,trim={64.39bp 142.98bp 64.39bp 1674.79bp},clip,width=0.94\textwidth,max height=0.38\textheight]{Pope.pdf}
\end{sourceblock}
在位置 Yo 处的单位点源的特定情况下，即

\begin{sourceblock}
\includegraphics[page=526,pagebox=mediabox,trim={64.39bp 1719.78bp 64.39bp 226.97bp},clip,width=0.94\textwidth,max height=0.38\textheight]{Pope.pdf}
\end{sourceblock}
\begin{sourceblock}
\includegraphics[page=526,pagebox=mediabox,trim={64.39bp 1584.80bp 64.39bp 300.96bp},clip,width=0.94\textwidth,max height=0.38\textheight]{Pope.pdf}
\end{sourceblock}
因此，由单位点源产生的平均守恒标量场由源自源的流体粒子的位置的 PDF 给出。

重要的是要认识到，只有就 @ 的空间传输而言，在这一发展中忽视分子扩散才是合理的。例如，即使在高雷诺数下，分子扩散也会对标量方差 (\#7) 和标量 PDF f,(p;x,t) 产生一阶效应。因此，即使这些量的表达式可以从等式中获得。 （12.141），它们没有物理有效性。 （Saffman (1960) 和 Pope (1998b) 考虑了分子扩散的影响。）

来自点源的色散

要考虑的最简单的情况是统计平稳的各向同性湍流中点源的色散。单位源（方程 (12.143)）位于原点 (Yo = 0)，释放发生在时间 t) = 0。各向同性湍流速度场的均值为零 ((U(x,t)) = 0)，并保持统计平稳（通过人工强迫）。均方根值速度为 u’，因此雷诺应力为 (uju;) = u7d,。

对于这种情况，平均标量场 (@(x, t)) 等于源自源的流体粒子位置 X*(t,0\} 的 PDF fx(x;1|0)（方程 (12.144)）。因此，色散可以通过该 PDF 的一阶矩和二阶矩来表征，即 X*(t,0) 的均值和方差。（如果 fx 是联合法线 - 正如 Langevin 模型预测 - 那么这些矩完全表征色散。）

流体-质点运动方程 0X* /ét = U* 可以积分，

\begin{sourceblock}
\includegraphics[page=526,pagebox=mediabox,trim={64.39bp 381.95bp 64.39bp 1516.81bp},clip,width=0.94\textwidth,max height=0.38\textheight]{Pope.pdf}
\end{sourceblock}
这样X*的统计量就可以用拉格朗日速度统计量来表示。平均值为零

\begin{sourceblock}
\includegraphics[page=526,pagebox=mediabox,trim={64.39bp 168.98bp 64.39bp 1733.78bp},clip,width=0.94\textwidth,max height=0.38\textheight]{Pope.pdf}
\end{sourceblock}
流体-粒子位置的协方差为

\begin{sourceblock}
\includegraphics[page=527,pagebox=mediabox,trim={64.39bp 1681.79bp 64.39bp 219.97bp},clip,width=0.94\textwidth,max height=0.38\textheight]{Pope.pdf}
\end{sourceblock}
考虑到湍流是平稳且各向同性的，二次拉格朗日速度相关性可以写成

\begin{sourceblock}
\includegraphics[page=527,pagebox=mediabox,trim={64.39bp 1498.81bp 64.39bp 439.94bp},clip,width=0.94\textwidth,max height=0.38\textheight]{Pope.pdf}
\end{sourceblock}
其中 p(s) 是拉格朗日速度自相关函数。因此协方差位置也是各向同性的，

\begin{sourceblock}
\includegraphics[page=527,pagebox=mediabox,trim={64.39bp 1310.83bp 64.39bp 627.92bp},clip,width=0.94\textwidth,max height=0.38\textheight]{Pope.pdf}
\end{sourceblock}
并由标准导数 ox(t) 表征。从这三个方程我们得到

\begin{sourceblock}
\includegraphics[page=527,pagebox=mediabox,trim={64.39bp 1087.86bp 64.39bp 809.90bp},clip,width=0.94\textwidth,max height=0.38\textheight]{Pope.pdf}
\end{sourceblock}
通过一个不平凡的操作（参见练习12.26），二重积分可以化简为单积分，得到

\begin{sourceblock}
\includegraphics[page=527,pagebox=mediabox,trim={64.39bp 868.89bp 64.39bp 1032.87bp},clip,width=0.94\textwidth,max height=0.38\textheight]{Pope.pdf}
\end{sourceblock}
对于非常短的时间 ( t TL )，p(s) 充分近似为
SO p(0) = 1) 方程中的积分。 (12.151) 是 - 1 t 2，从而得出结果

\begin{sourceblock}
\includegraphics[page=527,pagebox=mediabox,trim={64.39bp 690.91bp 64.39bp 1255.84bp},clip,width=0.94\textwidth,max height=0.38\textheight]{Pope.pdf}
\end{sourceblock}
对于所考虑的短时间内，流体粒子速度 (p(s) I) 的变化可以忽略不计，因此基本上存在直线运动

\begin{sourceblock}
\includegraphics[page=527,pagebox=mediabox,trim={64.39bp 501.94bp 64.39bp 1441.82bp},clip,width=0.94\textwidth,max height=0.38\textheight]{Pope.pdf}
\end{sourceblock}
（紧随其后的是方程（112.152））。

在另一个极端，对于非常大的时间（t TL），方程中的积分： (12.1 51) 可以近似为

\begin{sourceblock}
\includegraphics[page=527,pagebox=mediabox,trim={64.39bp 228.97bp 64.39bp 1668.79bp},clip,width=0.94\textwidth,max height=0.38\textheight]{Pope.pdf}
\end{sourceblock}
这导致平方根 s \textasciitilde{} read

\begin{sourceblock}
\includegraphics[page=527,pagebox=mediabox,trim={64.39bp 119.39bp 64.39bp 1836.77bp},clip,width=0.94\textwidth,max height=0.38\textheight]{Pope.pdf}
\end{sourceblock}
\begin{sourceblock}
\includegraphics[page=528,pagebox=mediabox,trim={64.39bp 1210.85bp 64.39bp 109.44bp},clip,width=\textwidth,max height=0.72\textheight]{Pope.pdf}
\par\vspace{0.45\baselineskip}
\small 图 12.10。由 Langevin 模型给出的点源色散的标准偏差 ox（方程（12.159））。
\end{sourceblock}
\begin{sourceblock}
\includegraphics[page=528,pagebox=mediabox,trim={64.39bp 939.88bp 64.39bp 1003.87bp},clip,width=0.94\textwidth,max height=0.38\textheight]{Pope.pdf}
\end{sourceblock}
\begin{sourceblock}
\includegraphics[page=528,pagebox=mediabox,trim={64.39bp 691.91bp 64.39bp 1245.84bp},clip,width=0.94\textwidth,max height=0.38\textheight]{Pope.pdf}
\end{sourceblock}
\begin{sourceblock}
\includegraphics[page=528,pagebox=mediabox,trim={64.39bp 508.94bp 64.39bp 1392.82bp},clip,width=0.94\textwidth,max height=0.38\textheight]{Pope.pdf}
\end{sourceblock}
图：朗之万模型给出的流体-粒子路径样本，显示为
(a) 中等时间和 (b) 长时间。虚线显示-
朗之万方程 流体粒子速度的朗之万模型可以对湍流扩散进行完整的预测。根据该模型，速度自相关函数是指数 p(s) exp(- TL)。

= sJ 使用 p(s) 的表达式，方程： (12.151) 可以积分得到

\begin{sourceblock}
\includegraphics[page=529,pagebox=mediabox,trim={64.39bp 246.97bp 64.39bp 1695.79bp},clip,width=0.94\textwidth,max height=0.38\textheight]{Pope.pdf}
\end{sourceblock}
图 12.10 显示了该公式给出的 ox(t)：分别在小时间和大时间的线性和平方根行为清晰可见。

速度是一个 Omstein-Uhlenbeck 过程，它是一个高斯过程。由此可见，流体质点位置（OU 过程的 mth e 积分）也是高斯过程。因此，Langevin模型预测的平均标量场是高斯分布

\begin{sourceblock}
\includegraphics[page=530,pagebox=mediabox,trim={64.39bp 1595.80bp 64.39bp 349.96bp},clip,width=0.94\textwidth,max height=0.38\textheight]{Pope.pdf}
\end{sourceblock}
ox(t) 由方程式给出。 （12.159）。图 12.11 显示了朗之万方程生成的流体粒子路径。

练习 12.26 来自方程式(12.150)，替换为 r = t" 和 s = t" - t'，并且
回想一下 p(s) 是偶函数，获得

\begin{sourceblock}
\includegraphics[page=530,pagebox=mediabox,trim={64.39bp 1215.85bp 64.39bp 681.91bp},clip,width=0.94\textwidth,max height=0.38\textheight]{Pope.pdf}
\end{sourceblock}
操纵表达式

表明

\begin{sourceblock}
\includegraphics[page=530,pagebox=mediabox,trim={64.39bp 901.89bp 64.39bp 991.87bp},clip,width=0.94\textwidth,max height=0.38\textheight]{Pope.pdf}
\end{sourceblock}
因此验证方程。 （12.151）。 12.2.7 证明扩散方程

\begin{sourceblock}
\includegraphics[page=530,pagebox=mediabox,trim={64.39bp 706.91bp 64.39bp 1190.85bp},clip,width=0.94\textwidth,max height=0.38\textheight]{Pope.pdf}
\end{sourceblock}
承认高斯解，方程。 (12.160), ox(t) 变化ac-
根据方程式(12.157): 12.28 证明指数近似 p(s) = e x p ( - s l / T - ) 得出

\begin{sourceblock}
\includegraphics[page=530,pagebox=mediabox,trim={64.39bp 500.94bp 64.39bp 1426.82bp},clip,width=0.94\textwidth,max height=0.38\textheight]{Pope.pdf}
\end{sourceblock}
和等式。 (12.159) 为 ～ (tV)eri\textasciitilde{}fy。短期和长期
这些表达式的行为。 12.29 L e t X * ( t ) 和 U*(t)b 流体位置和速度的分量
粒子根据朗之万模型演化，方程： （12.89）。展示
他们的协方差是

\begin{sourceblock}
\includegraphics[page=530,pagebox=mediabox,trim={64.39bp 201.97bp 64.39bp 1737.78bp},clip,width=0.94\textwidth,max height=0.38\textheight]{Pope.pdf}
\end{sourceblock}
确定相关系数 px.. 这在以下情况下表现如何

\begin{sourceblock}
\includegraphics[page=531,pagebox=mediabox,trim={64.39bp 1403.82bp 64.39bp 109.44bp},clip,width=\textwidth,max height=0.72\textheight]{Pope.pdf}
\par\vspace{0.45\baselineskip}
\small 图 12.12。热线源实验的草图，显示了湍流产生网格下游的加热线。
\end{sourceblock}
Taylor (1921) 提出了一个随机模型
流体粒子的位置 X*(t).A 根据模型，X*(t)i s
马尔可夫过程，并且在一个小的时间间隔 A 处增量
A X * )h 作为属性

\begin{sourceblock}
\includegraphics[page=531,pagebox=mediabox,trim={64.39bp 991.87bp 64.39bp 898.89bp},clip,width=0.94\textwidth,max height=0.38\textheight]{Pope.pdf}
\end{sourceblock}
其中方差 dZ 和相关系数 c 取决于 At。

表明，在适当指定 d2 和 c 的情况下，在极限内
在 t 0 时，该模型等效于 Langevin 方程，Eq. （12.89）。

（提示：表明 U*(t) A X * ( ) A tt 结束于扩散过程。）

网格湍流中的线源 用于研究湍流扩散的最简单的实验是图 12.12 中描绘的热尾流。细加热丝垂直于平均流量放置，位于网格尺寸的湍流生成网格下游一段距离
xo M。导线足够细，不会显着影响速度场。尽管电线可能很热，但经过它的空气温度会迅速下降到几度。因此，除了紧邻电线之外，多余温度 \$(x, t) 是守恒无源标量。

在实验中（例如，Warhaft (1984) 和 Stapountzis 等人 (1986)），发现平均分布呈高斯分布，

\begin{sourceblock}
\includegraphics[page=531,pagebox=mediabox,trim={64.39bp 119.39bp 64.39bp 1806.77bp},clip,width=0.94\textwidth,max height=0.38\textheight]{Pope.pdf}
\end{sourceblock}
特征宽度 oy(x,) 随着导线 x, 下游距离的增加而增加。 （（∨）在整个流中的积分是守恒的，并且，如方程（12.167）所暗示的，@被归一化，使得该积分是统一的。）

在实验室框架中，热尾流在统计上是固定的和二维的：统计数据仅取决于 x 和 y。为了更好地近似，在以平均速度移动的坐标系中，温度场在统计上是一维的：统计数据仅取决于 y 和 t（距源的平均飞行时间）。这个随时间演变的问题对应于平面源的色散（y = 0 且 t = 0），并且可以采用与点源的色散大致相同的方式进行分析。事实上，如果湍流是非衰减的，则上一小节中获得的 oy 表达式（方程（12.151）和（12.159））将适用于热尾流 oy 的厚度。

Anand 和 Pope (1985) 将 Langevin 模型应用于热尾流，考虑到湍流的衰减以及靠近源的分子扩散的不可忽略的影响。他们的结果是

\begin{sourceblock}
\includegraphics[page=532,pagebox=mediabox,trim={64.39bp 1036.87bp 64.39bp 852.89bp},clip,width=0.94\textwidth,max height=0.38\textheight]{Pope.pdf}
\end{sourceblock}
其中 Ly 是导线处的长度尺度 k*/*/¢，I 是分子扩散率，nis 是湍流衰减指数，常数 r 和 s 根据朗之万模型常数 Cp 给出：

\begin{sourceblock}
\includegraphics[page=532,pagebox=mediabox,trim={64.39bp 819.90bp 64.39bp 1127.86bp},clip,width=0.94\textwidth,max height=0.38\textheight]{Pope.pdf}
\end{sourceblock}
\begin{sourceblock}
\includegraphics[page=532,pagebox=mediabox,trim={64.39bp 732.91bp 64.39bp 1205.85bp},clip,width=0.94\textwidth,max height=0.38\textheight]{Pope.pdf}
\end{sourceblock}
图 12.13 将 oy 的 Langevin 模型预测（方程（12.168））与实验数据进行了比较。尽管针对这些数据选择了值 Cy = 2.1，但在相当大的范围内的一致性可能看起来非常好。

Cy 和 Cy 的值 考虑到 Langevin 方程与 Kolmogorov 假设之间的一致性，人们可能会认为从热尾流数据获得的 Cy = 2.1 将为 Kolmogorov 常数 Co 提供良好的估计。但事实并非如此，因为实验的雷诺数相当低，并且 2 的预测取决于惯性子区之外的时间间隔 s 的自相关函数 p(s)。尽管Co很难确定，但有迹象表明其价值大于

\begin{sourceblock}
\includegraphics[page=533,pagebox=mediabox,trim={64.39bp 1083.86bp 64.39bp 109.44bp},clip,width=\textwidth,max height=0.72\textheight]{Pope.pdf}
\par\vspace{0.45\baselineskip}
\small 图 12.13。热尾流 oy 的厚度（通过湍流长度尺度 Lo 标准化）作为导线下游距离 x 的函数（通过从网格到导线 xo 的距离标准化）。线，朗之万模型方程。 (12.168)；符号，Warhaft (1984) 的实验数据，xo/M = 20(e) 和 60(A)。
\end{sourceblock}
\begin{sourceblock}
\includegraphics[page=533,pagebox=mediabox,trim={64.39bp 483.94bp 64.39bp 1414.82bp},clip,width=0.94\textwidth,max height=0.38\textheight]{Pope.pdf}
\end{sourceblock}
\begin{sourceblock}
\includegraphics[page=533,pagebox=mediabox,trim={64.39bp 299.96bp 64.39bp 1590.80bp},clip,width=0.94\textwidth,max height=0.38\textheight]{Pope.pdf}
\end{sourceblock}
\section*{12.5 速度-频率联合 PDF}\phantomsection\addcontentsline{toc}{section}{12.5 速度-频率联合 PDF}
\subsection*{12.5.1 完整的PDF闭合}\phantomsection\addcontentsline{toc}{subsection}{12.5.1 完整的PDF闭合}
广义朗之万模型提供了速度 f ( V ; x, t) 的 PDF 的模型方程，方程 (12.26)。就其本身而言，该方程并不能提供完整的闭合，因为方程中的系数涉及数量
k / \textasciitilde{} )th at 在 f 方面未知。模型（例如，f 和（或 f 和 o \textasciitilde{} / k ）d o 的 and = 方程一起提供了完整的闭合，并且 Haworth 和 El Tahry (1991) 的计算提供了该方法的应用示例。

获得完全闭合的替代方法（此处描述）是将有关 PDF 框架内湍流尺度分布的信息纳入其中。具体来说，我们考虑速度和湍流频率（定义如下）的联合 PDF。该联合 PDF 的模型方程提供了完整的闭合。

除了提供完全封闭之外，描述湍流频率的分布还有其他好处。可以考虑耗散率的大波动（例如，内部间歇性）；并且，在非均匀流动中，可以根据湍流的起源和历史更全面地考虑湍流流体的不同行为。对于非湍流流体，湍流频率为零，而对于湍流流体，湍流频率严格为正。通过利用这一特性，可以（从速度频率联合 PDF）获得间歇因子 y (x, t)，以及湍流和非湍流条件 PDF。

我们重新定义o(x,t)为瞬时湍流频率

\begin{sourceblock}
\includegraphics[page=534,pagebox=mediabox,trim={64.39bp 697.91bp 64.39bp 1248.84bp},clip,width=0.94\textwidth,max height=0.38\textheight]{Pope.pdf}
\end{sourceblock}
其中 E \textasciitilde{} ( xt), i 是瞬时耗散率（方程（6.312））。它的意思是

\begin{sourceblock}
\includegraphics[page=534,pagebox=mediabox,trim={64.39bp 537.93bp 64.39bp 1387.83bp},clip,width=0.94\textwidth,max height=0.38\textheight]{Pope.pdf}
\end{sourceblock}
这与 k-o 模型中的第二个变量相同。 w(x, t) 的欧拉 PDF 表示为 f,,(0 ;x , t) （其中 0 是对应于 a)s o 的样本空间变量，其平均值也由下式给出

\begin{sourceblock}
\includegraphics[page=534,pagebox=mediabox,trim={64.39bp 251.97bp 64.39bp 1648.79bp},clip,width=0.94\textwidth,max height=0.38\textheight]{Pope.pdf}
\end{sourceblock}
因为根据定义，o 是非负的，所以样本空间的下限为零。

\subsection*{12.5.2 湍流频率的对数正态模型}\phantomsection\addcontentsline{toc}{subsection}{12.5.2 湍流频率的对数正态模型}
o ( x ,t ) - 表示为 J'(V O ,x , t )。和往常一样，边际 PDF 是

\begin{sourceblock}
\includegraphics[page=535,pagebox=mediabox,trim={64.39bp 1661.79bp 64.39bp 201.97bp},clip,width=0.94\textwidth,max height=0.38\textheight]{Pope.pdf}
\end{sourceblock}
\begin{sourceblock}
\includegraphics[page=535,pagebox=mediabox,trim={64.39bp 1552.80bp 64.39bp 352.96bp},clip,width=0.94\textwidth,max height=0.38\textheight]{Pope.pdf}
\end{sourceblock}
与速度一样，湍流频率模型也是从拉格朗日观点开发的。在接下来的两小节中，描述了随机模型 o * ( t )f 或跟随流体粒子的湍流频率 o f ( t , Y )。 Pope 和 Chen (1990) 提出的模型在均匀湍流中产生对数正态 PDF obf o，而 Jayesh 和 Pope (1995) 提出的模型产生伽玛分布。 U*(t) 和 o* ( t ) 的随机拉格朗日模型得出欧拉速度-频率联合 PDF, V, 8 的闭合模型方程； x，t）。

湍流频率的对数正态模型 正如第 6.7 节中所讨论的，瞬时耗散表现出间歇性行为，并且近似于对数正态分布。 Pope 和 Chen (1990) 为 w*(t) 设计了一个随机模型来整合这些属性。

首先，我们考虑统计上平稳的均匀湍流，因此平均湍流频率 (o) 是恒定的。 （对于这种情况，欧拉平均值 (o) 和拉格朗日平均值 (a*a)r e 相等。）如果 o* 是对数正态分布，则（根据定义）数量

\begin{sourceblock}
\includegraphics[page=535,pagebox=mediabox,trim={64.39bp 690.91bp 64.39bp 1254.84bp},clip,width=0.94\textwidth,max height=0.38\textheight]{Pope.pdf}
\end{sourceblock}
是正态分布的。因此，通过取 \textasciitilde{} " ( tto) 获得的 o*(ti)s 对数正态模型是 Ornstein-Uhlenbeck 过程（因为这会产生 \textasciitilde{} ' ( t ) 的正态分布）。

OU 进程为 \textasciitilde{} * ( tis) 编写

\begin{sourceblock}
\includegraphics[page=535,pagebox=mediabox,trim={64.39bp 366.95bp 64.39bp 1521.81bp},clip,width=0.94\textwidth,max height=0.38\textheight]{Pope.pdf}
\end{sourceblock}
其中要指定的系数是 x* 的方差及其自相关时间尺度 Tx 通过参考中等雷诺数 DNS 数据，这些系数被视为

\begin{sourceblock}
\includegraphics[page=535,pagebox=mediabox,trim={64.39bp 158.98bp 64.39bp 1778.78bp},clip,width=0.94\textwidth,max height=0.38\textheight]{Pope.pdf}
\end{sourceblock}
\begin{sourceblock}
\includegraphics[page=536,pagebox=mediabox,trim={64.39bp 1195.85bp 64.39bp 109.44bp},clip,width=\textwidth,max height=0.72\textheight]{Pope.pdf}
\par\vspace{0.45\baselineskip}
\small 图 12.14。湍流频率的对数正态随机模型的样本路径，方程： （12.181）。
\end{sourceblock}
X* 的平均值为 --o 的正态性（参见练习 12.32）。方程式中漂移项中的因子 1 2 (12.179) 确保 (x') 具有正确的值（在静止状态下）。

对应的随机微分方程 co* (t) = (w ) exp (t)] 从方程获得。 （12.179）通过使用 i 到 o 变换（方程（J.52）），是

\begin{sourceblock}
\includegraphics[page=536,pagebox=mediabox,trim={64.39bp 748.91bp 64.39bp 1147.86bp},clip,width=0.94\textwidth,max height=0.38\textheight]{Pope.pdf}
\end{sourceblock}
该过程的示例路径如图 12.14 所示。可以看出，该过程确实是高度可变且间歇性的。存在幅度较小的持续平静期；也有大幅波动的时期，达到平均值的十倍。

当它扩展到非平稳和非均匀流时（Pope 1991 a），co* ( t ) 的随机微分方程为

\begin{sourceblock}
\includegraphics[page=536,pagebox=mediabox,trim={64.39bp 300.96bp 64.39bp 1638.79bp},clip,width=0.94\textwidth,max height=0.38\textheight]{Pope.pdf}
\end{sourceblock}
扩散项保持不变，漂移项被重新表达（因此，即使不是对数正态分布，其均值也为零）；还包括涉及系数 h 和 S 的附加漂移项。根据方程。 (12.182)，如果 h 是
所有时间。因此，包含 h 中的术语是为了允许非湍流流体变成湍流：Pope (199 la) 给出了详细信息。

对于均匀湍流，h 中的项可以忽略不计，并且方程 (co) 隐含的平均值 (co) 的演变。 (12.182) 是

\begin{sourceblock}
\includegraphics[page=537,pagebox=mediabox,trim={64.39bp 1571.80bp 64.39bp 349.96bp},clip,width=0.94\textwidth,max height=0.38\textheight]{Pope.pdf}
\end{sourceblock}
因此，S 决定了均值的演变，并且为了与 k-E 和 k-co 模型保持一致，可以将其指定为
CW1 CC1 1 和 C(02 CE2 1，参见方程 (10.94)。

除了上述 co*(t) 模型之外，Pope 和 Chen (1990) 还提出了 Langevin 模型（针对速度 U*(t)），其中
细化系数取决于 w (t) （以及其他量）。特别地，扩散系数是Cokw*来代替Coe Cok(co)：即它基于瞬时耗散，与内部间歇性的思想和精炼的相似性假设一致。

o* 的对数正态模型与 U*(t) 的精化 Langevin 模型一起产生速度-频率联合 PDF 的闭合模型方程。该模型已被 Pope (1991a)、Anand、Pope 和 Mongia (19931，以及 Minier 和 Pozorski (1995)) 应用于许多流。作为一个例子，图 12.15 显示了自相似平面混合层中速度的偏度和峰度剖面。显然，这些较高矩的计算相当准确。

练习 12.32 考虑 w* 是对数正态分布，y* 定义为
等式。 (12.178) 有方差表明归一化 con-

\begin{sourceblock}
\includegraphics[page=537,pagebox=mediabox,trim={64.39bp 429.95bp 64.39bp 1515.81bp},clip,width=0.94\textwidth,max height=0.38\textheight]{Pope.pdf}
\end{sourceblock}
（源自方程 12.178）意味着 的平均值是

\begin{sourceblock}
\includegraphics[page=537,pagebox=mediabox,trim={64.39bp 281.96bp 64.39bp 1662.79bp},clip,width=0.94\textwidth,max height=0.38\textheight]{Pope.pdf}
\end{sourceblock}
也得到结果

\begin{sourceblock}
\includegraphics[page=537,pagebox=mediabox,trim={64.39bp 119.39bp 64.39bp 1791.77bp},clip,width=0.94\textwidth,max height=0.38\textheight]{Pope.pdf}
\end{sourceblock}
\begin{sourceblock}
\includegraphics[page=538,pagebox=mediabox,trim={64.39bp 1003.87bp 64.39bp 109.44bp},clip,width=\textwidth,max height=0.72\textheight]{Pope.pdf}
\par\vspace{0.45\baselineskip}
\small 图 12.15。自相似平面混合层中轴向 (u) 和横向 (v) 速度的偏斜度和平坦度分布。线，Minier 和 Pozorski (1995) 基于 Pope (1991a) 的对数正态/精炼 Langevin 模型的计算；符号、Wygnanski 和 Fiedler (1970) (e) 以及 Champagne 等人的实验数据。 （1976）（©）。横坐标是标准化的跨流坐标。 （来自米尼尔和
\end{sourceblock}
给定 W' 的 N 个独立样本 w("), n = l,2, ... N, 取自
相同的分布 f,,(B)，平均值是通过取
整体平均

\begin{sourceblock}
\includegraphics[page=538,pagebox=mediabox,trim={64.39bp 489.94bp 64.39bp 1391.82bp},clip,width=0.94\textwidth,max height=0.38\textheight]{Pope.pdf}
\end{sourceblock}
证明归一化均方根值这个近似值的误差

\begin{sourceblock}
\includegraphics[page=538,pagebox=mediabox,trim={64.39bp 309.96bp 64.39bp 1631.79bp},clip,width=0.94\textwidth,max height=0.38\textheight]{Pope.pdf}
\end{sourceblock}
被给予

\begin{sourceblock}
\includegraphics[page=538,pagebox=mediabox,trim={64.39bp 119.39bp 64.39bp 1798.77bp},clip,width=0.94\textwidth,max height=0.38\textheight]{Pope.pdf}
\end{sourceblock}
\subsection*{12.5.3 伽马分布模型}\phantomsection\addcontentsline{toc}{subsection}{12.5.3 伽马分布模型}
需要多少样本 N 才能使误差达到 1\%
(a) f, (8) 是具有 var [ln(w* (w))] 的对数正态分布
(b) f,(O) 是伽玛分布，其中 var(w*/(o)) = o 2 - - 1 7
伽马分布模型 虽然它在物理上是准确的（与 D N S 数据相比），但对数正态模型有几个不良特征：需要 h 中的临时项（在方程（12.182）中）以夹带非湍流流体；随机微分方程复杂，数值分析和实现困难；并且底层对数正态分布的非常长的尾部导致数值实现中的显着统计波动（参见练习12.33）。 Jayesh 和 Pope (1995) 开发了 w'(t) 的替代模型，旨在避免这些困难并尽可能简单。

再次从统计平稳的各向同性湍流开始，Jayesh-Pope 模型为

\begin{sourceblock}
\includegraphics[page=539,pagebox=mediabox,trim={64.39bp 881.89bp 64.39bp 1002.87bp},clip,width=0.94\textwidth,max height=0.38\textheight]{Pope.pdf}
\end{sourceblock}
其中时间尺度取为 T C3(w)，C3 为 1.0，并且是
w'/(w) 的 GO 方差，取为 o2 = 因为漂移项在 w* 中是线性的，所以自相关函数是指数 exp(-IsI/T,)。

w' 的平稳分布（参见练习 12.34）是
其中 a 1/02 4，T(a) 是伽马函数。这确实是均值 (w) 和方差 (w)2 a 2 的伽玛分布（参见练习 12.35）。

由方程式生成的 w*(t) 的样本路径(12.19 1)如图12.16所示。显然，这些并没有表现出间歇性，并且超过 3(w) 的偏移很少见。

对数正态分布和伽马 P D F 的比较如图 12.17 所示。乍一看，可能会令人惊讶的是，这些分布具有相同的均值，并且对数正态分布的方差是伽玛分布的六倍以上。解释在于对数正态分布的长尾：在 8f,(O) 和 O2f ,,(8) 的积分中，贡献分别为 13\% 和 45\%

\begin{sourceblock}
\includegraphics[page=540,pagebox=mediabox,trim={64.39bp 1209.85bp 64.39bp 109.44bp},clip,width=\textwidth,max height=0.72\textheight]{Pope.pdf}
\par\vspace{0.45\baselineskip}
\small 图 12.16。湍流频率的伽马分布模型的样本路径，方程。 （12.191）。
\end{sourceblock}
\begin{sourceblock}
\includegraphics[page=540,pagebox=mediabox,trim={64.39bp 439.95bp 64.39bp 850.89bp},clip,width=\textwidth,max height=0.72\textheight]{Pope.pdf}
\par\vspace{0.45\baselineskip}
\small 图 12.17。由对数正态模型（虚线）和伽玛分布模型（实线）给出的湍流频率的稳态 PDF。
\end{sourceblock}
将伽马分布模型扩展到非均匀流涉及直接添加源 S,,（正如在对数正态模型中，方程 (12.182) 一样），以及根据 Q 重新定义时间尺度 T。

条件平均湍流频率

\begin{sourceblock}
\includegraphics[page=541,pagebox=mediabox,trim={64.39bp 1190.85bp 64.39bp 109.44bp},clip,width=\textwidth,max height=0.72\textheight]{Pope.pdf}
\par\vspace{0.45\baselineskip}
\small 图 12.18。湍流频率 w*（在 € = 0 时通过 (w) 归一化）相对于自相似时间剪切层中归一化横向距离的散点图。虚线是无条件平均值 (w)。实线是条件平均值，方程 9。 （12.193）。 （摘自 Van Slooten、Jayesh 和 Pope (1998)。\}
\end{sourceblock}
图 12.18 将模型应用于时间混合层获得的 w* 散点图。请注意，w* 以对数刻度显示，并且跨流坐标 \& 的范围很好地延伸到间歇性湍流/非湍流区域。在层的中心，正如预期的那样，w”* 的 PDF 被发现是伽玛分布：没有 w” 小于 +(w) 的样本。在该层的边缘（例如 € = +0.8），情况完全不同。有一个 w* 值非常小的颗粒带，对应于夹带的非湍流流体，并且存在具有较高 w* 值的颗粒的散射，与层中心的颗粒相当，对应于完全湍流的流体。可以看出，平均值（w）的分布提供了非湍流样本和湍流样本之间的分离。

Jayesh 和 Pope (1995) 认为，模拟间歇区域湍流过程的适当速率 Q 是湍流流体的速率 Q。 （相反，无条件平均值 (w) 基于湍流和非湍流流体。）因此，Q 定义为

\begin{sourceblock}
\includegraphics[page=541,pagebox=mediabox,trim={64.39bp 155.98bp 64.39bp 1688.79bp},clip,width=0.94\textwidth,max height=0.38\textheight]{Pope.pdf}
\end{sourceblock}
\subsection*{12.5.4 模型联合PDF方程}\phantomsection\addcontentsline{toc}{subsection}{12.5.4 模型联合PDF方程}
在间歇区域，条件 w" 2 ( w ) 排除了非湍流粒子。因此，如图 12.18 所示，在层的边缘，R 约为其峰值的一半，并且比 (o)I 大两个数量级。在层的中心，条件 w' (w) 排除了一些湍流粒子，因此，如果 Cn 为 1，R 将超过 (w)。相反，指定 Cn 使得 R 和(a) 当 a' 的 PDF 是伽马分布时，a 相等（见练习 12.36）。

非均匀流模型（Van Slooten 和 Pope 1999）由方程 1 获得。 (12.191) 包括源项 S, 并重新定义

\begin{sourceblock}
\includegraphics[page=542,pagebox=mediabox,trim={64.39bp 1334.83bp 64.39bp 545.93bp},clip,width=0.94\textwidth,max height=0.38\textheight]{Pope.pdf}
\end{sourceblock}
因为 a* 的伽玛分布与观察到的瞬时耗散的近似对数正态分布有很大不同（见图 12.17），所以最好放弃 w' 的精确定义（方程（12.173））。相反，w* 可以被认为是湍流过程的不精确定义的特征速率，它与平均耗散有关

\begin{sourceblock}
\includegraphics[page=542,pagebox=mediabox,trim={64.39bp 961.88bp 64.39bp 985.88bp},clip,width=0.94\textwidth,max height=0.38\textheight]{Pope.pdf}
\end{sourceblock}
模型联合方程 方程的广义 Langevin 模型。 (12110) 和等式。 (12.194)

U*(t)，对于 a' (t)，得出以下速度-频率联合 PDF 方程

\begin{sourceblock}
\includegraphics[page=542,pagebox=mediabox,trim={64.39bp 418.95bp 64.39bp 1388.82bp},clip,width=0.94\textwidth,max height=0.38\textheight]{Pope.pdf}
\end{sourceblock}
从联合 PDF f(V,O;x, t) 可以确定 (U)、( a ) 、R k 和 (uiuj) 的字段，从而可以得到式 (1) 中的所有系数。 (12.196) 是已知的 f。因此，这个单一模型方程是封闭的。例如，Anand、Hsu 和 Pope (1997) 以及 Van Slooten 等人 (1998) 已经进行了基于该方程的湍流计算。

可以回忆一下第 383 页的 10.5 节，k-co 模型有
条件均值 Q，速度-频率联合 PDF 模型不会遇到这样的困难。

练习\_12.34 识别系数a@) 和! \textasciitilde{} ( 0 ) 通过写 Jayesh 定义\textasciitilde{}
一般形式的 Pope 模型（方程（12.191））

\begin{sourceblock}
\includegraphics[page=543,pagebox=mediabox,trim={64.39bp 1503.81bp 64.39bp 431.95bp},clip,width=0.94\textwidth,max height=0.38\textheight]{Pope.pdf}
\end{sourceblock}
使用方程式。 (5.21) 证明 o' 的平稳分布为

\begin{sourceblock}
\includegraphics[page=543,pagebox=mediabox,trim={64.39bp 1335.83bp 64.39bp 582.93bp},clip,width=0.94\textwidth,max height=0.38\textheight]{Pope.pdf}
\end{sourceblock}
其中 m = 1/02a 且 C 是常数。使用替换 x = a 0 / ( o )

表明归一化条件产生

\begin{sourceblock}
\includegraphics[page=543,pagebox=mediabox,trim={64.39bp 1084.86bp 64.39bp 777.90bp},clip,width=0.94\textwidth,max height=0.38\textheight]{Pope.pdf}
\end{sourceblock}
其中 r ( a ) 是伽玛函数。因此验证方程。 （12.192）。 12.35 使用第 51 页练习 3.10 的结果来表明矩
伽玛分布（方程（12.192））是

\begin{sourceblock}
\includegraphics[page=543,pagebox=mediabox,trim={64.39bp 846.89bp 64.39bp 1050.87bp},clip,width=0.94\textwidth,max height=0.38\textheight]{Pope.pdf}
\end{sourceblock}
将较高矩的值（例如，n 4）与那些值进行比较
对数正态分布（参见方程（6.333））。 12.36 f ,(0) 是伽马分布，方程： (12.192)，得到

\begin{sourceblock}
\includegraphics[page=543,pagebox=mediabox,trim={64.39bp 540.93bp 64.39bp 1291.84bp},clip,width=0.94\textwidth,max height=0.38\textheight]{Pope.pdf}
\end{sourceblock}
其中 Q(a, z) 是由下式定义的不完全伽玛函数
因此表明，如果指定了 Cn，则 (a)a 和 Q (Tq' . ( 12.193)) 相等
（对于 - 4，该等式得出 Cn 0.6893。）

\section*{12.6 拉格朗日粒子方法}\phantomsection\addcontentsline{toc}{section}{12.6 拉格朗日粒子方法}
I：在上一节中，获得了基于流体-粒子性质演化的随机模型的速度频率联合PDF（V，8；x，t）的模型方程（Eq.（12.196））。对于一般的三维流，在给定时间 t，该联合 PDF 是七个自变量（e V V V3、8、XI、x 2 和 x3）的函数。在计算上，通过七维 V-8-x 空间的离散化来精确表示联合 PDF 是不可行的，而这在有限差分、有限体积和有限元方法中是必需的。相反，模型 PDF 方程是通过粒子方法求解的。 （Hackney 和 Eastwood (1998) 的书介绍了粒子方法在各种应用中的使用。）

粒子方法的某些方面是简单且明显的。例如，存在大量粒子，每个粒子都根据随机模型方程演化。然而，还有其他一些不太明显的重要方面。本节的目的是描述粒子方法的概念基础。在大多数情况下，这是在速度 PDF f ( V ; x, t) 级别（与速度频率 PDF f 相反）最简单地完成的。

\subsection*{12.6.1 流体和颗粒系统}\phantomsection\addcontentsline{toc}{subsection}{12.6.1 流体和颗粒系统}
明确区分湍流和用于模拟湍流的粒子方法非常重要。我们分别将它们称为流体系统和粒子系统。

流体系统 湍流的基本表示形式是欧拉速度场 U(x, t)。这是由连续性和纳维-斯托克斯方程控制的，并且压力场 p(x, t) 由泊松方程 (Eq.1) 确定。 （2.42）。考虑的基本概率描述是单点一次性欧拉 PDF f ( V ;x , t)，其前两个矩是平均速度 (U(x, t)) 和雷诺应力 (uiuj) 的场。为简单起见，我们考虑体积为 V 的固定封闭区域中的流动，没有流过该区域的边界，因此流体的质量 M pV 是恒定的。

粒子系统 粒子系统的基本表示是大量 (N) 个粒子的属性。每个粒子代表一个质量 m M I N
统计上相同，因此在某些分析中考虑单个粒子就足够了，其属性用 X*(t),U *(t),e tc 表示。

X*(t)a 和 U*(t) 的一次性联合 PDF 表示为 f (; V, x ; t)。位置的边际 PDF 为

\begin{sourceblock}
\includegraphics[page=545,pagebox=mediabox,trim={64.39bp 1562.80bp 64.39bp 342.96bp},clip,width=0.94\textwidth,max height=0.38\textheight]{Pope.pdf}
\end{sourceblock}
以 X*(t) x 为条件的 U*(t) 的 PDF 为

\begin{sourceblock}
\includegraphics[page=545,pagebox=mediabox,trim={64.39bp 1419.82bp 64.39bp 506.94bp},clip,width=0.94\textwidth,max height=0.38\textheight]{Pope.pdf}
\end{sourceblock}
该 PDF 的一阶矩是条件平均粒子速度

\begin{sourceblock}
\includegraphics[page=545,pagebox=mediabox,trim={64.39bp 1264.84bp 64.39bp 638.92bp},clip,width=0.94\textwidth,max height=0.38\textheight]{Pope.pdf}
\end{sourceblock}
缩写为 (U*l x)。脉动粒子速度定义为

\begin{sourceblock}
\includegraphics[page=545,pagebox=mediabox,trim={64.39bp 1098.86bp 64.39bp 789.90bp},clip,width=0.94\textwidth,max height=0.38\textheight]{Pope.pdf}
\end{sourceblock}
使得粒子速度的条件协方差为

\begin{sourceblock}
\includegraphics[page=545,pagebox=mediabox,trim={64.39bp 873.89bp 64.39bp 972.88bp},clip,width=0.94\textwidth,max height=0.38\textheight]{Pope.pdf}
\end{sourceblock}
对应关系 重要的是要认识到流体系统和粒子系统之间存在根本差异。在流体系统中，存在一个潜在的瞬时新运动场 U(x, t)，速度梯度（例如，
ui/(8xj 8xk))、压力场p(x,t)以及多点、多时间统计(例如(Ui(x,t)uj(xr,ts)))可以被确定。另一方面，在粒子系统中，不存在潜在的瞬时速度场。与粒子相关的唯一场是平均场，例如 (U*(t)lx)。

考虑到这些固有的差异，粒子系统只能以有限的方式模拟流体系统。具体来说，我们寻求系统之间在单点、一次性统计水平上的对应关系。也就是说，如果模型是完美的，则粒子速度 f8(Vlx;t) 的条件 PDF 将等于流体速度 f ( V ;x , t) 的 PDF，因此粒子速度的矩 (U*I x) 和 (u*u*I x) 将等于平均速度 (U(x,t )) 和雷诺应力 (uiuj)+一些区别和流体系统和粒子系统之间的对应关系是

\begin{center}\small 表 12.1。流体系统和粒子系统之间的比较表明：基本表示形式根本不同；一点、一次统计水平的对应关系；以及这些统计数据的控制方程之间的对应关系\end{center}
\begin{sourceblock}
\includegraphics[page=546,pagebox=mediabox,trim={64.39bp 982.88bp 64.39bp 347.96bp},clip,width=\textwidth,max height=0.72\textheight]{Pope.pdf}
\end{sourceblock}
某些流体统计数据（例如 ui/axj aUk/8xl））在粒子系统中没有等效项。对于其他情况（例如，两点两次相关性 (Ui(x,t )Uj(x Y,t s))），可以定义粒子系统中的等效量 ( e m ,(U i(n )@ )U(mrn)( t +s)lX (n)( t) = x , X( m)( t + s) = x + Y )，其中 n 和 m 表示两个不同的粒子）。然而，这种多点、多时间统计数据的相等性不是必需的，它们并不旨在对应。

在大多数粒子方法中，粒子在统计上是独立的（或者更准确地说，它们具有：在极限 N a 中消失的弱依赖性）。因此，最好将粒子视为湍流的不同实现上的流体粒子（模型）。位于（或非常接近）同一位置 x 的两个粒子可能具有完全不同的速度，即来自 f ( V ;x , t) 的独立样本。鉴于这种独立性，上面介绍的二粒子相关性是

\begin{sourceblock}
\includegraphics[page=546,pagebox=mediabox,trim={64.39bp 242.97bp 64.39bp 1687.79bp},clip,width=0.94\textwidth,max height=0.38\textheight]{Pope.pdf}
\end{sourceblock}
这显然不符合两点、两次相关性

\subsection*{12.6.2 相应方程}\phantomsection\addcontentsline{toc}{subsection}{12.6.2 相应方程}
为了使流体和粒子统计数据对应，它们的控制方程也必须对应。例如，平均连续性方程为

\begin{sourceblock}
\includegraphics[page=547,pagebox=mediabox,trim={64.39bp 1627.79bp 64.39bp 315.96bp},clip,width=0.94\textwidth,max height=0.38\textheight]{Pope.pdf}
\end{sourceblock}
因此 ( U(x, t)) 和 ( U*( t) lx) 对应的必要条件

\begin{sourceblock}
\includegraphics[page=547,pagebox=mediabox,trim={64.39bp 1463.82bp 64.39bp 438.94bp},clip,width=0.94\textwidth,max height=0.38\textheight]{Pope.pdf}
\end{sourceblock}
诸如此类的方程的满足对粒子系统施加了约束，现在将对此进行讨论。

粒子方程 每个粒子像流体粒子一样以自己的速度运动

\begin{sourceblock}
\includegraphics[page=547,pagebox=mediabox,trim={64.39bp 1126.86bp 64.39bp 801.90bp},clip,width=0.94\textwidth,max height=0.38\textheight]{Pope.pdf}
\end{sourceblock}
而其速度则通过扩散过程而演变

\begin{sourceblock}
\includegraphics[page=547,pagebox=mediabox,trim={64.39bp 1003.87bp 64.39bp 941.88bp},clip,width=0.94\textwidth,max height=0.38\textheight]{Pope.pdf}
\end{sourceblock}
其中 a( V, x, t) 和 b(x, t) 是漂移和扩散系数。这种形式足够通用，可以包含广义朗之万模型；如果使用完全通用的扩散项，则下面获得的结果不会改变。从这些粒子方程获得的 f (; V, x ;t ) 的 Fokker-Planck 方程为
粒子位置密度 f ( V ; x , t ) 与 f a ( V l x ; t ) fL(V,x;t)/f;(x;t) 之间的对应关系不足以确定物理空间中的粒子位置密度 f;：条件 f = f * 满足 f = f f;对于 f; 的任何正规范。然而，正如现在所示，进化方程规定 fk(x, t) 必须是常数和均匀的，即

\begin{sourceblock}
\includegraphics[page=547,pagebox=mediabox,trim={64.39bp 296.96bp 64.39bp 1649.79bp},clip,width=0.94\textwidth,max height=0.38\textheight]{Pope.pdf}
\end{sourceblock}
当 f 的方程是从 f 的 Fokker-Planck 方程导出时，项

\begin{sourceblock}
\includegraphics[page=547,pagebox=mediabox,trim={64.39bp 119.39bp 64.39bp 1807.77bp},clip,width=0.94\textwidth,max height=0.38\textheight]{Pope.pdf}
\end{sourceblock}
出现，参见练习 12.38。流体系统中没有 f x 的等价物，
显然，为了使 f 和 f' 的方程相对应，SIO 项必须为零。当且仅当 f 时，才会发生这种情况；是均匀的。

还可以从 Fokker-Planck 方程（参见练习 12.37）获得 f i 的演化方程：

\begin{sourceblock}
\includegraphics[page=548,pagebox=mediabox,trim={64.39bp 1511.81bp 64.39bp 384.95bp},clip,width=0.94\textwidth,max height=0.38\textheight]{Pope.pdf}
\end{sourceblock}
为了满足粒子平均连续性方程（方程（12.211）），方程左侧是必要且充分的。 （12.217）为零。鉴于 f ;是一致的，则该条件要求 f 也为常数。 （归一化条件 f ; dx 1 也暗示了该结果。）

从定义方程。 (12.204)-(12.206)得到结果

\begin{sourceblock}
\includegraphics[page=548,pagebox=mediabox,trim={64.39bp 1096.86bp 64.39bp 806.90bp},clip,width=0.94\textwidth,max height=0.38\textheight]{Pope.pdf}
\end{sourceblock}
通过对 f Eq 的 Fokker-Planck 方程进行积分。 (12.214) 超过
所有 V，表明粒子位置密度 f k(x; t) 的演化过程为

\begin{sourceblock}
\includegraphics[page=548,pagebox=mediabox,trim={64.39bp 914.88bp 64.39bp 1008.87bp},clip,width=0.94\textwidth,max height=0.38\textheight]{Pope.pdf}
\end{sourceblock}
除以 f；验证方程。 (12.217)。

粒子条件 PDF 的演化方程
速度 f *( V lx; t) 由 Fokker-Planck 方程获得
f i ( V , x ; t) (方程 (12.214)) 除以 PDF
位置 f 假设粒子平均连续性方程
满足(Eq.(12.211))，得到结果
方程 粒子速度 '( lx ;t ) 的条件 PDF 的演化方程是通过将 Fokker-Planck 方程（方程 (12.214)）除以得到的 假设 f \& 是均匀的，结果很简单

\begin{sourceblock}
\includegraphics[page=548,pagebox=mediabox,trim={64.39bp 154.98bp 64.39bp 1743.78bp},clip,width=0.94\textwidth,max height=0.38\textheight]{Pope.pdf}
\end{sourceblock}
for f 对应于 f 的演化方程，Eq. (12 .9) 或等效方程。 （H.25）第 706 页。如果 U*(t) 的扩散模型是完美的，那么这些方程的右侧将相等。该条件提供了模型系数 a 和 b 与通常未知的流体速度场的条件统计量之间的关系。通过比较力矩方程可以提供更明确的信息。

平均动量方程

粒子平均动量方程是通过乘以方程得到的。 (12.221) 通过 V 并积分：

\begin{sourceblock}
\includegraphics[page=549,pagebox=mediabox,trim={64.39bp 1241.84bp 64.39bp 657.92bp},clip,width=0.94\textwidth,max height=0.38\textheight]{Pope.pdf}
\end{sourceblock}
其中 A(x, t ) 是条件平均漂移

\begin{sourceblock}
\includegraphics[page=549,pagebox=mediabox,trim={64.39bp 998.87bp 64.39bp 853.89bp},clip,width=0.94\textwidth,max height=0.38\textheight]{Pope.pdf}
\end{sourceblock}
这个粒子平均矩 um 方程对应于雷诺方程，（忽略粘性项）可以写成

\begin{sourceblock}
\includegraphics[page=549,pagebox=mediabox,trim={64.39bp 756.90bp 64.39bp 1139.86bp},clip,width=0.94\textwidth,max height=0.38\textheight]{Pope.pdf}
\end{sourceblock}
（第 12.7.2 节描述了近壁区域粘性效应的结合。）

雷诺方程 ( e . -V(p)/p) 的右侧是由标量 -(p)/p 的梯度给出的无旋矢量场。为了使粒子和流体平均动量 um 方程相对应，条件漂移 A 也必须是无旋的，因此可以写成标量的梯度。当.乘以 -p 时，该标量场表示为 P(x, t) 并称为粒子压力 Jield，因此 A 由下式给出

\begin{sourceblock}
\includegraphics[page=549,pagebox=mediabox,trim={64.39bp 254.97bp 64.39bp 1666.79bp},clip,width=0.94\textwidth,max height=0.38\textheight]{Pope.pdf}
\end{sourceblock}
显然，粒子和流体平均动量方程当且仅对应于平均流体压力 (p(x, t)) 和粒子压力
压力的泊松方程 A(x, t) 由方程给出。 （12.225），方程的散度。 (12.222) 是

\begin{sourceblock}
\includegraphics[page=550,pagebox=mediabox,trim={64.39bp 1655.79bp 64.39bp 264.97bp},clip,width=0.94\textwidth,max height=0.38\textheight]{Pope.pdf}
\end{sourceblock}
鉴于粒子平均连续性 y 方程（Eq. (12.2 1 1 )t），第一项为零，因此 P(x, t) 由泊松方程确定：

\begin{sourceblock}
\includegraphics[page=550,pagebox=mediabox,trim={64.39bp 1414.82bp 64.39bp 478.94bp},clip,width=0.94\textwidth,max height=0.38\textheight]{Pope.pdf}
\end{sourceblock}
该方程的 \textasciitilde{} 0 t shids 直接对应于泊松方程
在粒子系统中，没有潜在的瞬时场，但有定义为粒子属性的条件平均值的平均场（例如，(U'(t)lX*(t) x)）。对于压力来说，既不存在潜在的瞬时场，也不是可以从中获得条件平均场的粒子属性。因此，在描述粒子压力场 P(x, t) 时避免使用“平均值”一词，因为它不是底层随机对象的平均值。相反，它是作为泊松方程的解获得的。

从方程式可以清楚地看出。由式（12.226）可知，泊松方程的满足是粒子平均连续性方程成立的必要条件。因此，漂移系数的条件平均值 (A (alx) -VP/p) 由该条件确定。

广义朗之万模型 如先前所写的十（方程（12 .1 lo）），广义朗之万模型包含与流体系统和粒子系统两者相关的量。从粒子性质的角度来看，它是一致的：

dU;(t) = -- -d t Gij(Ur(t) - (U:(t)lX*(t))) dt

\begin{sourceblock}
\includegraphics[page=550,pagebox=mediabox,trim={64.39bp 456.94bp 64.39bp 1484.81bp},clip,width=0.94\textwidth,max height=0.38\textheight]{Pope.pdf}
\end{sourceblock}
12.39 对于流体系统，瞬时连续性方程为 V*U= O
由此，雅可比行列式 J(t, Y ) （方程（2.23））是统一的。

是否有一个对应于 V U 0 和一个量的方程
对应于粒子系统的 J(t, Y)？

比较关联拉格朗日和欧拉所涉及的步骤

12.40 确定漂移 a ( v , x, t) 和扩散 b(x,t) 系数
响应；到 GLM，方程。 （12.228）。验证条件均值
漂移为-VP/p。写出 f * ( VI x, t) 的演化方程并且
验证 (U*(t)lx)c 的一阶矩方程或对应于
雷诺方程。证明二阶矩方程为

\begin{sourceblock}
\includegraphics[page=551,pagebox=mediabox,trim={64.39bp 1371.83bp 64.39bp 386.95bp},clip,width=0.94\textwidth,max height=0.38\textheight]{Pope.pdf}
\end{sourceblock}
将该方程与精确的雷诺应力方程进行比较，
等式。 （7.1 94）。对应哪些术语？

\subsection*{12.6.3 均值估计}\phantomsection\addcontentsline{toc}{subsection}{12.6.3 均值估计}
到目前为止，本节的发展涉及条件手段，例如出现在广义 Langevin 模型中的 (U*\textasciitilde{} x ) 和 (utu*I x)。在数值实现中，对于大量 N 的粒子，必须根据粒子属性 X (n) (n) (t)、n 1 , 2 , ... 来估计这些条件均值。。。 ,N.通常是核估计器（如下所述）

(t)，使用U=。这里的目的不是描述特定数值方法的细节；但从概念上讲，条件均值的估计是 PDF 方法的粒子实现中的一个重要组成部分。

核估计 Q ( V ) 是给定的速度函数，我们考虑根据 N 个粒子的集合的属性来估计条件平均值 (Q(U*[t])lX*(t) x)。粒子的位置X(n)(t)是随机的，均匀分布在流域中。由于（概率为 1）在感兴趣点 x 处没有粒子，因此估计必须涉及 x 附近的粒子是不可避免的。核估计是 x 附近粒子的加权平均值，权重与指定的核函数 K(r, h) 成比例。核函数的简单示例
尺寸是

\begin{sourceblock}
\includegraphics[page=551,pagebox=mediabox,trim={64.39bp 119.39bp 64.39bp 1746.78bp},clip,width=0.94\textwidth,max height=0.38\textheight]{Pope.pdf}
\end{sourceblock}
其中 h 是指定的（正）带宽，r = lrl。一般来说，集成到unity需要一个核函数：

\begin{sourceblock}
\includegraphics[page=552,pagebox=mediabox,trim={64.39bp 1667.79bp 64.39bp 279.96bp},clip,width=0.94\textwidth,max height=0.38\textheight]{Pope.pdf}
\end{sourceblock}
对于方程给出的内核。 (12.230)，这个条件决定了情况 D 1,2 和 3 的系数 cq（参见练习 12m.41I）t。通常还要求 K 为非负数，并且有界支持 K 0
为了简单起见，我们考虑一个距离边界至少为 h 的点 x，并且我们考虑一个固定时间（在符号中没有明确指出）。核估计为

\begin{sourceblock}
\includegraphics[page=552,pagebox=mediabox,trim={64.39bp 1257.84bp 64.39bp 685.91bp},clip,width=0.94\textwidth,max height=0.38\textheight]{Pope.pdf}
\end{sourceblock}
\begin{sourceblock}
\includegraphics[page=552,pagebox=mediabox,trim={64.39bp 1088.86bp 64.39bp 828.90bp},clip,width=0.94\textwidth,max height=0.38\textheight]{Pope.pdf}
\end{sourceblock}
其中 Q(") 写为 Q(U [t])。可以看出，(Qlx)Nh 是粒子值 Q(") 的系综平均值，由以 x 为中心的核加权。事实上，等式。 (12.233) 也可以写成

\begin{sourceblock}
\includegraphics[page=552,pagebox=mediabox,trim={64.39bp 808.90bp 64.39bp 1072.86bp},clip,width=0.94\textwidth,max height=0.38\textheight]{Pope.pdf}
\end{sourceblock}
权重在哪里

\begin{sourceblock}
\includegraphics[page=552,pagebox=mediabox,trim={64.39bp 628.92bp 64.39bp 1258.84bp},clip,width=0.94\textwidth,max height=0.38\textheight]{Pope.pdf}
\end{sourceblock}
显然它们的总和为 1
估计误差 核估计中的误差

\begin{sourceblock}
\includegraphics[page=552,pagebox=mediabox,trim={64.39bp 357.95bp 64.39bp 1585.80bp},clip,width=0.94\textwidth,max height=0.38\textheight]{Pope.pdf}
\end{sourceblock}
可以分解为一个确定性部分，称为

\begin{sourceblock}
\includegraphics[page=552,pagebox=mediabox,trim={64.39bp 222.97bp 64.39bp 1718.78bp},clip,width=0.94\textwidth,max height=0.38\textheight]{Pope.pdf}
\end{sourceblock}
和随机部分，即

\begin{sourceblock}
\includegraphics[page=552,pagebox=mediabox,trim={64.39bp 119.39bp 64.39bp 1842.77bp},clip,width=0.94\textwidth,max height=0.38\textheight]{Pope.pdf}
\end{sourceblock}
其中 是标准化随机变量 ((é) = 0, = 1)，e 是

\begin{sourceblock}
\includegraphics[page=553,pagebox=mediabox,trim={64.39bp 1685.79bp 64.39bp 209.97bp},clip,width=0.94\textwidth,max height=0.38\textheight]{Pope.pdf}
\end{sourceblock}
参见练习 12.42。出现偏差是因为估计基于 x +r（对于 [|r| < h）处的粒子，其中平均值不同于 (Q|x)。如练习 12.43 所示，对于前导顺序，偏差随着 h’V?(Q|x) 变化。统计误差的产生是由于具有非零权重的样本数量有限。如练习 12.44 所示，对于 (Q|x)w 的方差的第一近似值变化为

\begin{sourceblock}
\includegraphics[page=553,pagebox=mediabox,trim={64.39bp 1294.84bp 64.39bp 607.92bp},clip,width=0.94\textwidth,max height=0.38\textheight]{Pope.pdf}
\end{sourceblock}
偏差随着 h 的增加而增加，而统计误差随着 h 的增加而减少。因此，如练习 12.45 所示，对于给定的 N，存在 h 的最佳选择，并且其变化为： 对于该选择，总误差变化为

\begin{sourceblock}
\includegraphics[page=553,pagebox=mediabox,trim={64.39bp 1038.87bp 64.39bp 897.89bp},clip,width=0.94\textwidth,max height=0.38\textheight]{Pope.pdf}
\end{sourceblock}
结论是核估计可用于根据粒子值估计平均场。随着粒子数 N 趋于无穷大，所涉及的误差趋于零，但趋于缓慢。有关核估计的更多信息可以在 Eubank (1988) 和 Hardle (1990) 中找到，其中比方程更好的核函数。 (12.230) 进行了讨论。

粒子和粒子网格方法

为了对广义 Langevin 方程（Eq. (12.228)）进行积分，必须在每个粒子位置 (n = 1,2,...,N) 评估诸如 (U(t)*|X(t)) 之类的方法。直接的方法是在每个粒子位置执行核估计。乍一看，对每个粒子执行这种估计需要 N* 级的运算。然而，可以按 N 个操作的顺序执行任务（Welton 和 Pope 1997，Welton 1998）。由此产生的方法是纯粒子方法，类似于平滑粒子流体动力学（SPH，参见例如 Monaghan (1992)）。

通常使用的替代方法是使用粒子网格方法，类似于细胞内云方法（Birdsall 和 Fuss 1969）。解域被网格覆盖，并且在网格节点处形成平均场的核估计。在这些值的基础上，整个场的平均场由线性样条表示，因此粒子位置处的值
12.6.4 以下几点总结了模型 PDF 方程数值求解的粒子方法的理论基础。

\subsection*{12.6.4 总结}\phantomsection\addcontentsline{toc}{subsection}{12.6.4 总结}
(i) 具有属性 X (n) (t)、U (n) (t) 的 N 个粒子的集合，
1,2,N可以一致地代表湍流。该党派-
cles 可以被认为是不同物体上的流体粒子（模型）

流程的实现。

(ii) 流体系统和颗粒系统对应于
一点、一次性欧拉 PDF（见表 12.1）。多粒子
统计数据与多点流体统计数据不对应。

Qiii) 使用粒子方法获得模型的数值解
PDF 方程。在这些方法中，模型方程（例如，方程（12.212），
(12.213) 和 (12.228)) 及时向前积分以确定
粒子性质的演变。

Qiv) 确定了粒子系统的几个重要属性
从演化方程之间的对应关系

流体和粒子 PDF，f ( V ; x , t) 和 f'(VIx, t)。

(a) 粒子位置f的PDF；是均匀的。最初的一个
如果粒子均值均匀分布保持均匀

满足连续性方程（见方程（12.217））。

(b) 在速度的随机模型方程（Eq. (12.213 )）中，
条件平均漂移 A（方程（12.223））由 A -VP /p 给出，
其中 P(x, t) 是粒子压力。

(c) 粒子压力（对应于平均流体
压力 ( \textasciitilde{} ( x , t ) ) ) 由泊松方程确定，
等式。 （12.227）。泊松方程的满足是必要的-
sary 的持续满足的粒子平均连续
维数方程。

(v) 方程 (12.228) 给出了广义 Langevin 模型令状 10 例
包括粒子特性。

(vi) 条件均值（例如，(U*(t)lx)）可以使用核来确定
估计。随着数量的增加，这种估计的误差趋于零。

粒子 N 趋于无穷大，但它们的速度相当缓慢（方程（12.241））。

对于方程给出的内核 K(r,h) （12.230），表明nor-

\begin{sourceblock}
\includegraphics[page=555,pagebox=mediabox,trim={64.39bp 1710.78bp 64.39bp 157.98bp},clip,width=0.94\textwidth,max height=0.38\textheight]{Pope.pdf}
\end{sourceblock}
证明核的初阶矩为零，

\begin{sourceblock}
\includegraphics[page=555,pagebox=mediabox,trim={64.39bp 1560.80bp 64.39bp 344.96bp},clip,width=0.94\textwidth,max height=0.38\textheight]{Pope.pdf}
\end{sourceblock}
第二个时刻是

\begin{sourceblock}
\includegraphics[page=555,pagebox=mediabox,trim={64.39bp 1306.84bp 64.39bp 487.94bp},clip,width=0.94\textwidth,max height=0.38\textheight]{Pope.pdf}
\end{sourceblock}
获取结果

\begin{sourceblock}
\includegraphics[page=555,pagebox=mediabox,trim={64.39bp 1190.85bp 64.39bp 744.91bp},clip,width=0.94\textwidth,max height=0.38\textheight]{Pope.pdf}
\end{sourceblock}
\begin{sourceblock}
\includegraphics[page=555,pagebox=mediabox,trim={64.39bp 999.87bp 64.39bp 829.90bp},clip,width=0.94\textwidth,max height=0.38\textheight]{Pope.pdf}
\end{sourceblock}
核估计的误差为

\begin{sourceblock}
\includegraphics[page=555,pagebox=mediabox,trim={64.39bp 878.89bp 64.39bp 1057.87bp},clip,width=0.94\textwidth,max height=0.38\textheight]{Pope.pdf}
\end{sourceblock}
从而证实分解BQ的有效性

（方程（12.238））。 12.43 回顾 X (t) 和 U (t) 的联合 PDF 为 fL(V,x,t) =
f ;(x; t) f ( V lx; t)，且 f;是均匀分布，显示
核估计中分子的期望

\begin{sourceblock}
\includegraphics[page=555,pagebox=mediabox,trim={64.39bp 385.95bp 64.39bp 1344.83bp},clip,width=0.94\textwidth,max height=0.38\textheight]{Pope.pdf}
\end{sourceblock}
证明方程中分母的期望。 (12.233) 是
团结。将 f '(Vlx - r ; t) 展开关于 x 的泰勒级数，得到

\begin{sourceblock}
\includegraphics[page=555,pagebox=mediabox,trim={64.39bp 119.39bp 64.39bp 1714.78bp},clip,width=0.94\textwidth,max height=0.38\textheight]{Pope.pdf}
\end{sourceblock}
其中内核的矩 Ki 和 Kij 定义为
等式。 （12.243）和（12.244）。由此可见，在极限 N m 下，
核估计的偏差是

\begin{sourceblock}
\includegraphics[page=556,pagebox=mediabox,trim={64.39bp 1587.80bp 64.39bp 303.96bp},clip,width=0.94\textwidth,max height=0.38\textheight]{Pope.pdf}
\end{sourceblock}
其中系数由方程定义。 （12.244）。

让 表示由方程式定义的任何一个权重。 （12.235）。显示
平均值是

\begin{sourceblock}
\includegraphics[page=556,pagebox=mediabox,trim={64.39bp 1365.83bp 64.39bp 576.93bp},clip,width=0.94\textwidth,max height=0.38\textheight]{Pope.pdf}
\end{sourceblock}
注意到，由于归一化条件 1，
权重不是相互独立的，得到结果

\begin{sourceblock}
\includegraphics[page=556,pagebox=mediabox,trim={64.39bp 1139.86bp 64.39bp 767.90bp},clip,width=0.94\textwidth,max height=0.38\textheight]{Pope.pdf}
\end{sourceblock}
考虑统计同质情况，其中 (Qlx) 是独立的
x 的下垂，其中 和 Q 是独立的。获取

\begin{sourceblock}
\includegraphics[page=556,pagebox=mediabox,trim={64.39bp 809.90bp 64.39bp 1131.86bp},clip,width=0.94\textwidth,max height=0.38\textheight]{Pope.pdf}
\end{sourceblock}
从方程。 (12.235) 认为，对于较大的 N，权重（这是一个
随机变量）可以写成十

\begin{sourceblock}
\includegraphics[page=556,pagebox=mediabox,trim={64.39bp 588.93bp 64.39bp 1358.83bp},clip,width=0.94\textwidth,max height=0.38\textheight]{Pope.pdf}
\end{sourceblock}
其中 是具有密度的均匀分布随机变量
1/V。由此得到

\begin{sourceblock}
\includegraphics[page=556,pagebox=mediabox,trim={64.39bp 360.95bp 64.39bp 1542.81bp},clip,width=0.94\textwidth,max height=0.38\textheight]{Pope.pdf}
\end{sourceblock}
其中系数由方程定义。 （12.246）。

作为解决方案的特征维度

域，表明核估计的方差为

\begin{sourceblock}
\includegraphics[page=556,pagebox=mediabox,trim={64.39bp 119.39bp 64.39bp 1791.77bp},clip,width=0.94\textwidth,max height=0.38\textheight]{Pope.pdf}
\end{sourceblock}
\section*{12.7 扩展}\phantomsection\addcontentsline{toc}{section}{12.7 扩展}
表明这可以重新表达为

\begin{sourceblock}
\includegraphics[page=557,pagebox=mediabox,trim={64.39bp 1689.79bp 64.39bp 223.97bp},clip,width=0.94\textwidth,max height=0.38\textheight]{Pope.pdf}
\end{sourceblock}
其中 Nj 是支撑内的预期粒子数
内核和 y；是一个常数。 12.45 上面的结果表明，对于首阶，均方误差
内核估计的形式为

\begin{sourceblock}
\includegraphics[page=557,pagebox=mediabox,trim={64.39bp 1414.82bp 64.39bp 494.94bp},clip,width=0.94\textwidth,max height=0.38\textheight]{Pope.pdf}
\end{sourceblock}
其中 C 和 C 是正系数（与 N 和 h 无关）

描述偏差和统计错误。表明选择h
最小化误差的是

\begin{sourceblock}
\includegraphics[page=557,pagebox=mediabox,trim={64.39bp 1161.85bp 64.39bp 723.91bp},clip,width=0.94\textwidth,max height=0.38\textheight]{Pope.pdf}
\end{sourceblock}
确定 h 的选择，从而确认等式： （12.241）。

在本节中，我们将描述基于速度和湍流频率的联合 PDF 的模型的各种扩展。与其他湍流模型的情况一样，在 PDF 方法中，可以使用以下方法处理壁

\subsection*{12.7.1 墙功能}\phantomsection\addcontentsline{toc}{subsection}{12.7.1 墙功能}
穿过粘性壁区域，并在壁上施加边界条件。这两种方法在前两小节中进行了描述。通过在 PDF 公式中添加随机单位向量（波向量）（第 12.7.3 节），可以更准确地表示快速扭曲。在第 12.7.4 节中，我们简要描述了 PDF 方法对化学反应流的扩展。这涉及到将成分变量添加到 PDF 公式中。

12.7.1 壁面函数 我们在通常的坐标系中考虑统计二维流，因此 U 和 V 是平行于和垂直于位于 y = 0 处的壁的速度分量。如第 11.7.6 节中所述，壁面函数方法的思想是在某个位置施加边界条件

\begin{sourceblock}
\includegraphics[page=558,pagebox=mediabox,trim={64.39bp 1398.82bp 64.39bp 109.44bp},clip,width=\textwidth,max height=0.72\textheight]{Pope.pdf}
\par\vspace{0.45\baselineskip}
\small 图 12.19。施加于以下位置的壁面函数的入射和反射粒子速度
\end{sourceblock}
\begin{sourceblock}
\includegraphics[page=558,pagebox=mediabox,trim={64.39bp 330.96bp 64.39bp 1608.80bp},clip,width=0.94\textwidth,max height=0.38\textheight]{Pope.pdf}
\end{sourceblock}
\begin{sourceblock}
\includegraphics[page=558,pagebox=mediabox,trim={64.39bp 119.39bp 64.39bp 1839.77bp},clip,width=0.94\textwidth,max height=0.38\textheight]{Pope.pdf}
\end{sourceblock}
此条件适当地确保 V - 为正，并且平均法向速度：y 为零。

在标准壁面函数中，轴向动量方程的边界条件相当于根据方程 1 的剪切应力 (uv) 的规范。例如，( 1 . 1 9 )。在PDF方法中，通过规范实现相同的效果

\begin{sourceblock}
\includegraphics[page=559,pagebox=mediabox,trim={64.39bp 1519.81bp 64.39bp 423.95bp},clip,width=0.94\textwidth,max height=0.38\textheight]{Pope.pdf}
\end{sourceblock}
对所得雷诺应力（练习 12.46）的直接分析表明，系数 a 的适当规格为

\begin{sourceblock}
\includegraphics[page=559,pagebox=mediabox,trim={64.39bp 1305.84bp 64.39bp 585.93bp},clip,width=0.94\textwidth,max height=0.38\textheight]{Pope.pdf}
\end{sourceblock}
产生剪切应力的原因是，在从壁反射时，颗粒的比轴向动量减少了 a1 V；如练习 12.47 所示，这些边界条件与 U* 呈联合正态分布一致。对于所考虑的统计二维流，展向速度在反射上守恒，即

\begin{sourceblock}
\includegraphics[page=559,pagebox=mediabox,trim={64.39bp 1005.87bp 64.39bp 937.88bp},clip,width=0.94\textwidth,max height=0.38\textheight]{Pope.pdf}
\end{sourceblock}
对湍流频率 o' 施加的壁面函数边界条件基于类似的想法，但细节不太简单。在对数定律区域中，随着接近壁面，平均值 (o)in 增加（y ），并且远离壁面存在正通量 ( v o )。因此，适当的边界条件会导致 a' 在反射 ( ) 上至少在平均值上增加。 Dreeben 和 Pope (1997b) 提出的条件是

\begin{sourceblock}
\includegraphics[page=559,pagebox=mediabox,trim={64.39bp 613.92bp 64.39bp 1282.84bp},clip,width=0.94\textwidth,max height=0.38\textheight]{Pope.pdf}
\end{sourceblock}
其中 是指定用于产生适当通量 (vo) 的无量纲（负）系数。指数形式与 U' 和 l no * 的联合正态分布一致。

壁面函数、边界条件已在用于求解模型 PDF 输运方程的粒子法中实现时进行了描述。就 U 和 o , f ( V , 8;x,t) 的联合 PDF 而言，边界条件是（参见练习 12.48），对于 V2 0，

\begin{sourceblock}
\includegraphics[page=559,pagebox=mediabox,trim={64.39bp 119.39bp 64.39bp 1799.77bp},clip,width=0.94\textwidth,max height=0.38\textheight]{Pope.pdf}
\end{sourceblock}
\begin{sourceblock}
\includegraphics[page=560,pagebox=mediabox,trim={64.39bp 1190.85bp 64.39bp 109.44bp},clip,width=\textwidth,max height=0.72\textheight]{Pope.pdf}
\par\vspace{0.45\baselineskip}
\small 图 12.20。 Re = 13,750 时完全发育的槽道流的湍动能剖面（以壁单位计）。 Kim 等人的符号、DNS 数据。 （1987）；使用壁面函数进行线、速度频率联合 PDF 计算（来自 Dreeben 和 Pope (1997b)）。
\end{sourceblock}
\begin{sourceblock}
\includegraphics[page=560,pagebox=mediabox,trim={64.39bp 270.97bp 64.39bp 1673.79bp},clip,width=0.94\textwidth,max height=0.38\textheight]{Pope.pdf}
\end{sourceblock}
\begin{sourceblock}
\includegraphics[page=560,pagebox=mediabox,trim={64.39bp 188.98bp 64.39bp 1752.78bp},clip,width=0.94\textwidth,max height=0.38\textheight]{Pope.pdf}
\end{sourceblock}
\begin{sourceblock}
\includegraphics[page=560,pagebox=mediabox,trim={64.39bp 119.39bp 64.39bp 1832.77bp},clip,width=0.94\textwidth,max height=0.38\textheight]{Pope.pdf}
\end{sourceblock}
\begin{sourceblock}
\includegraphics[page=561,pagebox=mediabox,trim={64.39bp 1191.85bp 64.39bp 109.44bp},clip,width=\textwidth,max height=0.72\textheight]{Pope.pdf}
\par\vspace{0.45\baselineskip}
\small 图 12.21。槽道流的表面摩擦系数 cp = tw/(\$pU8) 相对于雷诺数 (Re = 2U8/v)：符号，由 Dean (1978) 编制的实验数据：实线，使用壁面函数的速度-频率联合 PDF 计算（Dreeben 和 Pope 1997b）；虚线，使用椭圆松弛的近壁接头 PDF 计算（第 12.7.2 节，Dreeben 和 Pope (1998)）。
\end{sourceblock}
\begin{sourceblock}
\includegraphics[page=561,pagebox=mediabox,trim={64.39bp 813.90bp 64.39bp 1074.86bp},clip,width=0.94\textwidth,max height=0.38\textheight]{Pope.pdf}
\end{sourceblock}
\begin{sourceblock}
\includegraphics[page=561,pagebox=mediabox,trim={64.39bp 642.92bp 64.39bp 1305.84bp},clip,width=0.94\textwidth,max height=0.38\textheight]{Pope.pdf}
\end{sourceblock}
\begin{sourceblock}
\includegraphics[page=561,pagebox=mediabox,trim={64.39bp 559.93bp 64.39bp 1387.83bp},clip,width=0.94\textwidth,max height=0.38\textheight]{Pope.pdf}
\end{sourceblock}
V、W 和 mi。证明雅可比行列式的行列式
这些变量之间的变换是exp yp(o))]。因此
验证方程（12.267）。

\subsection*{12.7.2 近壁椭圆松弛模型}\phantomsection\addcontentsline{toc}{subsection}{12.7.2 近壁椭圆松弛模型}
壁面函数的替代方法是通过粘性近壁区域求解模型 PDF 方程，从而在壁上施加边界条件。我们在这里描述了 Dreeben 和 Pope (1998) 开发的速度-频率联合 PDF 模型的扩展，该模型解释了近壁效应。这三个扩展是结合粘度的直接影响，通过椭圆松弛方程确定 GLM 系数 Go 和 Gij，并修改与 (o)t o k 和
粘性输运 从纳维-斯托克斯方程推导出的速度 f ( V ; x, t) 的 PDF 的精确演化方程为

\begin{sourceblock}
\includegraphics[page=562,pagebox=mediabox,trim={64.39bp 955.88bp 64.39bp 941.88bp},clip,width=0.94\textwidth,max height=0.38\textheight]{Pope.pdf}
\end{sourceblock}
jX 其中 为压力波动和耗散项（在方程 (H.25) 中确定）。

前面几节中描述的 GLM 模型适用于高
vv2f 雷诺数区域，并且相应的模型 PDF 方程（方程 12.274）中缺少粘性传输项（方程 (12.274) 的右侧）。 （12.26）。然而，这个粘性项当然在粘性子层中至关重要，并且必须包含在成功的近壁模型中。 t
从微观角度看，粘性效应是由分子的随机运动产生的；因此，在 PDF 方法的粒子实现中包含粘性 \textasciitilde{}ATect 的一种自然而方便的方法是向粒子添加随机运动。具体来说，粒子位置X*(t)被重新定义，通过扩散过程演化

\begin{sourceblock}
\includegraphics[page=562,pagebox=mediabox,trim={64.39bp 299.96bp 64.39bp 1645.79bp},clip,width=0.94\textwidth,max height=0.38\textheight]{Pope.pdf}
\end{sourceblock}
其中 W'(t) 是维纳过程（独立于 GLM 中出现的过程）。根据这个定义，X*(t) 不再是流体粒子的位置：它是随着扩散而扩散的分子位置的模型
如前所述，U*(t) 是遵循以下公式的连续流体速度模型：

即，它是 4 的模型

\begin{sourceblock}
\includegraphics[page=563,pagebox=mediabox,trim={64.39bp 1680.79bp 64.39bp 264.97bp},clip,width=0.94\textwidth,max height=0.38\textheight]{Pope.pdf}
\end{sourceblock}
尽管 U(x, t) 是可微的，但 \textasciitilde{} * ( tis) 不是 ot，因为参数 X*( t) 是不可微的。相反，U*(t) 是一个具有无穷小的扩散过程

\begin{sourceblock}
\includegraphics[page=563,pagebox=mediabox,trim={64.39bp 1245.84bp 64.39bp 420.95bp},clip,width=0.94\textwidth,max height=0.38\textheight]{Pope.pdf}
\end{sourceblock}
（第一行是式（12.276）得到的增量，第二行是用式（12.275）代入dX*得到，最后的表达式是用Navier-Stokes方程代入DUi/Dt+得到）

在等式的右侧。 （12.277），速度可以分解为均值和波动。 U*( t) 的 GLM 通过包含所得的均值项进行扩展，即扩展模型为

\begin{sourceblock}
\includegraphics[page=563,pagebox=mediabox,trim={64.39bp 805.90bp 64.39bp 1053.87bp},clip,width=0.94\textwidth,max height=0.38\textheight]{Pope.pdf}
\end{sourceblock}
（请注意，这个方程是用第 12.1-12.5 节中使用的相对简单的符号编写的。在第 12.6 节更精确但麻烦的符号中，( U ) 被替换为 ( * ( t ) X * ( t ) )，并且 (p) 被粒子压力 P 替换。）

速度 P D F 的模型方程（从 X' 和 U* 的随机模型推导出来）正是从纳维-斯托克斯方程（方程 1）获得的形式。 （12.274）。特别是，它包含精确的粘性传输项 v（参见练习 12.50，方程（12.282））。因此，模型平均动量和雷诺应力方程都包含精确的粘性传递项（参见练习 12.51）。

练习 12.49 考虑层流中理想气体的混合物，以便
位置 U(x, t) 是非随机的。 a 分子的位置 X*(t) 4 请注意从 U(x, t) 获得的流体速度 u*-(t) 与由下式演变而来的流体速度 u*-(t) 之间的区别
特定物种通过随机微分方程建模

\begin{sourceblock}
\includegraphics[page=564,pagebox=mediabox,trim={64.39bp 1722.78bp 64.39bp 214.97bp},clip,width=0.94\textwidth,max height=0.38\textheight]{Pope.pdf}
\end{sourceblock}
其中 是正常数。该物种的浓度为
用 \$(x, t) 表示。这只是预期的数密度
物种的分子，它是总数的乘积
这些分子的概率密度函数和 X*(t) 的 PDF。来自福克-普朗克
该 PDF 的方程表明 \$(x, t) 的演化过程为

\begin{sourceblock}
\includegraphics[page=564,pagebox=mediabox,trim={64.39bp 1384.83bp 64.39bp 514.93bp},clip,width=0.94\textwidth,max height=0.38\textheight]{Pope.pdf}
\end{sourceblock}
评论分子运动模型的有效性，
等式。 (12+279),以及系数T。如何得到相同的结果
可以得到湍流？ 12.50 考虑六维扩散过程 Z*(t)，其中
Z* 的分量为 \{ Z , Z ;, Z;, 2 , Z ;, Z \} \{ X X i, Xi, U;, U;,
X8(t) 和 U*(t) 根据方程式演化。 (12.275) 和
(12.278)，证明 Z*(t)e 通过一般扩散方程旋转

\begin{sourceblock}
\includegraphics[page=564,pagebox=mediabox,trim={64.39bp 999.87bp 64.39bp 944.88bp},clip,width=0.94\textwidth,max height=0.38\textheight]{Pope.pdf}
\end{sourceblock}
其中 \{ W W W , W W,', WWQ'\} \{ WI, W W ,; W1、W2、W3）。

写出漂移向量扩散的显式表达式

矩阵 B i j 和扩散系数 Bij BikBjkH+ ence 使用
Fokker-Planck 方程（Eq. (5.56)）的一般形式表明
X*(t)a 和 U*(t),f L(V,x;t) 的联合 PDF 演化为

\begin{sourceblock}
\includegraphics[page=564,pagebox=mediabox,trim={64.39bp 487.94bp 64.39bp 1433.82bp},clip,width=0.94\textwidth,max height=0.38\textheight]{Pope.pdf}
\end{sourceblock}
将该方程积分到所有 V 上以获得方程
位置 f;(x; t) 的 PDF 的演变。证明这个方程是
与 f 一致，因此速度的 PDF
以 X*(t) x 为条件，f * ( Vl x; t) 也由等式演化。 （12.282）。 12.5 1 根据模型 PDF 方程，Eq. (12.282) 写为 f * (V lx; t)，

表明模型平均动量方程与

雷诺方程（包括粘性项）。

(uiuj) 与 GLM 给出的相同（方程（12.57）），但是
添加精确的粘性传输项 V2(uiuj)

右侧。

椭圆松弛 Dreeben 和 Pope (1998) 在 Durbin (1991, 1993) 方法的基础上开发了一种椭圆松弛模型来确定 GLM 系数 Gij 和 Co [第 11.8 节]。张量场的椭圆方程求解

\begin{sourceblock}
\includegraphics[page=565,pagebox=mediabox,trim={64.39bp 1338.83bp 64.39bp 548.93bp},clip,width=0.94\textwidth,max height=0.38\textheight]{Pope.pdf}
\end{sourceblock}
Gij 其中 是 Gij 的局部模型（与 LIPM 非常相似），长度尺度 LD 由等式定义。 (11.204)。然后使用解 gij 来指定 Gas

\begin{sourceblock}
\includegraphics[page=565,pagebox=mediabox,trim={64.39bp 1160.85bp 64.39bp 778.90bp},clip,width=0.94\textwidth,max height=0.38\textheight]{Pope.pdf}
\end{sourceblock}
可以看出，对于均匀湍流，对于 k(w)，这两个方程恢复为局部模型，即 Gij Gij+
让 表示局部模型中使用的 Co 值。通过构造，与每个一致的 GLM 一样，局部模型满足约束方程： (12.28)

Co Gij+ Co，关联其系数和 然而，对于 Co，模型 Gij 很可能不满足约束（均匀湍流除外）。通过通过约束确定 GLM 中使用的 Co 值（方程（12.278））可以克服这个困难

\begin{sourceblock}
\includegraphics[page=565,pagebox=mediabox,trim={64.39bp 677.91bp 64.39bp 1213.85bp},clip,width=0.94\textwidth,max height=0.38\textheight]{Pope.pdf}
\end{sourceblock}
在壁面处，gij+ 的每个分量都需要边界条件，发现满足条件的是

\begin{sourceblock}
\includegraphics[page=565,pagebox=mediabox,trim={64.39bp 509.94bp 64.39bp 1431.82bp},clip,width=0.94\textwidth,max height=0.38\textheight]{Pope.pdf}
\end{sourceblock}
其中 是壁面处的耗散率， 是单位法线，常数取 4.5。在壁位于 y = 0 的通常坐标系中，此条件产生 g22 = -yo\&,，并且所有其他分量均为零。

对所得模型雷诺应力方程的分析表明，R-e ynolds-st-res 分量非常接近壁尺度，如 (u2) y2、(v2) y3、(w2) y2 和 (uv) y3（前提是 yo 至少为 2）。这些
湍流频率和耗散 正如第 11.7 节中所讨论的，非常接近墙壁，柯尔莫哥洛夫时间尺度
是相关的物理时间尺度，而 k / i\textasciitilde{}s z ero 在
( v / \textasciitilde{} ) 1 / 2 T = 墙。因此，适当修改湍流频率的定义，使壁面处的 (o)i 为 1 阶（而不是根据关系 (o) \textasciitilde{} / k 为无限）。为此，Dreeben 和 Pope (1998) 将 (w) 与 (w) 联系起来

\begin{sourceblock}
\includegraphics[page=566,pagebox=mediabox,trim={64.39bp 1459.82bp 64.39bp 471.94bp},clip,width=0.94\textwidth,max height=0.38\textheight]{Pope.pdf}
\end{sourceblock}
与 CT 6。在高雷诺数区域，这会恢复到原始关系 ((o) elk)，而在墙壁处，它会产生 (o) C T \textasciitilde{} l，正如 Durbin 模型中的方程 1 所示。 (11.163)。壁上的 (o)a 值由 ( u ) \textasciitilde{} 表示，由等式 1 可知。 （11.160）和（12.287）：

\begin{sourceblock}
\includegraphics[page=566,pagebox=mediabox,trim={64.39bp 1155.85bp 64.39bp 746.91bp},clip,width=0.94\textwidth,max height=0.38\textheight]{Pope.pdf}
\end{sourceblock}
边界条件 在壁 (y 0) 处，需要 U 和 o, ( V, 6 ;x , t) 的连接 PDF 上的边界条件（除了已经讨论过的 gij 的边界条件，方程（12.286））。指定的边界条件是

\begin{sourceblock}
\includegraphics[page=566,pagebox=mediabox,trim={64.39bp 820.90bp 64.39bp 1112.86bp},clip,width=0.94\textwidth,max height=0.38\textheight]{Pope.pdf}
\end{sourceblock}
其中 f " 是伽马分布，平均值为 (w)，(Eq. (12.288)) 和方差 \textasciitilde{} \textasciitilde{} ( Tohe) 项 6 ( V ) 对应于确定性无滑移和不渗透条件 U(xl,0 , x3, t) -- 0。

检查非常靠近壁的粒子的行为（例如 y+1）是有启发性的。粒子位置 X*(t)（方程（12.275））的演变以扩散为主，流体速度引起的对流相比而言较小（对于 y+ 1）。

我们以墙单位考虑粒子距墙的距离

\begin{sourceblock}
\includegraphics[page=566,pagebox=mediabox,trim={64.39bp 381.95bp 64.39bp 1563.80bp},clip,width=0.94\textwidth,max height=0.38\textheight]{Pope.pdf}
\end{sourceblock}
作为墙单位时间的函数，

\begin{sourceblock}
\includegraphics[page=566,pagebox=mediabox,trim={64.39bp 244.97bp 64.39bp 1698.79bp},clip,width=0.94\textwidth,max height=0.38\textheight]{Pope.pdf}
\end{sourceblock}
忽略流体速度方程。 (12.275) 变换为

\begin{sourceblock}
\includegraphics[page=566,pagebox=mediabox,trim={64.39bp 119.39bp 64.39bp 1835.77bp},clip,width=0.94\textwidth,max height=0.38\textheight]{Pope.pdf}
\end{sourceblock}
\begin{sourceblock}
\includegraphics[page=567,pagebox=mediabox,trim={64.39bp 1178.85bp 64.39bp 109.44bp},clip,width=\textwidth,max height=0.72\textheight]{Pope.pdf}
\par\vspace{0.45\baselineskip}
\small 图 12.22。粒子距壁的距离 Y+(tt)（以壁单位为单位）作为时间的函数：反射布朗运动的样本路径，方程： （12.293）。对于给定的水平 yg，在 tj 处有一个下交叉，随后的上交叉在 cf 处。
\end{sourceblock}
其中维纳过程 W+ 满足 (dW (t+ )2 )

dt+。粒子从墙壁反射，因此 Y + 的随机过程可以写成

\begin{sourceblock}
\includegraphics[page=567,pagebox=mediabox,trim={64.39bp 867.89bp 64.39bp 1076.86bp},clip,width=0.94\textwidth,max height=0.38\textheight]{Pope.pdf}
\end{sourceblock}
该过程的示例路径称为反射布朗运动，如图 12.22 所示。

对于距墙的给定距离 y+ 0.2，图 12.22 显示了在 t l 处的下交叉（即，样本路径 Y (t ) 与线 y+ y 相交：从上方）以及随后在 t+ 处的上交叉。可以观察到，在时间间隔（t t ）内，粒子多次撞击壁。事实上，反射布朗运动的一个显着特性是，如果粒子在时间间隔 (t;, t+) 内撞击墙壁，那么（概率为 1）它会撞击墙壁
无限次（参见第 724 页的练习 5.2）。该属性 SO 适用于所有正 y+，无论多小。

事件属性U；和w；粒子的属性是碰撞壁之前的属性，反射属性 Uk 和 是碰撞壁之后的属性。反射速度指定为零 (Uk O)，而 w i 是从适当的伽马分布中采样的。关于事件属性的已知信息，U；和w；？由于反射布朗运动的行为，U；和w；在统计上与英国相同并且

\begin{sourceblock}
\includegraphics[page=568,pagebox=mediabox,trim={64.39bp 1202.85bp 64.39bp 109.44bp},clip,width=\textwidth,max height=0.72\textheight]{Pope.pdf}
\par\vspace{0.45\baselineskip}
\small 图 12.23。 Re = 13,750 时完全发展的湍流槽道流中的雷诺应力。 Kim 等人的符号、DNS 数据。 (1987)，A (u’)，© (v?)，© (w*)，好的；线、近壁速度-频率联合 PDF 计算（来自 Dreeben 和 Pope (1998)）。 ;
\end{sourceblock}
它也会在任意小的先前时间间隔内撞墙，此时粒子属性已设置为 Uk 和 m i。因此，壁上接头 PDF 的边界条件（方程 (12.289)）由反射属性 Uk 0 和 的规范强制执行。 （请注意，在壁面函数处理中，y 处的联合 PDF 与反射粒子属性的联合 PDF 不同。）

通道 JEow 的应用 图 12 .23 说明了近壁速度-频率联合 PDF 模型在完全发展的湍流槽道流中的应用。可以看出，整个通道的雷诺应力计算相当准确，并且与粘性近壁区域的 DNS 数据非常一致。模型给出的皮肤摩擦系数如图12.21所示。

\subsection*{12.7.3 波矢模型}\phantomsection\addcontentsline{toc}{subsection}{12.7.3 波矢模型}
正如第 11 章所讨论的，不可能构建一个对于均匀湍流的任意快速扭曲准确的雷诺应力模型。这是因为快速压力-应变率率张量取决于不受约束的速度相关信息的方向。

和 Reynolds (1994) 以及 Reynolds 和 Kassinos (1995)、Van Slooten 和 Pope (1997) 开发了一个 PDF 模型，该模型可以精确计算快速扭曲的均匀湍流中的应力演化。该模型适用于速度、湍流频率和单位向量的联合 PDF
包含描述快速扭曲所需的方向信息的波矢。现在描述模型的开发，首先针对快速扭曲，然后针对一般情况。

均匀湍流的快速畸变我们回顾第 11 .4 节，对于所考虑的情况，脉动速度场可以表示为傅立叶模态之和：

\begin{sourceblock}
\includegraphics[page=569,pagebox=mediabox,trim={64.39bp 1288.84bp 64.39bp 642.92bp},clip,width=0.94\textwidth,max height=0.38\textheight]{Pope.pdf}
\end{sourceblock}
每个傅里叶模式独立演化：随时间变化的波数向量 k(t) 通过常微分方程 Eq 演化。 (11.80) 从
击倒；以及相应的傅立叶系数 i i ( \textasciitilde{} "t)，初始条件 k(0) 由方程 (11 .83) 演化。

ii(rc0,t ) 的演化取决于波数向量的方向，由单位波向量给出

\begin{sourceblock}
\includegraphics[page=569,pagebox=mediabox,trim={64.39bp 884.89bp 64.39bp 1015.87bp},clip,width=0.94\textwidth,max height=0.38\textheight]{Pope.pdf}
\end{sourceblock}
但它不取决于其大小 k(t) lk(t)l。两个量
t) 和 2(t) 根据常微分 \&(KO,

\begin{sourceblock}
\includegraphics[page=569,pagebox=mediabox,trim={64.39bp 507.94bp 64.39bp 1206.85bp},clip,width=0.94\textwidth,max height=0.38\textheight]{Pope.pdf}
\end{sourceblock}
（见方程（1·1.85）和（11.86））。由于连续性方程， ii 与 k 和 e 正交：

\begin{sourceblock}
\includegraphics[page=569,pagebox=mediabox,trim={64.39bp 345.96bp 64.39bp 1600.80bp},clip,width=0.94\textwidth,max height=0.38\textheight]{Pope.pdf}
\end{sourceblock}
雷诺应力张量由下式给出

\begin{sourceblock}
\includegraphics[page=569,pagebox=mediabox,trim={64.39bp 153.98bp 64.39bp 1725.78bp},clip,width=0.94\textwidth,max height=0.38\textheight]{Pope.pdf}
\end{sourceblock}
雷诺应力可以通过求解 fi 和 2 的常微分方程来确定。

对快速畸变效应的描述是在波数空间中，而 PDF 方法是基于物理空间。因此，Van Slooten 和 Pope (1997) 的速度波矢 PDF 方法并不是先前发展的直接实现：它是基于等式 1 的精确数学类比。 (12.296)-(12.299)。

在PDF方法的粒子实现中，一般粒子具有脉动速度u*(t)和单位向量（称为波向量）

根据方程。 (12.298)，这些向量是正交的，e' (t)。

\begin{sourceblock}
\includegraphics[page=570,pagebox=mediabox,trim={64.39bp 1342.83bp 64.39bp 604.92bp},clip,width=0.94\textwidth,max height=0.38\textheight]{Pope.pdf}
\end{sourceblock}
并且它们被规定根据方程式进化。 （12.296）和（12.297）：

\begin{sourceblock}
\includegraphics[page=570,pagebox=mediabox,trim={64.39bp 1172.85bp 64.39bp 726.91bp},clip,width=0.94\textwidth,max height=0.38\textheight]{Pope.pdf}
\end{sourceblock}
这是从等式之间的精确类比得出的。 (12.296)-(12.298) 对于
t) 和 2(t) 以及等式。 (12.300)-(12.302) 对于 u*(t) 和 e*(t)（给定 i i ( \textasciitilde{} O，适当的初始条件），从粒子 ( e , ( u f u ? ) ) 确定的雷诺应力张量与从傅里叶系数获得的相同（方程 (12.299)）。因此，模型给出的雷诺应力 (urut) 的演变对于均匀湍流的任意快速扭曲。

方程(12.301)可以重写

\begin{sourceblock}
\includegraphics[page=570,pagebox=mediabox,trim={64.39bp 440.94bp 64.39bp 1308.83bp},clip,width=0.94\textwidth,max height=0.38\textheight]{Pope.pdf}
\end{sourceblock}
在快速畸变极限下，广义 Langevin 模型中的扩散消失，因此 u*(t) 的 GLM 方程，Eq. (12.126), 变为
显然，波矢模型和GLM具有相似的形式。然而，
在这里，以及本节的其余部分，上标 ' 表示粒子属性。在等式中。 (12.299)

产生快速扭曲的正确行为的系数 Gb 取决于 e*(t)；而在 GLM 中，G Mism 根据 (uiuj) 建模，并且
均匀湍流 对于均匀湍流（不在快速畸变极限）的情况，Van Slooten 和 Pope (1997) 提出的速度模型可以写成

\begin{sourceblock}
\includegraphics[page=571,pagebox=mediabox,trim={64.39bp 1423.82bp 64.39bp 471.94bp},clip,width=0.94\textwidth,max height=0.38\textheight]{Pope.pdf}
\end{sourceblock}
与 Co 2.1 和 a2 2.0。在等式中。 （12.306），平均速度梯度中的术语是准确的，并且对应于生产。在 G* 的表达式中
ij' 平均速度梯度中的项表示快速压力（参见方程（12.304））；其余项代表慢压力，对应于非线性向各向同性回归模型（对于非零 a2）。

在abse:nceo f 平均速度梯度中，方程。 (12.306) 变成一个不依赖于 e*(t) 的 Langevin 方程；而在快速畸变极限下，所有涉及的项都可以忽略不计，方程正确地恢复为方程 (12.306)。 (12.303)。

e*(t)c 的一般扩散过程可以写成

\begin{sourceblock}
\includegraphics[page=571,pagebox=mediabox,trim={64.39bp 783.90bp 64.39bp 1160.85bp},clip,width=0.94\textwidth,max height=0.38\textheight]{Pope.pdf}
\end{sourceblock}
其中维纳过程 W'(t) 独立于速度方程 W(t) 以及系数 Ei、Bij 和 B!。可以同时依赖于 u*(t)

IJ 和 e*(t)。Van Slooten 和 Pope (1997) 的模型对应于与 e*(t) 一致的最简单的系数规范，e*(t) 是与 u*(t) 正交的单位向量。所得系数（参见练习 12.52）为

\begin{sourceblock}
\includegraphics[page=571,pagebox=mediabox,trim={64.39bp 281.96bp 64.39bp 1637.79bp},clip,width=0.94\textwidth,max height=0.38\textheight]{Pope.pdf}
\end{sourceblock}
\begin{sourceblock}
\includegraphics[page=572,pagebox=mediabox,trim={64.39bp 1276.84bp 64.39bp 109.44bp},clip,width=\textwidth,max height=0.72\textheight]{Pope.pdf}
\par\vspace{0.45\baselineskip}
\small 图 12.24。 (Sk/e)y = 2.36 的均匀剪切流中雷诺应力各向异性的演变。 Van Slooten 和 Pope (1997) 的速度波矢 PDF 模型计算（线）与 Rogers 和 Moin (1987) 的 DNS 数据（符号）相比：(- ,@), bi23 (175A),
\end{sourceblock}
Van Slooten 和 Pope (1997) 将此模型应用于各种均匀湍流，并发现与实验和 DNS 数据非常一致。作为示例，图 12.24 显示了在均匀剪切流中计算的雷诺应力各向异性。

Van Slooten 等人以一种直接的方式。 (1998) 将模型扩展到非均匀流，并将其与湍流频率相结合，以获得速度、波矢和湍流频率的联合 PDF 的模型方程。该模型已应用于多种流动，包括旋流射流（Van Slooten dnd Pope 1999）。虽然这些流动的模型的性能非常好，但更简单的 LIPM 的性能也相当好。尽管如此，波矢的优点是对于均匀湍流的快速畸变来说是精确的，并且对于接近该极限的流动，其性能预计会更优越。

C \textasciitilde{} \textasciitilde{} \textasciitilde{} s itdhee gr e neral 模型为方程； （12.308），以及模型
u*(t)，方程。 (12 .306)，（为简单起见）重写为

\begin{sourceblock}
\includegraphics[page=572,pagebox=mediabox,trim={64.39bp 171.98bp 64.39bp 1770.78bp},clip,width=0.94\textwidth,max height=0.38\textheight]{Pope.pdf}
\end{sourceblock}
以下对系数的约束

\begin{sourceblock}
\includegraphics[page=573,pagebox=mediabox,trim={64.39bp 1717.78bp 64.39bp 222.97bp},clip,width=0.94\textwidth,max height=0.38\textheight]{Pope.pdf}
\end{sourceblock}
\begin{sourceblock}
\includegraphics[page=573,pagebox=mediabox,trim={64.39bp 1631.79bp 64.39bp 290.96bp},clip,width=0.94\textwidth,max height=0.38\textheight]{Pope.pdf}
\end{sourceblock}
类似地，证明 u'(t) 和 e'(t) 的正交性导致

\begin{sourceblock}
\includegraphics[page=573,pagebox=mediabox,trim={64.39bp 1466.82bp 64.39bp 471.94bp},clip,width=0.94\textwidth,max height=0.38\textheight]{Pope.pdf}
\end{sourceblock}
\begin{sourceblock}
\includegraphics[page=573,pagebox=mediabox,trim={64.39bp 1386.83bp 64.39bp 556.93bp},clip,width=0.94\textwidth,max height=0.38\textheight]{Pope.pdf}
\end{sourceblock}
表明Bij和B'的规格。由方程式(12.310) 和
(12.311) 满足方程。 （12.31 3）和（12.315）。从这些规格来看，
获得结果

\begin{sourceblock}
\includegraphics[page=573,pagebox=mediabox,trim={64.39bp 1139.86bp 64.39bp 755.90bp},clip,width=0.94\textwidth,max height=0.38\textheight]{Pope.pdf}
\end{sourceblock}
\begin{sourceblock}
\includegraphics[page=573,pagebox=mediabox,trim={64.39bp 1041.87bp 64.39bp 872.89bp},clip,width=0.94\textwidth,max height=0.38\textheight]{Pope.pdf}
\end{sourceblock}
和 Bii 0。因此表明约束方程。 (12.314) 和
(12.316) 可以重写

\begin{sourceblock}
\includegraphics[page=573,pagebox=mediabox,trim={64.39bp 738.91bp 64.39bp 1204.85bp},clip,width=0.94\textwidth,max height=0.38\textheight]{Pope.pdf}
\end{sourceblock}
并且这些可以通过 Ei 的规范来满足（方程（12.309））。

在 e*(t) 的一般扩散过程中，方程： (12.308), 取 Bij 0,
并确定系数 Ei 和 B!对应于 e'(t) 为
单位球面上的各向同性扩散。

\subsection*{12.7.4 混合和反应}\phantomsection\addcontentsline{toc}{subsection}{12.7.4 混合和反应}
PDF 方法对于处理湍流反应流特别有吸引力，因为非线性化学反应在 PDF 方程中以封闭形式出现。 Pope (19 85)、Kuznetsov 和 Sabel'nikov (19 90) 以及 Dopazo (19 94) 提供了对反应流 PDF 方法的评论。这里首先介绍单个标量 \$(x, t) 的 PDF 的主要思想，然后介绍
复合 PDF 在最简单的情况下，反应标量 +(x, 的守恒方程为

\begin{sourceblock}
\includegraphics[page=574,pagebox=mediabox,trim={64.39bp 1648.79bp 64.39bp 258.97bp},clip,width=0.94\textwidth,max height=0.38\textheight]{Pope.pdf}
\end{sourceblock}
其中 S 是化学反应的来源，它是 4 的已知函数。在许多应用中，特别是燃烧，是一种高度非线性的
S((4)) 函数，因此基于均值的源无法提供均值源 (S(4)) 的实际近似值。因此，化学源项导致平均流和二阶矩闭合中出现严重的闭合问题。

+(x, t) 的单点一次性欧拉 PDF 表示为 f 4 ( y,
X,t )，其中 y 是样本空间变量。 f 的演化方程由方程推导出来。 (12.321) 是

\begin{sourceblock}
\includegraphics[page=574,pagebox=mediabox,trim={64.39bp 1005.87bp 64.39bp 745.91bp},clip,width=0.94\textwidth,max height=0.38\textheight]{Pope.pdf}
\end{sourceblock}
\begin{sourceblock}
\includegraphics[page=574,pagebox=mediabox,trim={64.39bp 804.90bp 64.39bp 1066.87bp},clip,width=0.94\textwidth,max height=0.38\textheight]{Pope.pdf}
\end{sourceblock}
1 I+(X，其中 (ui y ) 表示 (u!(x, t) t) y)，等等（参见练习 12.54）。最重要的观察是，由于关系

\begin{sourceblock}
\includegraphics[page=574,pagebox=mediabox,trim={64.39bp 645.92bp 64.39bp 1299.84bp},clip,width=0.94\textwidth,max height=0.38\textheight]{Pope.pdf}
\end{sourceblock}
化学源以封闭形式出现。

脉动速度场导致湍流对流通量-f 4 (ul y )，它不是封闭形式的。梯度扩散模型（具有常见的缺陷）是

\begin{sourceblock}
\includegraphics[page=574,pagebox=mediabox,trim={64.39bp 387.95bp 64.39bp 1559.80bp},clip,width=0.94\textwidth,max height=0.38\textheight]{Pope.pdf}
\end{sourceblock}
其中 rT(x,t ) 是湍流扩散率。

分子扩散的影响也不是封闭形式的，

\textasciitilde{}\textasciitilde{} \$ 1 以条件拉普拉斯 y ) 的形式出现，或者以条件标量耗散 (2rVqb V4ly) 的形式出现。分子扩散效应的模型称为混合模型6，并进行了讨论

下面。早期且简单的混合模型是

\begin{sourceblock}
\includegraphics[page=575,pagebox=mediabox,trim={64.39bp 1712.78bp 64.39bp 231.97bp},clip,width=0.94\textwidth,max height=0.38\textheight]{Pope.pdf}
\end{sourceblock}
其中模型常数 C+ 通常取 2.0 是机械-
TIT+ 与标量时间尺度之比（参见第 5.5 节和练习 12.55）。这种线性确定性模型称为 IEM 模型（通过与均值交换进行交互，Villermaux 和 Devillon (1972)）或 LMSE 模型（线性均方估计，Dopazo 和 O'Brien (1974)）。

随着这些模型的结合，组合-PDF 方程（方程（12.322））变为
rT 假设 (U)、k 和 是通过湍流模型计算得知的，则方程 (1) (12.327) 是一个闭合模型方程，可以通过求解来确定 f+(y; x , t)。

该方程很容易扩展到一组成分，并且许多反应流计算都基于它，例如Chen Dibble 和 Bilger (1990)、Roekaerts (1991)、Hsu、Anand 和 Razdan (1990) 以及 Jones 和 Kakhi (1997)。这些计算采用基于节点的粒子方法（Pope (1981a)），其中有限差分网格的每个节点处都有一个粒子集合。或者，也使用了类似于第 12.6 节中描述的分布式粒子算法（Pope (1985)，参见练习 12.56），例如 Tsai 和 Fox (1996) 以及 Colucci 等人。 （1998）。

\begin{sourceblock}
\includegraphics[page=575,pagebox=mediabox,trim={64.39bp 552.93bp 64.39bp 1392.82bp},clip,width=0.94\textwidth,max height=0.38\textheight]{Pope.pdf}
\end{sourceblock}
作为合成的细粒度PDF，获得结果

\begin{sourceblock}
\includegraphics[page=575,pagebox=mediabox,trim={64.39bp 220.97bp 64.39bp 1666.79bp},clip,width=0.94\textwidth,max height=0.38\textheight]{Pope.pdf}
\end{sourceblock}
\begin{sourceblock}
\includegraphics[page=575,pagebox=mediabox,trim={64.39bp 119.39bp 64.39bp 1800.77bp},clip,width=0.94\textwidth,max height=0.38\textheight]{Pope.pdf}
\end{sourceblock}
（提示：请参阅第 703 页上的方程式 (H.8)。） 与方程式类似。 （12.3）第 464 页，
得到结果

\begin{sourceblock}
\includegraphics[page=576,pagebox=mediabox,trim={64.39bp 1692.79bp 64.39bp 250.97bp},clip,width=0.94\textwidth,max height=0.38\textheight]{Pope.pdf}
\end{sourceblock}
其中 (ul y ) 是 (u(x, t)I +(x, t) = y ) 的缩写。因此验证
等式。 （12.322）。微分方程。 （12.330）关于获得

\begin{sourceblock}
\includegraphics[page=576,pagebox=mediabox,trim={64.39bp 1480.81bp 64.39bp 438.94bp},clip,width=0.94\textwidth,max height=0.38\textheight]{Pope.pdf}
\end{sourceblock}
通过取该方程的平均值，验证方程。 （12.323）。 12.55 证明均值和的演化方程
从模型 PDF 方程获得的方差（方程（12.327））

\begin{sourceblock}
\includegraphics[page=576,pagebox=mediabox,trim={64.39bp 1199.85bp 64.39bp 660.92bp},clip,width=0.94\textwidth,max height=0.38\textheight]{Pope.pdf}
\end{sourceblock}
\begin{sourceblock}
\includegraphics[page=576,pagebox=mediabox,trim={64.39bp 946.88bp 64.39bp 914.88bp},clip,width=0.94\textwidth,max height=0.38\textheight]{Pope.pdf}
\end{sourceblock}
粒子的位置 +'(t) 演化为

\begin{sourceblock}
\includegraphics[page=576,pagebox=mediabox,trim={64.39bp 747.91bp 64.39bp 1101.86bp},clip,width=0.94\textwidth,max height=0.38\textheight]{Pope.pdf}
\end{sourceblock}
其中 a、b 和 c 是系数，W(t) 是各向同性维纳
过程。显示 X* 和 * x 的联合 PDF； t) - 进化
。 - - -。 2 斧头, 斧头, t)l
在斧头，

\begin{sourceblock}
\includegraphics[page=576,pagebox=mediabox,trim={64.39bp 405.95bp 64.39bp 1494.81bp},clip,width=0.94\textwidth,max height=0.38\textheight]{Pope.pdf}
\end{sourceblock}
对所有 y 积分该方程，得到演化方程
为粒子位置密度 f i ( x ; t )。表明，如果最初是
均匀，那么它将保持均匀，只要系数

\begin{sourceblock}
\includegraphics[page=576,pagebox=mediabox,trim={64.39bp 165.98bp 64.39bp 1724.78bp},clip,width=0.94\textwidth,max height=0.38\textheight]{Pope.pdf}
\end{sourceblock}
f;(x;t)。取f；保持一致，并证明 f；演变为
模型 PDF 方程(12.327) 前提是系数是

\begin{sourceblock}
\includegraphics[page=577,pagebox=mediabox,trim={64.39bp 1635.79bp 64.39bp 306.96bp},clip,width=0.94\textwidth,max height=0.38\textheight]{Pope.pdf}
\end{sourceblock}
\begin{sourceblock}
\includegraphics[page=577,pagebox=mediabox,trim={64.39bp 1533.81bp 64.39bp 410.95bp},clip,width=0.94\textwidth,max height=0.38\textheight]{Pope.pdf}
\end{sourceblock}
\begin{sourceblock}
\includegraphics[page=577,pagebox=mediabox,trim={64.39bp 1427.82bp 64.39bp 518.93bp},clip,width=0.94\textwidth,max height=0.38\textheight]{Pope.pdf}
\end{sourceblock}
各向同性湍流中的保守标量 P D F 组合 PDF 方法的中心问题是分子扩散效应的建模。要考虑的最简单的情况是均匀各向同性湍流中单个守恒无源标量的统计均匀场的衰减（从指定的初始条件）。 Eswaran 和 Pope (1988) 描述了这种情况的 DNS 研究，初始条件对应于统计上相同的 \$0 和 \$1 的流体斑点，其初始 PDF 为

\begin{sourceblock}
\includegraphics[page=577,pagebox=mediabox,trim={64.39bp 840.89bp 64.39bp 1098.86bp},clip,width=0.94\textwidth,max height=0.38\textheight]{Pope.pdf}
\end{sourceblock}
随着场 \$(x,t) 的演变，均值守恒，方差根据以下公式衰减

\begin{sourceblock}
\includegraphics[page=577,pagebox=mediabox,trim={64.39bp 620.92bp 64.39bp 1283.84bp},clip,width=0.94\textwidth,max height=0.38\textheight]{Pope.pdf}
\end{sourceblock}
（参见第 163 页的方程（5.281）），演化 PDF f\$(y,t ) 采用的形状如图 12.25(a) 所示。 DNS 数据显示 PDF f\$(y,t ) 在大部分时间内趋于高斯分布。

对于统计齐次守恒标量场的这种情况，复合 PDF 方程（方程 (12 .322) 和 (12 .323)）简化为

\begin{sourceblock}
\includegraphics[page=577,pagebox=mediabox,trim={64.39bp 252.97bp 64.39bp 1640.79bp},clip,width=0.94\textwidth,max height=0.38\textheight]{Pope.pdf}
\end{sourceblock}
其中 是条件标量耗散

\begin{sourceblock}
\includegraphics[page=577,pagebox=mediabox,trim={64.39bp 119.39bp 64.39bp 1818.77bp},clip,width=0.94\textwidth,max height=0.38\textheight]{Pope.pdf}
\end{sourceblock}
\begin{sourceblock}
\includegraphics[page=578,pagebox=mediabox,trim={64.39bp 1341.83bp 64.39bp 109.44bp},clip,width=\textwidth,max height=0.72\textheight]{Pope.pdf}
\par\vspace{0.45\baselineskip}
\small 图 12.25。各向同性湍流中守恒无源标量的 PDF f,(y;t) 从双δ函数初始条件的演变：(a) Eswaran 和 Pope (1988) 的 DNS； (b) 根据映射闭合计算得出（Pope 1991b）。
\end{sourceblock}
高斯场模型

研究基于高斯场的三种不同模型非常有用。首先，\$(x, t) 是高斯场的假设意味着条件拉普拉斯恰好是由 IEM 模型给出的，方程 1。 (12.326)

参见练习 12.57。该模型直接类似于扩散系数为零的简并 Langevin 方程（第 471 页的方程（12.32））。我们已经看到这个模型在保留 PDF 形状方面存在严重缺陷。从高斯初始条件它正确地给出了高斯演化。然而，从双δ函数初始条件来看，没有松弛到高斯。相反，δ 函数持续存在并在组合空间中相互靠近。

考虑的第二个模型也遵循 4(x, t) 是高斯场的假设（参见练习 12-57），即条件标量耗散和无条件标量耗散相等。然后 PDF 方程（方程（12.345））简化为

\begin{sourceblock}
\includegraphics[page=578,pagebox=mediabox,trim={64.39bp 430.95bp 64.39bp 1469.81bp},clip,width=0.94\textwidth,max height=0.38\textheight]{Pope.pdf}
\end{sourceblock}
这是扩散方程，但扩散系数为负（即
这是不稳定微分方程的典型例子：方程的解析解。 (12.347)（参见第 473 页的练习 12.9）表明，解的傅里叶系数随时间呈指数发散，对于最高模态发散最快。除非从高斯的特定初始条件开始，PDF f \$ ( y ; t ) 变为负数，并且

\$(x, t) 的高斯场假设导致基于条件拉普拉斯和条件耗散的不同有缺陷的混合模型。对于高斯初始条件，两个模型都正确预测了高斯演化。然而，对于其他初始条件，这两个模型表现出不同类型的令人不满意的行为。

练习12.57 令\#(x) 为统计上均匀的高斯场。然后，在
每个地点的数量 8+/8xi 和 V2\$ 都是正常的
随机变量。证明 和 之间的协方差为

\begin{sourceblock}
\includegraphics[page=579,pagebox=mediabox,trim={64.39bp 1326.83bp 64.39bp 611.92bp},clip,width=0.94\textwidth,max height=0.38\textheight]{Pope.pdf}
\end{sourceblock}
因此，使用方程。 (3.119) 第 63 页，表明条件

\begin{sourceblock}
\includegraphics[page=579,pagebox=mediabox,trim={64.39bp 1124.86bp 64.39bp 727.91bp},clip,width=0.94\textwidth,max height=0.38\textheight]{Pope.pdf}
\end{sourceblock}
证明这与 IEM 模型相同，方程 1。 （12.326）。显示
那和 8\$/8xi 不相关，因此是独立的（因为
它们都是正常的）。因此表明条件标量
耗散（方程（12.346））等于无条件标量
耗散（需要强调的是，这些结果特定于
高斯场，一般不适用。）

映射闭合 对于非高斯初始 PDF f\#(y;O)，\$(x,t) 是高斯场的假设显然是不一致的。相反，映射闭合（Chen、Chen 和 Kraichnan 1989、Pope 1991b、Gao 和 O'Brien 1991）提供了一致的高斯假设并得出令人满意的模型。

令 8(x) 为统计齐次高斯场，令 q 为对应于 0 的样本空间变量。在每个时间 t 都有一个唯一定义的非递减函数 X(q, t) - “映射” - 使得代理 \$eld

\begin{sourceblock}
\includegraphics[page=579,pagebox=mediabox,trim={64.39bp 297.96bp 64.39bp 1648.79bp},clip,width=0.94\textwidth,max height=0.38\textheight]{Pope.pdf}
\end{sourceblock}
与合成域 \$(x,t) 具有相同的单点 PDF。通过假设组合字段 \$@, t) 的统计数据与（非高斯）代理项的统计数据相同来实现闭合
根据映射已知。这个假设产生了映射 X ( q ,t ) 演化的方程，它隐式地确定了 PDF。

对于双δ函数初始条件的测试用例，图12.25（b）显示了映射闭合预测的PDF形状的演变。显然，与 DNS 数据非常吻合。在大多数情况下，映射闭合给出的 PDF 趋于高斯分布。对于单个
Pope (1991b) 描述了映射闭合的粒子实现。

速度合成 P DF
发生化学反应的湍流中流体的热化学成分可以用标量 \# 42,... , v \} 来描述，具体取决于所涉及的化学性质，标量的数量可以是两个或三个，大约十个，甚至是 100 个量级。我们现在将前面的考虑扩展到速度和一组标量的联合 PDF。为简单起见，我们采用标量根据守恒方程演化

\begin{sourceblock}
\includegraphics[page=580,pagebox=mediabox,trim={64.39bp 1020.87bp 64.39bp 888.89bp},clip,width=0.94\textwidth,max height=0.38\textheight]{Pope.pdf}
\end{sourceblock}
其中反应源 S 是组合物的已知函数。对于每种组合物，分子扩散被认为是恒定的、均匀的且相同的。 （实际上，物种的差异扩散可能很重要；参见 Yeung (1998)、Nilsen 和 Kosaly (1999) 以及其中的参考文献。）

速度 U(x, t) 和复合 \#(x, t) 的单点一次性欧拉联合 P D F 表示为 f ( V , e ; x , t )，其中 = , y \} a re

\{yl,y2,... 组合的样本空间变量。 f ( V, ;x , t) 的演化方程由纳维-斯托克斯方程和方程推导出来。 (12.351) 是

\begin{sourceblock}
\includegraphics[page=580,pagebox=mediabox,trim={64.39bp 386.95bp 64.39bp 1509.81bp},clip,width=0.94\textwidth,max height=0.38\textheight]{Pope.pdf}
\end{sourceblock}
左侧的项是封闭形式，最显着的是对流和反应。因此，避免了组合 PDF 模型中使用的梯度扩散模型，同时保留了封闭形式反应的优点。

对于右侧未封闭的项，期望是 con-
位置是被动的，它们的值对 p、v、U 和 p' 没有影响。因此，右侧第一项可以通过广义朗之万模型独立于 进行建模。

最简单的混音模型还是 IEM

\begin{sourceblock}
\includegraphics[page=581,pagebox=mediabox,trim={64.39bp 1559.80bp 64.39bp 348.96bp},clip,width=0.94\textwidth,max height=0.38\textheight]{Pope.pdf}
\end{sourceblock}
cf \textasciitilde{} q (1.2. 326)。根据该模型，条件拉普拉斯算子与 V 无关，这是不正确的（Pope 1998b）。已经提出了其他几个包含对 V 的依赖性的模型（例如，P ope (1985) 和 Fox (1996)）。

这种方法很容易扩展到速度、湍流频率和成分的联合 PDF，以生成湍流反应流的简单闭合模型 PDF 方程。 12.6 节中描述的粒子方法允许求解该方程，即使联合 PDF 是许多自变量的函数。例如，Saxena 和 Pope (19 98) 描述了基于速度、湍流频率和 14 种成分（即 f 、
写下速度复合联合 PDF 方程 (Eq.

12.352))，使用广义 Langevin 模型（参见方程 (12 .26)）和
IEM 模型（方程（12.353））替代了两个未封闭项
在右侧。乘以适当的样本空间
变量并积分以获得平均演化方程

\begin{sourceblock}
\includegraphics[page=581,pagebox=mediabox,trim={64.39bp 394.95bp 64.39bp 1497.81bp},clip,width=0.94\textwidth,max height=0.38\textheight]{Pope.pdf}
\end{sourceblock}
将模型联合 PDF 方程乘以 V,yp 并积分为

\begin{sourceblock}
\includegraphics[page=581,pagebox=mediabox,trim={64.39bp 119.39bp 64.39bp 1723.78bp},clip,width=0.94\textwidth,max height=0.38\textheight]{Pope.pdf}
\end{sourceblock}
由此得到模型标量通量方程

\begin{sourceblock}
\includegraphics[page=582,pagebox=mediabox,trim={64.39bp 1607.80bp 64.39bp 218.97bp},clip,width=0.94\textwidth,max height=0.38\textheight]{Pope.pdf}
\end{sourceblock}
混合模型 由于IEM和其他简单模型的缺陷，人们进行了许多尝试来开发更好的混合模型（其中大多数由Dopazo（1994）审查）。尽管有所改进，但还没有一个模型在各方面都令人满意。

理想混合模型的一些行为和属性可以从守恒方程中推导出来，其他行为和属性可以从实验和 DNS 数据中推导出来。对于一组具有相等扩散率的守恒标量，以下应该成立。

(i) 该方法不受混合的直接影响。

(ii) 混合导致方差 ) 随着时间的推移而减小。 （一
更强的条件是它导致协方差的特征值

ance 矩阵减小。）

(iii) 有界条件的推广（第 (2.55) 页）

21) 是样本空间中由
标量随时间减少。令 C(t) 表示最小凸区域
在时间 t 时组合物占据的样本空间中，因此
对于 C(t) 之外的所有组合，PDF 为零。然后，对于 t2 > tl，
C(t2)包含在C(t1)内。

(iv) 理想的混合模型满足线性不变性
和守恒定律所隐含的独立性

(v) 对于均匀湍流中的统计均匀标量场
因此，联合 PDF f \& (@t,) 趋向于联合法线（Juneja 和
在更一般的情况下，理想的混合模型将准确地考虑
(六) 差异扩散
(vii) 标量场长度尺度的影响，以及

\section*{12.8 讨论}\phantomsection\addcontentsline{toc}{section}{12.8 讨论}
\begin{center}\small 表 12.2。不同级别的PDF模型\end{center}
\begin{sourceblock}
\includegraphics[page=583,pagebox=mediabox,trim={64.39bp 1230.85bp 64.39bp 220.97bp},clip,width=\textwidth,max height=0.72\textheight]{Pope.pdf}
\end{sourceblock}
如前所述，IEM 模型保留了 PDF 的形状，与 (v) 相反。这一缺陷可以通过添加扩散项 来弥补，这样 ( ) 就可以通过 Langevin 方程建模，但这样就会违反有界条件 (iii)。对于简单边界，可以构造满足有界性的扩散过程（例如，Valifio 和 Dopazo (1991) 以及 Fox (1992,1994)），但一般情况仍然存在问题。对于一般情况，令人满意的、易于处理的映射闭合扩展也被证明是难以捉摸的。

PDF 方法的粒子实现允许使用基于粒子相互作用的模型（Curl 1963、Janicka、Kolbe 和 Kollmann 1977、Pope 1985、Subramaniam 和 Pope 1998）。这些模型很容易满足一般有界条件 (iii)，并提供 PDF 形状的松弛。为了解决 (vii)，Fox (1997) 开发了一个模型，其中包含有关标量场的光谱信息。

表 12.2 显示了本章介绍的一些 PDF 模型。为了简单起见，仅在这些模型中的第一个中示出了组合物，但是应该理解的是，所有其他模型也可以容易地扩展以包括组合物。本节将根据第 8 章中描述的评估模型的标准讨论这些 PDF 模型。在本讨论中，很自然地将 PDF 方法与雷诺应力进行比较
描述级别 各种联合 PDF 提供了所考虑的流量的完整的单点、一次性统计描述。一般来说，这比均值和协方差中包含的信息多得多。对于某些流动（例如均匀湍流），这种更高级别的描述可能比二阶矩闭合没有什么好处。然而，一般来说，更全面的描述有利于允许准确地处理更多过程，并提供更多可用于构造闭合模型的信息。具体来说，
(i) 在包含速度的 PDF 模型中，湍流对流处于封闭状态
形成梯度扩散模型（例如，对于 (uiujuk) 和
被回避；

[0070] (ii)当包含组合物时，反应是封闭形式的，
与 moment 遇到的闭合问题形成鲜明对比
(iii) 在波矢模型中，处理压力-应变率率张量
完全处于快速畸变极限（对于均匀湍流）；

(iv) 根据湍流频率的分布，条件
可以得到平均湍流频率 R 并将其用于 mod-
elling，从而消除自由流边界条件问题
在 k-w' 模型中遇到。

完整性 所有考虑的 PDF 模型都是完整的。

成本和易用性与雷诺应力模型相比，PDF 方法在此处于劣势。 12.6 节中描述的粒子方法在算法和代码的可用性方面都不太成熟
与通常用于 k-E 和雷诺应力计算的有限体积方法相比。在典型的粒子方法计算中，粒子的数量可能比典型的有限体积计算中的单元数量多 100 倍。因此，毫不奇怪，粒子方法的 CPU 时间要求较大，但不到 100 倍。例如，Hsu 等人。 (1990) 在比较组合 PDF 计算与 k--E 计算时引用了因子 5。

然而，与 LES 相比，PDF 方法的计算成本
远低于积分尺度的问题不需要解决，并且可以利用统计对称性来降低问题的维数。

适用范围 与 k-E 和雷诺应力模型一样，PDF 方法可以应用于任何湍流（在本书考虑的类别内）。 PDF 方法在反应流中得到了最广泛的应用，因为非线性化学
可以在没有闭合近似的情况下进行处理

精度 用于处理均匀湍流的 PDF 方法可以构造为等同于任何可实现的雷诺应力闭合，因此（对于这些流动）两种方法的精度是相当的。对于非均匀流，速度-PDF方法的理论优势是湍流对流传输是封闭形式的；波矢模型对于接近快速畸变的流动具有理论优势。然而，由于缺乏对雷诺应力和 PDF 模型计算的比较研究，很难评估 PDF 方法在处理各种流动时所获得的精度增益。

大涡流

\section*{13.1 引言}\phantomsection\addcontentsline{toc}{section}{13.1 引言}
\chapter{大涡模拟}
直接表示大尺度运动，而对较小尺度运动的影响进行建模。在计算成本方面，LES 介于雷诺应力模型和 DNS 之间，并且它的动机是这些方法的局限性。由于大规模非定常运动被明确表示，因此对于大规模非定常性显着的流动（例如涉及非定常分离和涡旋脱落的流过钝体的流动），LES 预计比雷诺应力模型更准确、更可靠。

正如第9章所讨论的，DNS的计算成本很高，并且随着雷诺数的立方而增加，因此DNS不适用于高雷诺数流。 DNS 中几乎所有的计算工作都花费在最小的耗散运动上（参见第 351 页的图 9.4），而能量和各向异性主要包含在较大尺度的运动中。在 LES 中，较大尺度运动（受流几何影响且不具有普遍性）的动力学被明确计算，较小尺度（在某种程度上具有普遍性）的影响由简单模型表示。因此，与 DNS 相比，可以避免显式表示小规模运动的巨大计算成本。

LES 有四个概念步骤。

(i) 定义过滤操作将速度 U(x, t) 分解为
过滤（或解析）分量 U(x, t ) 和残差的总和
（或亚格子尺度，SGS）分量 J(x, t)。滤波后的速度场
U(x, t) 是三维且与时间相关的 - 表示
大涡流的运动。

\begin{center}\small 表 13.1。 DNS 和 LES 的某些变体中的解析\end{center}
\begin{sourceblock}
\includegraphics[page=588,pagebox=mediabox,trim={64.39bp 1207.85bp 64.39bp 214.97bp},clip,width=\textwidth,max height=0.72\textheight]{Pope.pdf}
\end{sourceblock}
(ii) 使用有限差分法（在物理上）的各向同性湍流
(iii) 使用均匀矩形网格的自由剪切流，
(iv) 使用非均匀矩形充分发展湍流槽道流
喉咙网格，
(v) 使用矩形网格的大气边界层，
(vi) 使用结构化网格流过钝体的流动，以及
(vi:i) 使用非结构化网格在复杂几何形状中流动。

LES 变体之间存在重要区别，表 13.1 中总结了这些区别。首先，考虑到远离壁面的流动，我们对 LES 和 VLES（超大涡模拟）进行了区分。在 LES 中，过滤后的速度场占流场中各处湍动能的大部分（例如 80\%）。在 VLES 中，网格和滤波器太大，无法解析含能量的运动，而是大部分能量存在于残余运动中。尽管 VLES 可以在较粗糙的网格上执行，因此成本较低，但模拟更加依赖于残余运动的建模。在实践中，很少估计所解析的能量比例，因此并不总是清楚特定的模拟是 LES 还是 VLES。

对于壁面边界流动，LES 和 VLES 之间的相同区别也适用于远离壁面的流动；进一步的区别取决于近壁运动的处理。对于光滑壁，近壁运动的粘性长度尺度为 6（与流动尺度 6 相比，其随着雷诺数的减小而减小）。如果过滤器和网格

\section*{13.2 滤波}\phantomsection\addcontentsline{toc}{section}{13.2 滤波}
解决了，那么结果就是近壁分辨的大涡模拟，LES-NWR。这需要靠近壁的非常精细的网格，并且计算成本随着雷诺数的幂而增加（第 13.4.5 节），因此，与 DNS 一样，LES-NWR 对于高雷诺数流是不可行的。另一种方法是使用近壁建模的大涡模拟，LES-NWM，其中滤波器和网格太粗糙，无法解析近壁运动，因此可以对它们的影响进行建模（显式或隐式）。 LES 的气象应用是 LES-NWM 或 VLES。

本章概述 在接下来的两节中，通过检查滤波操作及其在纳维-斯托克斯方程中的应用来开发 LES 公式。对于在物理空间中执行的 LES，基本残余应力模型是 Smagorinsky (1963) 提出的涡粘模型，该模型在 13.4 节中进行了描述和检验。在考虑波数空间中的 LES（第 13.5 节）后，进一步的残余应力模型将在第 13.,6 节中描述。本章最后对 LES 及其变体进行了评估。

1 3.2 过滤 在 DNS 中，速度场 U(x, t ) 必须在长度尺度上解析到柯尔莫哥洛夫尺度 q-。在 LES 中，执行低通滤波操作，以便可以在相对粗糙的网格上充分解析所得的滤波速度场 U(x, t)。具体来说，所需的网格间距 1s 与指定的滤波器宽度成正比。在理想情况下，滤波器宽度
eEI 比最小的含能运动的尺寸稍小（参见第 188 页的图 6.2）。因此，在解决含能运动的条件下，网格间距尽可能大。

\subsection*{13.2.1 一般定义}\phantomsection\addcontentsline{toc}{subsection}{13.2.1 一般定义}
一般过滤操作（由 Leonard (1974) 引入）定义为

\begin{sourceblock}
\includegraphics[page=589,pagebox=mediabox,trim={64.39bp 279.96bp 64.39bp 1623.80bp},clip,width=0.94\textwidth,max height=0.38\textheight]{Pope.pdf}
\end{sourceblock}
其中积分是在整个流域上进行的，并且指定的过滤函数 G 满足归一化条件

\begin{sourceblock}
\includegraphics[page=589,pagebox=mediabox,trim={64.39bp 119.39bp 64.39bp 1808.77bp},clip,width=0.94\textwidth,max height=0.38\textheight]{Pope.pdf}
\end{sourceblock}
在最简单的情况下，滤波器函数是齐次的，即独立于
残差场定义为

\begin{sourceblock}
\includegraphics[page=590,pagebox=mediabox,trim={64.39bp 1638.79bp 64.39bp 303.96bp},clip,width=0.94\textwidth,max height=0.38\textheight]{Pope.pdf}
\end{sourceblock}
使得速度场有分解

\begin{sourceblock}
\includegraphics[page=590,pagebox=mediabox,trim={64.39bp 1501.81bp 64.39bp 441.94bp},clip,width=0.94\textwidth,max height=0.38\textheight]{Pope.pdf}
\end{sourceblock}
这看起来类似于雷诺分解。然而，重要的区别是 U(x, t) 是一个随机场，并且（通常）过滤后的残差不为零：

\begin{sourceblock}
\includegraphics[page=590,pagebox=mediabox,trim={64.39bp 1288.84bp 64.39bp 656.92bp},clip,width=0.94\textwidth,max height=0.38\textheight]{Pope.pdf}
\end{sourceblock}
由方程式可知(13.1) 滤波和微分运算
就通勤时间而言，饮食是.,

\begin{sourceblock}
\includegraphics[page=590,pagebox=mediabox,trim={64.39bp 1002.87bp 64.39bp 900.89bp},clip,width=0.94\textwidth,max height=0.38\textheight]{Pope.pdf}
\end{sourceblock}
表明过滤和取平均通勤的操作，

\begin{sourceblock}
\includegraphics[page=590,pagebox=mediabox,trim={64.39bp 831.89bp 64.39bp 1056.87bp},clip,width=0.94\textwidth,max height=0.38\textheight]{Pope.pdf}
\end{sourceblock}
微分方程。 （13.1）关于获得结果

表明过滤和微分操作具有重新

观点到位置一般不会交换，但这样做是为了同质化
新过滤器。

13.2.2 各种滤波器的特性最简单地在一维上进行检查。因此，我们考虑为所有 x < x < m) 定义的随机标量函数 U(x)。三维向量情况的推广

\subsection*{13.2.2 一维过滤}\phantomsection\addcontentsline{toc}{subsection}{13.2.2 一维过滤}
（-CCI 很简单；考虑到这一点，我们将 U(x) 称为“速度”

\begin{center}\small 表 13.2。一维滤波器的滤波器函数和传递函数：盒式滤波器函数具有与高斯相同的二阶矩（+，A？）；其他滤波器在特征波数 x, = 7/A 处具有相同的传递函数值，即 G(k,) =\end{center}
\clearpage
\begin{sourceblock}
\includegraphics[page=591,pagebox=mediabox,trim={64.39bp 269.97bp 64.39bp 345.96bp},clip,width=\textwidth,max height=0.72\textheight]{Pope.pdf}
\end{sourceblock}
图 13.2 显示了使用 A 0.35 的高斯滤波器获得的样本速度场 U(x) 和相应的滤波场 V(x)。它是
V ( x ) 显然遵循 U(x) 的总体趋势，但短期

\begin{sourceblock}
\includegraphics[page=592,pagebox=mediabox,trim={64.39bp 1074.87bp 64.39bp 109.44bp},clip,width=\textwidth,max height=0.72\textheight]{Pope.pdf}
\par\vspace{0.45\baselineskip}
\small 图 13.1。滤波器G(r)：箱式滤波器，虚线；高斯滤波器，实线；锐利的光谱滤光片，点虚线。
\end{sourceblock}
\begin{sourceblock}
\includegraphics[page=592,pagebox=mediabox,trim={64.39bp 277.97bp 64.39bp 978.88bp},clip,width=\textwidth,max height=0.72\textheight]{Pope.pdf}
\par\vspace{0.45\baselineskip}
\small 图 13.2。上曲线：速度场 U(x) 和相应的滤波场 U(x)（粗线）的样本，使用 A \textasciitilde{} 0,35 的高斯滤波器。下方曲线：残差场 u’(x) 和滤波后的残差场 w’(x)（粗线）。
\end{sourceblock}
\begin{sourceblock}
\includegraphics[page=593,pagebox=mediabox,trim={64.39bp 1114.86bp 64.39bp 109.44bp},clip,width=\textwidth,max height=0.72\textheight]{Pope.pdf}
\par\vspace{0.45\baselineskip}
\small 图 13.3。滤波器传递函数G(x)：盒式滤波器，虚线；高斯滤波器，实线；锐利的光谱滤光片，点划线。
\end{sourceblock}
字段 ut(x) 也如图 13.2 所示。如前所述，如可以观察到的，过滤后的残差是非零的。

\subsection*{13.2.3 谱表示}\phantomsection\addcontentsline{toc}{subsection}{13.2.3 谱表示}
滤波器的效果在波数空间中最清楚地显现出来。假设 U ( x ) 有傅里叶变换

\begin{sourceblock}
\includegraphics[page=593,pagebox=mediabox,trim={64.39bp 612.92bp 64.39bp 1330.83bp},clip,width=0.94\textwidth,max height=0.38\textheight]{Pope.pdf}
\end{sourceblock}
那么滤波后的速度的傅里叶变换是

\begin{sourceblock}
\includegraphics[page=593,pagebox=mediabox,trim={64.39bp 420.95bp 64.39bp 1461.82bp},clip,width=0.94\textwidth,max height=0.38\textheight]{Pope.pdf}
\end{sourceblock}
其中传递函数 G(K) 是滤波器傅里叶传递的 2n 倍：

\begin{sourceblock}
\includegraphics[page=593,pagebox=mediabox,trim={64.39bp 246.97bp 64.39bp 1659.79bp},clip,width=0.94\textwidth,max height=0.38\textheight]{Pope.pdf}
\end{sourceblock}
该结果直接由滤波方程（方程（13.9））和卷积定理（第 681 页的方程（D.15））得出。

如图13.3所示。由于考虑的滤波器函数 G(r) 是实数且偶数，因此传递函数也是实数和偶数。在原点处，传递函数是统一的，因为归一化条件是

\begin{sourceblock}
\includegraphics[page=594,pagebox=mediabox,trim={64.39bp 1598.80bp 64.39bp 303.96bp},clip,width=0.94\textwidth,max height=0.38\textheight]{Pope.pdf}
\end{sourceblock}
尖锐光谱滤波器的重要性现在显而易见：它消除了波数大于截止波数的所有傅里叶模式

\begin{sourceblock}
\includegraphics[page=594,pagebox=mediabox,trim={64.39bp 1416.82bp 64.39bp 531.93bp},clip,width=0.94\textwidth,max height=0.38\textheight]{Pope.pdf}
\end{sourceblock}
而它对较低波数模式没有影响。然而，尽管它在波数空间中是尖锐的，但尖锐的光谱滤波器在物理空间中显然是非局域的。对于箱式过滤器来说，情况正好相反。在所考虑的滤波器中，只有高斯滤波器在物理空间和波数空间中都相当紧凑。

（表 13.2 中定义的柯西滤波器和 Pao 滤波器并未在实践中使用，但它们具有第 13.4 节中揭示的理论意义。它们的传递函数和滤波器函数与下面图 13.8 和 13.9 中的高斯滤波器函数进行了比较。）

让U(x)进行傅里叶变换(方程(13.0))，使得U(x)

FI(LC)

具有傅立叶变换 G ( I c ) \textasciitilde{}\textasciitilde{} ( I(cE)q . (13.11))。表明
残差 u'(x) 的傅立叶变换为

\begin{sourceblock}
\includegraphics[page=594,pagebox=mediabox,trim={64.39bp 687.91bp 64.39bp 1248.84bp},clip,width=0.94\textwidth,max height=0.38\textheight]{Pope.pdf}
\end{sourceblock}
滤波残差 u' 的傅立叶变换为

\begin{sourceblock}
\includegraphics[page=594,pagebox=mediabox,trim={64.39bp 545.93bp 64.39bp 1385.83bp},clip,width=0.94\textwidth,max height=0.38\textheight]{Pope.pdf}
\end{sourceblock}
双重滤波场 U(x) 的傅立叶变换为

\begin{sourceblock}
\includegraphics[page=594,pagebox=mediabox,trim={64.39bp 379.95bp 64.39bp 1533.81bp},clip,width=0.94\textwidth,max height=0.38\textheight]{Pope.pdf}
\end{sourceblock}
表明两个方程。 (13.4) 和上面的方程得出结果

\begin{sourceblock}
\includegraphics[page=594,pagebox=mediabox,trim={64.39bp 243.97bp 64.39bp 1690.79bp},clip,width=0.94\textwidth,max height=0.38\textheight]{Pope.pdf}
\end{sourceblock}
对于尖锐的光谱滤波器，表明

\begin{sourceblock}
\includegraphics[page=594,pagebox=mediabox,trim={64.39bp 119.39bp 64.39bp 1831.77bp},clip,width=0.94\textwidth,max height=0.38\textheight]{Pope.pdf}
\end{sourceblock}
Hencle 获得结果（对于锐光谱滤波器）

\begin{sourceblock}
\includegraphics[page=595,pagebox=mediabox,trim={64.39bp 1721.78bp 64.39bp 216.97bp},clip,width=0.94\textwidth,max height=0.38\textheight]{Pope.pdf}
\end{sourceblock}
\begin{sourceblock}
\includegraphics[page=595,pagebox=mediabox,trim={64.39bp 1642.79bp 64.39bp 301.96bp},clip,width=0.94\textwidth,max height=0.38\textheight]{Pope.pdf}
\end{sourceblock}
更一般地，表明过滤后的残差为零（方程（13.21））

当且仅当过滤操作是投影，即它产生
让 ( ) 表示应用宽度高斯滤波器的操作

\begin{sourceblock}
\includegraphics[page=595,pagebox=mediabox,trim={64.39bp 1286.84bp 64.39bp 603.92bp},clip,width=0.94\textwidth,max height=0.38\textheight]{Pope.pdf}
\end{sourceblock}
并令 G(K;A ) 表示相应的传递函数，即

\begin{sourceblock}
\includegraphics[page=595,pagebox=mediabox,trim={64.39bp 1134.86bp 64.39bp 813.90bp},clip,width=0.94\textwidth,max height=0.38\textheight]{Pope.pdf}
\end{sourceblock}
获取结果

\begin{sourceblock}
\includegraphics[page=595,pagebox=mediabox,trim={64.39bp 922.88bp 64.39bp 950.88bp},clip,width=0.94\textwidth,max height=0.38\textheight]{Pope.pdf}
\end{sourceblock}
\begin{sourceblock}
\includegraphics[page=595,pagebox=mediabox,trim={64.39bp 857.89bp 64.39bp 1083.86bp},clip,width=0.94\textwidth,max height=0.38\textheight]{Pope.pdf}
\end{sourceblock}
因此表明

\begin{sourceblock}
\includegraphics[page=595,pagebox=mediabox,trim={64.39bp 673.91bp 64.39bp 1206.85bp},clip,width=0.94\textwidth,max height=0.38\textheight]{Pope.pdf}
\end{sourceblock}
\begin{sourceblock}
\includegraphics[page=595,pagebox=mediabox,trim={64.39bp 540.93bp 64.39bp 1344.83bp},clip,width=0.94\textwidth,max height=0.38\textheight]{Pope.pdf}
\end{sourceblock}
因此验证盒子和高斯滤波器的二阶矩

是A2/12。 .评论对锐光谱和
柯西滤波器。

一组 N 个基函数 cpn(x), n 1,2, ,N , 定义在
间隔 0 5 x 5 给定速度 U(x)，滤波后的速度
U(x) 定义为 U(x) 在基底上的最小二乘投影
功能。也就是说，U(x) 由下式给出

\begin{sourceblock}
\includegraphics[page=595,pagebox=mediabox,trim={64.39bp 128.98bp 64.39bp 1786.77bp},clip,width=0.94\textwidth,max height=0.38\textheight]{Pope.pdf}
\end{sourceblock}
其中基函数系数 a, 由条件确定
重刑他们最小化均方残差

\begin{sourceblock}
\includegraphics[page=596,pagebox=mediabox,trim={64.39bp 1663.79bp 64.39bp 264.97bp},clip,width=0.94\textwidth,max height=0.38\textheight]{Pope.pdf}
\end{sourceblock}
通过代入方程。 （13.28）代入等式。 (13.29) 并微分于
关于 a,, 表明系数满足矩阵方程

\begin{sourceblock}
\includegraphics[page=596,pagebox=mediabox,trim={64.39bp 1406.82bp 64.39bp 475.94bp},clip,width=0.94\textwidth,max height=0.38\textheight]{Pope.pdf}
\end{sourceblock}
\begin{sourceblock}
\includegraphics[page=596,pagebox=mediabox,trim={64.39bp 1119.86bp 64.39bp 775.90bp},clip,width=0.94\textwidth,max height=0.38\textheight]{Pope.pdf}
\end{sourceblock}
证明矩阵 B 是对称正定矩阵。表明
U(x) 可以表示为滤波操作的结果

\begin{sourceblock}
\includegraphics[page=596,pagebox=mediabox,trim={64.39bp 907.89bp 64.39bp 993.87bp},clip,width=0.94\textwidth,max height=0.38\textheight]{Pope.pdf}
\end{sourceblock}
过滤器在哪里

\begin{sourceblock}
\includegraphics[page=596,pagebox=mediabox,trim={64.39bp 722.91bp 64.39bp 1157.85bp},clip,width=0.94\textwidth,max height=0.38\textheight]{Pope.pdf}
\end{sourceblock}
B-'表示矩阵B的逆矩阵的元素。

争论，都来自方程。 (13.33) 并根据 dx/aam,t 方程
过滤后的残差 ti(x) 为零。

\subsection*{13.2.4 滤波后的能谱}\phantomsection\addcontentsline{toc}{subsection}{13.2.4 滤波后的能谱}
为了考察滤波对频谱影响的重要问题，我们现在将U(x)进行统计
齐次并考虑脉动场 u(x) r U(x) - (U)。我们可以将 u(x) 视为沿统计均匀湍流中的一条线的速度 u(x) 的分量，例如，
ul(elx).自协方差为 U(X) =

\begin{sourceblock}
\includegraphics[page=596,pagebox=mediabox,trim={64.39bp 119.39bp 64.39bp 1836.77bp},clip,width=0.94\textwidth,max height=0.38\textheight]{Pope.pdf}
\end{sourceblock}
\begin{sourceblock}
\includegraphics[page=597,pagebox=mediabox,trim={64.39bp 1128.86bp 64.39bp 109.44bp},clip,width=\textwidth,max height=0.72\textheight]{Pope.pdf}
\par\vspace{0.45\baselineskip}
\small 图 13.4。箱式滤波器的衰减系数，虚线；高斯滤波器，实线；锐利的光谱滤光片，点划线。
\end{sourceblock}
并且，与第 6 章中的符号一致，定义频谱 E
I ( K ) 为其普里叶变换的两倍：

\begin{sourceblock}
\includegraphics[page=597,pagebox=mediabox,trim={64.39bp 828.90bp 64.39bp 1071.86bp},clip,width=0.94\textwidth,max height=0.38\textheight]{Pope.pdf}
\end{sourceblock}
由此可见，过滤后的波动在统计上也是平稳的；和，

\textasciitilde{} ( x ) 如练习 13.7 所示，其频谱为

\begin{sourceblock}
\includegraphics[page=597,pagebox=mediabox,trim={64.39bp 654.92bp 64.39bp 1292.84bp},clip,width=0.94\textwidth,max height=0.38\textheight]{Pope.pdf}
\end{sourceblock}
各种滤波器的衰减系数 G (K) \textasciitilde{} 如图 13.4 所示。显然，盒式滤波器在衰减高波数模式方面不是很有效。

为了说明滤波对湍流频谱的影响，图 13.5 显示了雷诺数 Ri 500 处的一维频谱 E ( 从模型频谱 (方程 (6.246)) 获得，相应的滤波频谱 E l l ( )T his 是使用高斯滤波器从方程 (13.37) 获得的，其中 A = f E r = 1 \textasciitilde{} whe \textasciitilde{} re L \textasciitilde{} I 1i s，积分长度标尺

\begin{sourceblock}
\includegraphics[page=597,pagebox=mediabox,trim={64.39bp 188.98bp 64.39bp 1734.78bp},clip,width=0.94\textwidth,max height=0.38\textheight]{Pope.pdf}
\end{sourceblock}
\begin{sourceblock}
\includegraphics[page=598,pagebox=mediabox,trim={64.39bp 1123.86bp 64.39bp 109.44bp},clip,width=\textwidth,max height=0.72\textheight]{Pope.pdf}
\par\vspace{0.45\baselineskip}
\small 图 13.5。为 R, = 500 处的模型光谱获得的一维光谱（实线）；以及使用 A = 1111 的高斯滤波器获得的滤波频谱 E1,(«)（虚线）。
\end{sourceblock}
是总能量l(u2)的92\%。 （请注意，该结果基于一维滤波和一维频谱。在更相关的三维情况下，使用宽度为 A tEIa 的相应三维滤波器，大约 80\% 的能量被解析，请参见练习 13.10。）

表明过滤波动的自协方差为

\begin{sourceblock}
\includegraphics[page=598,pagebox=mediabox,trim={64.39bp 429.95bp 64.39bp 1421.82bp},clip,width=0.94\textwidth,max height=0.38\textheight]{Pope.pdf}
\end{sourceblock}
证明 的谱可以写成

\begin{sourceblock}
\includegraphics[page=598,pagebox=mediabox,trim={64.39bp 119.39bp 64.39bp 1838.77bp},clip,width=0.94\textwidth,max height=0.38\textheight]{Pope.pdf}
\end{sourceblock}
对于固定 y 和 z 显示

\begin{sourceblock}
\includegraphics[page=599,pagebox=mediabox,trim={64.39bp 1687.79bp 64.39bp 219.97bp},clip,width=0.94\textwidth,max height=0.38\textheight]{Pope.pdf}
\end{sourceblock}
从而得到结果

\begin{sourceblock}
\includegraphics[page=599,pagebox=mediabox,trim={64.39bp 1559.80bp 64.39bp 366.95bp},clip,width=0.94\textwidth,max height=0.38\textheight]{Pope.pdf}
\end{sourceblock}
\subsection*{13.2.5 过滤字段的分辨率}\phantomsection\addcontentsline{toc}{subsection}{13.2.5 过滤字段的分辨率}
考虑区间 O < x < C 中的统计齐次场 u(x)，并让滤波场 E(x) 由其在间距为 h N 的均匀网格的 N 个节点处的值离散表示。 需要多少网格间距 h（或等效地需要多少个节点）才能同等地解析 \textasciitilde{} ( x a)d？当然，这是一个重要的问题，因为 LES 计算的成本随着节点数量的增加而增加（至少呈线性）。答案取决于滤波器的选择以及从 \textasciitilde{}(x) 中提取的信息。

为了解决滤波和分辨率问题，我们假设 u(x) 是周期性的（周期为 C），并考虑 u(x) 和 \textasciitilde{} ( x )T 的傅里叶级数。u(x) 的傅里叶级数以 n 离散傅里叶逆变换的形式编写（参见第 693 页的方程（F.6））

\begin{sourceblock}
\includegraphics[page=599,pagebox=mediabox,trim={64.39bp 842.89bp 64.39bp 1104.86bp},clip,width=0.94\textwidth,max height=0.38\textheight]{Pope.pdf}
\end{sourceblock}
其中c是第n个波数对应的傅里叶系数

\begin{sourceblock}
\includegraphics[page=599,pagebox=mediabox,trim={64.39bp 636.92bp 64.39bp 1269.84bp},clip,width=0.94\textwidth,max height=0.38\textheight]{Pope.pdf}
\end{sourceblock}
偶数 Nmaxd 确定最大分辨波数，

\begin{sourceblock}
\includegraphics[page=599,pagebox=mediabox,trim={64.39bp 498.94bp 64.39bp 1418.82bp},clip,width=0.94\textwidth,max height=0.38\textheight]{Pope.pdf}
\end{sourceblock}
选择足够大的值来求解 u(x)。 （如第 9 章所述，各向同性湍流可通过 1.5 充分解决。）

离散傅里叶变换（参见附录 F）提供了这些傅里叶系数与具有 Nmaxn 个节点的均匀网格上的 u(x) 值之间的一对一映射，即 u(nhma,), n = 0, 1,。 - 1，其中

\begin{sourceblock}
\includegraphics[page=599,pagebox=mediabox,trim={64.39bp 144.98bp 64.39bp 1749.78bp},clip,width=0.94\textwidth,max height=0.38\textheight]{Pope.pdf}
\end{sourceblock}
滤波场的傅里叶级数 \textasciitilde{} (xc)a n 可以写成

\begin{sourceblock}
\includegraphics[page=600,pagebox=mediabox,trim={64.39bp 1668.79bp 64.39bp 225.97bp},clip,width=0.94\textwidth,max height=0.38\textheight]{Pope.pdf}
\end{sourceblock}
其中系数是由传递函数衰减的 u(x) 的系数
过滤器的En：

\begin{sourceblock}
\includegraphics[page=600,pagebox=mediabox,trim={64.39bp 1452.82bp 64.39bp 495.94bp},clip,width=0.94\textwidth,max height=0.38\textheight]{Pope.pdf}
\end{sourceblock}
参见等式。 （13.11）。

锐光谱滤波器 我们首先考虑截止波数 < K,,, 的锐光谱滤波器。为简单起见，我们假设这样选择

\begin{sourceblock}
\includegraphics[page=600,pagebox=mediabox,trim={64.39bp 1074.86bp 64.39bp 833.89bp},clip,width=0.94\textwidth,max height=0.38\textheight]{Pope.pdf}
\end{sourceblock}
是偶数。然后是由锐谱产生的系数 T
A+-\% I K I ) 滤波器与 G(K) H(K, 是

\begin{sourceblock}
\includegraphics[page=600,pagebox=mediabox,trim={64.39bp 797.90bp 64.39bp 1053.87bp},clip,width=0.94\textwidth,max height=0.38\textheight]{Pope.pdf}
\end{sourceblock}
因此 E(x) 的傅立叶级数（方程（13.47））可以重新表示为

\begin{sourceblock}
\includegraphics[page=600,pagebox=mediabox,trim={64.39bp 582.93bp 64.39bp 1288.84bp},clip,width=0.94\textwidth,max height=0.38\textheight]{Pope.pdf}
\end{sourceblock}
相应地，在不丢失信息的情况下，离散傅立叶变换允许在物理空间中的间距网格上表示

\begin{sourceblock}
\includegraphics[page=600,pagebox=mediabox,trim={64.39bp 360.95bp 64.39bp 1587.80bp},clip,width=0.94\textwidth,max height=0.38\textheight]{Pope.pdf}
\end{sourceblock}
正是由于这个原因，特征滤波器宽度被定义为A n/u,。

总结一下：锐光谱滤波器给出的滤波场可以用 N K 、 C / \textasciitilde{}Fourier 模式精确表示，或者等效地用 N 个值 a(nh), n 0, 1, ... 表示。 N - 1，在间距 h A n 的网格上： / \textasciitilde{} , .更多
我们现在将注意力转向高斯滤波器，其传递函数 G(K) j.s 对所有参数严格为正。因此，u(x) 和 \textasciitilde{} ( x(E) q . (13.47)) 的傅立叶级数包含相同的信息：给定 \textasciitilde{} ( x )th 的 \textasciitilde{}our i e cro 效率，u(x) 的 ose 可以恢复为

\begin{sourceblock}
\includegraphics[page=601,pagebox=mediabox,trim={64.39bp 1544.81bp 64.39bp 390.95bp},clip,width=0.94\textwidth,max height=0.38\textheight]{Pope.pdf}
\end{sourceblock}
这是一个条件良好的操作。另一方面，如果 \textasciitilde{} ( xi) 在物理空间中由 u(nh,,,) 的 N,,, 值表示，则 n = 0,1,。 .. , Nmax- 1, 那么恢复过程 \textasciitilde{}(nh,,,) （或者条件很差。

在 LES 中，过滤的目标是允许使用比解析 u(x) 所需的更少的模式或节点来充分解析过滤后的场 E(x)。为了检查所涉及的问题，我们考虑修改后的系列

\begin{sourceblock}
\includegraphics[page=601,pagebox=mediabox,trim={64.39bp 1130.86bp 64.39bp 739.91bp},clip,width=0.94\textwidth,max height=0.38\textheight]{Pope.pdf}
\end{sourceblock}
对于一些偶数 N ( N N,,,)，以及在间距为 h C/N 的网格上的物理空间中的相应表示。最高分辨模式的波数为
要考虑的适当系数取决于 E(x) 在波数空间还是物理空间中是否已知。如果波数空间中已知 \textasciitilde{} ( x i)，则可以通过 In1 < L N 且 = 0 获得最佳近似（以确保 O(x) 为实数）。那么频谱
El ( ) fi(x) 与 \textasciitilde{} ( x )E、ll ( u ) 、f 或 < u 相同，并且在较高波数处为零。

分辨率用h/AK,/K,来表征。对于模型频谱，图 13.6 显示了滤波后的速度导数的频谱
K\textasciitilde{}E\textasciitilde{}\textasciitilde{}(K)d\textasciitilde{}(x)/dxR。解析到波数意味着忽略 u > u 的所有贡献。如练习 13.8 所示，分辨率 = h / A 且
K,/K, = 1 导致分别忽略 2\% 和 28\% 对 ((d\textasciitilde{}ildx)') 的贡献。因此 h A 被认为产生良好的分辨率，而 h A 产生较差的分辨率。

另一方面，如果 \textasciitilde{} ( x) 在物理空间的网格上已知，则系数 Z., 可以通过 ii(nh) = E(nh) 的离散傅里叶变换获得，其中 n = 0, I,..., N 1。在这种情况下，\textasciitilde{} ( xa)t 中的傅里叶模式
（ I K I 比网格解析的波数更高 > 别名为

\begin{sourceblock}
\includegraphics[page=602,pagebox=mediabox,trim={64.39bp 1098.86bp 64.39bp 109.44bp},clip,width=\textwidth,max height=0.72\textheight]{Pope.pdf}
\par\vspace{0.45\baselineskip}
\small 图 13.6。 di(x)/dx 的频谱，（虚线）； h/A = 4 和 1 时的混叠谱 x?E1(x)（虚线）；以及 h/A = i 和 1（实线）、模型频谱 (R; = 500) 和高斯滤波器的有限差分近似 (Eq. (13.56)) 的频谱 D,(x) (Eq. (13.57))。 （光谱经过缩放，以便积分统一。）
\end{sourceblock}
Bll(rc) 显示对应于 h/A rc、/rc 和 I 的两个 N 选择的混叠光谱（乘以 rc2）。显然，对于 h/A 1，混叠光谱
El 谱 (rc) 与 E (rc) 的近似值很差。

如果在物理空间中求解 LES 方程，则必须近似空间导数，并且近似的精度取决于分辨率。通过研究简单的有限差分公式的导数 d \textasciitilde{} ( x ) / d x 的近似值，可以获得所涉及的误差的指示

\begin{sourceblock}
\includegraphics[page=602,pagebox=mediabox,trim={64.39bp 471.94bp 64.39bp 1466.82bp},clip,width=0.94\textwidth,max height=0.38\textheight]{Pope.pdf}
\end{sourceblock}
dti(x)/dx 的频谱为 rc2Ell(rc)，而 Dhii(x)i 的频谱为

\begin{sourceblock}
\includegraphics[page=602,pagebox=mediabox,trim={64.39bp 291.96bp 64.39bp 1605.80bp},clip,width=0.94\textwidth,max height=0.38\textheight]{Pope.pdf}
\end{sourceblock}
（参见练习 13.9）。对于模型频谱，图 13.6 显示了 de(x)/dx、rc 2 E l l ( ) 的频谱，以及 h/A 和 h/A 1 的近似值。

fih(rc) 2 可以清楚地看出，近似误差在高值时最大
精度的方差是 Dhfi(x)r 相对于 dti(x)/dx 的方差（其中
并且这些值是光谱积分的比率）。对于 h/A 1，
= Z, 4, 该比率分别为 0.60、0.86 和 0.96。

在具有固定过滤器宽度 A 的湍流的 LES 计算中，计算成本大致与 ( A ) 成正比 - 因为时间步长和每个方向上的节点数按 h 缩放。因此，良好的分辨率需要付出高昂的代价：将分辨率加倍-例如从 h/A = 1 到 h/A i，成本增加 16 倍。

假设 u(x) 具有由下式给出的柯尔莫哥洛夫谱 Ell@)

等式。 (6.240) 第 231 页。表明，对于高斯滤波器，
dti(x)/dx 的频谱为

\begin{sourceblock}
\includegraphics[page=603,pagebox=mediabox,trim={64.39bp 1161.85bp 64.39bp 732.91bp},clip,width=0.94\textwidth,max height=0.38\textheight]{Pope.pdf}
\end{sourceblock}
如果 u(x) 用其高达波数的傅里叶系数表示
K,, 表明 [dii(x)/dxI2 解析的分数为

\begin{sourceblock}
\includegraphics[page=603,pagebox=mediabox,trim={64.39bp 889.89bp 64.39bp 1057.87bp},clip,width=0.94\textwidth,max height=0.38\textheight]{Pope.pdf}
\end{sourceblock}
其中 P 是不完全伽玛函数。因此表明，对于
h/A 和 1，该分数分别为 0.98 和 0.72。

证明 B(x Lh) 的傅立叶变换为

\begin{sourceblock}
\includegraphics[page=603,pagebox=mediabox,trim={64.39bp 685.91bp 64.39bp 1250.84bp},clip,width=0.94\textwidth,max height=0.38\textheight]{Pope.pdf}
\end{sourceblock}
因此表明 Dhn(x)( 方程 (13.56)) 的傅立叶变换为

\begin{sourceblock}
\includegraphics[page=603,pagebox=mediabox,trim={64.39bp 405.95bp 64.39bp 1485.81bp},clip,width=0.94\textwidth,max height=0.38\textheight]{Pope.pdf}
\end{sourceblock}
并验证方程。 （13.57）。

\subsection*{13.2.6 三维过滤}\phantomsection\addcontentsline{toc}{subsection}{13.2.6 三维过滤}
对于一般情况，过滤后的 .*elacity - 场 U(x, t) 由等式定义。 (13.1)

不依赖于 x) 或者它可能不均匀。它可能是各向同性的（即，
仅通过 r = |r| 取决于 r，或者它可能是各向异性的。

所有在一维中引入的滤波器都可以用来生成均匀的、各向同性的三维滤波器。箱式滤波器成为半径为 3A 的球体上的体积平均值；高斯滤波器函数是均值为零、协方差为 6,A?/12 的联合法线；尖锐的光谱滤波器消除了所有具有 |x| 的傅里叶模式> x..在一维中获得的大多数结果可以直接用于这些三维滤波器。特别是，对于均匀湍流，滤波后的能谱函数为

\begin{sourceblock}
\includegraphics[page=604,pagebox=mediabox,trim={64.39bp 1342.83bp 64.39bp 598.92bp},clip,width=0.94\textwidth,max height=0.38\textheight]{Pope.pdf}
\end{sourceblock}
这个关系可以用来估计（见练习13.10），在高雷诺数各向同性湍流中，如果滤波器宽度取为A = 1.24（对于锐光谱滤波器）或A = 0.84（对于高斯滤波器），则80％的动能被解析。

对于一些非均匀流，LES方程可以在均匀各向异性矩形网格上求解，在三个坐标方向上网格间距为hy、h2和h。然后，使用各向异性滤波器，滤波器特征宽度为A、A、A；一般与 hy、ha 和 h 成正比；。例如，各向异性高斯滤波器函数是协方差为 12 的联合法线。下面描述的一些模型涉及特征滤波器宽度 A，对于各向异性滤波器，通常将其视为 A = (A,A2A3)!7，如 Deardorff (1970) 所建议的。 Scotti、Meneveau 和 Lilly (1993) 为这一选择提供了一些理论支持，Scotti、Meneveau 和 Fatica (1997) 进一步讨论了这个问题。

从历史上看，在 Leonard (1974) 引入一般过滤形式主义之前，LES 的一种观点（例如 Deardorff (1970)）是表示矩形网格单元上的平均速度 U(x,t)。也就是说，对于以 x 为中心的单元格，U(x, t) 为

\begin{sourceblock}
\includegraphics[page=604,pagebox=mediabox,trim={64.39bp 371.95bp 64.39bp 1507.81bp},clip,width=0.94\textwidth,max height=0.38\textheight]{Pope.pdf}
\end{sourceblock}
对应的过滤器是

\begin{sourceblock}
\includegraphics[page=604,pagebox=mediabox,trim={64.39bp 205.97bp 64.39bp 1677.79bp},clip,width=0.94\textwidth,max height=0.38\textheight]{Pope.pdf}
\end{sourceblock}
其中 A = h，即使网格间距相同（即，
在针对壁面边界流的 LES-NWR 中，靠近壁面的网格间距被认为更细，以便可以解析近壁面结构（参见 Piomelli (1993)）。这会导致滤波器不均匀（如果滤波器宽度与网格间距成正比）。对于非均匀滤波器，滤波和空间微分运算不会交换（参见练习 3），Ghosal 和 Main (1995) 以及 Vasilyev、Lund 和 Moin (1998) 进一步讨论了这个问题。这些作者还解决了边界附近的过滤问题以及过滤操作的离散实现。为了避免这些困难，在槽道流的 LES 中，通常的做法是在壁法线 (y) 方向上使用非均匀网格，但仅在 x-z 平面上进行过滤（参见练习 13-13）。

练习 13.10 考虑高雷诺数均匀各向同性湍流
以及宽度为 A 和特征波数的各向同性滤波器
A 剩余运动的动能为

\begin{sourceblock}
\includegraphics[page=605,pagebox=mediabox,trim={64.39bp 1035.87bp 64.39bp 864.89bp},clip,width=0.94\textwidth,max height=0.38\textheight]{Pope.pdf}
\end{sourceblock}
对于在惯性子区内截止的锐光谱滤波器，
使用柯尔莫哥洛夫谱获得分数的估计

剩余运动的能量

\begin{sourceblock}
\includegraphics[page=605,pagebox=mediabox,trim={64.39bp 792.90bp 64.39bp 1112.86bp},clip,width=0.94\textwidth,max height=0.38\textheight]{Pope.pdf}
\end{sourceblock}
其中长度尺度为 L k 3/2 表明 80\% 的能量为
已解决 ( e ( k ; ) / k 0.1) 如果选择

\begin{sourceblock}
\includegraphics[page=605,pagebox=mediabox,trim={64.39bp 615.92bp 64.39bp 1321.83bp},clip,width=0.94\textwidth,max height=0.38\textheight]{Pope.pdf}
\end{sourceblock}
使用第 6.5 节中的关系 LLl1 和 LI1 0.43L，

表明相应的滤波器宽度为

对高斯滤波器进行同样的分析，得到

\begin{sourceblock}
\includegraphics[page=605,pagebox=mediabox,trim={64.39bp 249.97bp 64.39bp 1655.79bp},clip,width=0.94\textwidth,max height=0.38\textheight]{Pope.pdf}
\end{sourceblock}
其中 I. 由下式给出

\begin{sourceblock}
\includegraphics[page=605,pagebox=mediabox,trim={64.39bp 119.39bp 64.39bp 1828.77bp},clip,width=0.94\textwidth,max height=0.38\textheight]{Pope.pdf}
\end{sourceblock}
并表明 80\% 的能量被解析为 rc，n/A 给出：

\begin{sourceblock}
\includegraphics[page=606,pagebox=mediabox,trim={64.39bp 1677.79bp 64.39bp 213.97bp},clip,width=0.94\textwidth,max height=0.38\textheight]{Pope.pdf}
\end{sourceblock}
考虑通过伪谱方法求解 LES 方程
解析了哪些 1x1 的波数。认为，对于
高斯滤波器和（对应于较差的空间分辨率）

克尔=Kc
解），所需节点数为 (54/38)3 2.9 个
如果使用锐光谱滤波器（带有 rc，选择为
每种情况下分解 80\% 的能量）。表明对应的
如果使用更好的分辨率 2rc (h i d )，则因子为 23。

rc, = = 13 .12 考虑高雷诺数各向同性湍流的 LES us-
伪谱方法和带截止的锐光谱滤波器

波数 rc，在惯性子区内。表明，如果 90\%
能量要被解决，那么一个因素
如果只有 80\%，则需要更多的模式
能源问题亟待解决。 13 .13 用能谱函数 E (rc) 考虑各向同性湍流。

具有截止波数 rc 的锐光谱滤波器用于 1-2
平面，但在 3 方向上不进行滤波。也就是说，过滤器
使用 rc: rc: 2 rc; 消灭模式。证明能量谱
滤波场的函数 E(rc) 为

\begin{sourceblock}
\includegraphics[page=606,pagebox=mediabox,trim={64.39bp 668.92bp 64.39bp 1207.85bp},clip,width=0.94\textwidth,max height=0.38\textheight]{Pope.pdf}
\end{sourceblock}
表明衰减因子具有高波数渐近

tote l ( r c , / \textasciitilde{} ) \textasciitilde{} .S 如何，如果在一个方向上执行过滤
只是，对应的渐近线是(rc,/rc)。

滤波后的应变率 l+xn 滤波后的速度场 n ( x , t ) 我们可以形成滤波后的速度

\subsection*{13.2.7 应变过滤率}\phantomsection\addcontentsline{toc}{subsection}{13.2.7 应变过滤率}
\begin{sourceblock}
\includegraphics[page=606,pagebox=mediabox,trim={64.39bp 162.98bp 64.39bp 1736.78bp},clip,width=0.94\textwidth,max height=0.38\textheight]{Pope.pdf}
\end{sourceblock}
在此基础上，应变的特征过滤率定义为

\begin{sourceblock}
\includegraphics[page=607,pagebox=mediabox,trim={64.39bp 1716.78bp 64.39bp 228.97bp},clip,width=0.94\textwidth,max height=0.38\textheight]{Pope.pdf}
\end{sourceblock}
\textasciitilde{} 0 t ohf 这些量在下面描述的 SGS 模型中很突出，检查它们对过滤器类型和过滤器宽度的依赖性是有启发性的。

我们考虑高雷诺数各向同性湍流，因此存在一个惯性子区，其中能谱函数 E(K) 采用柯尔莫哥洛夫形式 CE2 /3rc-5/3（方程（6.239））。 S 的均方由频谱 E(K) 和滤波器传递函数确定：

\begin{sourceblock}
\includegraphics[page=607,pagebox=mediabox,trim={64.39bp 1131.86bp 64.39bp 678.91bp},clip,width=0.94\textwidth,max height=0.38\textheight]{Pope.pdf}
\end{sourceblock}
第二行是各向同性速度场的均方应变率与其频谱 E(rc) 之间的标准关系；最终的表达式来自等式： （13.62）。

rc2 因子将积分向高波数加权，而因子 G ( K ) \textasciitilde{}at 将其衰减到截止波数之外因此，对 ( S ) 的贡献主要来自周围的波数
如果截止波数 rc 在惯性子区内，则式（1）中的积分为：通过用柯尔莫哥洛夫谱代替 E(rc) 可以很好地近似 (13.75)。这导致估计

\begin{sourceblock}
\includegraphics[page=607,pagebox=mediabox,trim={64.39bp 516.93bp 64.39bp 1336.83bp},clip,width=0.94\textwidth,max height=0.38\textheight]{Pope.pdf}
\end{sourceblock}
其中常数 af 定义为

\begin{sourceblock}
\includegraphics[page=607,pagebox=mediabox,trim={64.39bp 331.96bp 64.39bp 1569.80bp},clip,width=0.94\textwidth,max height=0.38\textheight]{Pope.pdf}
\end{sourceblock}
和 取决于滤波器类型，但独立于（af 的值在练习 13.14 中推导。）

这种类型的分析起源于 Lilly (1967)。最重要的推论是均方根值。 S 的尺度为 A（对于惯性系中的 A

\begin{sourceblock}
\includegraphics[page=608,pagebox=mediabox,trim={64.39bp 1645.79bp 64.39bp 201.97bp},clip,width=0.94\textwidth,max height=0.38\textheight]{Pope.pdf}
\end{sourceblock}
对于高斯滤波器来说是

\begin{sourceblock}
\includegraphics[page=608,pagebox=mediabox,trim={64.39bp 1511.81bp 64.39bp 425.95bp},clip,width=0.94\textwidth,max height=0.38\textheight]{Pope.pdf}
\end{sourceblock}
\begin{sourceblock}
\includegraphics[page=608,pagebox=mediabox,trim={64.39bp 1391.82bp 64.39bp 498.94bp},clip,width=0.94\textwidth,max height=0.38\textheight]{Pope.pdf}
\end{sourceblock}
作为高波数谱的模型（参见方程（6.249）和
第 237 页的练习 6.33）。对于这个谱，以及柯西
滤波器（见表 13.2），表明方程。 (13.75) 导致

\begin{sourceblock}
\includegraphics[page=608,pagebox=mediabox,trim={64.39bp 1040.87bp 64.39bp 774.90bp},clip,width=0.94\textwidth,max height=0.38\textheight]{Pope.pdf}
\end{sourceblock}
（对于 C 1.5）。显示柯西滤波器

\begin{sourceblock}
\includegraphics[page=608,pagebox=mediabox,trim={64.39bp 816.90bp 64.39bp 1022.87bp},clip,width=0.94\textwidth,max height=0.38\textheight]{Pope.pdf}
\end{sourceblock}
\begin{sourceblock}
\includegraphics[page=608,pagebox=mediabox,trim={64.39bp 735.91bp 64.39bp 1202.85bp},clip,width=0.94\textwidth,max height=0.38\textheight]{Pope.pdf}
\end{sourceblock}
和 Pao 过滤器

\begin{sourceblock}
\includegraphics[page=608,pagebox=mediabox,trim={64.39bp 567.93bp 64.39bp 1328.83bp},clip,width=0.94\textwidth,max height=0.38\textheight]{Pope.pdf}
\end{sourceblock}
等式。 (13.75) 导致

\begin{sourceblock}
\includegraphics[page=608,pagebox=mediabox,trim={64.39bp 291.96bp 64.39bp 1599.80bp},clip,width=0.94\textwidth,max height=0.38\textheight]{Pope.pdf}
\end{sourceblock}
（对于 C 1.5）。表明，对于 Pao 过滤器，

\begin{sourceblock}
\includegraphics[page=608,pagebox=mediabox,trim={64.39bp 167.98bp 64.39bp 1770.78bp},clip,width=0.94\textwidth,max height=0.38\textheight]{Pope.pdf}
\end{sourceblock}
\section*{13.3 滤波守恒方程}\phantomsection\addcontentsline{toc}{section}{13.3 滤波守恒方程}
通过将滤波操作应用于纳维-斯托克斯方程，获得控制滤波速度场 U(x, t) 的守恒方程。我们考虑空间均匀滤波器，以便过滤和微分
滤波后的连续性方程为

\begin{sourceblock}
\includegraphics[page=609,pagebox=mediabox,trim={64.39bp 1449.82bp 64.39bp 444.94bp},clip,width=0.94\textwidth,max height=0.38\textheight]{Pope.pdf}
\end{sourceblock}
从中我们得到

\begin{sourceblock}
\includegraphics[page=609,pagebox=mediabox,trim={64.39bp 1287.84bp 64.39bp 610.92bp},clip,width=0.94\textwidth,max height=0.38\textheight]{Pope.pdf}
\end{sourceblock}
因此，滤波磁场 U 和剩余磁场 u' 都是螺线管磁场。

\subsection*{13.3.1 动量守恒}\phantomsection\addcontentsline{toc}{subsection}{13.3.1 动量守恒}
滤波后的动量方程（以保守形式编写）很简单

\begin{sourceblock}
\includegraphics[page=609,pagebox=mediabox,trim={64.39bp 933.88bp 64.39bp 963.88bp},clip,width=0.94\textwidth,max height=0.38\textheight]{Pope.pdf}
\end{sourceblock}
其中 P(x, t) 是滤波后的压力场。该方程与 Navier-S tokes 方程不同，因为滤波后的乘积 U;U j 不同于
giVB。滤波速度的乘积 差值是残余应力张量，定义为

\begin{sourceblock}
\includegraphics[page=609,pagebox=mediabox,trim={64.39bp 678.91bp 64.39bp 1257.84bp},clip,width=0.94\textwidth,max height=0.38\textheight]{Pope.pdf}
\end{sourceblock}
类似于雷诺应力张量

\begin{sourceblock}
\includegraphics[page=609,pagebox=mediabox,trim={64.39bp 541.93bp 64.39bp 1399.82bp},clip,width=0.94\textwidth,max height=0.38\textheight]{Pope.pdf}
\end{sourceblock}
（严格来说，应力张量是 -PTR 和 -p(uiuj)+）残余动力学

\begin{sourceblock}
\includegraphics[page=609,pagebox=mediabox,trim={64.39bp 366.95bp 64.39bp 1514.81bp},clip,width=0.94\textwidth,max height=0.38\textheight]{Pope.pdf}
\end{sourceblock}
各向异性残余应力张量定义为

\begin{sourceblock}
\includegraphics[page=609,pagebox=mediabox,trim={64.39bp 229.97bp 64.39bp 1707.78bp},clip,width=0.94\textwidth,max height=0.38\textheight]{Pope.pdf}
\end{sourceblock}
各向同性残余应力包含在修正的过滤压力中

\begin{sourceblock}
\includegraphics[page=609,pagebox=mediabox,trim={64.39bp 119.39bp 64.39bp 1839.77bp},clip,width=0.94\textwidth,max height=0.38\textheight]{Pope.pdf}
\end{sourceblock}
通过这些定义，滤波后的动量方程 ( \textasciitilde{} q(1. 3 .9 1)) 可以是

\begin{sourceblock}
\includegraphics[page=610,pagebox=mediabox,trim={64.39bp 1658.79bp 64.39bp 201.97bp},clip,width=0.94\textwidth,max height=0.38\textheight]{Pope.pdf}
\end{sourceblock}
其中基于滤波速度的实质性导数是

\begin{sourceblock}
\includegraphics[page=610,pagebox=mediabox,trim={64.39bp 1500.81bp 64.39bp 404.95bp},clip,width=0.94\textwidth,max height=0.38\textheight]{Pope.pdf}
\end{sourceblock}
像往常一样，动量方程的散度产生修正压力的泊松方程
与 (U) 的雷诺方程一样，U 的滤波方程（方程 (13.87) 和 (13.95)）是未封闭的。通过对残余（或 SGS）应力张量 z! 进行建模来实现闭合。这是第 13.4 节和
然而，在几个重要方面，滤波方程与雷诺方程有很大不同。涉及的字段 - U(x, t)、p(x, t) 和 zf。 (x, t) - 是随机的、三维的和不稳定的，即使流量在统计上是平稳的或均匀的。此外，要建模的应力张量取决于滤波器的类型和宽度的规格。 （事实上，只有对于正滤波器，G 2 0，才能保证k和zR是
IJ 正半定，参见练习 13 .2 1。）

zf .(x, t) 由残余应力模型给出，方程： (13 .87) 和
可以求解 IJ (13.95) 以确定 U(x, t ) 和 P(x, t)。滤波后的速度场取决于滤波器的类型和滤波器宽度 A，但这些量并不直接出现在方程中，它们只是通过 zT.(x, t) 模型间接出现。

练习 - 13.17 考虑高雷诺数均匀湍流
惯性子区内的锐光谱滤波器。使用柯尔莫哥洛夫
频谱以获得平均剩余动能的估计

\begin{sourceblock}
\includegraphics[page=610,pagebox=mediabox,trim={64.39bp 347.96bp 64.39bp 1534.81bp},clip,width=0.94\textwidth,max height=0.38\textheight]{Pope.pdf}
\end{sourceblock}
13.3.2 LES 中的闭合问题是由非线性对流项引起的

\subsection*{13.3.2 残余应力分解}\phantomsection\addcontentsline{toc}{subsection}{13.3.2 残余应力分解}
可以考虑其他分解。伦纳德 (19 74) 引入了一种去-
.zR 分解为三个分量应力，如 Exer- 中所述
ij 西塞 13.18。然而，由于其中两个分量应力不是伽利略不变的（Speziale (19 85)，练习 13.19），所以首选的伽利略不变分解是由 Germano (1986) 提出的：

\begin{sourceblock}
\includegraphics[page=611,pagebox=mediabox,trim={64.39bp 1417.82bp 64.39bp 518.93bp},clip,width=0.94\textwidth,max height=0.38\textheight]{Pope.pdf}
\end{sourceblock}
他们在哪里

伦纳德强调

\begin{sourceblock}
\includegraphics[page=611,pagebox=mediabox,trim={64.39bp 1279.84bp 64.39bp 652.92bp},clip,width=0.94\textwidth,max height=0.38\textheight]{Pope.pdf}
\end{sourceblock}
是

交叉应力

\begin{sourceblock}
\includegraphics[page=611,pagebox=mediabox,trim={64.39bp 1137.86bp 64.39bp 792.90bp},clip,width=0.94\textwidth,max height=0.38\textheight]{Pope.pdf}
\end{sourceblock}
SGS 是
雷诺强调

\begin{sourceblock}
\includegraphics[page=611,pagebox=mediabox,trim={64.39bp 996.87bp 64.39bp 943.88bp},clip,width=0.94\textwidth,max height=0.38\textheight]{Pope.pdf}
\end{sourceblock}
这些应力的重要性在第 13.5.2 节中讨论。

使用尖锐的光谱滤波器，更好的分解是

\begin{sourceblock}
\includegraphics[page=611,pagebox=mediabox,trim={64.39bp 811.90bp 64.39bp 1124.86bp},clip,width=0.94\textwidth,max height=0.38\textheight]{Pope.pdf}
\end{sourceblock}
\begin{sourceblock}
\includegraphics[page=611,pagebox=mediabox,trim={64.39bp 623.92bp 64.39bp 1265.84bp},clip,width=0.94\textwidth,max height=0.38\textheight]{Pope.pdf}
\end{sourceblock}
这是优选的，因为每个分量应力都是经过滤波的量，因此可以用解析模式（即波数为 1x1 rc 的傅立叶模式）精确表示。这种分解对于其他投影（例如最小二乘投影（练习 13.6））也是更可取的（出于同样的原因）。

从分解 U U d 可以看出（如最初亲
由 Leonard (1974) 提出）残余应力张量可以分解为

\begin{sourceblock}
\includegraphics[page=611,pagebox=mediabox,trim={64.39bp 119.39bp 64.39bp 1787.77bp},clip,width=0.94\textwidth,max height=0.38\textheight]{Pope.pdf}
\end{sourceblock}
他们在哪里

是

交叉应力

\begin{sourceblock}
\includegraphics[page=612,pagebox=mediabox,trim={64.39bp 1603.80bp 64.39bp 332.96bp},clip,width=0.94\textwidth,max height=0.38\textheight]{Pope.pdf}
\end{sourceblock}
SGS 是
雷诺强调

\begin{sourceblock}
\includegraphics[page=612,pagebox=mediabox,trim={64.39bp 1476.81bp 64.39bp 465.94bp},clip,width=0.94\textwidth,max height=0.38\textheight]{Pope.pdf}
\end{sourceblock}
（请注意，虽然使用相同的名称，但这三个强调
与等式定义的不同。 (13.100)-(13.102).) 证明
并因伦纳德应力而异，即
杰杰伊杰

\begin{sourceblock}
\includegraphics[page=612,pagebox=mediabox,trim={64.39bp 1259.84bp 64.39bp 676.91bp},clip,width=0.94\textwidth,max height=0.38\textheight]{Pope.pdf}
\end{sourceblock}
13.19 为了考虑各种应力的伽利略不变性
在原始的伦纳德分解中，方程。 （13.105），考虑
变换后的速度场

\begin{sourceblock}
\includegraphics[page=612,pagebox=mediabox,trim={64.39bp 1050.87bp 64.39bp 896.89bp},clip,width=0.94\textwidth,max height=0.38\textheight]{Pope.pdf}
\end{sourceblock}
其中 V 是恒定速度差。获取结果

\begin{sourceblock}
\includegraphics[page=612,pagebox=mediabox,trim={64.39bp 913.88bp 64.39bp 1010.87bp},clip,width=0.94\textwidth,max height=0.38\textheight]{Pope.pdf}
\end{sourceblock}
表明残余应力是伽利略不变量；还有

\begin{sourceblock}
\includegraphics[page=612,pagebox=mediabox,trim={64.39bp 778.90bp 64.39bp 1158.85bp},clip,width=0.94\textwidth,max height=0.38\textheight]{Pope.pdf}
\end{sourceblock}
表明（一般而言）伦纳德应力不是伽利略应力-
不变的。对于方程中的其他应力。 （13.105），证明 Rij 是
伽利略不变量，而（一般来说）Cij 不是。 13.20 验证 Germano 分解 Eqw( 13 .99) 的有效性。显示
Leonard 压力、交叉压力和 SGS Reynolds 压力
由方程式定义。 ( 13 .100)-( 13.102) 是伽利略不变量。表明，如果
滤波器是一个投影，然后是 Germano 的分解（方程（13 .99）-
(13 .102)) 与 Leonard 的相同（方程 (13 .105)-(13.108)）。

对于一般滤波器G(r，满足归一化条件
等式。 (13 .2)，(Pope 1990) 的定义为
抖动密度函数

\begin{sourceblock}
\includegraphics[page=612,pagebox=mediabox,trim={64.39bp 230.97bp 64.39bp 1674.79bp},clip,width=0.94\textwidth,max height=0.38\textheight]{Pope.pdf}
\end{sourceblock}
获取结果

\begin{sourceblock}
\includegraphics[page=612,pagebox=mediabox,trim={64.39bp 119.39bp 64.39bp 1803.77bp},clip,width=0.94\textwidth,max height=0.38\textheight]{Pope.pdf}
\end{sourceblock}
\begin{sourceblock}
\includegraphics[page=613,pagebox=mediabox,trim={64.39bp 1748.78bp 64.39bp 149.98bp},clip,width=0.94\textwidth,max height=0.38\textheight]{Pope.pdf}
\end{sourceblock}
\begin{sourceblock}
\includegraphics[page=613,pagebox=mediabox,trim={64.39bp 1668.79bp 64.39bp 278.96bp},clip,width=0.94\textwidth,max height=0.38\textheight]{Pope.pdf}
\end{sourceblock}
其中积分针对所有 V；和 U、Ui Uj 和 进行评估
在 x，t。证明 f 满足归一化条件 Eq. (12.1)

并且，如果过滤器到处都是非负的，那么 f 也是非-
负，因此具有联合 PDF 的属性。认为，
对于这样的正滤波器，残余应力是半正定的
（高和奥布莱恩 1993）。使用类似的推理来证明
R”（方程（13.100）和（13.102））对于
积极的过滤器。

13.3.3 一个重要的问题是经过滤波的速度场和剩余运动之间的动能传递。滤波后的动能E(x,t)是对动能场E(x,t)-=UU进行滤波得到的，即

\subsection*{13.3.3 能量守恒}\phantomsection\addcontentsline{toc}{subsection}{13.3.3 能量守恒}
\begin{sourceblock}
\includegraphics[page=613,pagebox=mediabox,trim={64.39bp 922.88bp 64.39bp 1024.87bp},clip,width=0.94\textwidth,max height=0.38\textheight]{Pope.pdf}
\end{sourceblock}
这可以分解为

\begin{sourceblock}
\includegraphics[page=613,pagebox=mediabox,trim={64.39bp 793.90bp 64.39bp 1149.85bp},clip,width=0.94\textwidth,max height=0.38\textheight]{Pope.pdf}
\end{sourceblock}
\begin{sourceblock}
\includegraphics[page=613,pagebox=mediabox,trim={64.39bp 668.92bp 64.39bp 1217.85bp},clip,width=0.94\textwidth,max height=0.38\textheight]{Pope.pdf}
\end{sourceblock}
是滤波速度场的动能，k 是剩余动能

\begin{sourceblock}
\includegraphics[page=613,pagebox=mediabox,trim={64.39bp 505.94bp 64.39bp 1388.82bp},clip,width=0.94\textwidth,max height=0.38\textheight]{Pope.pdf}
\end{sourceblock}
Ef 的守恒方程是通过乘以方程得到的。 (13.95) by Uj- 结果可以写成
其中 cf 和 P 定义为

\begin{sourceblock}
\includegraphics[page=613,pagebox=mediabox,trim={64.39bp 168.98bp 64.39bp 1772.78bp},clip,width=0.94\textwidth,max height=0.38\textheight]{Pope.pdf}
\end{sourceblock}
\begin{sourceblock}
\includegraphics[page=613,pagebox=mediabox,trim={64.39bp 119.39bp 64.39bp 1846.77bp},clip,width=0.94\textwidth,max height=0.38\textheight]{Pope.pdf}
\end{sourceblock}
等式左侧的项。 (13.121) 代表运输；但最有趣的是右侧的源项或汇项。水槽

表示直接来自滤波速度场的粘性耗散。对于高雷诺数流，滤波器宽度远大于柯尔莫哥洛夫尺度，此项相对较小（参见练习 13.23）。

方程中的最后一项。 ( 11 2 1 , Pr 是残余动能的产生速率。术语 e r n 在 Ef 方程中显示为汇 ( P )，在 k, 方程中显示为源 (+Pr)。因此，它表示能量从过滤运动到残余运动的传递速率。有时 pr 称为 SGS 耗散并表示为 该术语是不合适的，因为与真正的耗散不同，Pr 完全归因于无粘性，惯性过程，并且它可以是负的。

在高雷诺数下，滤波器处于惯性子区内，滤波后的速度场几乎占据了所有的动能，即

\begin{sourceblock}
\includegraphics[page=614,pagebox=mediabox,trim={64.39bp 1160.85bp 64.39bp 781.90bp},clip,width=0.94\textwidth,max height=0.38\textheight]{Pope.pdf}
\end{sourceblock}
(E) 方程中的主要汇是 ( )，而 (E) 方程（第 124 页上的方程 (5.125)）是动能耗散率，因此，在所考虑的情况下，这两个量几乎相等：

\begin{sourceblock}
\includegraphics[page=614,pagebox=mediabox,trim={64.39bp 902.89bp 64.39bp 1041.87bp},clip,width=0.94\textwidth,max height=0.38\textheight]{Pope.pdf}
\end{sourceblock}
该结果的等效观点（源自 Lilly (1967)）是，在平均剩余动能 ( k ) 的方程中，产生 (Fr) 和耗散之间存在密切平衡
虽然平均而言，能量是从大尺度转移的，但局部可能存在反向散射，即将能量从残余运动转移到滤波后的速度场 Pr 0。这将在 13.5 节中讨论。

练习 13.22 对于各向同性湍流，验证是否给出 (Ef)、( ) 和 d (cf)

通过频谱 E (rc) 的所有 rc 的积分乘以 1，G ( K ) \textasciitilde{}，
- K \textasciitilde{} G ( K ) \textasciitilde{} , 13.23 考虑具有过滤器宽度的高雷诺数流的 LES
A 小于 使用练习 13.15 的结果来估计

\begin{sourceblock}
\includegraphics[page=614,pagebox=mediabox,trim={64.39bp 151.98bp 64.39bp 1748.78bp},clip,width=0.94\textwidth,max height=0.38\textheight]{Pope.pdf}
\end{sourceblock}
\section*{13.4 Smagorinsky 模型}\phantomsection\addcontentsline{toc}{section}{13.4 Smagorinsky 模型}
13.4 Smagorinsky 模型 为了闭合滤波速度方程，需要建立一个模型
需要残余应力张量。最简单的模型是 Smagorinsky (1963) 提出的模型，它也构成了第 13.6 节中描述的几个更高级模型的基础。

\subsection*{13.4.1 模型的定义}\phantomsection\addcontentsline{toc}{subsection}{13.4.1 模型的定义}
该模型可以分为两部分。一、线性涡粘模型

\begin{sourceblock}
\includegraphics[page=615,pagebox=mediabox,trim={64.39bp 1334.83bp 64.39bp 608.92bp},clip,width=0.94\textwidth,max height=0.38\textheight]{Pope.pdf}
\end{sourceblock}
用于将残余应力重新调整为过滤后的应变率。比例系数 vr(x, t) 是剩余运动的涡粘性。其次，通过类比混合长度假设（方程（10.20）），涡黏度建模为

\begin{sourceblock}
\includegraphics[page=615,pagebox=mediabox,trim={64.39bp 998.87bp 64.39bp 890.89bp},clip,width=0.94\textwidth,max height=0.38\textheight]{Pope.pdf}
\end{sourceblock}
其中 S 是特征过滤应变率（方程 (13.74)），lsi 是 Smagorinsky 长度尺度（类似于混合长度），通过 Smagorinsky 系数 Cs，它与过滤器宽度 A m 成正比
根据涡粘模型（方程（13 .127）），能量传递到残余运动的速率为

\begin{sourceblock}
\includegraphics[page=615,pagebox=mediabox,trim={64.39bp 662.92bp 64.39bp 1277.84bp},clip,width=0.94\textwidth,max height=0.38\textheight]{Pope.pdf}
\end{sourceblock}
对于 Smagorinsky 模型（或任何其他具有 vr O 的涡粘模型），这种能量传递无处不在，从滤波运动到残余运动：不存在反向散射。

\subsection*{13.4.2 惯性子区内的行为}\phantomsection\addcontentsline{toc}{subsection}{13.4.2 惯性子区内的行为}
研究应用于高雷诺数湍流的 Smagorinsky 模型非常有用，滤波器宽度 A 在惯性子区内 (ire.,
> A > ) 对于这种情况，平均而言，剩余运动 (Pr) 的能量传递通过耗散来平衡

\begin{sourceblock}
\includegraphics[page=615,pagebox=mediabox,trim={64.39bp 119.39bp 64.39bp 1836.77bp},clip,width=0.94\textwidth,max height=0.38\textheight]{Pope.pdf}
\end{sourceblock}
\begin{sourceblock}
\includegraphics[page=616,pagebox=mediabox,trim={64.39bp 1033.87bp 64.39bp 109.44bp},clip,width=\textwidth,max height=0.72\textheight]{Pope.pdf}
\par\vspace{0.45\baselineskip}
\small 图 13.7。高雷诺数湍流惯性子区内锐谱滤波器的归一化滤波量和剩余量的估计值作为滤波器宽度 A 的函数。表 13,3 给出了标准化和估计值。
\end{sourceblock}
使用从 Kolmogorov 谱获得的 ( S ) 估计（方程 (13.76)），可以求解该方程以获得 Smagorinsky 长度：

\begin{sourceblock}
\includegraphics[page=616,pagebox=mediabox,trim={64.39bp 664.92bp 64.39bp 1216.85bp},clip,width=0.94\textwidth,max height=0.38\textheight]{Pope.pdf}
\end{sourceblock}
其中 C 是柯尔莫哥洛夫常数，af 是由方程定义的滤波器相关常数。 （13.77）。该分析最初由 Lilly (1967) 提出，他使用
-3 = -2 3/2 锐光谱滤波器（参见方程 (13 .78)）和近似 ( S ) ( S ) 以获得结果

\begin{sourceblock}
\includegraphics[page=616,pagebox=mediabox,trim={64.39bp 363.95bp 64.39bp 1560.80bp},clip,width=0.94\textwidth,max height=0.38\textheight]{Pope.pdf}
\end{sourceblock}
因此，Smagorinsky 模型与 A 成比例的假设对于 - 惯性子区中的 A 来说得到了证实。

根据 ls A 和 S E ' / \textasciitilde{} A 的标度 - 其他量的 s\textasciitilde{}c,alings

\begin{center}\small 表 13.3。高雷诺数湍流惯性子区内锐光谱滤波器的滤波量和残余量的估计。数量由 k 标准化，L = k*/? /¢, 后三个量的估计基于 Smagorinsky 模型\end{center}
\clearpage
\begin{sourceblock}
\includegraphics[page=617,pagebox=mediabox,trim={64.39bp 558.93bp 64.39bp 349.96bp},clip,width=\textwidth,max height=0.72\textheight]{Pope.pdf}
\end{sourceblock}
\begin{sourceblock}
\includegraphics[page=617,pagebox=mediabox,trim={64.39bp 704.91bp 64.39bp 623.92bp},clip,width=\textwidth,max height=0.72\textheight]{Pope.pdf}
\par\vspace{0.45\baselineskip}
\small 图 13.7。
\end{sourceblock}
假设 是对数正态分布 var[ln(S /(S
获取结果

\begin{sourceblock}
\includegraphics[page=617,pagebox=mediabox,trim={64.39bp 363.95bp 64.39bp 1542.81bp},clip,width=0.94\textwidth,max height=0.38\textheight]{Pope.pdf}
\end{sourceblock}
表明，如果该估计与 o2 1 一起使用，则 Smagorinsky
Lilly 分析给出的常数减少了约 20\%。讨论
o2 对 A 的依赖性，因此弱的可能性
由于内部间歇性，Cs 的依赖性（诺维科夫
13.25 表明基于Pao滤波器的Lilly分析（见表13.2）

\begin{sourceblock}
\includegraphics[page=618,pagebox=mediabox,trim={64.39bp 1690.79bp 64.39bp 199.97bp},clip,width=0.94\textwidth,max height=0.38\textheight]{Pope.pdf}
\end{sourceblock}
\subsection*{13.4.3 斯马戈林斯基滤波器}\phantomsection\addcontentsline{toc}{subsection}{13.4.3 斯马戈林斯基滤波器}
通过指定 Smagorinsky 长度尺度 és，可以求解使用 Smagorinsky 模型的 LES 方程以确定滤波后的速度场 U(x,t)。这可以在没有明确指定滤波器的情况下完成：因为滤波器既不会出现在 LES 方程（方程（13.87）和（13.95））中，也不会出现在残余粘度指定为 v 的 Smagorinsky 模型中，=
对于均匀各向同性湍流，有一个独特的隐含滤波器-Smagorinsky 滤波器-计算出的光谱是一致的（Pope 1998a）。假设各向同性湍流是使用 DNS 计算的，它产生（在特定时间）能量谱函数 E(x)；使用 LES（使用 Smagorinsky 模型和指定的 £s）计算相同的流量，从而产生滤波场 E(x) 的能谱。因此，原则上（并且实际上在中等雷诺数下）可以计算光谱 E(x) 和 E(x)。根据方程式，这些光谱通过滤波器传递函数 G(x) 相关。 （13.62），可以反转得到

\begin{sourceblock}
\includegraphics[page=618,pagebox=mediabox,trim={64.39bp 751.90bp 64.39bp 1159.85bp},clip,width=0.94\textwidth,max height=0.38\textheight]{Pope.pdf}
\end{sourceblock}
这就是 Smagorinsky 滤波器的传递函数：只有使用该滤波器，滤波后的 DNS 速度场才具有与从 LES 获得的滤波后速度场相同的频谱。 （虽然使用此滤波器，Smagorinsky 模型可以生成正确的频谱，但并不意味着它相对于其他统计数据是准确的。）

对于高雷诺数各向同性湍流和小于积分长度尺度的 és，可以粗略估计 Smagorinsky 滤波器
<1）。这是通过估计 E(x) 和 E(x) 来完成的。对于 E(x)，我们假设耗散模型谱具有 Pao 形式（方程 (6.254)）

\begin{sourceblock}
\includegraphics[page=618,pagebox=mediabox,trim={64.39bp 210.97bp 64.39bp 1669.79bp},clip,width=0.94\textwidth,max height=0.38\textheight]{Pope.pdf}
\end{sourceblock}
（使用 Pao 的形式，而不是方程（6.248）来简化代数。）

由 (SL) 则残余粘度

\begin{sourceblock}
\includegraphics[page=619,pagebox=mediabox,trim={64.39bp 1716.78bp 64.39bp 229.97bp},clip,width=0.94\textwidth,max height=0.38\textheight]{Pope.pdf}
\end{sourceblock}
是非随机且均匀的，LES 动量方程 (Eq, (13.95)) 变为
这与 Navier-Stokes 方程相同，只是有效粘度为 v v，即方程的 LES 解。 (13.138) 与
r\textasciitilde{} SO 雷诺数低 (1 v,/v) 时纳维-斯托克斯方程的 DNS 解。有效的柯尔莫哥洛夫量表定义为

\begin{sourceblock}
\includegraphics[page=619,pagebox=mediabox,trim={64.39bp 1217.85bp 64.39bp 670.92bp},clip,width=0.94\textwidth,max height=0.38\textheight]{Pope.pdf}
\end{sourceblock}
（参见练习 13.26）。只要与积分尺度相比较小，则与式（1）一致。 (13.136) 假设 LES 产生

\begin{sourceblock}
\includegraphics[page=619,pagebox=mediabox,trim={64.39bp 1020.87bp 64.39bp 876.89bp},clip,width=0.94\textwidth,max height=0.38\textheight]{Pope.pdf}
\end{sourceblock}
根据 E(Ic) 和 E(Ic) 的这些估计，可以从等式推导出 Smagorinsky 滤波器的传递函数。 (13.135) 是

\begin{sourceblock}
\includegraphics[page=619,pagebox=mediabox,trim={64.39bp 838.89bp 64.39bp 1087.86bp},clip,width=0.94\textwidth,max height=0.38\textheight]{Pope.pdf}
\end{sourceblock}
这对应于表 13.2 中定义的 Pao 过滤器，过滤器宽度为 A es/Cs，其中

\begin{sourceblock}
\includegraphics[page=619,pagebox=mediabox,trim={64.39bp 448.94bp 64.39bp 1370.83bp},clip,width=0.94\textwidth,max height=0.38\textheight]{Pope.pdf}
\end{sourceblock}
重要的一点是，有一个由等式给出的独特滤波器。 (13.135)。从前面的参数推导出来的特定滤波器是一个近似值（因为忽略了波动），并且它取决于耗散谱假设的形式。如图所示，Pao 谱通向 Pao 滤波器，而指数谱 exp(-poicq) 通向 Cauchy 滤波器，模型谱通向滤波器 Eq. (13.150) 由练习 13.27 推导出来。这些滤波器传递函数如图 13.8 所示，并且

\begin{sourceblock}
\includegraphics[page=620,pagebox=mediabox,trim={64.39bp 1116.86bp 64.39bp 109.44bp},clip,width=\textwidth,max height=0.72\textheight]{Pope.pdf}
\par\vspace{0.45\baselineskip}
\small 图 13.8。滤波器传递函数：高斯，实线；柯西虚线； Pao，点划线；由模型谱暗示（方程（13.150）），虚线，
\end{sourceblock}
\begin{sourceblock}
\includegraphics[page=620,pagebox=mediabox,trim={64.39bp 119.39bp 64.39bp 1815.77bp},clip,width=0.94\textwidth,max height=0.38\textheight]{Pope.pdf}
\end{sourceblock}
\begin{sourceblock}
\includegraphics[page=621,pagebox=mediabox,trim={64.39bp 848.89bp 64.39bp 1086.86bp},clip,width=0.94\textwidth,max height=0.38\textheight]{Pope.pdf}
\end{sourceblock}
\begin{sourceblock}
\includegraphics[page=621,pagebox=mediabox,trim={64.39bp 699.91bp 64.39bp 1247.84bp},clip,width=0.94\textwidth,max height=0.38\textheight]{Pope.pdf}
\end{sourceblock}
\begin{sourceblock}
\includegraphics[page=621,pagebox=mediabox,trim={64.39bp 551.93bp 64.39bp 1395.82bp},clip,width=0.94\textwidth,max height=0.38\textheight]{Pope.pdf}
\end{sourceblock}
\begin{sourceblock}
\includegraphics[page=621,pagebox=mediabox,trim={64.39bp 316.96bp 64.39bp 1518.81bp},clip,width=0.94\textwidth,max height=0.38\textheight]{Pope.pdf}
\end{sourceblock}
\subsection*{13.4.4 限制行为}\phantomsection\addcontentsline{toc}{subsection}{13.4.4 限制行为}
从方程式(13 .145)，并由 v, = \&\$ (S ) 得出结果

\begin{sourceblock}
\includegraphics[page=622,pagebox=mediabox,trim={64.39bp 1674.79bp 64.39bp 216.97bp},clip,width=0.94\textwidth,max height=0.38\textheight]{Pope.pdf}
\end{sourceblock}
13-27 使用以下公式考虑均匀各向同性湍流的 LES
Smagorinsky 模型的滤波器位于惯性子区内，因此
Ys 比柯尔莫哥洛夫尺度 (ls q) 大得多。显示
如果 E(K)i 由模型谱近似（方程（6.246）

和 (6.248))，则隐含的 Srnagorinsky 滤波器具有传递

\begin{sourceblock}
\includegraphics[page=622,pagebox=mediabox,trim={64.39bp 1322.83bp 64.39bp 611.92bp},clip,width=0.94\textwidth,max height=0.38\textheight]{Pope.pdf}
\end{sourceblock}
如果滤波器被缩放使得 G(K,) exp(-n2/24) （就像对于
高斯、柯西和 Pao 滤波器），表明隐含值
斯马戈林斯基常数是

\begin{sourceblock}
\includegraphics[page=622,pagebox=mediabox,trim={64.39bp 1028.87bp 64.39bp 841.89bp},clip,width=0.94\textwidth,max height=0.38\textheight]{Pope.pdf}
\end{sourceblock}
对于 5.2 和 c，--0.4。

13.4.4 限制行为 LES 的一个理想应用是过滤器处于惯性子区内的高雷诺数湍流。然而，在实践中，偏离这一理想是常态。通过考虑极端偏离可以获得一些见解。具体来说，我们考虑滤波器处于耗散区内的三种情况；滤波器与积分尺度相比较大；和层流。

非常 srnal1Jilter 宽度 如果滤波器处于远耗散区 ( i , A I) 中，则可以执行泰勒级数分析以表明残余应力张量为（至前导阶）

\begin{sourceblock}
\includegraphics[page=622,pagebox=mediabox,trim={64.39bp 216.97bp 64.39bp 1708.78bp},clip,width=0.94\textwidth,max height=0.38\textheight]{Pope.pdf}
\end{sourceblock}
（参见练习 13.28）。该结果适用于具有有限矩的滤波器，例如
模型，该表达式在 A 和滤波速度中是二次的
zR梯度，但张量形式不同。

尽管存在这种差异，但可以确定斯马戈林斯基系数的值，使得能量转移到残余尺度的平均速率
与方程一致。 (13.15'2)（参见练习 13.28 和方程（13.160））。对于 ( p r > 高雷诺数湍流，Cs 0.13 的最终估计值略低于 Cs 0.17 的惯性范围值。据推测，在 A/q = 1 附近，斯马戈林斯基系数在惯性范围和远耗散区的常数值之间存在过渡。

在考虑的极限条件下（A/q，应变过滤率 Sij 趋于 su, 6
因此 ( S ) 趋于 (SijSij) (sijsij) 因此，Smagorinsky 模型得出

\begin{sourceblock}
\includegraphics[page=623,pagebox=mediabox,trim={64.39bp 1184.85bp 64.39bp 705.91bp},clip,width=0.94\textwidth,max height=0.38\textheight]{Pope.pdf}
\end{sourceblock}
在远耗散区内，A 的下降速度比 vr 的下降速度更快

\textasciitilde{} 4 / 3 在惯性子区内。

与此处获得的结果相反，第 13.4.3 节中的分析预测 Cs 随着 A/q 的减小而增加（方程（13.142）和（13.147））。另一方面，Voke (1996) 以及 Meneveau 和 Lund (1997) 进行的 Lilly 型分析（使用锐光谱滤波器）表明，Cs 随着 A 的增加而降低。这些差异部分归因于 Smagorinsky 模型对残余应力的详细描述很差，如下所述。因此，当使用模型行为的不同方面来估计 Cs 时，可以获得不同的结果。从实践的角度来看，Smagorinsky 模型的这种限制行为的不确定性并不是很重要，因为对于较小的 A 来说，粘性应力在残余应力中占主导地位。

考虑各向同性滤波器（例如高斯滤波器），其具有
第二时刻

\begin{sourceblock}
\includegraphics[page=623,pagebox=mediabox,trim={64.39bp 267.97bp 64.39bp 1636.79bp},clip,width=0.94\textwidth,max height=0.38\textheight]{Pope.pdf}
\end{sourceblock}
和有限的更高的\textasciitilde{}\textasciitilde{}\textasciitilde{} o m e nTthsis。应用于湍流，其中
滤波器宽度 A 远小于柯尔莫哥洛夫尺度 q。

U(x r) 可以展开泰勒级数

\begin{sourceblock}
\includegraphics[page=624,pagebox=mediabox,trim={64.39bp 1689.79bp 64.39bp 211.97bp},clip,width=0.94\textwidth,max height=0.38\textheight]{Pope.pdf}
\end{sourceblock}
其中速度梯度在 x 处计算。通过乘法
等式。 (13.155) 通过 G(r) 并积分，表明滤波后的速度
ity 字段是（按前导顺序）

\begin{sourceblock}
\includegraphics[page=624,pagebox=mediabox,trim={64.39bp 1447.82bp 64.39bp 456.94bp},clip,width=0.94\textwidth,max height=0.38\textheight]{Pope.pdf}
\end{sourceblock}
按照类似的过程来证明残余应力张量为

\begin{sourceblock}
\includegraphics[page=624,pagebox=mediabox,trim={64.39bp 1279.84bp 64.39bp 618.92bp},clip,width=0.94\textwidth,max height=0.38\textheight]{Pope.pdf}
\end{sourceblock}
显示剩余运动的能量转移的平均速率

\begin{sourceblock}
\includegraphics[page=624,pagebox=mediabox,trim={64.39bp 1074.86bp 64.39bp 768.90bp},clip,width=0.94\textwidth,max height=0.38\textheight]{Pope.pdf}
\end{sourceblock}
使用第 205 页练习 6.1 1 的结果表明，在高
雷诺数，局部各向同性湍流（对于 A I），

\begin{sourceblock}
\includegraphics[page=624,pagebox=mediabox,trim={64.39bp 862.89bp 64.39bp 1032.87bp},clip,width=0.94\textwidth,max height=0.38\textheight]{Pope.pdf}
\end{sourceblock}
其中 S 是速度导数 s.kewness（为负）。由
将该结果与方程进行比较。 （13.130），表明它是一致的
与 Smagorinsky 模型

\begin{sourceblock}
\includegraphics[page=624,pagebox=mediabox,trim={64.39bp 597.92bp 64.39bp 1290.84bp},clip,width=0.94\textwidth,max height=0.38\textheight]{Pope.pdf}
\end{sourceblock}
取 (S ) （S 和 S -0.7（参见第 262 页图 6.33）

\begin{sourceblock}
\includegraphics[page=624,pagebox=mediabox,trim={64.39bp 426.95bp 64.39bp 1520.81bp},clip,width=0.94\textwidth,max height=0.38\textheight]{Pope.pdf}
\end{sourceblock}
非常大的Blter宽度相反的极端，即滤波器宽度A与湍流积分尺度L相比较大，可以针对均匀湍流的情况进行分析。当 A 趋于无穷大时，滤波后的速度 V(x, t)

波动 u(x, t)。因此，残余应力张量 zR 趋向于雷诺应力张量 ) 对于均匀湍流剪切流，Smagorinsky 模型给出的残余剪切应力为

\begin{sourceblock}
\includegraphics[page=625,pagebox=mediabox,trim={64.39bp 1599.80bp 64.39bp 298.96bp},clip,width=0.94\textwidth,max height=0.38\textheight]{Pope.pdf}
\end{sourceblock}
而混合长度模型给出的雷诺剪切应力为

\begin{sourceblock}
\includegraphics[page=625,pagebox=mediabox,trim={64.39bp 1431.82bp 64.39bp 462.94bp},clip,width=0.94\textwidth,max height=0.38\textheight]{Pope.pdf}
\end{sourceblock}
显然，在所考虑的极限（A/L a）下，残余涡黏度 vr 为湍流黏度，Smagorinsky 长度 ls 为混合
VT，长度 混合长度可计算为

\begin{sourceblock}
\includegraphics[page=625,pagebox=mediabox,trim={64.39bp 1198.85bp 64.39bp 738.91bp},clip,width=0.94\textwidth,max height=0.38\textheight]{Pope.pdf}
\end{sourceblock}
当然，它与 A 无关。因此，当 AIL 趋于无穷大时，Smagorinsky 系数趋于零，如下所示

\begin{sourceblock}
\includegraphics[page=625,pagebox=mediabox,trim={64.39bp 977.88bp 64.39bp 930.88bp},clip,width=0.94\textwidth,max height=0.38\textheight]{Pope.pdf}
\end{sourceblock}
流动 对于层流，雷诺方程恢复为纳维-斯托克斯方程，并且雷诺应力为零。相反，一般来说，层流中的残余应力张量不为零。对于与层流速度场的长度尺度相比较小的滤波器宽度 A，练习 13.28 的泰勒-斯克里分析是有效的，并且它得出方程： (13.157) 为残差
fi 应力张量
尽管有这样的一般结果，但对于几个重要的流动来说，残余应力基本上为零。例如，考虑单向流，其中速度的唯一分量 U1 仅取决于 x2 和 x3。从等式可以看出。 （13.157）唯一的非零残余应力是 zfl，而影响动量传递的剪切应力为零。对于二维和三维边界层流（在边界层近似内）也是如此。

对于残余剪应力为零的层流剪切流，Smagorinsky 系数的适当值为 Cs 0：Cs 的非零值会错误地导致残余剪应力
对于层流，Cs 值不正确。在某些实现中（例如
我的 Sclhumann (1975) 和 Moin 和 Kim (1.982)) 的 Smagorinsky 模型基于 Sij (Sij)，而不是基于层流中为零的 Sij。

SO V，但是请注意，这意味着 (Sij) 在 LES 中不容易获得，但必须通过计算适当类型的平均值来估计。

13.4.5 近壁分辨
正如本章介绍中所讨论的，在 LES 中，有两种不同的方法来治疗近壁区域（参见第 560 页的表 13.1）。在 LES-

\subsection*{13.4.5 近壁分辨}\phantomsection\addcontentsline{toc}{subsection}{13.4.5 近壁分辨}
精细解析各处 80\% 的能量，包括粘性壁区域；而在 LES-NWM（具有近壁建模的 LES）中，近壁运动未得到解决。在本节中，我们将研究这两种方法的计算成本，并考虑在 LES-NWR 中使用 Smagorinsky 模型。

首先，我们回顾第七章，在边界层型流动中，粘性壁区域非常重要：产生、耗散、动能和雷诺应力各向异性都在小于 20 时达到峰值，即在壁的 20 个粘性长度尺度 6 内。相对于流动长度尺度 6，粘性长度尺度 6, 较小，并且它随雷诺数减小，大约为 6,/6 Re。

在 LES-NWR 中，为了解决近壁运动，粘性近壁区域中的滤波器宽度和网格间距必须约为 6。由此可以估计，所需的网格节点数量随着 Re 的增加而增加
（参见 Chapman (1979) 和练习 13.29）。因此，LES-NWR 对于高雷诺数流（例如航空和气象应用中出现的流）不可行。

相反，在 LES-NWM 中，过滤器宽度和网格间距比例为流长度比例 6，因此计算要求与雷诺数无关（参见练习 13.29），但是，重要的近壁过程并未得到解决，而是被建模。！ LES-NWM 中使用的近壁处理的 km \textasciitilde{} eo 在第 13.6.5 节中描述。

在近壁区域（就像其他地方一样），残余应力取决于过滤器的类型和宽度。在槽道流的 LES 中，通常的做法是仅在平面上进行过滤。如果这是使用均匀滤波器（特别是使用独立于 y 的 A）完成的，则可以证明，非常接近
雷诺应力 (u-u*)（参见练习 13.30）。特别是，剪切应力
变化为 y3（按领先顺序）。

在 Smagorinsky 模型中，规范 ls CsA（对于常数 Cs）对于高雷诺数湍流惯性子区中的 A 是合理的。这些情况不属于粘性壁区域，并且这种特性会错误地导致壁上的非零残余粘度和剪切应力（参见练习 13.31）。相反，Moin 和 Kim (1982) 使用 van Driest 阻尼函数将 ls 指定为

\begin{sourceblock}
\includegraphics[page=627,pagebox=mediabox,trim={64.39bp 1421.82bp 64.39bp 522.93bp},clip,width=0.94\textwidth,max height=0.38\textheight]{Pope.pdf}
\end{sourceblock}
本练习的目的是调查计算成本

LES-NWR 和 LES-NWM 用于壁面流动，特别是
确定所需的网格点数量 N 如何缩放
雷诺数。考虑完全发展的湍流的 LES
使用四种不同网格的槽道流量。

(A) 使用结构化网格的传统 LES-NWR

\begin{sourceblock}
\includegraphics[page=627,pagebox=mediabox,trim={64.39bp 919.88bp 64.39bp 1026.87bp},clip,width=0.94\textwidth,max height=0.38\textheight]{Pope.pdf}
\end{sourceblock}
\begin{sourceblock}
\includegraphics[page=627,pagebox=mediabox,trim={64.39bp 833.89bp 64.39bp 1105.86bp},clip,width=0.94\textwidth,max height=0.38\textheight]{Pope.pdf}
\end{sourceblock}
其中a、、a、、a、b和c是正常数，b是c。

请注意，Ax 和 Az 的缩放比例为 6，允许分辨率
近壁结构，但导致超分辨率
墙。

(B) 使用不规则网格的最佳分辨率 LES-NWR（第 1 章）

man 1979）：Ay 再次由方程指定。 (13.168)，但是斧头和
Az 与 Ay 成比例

\begin{sourceblock}
\includegraphics[page=627,pagebox=mediabox,trim={64.39bp 435.94bp 64.39bp 1510.81bp},clip,width=0.94\textwidth,max height=0.38\textheight]{Pope.pdf}
\end{sourceblock}
(C) LES-NWM 具有近壁平均轮廓的分辨率：Ay 是
再次由方程指定。 (13.168) 使得平均值的轮廓
速度、雷诺应力等在粘性中得到解决
近壁区域，但 Ax 和 Az 指定为

\begin{sourceblock}
\includegraphics[page=627,pagebox=mediabox,trim={64.39bp 167.98bp 64.39bp 1777.78bp},clip,width=0.94\textwidth,max height=0.38\textheight]{Pope.pdf}
\end{sourceblock}
小），因为雷诺数趋于无穷大（ 6 O），
近壁结构尚未解决。

(D) 没有近壁平均轮廓分辨率的 LES-NWM：Ax
和 Az 由方程式指定。 (13 .170)，Ay 指定为

\begin{sourceblock}
\includegraphics[page=628,pagebox=mediabox,trim={64.39bp 1597.80bp 64.39bp 340.96bp},clip,width=0.94\textwidth,max height=0.38\textheight]{Pope.pdf}
\end{sourceblock}
对于这四个规格中的每一个，评估近似的数字
网格点误码率为

\begin{sourceblock}
\includegraphics[page=628,pagebox=mediabox,trim={64.39bp 1336.83bp 64.39bp 549.93bp},clip,width=0.94\textwidth,max height=0.38\textheight]{Pope.pdf}
\end{sourceblock}
用 Re, 6/6, 表示结果。使用经验
关系（对于高雷诺数）Re, 0.09 Re 获得
四个网格的缩放比例如下：

\begin{sourceblock}
\includegraphics[page=628,pagebox=mediabox,trim={64.39bp 978.88bp 64.39bp 799.90bp},clip,width=0.94\textwidth,max height=0.38\textheight]{Pope.pdf}
\end{sourceblock}
（请注意，最大的区别在于 B 和 C 之间，这导致
来自不同的近壁 x-z 分辨率，其比例为 6，对于
B，C 为 6。） 13.30 考虑通常坐标系中的湍流槽道流，因此
y = x2 是距墙壁的正常距离。非常接近
wall，对于 y 的前导阶，速度场可以写成

\begin{sourceblock}
\includegraphics[page=628,pagebox=mediabox,trim={64.39bp 635.92bp 64.39bp 1311.83bp},clip,width=0.94\textwidth,max height=0.38\textheight]{Pope.pdf}
\end{sourceblock}
\begin{sourceblock}
\includegraphics[page=628,pagebox=mediabox,trim={64.39bp 561.93bp 64.39bp 1385.83bp},clip,width=0.94\textwidth,max height=0.38\textheight]{Pope.pdf}
\end{sourceblock}
\begin{sourceblock}
\includegraphics[page=628,pagebox=mediabox,trim={64.39bp 488.94bp 64.39bp 1458.82bp},clip,width=0.94\textwidth,max height=0.38\textheight]{Pope.pdf}
\end{sourceblock}
其中 a、b 和 是随机函数，参见。等式。 (7.56)-(7.58)。过滤
仅应用于平面，使用宽度为 A 的均匀滤波器。 显示
滤波后的速度场为

\begin{sourceblock}
\includegraphics[page=628,pagebox=mediabox,trim={64.39bp 275.96bp 64.39bp 1664.79bp},clip,width=0.94\textwidth,max height=0.38\textheight]{Pope.pdf}
\end{sourceblock}
残余应力为

\begin{sourceblock}
\includegraphics[page=628,pagebox=mediabox,trim={64.39bp 119.39bp 64.39bp 1783.78bp},clip,width=0.94\textwidth,max height=0.38\textheight]{Pope.pdf}
\end{sourceblock}
\subsection*{13.4.6 模型性能测试}\phantomsection\addcontentsline{toc}{subsection}{13.4.6 模型性能测试}
确认（对于所考虑的滤波器）残余应力 zR 变化
y 的幂与雷诺应力 (uiuj) 相同。

论证在 LES 中，由于 U 是螺线管并且在
墙，其前导行为由方程式给出。 (13.180)，而
zR 的行为取决于所采用的残余应力模型。

进一步练习 13.30，考虑 Smagorin- 的性能
湍流槽道流 LES 中的天空模型。表明 aVl/dx2
和 d U3/dx2 c 是唯一的前阶速度梯度。

表明建模的残余剪应力为

\begin{sourceblock}
\includegraphics[page=629,pagebox=mediabox,trim={64.39bp 1373.83bp 64.39bp 566.93bp},clip,width=0.94\textwidth,max height=0.38\textheight]{Pope.pdf}
\end{sourceblock}
因此需要变化 ls y 才能产生
正确的行为。相比之下，表明标准规范
ls = CsA 导致 v，并且独立于 y（导致
订单）。通过忽略波动并获得估计

\begin{sourceblock}
\includegraphics[page=629,pagebox=mediabox,trim={64.39bp 1074.86bp 64.39bp 828.90bp},clip,width=0.94\textwidth,max height=0.38\textheight]{Pope.pdf}
\end{sourceblock}
对于 A >> 6、v/(z,/p) 1/2 有效
rtiodel性能测试 一般而言，可以通过两种截然不同的方式来测试模型。在 LES 中，这两种方法被称为先验测试和后验测试。先验测试使用实验或 DNS 数据来直接测量建模假设的准确性，例如，Smagorinsky 模型给出的残余应力张量 zr 的关系（方程（13.127）和（13.128））。在后验测试中，模型用于执行湍流计算，并再次参考实验或 DNS 数据来评估计算统计数据（例如（U）和（uiuj））的准确性。进行先验测试来直接评估所做出的近似值的有效性和准确性是自然且适当的。然而，为了使 LES 方法有用，需要在后验测试中取得成功。

对于均匀湍流，Clark、Ferziger 和 Reynolds (1979) 以及 McMillan 和 Ferziger (1979) 报告了 Smagorinsky 模型的先验测试。对于带有 Riw 38 的 643 DNS，提取滤波后的速度场 U(x, t)，并对残余应力进行 Smagorinsky 预测
是由此确定的。然后将 Thls 与残余应力 Ti i 进行比较
和 Ferziger (1979) 认为 z 和 zij 之间的相关性最大
ir,jS mag r,DNS 当 Smagorinsky 系数取 Cs 0.17 时，与 Lilly 分析一致。然而，相关系数很小，大约
发现标量 1。较高的相关系数，大约

\begin{sourceblock}
\includegraphics[page=630,pagebox=mediabox,trim={64.39bp 1546.81bp 64.39bp 346.96bp},clip,width=0.94\textwidth,max height=0.38\textheight]{Pope.pdf}
\end{sourceblock}
这与能量向剩余运动 pr 的转移有关。

在这些早期工作之后，出现了许多对 Smagorinsky 模型进行测试的研究，其中
先验支持所得出的结论。例如，根据 RA 310 的圆射流实验，Liu、Meneveau 和 Katz (1994) 表明，Smagorinsky 模型应力和测量的残余应力之间的相关系数不大于 13.5.6 节中讨论的，测试的重要性（尤其是相关系数的测量）为
先验并不像看起来那么明显。

在测试中，比较 LES 计算的统计数据
后验与从实验或 DNS 获得的结果。对于统计平稳流的 LES，可以执行长时间平均以获得平均滤波速度场 (U) 和
解决雷诺应力

\begin{sourceblock}
\includegraphics[page=630,pagebox=mediabox,trim={64.39bp 899.89bp 64.39bp 1036.87bp},clip,width=0.94\textwidth,max height=0.38\textheight]{Pope.pdf}
\end{sourceblock}
在测试中，从 LES 到 (U) 和 RS+ 是很常见的
后验与实验或 DNS 中的 (U) 和 (uiuj) 进行比较。然而，应该认识到，这些量并不直接相等（参见练习 13.32），尽管它们之间的差异随着滤波器宽度的减小而减小。 DNS 数据可以被过滤以获得之间的直接比较
qf 过滤的数量，例如
对于非均匀流动，测试的一般结论
后验的问题是 Smagorinsky 模型耗散太大，也就是说，它向残余运动传递了太多能量。在河道流量计算中，Smagorinsky 系数通常会降低，例如降低至 Cs 0.1（Deardorff 1970、Piomelli、Moin 和 Ferziger（1988））或 Cs 0.065（Moin 和 Kim 1982）。由于 Smagorinsky 模型在层流中产生虚假残余应力，因此该模型在过渡 fl \textasciitilde{} w sA 中又过于耗散。详细介绍了 Vreman 对湍流混合层的研究，
后验 Geurts 和 Kuerten (1997) 揭示了类似的缺陷。

重要的是要认识到残余应力模型的性能取决于雷诺 r \textasciitilde{} m 轴承类型的选择和
数字很低。因此，尺度的能量包含范围和耗散区之间存在相当大的重叠，并且过滤器很可能被放置在该重叠区域中：残余粘度通常约为分子粘度的两倍。这种情况与高雷诺数流的理想情况完全不同，过滤器被放置在靠近大惯性子区的开始处。对于 LES 在高雷诺数自由剪切流中的这种理想应用，我们有充分的理由认为 Smagorinsky 模型是令人满意的；所得到的方程类似于较低雷诺数下的纳维-斯托克斯方程，并且已知自由剪切流的单点统计对雷诺数不敏感。

大气边界层的S-NWM，滤波器宽度和网格间距与边界层厚度成比例，因此，在地面上，它们与含能湍流运动相比较大。与 LES-NWR 一样，靠近墙壁时衰减（但方式不同）。 Mason 和 Thomson (1992) 研究了 CsA 的各种混合和混合长度 t,,，但得出的结论是，Smagorinsky 模型本质上无法产生正确的对数速度分布。

在 LES 中，残余应力模型的主要功能是以适当的速率从解析尺度中移除能量，该速率为 Pr = -zf Sij（方程（13.123））。根据 Smagorinsky 模型，-zr+ 和 Sij 几乎完全相关（参见练习 13 .3 3），而先验测试表明，事实上相关性实际上较弱。因此，正如 Jiminez 和 Moser (1998) 所观察到的，如果将 Smagorinsky 系数设置为产生适当的 (P) 值，则建模的残余应力的大小就会太低。靠近墙壁，如果过滤器宽度与含能尺度相比不小，则剪切应力的很大一部分来自残余应力模型。因此，可以理解的是，在这种情况下，Smagorinsky 模型表现不佳，因为 Cs 的选择无法产生 (pr) 和 (zi2) 的正确水平。

平均量 ( Q ( x , t)) （例如，(U) 和 (uiuj)）在长度尺度上变化
L 与含能运动相当。争论说，
对于 A L 来说，这种均值受过滤的影响很小，即

\begin{sourceblock}
\includegraphics[page=631,pagebox=mediabox,trim={64.39bp 243.97bp 64.39bp 1695.79bp},clip,width=0.94\textwidth,max height=0.38\textheight]{Pope.pdf}
\end{sourceblock}
因此认为，在统计平稳流的 LES 中，(U) 可以
准确地近似为 U 的长期平均值。

\section*{13.5 波数空间中的 LES}\phantomsection\addcontentsline{toc}{section}{13.5 波数空间中的 LES}
\subsection*{13.5.1 滤波方程}\phantomsection\addcontentsline{toc}{subsection}{13.5.1 滤波方程}
\begin{sourceblock}
\includegraphics[page=632,pagebox=mediabox,trim={64.39bp 1680.79bp 64.39bp 201.97bp},clip,width=0.94\textwidth,max height=0.38\textheight]{Pope.pdf}
\end{sourceblock}
\begin{sourceblock}
\includegraphics[page=632,pagebox=mediabox,trim={64.39bp 1350.83bp 64.39bp 571.93bp},clip,width=0.94\textwidth,max height=0.38\textheight]{Pope.pdf}
\end{sourceblock}
\begin{sourceblock}
\includegraphics[page=632,pagebox=mediabox,trim={64.39bp 1220.85bp 64.39bp 713.91bp},clip,width=0.94\textwidth,max height=0.38\textheight]{Pope.pdf}
\end{sourceblock}
\begin{sourceblock}
\includegraphics[page=632,pagebox=mediabox,trim={64.39bp 1021.87bp 64.39bp 915.88bp},clip,width=0.94\textwidth,max height=0.38\textheight]{Pope.pdf}
\end{sourceblock}
三：坐标方向。那么速度场 u(x, t) 具有傅里叶级数

\begin{sourceblock}
\includegraphics[page=633,pagebox=mediabox,trim={64.39bp 1688.79bp 64.39bp 255.97bp},clip,width=0.94\textwidth,max height=0.38\textheight]{Pope.pdf}
\end{sourceblock}
其中波数向量是 2z/L 的整数倍
K KO (方程(6.105)-(6.109))。傅里叶系数 h ( \textasciitilde{}t) 满足共轭对称性（方程 (6.121)）并且正交于
对于具有截止波数的锐谱滤波器，滤波速度场的傅里叶系数为

\begin{sourceblock}
\includegraphics[page=633,pagebox=mediabox,trim={64.39bp 1352.83bp 64.39bp 585.93bp},clip,width=0.94\textwidth,max height=0.38\textheight]{Pope.pdf}
\end{sourceblock}
其中 H 是赫维赛德函数，1x1 是波数的大小。滤波后速度的傅里叶级数为

\begin{sourceblock}
\includegraphics[page=633,pagebox=mediabox,trim={64.39bp 1125.86bp 64.39bp 764.90bp},clip,width=0.94\textwidth,max height=0.38\textheight]{Pope.pdf}
\end{sourceblock}
在 LES 计算中，所表示的 N3 波数位于边长 2 \textasciitilde{} i,n 波数空间的立方体内的均匀网格上。在每个坐标方向上，表示的波数为-(LA/l)x0,-(\$N Z)KO.,。。 , --KO0,, KO,. , ('N l)lco,\{ N K \textasciitilde{}T.h e 截止波数和模式数 N3 的关系为

\begin{sourceblock}
\includegraphics[page=633,pagebox=mediabox,trim={64.39bp 811.90bp 64.39bp 1125.86bp},clip,width=0.94\textwidth,max height=0.38\textheight]{Pope.pdf}
\end{sourceblock}
请注意，滤波器被视为各向同性，因此等式中的第二个和： (13.193) 仅限于半径球内的波数（在边 2 \textasciitilde{} \textasciitilde{}E 的立方体内）qu。 i价地，所表示的系数t)中非零的分数是：大约-3 Z K； / ( \textasciitilde{} K , ) \textasciitilde{} =
波数空间中的纳维-斯托克斯方程（第 214 页的方程（6.145））可以写成

\begin{sourceblock}
\includegraphics[page=633,pagebox=mediabox,trim={64.39bp 372.95bp 64.39bp 1516.81bp},clip,width=0.94\textwidth,max height=0.38\textheight]{Pope.pdf}
\end{sourceblock}
其中 Pjk 是投影张量（方程（6.133））。克罗内克三角洲是三元相互作用的统一，即。

\begin{sourceblock}
\includegraphics[page=633,pagebox=mediabox,trim={64.39bp 204.97bp 64.39bp 1742.78bp},clip,width=0.94\textwidth,max height=0.38\textheight]{Pope.pdf}
\end{sourceblock}
否则为零。方程（13.195）是一组普通微分的耦合组
从方程获得的滤波方程。 (13.195) 适用于非零系数 <

\begin{sourceblock}
\includegraphics[page=634,pagebox=mediabox,trim={64.39bp 1596.80bp 64.39bp 350.96bp},clip,width=0.94\textwidth,max height=0.38\textheight]{Pope.pdf}
\end{sourceblock}
这是一组有限（大约 7rN3/6）常微分方程，但它是非封闭的。出现闭合问题是因为非线性项（方程 (13 .197) 的右侧）包含未知的傅立叶系数，即 i i ( \textasciitilde{} ' a) nd ;(xu) 对于 r 2 和 ( 2
考虑高雷诺数均匀各向同性湍流的 LES
使用锐光谱滤波器和伪谱方法的 bulence
具有最大解析波数 正如在 DNS 中一样，足够
大尺度的分辨率要求域的大小 L 为
八个积分长度尺度，即 LC：8Lll（见方程（9.5））。这样就没有
超过 20\% 的能量存在于残余运动中，即截止
波数取 lc,Lll 15（参见第 240 页表 6.2）。

证明这些要求决定了

\begin{sourceblock}
\includegraphics[page=634,pagebox=mediabox,trim={64.39bp 906.89bp 64.39bp 1014.87bp},clip,width=0.94\textwidth,max height=0.38\textheight]{Pope.pdf}
\end{sourceblock}
第383章

\subsection*{13.5.2 三波相互作用}\phantomsection\addcontentsline{toc}{subsection}{13.5.2 三波相互作用}
当然，方程中的非线性项之间存在直接联系。 (13 .197)，物理空间中的速度乘积和各种应力
ukup，由它组成（例如，ukue）。从方程。 (13 .19 1) 我们得到

\begin{sourceblock}
\includegraphics[page=634,pagebox=mediabox,trim={64.39bp 315.96bp 64.39bp 1521.81bp},clip,width=0.94\textwidth,max height=0.38\textheight]{Pope.pdf}
\end{sourceblock}
图 13.10 是不同性质的三波相互作用的草图。这些被称为类型 ( a - ( d ) 并在表 13.4 中定义。就 LES (< K ) 中表示的模式而言，只有类型的相互作用

\begin{sourceblock}
\includegraphics[page=635,pagebox=mediabox,trim={64.39bp 546.93bp 64.39bp 291.76bp},clip,width=\textwidth,max height=0.72\textheight]{Pope.pdf}
\par\vspace{0.45\baselineskip}
\small 图 13.10。表 13.4 中定义的各种类型的三元相互作用的草图：(a) 解析，(b) Leonard，(c) 交叉，(d) SGS。
\end{sourceblock}
仅限于这些波数。限制 < 和 < UkUt K' K, K" K,
ak 相当于替换为 gt 和向上替换为 gt，并且限制 < 金额
uk K K，过滤结果。因此，类型 (a) 的相互作用的总和产生

\begin{sourceblock}
\includegraphics[page=635,pagebox=mediabox,trim={64.39bp 119.39bp 64.39bp 1800.77bp},clip,width=0.94\textwidth,max height=0.38\textheight]{Pope.pdf}
\par\vspace{0.45\baselineskip}
\small 表 13.4。图 13.10 中描绘的三元相互作用类型的定义，以及它们对（方程（13.199））的贡献（包括 kK < Kk 的所有可能的相互作用，但还有 k > k 的其他相互作用）
\end{sourceblock}
\begin{sourceblock}
\includegraphics[page=636,pagebox=mediabox,trim={64.39bp 1036.87bp 64.39bp 300.96bp},clip,width=\textwidth,max height=0.72\textheight]{Pope.pdf}
\end{sourceblock}
给定傅立叶系数 i ( \textasciitilde{}，类型 (b) 的三重相互作用是已知的，但它们产生超出滤波器截止的波数贡献 类型 (a) 和 (b) 对 in 的贡献之和
(KK) 方程(13 .199) 是 akaf：因此仅来自类型 (b) 的贡献是

\begin{sourceblock}
\includegraphics[page=636,pagebox=mediabox,trim={64.39bp 793.90bp 64.39bp 1152.85bp},clip,width=0.94\textwidth,max height=0.38\textheight]{Pope.pdf}
\end{sourceblock}
即伦纳德应力的负值（方程（13.100））。 （回想一下，使用锐光谱滤波器，伦纳德应力、十字应力和 SGS 雷诺应力的两种不同定义是等效的。）

如前所述，对于波数空间中的 LES，残余应力张量最好定义为

\begin{sourceblock}
\includegraphics[page=636,pagebox=mediabox,trim={64.39bp 498.94bp 64.39bp 1447.82bp},clip,width=0.94\textwidth,max height=0.38\textheight]{Pope.pdf}
\end{sourceblock}
使得过滤后的产物分解为

\begin{sourceblock}
\includegraphics[page=636,pagebox=mediabox,trim={64.39bp 367.95bp 64.39bp 1574.80bp},clip,width=0.94\textwidth,max height=0.38\textheight]{Pope.pdf}
\end{sourceblock}
这样，方程中的每一项都可以得到。 (13.203) 是一个滤波量，因此波数没有大于 的贡献。相反，在

\begin{sourceblock}
\includegraphics[page=636,pagebox=mediabox,trim={64.39bp 164.98bp 64.39bp 1722.78bp},clip,width=0.94\textwidth,max height=0.38\textheight]{Pope.pdf}
\end{sourceblock}
ZK，；并且（因为不包含此类模式），需要残余应力模型来消除这些贡献，这是一个很高的要求！

(c) 类型的三元交互是表示模式（例如 < 和产生贡献的残差模式）之间的交互
<K,.这种相互作用会产生交叉应力（方程（13.101））。

在类型 (d) 的相互作用中，两个残差模式 > 和 >- ) 给出
(K' K, K" K 上升到代表贡献 < K )。这些相互作用产生

\$ui。 SGS Reynolds 强调，类型 (c) 和 (d) 的三重相互作用共同产生（方程 (13.202)）。

过滤后的纳维-斯托克斯方程（方程（13 .197））可以根据这些不同类型的三波相互作用重写（对于 < 为

\begin{sourceblock}
\includegraphics[page=637,pagebox=mediabox,trim={64.39bp 1243.84bp 64.39bp 645.92bp},clip,width=0.94\textwidth,max height=0.38\textheight]{Pope.pdf}
\end{sourceblock}
其中 F' 来自类型 (a) 已解决的相互作用，

\begin{sourceblock}
\includegraphics[page=637,pagebox=mediabox,trim={64.39bp 1111.86bp 64.39bp 828.90bp},clip,width=0.94\textwidth,max height=0.38\textheight]{Pope.pdf}
\end{sourceblock}
是 iri 闭合形式；而 F?(K,t ) 由类型 (c) 的相互作用产生

\begin{sourceblock}
\includegraphics[page=637,pagebox=mediabox,trim={64.39bp 859.89bp 64.39bp 974.88bp},clip,width=0.94\textwidth,max height=0.38\textheight]{Pope.pdf}
\end{sourceblock}
并且必须建模。

\subsection*{13.5.3 光谱能量平衡}\phantomsection\addcontentsline{toc}{subsection}{13.5.3 光谱能量平衡}
各种能量和谱方程可以从方程导出。 (13.205)。最简单的是，让

\begin{sourceblock}
\includegraphics[page=637,pagebox=mediabox,trim={64.39bp 515.93bp 64.39bp 1407.82bp},clip,width=0.94\textwidth,max height=0.38\textheight]{Pope.pdf}
\end{sourceblock}
表示波数K模的瞬时能量。则由式： （13.205），我们有

\begin{sourceblock}
\includegraphics[page=637,pagebox=mediabox,trim={64.39bp 306.96bp 64.39bp 1591.80bp},clip,width=0.94\textwidth,max height=0.38\textheight]{Pope.pdf}
\end{sourceblock}
其中 T f ( \textasciitilde{}t)、a 和 Tr(x,t ) 定义为

\begin{sourceblock}
\includegraphics[page=637,pagebox=mediabox,trim={64.39bp 181.98bp 64.39bp 1763.78bp},clip,width=0.94\textwidth,max height=0.38\textheight]{Pope.pdf}
\end{sourceblock}
\begin{sourceblock}
\includegraphics[page=637,pagebox=mediabox,trim={64.39bp 119.39bp 64.39bp 1850.77bp},clip,width=0.94\textwidth,max height=0.38\textheight]{Pope.pdf}
\end{sourceblock}
等式右侧的第一项。 (13.209) 是解析运动的分子耗散，并且始终为负。第二项 Tf(x,t ) 表示通过类型 (a) 的三波相互作用从其他解析模式转移到波数的能量速率。在这种相互作用中，总能量是守恒的（参见第 214 页的练习 6.17），因此当对所有解析模式求和时，该项消失：

\begin{sourceblock}
\includegraphics[page=638,pagebox=mediabox,trim={64.39bp 1501.81bp 64.39bp 428.95bp},clip,width=0.94\textwidth,max height=0.38\textheight]{Pope.pdf}
\end{sourceblock}
方程中的最后一项。 (13 .209) Tr(K,t ) 表示通过类型 (c) 和 (d) 的三波相互作用从残余运动获得能量的速率。正值 (T, 0) 对应于反向散射，而该项主要为负值（前向散射）。能量传递到剩余运动的预期速率为

\begin{sourceblock}
\includegraphics[page=638,pagebox=mediabox,trim={64.39bp 1131.86bp 64.39bp 744.91bp},clip,width=0.94\textwidth,max height=0.38\textheight]{Pope.pdf}
\end{sourceblock}
13.5.4 为了进行大涡模拟，式中的非闭项F'(K,t) （13.205）必须建模。最简单且常用的模型类型是基于谱涡粘性的。残余运动的净效应应该与分子耗散的净效应类似，但光谱粘度取决于波数

\subsection*{13.5.4 谱涡粘性}\phantomsection\addcontentsline{toc}{subsection}{13.5.4 谱涡粘性}
并在截止点上因此模型为 rc

\begin{sourceblock}
\includegraphics[page=638,pagebox=mediabox,trim={64.39bp 677.91bp 64.39bp 1253.84bp},clip,width=0.94\textwidth,max height=0.38\textheight]{Pope.pdf}
\end{sourceblock}
从 Kraichnan (1976) 的工作开始，各种湍流理论已被用来估计 Chollet 和 Lesieur (1981) 提出了

\begin{sourceblock}
\includegraphics[page=638,pagebox=mediabox,trim={64.39bp 432.94bp 64.39bp 1415.82bp},clip,width=0.94\textwidth,max height=0.38\textheight]{Pope.pdf}
\end{sourceblock}
这是“基于截止波数 E(K,, t) 处的能谱函数的值； Chollet (1984) 提出的无量纲函数 v + 的形式为

\begin{sourceblock}
\includegraphics[page=638,pagebox=mediabox,trim={64.39bp 203.97bp 64.39bp 1715.78bp},clip,width=0.94\textwidth,max height=0.38\textheight]{Pope.pdf}
\end{sourceblock}
如图13.11所示。 AS 首先由 Kraichnan (1976) 预测，

\begin{sourceblock}
\includegraphics[page=639,pagebox=mediabox,trim={64.39bp 1258.84bp 64.39bp 109.44bp},clip,width=\textwidth,max height=0.72\textheight]{Pope.pdf}
\par\vspace{0.45\baselineskip}
\small 图 13.14。无量纲谱粘度方程。 (13.216)。
\end{sourceblock}
随着接近 rc，。根据谱涡粘模型，仅存在前向散射，E(K,t ) 以以下速率减小

\begin{sourceblock}
\includegraphics[page=639,pagebox=mediabox,trim={64.39bp 1008.87bp 64.39bp 922.88bp},clip,width=0.94\textwidth,max height=0.38\textheight]{Pope.pdf}
\end{sourceblock}
由于这个建模术语。

13.5.5 预期能量转移率 (P,) 可分解为前向散射和后向散射贡献：

\subsection*{13.5.5 反向散射}\phantomsection\addcontentsline{toc}{subsection}{13.5.5 反向散射}
\begin{sourceblock}
\includegraphics[page=639,pagebox=mediabox,trim={64.39bp 631.92bp 64.39bp 1313.83bp},clip,width=0.94\textwidth,max height=0.38\textheight]{Pope.pdf}
\end{sourceblock}
\begin{sourceblock}
\includegraphics[page=639,pagebox=mediabox,trim={64.39bp 453.94bp 64.39bp 1470.81bp},clip,width=0.94\textwidth,max height=0.38\textheight]{Pope.pdf}
\end{sourceblock}
\begin{sourceblock}
\includegraphics[page=639,pagebox=mediabox,trim={64.39bp 282.96bp 64.39bp 1592.80bp},clip,width=0.94\textwidth,max height=0.38\textheight]{Pope.pdf}
\end{sourceblock}
Leslie 和 Quarini (1979) 发现反向散射可能非常显着，其比率 ( p r b )( P,) 大于 1。

通过添加白光，可以将反向散射包含在 F t) 的模型中
其控制随机微分方程可以写成

\begin{sourceblock}
\includegraphics[page=640,pagebox=mediabox,trim={64.39bp 1666.79bp 64.39bp 218.97bp},clip,width=0.94\textwidth,max height=0.38\textheight]{Pope.pdf}
\end{sourceblock}
其中各向同性维纳过程 W(x, ¢) 对于每种模式都是独立的，

\begin{sourceblock}
\includegraphics[page=640,pagebox=mediabox,trim={64.39bp 1542.81bp 64.39bp 398.95bp},clip,width=0.94\textwidth,max height=0.38\textheight]{Pope.pdf}
\end{sourceblock}
\&, 是指定系数。从方程 (13.221) 推导出来的模式 (E(x, 0)) 的预期能量演化为

\begin{sourceblock}
\includegraphics[page=640,pagebox=mediabox,trim={64.39bp 1332.83bp 64.39bp 566.93bp},clip,width=0.94\textwidth,max height=0.38\textheight]{Pope.pdf}
\end{sourceblock}
因此 E 被确定为反向散射能量添加的平均速率（见练习 13.35）。

练习 13.35 验证方程。 (13.221) 与连续性方程一致
其中 x, dii,(x, t) 完全为零。回顾 P 的定义，
(式(6.133))，得到结果

\begin{sourceblock}
\includegraphics[page=640,pagebox=mediabox,trim={64.39bp 938.88bp 64.39bp 995.87bp},clip,width=0.94\textwidth,max height=0.38\textheight]{Pope.pdf}
\end{sourceblock}
因此验证方程。 (13.223)。

\subsection*{13.5.6 LES 的统计视图}\phantomsection\addcontentsline{toc}{subsection}{13.5.6 LES 的统计视图}
残余运动的理想 LES 模型是什么？也就是说，在物理空间中，残余应力 T 的理想模型是什么？或者，在波数空间中，F* 的理想模型是什么？

对于前面章节中考虑的湍流模型，相应的问题很容易回答。例如，理想的湍流黏度模型会产生正确的雷诺应力场（发生在所考虑的流动中）。然而，对于 LES 来说，这个问题更难回答。

为了检查所涉及的问题，请考虑使用具有 Nang 模式的伪谱方法执行均匀各向同性湍流的精确 DNS。我们还考虑使用截止波数 kK 的相同流的 LES，因此表示的模式数量显着减少 (Nies < Néns)。 AS 在实践中经常使用，DNS 的初始条件（在 ¢ = 0 处）是通过指定傅里叶幅度 [é(x,0)| 设置的。

\begin{sourceblock}
\includegraphics[page=641,pagebox=mediabox,trim={64.39bp 1046.87bp 64.39bp 109.44bp},clip,width=\textwidth,max height=0.72\textheight]{Pope.pdf}
\par\vspace{0.45\baselineskip}
\small 图 13.12。各向同性湍流的 DNS 和 LES 样本路径草图：(a) DNS, ty\}； (b) 过滤后的 DNS，\{a(<,t)\}； (c) 具有确定性残余应力模型的 LES，\{u(x,t)\}； (d) 具有反向散射模型的 LES，\{ua x,t)\}。 (a)中的路径在状态空间中； (b)-(d)中的那些处于Npc维状态空间中。
\end{sourceblock}
阶段是随机设置的。我们考虑 DNS 的不同实现，其中初始条件 0 相同，但不同
（正如 Piomelli 和 Chasnov (1996) 所建议的那样）。

图 13.12(a) 是 DNS 不同实现上的流程（在 NhNS 维状态空间中）演化的草图。初始状态不同（由于 > I C 的阶段不同，因此样本路径不同。（状态 \{ f ( \textasciitilde{}t)\} 根据纳维-斯托克斯方程确定性演化，因此样本路径在 NAN 维状态空间中不会交叉。）

图13.12(b)显示了过滤后的DNS字段在N\&维状态空间中的演化。这些样本路径是图 13.12(a) 中所示的 DNS 样本路径在低维 LES 状态空间上的投影。初始状态 \{Z(K0, )) 在结构上是相同的。然而，随着时间的推移，不同证悟的状态会分散。

对于确定性 LES 模型，例如谱涡粘性模型，
初始条件，如图13.12(c)所示。对于随机 LES 模型，例如 Eq. （13.221）我们有不同实现的状态分散（图13.12（d））；但分散的原因（白噪声）与过滤 DNS 字段的情况完全不同。

从这些考虑可以得出以下结论。

(i) 过滤后的 DNS 字段 t)\} 不是由
过滤初始条件O)\}。

(ii) 因此，不可能构建一个产生过滤的 LES 模型
速度场与 DNS 实现中的速度场相匹配。

(iii) 相反，可以实现的最好结果是统计对应
LES 和 DNS 过滤字段之间。

(iv) 通信（原则上）可以实现一次性级别
通过确定性模型进行统计，其中 F'(K, t) 被替换为
它的条件期望

\begin{sourceblock}
\includegraphics[page=642,pagebox=mediabox,trim={64.39bp 1094.86bp 64.39bp 849.89bp},clip,width=0.94\textwidth,max height=0.38\textheight]{Pope.pdf}
\end{sourceblock}
（这个结果直接来自进化方程
傅立叶系数 t)\} 的一次性联合 PDF，并且
取决于初始联合 PDF 是连续的，而不是
上面考虑的 0)) 上的急剧初始条件：参见 Adrian
(19 90) 和练习 13 .36。)

(v) 模型和模型之间相关系数的先验检验
测量的残余应力的意义值得怀疑。随着
理想模型，方程。 (13.225)，获得 F' 的相关系数
来自 DNS 和来自模型的结果可能明显小于统一。

匹配仅适用于条件均值，即均值 con-
涉及整个领域。因此，更有意义的是
基于条件统计的测试，由 Piomelli、Yu 执行，
例如，Adrian (1996)。

考虑一个随机向量过程 u(t) = \{ul (t), u2(t), ..., uN(t)\}
根据普通的随机初始条件 u(0) 演化而来

\begin{sourceblock}
\includegraphics[page=642,pagebox=mediabox,trim={64.39bp 169.98bp 64.39bp 1684.79bp},clip,width=0.94\textwidth,max height=0.38\textheight]{Pope.pdf}
\end{sourceblock}
\textasciitilde{} ( t ) f , ( u ; t)，演变为

\begin{sourceblock}
\includegraphics[page=643,pagebox=mediabox,trim={64.39bp 1697.79bp 64.39bp 222.97bp},clip,width=0.94\textwidth,max height=0.38\textheight]{Pope.pdf}
\end{sourceblock}
\begin{sourceblock}
\includegraphics[page=643,pagebox=mediabox,trim={64.39bp 1543.81bp 64.39bp 397.95bp},clip,width=0.94\textwidth,max height=0.38\textheight]{Pope.pdf}
\end{sourceblock}
考虑第二个向量过程 ii(t) \{iii(t), ii2(t), dm6,。 7(t)\} 那
从与 u ( e k( O) u(O)) 相同的初始条件演化而来，符合-
求解确定性方程

\begin{sourceblock}
\includegraphics[page=643,pagebox=mediabox,trim={64.39bp 1288.84bp 64.39bp 631.92bp},clip,width=0.94\textwidth,max height=0.38\textheight]{Pope.pdf}
\end{sourceblock}
证明 ii(t) 的一次性 PDF 与 u(t) 的一次性 PDF 相同。

论证这一观察结果证实了第(iv)项中提出的主张

\subsection*{13.5.7 分辨率和建模}\phantomsection\addcontentsline{toc}{subsection}{13.5.7 分辨率和建模}
使用尖锐的光谱滤波器（具有截止波数），分辨率和建模问题变得清晰且分开。必须表示所有波数，并且表示更高波数模式没有任何好处。过滤后的速度场用完整的表示
2 ( \textasciitilde{} , I K I 傅里叶模式 t 的精度) 滤波后的速度场不包含有关残余运动 2 的直接信息2，因此必须对这些残余运动的影响进行建模。

对于高斯滤波器或具有严格正传递函数 G(K) 的任何其他滤波器，分辨率和建模问题不太清晰和独立。使用伪谱方法考虑均匀湍流（具有 (U) 0）的 DNS 和 LES。在 DNS 中，所有波数都被表示，从而确保瞬时的足够分辨率
KDNS 正切速度场被选择为大于 1.5。

KDNS K D N S \textasciitilde{} 从给定的初始条件，纳维-斯托克斯方程 (Eq. (13.195)) 及时向前积分以确定傅立叶系数 i i ( \textasciitilde{}t)，为
LES 是使用具有特征波的高斯滤波器执行的
n-传递函数为的数字

\begin{sourceblock}
\includegraphics[page=644,pagebox=mediabox,trim={64.39bp 1661.79bp 64.39bp 217.97bp},clip,width=0.94\textwidth,max height=0.38\textheight]{Pope.pdf}
\end{sourceblock}
正如第 13.2 节中所讨论的，所得到的滤波速度场 t) 的光谱含量超出了因此，在 LES 中，模态表示为更高的波数，例如 2 \textasciitilde{} ,。

克莱斯, 克莱斯 =
LES 方程的形式取决于 uiujm 分解的选择，其中分解 uiuj uiuj + z R 方程为（对于

\begin{sourceblock}
\includegraphics[page=644,pagebox=mediabox,trim={64.39bp 1247.84bp 64.39bp 666.92bp},clip,width=0.94\textwidth,max height=0.38\textheight]{Pope.pdf}
\end{sourceblock}
这与纳维-斯托克斯方程（Eq. (13.195)）相同，只是它是针对滤波场的傅立叶系数而写的；波数
和 被限制在半径范围内，并且 IC、K'、K" ICLES；已添加代表残差模型的术语 F?(Kt,)。

在典型的应用中，波数的相对大小可能
1 hDNS, be 和 2 \textasciitilde{} , LES 需要较少的模式
K, = -8K DNS KLES == 4 SO 乘以 43 64 倍。这种节省是以 F'(K, t) 建模中涉及的不确定性为代价的。然而，应该指出的是，需要对 F' 进行建模的明显闭合问题也可以被视为解决问题。因为，如果过滤字段完全解析，则可以确定 F'，如现在所示。

瞬时 u(x, t) 和滤波后的 t) ve:locity 场的傅立叶系数通过以下关系相关：

\begin{sourceblock}
\includegraphics[page=644,pagebox=mediabox,trim={64.39bp 574.93bp 64.39bp 1371.83bp},clip,width=0.94\textwidth,max height=0.38\textheight]{Pope.pdf}
\end{sourceblock}
由于（对于高斯滤波器）严格为正，因此可以反转该关系以得出

\begin{sourceblock}
\includegraphics[page=644,pagebox=mediabox,trim={64.39bp 345.96bp 64.39bp 1571.80bp},clip,width=0.94\textwidth,max height=0.38\textheight]{Pope.pdf}
\end{sourceblock}
通过将 Navier-Stokes 方程（方程 (13.195)）乘以 GA- (K) 可以获得系数 t) 的精确闭方程组：

\begin{sourceblock}
\includegraphics[page=644,pagebox=mediabox,trim={64.39bp 119.39bp 64.39bp 1770.78bp},clip,width=0.94\textwidth,max height=0.38\textheight]{Pope.pdf}
\end{sourceblock}
通过将该方程与 LES 方程（方程（13.231））进行比较，可以获得以下剩余运动的精确方程：

仅当过滤后的速度场完全解析以使所考虑的波数延伸至或超出时，这些结果才是准确的。如果，相反
KDNS 过滤后的速度系数仅表示到那时
K L E \textasciitilde{} KDNS, F t) (可分解为

\begin{sourceblock}
\includegraphics[page=645,pagebox=mediabox,trim={64.39bp 1342.83bp 64.39bp 578.93bp},clip,width=0.94\textwidth,max height=0.38\textheight]{Pope.pdf}
\end{sourceblock}
其中 t) 由等式右侧定义。 (13.235) 但随着
F!(JK，总和限制为 KI' 位于半径范围内 因此
K'KL \textasciitilde{} s。 F' 根据所表示的系数已知，而 F 必须建模。随着增加，F 的相对贡献减少，并且
KLES >> 变得可以忽略不计
KLES = KDNS。

对于大值的传递函数 G(K) 非常小，这引发了方程 1 的病态问题。 （13.234）和（13.235）。如果
2（\textasciitilde{}，滤波后的速度场由其傅立叶系数 t 表示），则可以避免病态（参见练习 13.38），但如果 t） 在物理空间中的网格上表示，就像在非均匀流的 LES 中一样，那么从 u(x, t) 确定的任务肯定是病态的，无论网格多么精细。因此，如果类推等式。 (13 .236)，分解为
然后是贡献zf。无法可靠地确定波长远大于 的运动（参见练习 13.37 和 13.39）。

我们现在从这些考虑中得出重要的结论。对于使用锐光谱滤波器（或其他投影）的 LES，结论如下。

(i) 滤波后的速度场完全准确地表示为
有限模式集。

(ii) 滤波后的速度场没有提供有关速度场的直接信息
残余运动。

(iii) 建模和解析问题是分开的。

对于使用高斯滤波器（或其他可逆滤波器，G(K) > 0）的 LES
(iv) 滤波后的速度场不能用完整的精确表示
由一组有限的模式来确定。准确度随
(v) 滤波后的速度场包含残差信息
议案。该信息可在多大程度上用于
代表残余运动影响的 LES 取决于
过滤字段的分辨率，以及它们是否是已知的
物理空间或波数空间中的网格。

(vi) 分辨率和建模本质上是相关的：如果过滤
场在波数空间中完全解析（以相同的代价
作为 DNS），那么残余运动的影响是已知的，并且没有
需要建模。

13.37 证明对于高斯滤波器（方程（13.230）），传递函数
在等式中(13.234) 可以重新表示为
如果 和 位于半径为 2 \textasciitilde{} , ,s 的球体内
指数项大于0.037；并表明，对于一个球体
半径为 4 \textasciitilde{}，t,h e 对应的界限为 1.9 x 10- 6.13 .38 涉及 v\% 量值差异很大的变量的计算
可能会生病。这个问题可以通过缩放来缓解
变量，以便它们的大小具有可比性。表明
求解缩放变量的 LES 方程（方程 (13.234)）

\textasciitilde{} ) / G ( K ) 相当于求解 DNS 方程，Eq. (13.231)。 13.39 给定用以下方法获得的完全解析的滤波速度场 B(x, t)

高斯滤波器，s h \textasciitilde{} wth 处的残余应力可以（原则上）

得到为

\begin{sourceblock}
\includegraphics[page=646,pagebox=mediabox,trim={64.39bp 504.94bp 64.39bp 1415.82bp},clip,width=0.94\textwidth,max height=0.38\textheight]{Pope.pdf}
\end{sourceblock}
对于>H.. (K)i 定义为
集成电路，，

\begin{sourceblock}
\includegraphics[page=646,pagebox=mediabox,trim={64.39bp 378.95bp 64.39bp 1558.80bp},clip,width=0.94\textwidth,max height=0.38\textheight]{Pope.pdf}
\end{sourceblock}
证明 r f w(E q. (13.237)) - 模态对 zR 的贡献
IJlJ

小于 - 的波数可以从 u(x, t) 获得：

\begin{sourceblock}
\includegraphics[page=646,pagebox=mediabox,trim={64.39bp 203.97bp 64.39bp 1730.78bp},clip,width=0.94\textwidth,max height=0.38\textheight]{Pope.pdf}
\end{sourceblock}
对这些操作的条件进行评论。

\section*{13.6 其他残余应力模型}\phantomsection\addcontentsline{toc}{section}{13.6 其他残余应力模型}
现在我们回到物理空间来考虑适用于非均匀湍流的残余应力模型。这些模型是由 Smagorinsky 模型的缺点引发的。

\subsection*{13.6.1 动态模型}\phantomsection\addcontentsline{toc}{subsection}{13.6.1 动态模型}
Smagorinsky模型的问题之一是系数Cs的合适值在不同的流态下是不同的。特别是，它在层流中为零，并且与高雷诺数自由湍流中的值 (Cs 0.15) 相比，它在壁附近衰减。动态模型提供了一种用于确定 Smagorinsky 系数的适当局部值的方法。该模型由 German0 等人提出。 (19 91)，Lilly (1992) 和 Meneveau 等人提供了重要的修改和扩展。 （1996）。

网格和测试滤波器动态模型涉及不同滤波器宽度的滤波器。为了简单起见，我们考虑均匀各向同性滤波器，并用 G A) 和 G(K;A ) 表示宽度为 A 的滤波器及其传递函数。

网格滤波器的滤波器宽度为 A，其与网格间距 h 成正比（例如+、A h，或者为了获得更好的分辨率，A 2h）。网格过滤的操作由上划线表示，例如，

\begin{sourceblock}
\includegraphics[page=647,pagebox=mediabox,trim={64.39bp 716.91bp 64.39bp 1185.85bp},clip,width=0.94\textwidth,max height=0.38\textheight]{Pope.pdf}
\end{sourceblock}
尽管未明确执行此过滤操作，但 LES 方程被认为已求解 U(x, t)。

测试过滤器的过滤器宽度为 A，通常取为 A 的两倍，并且测试过滤由波形符表示，例如，

\begin{sourceblock}
\includegraphics[page=647,pagebox=mediabox,trim={64.39bp 416.95bp 64.39bp 1485.81bp},clip,width=0.94\textwidth,max height=0.38\textheight]{Pope.pdf}
\end{sourceblock}
事实上，由于 U 在 LES 计算中是未知的，因此考虑应用于 U 的测试滤波器以产生双重滤波量 U 更为相关。

请注意，如果使用锐光谱滤波器，则与传递函数的乘积相同

\begin{sourceblock}
\includegraphics[page=647,pagebox=mediabox,trim={64.39bp 119.39bp 64.39bp 1790.77bp},clip,width=0.94\textwidth,max height=0.38\textheight]{Pope.pdf}
\end{sourceblock}
（对于 A 2 另一方面，对于高斯滤波器，我们有

\begin{sourceblock}
\includegraphics[page=648,pagebox=mediabox,trim={64.39bp 1729.78bp 64.39bp 207.97bp},clip,width=0.94\textwidth,max height=0.38\textheight]{Pope.pdf}
\end{sourceblock}
请参见第 567 页的练习 13.4。因此，对于两个过滤器，双重过滤操作可以写成

\begin{sourceblock}
\includegraphics[page=648,pagebox=mediabox,trim={64.39bp 1439.82bp 64.39bp 376.95bp},clip,width=0.94\textwidth,max height=0.38\textheight]{Pope.pdf}
\end{sourceblock}
其中有效双滤波器宽度为

\begin{sourceblock}
\includegraphics[page=648,pagebox=mediabox,trim={64.39bp 1241.84bp 64.39bp 618.92bp},clip,width=0.94\textwidth,max height=0.38\textheight]{Pope.pdf}
\end{sourceblock}
最小解析运动 速度 U = U u' 的分解很容易扩展到

\begin{sourceblock}
\includegraphics[page=648,pagebox=mediabox,trim={64.39bp 1037.87bp 64.39bp 884.89bp},clip,width=0.94\textwidth,max height=0.38\textheight]{Pope.pdf}
\end{sourceblock}
其中三个贡献代表（松散地）大于 A 的运动，即 A 和 A 之间的运动；最小解析运动 U U 可以从 U 确定，并且在包括动态模型在内的多个模型中在概念上很重要。

在物理空间中，U U 由 U 与滤波器差值 G(r; A) - G(r; A) 的卷积给出。对于均匀湍流，U - U 的傅里叶系数是 U 的傅里叶系数乘以传递函数差 G(lc;K) - G； A).这种差异如图 13.13 所示
具有 A 2 A 的 1 rv 锐谱滤波器和高斯滤波器：可以看出，这些滤波器分别对应于完美带通滤波器和不完美带通滤波器。

如果测试过滤器和网格过滤器被视为相同（ i， A LA- ），那么我们

\begin{sourceblock}
\includegraphics[page=648,pagebox=mediabox,trim={64.39bp 392.95bp 64.39bp 1488.81bp},clip,width=0.94\textwidth,max height=0.38\textheight]{Pope.pdf}
\end{sourceblock}
对于锐光谱滤波器为零；但是，对于高斯滤波器，相应的传递函数差异在形状上与 rv A 2A 的相似（见图 13.13）。

我们将 U - CULe 解释为由间距网格解析的最小解析运动 同样，C U L- - U 代表最大的运动

\begin{sourceblock}
\includegraphics[page=649,pagebox=mediabox,trim={64.39bp 1098.86bp 64.39bp 109.44bp},clip,width=\textwidth,max height=0.72\textheight]{Pope.pdf}
\par\vspace{0.45\baselineskip}
\small 图 13.13。与最小解析运动相对应的滤波器传递函数差异。实线，高斯滤波器，A = 2A；点-虚线，锐利光谱滤波器，A = 2A；虚线，高斯滤波器，A = A。
\end{sourceblock}
考虑具有锐角的高雷诺数各向同性湍流

在惯性子区内应用光谱网格和测试滤波器。

使用柯尔莫哥洛夫谱来估计

最小解析尺度的能量到剩余 k-in etic 能量
是 1。表明 A/A 2 的比率为 0.59 和 1，并且
Gerrnano 的恒等式 基于单次和双次过滤操作的残余应力为

\begin{sourceblock}
\includegraphics[page=649,pagebox=mediabox,trim={64.39bp 360.95bp 64.39bp 1522.81bp},clip,width=0.94\textwidth,max height=0.38\textheight]{Pope.pdf}
\end{sourceblock}
\begin{sourceblock}
\includegraphics[page=649,pagebox=mediabox,trim={64.39bp 273.96bp 64.39bp 1667.79bp},clip,width=0.94\textwidth,max height=0.38\textheight]{Pope.pdf}
\end{sourceblock}
通过将测试过滤器应用于方程（1992），可以获得由 Germano（1992）给出的恒等式。 (13.250) 并从方程中减去结果。 （13.251）：

\begin{sourceblock}
\includegraphics[page=649,pagebox=mediabox,trim={64.39bp 119.39bp 64.39bp 1835.77bp},clip,width=0.94\textwidth,max height=0.38\textheight]{Pope.pdf}
\end{sourceblock}
这一恒等式的重要性在于，Cij（称为解析应力）以 U 为单位，而 Ti 和 z 为已知；不是。

等式的右侧。 (13.252) 是基于网格滤波速度的残余应力（相对于测试滤波器）的形式。因此，对于测试滤波器，Lij 可以粗略地解释为贡献
rv 来自最大未解决运动的残余应力。另外，对于 A 4，Cij 在此由等式定义。 (13.252) 与方程 1 定义的 Leonard 应力 Lo 相同。 (13.100)。

Srnagorinsky 系数 Smagorinsky 模型（方程 (13.127) 和 (13.128)）的偏差部分
fi 可以写成

\begin{sourceblock}
\includegraphics[page=650,pagebox=mediabox,trim={64.39bp 1243.84bp 64.39bp 693.91bp},clip,width=0.94\textwidth,max height=0.38\textheight]{Pope.pdf}
\end{sourceblock}
该系数已被重新定义为 cs （代替，以允许负值的可能性，对应于反向散射。为宽度 A 的滤波器编写的相同模型方程为

\begin{sourceblock}
\includegraphics[page=650,pagebox=mediabox,trim={64.39bp 1008.87bp 64.39bp 905.89bp},clip,width=0.94\textwidth,max height=0.38\textheight]{Pope.pdf}
\end{sourceblock}
其中 Sij 和 S 是根据 U 直接定义的，与等式 1 类似。 （13.73）和（13.74）。取cs统一，并定义

\begin{sourceblock}
\includegraphics[page=650,pagebox=mediabox,trim={64.39bp 812.90bp 64.39bp 1116.86bp},clip,width=0.94\textwidth,max height=0.38\textheight]{Pope.pdf}
\end{sourceblock}
等式。 (13.253) 和 (13.254) 导致

\begin{sourceblock}
\includegraphics[page=650,pagebox=mediabox,trim={64.39bp 679.91bp 64.39bp 1256.84bp},clip,width=0.94\textwidth,max height=0.38\textheight]{Pope.pdf}
\end{sourceblock}
这是 Cij 偏差部分的 Smagorinsky 模型（方程 (13.252)），
在 LES 中，Mij 和 Ld。 （它们是 x 和 t 的函数）以 U(x, t) 形式已知。该信息可用于确定 Smagorinsky 系数 cs 的值，对于该值，模型 LE 提供 LdaO 的最佳近似值。 f 当然，不能选择单个系数 cs 来匹配 L、S/ 和 的五个独立分量。但是，如 Lilly (1992) 所示，均方误差通过规范最小化

\begin{sourceblock}
\includegraphics[page=650,pagebox=mediabox,trim={64.39bp 206.97bp 64.39bp 1736.78bp},clip,width=0.94\textwidth,max height=0.38\textheight]{Pope.pdf}
\end{sourceblock}
参见练习 13.41。 （请注意，由于 Mij 是偏数，因此 MvLij 等于
'偏应力 Ld 之间的均方误差。和
Smagorinsky 模型预测 LS（方程（13.256））定义为

\begin{sourceblock}
\includegraphics[page=651,pagebox=mediabox,trim={64.39bp 1635.79bp 64.39bp 299.96bp},clip,width=0.94\textwidth,max height=0.38\textheight]{Pope.pdf}
\end{sourceblock}
对cs求导可得

\begin{sourceblock}
\includegraphics[page=651,pagebox=mediabox,trim={64.39bp 1477.81bp 64.39bp 423.95bp},clip,width=0.94\textwidth,max height=0.38\textheight]{Pope.pdf}
\end{sourceblock}
\begin{sourceblock}
\includegraphics[page=651,pagebox=mediabox,trim={64.39bp 1360.83bp 64.39bp 534.93bp},clip,width=0.94\textwidth,max height=0.38\textheight]{Pope.pdf}
\end{sourceblock}
因此，根据方程 cs 的规范。 (13.258)

考虑动态模型在层流中的应用。

在滤波器宽度 A 量级的距离上，速度可以为
假设随 x 线性变化。使用练习 13.28 的结果

表明各种应力是

\begin{sourceblock}
\includegraphics[page=651,pagebox=mediabox,trim={64.39bp 917.88bp 64.39bp 936.88bp},clip,width=0.94\textwidth,max height=0.38\textheight]{Pope.pdf}
\end{sourceblock}
其中 a 是过滤器宽度比 a A。证明张量 Mij

\begin{sourceblock}
\includegraphics[page=651,pagebox=mediabox,trim={64.39bp 751.90bp 64.39bp 1137.86bp},clip,width=0.94\textwidth,max height=0.38\textheight]{Pope.pdf}
\end{sourceblock}
因此从方程可以看出。 (13.258) 动态模型得出

\begin{sourceblock}
\includegraphics[page=651,pagebox=mediabox,trim={64.39bp 579.93bp 64.39bp 1313.83bp},clip,width=0.94\textwidth,max height=0.38\textheight]{Pope.pdf}
\end{sourceblock}
验证以下均匀层流 de- 的 cs 值
形成（参见第 41 5 页表 11 .2）：轴对称收缩，
-0.012；轴对称膨胀； CS 0.012；平面应变，
CS 0;和剪切力，cs 0。

通道 JEow 的应用鉴于先验测试中观察到的应力和应变率之间的相关性较低，从方程 1 获得的 cs 值并不奇怪。 (13 .258) 表现出非常大的波动，包括负值。作为一个
为了克服这个问题，Germano 等人在计算完全发展的湍流槽道流时。 (1991) 和 Piomelli (1993) 对等式中的分子和分母进行平均。 （13.258）。？因此，将该平均值表示为
,,,,) Smagorinsky 系数的计算公式为

\begin{sourceblock}
\includegraphics[page=652,pagebox=mediabox,trim={64.54bp 1624.59bp 64.54bp 311.92bp},clip,width=0.94\textwidth,max height=0.38\textheight]{Pope.pdf}
\end{sourceblock}
通过这种实现，只要网格足够细以解析近壁含能运动，动态模型就可以很好地计算过渡和完全湍流槽道流。因此，不需要阻尼功能或其他特殊的近壁处理。在某种程度上，模型的成功可以归因于这样一个事实：它正确地产生了层流剪切流的 cs 0（参见练习 13.42），并且在靠近壁的地方，平均残余剪切应力 (zE) 正确地随着距壁的距离的立方而变化（参见练习 13.43）。

相同的方法已成功应用于时间演化的自由剪切流（例如，Ghosal 和 Rogers（1997）和 Vreman 等人（1997）），也应用于统计轴对称流（Akselvoll 和 Moin 1996），在这种情况下，方程 1 中的平均值为： (13.265) 是在圆周方向上执行的。在一些实现中，通过将 cs 的负值设置为零来执行“削波”，从而消除反向散射（Vreman 等人 1997，Zang、Street 和 Koseff 1993）。

练习 13.\#43 考虑通常坐标系中的湍流槽道流，因此
y x2 是距墙壁的正常距离。如图所示
第 600 页的练习 13.30，LES 速度场非常接近
wall（y 中的前导顺序）可以写成
其中 if、b 和 是 x、z 和 t 的随机函数。贯穿始终
在这个练习中，我们认为该区域非常靠近墙，这样
等式。 (13.266) 是有效的，并且只需要保留 y 中的前导项。

过滤仅应用于 x-z 平面，使用宽度均匀的过滤器
d.和 A. 来自方程。 (13.252) 得到结果

\begin{sourceblock}
\includegraphics[page=652,pagebox=mediabox,trim={64.54bp 297.92bp 64.54bp 1626.59bp},clip,width=0.94\textwidth,max height=0.38\textheight]{Pope.pdf}
\end{sourceblock}
并表明 Li 的分量随相同功率而变化
y 与旧应力和真实残余应力一样，
German0 \textasciitilde{}t a/.( 1991) 使用的 cs 方程与方程不完全相同。 （13.258）。

等式。 (13.181)。证明前阶速度梯度
独立于 y，因此 Mj 也独立于 y，因此，方程 1 也独立于 y。 （13.255）。因此

表明，根据动态模型，cg、v 和 t\{ 都变化
正确为 y3。比较模型的行为和正确的行为
tj的其他组成部分。

局部动态模型

对于一般流动（没有统计同质性方向），需要一种替代方法以获得 Smagorinsky 系数 cs 的稳定场。戈萨尔等人。 （1995）从动态模型的标准推导中存在不一致的观察出发，开发了这样一种方法。具体来说，由于 cs 在空间上变化，因此没有理由假设它是均匀的以进行评估（参见方程（13.253）-（13.256））。

如果分析时考虑cs的空间变化，则C4和C7之间的差异（误差）不再由cs的局部值决定。相反，可以解决约束变分问题来确定非负场 cs(x)，从而最小化 Li 和 LS 之间的全局差异。以这种方式确定cg(x)的Smagorinsky模型称为局域动态模型。 Piomelli 和 Liu (1995) 基于类似的想法描述了一个更简单的模型。

拉格朗日动态模型 对于一般的非均匀流动，可以考虑使用具有方程 1 中的平均值的动态模型。 (13.265) 对于 cs(x) 在 x 周围的某个体积上执行。然而，指定体积（尤其是靠近墙壁）存在一定的随意性和困难，并且评估每个网格点的体积平均值会产生计算成本。相反，Meneveau 等人。 (1996) 开发了拉格朗日动力学模型，其中加权平均值沿着流体粒子路径（基于过滤的速度场）在时间上向后形成。

根据该模型，Smagorinsky 系数计算为

\begin{sourceblock}
\includegraphics[page=653,pagebox=mediabox,trim={64.39bp 311.96bp 64.39bp 1634.79bp},clip,width=0.94\textwidth,max height=0.38\textheight]{Pope.pdf}
\end{sourceblock}
其中 Am 和 Jum 代表平均值和 (MijMjj)ave - 简单松弛方程

\begin{sourceblock}
\includegraphics[page=653,pagebox=mediabox,trim={64.39bp 119.39bp 64.39bp 1811.77bp},clip,width=0.94\textwidth,max height=0.38\textheight]{Pope.pdf}
\end{sourceblock}
AM 的求解，其中 T 是指定的弛豫时间。这相当于沿着粒子路径进行平均，在较早时间 t' 处使用相对权重 exp [-(t - t')/ TI。求解的类似方程为

\begin{sourceblock}
\includegraphics[page=654,pagebox=mediabox,trim={64.39bp 1598.80bp 64.39bp 302.96bp},clip,width=0.94\textwidth,max height=0.38\textheight]{Pope.pdf}
\end{sourceblock}
其中指标函数

\begin{sourceblock}
\includegraphics[page=654,pagebox=mediabox,trim={64.39bp 1475.81bp 64.39bp 471.94bp},clip,width=0.94\textwidth,max height=0.38\textheight]{Pope.pdf}
\end{sourceblock}
防止（并因此 cs）变为负值。

拉格朗日动力学模型已被其创始人应用于过渡和完全湍流的河道流，并取得了与原始平面平均动力学模型类似的成功。 Haworth 和 Jansen (2000) 对火花点火发动机气缸内的流量进行了 LES 计算，清楚地证明了其普遍适用性。

后验测试的一般结论是，只要解决了含能运动，动态模型就相当成功。该附带条件需要分辨率达到接近壁的粘性长度尺度 6 的量级（例如 L ES-NWR）。

鉴于这一成功，回顾模型两个组成部分的物理基础是有用的。第一个组成部分是基本 Smagorinsky 模型假设，即残余应力 zfi 与滤波后的应力一致
ZJ 应变率 Sij，Eq. (13.253)。正如第 13.4.6 节中所讨论的，先验测试表明 z' 和 Sij 之间的相关性比 Smagorinsky 模型所暗示的近乎完美的相关性要弱得多。因此，斯马戈林斯基系数 cs 的任何值都无法产生正确的 zf 水平和 P、-zr*Sij 的能量转移率。

模型的第二个组成部分动态方面基于 Smagorinsky 系数与滤波器宽度无关的假设（参见方程（13 .253）和（1 3.254））。同样，假设斯马戈林斯基长度尺度 ts 与过滤器宽度 A 成正比。对于过滤器位于高雷诺数湍流的惯性子区内的情况，该假设是有根据的，如方程 1 所示。 （13.131）。然而，该模型的成功归因于其流态行为，其假设缺乏层流、过渡流和粘性壁区域流的合理性。这样的考虑
LES [..] 中动态 Smagorinsky 子网格模型的良好后验性能似乎与它们正确表示子网格物理的能力关系不大。

Porte-Agel、Meneveau 和 Parlange (2000) 描述了一个与尺度相关的动态模型，其中假定 Smagorinsky 系数 cg 作为滤波器宽度 A 的未知幂而变化。通过在三个级别进行滤波，可以推导出假定幂律中的指数，从而确定 cg。

\subsection*{13.6.2 混合模型和变体}\phantomsection\addcontentsline{toc}{subsection}{13.6.2 混合模型和变体}
多年来，许多残余应力模型基于滤波速度梯度 0U,/dx；并提出了双重滤波速度 U 或 U。这些模型通常与 Smagorinsky 模型结合使用，形成所谓的混合模型。动态过程可用于确定模型系数。 Liu 等人对其中几个模型及其变体进行了先验测试。 （1994）。

巴迪纳模型

Bardina、Ferziger 和 Reynolds (1980) 提出的模型（也称为尺度相似模型）最简单地解释为明确纳入 Leonard 应力（尽管该模型的创始人和其他人提出了对该模型的不同物理解释）。正如第 13.3 节中所讨论的，残余应力可以通过不同的方式分解为伦纳德应力、交叉应力和 SGS 雷诺应力（参见方程（13.105）和（13.99））。因为分解的每个分量都是伽利略不变量，所以首选分解是由下式提出的：

\begin{sourceblock}
\includegraphics[page=655,pagebox=mediabox,trim={64.39bp 581.93bp 64.39bp 1317.83bp},clip,width=0.94\textwidth,max height=0.38\textheight]{Pope.pdf}
\end{sourceblock}
其中伦纳德应力为

\begin{sourceblock}
\includegraphics[page=655,pagebox=mediabox,trim={64.39bp 450.94bp 64.39bp 1485.81bp},clip,width=0.94\textwidth,max height=0.38\textheight]{Pope.pdf}
\end{sourceblock}
在 LES 中，U 是已知的，因此可以评估 Ly。

之前，我们采用 Smagorinsky 模型作为（方程（13.127））的偏数部分的模型。相反，如果我们将 Smagorinsky 模型作为 Cz +R) 的模型，并明确包含 Leonard 应力，则所得模型为

\begin{sourceblock}
\includegraphics[page=655,pagebox=mediabox,trim={64.39bp 149.98bp 64.39bp 1788.77bp},clip,width=0.94\textwidth,max height=0.38\textheight]{Pope.pdf}
\end{sourceblock}
Zang 等人报道了使用这种混合模型的成功 LES 计算，其中使用动态过程来确定 cs。 (1993) 和 Vreman 等人。 （1997）。事实上，对于湍流混合层，这是 Vreman 等人评估的六个模型中最成功的一个。 （1997）。

克拉克模型正如克拉克等人所观察到的。 (1979)，对于高斯滤波器等滤波器，伦纳德应力的近似值为
（这可以通过第 595 页练习 13.28 中执行的泰勒级数分析的改编来显示。）如果在 Eqw(13.274) 中使用此近似，则结果是混合模型

\begin{sourceblock}
\includegraphics[page=656,pagebox=mediabox,trim={64.39bp 1119.86bp 64.39bp 770.90bp},clip,width=0.94\textwidth,max height=0.38\textheight]{Pope.pdf}
\end{sourceblock}
右侧第一项是克拉克模型。弗雷曼等人。 (1997)将该模型（使用动态过程来确定cs）应用于湍流混合层，并获得了与 Bardina 混合模型相当的良好结果。

替代残余速度尺度 Smagorinsky 模型给出的涡黏度可以写成

\begin{sourceblock}
\includegraphics[page=656,pagebox=mediabox,trim={64.39bp 723.91bp 64.39bp 1217.85bp},clip,width=0.94\textwidth,max height=0.38\textheight]{Pope.pdf}
\end{sourceblock}
其中速度标度是

\begin{sourceblock}
\includegraphics[page=656,pagebox=mediabox,trim={64.39bp 598.92bp 64.39bp 1342.83bp},clip,width=0.94\textwidth,max height=0.38\textheight]{Pope.pdf}
\end{sourceblock}
已经提出了几种基于不同速度尺度定义的涡粘模型。

可以说，速度尺度应该代表残余运动。使用高斯滤波器等滤波器，残余速度的估计由恒等式提供

\begin{sourceblock}
\includegraphics[page=656,pagebox=mediabox,trim={64.39bp 284.96bp 64.39bp 1662.79bp},clip,width=0.94\textwidth,max height=0.38\textheight]{Pope.pdf}
\end{sourceblock}
基于这一观察，Bardina 等人。 (1980) 定义速度

\begin{sourceblock}
\includegraphics[page=656,pagebox=mediabox,trim={64.39bp 119.39bp 64.39bp 1763.78bp},clip,width=0.94\textwidth,max height=0.38\textheight]{Pope.pdf}
\end{sourceblock}
重心 = 0.126。 Colucci 等人描述了使用该模型的伽利略不变版本的 LES 计算。 （1998）。

结构函数模型 在 Métais 和 Lesieur (1992) 的结构函数模型中，速度尺度是从二阶结构函数获得的，定义为

\begin{sourceblock}
\includegraphics[page=657,pagebox=mediabox,trim={64.39bp 1476.81bp 64.39bp 421.95bp},clip,width=0.94\textwidth,max height=0.38\textheight]{Pope.pdf}
\end{sourceblock}
这是以 x 为中心、半径为 A 的球体表面上距离 A 上速度差的平方的面积平均值。用于定义 v 的速度标度（方程 13.277）为

\begin{sourceblock}
\includegraphics[page=657,pagebox=mediabox,trim={64.39bp 1252.84bp 64.39bp 678.91bp},clip,width=0.94\textwidth,max height=0.38\textheight]{Pope.pdf}
\end{sourceblock}
其中 c = 0.063。 Lesieur 和 Métais (1998) 对结构函数模型的理论基础进行了描述，并对扩展和应用进行了回顾。

练习 13.44 证明 F 的一阶估计（方程（13.281））基于
U(x +r) 的泰勒级数展开式为
其中Q；是滤波速度场的旋转率。如果
应变率和旋转率不变量相等，表明
结构-功能模型（基于方程（13.283））简化为
斯马戈林斯基模型

\begin{sourceblock}
\includegraphics[page=657,pagebox=mediabox,trim={64.39bp 591.92bp 64.39bp 1353.83bp},clip,width=0.94\textwidth,max height=0.38\textheight]{Pope.pdf}
\end{sourceblock}
\subsection*{13.6.3 输运方程模型}\phantomsection\addcontentsline{toc}{subsection}{13.6.3 输运方程模型}
上面考虑的所有残余应力张量模型都涉及

到在同一时间 t 且位于 x 附近的滤波速度场 U(x',t)。为了追求更准确的模型，很自然地通过 1 的传输方程以及与剩余运动相关的其他量来纳入历史和非局部效应。这种方法与传统的湍流建模并行：就像 Smagorinsky 模型一样
与一方程模型类似，已经提出了代数应力模型、雷诺应力模型和 PDF 模型。然而，不同之处在于，在每种情况下，滤波器宽度 A 都被视为残余运动的特征长度尺度，因此不使用尺度方程（类似于 或 的模型方程。

Deardorff (1974) 的早期工作是 LES 中求解残余应力模型输运方程的少数例子之一。的
zR 模型方程与雷诺应力闭合的形式相同：生产是闭式的；再分配由罗塔模型和简单快速模型建模；运输是通过梯度扩散进行的；并且耗散被认为是各向同性的。残余动能的瞬时耗散率 k, = 1 R 与 k 相关，滤波器宽度 A 为

\begin{sourceblock}
\includegraphics[page=658,pagebox=mediabox,trim={64.39bp 1202.85bp 64.39bp 727.91bp},clip,width=0.94\textwidth,max height=0.38\textheight]{Pope.pdf}
\end{sourceblock}
其中常数CE取为CE 0.7。如练习13.45所示，在高雷诺数和惯性子区内的尖锐光谱滤波器的理想条件下，这种形式是合理的，并且常数CE可以与柯尔莫哥洛夫常数C相关。

为了降低计算成本，随后 Deardorff (1980) 恢复为各向同性涡粘模型，其中涡粘粘度为

\begin{sourceblock}
\includegraphics[page=658,pagebox=mediabox,trim={64.39bp 822.90bp 64.39bp 1124.86bp},clip,width=0.94\textwidth,max height=0.38\textheight]{Pope.pdf}
\end{sourceblock}
并使用模型输运方程来确定残余动能。同样，在理想条件下，v 的这种形式可以被证明是合理的，并且可以估计常数 C 0.1 的值（参见练习 13.45）。这个
NN 类型的模型已被气象界广泛使用，例如 Mason (1989)、Schmidt 和 Schumann (1989) 以及 Sullivan、McWilliams 和 Moeng (19 94)。 Schmidt 和 Schumann (19 89) 没有使用涡粘模型，而是通过类似于代数应力模型的输运方程的代数近似来获得残余应力。

Wong (1992) 描述了（但没有演示）一种方法，其中求解对应于两个不同滤波器宽度和 A 的残余 k-i 网络能量和 k, 的方程。动态过程可用于确定两个系数 CE 和 C,。 Ghosal 等人描述了 k 方程与局部动态模型相结合的不同用途。 （1995）。

定义为，

\begin{sourceblock}
\includegraphics[page=659,pagebox=mediabox,trim={64.39bp 1720.78bp 64.39bp 219.97bp},clip,width=0.94\textwidth,max height=0.38\textheight]{Pope.pdf}
\end{sourceblock}
对于正滤波器 (G 2 O)，f 具有 PDF 的所有属性（参见第 584 页的练习 13.21）。的一阶矩是滤波速度 U 的分量；第二个力矩是残余应力 zR
ijw Gicquel 等人。 (1998) 使用 f ( V ; x, t) 的模型方程进行 LES 计算，类似于 Langevin 模型。

结合 k 模型方程的 LES 方法已被证明在存在其他效应（例如浮力）的气象应用中具有优势。基于过滤密度函数的方法对于反应流也是有利的，其中混合和反应的重要过程发生在亚格子尺度上。然而，除了这些例外，一般经验是，在 LES 中，求解额外输运方程所涉及的额外计算成本和复杂性并不能保证精度的提高。使用额外输运方程的好处在 VLES 中更为明显，其中更大一部分能量存在于亚格子尺度的运动中。

考虑高雷诺数均匀湍流

惯性子区内的锐光谱滤波器。使用科尔姆奥戈洛夫
频谱（C = 1.5）和关系式 ( c , ) w c 以获得
CE 估计（方程 (13.285)）

\begin{sourceblock}
\includegraphics[page=659,pagebox=mediabox,trim={64.39bp 705.91bp 64.39bp 1241.84bp},clip,width=0.94\textwidth,max height=0.38\textheight]{Pope.pdf}
\end{sourceblock}
涡黏度由方程式给出。 （13.286），表明近似

\begin{sourceblock}
\includegraphics[page=659,pagebox=mediabox,trim={64.39bp 539.93bp 64.39bp 1351.83bp},clip,width=0.94\textwidth,max height=0.38\textheight]{Pope.pdf}
\end{sourceblock}
导致估计

\begin{sourceblock}
\includegraphics[page=659,pagebox=mediabox,trim={64.39bp 405.95bp 64.39bp 1525.81bp},clip,width=0.94\textwidth,max height=0.38\textheight]{Pope.pdf}
\end{sourceblock}
13.6.4 在 LES 动量方程的数值解中， （13.95），会产生各种数值误差，最重要的是空间误差

\subsection*{13.6.4 隐式数值滤波器}\phantomsection\addcontentsline{toc}{subsection}{13.6.4 隐式数值滤波器}
方程，它是由 nurneri-\textasciitilde{}a 解满足的偏微分方程。对应于 LES 动量方程的修正方程可以写成 10

\begin{sourceblock}
\includegraphics[page=660,pagebox=mediabox,trim={64.39bp 1617.80bp 64.39bp 302.96bp},clip,width=0.94\textwidth,max height=0.38\textheight]{Pope.pdf}
\end{sourceblock}
以便空间截断误差显示为额外的数值应力，
这取决于网格间距h。若空间离散化为 pthzri, eider
准确，则为 hP 量级。 （练习 13.46 中给出了一阶精确方案的简单示例，其中数值应力可以表示为附加粘性项，数值粘性与 h 成正比。）

对于 LES 中数值应力的作用存在不同的观点。最简单的（在前面的讨论中隐含地假设）是应该准确求解 LES 方程。也就是说，对于给定的过滤器宽度 A，网格间距 h 应选择足够小，使得与建模残余应力相比，数值应力 \textasciitilde{} h i+s 可以忽略不计
鲍里斯等人提出了相反的观点。 (1992) 的观点是，不应执行显式过滤，也不应使用显式残余应力模型 (zfi 0)。相反，使用适当的数值方法来尝试求解 U 的纳维-斯托克斯方程。由于网格不够精细，无法求解纳维-斯托克斯方程的解，因此显着的数值应力 zh+a 会上升。因此，滤波和残余应力建模是通过数值方法隐式执行的。

数值应力 zh+d 取决于所使用的数值方法的类型，因此这种选择至关重要。理想的方案对于 U(x, t)（相对于 h 的长波长）的良好分辨贡献是准确的（e 很小），同时它削弱了解析不良的较短波长 h 贡献（并防止混叠误差污染良好分辨模式）。能量通过数值耗散从已解析的运动中转化为

\begin{sourceblock}
\includegraphics[page=660,pagebox=mediabox,trim={64.39bp 378.95bp 64.39bp 1558.80bp},clip,width=0.94\textwidth,max height=0.38\textheight]{Pope.pdf}
\end{sourceblock}
埃克维斯等人。 (1992) 使用所谓的单调数值方法，并将其方法称为 MILES 单调集成大涡模拟。除了鲍里斯等人引用的作品。 (1992)，Tamura 和 Kuwahara (1989)、Knight 等人描述了使用这种方法的代表性 LES 计算。 （1998），以及奥孔戈和奈特（1998）。

这种方法有优点也有缺点，有批评者也有批评者。优点是（对于给定的网格大小）尽可能多的湍流运动由 LES 速度场 U(x,t) 明确表示，并且仅在必要时从 U 中移除能量。有人认为，只要有一种机制可以从最小的解析尺度中去除能量而不污染更大的尺度，如何去除能量的细节并不重要。 （这类似于用于证明使用在先验测试中表现不佳的残余应力模型的论点。）另一个优点是消除了开发和测试残余应力模型所需的时间和精力。

主要缺点是建模和数值是不可分割地耦合的。有时该方法被称为“无模型”，但应该意识到这是一个不充分的描述：对于给定的流，模拟结果取决于数值方法和所使用的网格。不可能细化网格以获得独立于网格的解决方案（缺少执行 DNS）。另一个缺点是没有可用于去滤波或用于其他亚格子尺度过程的模型中的亚格子尺度运动的表示或估计。

EX E R C I S E 13.46 考虑 Burgers 方程的数值解（Burgers

\begin{sourceblock}
\includegraphics[page=661,pagebox=mediabox,trim={64.39bp 806.90bp 64.39bp 1140.86bp},clip,width=0.94\textwidth,max height=0.38\textheight]{Pope.pdf}
\end{sourceblock}
对于正速度 u(x, t) > 0。粗略的有限差分法
使用，其中 x 处的空间导数近似为

\begin{sourceblock}
\includegraphics[page=661,pagebox=mediabox,trim={64.39bp 591.92bp 64.39bp 1310.83bp},clip,width=0.94\textwidth,max height=0.38\textheight]{Pope.pdf}
\end{sourceblock}
\begin{sourceblock}
\includegraphics[page=661,pagebox=mediabox,trim={64.39bp 481.94bp 64.39bp 1406.82bp},clip,width=0.94\textwidth,max height=0.38\textheight]{Pope.pdf}
\end{sourceblock}
其中 h 是网格间距。使用泰勒级数分析来确定
这些近似值中的前导截断误差。因此

表明方程已解，即

\begin{sourceblock}
\includegraphics[page=661,pagebox=mediabox,trim={64.39bp 243.97bp 64.39bp 1664.79bp},clip,width=0.94\textwidth,max height=0.38\textheight]{Pope.pdf}
\end{sourceblock}
等价于修改后的方程

\begin{sourceblock}
\includegraphics[page=661,pagebox=mediabox,trim={64.39bp 119.39bp 64.39bp 1813.77bp},clip,width=0.94\textwidth,max height=0.38\textheight]{Pope.pdf}
\end{sourceblock}
其中数值粘度是

\begin{sourceblock}
\includegraphics[page=662,pagebox=mediabox,trim={64.39bp 1728.78bp 64.39bp 219.97bp},clip,width=0.94\textwidth,max height=0.38\textheight]{Pope.pdf}
\end{sourceblock}
（因此前导截断误差相当于额外的
粘性项。）

\subsection*{13.6.5 近壁处理}\phantomsection\addcontentsline{toc}{subsection}{13.6.5 近壁处理}
在具有近壁分辨的 LES 中，无滑移和不渗透边界条件被应用于壁处的滤波速度。靠近墙的斯马戈林斯基长度尺度 ls 的正确行为是 ls y，这导致 y3 和 zT2 w y3 （参见第 601 页的练习 13.31）。相反，对于常数 Cs，Smagorinsky 模型会错误地生成 、v 和 T 的常数值； \textasciitilde{}O。克服这一缺陷的新方法是通过例如 van Driest 阻尼函数（方程（13.166））来阻尼 ls。更令人满意的是从动态模型确定cs。这产生了正确的近壁缩放，并且对于槽道流动而言相当令人满意的结果（参见egg.，Piomelli (19 93)）。

如果不解决粘性壁区域，则会节省大量计算量，并且事实上，这是高雷诺数下唯一可行的方法（参见第 13.4.5 节）。在 LES-NWM 的情况下，未解析运动的影响通常通过使用类似于湍流黏度和雷诺应力模型中使用的壁面函数的边界条件来建模（第 11 .7.6 节）。这些条件应用于墙壁 (y = 0) 或远离墙壁的第一个网格节点 (y = yp)。

不渗透条件 U2 = 0 应用于法向速度。 U 切向分量上的边界条件通过剪切应力的指定隐式应用：

\begin{sourceblock}
\includegraphics[page=662,pagebox=mediabox,trim={64.39bp 508.94bp 64.39bp 1430.82bp},clip,width=0.94\textwidth,max height=0.38\textheight]{Pope.pdf}
\end{sourceblock}
速度尺度 u* 是通过假设 U 瞬间满足对数定律（Grotzbach 1981，Mason 和 Callen 1986）或幂律（Werner 和 Wengle 1989）而获得的。 Piomelli、Ferziger 和 Moin (1989) 建议的修改是将 qx 处的剪切应力基于 x 下游一定距离 A 处的滤波速度 Ui(x A,, y,, z)。

规范中包含的假设以及剪切应力与切向速度一致的假设当然会受到批评。另外，抗渗条件U2 0 较小

\section*{13.7 讨论}\phantomsection\addcontentsline{toc}{section}{13.7 讨论}
代入已解决的雷诺应力 R!。 （方程（13.185））和平均残余雷诺应力以及 LES 中理想的 R！。是主要贡献。

然而，规范 (U2)p 0 迫使 (RL2) 为零，因此剪应力 (uiu2) 完全由建模的边界条件方程产生。 （13.299）。

对于高达 Re, 5,000 的河道流量，Balaras、Benocci 和 Piomelli (1995) 使用动态模型与此类近壁处理相结合，获得了令人满意的结果。然而，人们普遍认为，对于具有分离、重新附着、冲击等的复杂流动，这些壁处理是 LES-NWM 最不令人满意的方面。

13.7 讨论 本节讨论了第 8 章中描述的评估模型标准的 LES 方法；然后给出了湍流模拟和建模的最终观点。

\subsection*{13.7.1 LES 评估}\phantomsection\addcontentsline{toc}{subsection}{13.7.1 LES 评估}
描述水平 由于滤波后的速度场提供了含能运动的直接表示，因此 LES 中的描述水平对于大多数目的来说已经足够了。由于解析的字段已被过滤，因此在提取未过滤的单点统计数据时涉及一些近似，例如 (U) 和
（参见练习 13.32）；但实际上这并不是一个主要问题。 ( u w ) 与RANS相比，LES的优点是可以描述非定常、大尺度的湍流结构，因此可以用来研究结构上的非定常气动载荷和声音的产生等现象。

完整性 混合长度模型是不完整的（因为必须指定），而 k-E 和雷诺应力模型是完整的。 LES 完整吗？或者是否需要依赖于流量的规格？答案取决于滤波器宽度 A 的规格，通常作为位置和时间的函数。

将 LES 应用于复杂流时，通常的做法是指定 numerics1 解中要使用的网格，然后（隐式或显式）使滤波器宽度 A 与局部网格大小 h 成比例。如果
不完整。经验表明，如果含能运动解析不佳（即 VLES），LES 计算可能会依赖于网格。因此，为了完整起见，该方法需要先了解流程
可以适当地指定网格（和滤波器），以确保大部分湍流能量在整个域中得到解析。 LES 的空闲数值方法将包括自适应网格化，以自动确保网格以及滤波器在任何地方都足够精细以解析含能量的运动。

成本和易用性 LES 的计算成本需要仔细考虑，特别是对于壁面流动。对于均匀各向同性湍流，403 模态足以解析 80\% 的能量，即使在高雷诺数下也是如此（参见练习 13.34）。 （有时使用较少的模式（例如，3 23），因此在高波数或低波数时光谱的分辨率较低。）

对于自由剪切流，通常使用的模态或网格点的数量是可比较的。因此，此类流的 LES 是相当实用的；并且，与 DNS 形成鲜明对比的是，计算成本与雷诺数无关。另一方面，与雷诺应力模型相比，对于统计同质性的每个方向，LES 的成本要高出大约两个数量级。

对于壁界流的 LES，分辨率要求在本质上是不同的，因为重要的近壁运动的大小与粘性长度尺度 6 成比例（粘性长度尺度 6 随雷诺数的增加而减小）。 Chapman (1979) 估计，空气动力学应用中具有近壁分辨的 LES (LES-NWR) 的成本随着雷诺数 Re 1.8 的增加而增加；也就是说，雷诺数每增加 10 年，成本就会增加 60 倍。 （在第 599 页的练习 13.29 中获得了基本相同的结果。）因此，虽然 LES-N W R 已成功应用于低雷诺数和中雷诺数的边界层和槽道流，但其应用于高雷诺数流（例如飞机机翼、船体和大气边界层中的流）是不可行的。

另一种方法是采用近壁建模的 LES (LES-NWM)，其中未解析近壁含能运动。然后，根据实现的不同，计算成本要么与雷诺数无关，要么像 Re 中那样微弱增加（参见练习 13.29）。

由于 LES 即使对于各向同性湍流也需要依赖时间的三维计算，因此与雷诺计算相比，它的成本损失更大
流动（雷诺应力模型计算也是三维且不稳定的）。因此，例如，对于火花点火发动机中固有的三维和不稳定流动，LES 是一种有吸引力的方法（Haworth 和 Jansen 2000）。对于统计平稳的三维流（例如 Rodi 等人（1997）），LES 计算通常针对多个流传输时间执行，以便可以使用时间平均来估计平均值。相对于雷诺应力计算，这种对时间平均的需要增加了计算成本。

对于工程应用中通常遇到的三维流动，求解 LES 方程的数值任务与求解纳维-斯托克斯或湍流模型方程的数值任务基本相同。 4 因此，执行 LES 的代码的可用性与求解湍流模型的代码相当。许多商业 CFD 代码都包含“LES 选项”。然而，应该意识到，在通常使用的网格上，计算是 VLES，并且空间截断误差可能很大（与建模的残余应力相比）。

适用范围 对于恒定密度惰性流，LES 通常适用 - 尽管如上所述，计算成本决定了对于高雷诺数壁面边界流使用 LES-NWM 而不是 LES-NWR。与 RANS 相比，LES 提供的更高水平的描述增加了其对气动声学和与非定常湍流运动相关的其他现象的适用范围。

基于 LES 的方法也已被开发用于更复杂的湍流现象，例如高速可压缩流（参见 Hussaini (1998)）和反应流（例如 Cook 和 Riley (1994) 以及 Colucci 等人 (1998)）。所涉及的一些重要过程（例如分子混合和反应）可能主要发生在亚格子尺度上，因此需要统计建模。

准确性 与其他湍流方法一样，评估各种 LES 方法准确性的理想方法是比较它们在广泛的测试流中的性能。无关误差（参见第 342 页图 8.2）必须很小；测试流程应与模型开发中使用的流程不同；并且测试不应由模型开发人员执行。 LES 特有的一个附加要求是应针对变化的滤波器宽度进行计算

\subsection*{13.7.2 最终观点}\phantomsection\addcontentsline{toc}{subsection}{13.7.2 最终观点}
在一定范围内，可以评估计算对滤波器规格的敏感性。尚未进行满足这些要求的测试，因此 LES 准确性的评估存在相当大的不确定性。尽管如此，Fureby 等人还是提出了一些有用的比较结果。 (1997)（各向同性湍流），Vreman 等人。 (1997)（对于混合层），由 Rodi 等人提出。 (1997)（两个流过方形障碍物）。

Vreman 等人的混合层比较。 （1997）证明，通过使用例如动态混合模型（方程（13.274））可以获得相当准确的结果。这些和其他计算支持这样的观点，即 LES 对于自由剪切流可以相当准确。

对于中等雷诺数的简单边界层型流动，有许多基于动态模型的 LES-NWR 计算支持这些流动的类似结论（例如，P iomelli (1993)）。然而，应该承认，不同流的数量（用于级别操作和测试）并不大。

对于更复杂和更高雷诺数的流动，有限的证据表明 LES-NWM 的准确性不太确定。例如，基于对越过障碍物的流的各种 LES-NWM 计算的比较研究，Rodi 等人。 （1997）得出的结论是“结果对网格和数值方法很敏感”，并且所使用的壁处理“在分离流中使用并不可靠”。

13.7.2 研究观点 在湍流背景下，研究的一个持续的主要挑战是开发方法（使用不断增加的可用计算机能力）来计算工程、大气科学和其他领域实际相关的流动和湍流特性。这是本书第一部分的主题。

湍流问题种类繁多，几何复杂性各异，许多涉及额外的物理和化学过程，需要不同程度的描述和准确性。因此，拥有具有各种属性的广泛方法是很有价值的，软管计算成本变化多达 lo6 倍。第 10 章至第 13 章中描述的每种方法都有其自己的位置：没有一种方法可以取代所有其他方法。

此外，这些方法之间存在有用的协同作用。虽然 DNS 的计算成本使其无法用于高雷诺数
其他方法。 RANS 建模预计将有助于解决 LES 中高雷诺数流中近壁区域建模的主要突出问题。统计方法（例如 PDF 方法）可以与 LES 结合使用，用于处理湍流无功流和其他在亚格子尺度上发生重要过程的流。

对于某些湍流问题，现有方法的准确度可以接受，并且计算成本也可以接受。对于许多其他人来说，需要取得重大进展。因此，在未来一段时间内，湍流仍将是一个有价值且具有挑战性的研究领域。

\part{附录}
\appendix
\titleformat{\chapter}[display]{\bfseries\Huge\centering}{附录\thechapter}{0.75em}{}
\chapter{笛卡尔张量}
从向量微积分中我们熟悉标量和向量。标量具有单个值，该值在任何坐标系中都相同。向量具有大小和方向，并且（在任何给定的坐标系中）它具有三个分量。使用笛卡尔张量，我们不仅可以表示标量和向量，还可以表示具有更多方向的量。具体来说，N 阶张量 ( N 0) 有 N 个与其关联的方向，并且（在给定的笛卡尔坐标系中）它有 3N 个分量。零阶张量是标量，一阶张量是向量。在定义高阶张量之前，我们简要回顾一下笛卡尔坐标中向量的表示。

\subsection*{A.l 笛卡尔坐标和向量}\phantomsection\addcontentsline{toc}{subsection}{A.l 笛卡尔坐标和向量}
流体流动（以及经典力学中的其他现象）发生在三维欧几里得物理空间中。如图 A.1 所示，令 E 表示物理空间中的笛卡尔坐标系。这是由原点 0 的位置和三个相互垂直的轴的方向定义的。三个坐标方向上的单位向量用el、ez和e3表示。我们写 ei 来指代其中任何一个，并理解 subx i（或任何其他后缀）采用值 1、2 或 3。

单位向量的基本属性可以用以下术语简洁地表达：

克罗内克三角洲 dij 的 ei。这是由以下定义的
我们在这里考虑三维空间。张量扩展到不同维度的空间
很简单。

图A.1。 E 坐标系草图，显示原点 0、三个
正交基向量 ei 和位置 x = xlel xzez 的一般点 P
或者，用矩阵表示法，
从它们的构造中可以明显看出，向量 ei 具有正交性：

对于，如果 i = j，则点积为 1（因为它们是单位向量），而对于 i j，点积为零（因为向量是正交的）。因此e1、e 2 和e3 形成一组正交基向量。

参见图A.1，'是位置向量，给出一般点P相对于原点0的位置。这个向量可以写成

\begin{sourceblock}
\includegraphics[page=672,pagebox=mediabox,trim={64.39bp 519.93bp 64.39bp 1427.82bp},clip,width=0.94\textwidth,max height=0.38\textheight]{Pope.pdf}
\end{sourceblock}
其中 x1、x2 和 x3 是 E 坐标系中的分量。

求和约定 通过使用爱因斯坦求和约定，笛卡尔张量的表示法大大简化。根据此约定，如果后缀重复（例如，eixi 中的后缀 i），则隐含对后缀的所有三个值（i= 1 2 和 3）求和。因此，通过这个隐含求和，方程。 (A.4) 写为

\begin{sourceblock}
\includegraphics[page=672,pagebox=mediabox,trim={64.39bp 119.39bp 64.39bp 1843.77bp},clip,width=0.94\textwidth,max height=0.38\textheight]{Pope.pdf}
\end{sourceblock}
请注意，重复或虚拟后缀的符号是不相关的（即 eiXi ejxj）。

在张量方程中， - 和
符号分隔不同的表达方式。在给定的表达式中，不重复的后缀称为自由后缀。与求和约定相关的一条重要规则是后缀在表达式中出现的次数不能超过两次。

通过显式展开求和，显示替换规则
获取结果

\begin{sourceblock}
\includegraphics[page=673,pagebox=mediabox,trim={64.39bp 1259.84bp 64.39bp 686.91bp},clip,width=0.94\textwidth,max height=0.38\textheight]{Pope.pdf}
\end{sourceblock}
证明向量的分量可以通过以下方式获得

令 u 和 v 为向量，其分量为 (u1,u2,u3) 和 (vl,v2,v3)。

证明点积为

\begin{sourceblock}
\includegraphics[page=673,pagebox=mediabox,trim={64.39bp 880.89bp 64.39bp 1066.87bp},clip,width=0.94\textwidth,max height=0.38\textheight]{Pope.pdf}
\end{sourceblock}
坐标变换 张量是根据其变换属性来定义的。我们现在介绍第二个笛卡尔坐标系（由 表示并确定向量分量的变换规则。

E系统是由E通过轴的旋转和反射的任意组合获得的；图 A.2 给出了一个例子。需要注意的是，如果反射次数为奇数，则E和E的旋向性会不同。即，如果E是右手系，并且E是通过反射一个轴得到的，则E是左手系。在笛卡尔张量中，右手系统和左手系统都是有效的，并且没有必要区分它们。

令el、Z2 和Z3 为E 系统的正交基向量。由点积定义的九个量

\begin{sourceblock}
\includegraphics[page=673,pagebox=mediabox,trim={64.39bp 154.98bp 64.39bp 1788.77bp},clip,width=0.94\textwidth,max height=0.38\textheight]{Pope.pdf}
\end{sourceblock}
图；。 E（实线）和E（虚线）坐标In的草图
在这个具体例子中，E是通过Z3轴的反射从E获得的，并且
平面内的旋转。

E 系统和 E 系统中的 j 轴。此外，aij 是E 系统中的第i 个分量（见式（A.9））。从这个解释我们得到一个
ek 方向余弦的基本属性： 和 的点积
Zj 为（见式（A.10））

因此，由于向量是正交的，我们得到

\begin{sourceblock}
\includegraphics[page=674,pagebox=mediabox,trim={64.39bp 674.91bp 64.39bp 1269.84bp},clip,width=0.94\textwidth,max height=0.38\textheight]{Pope.pdf}
\end{sourceblock}
类似地可以证明

\begin{sourceblock}
\includegraphics[page=674,pagebox=mediabox,trim={64.39bp 535.93bp 64.39bp 1407.82bp},clip,width=0.94\textwidth,max height=0.38\textheight]{Pope.pdf}
\end{sourceblock}
E 和 E 系统中的向量相同，但其分量 (xl, x2, x3) 和 (XI,, 不同：

\begin{sourceblock}
\includegraphics[page=674,pagebox=mediabox,trim={64.39bp 354.95bp 64.39bp 1590.80bp},clip,width=0.94\textwidth,max height=0.38\textheight]{Pope.pdf}
\end{sourceblock}
通过将该方程与 ek 和 ek 进行点积来获得变换规则：

\begin{sourceblock}
\includegraphics[page=674,pagebox=mediabox,trim={64.39bp 189.98bp 64.39bp 1758.78bp},clip,width=0.94\textwidth,max height=0.38\textheight]{Pope.pdf}
\end{sourceblock}
\begin{sourceblock}
\includegraphics[page=674,pagebox=mediabox,trim={64.39bp 119.39bp 64.39bp 1837.77bp},clip,width=0.94\textwidth,max height=0.38\textheight]{Pope.pdf}
\end{sourceblock}
\section*{A.2 笛卡尔张量的定义}\phantomsection\addcontentsline{toc}{section}{A.2 笛卡尔张量的定义}
导数 x 的分量是 E 系中的坐标。因此，我们可以对物理空间中定义的 xi 函数进行微分。以下练习中包含一些有用的结果。

得到以下结果：

A.2 笛卡尔张量的定义 零阶张量是标量。它有 3O 1 个分量，在每个坐标系中具有相同的值。例如密度、温度和压力等物理量以及向量的点积。

一阶张量 u，
有 3l 3 个分量，其变换为
（参见方程式（A.17））。正如接下来的两个练习中所说明的，确定一个量是否为等式1的形式的方法。 （A.19）是一阶张量，就是确定它的变换规则，然后将该规则与式（A.19）进行比较。 （A.20）。

设 X(t) d 和 U(t) dX(t)/dt 为 a 的位置和速度
颗粒。确定组件之间的变换规则
Uj(t) 和 Uj(t)，因此表明速度是一阶张量。

令 \$(x) 为标量场。确定之间的变换规则
d\$/axi 和 d\$/dxj，因此表明标量梯度

\begin{sourceblock}
\includegraphics[page=675,pagebox=mediabox,trim={64.39bp 205.97bp 64.39bp 1740.78bp},clip,width=0.94\textwidth,max height=0.38\textheight]{Pope.pdf}
\end{sourceblock}
是一阶张量。

二阶张量 A 二阶张量，b，

\begin{sourceblock}
\includegraphics[page=676,pagebox=mediabox,trim={64.39bp 1679.79bp 64.39bp 268.97bp},clip,width=0.94\textwidth,max height=0.38\textheight]{Pope.pdf}
\end{sourceblock}
有 32 9 个分量，其变换为

\begin{sourceblock}
\includegraphics[page=676,pagebox=mediabox,trim={64.39bp 1555.80bp 64.39bp 388.95bp},clip,width=0.94\textwidth,max height=0.38\textheight]{Pope.pdf}
\end{sourceblock}
关于等式 1 中的表达式 eiejbij，需要进行一些观察。 （A.22）。 ei 和 e 之间没有运算符（例如点），并且数量 e ne m 在标准向量表示法中没有等效项。与往常一样，用于虚拟后缀的符号是无关紧要的，即

\begin{sourceblock}
\includegraphics[page=676,pagebox=mediabox,trim={64.39bp 1301.84bp 64.39bp 646.92bp},clip,width=0.94\textwidth,max height=0.38\textheight]{Pope.pdf}
\end{sourceblock}
但单位向量的排序很重要，即

\begin{sourceblock}
\includegraphics[page=676,pagebox=mediabox,trim={64.39bp 1172.85bp 64.39bp 771.90bp},clip,width=0.94\textwidth,max height=0.38\textheight]{Pope.pdf}
\end{sourceblock}
事实上，等式的右侧。 (A.25) 定义了一个不同的二阶张量，即 b 的转置：

\begin{sourceblock}
\includegraphics[page=676,pagebox=mediabox,trim={64.39bp 1000.87bp 64.39bp 930.88bp},clip,width=0.94\textwidth,max height=0.38\textheight]{Pope.pdf}
\end{sourceblock}
确定一个量是否符合方程（1）的形式。 （A.22）是二阶张量，需要确定变换规则并与式（A.22）进行比较。 （A.23）。例如，考虑由下式定义的量 cij 和 Zke

\begin{sourceblock}
\includegraphics[page=676,pagebox=mediabox,trim={64.39bp 792.90bp 64.39bp 1152.85bp},clip,width=0.94\textwidth,max height=0.38\textheight]{Pope.pdf}
\end{sourceblock}
其中 u 和 v 是向量。这些是二阶张量的组成部分吗？为了确定答案，我们需要确定 zke 的变换规则，这是我们从 Zik 和 Be 的变换规则获得的（式（A-19））：

\begin{sourceblock}
\includegraphics[page=676,pagebox=mediabox,trim={64.39bp 528.93bp 64.39bp 1368.83bp},clip,width=0.94\textwidth,max height=0.38\textheight]{Pope.pdf}
\end{sourceblock}
这确实是变换规则 Eq. (A.23)：因此 eiejcij ekeecke 是二阶张量。

证明以下是二阶张量的组成部分：

其中 \$(x) 是标量场，u(x) 是向量场。

\section*{A.3 张量运算}\phantomsection\addcontentsline{toc}{section}{A.3 张量运算}
高阶张量 扩展到高阶张量很简单。三阶张量

\begin{sourceblock}
\includegraphics[page=677,pagebox=mediabox,trim={64.39bp 1672.79bp 64.39bp 267.97bp},clip,width=0.94\textwidth,max height=0.38\textheight]{Pope.pdf}
\end{sourceblock}
有 33 27 个分量，其变换为

\begin{sourceblock}
\includegraphics[page=677,pagebox=mediabox,trim={64.39bp 1544.81bp 64.39bp 394.95bp},clip,width=0.94\textwidth,max height=0.38\textheight]{Pope.pdf}
\end{sourceblock}
N 阶张量 ( N 0) 的形式为

\begin{sourceblock}
\includegraphics[page=677,pagebox=mediabox,trim={64.39bp 1416.82bp 64.39bp 523.93bp},clip,width=0.94\textwidth,max height=0.38\textheight]{Pope.pdf}
\end{sourceblock}
并有 3N 个分量，其变换为

\begin{sourceblock}
\includegraphics[page=677,pagebox=mediabox,trim={64.39bp 1278.84bp 64.39bp 649.92bp},clip,width=0.94\textwidth,max height=0.38\textheight]{Pope.pdf}
\end{sourceblock}
A.3 张量运算 我们现在描述可以用张量执行的运算。如果操作数是张量，那么结果也是张量。留给读者来证明这一事实。

加法 两个同阶张量可以相加或相减。例如，如果 b 和 c 是二阶张量，则它们的和 s 由下式给出

\begin{sourceblock}
\includegraphics[page=677,pagebox=mediabox,trim={64.39bp 688.91bp 64.39bp 1203.85bp},clip,width=0.94\textwidth,max height=0.38\textheight]{Pope.pdf}
\end{sourceblock}
因此，就组件而言，我们简单地说，
张量积 例如，一阶张量 u 和二阶张量 b 的张量积是三阶张量 d：

\begin{sourceblock}
\includegraphics[page=677,pagebox=mediabox,trim={64.39bp 212.97bp 64.39bp 1621.80bp},clip,width=0.94\textwidth,max height=0.38\textheight]{Pope.pdf}
\end{sourceblock}
因此，就 ol 组件而言
一般来说，N阶张量和M阶张量的张量积是N阶张量。

标量（ e 是零阶张量）和 N 阶张量的乘积是 N 阶张量。特别是（将标量设为 -1）如果 bij 是张量，则 -bij 也是张量，因此减法是有效的运算，例如，

\begin{sourceblock}
\includegraphics[page=678,pagebox=mediabox,trim={64.39bp 1549.80bp 64.39bp 393.95bp},clip,width=0.94\textwidth,max height=0.38\textheight]{Pope.pdf}
\end{sourceblock}
方程 (A.34)、(A.36) 和 (A.37) 是克伦分量形式的张量方程的示例。在这些方程中，各个表达式由 , 或 = 分隔。在任何张量方程中，每个表达式都具有相同的自由后缀，例如方程 1 中的 i、j 和 k。 （A.36）。后缀的顺序很重要，但不需要在每个表达式中都相同。因此，例如，

\begin{sourceblock}
\includegraphics[page=678,pagebox=mediabox,trim={64.39bp 1208.85bp 64.39bp 677.91bp},clip,width=0.94\textwidth,max height=0.38\textheight]{Pope.pdf}
\end{sourceblock}
\begin{sourceblock}
\includegraphics[page=678,pagebox=mediabox,trim={64.39bp 1133.86bp 64.39bp 806.90bp},clip,width=0.94\textwidth,max height=0.38\textheight]{Pope.pdf}
\end{sourceblock}
都是有效的张量方程，但它们定义了不同的张量 d、d' 和 d"。

收缩 给定三阶张量的分量 dijk，考虑
wl、WZa 和 nd w3 是一阶张量的分量吗？在E系统中我们有

\begin{sourceblock}
\includegraphics[page=678,pagebox=mediabox,trim={64.39bp 635.92bp 64.39bp 1302.84bp},clip,width=0.94\textwidth,max height=0.38\textheight]{Pope.pdf}
\end{sourceblock}
现在在等式中。 （A.42），设置 rn 我们得到
这样我们就得到

\begin{sourceblock}
\includegraphics[page=678,pagebox=mediabox,trim={64.39bp 201.97bp 64.39bp 1743.78bp},clip,width=0.94\textwidth,max height=0.38\textheight]{Pope.pdf}
\end{sourceblock}
这确实是一阶张量的变换规则。

一般来说，当 N 阶张量 (N 2) 的任意两个索引收缩时，结果是 N 2 阶张量。请注意，diik、dijia 和 nd dijj 都是 dijk 的有效（但不同）收缩。

内积如果 b 和 c 分别是二阶张量和三阶张量，则它们的内积是三阶张量 d：

\begin{sourceblock}
\includegraphics[page=679,pagebox=mediabox,trim={64.39bp 1280.84bp 64.39bp 606.92bp},clip,width=0.94\textwidth,max height=0.38\textheight]{Pope.pdf}
\end{sourceblock}
'l'hus，就组成部分而言，
一般来说，对 N 阶张量和 M 阶张量 (N, M 2 1) 求内积会得到 N M - 2 阶张量。最熟悉的内积当然是两个向量的点积。

除法 没有与除法对应的张量运算。 （但是，正如附录 B 中所讨论的，二阶张量可能有逆张量。）

梯度 梯度是通过应用梯度算子获得的

\begin{sourceblock}
\includegraphics[page=679,pagebox=mediabox,trim={64.39bp 544.93bp 64.39bp 1378.83bp},clip,width=0.94\textwidth,max height=0.38\textheight]{Pope.pdf}
\end{sourceblock}
到张量场。例如，如果 b(x) 是二阶张量场，

\begin{sourceblock}
\includegraphics[page=679,pagebox=mediabox,trim={64.39bp 405.95bp 64.39bp 1536.81bp},clip,width=0.94\textwidth,max height=0.38\textheight]{Pope.pdf}
\end{sourceblock}
那么 b 的梯度就是三阶张量 h，由下式给出

\begin{sourceblock}
\includegraphics[page=679,pagebox=mediabox,trim={64.39bp 119.39bp 64.39bp 1826.77bp},clip,width=0.94\textwidth,max height=0.38\textheight]{Pope.pdf}
\end{sourceblock}
或者说，从成分来看，
有时有用的速记符号是将 db -m /dxk 写为 bij,k。

从式（1）可以看出。 (A.50) b 的梯度类似于 V 和 b 的张量积。一般来说，N 阶张量 (N 2 0) 的梯度是 N 1 阶张量。例如，重复应用梯度算子 V 会依次产生更高的导数 d 2 b i j / ( \textasciitilde{} xakx e)。

张量 b 的散度是 V 和 b 的内积：

\begin{sourceblock}
\includegraphics[page=680,pagebox=mediabox,trim={64.39bp 1092.86bp 64.39bp 722.91bp},clip,width=0.94\textwidth,max height=0.38\textheight]{Pope.pdf}
\end{sourceblock}
一般来说，N 阶张量 (N 1) 的散度是 N 1 阶张量。特别地，向量的散度是标量。

拉普拉斯算子 拉普拉斯算子是标量运算符

\begin{sourceblock}
\includegraphics[page=680,pagebox=mediabox,trim={64.39bp 616.92bp 64.39bp 1305.84bp},clip,width=0.94\textwidth,max height=0.38\textheight]{Pope.pdf}
\end{sourceblock}
泰勒级数 如果 b(x) 是平滑二阶张量场，则可以从泰勒级数的 b 及其在 y 处的导数获得 b 在位置 y r 处的值。就组成部分而言，该系列是

\begin{sourceblock}
\includegraphics[page=680,pagebox=mediabox,trim={64.39bp 162.98bp 64.39bp 1725.78bp},clip,width=0.94\textwidth,max height=0.38\textheight]{Pope.pdf}
\end{sourceblock}
高斯定理这适用于任何阶的（一旦连续可微的）张量场。我们以二阶张量 b(x) 为例。

令 A 为包围体积 V 的分段平滑闭合可定向曲面，并令 表示 A 上的向外指向的单位法线（参见第 85 页的图 4.1）。然后，

\begin{sourceblock}
\includegraphics[page=681,pagebox=mediabox,trim={64.39bp 1436.82bp 64.39bp 437.94bp},clip,width=0.94\textwidth,max height=0.38\textheight]{Pope.pdf}
\end{sourceblock}
其中 dA 和 dV 是面积和体积的元素。如果 i 和 k 收缩，结果就是散度定理。

符号 涉及张量的方程可以用直接符号书写，例如 b c，也可以用 su@x 符号书写，例如 b,，我们遵循流体力学的传统
cij。对二阶或更高阶张量使用后缀表示法，以及对一阶张量互换使用这两种表示法（例如， V u a \textasciitilde{} \textasciitilde{} / a \textasciitilde{} \textasciitilde{} ）。

在使用后缀表示法时，我们遵循引用“张量 bij”的正常做法。然而，应该理解的是，更正确地讲，bij 是张量 b 在 E 坐标系中的 i-j 分量。

u(x) 是向量场，\$(x) 和 0(x) 是标量场，
以笛卡尔张量后缀表示法写入

(c) 由 (a) 和 (b) 推导出的关于 0 的方程。

量 Aij、Bijkl 和 Cijkea 是给定的张量（不是张量
字段）。在 ( a - ( e ) 的每一个中，使用这些张量中的每一个来完成
指定数量的有效方程：

(b) 一阶张量：v j
(c) 二阶张量：Tke
(d) 三阶张量：F,,,
(e) 四阶张量：I/,jke
（例如，(a) 的可能答案是 - - .A ii uiBijj Ciijj）。

令 u(x, t ) 为速度场，令 n 为给定的单位向量。使用
(a) 流体加速度，
(b) n 方向上的加速度分量，以及
(c) 垂直于 n 的平面内的速度分量。

考虑两个移动点 A 和 B，其位置为 和
其速度为 d ( t ) 和 uB(t)。令 s(t) 表示
A 和 B 之间的距离，设 n(t) 为单位向量
xA(t - ) xB(t)U.s 的方向 笛卡尔张量后缀表示法，
获得表达式

(a) s 以 XA 和 XB 表示
(b) ni 表示 Xt, X!和，
(c) 用 ni、uA 和 U 表示？和
以 ni、uf、uB 和 s 表示。

\section*{A.4 向量叉积}\phantomsection\addcontentsline{toc}{section}{A.4 向量叉积}
在开发笛卡尔张量时，我们引入了许多在向量微积分标量、向量、点积、梯度和散度中具有直接等价物的量。然而，对于笛卡尔张量，没有叉积或旋度等价物。围绕这一观察的考虑有些微妙，而且常常不被重视或被理解。在湍流建模的背景下，这是一个重要的问题，我们将在本节中进行说明和澄清。

叉积可以通过使用由下式定义的交替符号（或 Levi-Civita 符号）以后缀表示法编写

\begin{sourceblock}
\includegraphics[page=682,pagebox=mediabox,trim={64.39bp 456.94bp 64.39bp 1404.82bp},clip,width=0.94\textwidth,max height=0.38\textheight]{Pope.pdf}
\end{sourceblock}
循环顺序为123、231、3 12；反循环顺序为 321、132 和 213；否则两个或多个后缀是相同的。两个向量 u 和 v 的叉积为

\begin{sourceblock}
\includegraphics[page=682,pagebox=mediabox,trim={64.39bp 119.39bp 64.39bp 1719.78bp},clip,width=0.94\textwidth,max height=0.38\textheight]{Pope.pdf}
\end{sourceblock}
因此组件是

\begin{sourceblock}
\includegraphics[page=683,pagebox=mediabox,trim={64.39bp 1735.78bp 64.39bp 212.97bp},clip,width=0.94\textwidth,max height=0.38\textheight]{Pope.pdf}
\end{sourceblock}
并且，例如，

\begin{sourceblock}
\includegraphics[page=683,pagebox=mediabox,trim={64.39bp 1550.80bp 64.39bp 341.96bp},clip,width=0.94\textwidth,max height=0.38\textheight]{Pope.pdf}
\end{sourceblock}
类似地，如果 U(x) 是矢量场（例如速度），则其旋度 w(x) V x U ( e . v orticity) 为

\begin{sourceblock}
\includegraphics[page=683,pagebox=mediabox,trim={64.39bp 1364.83bp 64.39bp 543.93bp},clip,width=0.94\textwidth,max height=0.38\textheight]{Pope.pdf}
\end{sourceblock}
带组件

Cr)i Eijk
r 和 o 是一阶张量吗？从等式可以看出。 (A.58) 和 (A.61) 答案取决于是否是三阶张量。一个
Eijk 简单的例子足以表明它不是一个张量。设E系统为
通过各轴的反射从系统中获得Eijk，即-ei。的
Zi = 方向余弦则为

\begin{sourceblock}
\includegraphics[page=683,pagebox=mediabox,trim={64.39bp 723.91bp 64.39bp 1163.85bp},clip,width=0.94\textwidth,max height=0.38\textheight]{Pope.pdf}
\end{sourceblock}
由于负号，它与三阶张量的变换规则不同。 （可以证明，如果存在偶数个反射（ e 、 0 或 2 ），则变换为张量，但是，如果存在
Eijk 奇数（ i 、 1 或 3），根据式（1）进行变换。 （A.63）。）

然后，我们得出结论，叉积 r = u x v 和旋度 w = V x U 不是一阶张量：它们更合适地称为伪向量。

惯用手的概念是围绕向量和伪向量问题的核心。在矢量微积分中，我们坚持使用固定惯用手坐标系，通常是右手坐标系，如图 A.1 所示。物体、运动和流动也可以表现出旋手性，主要的例子是右手和左手螺丝。

图 A.3 显示（在右手坐标系中）一个旋转箭头远离观察者朝镜子移动。箭头的速度为

\begin{sourceblock}
\includegraphics[page=684,pagebox=mediabox,trim={64.39bp 1310.84bp 64.39bp 109.44bp},clip,width=\textwidth,max height=0.72\textheight]{Pope.pdf}
\par\vspace{0.45\baselineskip}
\small 图 A.3。右手坐标系中的草图，其中旋转箭头（左下）朝镜子移动，其图像在镜子中（右上）。速度矢量改变方向，但旋转伪矢量 @ 不改变。
\end{sourceblock}
\begin{sourceblock}
\includegraphics[page=684,pagebox=mediabox,trim={64.39bp 638.92bp 64.39bp 781.90bp},clip,width=\textwidth,max height=0.72\textheight]{Pope.pdf}
\par\vspace{0.45\baselineskip}
\small 图 A.4。左手坐标系中的草图，具有与图 A.3 中所示相同的旋转箭头和镜子。注意旋转伪向量 Qa 的方向。
\end{sourceblock}
u 和 w 都平行于箭头的轴，从尾部指向头部。乘积 u w 为正，对应于右旋螺杆的运动。图 A.3 中还显示了箭头在镜子中的反射。反射图像是一个箭头以速度 ur - -u 向观察者移动，并以速率 w、w 旋转。这说明了矢量和伪矢量在反射下的行为之间的差异。请注意，对于图像箭头，u、ori 为负，相应地运动为左旋螺杆。

图A.5。螺旋运动草图 (a) 具有正螺旋度 (rr w > 0) 和 (b) 具有
负螺旋度 (rr w < 0)。

w 的方向相反，箭头的右旋螺旋运动现在对应于 u w 的负值。

这些都是重要考虑因素的原因是，力学定律和牛顿流体的特性都不偏向右手或左手运动2为了说明相反的行为，我们考虑以下方程

\begin{sourceblock}
\includegraphics[page=685,pagebox=mediabox,trim={64.39bp 847.89bp 64.39bp 1099.86bp},clip,width=0.94\textwidth,max height=0.38\textheight]{Pope.pdf}
\end{sourceblock}
旨在模拟流体行为的某些方面（u(x, t) 和 w(x, t) 是速度场和涡度场）。除了缺乏物理基础之外，这个方程还是一个糟糕的模型，因为它在维度上不共交，并且违反了伽利略不变性原理（参见第 2.9 节）。与当前讨论相关的是，该方程错误地将隐含运动偏向于特定的惯用手。为了看到这一点，选择一个局部右手坐标系，使得（在观察的时刻和位置）速度位于 el 方向。然后，根据方程。 (A.64)，w 的组成部分演化为

\begin{sourceblock}
\includegraphics[page=685,pagebox=mediabox,trim={64.39bp 182.98bp 64.39bp 1717.78bp},clip,width=0.94\textwidth,max height=0.38\textheight]{Pope.pdf}
\end{sourceblock}
{\small 3 流动可能会偏向某一方向或另一方向，但这是由初始和边界引起的\par}
1t可由式(1)看出。 (A.65) 该模型意味着 w 在 u 方向上的分量总是增加，因此偏向于正值；或者，等效地，螺旋度 u o 偏置为对应于右旋螺旋运动的正值，如图 A.5(a) 所示。如果改用左手坐标系，则方程： (A.64) 意味着偏向于左手螺旋运动。

由于力学定律和牛顿流体的特性在旋向性方面都是无偏的，因此有偏的方程（例如方程式） (A.64) 是根本错误的。在构建描述湍流的模型方程时，我们希望确保不存在旋向偏差。这可以通过要求用笛卡尔张量表示法编写方程来简单地实现。因此，在右手坐标系和左手坐标系中隐含的运动是相同的；因此，由于这种对称性，就不会有偏差。笛卡尔张量的旋向对称性源于其所需的变换属性中包含反射。

结论总结如下。

(i) 交替符号不是三阶张量，且 v = u x v
和 w = V x u 不是一阶张量。它们是伪向量
在（奇数）的反射下不变换为张量
坐标轴。

(ii) 方程式十所隐含的物理行为（合法地）

笛卡尔张量对于旋手性是无偏的，并且是
与坐标系的惯用手无关。

a.. 与 (ii) 相比，用向量表示法编写的方程涉及
伪向量，或等效地用后缀表示法编写的方程
涉及向量或伪向量，可能会错误地暗示不对称性
用手习惯 身体行为取决于
坐标系的惯用手。

虽然（考虑到上面的（iii））和伪向量是要避免的
由于建立了描述流体行为的方程，因此它们可以合法且有效地应用于其他情况。特别是，涡度和涡度方程是我们理解流体力学中许多现象的核心。

还可以观察到，包含偶数次交替符号的表达式像张量一样变换，因为减号抵消了。因此，这样的表达式可以重新表示为

\section*{A.5 笛卡尔张量后缀表示法总结}\phantomsection\addcontentsline{toc}{section}{A.5 笛卡尔张量后缀表示法总结}
显示以下内容：

\begin{sourceblock}
\includegraphics[page=687,pagebox=mediabox,trim={64.39bp 1680.79bp 64.39bp 266.97bp},clip,width=0.94\textwidth,max height=0.38\textheight]{Pope.pdf}
\end{sourceblock}
\begin{sourceblock}
\includegraphics[page=687,pagebox=mediabox,trim={64.39bp 1604.80bp 64.39bp 340.96bp},clip,width=0.94\textwidth,max height=0.38\textheight]{Pope.pdf}
\end{sourceblock}
\begin{sourceblock}
\includegraphics[page=687,pagebox=mediabox,trim={64.39bp 1526.81bp 64.39bp 416.95bp},clip,width=0.94\textwidth,max height=0.38\textheight]{Pope.pdf}
\end{sourceblock}
\begin{sourceblock}
\includegraphics[page=687,pagebox=mediabox,trim={64.39bp 1450.82bp 64.39bp 494.94bp},clip,width=0.94\textwidth,max height=0.38\textheight]{Pope.pdf}
\end{sourceblock}
使用方程式。 （A.61）表明

\begin{sourceblock}
\includegraphics[page=687,pagebox=mediabox,trim={64.39bp 1286.84bp 64.39bp 634.92bp},clip,width=0.94\textwidth,max height=0.38\textheight]{Pope.pdf}
\end{sourceblock}
哪里是

旋转率张量

\begin{sourceblock}
\includegraphics[page=687,pagebox=mediabox,trim={64.39bp 1123.86bp 64.39bp 773.90bp},clip,width=0.94\textwidth,max height=0.38\textheight]{Pope.pdf}
\end{sourceblock}
并为 wl 和 co3 编写显式表达式。

\begin{sourceblock}
\includegraphics[page=687,pagebox=mediabox,trim={64.39bp 962.88bp 64.39bp 924.88bp},clip,width=0.94\textwidth,max height=0.38\textheight]{Pope.pdf}
\end{sourceblock}
并重写这个方程，将每一边明确表示为 3 x 3
验证一下

\begin{sourceblock}
\includegraphics[page=687,pagebox=mediabox,trim={64.39bp 690.91bp 64.39bp 1199.85bp},clip,width=0.94\textwidth,max height=0.38\textheight]{Pope.pdf}
\end{sourceblock}
A.5 笛卡尔张量后缀表示法概述 使用后缀表示法的笛卡尔张量涉及的定义、规则和运算可以概括如下。

1. 在张量方程中， 之间用 - 或 = 分隔。

张量表达式

\begin{sourceblock}
\includegraphics[page=687,pagebox=mediabox,trim={64.39bp 244.97bp 64.39bp 1644.79bp},clip,width=0.94\textwidth,max height=0.38\textheight]{Pope.pdf}
\end{sourceblock}
2. 在张量表达式中，出现的后缀.once 是 a（例如，
自由后缀

方程中的 i 和 j （A.75））。

或虚拟后缀（例如，dijkk 中的后缀 k）。用于表示的符号
虚拟后缀并不重要，即 diikk dijpp*
4. 求和约定：重复的后缀意味着求和，即

\begin{sourceblock}
\includegraphics[page=688,pagebox=mediabox,trim={64.39bp 1589.80bp 64.39bp 292.96bp},clip,width=0.94\textwidth,max height=0.38\textheight]{Pope.pdf}
\end{sourceblock}
5. 在张量表达式中，后缀不能出现超过两次。对于
例如，表达式无效。

6. 具有 N 个自由后缀的张量表达式是（或者更正确地说，表示

表示 N 阶张量的分量。例如，每个

表达式为(A.75) 是二阶张量。

7. 张量方程中的每个表达式必须是相同的张量
顺序，具有相同的自由后缀（不一定是相同的顺序）。

方程（A.75）有效，而 bij dijk 和 bij 都是
8. 克罗内克三角洲的定义为

\begin{sourceblock}
\includegraphics[page=688,pagebox=mediabox,trim={64.39bp 1045.87bp 64.39bp 859.89bp},clip,width=0.94\textwidth,max height=0.38\textheight]{Pope.pdf}
\end{sourceblock}
它是一个二阶张量。请注意 3.

9. 方程中的交替符号。 (A.56) 不是张量。

10. 加法，例如 bijk dikj。每个张量必须具有相同的阶数
具有相同的自由后缀。

11. N阶张量和5阶张量的张量积
是 (N M) 阶张量，例如 bijkem cijdkem+
12. N阶张量（AT 2 2）可以通过改变两个自由来收缩
后缀变成重复后缀。结果是 N - 2 阶张量。

dijk 的不同缩写是 diik、d iji 和 dijj+
13. N阶张量和5阶张量的内积
(N 2 1, M 2 1) 是 N M - 2 阶张量：例如，f = cijdjne-
14. 替换规则是与克罗内克的内积
Delta 是，例如，

\begin{sourceblock}
\includegraphics[page=688,pagebox=mediabox,trim={64.39bp 368.95bp 64.39bp 1531.81bp},clip,width=0.94\textwidth,max height=0.38\textheight]{Pope.pdf}
\end{sourceblock}
16. 张量的梯度是高一阶的张量，例如+,d
17. N阶ta1\textasciitilde{}0.r( N 2 1)的散度是阶张量
18. 没有对应uector c:ross的张量运算

\chapter{二阶张量的性质}
在湍流研究中，我们遇到几个二阶张量，特别是速度梯度张量 aUi/axj、粘性应力张量 zij 和雷诺应力张量 (uiuj) - 本附录的目的是回顾本文中使用的这些张量的一些属性。

hfatrix 表示法 第一个观察结果是二阶张量 b 的分量 bij Iorm a x 矩阵，

\begin{sourceblock}
\includegraphics[page=689,pagebox=mediabox,trim={64.39bp 879.89bp 64.39bp 981.88bp},clip,width=0.94\textwidth,max height=0.38\textheight]{Pope.pdf}
\end{sourceblock}
类似地，一阶张量 u 的分量 ui 形成一个列向量

\begin{sourceblock}
\includegraphics[page=689,pagebox=mediabox,trim={64.39bp 627.92bp 64.39bp 1242.84bp},clip,width=0.94\textwidth,max height=0.38\textheight]{Pope.pdf}
\end{sourceblock}
因此，线性代数研究中熟悉的矩阵性质直接适用于二阶张量。

在继续之前，我们先澄清符号和术语的两个方面。首先，在本附录中，我们使用小写字母（例如 b 和 u）表示张量，使用大写字母表示相应的矩阵和列向量（e、B 和 U）。其次，术语“向量”在笛卡尔张量和矩阵分析中都使用。必要时，可以使用术语“一阶张量”和“列向量”来区分这两种不同的含义。表 B.1 显示了以所考虑的三种符号表示的各种运算。

\begin{center}\small 表 B.1。涉及以各种符号表示的一阶和二阶张量的运算\end{center}
\begin{sourceblock}
\includegraphics[page=690,pagebox=mediabox,trim={64.39bp 1673.79bp 64.39bp 254.97bp},clip,width=\textwidth,max height=0.72\textheight]{Pope.pdf}
\end{sourceblock}
直接表示法 后缀表示法 矩阵表示法
内积

两个向量的

张量积

两个向量的

的内积为

二阶张量

和矢量

二阶转置内积

张量和向量

两个的内积

二阶张量

添加两个

二阶张量

二阶张量的唯一收缩bii对应于矩阵B的迹trace(B)，即对角线分量之和
该迹可用于将张量分解为各向同性部分和偏向部分。根据定义，各向同性张量具有相同的分量
啊每个坐标系。不存在奇数阶各向同性张量，而各向同性二阶张量是克罗内克三角洲的标量倍数。 b 的各向同性部分是
根据定义，偏张量的迹为零，因此 b 的偏部分为

\begin{sourceblock}
\includegraphics[page=690,pagebox=mediabox,trim={64.39bp 248.97bp 64.39bp 1685.79bp},clip,width=0.94\textwidth,max height=0.38\textheight]{Pope.pdf}
\end{sourceblock}
因此，b 具有分解

\begin{sourceblock}
\includegraphics[page=690,pagebox=mediabox,trim={64.39bp 119.39bp 64.39bp 1836.77bp},clip,width=0.94\textwidth,max height=0.38\textheight]{Pope.pdf}
\end{sourceblock}
或者用矩阵表示法，I 是 3 x 3 单位矩阵（对应于 ），分解为：

\begin{sourceblock}
\includegraphics[page=691,pagebox=mediabox,trim={64.39bp 1677.79bp 64.39bp 258.97bp},clip,width=0.94\textwidth,max height=0.38\textheight]{Pope.pdf}
\end{sourceblock}
\begin{sourceblock}
\includegraphics[page=691,pagebox=mediabox,trim={64.39bp 1556.80bp 64.39bp 333.96bp},clip,width=0.94\textwidth,max height=0.38\textheight]{Pope.pdf}
\end{sourceblock}
偏差部分可以进一步分解为对称（S）和反对称（R）部分。

\begin{sourceblock}
\includegraphics[page=691,pagebox=mediabox,trim={64.39bp 1288.84bp 64.39bp 640.92bp},clip,width=0.94\textwidth,max height=0.38\textheight]{Pope.pdf}
\end{sourceblock}
或者，用矩阵表示法，

\begin{sourceblock}
\includegraphics[page=691,pagebox=mediabox,trim={64.39bp 1091.86bp 64.39bp 769.90bp},clip,width=0.94\textwidth,max height=0.38\textheight]{Pope.pdf}
\end{sourceblock}
因此，二阶张量 b 可以分解为各向同性部分 betaij、对称偏量部分 sij 和反对称偏量部分之和

\begin{sourceblock}
\includegraphics[page=691,pagebox=mediabox,trim={64.39bp 876.89bp 64.39bp 1008.87bp},clip,width=0.94\textwidth,max height=0.38\textheight]{Pope.pdf}
\end{sourceblock}
由于是对称的且 rij 是反对称的，因此使用两个后缀
和矩阵表示法表明

(a) 的对角线分量为零，
(b) rikrkj 和 sikrkj Sjkrki 是对称的，
(c) Sjkrki 是反对称的，并且
酉变换 方向余弦矩阵 A 是酉矩阵。回顾
aij 定义 ail = - ej，我们可以观察到 A 的第 j 列是
ei 列向量对应于 E 坐标系中 Zj 的分量。因此 A 的列是相互正交的单位向量。因此，根据方程式。 （A.13）和\textbackslash\{\}A.14），我们得到

\begin{sourceblock}
\includegraphics[page=691,pagebox=mediabox,trim={64.39bp 119.39bp 64.39bp 1845.77bp},clip,width=0.94\textwidth,max height=0.38\textheight]{Pope.pdf}
\end{sourceblock}
因此，一阶和二阶张量的变换规则（方程（A.17）和（A.23））对应于酉变换

\begin{sourceblock}
\includegraphics[page=692,pagebox=mediabox,trim={64.39bp 1677.79bp 64.39bp 267.97bp},clip,width=0.94\textwidth,max height=0.38\textheight]{Pope.pdf}
\end{sourceblock}
\begin{sourceblock}
\includegraphics[page=692,pagebox=mediabox,trim={64.39bp 1582.80bp 64.39bp 363.95bp},clip,width=0.94\textwidth,max height=0.38\textheight]{Pope.pdf}
\end{sourceblock}
主轴

我们现在进一步考虑对称（不一定是偏向）二阶张量的属性。从线性代数（参见例如 Franklin (1968), p. 100）我们得到以下重要结果。如果 S 是实对称矩阵，则存在对角化 S 的酉矩阵：

\begin{sourceblock}
\includegraphics[page=692,pagebox=mediabox,trim={64.39bp 1141.86bp 64.39bp 720.91bp},clip,width=0.94\textwidth,max height=0.38\textheight]{Pope.pdf}
\end{sourceblock}
对角线分量 A 和 A3（均为实数）是特征值
S 的 ;II 和 的列（用 Z1、Z2 和 Z3 表示）是相应的特征向量。根据其定义，特征值和特征向量满足方程

\begin{sourceblock}
\includegraphics[page=692,pagebox=mediabox,trim={64.39bp 870.89bp 64.39bp 1073.86bp},clip,width=0.94\textwidth,max height=0.38\textheight]{Pope.pdf}
\end{sourceblock}
其中括号内的后缀被排除在求和约定之外。特征向量 PI、Z2？ Z3 为称为 S 主轴的特殊坐标系提供正交基。

如果特征值不同，则单位特征向量根据其排序和+符号确定。因此，主轴根据轴的反射和 90" 旋转确定。如果两个特征值相等，
= A2 f A3 说，那么与 Z3 正交的每个向量都是特征向量。因此，任何包含 Z3 作为基向量的坐标系都提供 S 的主轴。如果所有三个特征值都相等，则 A1 A2 A3 A，则 S AI，每个
= = = =向量是特征向量，每个坐标系都提供主轴。

对角线和非对角线分量之间的区别经常存在。例如，对于应力张量，+对角线分量是法向应力，非对角线分量是剪切应力。上述结果表明，这种区别不是固有的，而是完全取决于坐标系：在主轴中，非对角线分量为零。本质区别在于各向同性部分之间

\begin{sourceblock}
\includegraphics[page=693,pagebox=mediabox,trim={64.39bp 1615.80bp 64.39bp 199.97bp},clip,width=0.94\textwidth,max height=0.38\textheight]{Pope.pdf}
\end{sourceblock}
\begin{sourceblock}
\includegraphics[page=693,pagebox=mediabox,trim={64.39bp 1410.82bp 64.39bp 447.94bp},clip,width=0.94\textwidth,max height=0.38\textheight]{Pope.pdf}
\end{sourceblock}
最后的结果有什么意义呢？

从方程。 (B.17) 获取结果

\begin{sourceblock}
\includegraphics[page=693,pagebox=mediabox,trim={64.39bp 1077.86bp 64.39bp 777.90bp},clip,width=0.94\textwidth,max height=0.38\textheight]{Pope.pdf}
\end{sourceblock}
\begin{sourceblock}
\includegraphics[page=693,pagebox=mediabox,trim={64.39bp 1002.87bp 64.39bp 939.88bp},clip,width=0.94\textwidth,max height=0.38\textheight]{Pope.pdf}
\end{sourceblock}
\begin{sourceblock}
\includegraphics[page=693,pagebox=mediabox,trim={64.39bp 927.88bp 64.39bp 1013.87bp},clip,width=0.94\textwidth,max height=0.38\textheight]{Pope.pdf}
\end{sourceblock}
因此认为 S 的主轴也是
S2，以及（通过归纳）S3、S4、...。 ..

令 dB) a 、 1,2 和 3 为相互正交的向量，并定义
对称二阶张量 s 为

\begin{sourceblock}
\includegraphics[page=693,pagebox=mediabox,trim={64.39bp 663.92bp 64.39bp 1271.84bp},clip,width=0.94\textwidth,max height=0.38\textheight]{Pope.pdf}
\end{sourceblock}
通过考虑主轴上的 s，证明其特征值是
（对于 0）相应的单位特征向量是

\begin{sourceblock}
\includegraphics[page=693,pagebox=mediabox,trim={64.39bp 414.95bp 64.39bp 1527.81bp},clip,width=0.94\textwidth,max height=0.38\textheight]{Pope.pdf}
\end{sourceblock}
（因此每个对称正半定二阶张量
有分解方程。 （B.25）。）

令 f、g 和 h 为相互正交的单位向量。使用结果
练习 B.4 的结果表明

\begin{sourceblock}
\includegraphics[page=693,pagebox=mediabox,trim={64.39bp 156.98bp 64.39bp 1788.77bp},clip,width=0.94\textwidth,max height=0.38\textheight]{Pope.pdf}
\end{sourceblock}
对于任何二阶张量 b，幂定义为

\begin{sourceblock}
\includegraphics[page=694,pagebox=mediabox,trim={64.39bp 1655.79bp 64.39bp 287.96bp},clip,width=0.94\textwidth,max height=0.38\textheight]{Pope.pdf}
\end{sourceblock}
等等，对应矩阵表达式

\begin{sourceblock}
\includegraphics[page=694,pagebox=mediabox,trim={64.39bp 1498.81bp 64.39bp 444.94bp},clip,width=0.94\textwidth,max height=0.38\textheight]{Pope.pdf}
\end{sourceblock}
就组件bij而言，等式： (B.28) 是

\begin{sourceblock}
\includegraphics[page=694,pagebox=mediabox,trim={64.39bp 1333.83bp 64.39bp 600.92bp},clip,width=0.94\textwidth,max height=0.38\textheight]{Pope.pdf}
\end{sourceblock}
请注意，b2 表示 b2 的 i-j 分量，它与 b2 的 i-j 分量的平方（即 (bij)2）不同。 （九个量 (bij)2 不是张量的组成部分。）

从张量获得的任何标量（例如 bii, b? b b n * b 3 bbijbjkbki）i s
t 1 JU Ir 称为不变量，因为它的值在任何坐标系中都是相同的。对于对称二阶张量 s（具有相应的矩阵 S），三个主不变量定义为

\begin{sourceblock}
\includegraphics[page=694,pagebox=mediabox,trim={64.39bp 920.88bp 64.39bp 1026.87bp},clip,width=0.94\textwidth,max height=0.38\textheight]{Pope.pdf}
\end{sourceblock}
\begin{sourceblock}
\includegraphics[page=694,pagebox=mediabox,trim={64.39bp 815.90bp 64.39bp 1123.86bp},clip,width=0.94\textwidth,max height=0.38\textheight]{Pope.pdf}
\end{sourceblock}
\begin{sourceblock}
\includegraphics[page=694,pagebox=mediabox,trim={64.39bp 709.91bp 64.39bp 1227.85bp},clip,width=0.94\textwidth,max height=0.38\textheight]{Pope.pdf}
\end{sourceblock}
其中 det 表示行列式。

根据定义，不变量在每个坐标系中都是相同的。在评估主轴的主不变量时，我们得到

\begin{sourceblock}
\includegraphics[page=694,pagebox=mediabox,trim={64.39bp 359.95bp 64.39bp 1582.80bp},clip,width=0.94\textwidth,max height=0.38\textheight]{Pope.pdf}
\end{sourceblock}
\begin{sourceblock}
\includegraphics[page=694,pagebox=mediabox,trim={64.39bp 260.97bp 64.39bp 1684.79bp},clip,width=0.94\textwidth,max height=0.38\textheight]{Pope.pdf}
\end{sourceblock}
可能会注意到，I E 和 III 是特征值的最简单可能的线性、二次和三次组合，不依赖于它们的特征值
特征方程 特征值-4gen向量方程（Eq. (B.1 8)）可以写成

\begin{sourceblock}
\includegraphics[page=695,pagebox=mediabox,trim={64.39bp 1671.79bp 64.39bp 273.96bp},clip,width=0.94\textwidth,max height=0.38\textheight]{Pope.pdf}
\end{sourceblock}
从线性方程组理论我们知道，只有当 (S-21) 的行列式为零时，才存在非平凡解。行列式的计算（练习 B.6）得出三次特征方程

\begin{sourceblock}
\includegraphics[page=695,pagebox=mediabox,trim={64.39bp 1455.82bp 64.39bp 483.94bp},clip,width=0.94\textwidth,max height=0.38\textheight]{Pope.pdf}
\end{sourceblock}
因此，主要不变量的重要性。

评估行列式

从而验证方程。 （B.38）。评估 ( A 2)(R2 和
讨论其意义。

凯莱-汉密尔顿定理 在本文中，重要的凯莱-汉密尔顿定理指出矩阵 S 满足其自身的特征方程：

\begin{sourceblock}
\includegraphics[page=695,pagebox=mediabox,trim={64.39bp 740.91bp 64.39bp 1202.85bp},clip,width=0.94\textwidth,max height=0.38\textheight]{Pope.pdf}
\end{sourceblock}
（参见练习 B.7）。该方程表明 S3 可以表示为 S2、S 和 I 的线性组合。

进一步，将方程预乘S，我们发现也可以用同样的方式表示。因此，通过归纳，Sn (n 2 3) 可以表示为 S2、S 和 I 的线性组合，其中这些矩阵的三个系数是 S 的不变量。这些观察结果是表示定理的基础。

如果所有特征值 A 、 A2 和 A3 均非零，则 det(S) 也非零，则 S 是非奇异的。在这种情况下，存在一个逆S-'，使得S-'S I。当它变换到主轴时，S成为特征值的对角矩阵（方程（B.17）），那么显然S-'是

\begin{sourceblock}
\includegraphics[page=695,pagebox=mediabox,trim={64.39bp 119.39bp 64.39bp 1768.78bp},clip,width=0.94\textwidth,max height=0.38\textheight]{Pope.pdf}
\end{sourceblock}
S-' 的显式表达式是通过预乘等式获得的。 (B.39) 由
就反 s-I 的张量分量而言，结果为。 s-ijI

\begin{sourceblock}
\includegraphics[page=696,pagebox=mediabox,trim={64.39bp 1673.79bp 64.39bp 263.97bp},clip,width=0.94\textwidth,max height=0.38\textheight]{Pope.pdf}
\end{sourceblock}
另一方面，如果一个或多个特征值为零，则 S 是奇异的，并且 S-I 是未定义的。 S 的秩不足等于零特征值的数量。

张量 s 是正定的，如果对于每个非零向量 v，数量

\begin{sourceblock}
\includegraphics[page=696,pagebox=mediabox,trim={64.39bp 1408.82bp 64.39bp 537.93bp},clip,width=0.94\textwidth,max height=0.38\textheight]{Pope.pdf}
\end{sourceblock}
是严格正数。由于 Q ( v ) 是标量，因此其值不受坐标变化的影响。在 s 的主轴上， fii 是 v 的分量，

\begin{sourceblock}
\includegraphics[page=696,pagebox=mediabox,trim={64.39bp 1191.85bp 64.39bp 696.91bp},clip,width=0.94\textwidth,max height=0.38\textheight]{Pope.pdf}
\end{sourceblock}
显然，当且仅当所有特征值都严格为正时，s 才是正定的。如果特征值为非负 (Ai 2 0)，则 Q ( v ) 为非负，并且 s 为正半定。

变换方程。 (B.39) 通过预乘 A- T 到主轴
通过用 A 的项表示结果来进行后乘。比较
的结果与特征方程（Eq. B.38），从而验证
凯莱-汉密尔顿定理。

对于正半定 S，定义

\begin{sourceblock}
\includegraphics[page=696,pagebox=mediabox,trim={64.39bp 554.93bp 64.39bp 1297.84bp},clip,width=0.94\textwidth,max height=0.38\textheight]{Pope.pdf}
\end{sourceblock}
\begin{sourceblock}
\includegraphics[page=696,pagebox=mediabox,trim={64.39bp 463.94bp 64.39bp 1484.81bp},clip,width=0.94\textwidth,max height=0.38\textheight]{Pope.pdf}
\end{sourceblock}
证明 sl/*i 是具有以下属性的对称张量

\begin{sourceblock}
\includegraphics[page=696,pagebox=mediabox,trim={64.39bp 322.96bp 64.39bp 1624.80bp},clip,width=0.94\textwidth,max height=0.38\textheight]{Pope.pdf}
\end{sourceblock}
即 s1l2i 是 s 的（对称）平方根。

s-I 是对称正定二阶的逆
阶张量 s 考虑方程

\begin{sourceblock}
\includegraphics[page=696,pagebox=mediabox,trim={64.39bp 119.39bp 64.39bp 1846.77bp},clip,width=0.94\textwidth,max height=0.38\textheight]{Pope.pdf}
\end{sourceblock}
通过考虑主轴，表明该方程定义了
主轴与 s 的特征向量对齐的椭球体
并且轴的长度等于

s 的特征值。

德尔塔函数

\chapter{狄拉克 δ 函数}
\subsection*{C.l 狄拉克的定义}\phantomsection\addcontentsline{toc}{subsection}{C.l 狄拉克的定义}
在本文中，我们广泛使用了 delta 函数 6(x) 和相关量，例如 Heaviside 函数 H(x)。在本附录中，回顾了这些广义函数的基本性质，并得出了一些特别有用的结果。

I 6(x) 的定义 狄拉克δ函数具有以下性质

\begin{sourceblock}
\includegraphics[page=698,pagebox=mediabox,trim={64.39bp 793.90bp 64.39bp 1152.85bp},clip,width=0.94\textwidth,max height=0.38\textheight]{Pope.pdf}
\end{sourceblock}
然而，应当理解，6(x)不是通常意义上的函数，并且上述方程没有提供充分的定义。相反，6(x) 是一个广义函数。因此，它的定义和属性取决于积分

\begin{sourceblock}
\includegraphics[page=698,pagebox=mediabox,trim={64.39bp 483.94bp 64.39bp 1416.82bp},clip,width=0.94\textwidth,max height=0.38\textheight]{Pope.pdf}
\end{sourceblock}
其中 g (x) 是任何测试函数（具有适当的可微性和无穷大行为）。

作为接近狄拉克 δ 函数的一种方法，请考虑以下函数
对于 n 1,2,3, .... 这是 delta 序列的示例。可能会被认可

\begin{sourceblock}
\includegraphics[page=699,pagebox=mediabox,trim={64.39bp 1177.85bp 64.39bp 109.44bp},clip,width=\textwidth,max height=0.72\textheight]{Pope.pdf}
\par\vspace{0.45\baselineskip}
\small 图 C.1。函数 D,(x)（方程（C.4）），n = 1,4 和 16。
\end{sourceblock}
D,(x) 是均值为零、标准差为 l / n 的高斯分布（参见方程（3.41））。因此我们有

\begin{sourceblock}
\includegraphics[page=699,pagebox=mediabox,trim={64.39bp 906.89bp 64.39bp 996.87bp},clip,width=0.94\textwidth,max height=0.38\textheight]{Pope.pdf}
\end{sourceblock}
更一般地，对于 rn 2 0，

\begin{sourceblock}
\includegraphics[page=699,pagebox=mediabox,trim={64.39bp 745.91bp 64.39bp 1157.85bp},clip,width=0.94\textwidth,max height=0.38\textheight]{Pope.pdf}
\end{sourceblock}
其中 \{ f i o , f i 1 , h , h , f i 4 , .。。 \} \{ l O l O , 3 .是标准化高斯分布的矩。图 C.1 显示了 D,(x) 的几个值
对于测试函数 g(x)，使用泰勒级数
（其中g(m?(x)表示g(x)的m次导数），对于式（C.3）形式的积分，我们得到

\begin{sourceblock}
\includegraphics[page=699,pagebox=mediabox,trim={64.39bp 119.39bp 64.39bp 1813.77bp},clip,width=0.94\textwidth,max height=0.38\textheight]{Pope.pdf}
\end{sourceblock}
\section*{C.2 狄拉克的性质}\phantomsection\addcontentsline{toc}{section}{C.2 狄拉克的性质}
虽然极限 limn+* D,(x) 不存在，但显然我们有
根据定义，狄拉克δ函数是恰好具有此性质的广义函数。

C.2 6 ( x ) 的性质 现在逐项列出 6(x) 的基本性质。

\begin{sourceblock}
\includegraphics[page=700,pagebox=mediabox,trim={64.39bp 1128.86bp 64.39bp 687.91bp},clip,width=0.94\textwidth,max height=0.38\textheight]{Pope.pdf}
\end{sourceblock}
这就是 delta 函数的筛选性质：6(x - a) 被引用
作为“a 处的 delta 函数”，并且从函数 g(x) 中积分
在等式中(C.11) 筛选出特定值 g(a)。

(ii) 更一般地，如果函数 f(x) 在 a 处连续，则对于 el 0

\begin{sourceblock}
\includegraphics[page=700,pagebox=mediabox,trim={64.39bp 790.90bp 64.39bp 1050.87bp},clip,width=0.94\textwidth,max height=0.38\textheight]{Pope.pdf}
\end{sourceblock}
如果 el 或 为零，或者 iff (x) 不连续，则该积分未定义
(iii) 合法运算包括乘以函数，例如。例如，f(x)b(x-
a)（假设 f (x) 在 a 处连续），以及加法，例如 f (x)6(x
a) h(x)b(x - b)。非法操作包括将两个 delta 相乘
同一变量的函数，例如 6(x a)6(x - b)，以及除以
德尔塔函数。

(iv) 涉及 delta 函数的表达式的唯一意义，例如，
f (x)6(x - a)，是乘以测试得到的值
函数 g(x) 并积分。例如，我们有

\begin{sourceblock}
\includegraphics[page=700,pagebox=mediabox,trim={64.39bp 119.39bp 64.39bp 1661.79bp},clip,width=0.94\textwidth,max height=0.38\textheight]{Pope.pdf}
\end{sourceblock}
\section*{C.3 狄拉克的导数}\phantomsection\addcontentsline{toc}{section}{C.3 狄拉克的导数}
因此，方程

\begin{sourceblock}
\includegraphics[page=701,pagebox=mediabox,trim={64.39bp 1722.78bp 64.39bp 222.97bp},clip,width=0.94\textwidth,max height=0.38\textheight]{Pope.pdf}
\end{sourceblock}
是正确的。 （回想一下，除以 6(x a) 是不合法的，所以有
并不意味着 f (x) 等于 f (a)。）

(v) 从归一化条件 6(x) dx 1 可以明显看出
6(x) 是密度，PDF 也是如此。所以，如果 x 的尺寸为长度
6(x) 具有长度倒数的尺寸。独立变革下
变量，6(x) 变换为密度。例如，考虑 6 [ ( x
a)/b] 对于 b 0。在设置 y (x a)/b 时，我们得到

\begin{sourceblock}
\includegraphics[page=701,pagebox=mediabox,trim={64.39bp 1162.85bp 64.39bp 606.92bp},clip,width=0.94\textwidth,max height=0.38\textheight]{Pope.pdf}
\end{sourceblock}
通过比较该方程中的第一个和最后一个表达式，并通过
重复 b 0 的分析，我们得出结论：变换
规则是（对于 b 0）

\begin{sourceblock}
\includegraphics[page=701,pagebox=mediabox,trim={64.39bp 937.88bp 64.39bp 978.88bp},clip,width=0.94\textwidth,max height=0.38\textheight]{Pope.pdf}
\end{sourceblock}
请注意，对于 b -1，该方程显示 delta 函数为

\begin{sourceblock}
\includegraphics[page=701,pagebox=mediabox,trim={64.39bp 777.90bp 64.39bp 1110.86bp},clip,width=0.94\textwidth,max height=0.38\textheight]{Pope.pdf}
\end{sourceblock}
C.3 6(x) 的导数 δ 序列 D,(x)（方程（C.4））的导数定义为

\begin{sourceblock}
\includegraphics[page=701,pagebox=mediabox,trim={64.39bp 528.93bp 64.39bp 1416.82bp},clip,width=0.94\textwidth,max height=0.38\textheight]{Pope.pdf}
\end{sourceblock}
如图 C.2 所示，n 的三个值。通过分部积分，我们得到

\begin{sourceblock}
\includegraphics[page=701,pagebox=mediabox,trim={64.39bp 119.39bp 64.39bp 1788.77bp},clip,width=0.94\textwidth,max height=0.38\textheight]{Pope.pdf}
\end{sourceblock}
\begin{sourceblock}
\includegraphics[page=702,pagebox=mediabox,trim={64.39bp 1305.84bp 64.39bp 109.44bp},clip,width=\textwidth,max height=0.72\textheight]{Pope.pdf}
\par\vspace{0.45\baselineskip}
\small 图 C.2。 n = 1,2 和 4 时的函数（方程（C.19））。
\end{sourceblock}
（测试函数满足 D,,(x)g(x) 在无穷远处消失所需的弱条件。）然后定义 delta 函数 6 (1) 的导数，使得
6 (1) (x)g(x) dx lim

\begin{sourceblock}
\includegraphics[page=702,pagebox=mediabox,trim={64.39bp 786.90bp 64.39bp 1161.85bp},clip,width=0.94\textwidth,max height=0.38\textheight]{Pope.pdf}
\end{sourceblock}
然后可以观察到 -6(')(x a) 筛选 a 处的导数，即 g (1) (a)；并且 6 (1) 是反对称的：

\begin{sourceblock}
\includegraphics[page=702,pagebox=mediabox,trim={64.39bp 588.93bp 64.39bp 1352.83bp},clip,width=0.94\textwidth,max height=0.38\textheight]{Pope.pdf}
\end{sourceblock}
（另请参见图 C.2）。

高阶导数的定义类似。对于 m 阶导数，我们得到

\begin{sourceblock}
\includegraphics[page=702,pagebox=mediabox,trim={64.39bp 354.95bp 64.39bp 1543.81bp},clip,width=0.94\textwidth,max height=0.38\textheight]{Pope.pdf}
\end{sourceblock}
变换规则为（对于b 0）

\begin{sourceblock}
\includegraphics[page=702,pagebox=mediabox,trim={64.39bp 192.98bp 64.39bp 1725.78bp},clip,width=0.94\textwidth,max height=0.38\textheight]{Pope.pdf}
\end{sourceblock}
\section*{C.4 泰勒级数}\phantomsection\addcontentsline{toc}{section}{C.4 泰勒级数}
（参见练习1）因此

\begin{sourceblock}
\includegraphics[page=703,pagebox=mediabox,trim={64.39bp 1691.79bp 64.39bp 213.97bp},clip,width=0.94\textwidth,max height=0.38\textheight]{Pope.pdf}
\end{sourceblock}
按照方程中采用的方法。 （C.16），表明反式
6 (1) 的形成规则是 (对于 b 0)

\begin{sourceblock}
\includegraphics[page=703,pagebox=mediabox,trim={64.39bp 1453.82bp 64.39bp 462.94bp},clip,width=0.94\textwidth,max height=0.38\textheight]{Pope.pdf}
\end{sourceblock}
C.4 泰勒级数 一个可能令人惊讶的观察结果是 6(x) 具有泰勒级数展开式。测试函数 g(x) 的泰勒级数为
就 delta 函数而言，左边是

\begin{sourceblock}
\includegraphics[page=703,pagebox=mediabox,trim={64.39bp 936.88bp 64.39bp 966.88bp},clip,width=0.94\textwidth,max height=0.38\textheight]{Pope.pdf}
\end{sourceblock}
而右侧是

因此我们得到

\begin{sourceblock}
\includegraphics[page=703,pagebox=mediabox,trim={64.39bp 591.92bp 64.39bp 1320.83bp},clip,width=0.94\textwidth,max height=0.38\textheight]{Pope.pdf}
\end{sourceblock}
\section*{C.5 海维赛函数}\phantomsection\addcontentsline{toc}{section}{C.5 海维赛函数}
对应于增量序列 D,(x)（方程（C.4）），序列 S,(x) 定义为
如图C.3所示。显然，当 n 趋于无穷大时，S,(x) 趋于单位阶跃函数

\begin{sourceblock}
\includegraphics[page=703,pagebox=mediabox,trim={64.39bp 119.39bp 64.39bp 1804.77bp},clip,width=0.94\textwidth,max height=0.38\textheight]{Pope.pdf}
\end{sourceblock}
\begin{sourceblock}
\includegraphics[page=704,pagebox=mediabox,trim={64.39bp 1152.86bp 64.39bp 109.44bp},clip,width=\textwidth,max height=0.72\textheight]{Pope.pdf}
\par\vspace{0.45\baselineskip}
\small 图 C.3。 n = 1,4 和 16 时的函数 S,(x)（方程（C.31））。
\end{sourceblock}
Heaviside 函数 H(x) 被定义为具有以下性质的广义函数

\begin{sourceblock}
\includegraphics[page=704,pagebox=mediabox,trim={64.39bp 862.89bp 64.39bp 1037.87bp},clip,width=0.94\textwidth,max height=0.38\textheight]{Pope.pdf}
\end{sourceblock}
Jx H(x)的导数为6(x)；反之，H(x) 是积分 6(y) dy 这从上面的定义可以明显看出，或者，从
证明 Heaviside 函数的变换规则为
{\small 1 H(x -a)，对于 b Om\par}
考虑函数的顺序

\begin{sourceblock}
\includegraphics[page=704,pagebox=mediabox,trim={64.39bp 119.39bp 64.39bp 1839.77bp},clip,width=0.94\textwidth,max height=0.38\textheight]{Pope.pdf}
\end{sourceblock}
\section*{C.6 多维度}\phantomsection\addcontentsline{toc}{section}{C.6 多维度}
S,(x) 和 S,(x) 趋于无穷大时有何不同？表明

\begin{sourceblock}
\includegraphics[page=705,pagebox=mediabox,trim={64.39bp 1686.79bp 64.39bp 216.97bp},clip,width=0.94\textwidth,max height=0.38\textheight]{Pope.pdf}
\end{sourceblock}
C.6 多维 设 y 为三维欧几里德空间中的一个点。 y 处的三维 Delta 函数写作 6(x y )，它只是三个一维 Delta 函数的乘积：

\begin{sourceblock}
\includegraphics[page=705,pagebox=mediabox,trim={64.39bp 1320.83bp 64.39bp 624.92bp},clip,width=0.94\textwidth,max height=0.38\textheight]{Pope.pdf}
\end{sourceblock}
具有筛分特性

其中 dx 为 dxl dx2d x3+（请注意，形成不同变量的 delta 函数的乘积是合法的，例如 6(xl a)6(x2 b)，但不能形成同一变量的 delta 函数的乘积，例如 6(x a)6(x b)。）

傅立叶变换

本附录的目的是提供定义和总结

\chapter{傅里叶变换}
进一步的解释和结果，读者可以参考标准文本，例如，Bracewell (1965)、Lighthill (1970) 和 Priestley (1981)。

给定函数 f (t)，其傅里叶变换为

\begin{sourceblock}
\includegraphics[page=706,pagebox=mediabox,trim={64.39bp 927.88bp 64.39bp 976.88bp},clip,width=0.94\textwidth,max height=0.38\textheight]{Pope.pdf}
\end{sourceblock}
逆变换是

\begin{sourceblock}
\includegraphics[page=706,pagebox=mediabox,trim={64.39bp 735.91bp 64.39bp 1165.85bp},clip,width=0.94\textwidth,max height=0.38\textheight]{Pope.pdf}
\end{sourceblock}
为了使 f (t) 和 g ( o ) 形成傅里叶变换对，上述积分必须收敛，至少作为广义函数。这里显示的是时域（t）和频域（a）之间的变换；物理空间之间变换的相应公式
(IC) 和波数空间是明显的。第 679 页的表 D.1 给出了一些有用的傅立叶变换对，Erddlyi、Oberhettinger 和 Tricomi (1954) 提供了全面的汇编。

傅立叶变换的定义没有唯一的约定。在某些定义中，负指数 -iot 出现在逆变换中，并且 2n 的因子可以以不同的方式在 3 和 3 之间分割。这里使用的约定与 Batchelor (1953)、Monin 和 Yaglom (1975) 以及 Termekes 和 Lumley (1972) 使用的约定相同。

\begin{center}\small 表 D.1。傅立叶变换对（a、b 和 v 是实常数，其中 b > 0 且 y > -\}）\end{center}
\clearpage
\begin{sourceblock}
\includegraphics[page=707,pagebox=mediabox,trim={64.39bp 119.39bp 64.39bp 261.97bp},clip,width=\textwidth,max height=0.72\textheight]{Pope.pdf}
\end{sourceblock}
证明 g ( o ) 也是实数且偶数。逆变换是

\begin{sourceblock}
\includegraphics[page=708,pagebox=mediabox,trim={64.39bp 1683.79bp 64.39bp 218.97bp},clip,width=0.94\textwidth,max height=0.38\textheight]{Pope.pdf}
\end{sourceblock}
方程（D.6）和（D.7）定义了余弦傅里叶变换及其逆变换。

在考虑光谱时，有时考虑实偶函数 f (t) 的两次傅立叶变换是很方便的，即

\begin{sourceblock}
\includegraphics[page=708,pagebox=mediabox,trim={64.39bp 1371.83bp 64.39bp 552.93bp},clip,width=0.94\textwidth,max height=0.38\textheight]{Pope.pdf}
\end{sourceblock}
从而得到反演公式

\begin{sourceblock}
\includegraphics[page=708,pagebox=mediabox,trim={64.39bp 1205.85bp 64.39bp 696.91bp},clip,width=0.94\textwidth,max height=0.38\textheight]{Pope.pdf}
\end{sourceblock}
不包含因子 2（参见方程 (D.7)）。

δ 函数 δ 函数 6(t a) 的傅立叶变换为（根据方程 (D.1) 并调用筛选特性方程 (C . l l ) ）

\begin{sourceblock}
\includegraphics[page=708,pagebox=mediabox,trim={64.39bp 851.89bp 64.39bp 1056.87bp},clip,width=0.94\textwidth,max height=0.38\textheight]{Pope.pdf}
\end{sourceblock}
并且，特别是，

\begin{sourceblock}
\includegraphics[page=708,pagebox=mediabox,trim={64.39bp 696.91bp 64.39bp 1211.85bp},clip,width=0.94\textwidth,max height=0.38\textheight]{Pope.pdf}
\end{sourceblock}
设 g ( o ) (1/2n)e-'"'，则反演公式(Eq.(D.2))得

\begin{sourceblock}
\includegraphics[page=708,pagebox=mediabox,trim={64.39bp 519.93bp 64.39bp 1383.83bp},clip,width=0.94\textwidth,max height=0.38\textheight]{Pope.pdf}
\end{sourceblock}
这是一个了不起且有价值的成果。然而，由于方程中的积分。 (D.12) 显然是发散的，它像 6(t a) 一样必须被视为广义函数。也就是说，以 G(t) 为测试函数，方程： (D.12) 的含义

\begin{sourceblock}
\includegraphics[page=708,pagebox=mediabox,trim={64.39bp 168.98bp 64.39bp 1773.78bp},clip,width=0.94\textwidth,max height=0.38\textheight]{Pope.pdf}
\end{sourceblock}
卷积 给定两个函数 fa(t) 和 fb(t)（两者都有傅里叶变换），它们的卷积定义为

\begin{sourceblock}
\includegraphics[page=709,pagebox=mediabox,trim={64.39bp 1597.80bp 64.39bp 347.96bp},clip,width=0.94\textwidth,max height=0.38\textheight]{Pope.pdf}
\end{sourceblock}
代入 r = t - s，卷积的傅立叶变换为

\begin{sourceblock}
\includegraphics[page=709,pagebox=mediabox,trim={64.39bp 1149.85bp 64.39bp 575.93bp},clip,width=0.94\textwidth,max height=0.38\textheight]{Pope.pdf}
\end{sourceblock}
也就是说，卷积的傅里叶变换等于 2n: 和函数的傅里叶变换的乘积。

帕塞瓦尔定理

我们考虑两个函数和 fb(t) 的乘积的积分
fa@) 有傅立叶变换 g,(a) 和 g b ( a )。通过写和逆
f a f b 傅里叶变换，我们得到

\begin{sourceblock}
\includegraphics[page=709,pagebox=mediabox,trim={64.39bp 591.92bp 64.39bp 1308.83bp},clip,width=0.94\textwidth,max height=0.38\textheight]{Pope.pdf}
\end{sourceblock}
对所有 t 的指数项积分得到 2n:6(a - a"),s ee Eq. （D.12），以便轻松执行所有 o" 的积分，产生

\begin{sourceblock}
\includegraphics[page=709,pagebox=mediabox,trim={64.39bp 359.95bp 64.39bp 1535.81bp},clip,width=0.94\textwidth,max height=0.38\textheight]{Pope.pdf}
\end{sourceblock}
这是帕塞瓦尔第二定理。

对于 和 是相同函数的情况（ i . f 和
f a f b f a f b 对应 g, gb g)，式(D. 18) 变为 Parseval 的 Jirst 定理：

\begin{sourceblock}
\includegraphics[page=709,pagebox=mediabox,trim={64.39bp 119.39bp 64.39bp 1806.77bp},clip,width=0.94\textwidth,max height=0.38\textheight]{Pope.pdf}
\end{sourceblock}
如果 为实数，则可以重新表示为

\begin{sourceblock}
\includegraphics[page=710,pagebox=mediabox,trim={64.39bp 1575.80bp 64.39bp 236.97bp},clip,width=0.94\textwidth,max height=0.38\textheight]{Pope.pdf}
\end{sourceblock}
其中 f(t) 是傅里叶变换 g(o) 的可微函数，
得到如下结果：

\begin{sourceblock}
\includegraphics[page=710,pagebox=mediabox,trim={64.39bp 1259.84bp 64.39bp 642.92bp},clip,width=0.94\textwidth,max height=0.38\textheight]{Pope.pdf}
\end{sourceblock}
\begin{sourceblock}
\includegraphics[page=710,pagebox=mediabox,trim={64.39bp 1146.86bp 64.39bp 756.90bp},clip,width=0.94\textwidth,max height=0.38\textheight]{Pope.pdf}
\end{sourceblock}
\begin{sourceblock}
\includegraphics[page=710,pagebox=mediabox,trim={64.39bp 1021.87bp 64.39bp 866.89bp},clip,width=0.94\textwidth,max height=0.38\textheight]{Pope.pdf}
\end{sourceblock}
重新表达 f 为实数情况的右侧。

令 为具有标准 de- 的零均值高斯分布

\begin{sourceblock}
\includegraphics[page=710,pagebox=mediabox,trim={64.39bp 773.90bp 64.39bp 1068.87bp},clip,width=0.94\textwidth,max height=0.38\textheight]{Pope.pdf}
\end{sourceblock}
令 f b ( t )= ;0 , b2)，其中 a 和 b 为正常数。显示
和 的卷积为 ;0 , a2 b2)。

\chapter{平稳随机过程的谱表示}
平稳的谱表示

随机过程

本附录的目的是显示统计平稳随机过程 U(t)、其频谱表示（以傅立叶模式表示）、其频谱 E (a) 及其自相关函数之间的联系
统计平稳随机过程具有恒定方差，因此不会随着 It1 趋于无穷大而衰减至零。因此，U(t) 的傅立叶变换不存在。这一事实导致了重大的技术困难，只能使用比此处适用的更复杂的数学工具来克服。我们通过首先发展周期函数的想法，然后将结果扩展到感兴趣的非周期函数来规避这个困难。

\subsection*{E.l 傅立叶级数}\phantomsection\addcontentsline{toc}{subsection}{E.l 傅立叶级数}
我们首先考虑时间间隔 0 5 t T 内的非随机实际过程 U(t)。通过定义周期性地继续该过程

\begin{sourceblock}
\includegraphics[page=711,pagebox=mediabox,trim={64.39bp 619.92bp 64.39bp 1308.83bp},clip,width=0.94\textwidth,max height=0.38\textheight]{Pope.pdf}
\end{sourceblock}
对于所有非零整数 N。

一段时间内 U(t) 的时间平均值定义为
其他量的时间平均值也以类似的方式定义。 U(t) 的波动定义为
显然它的时间平均值 (u(t))= 为零。

对于每个整数 n，频率 « 定义为

\begin{sourceblock}
\includegraphics[page=712,pagebox=mediabox,trim={64.39bp 1730.78bp 64.39bp 213.97bp},clip,width=0.94\textwidth,max height=0.38\textheight]{Pope.pdf}
\end{sourceblock}
我们考虑正 n 和负 n，并观察到

\begin{sourceblock}
\includegraphics[page=712,pagebox=mediabox,trim={64.39bp 1605.80bp 64.39bp 341.96bp},clip,width=0.94\textwidth,max height=0.38\textheight]{Pope.pdf}
\end{sourceblock}
第 n 个复傅立叶模是
恩！ - cos(w,t) + isin(@,t)

\begin{sourceblock}
\includegraphics[page=712,pagebox=mediabox,trim={64.39bp 1421.82bp 64.39bp 522.93bp},clip,width=0.94\textwidth,max height=0.38\textheight]{Pope.pdf}
\end{sourceblock}
其平均时间为

\begin{sourceblock}
\includegraphics[page=712,pagebox=mediabox,trim={64.39bp 1240.84bp 64.39bp 651.92bp},clip,width=0.94\textwidth,max height=0.38\textheight]{Pope.pdf}
\end{sourceblock}
因此模式满足正交条件

\begin{sourceblock}
\includegraphics[page=712,pagebox=mediabox,trim={64.39bp 1056.87bp 64.39bp 879.89bp},clip,width=0.94\textwidth,max height=0.38\textheight]{Pope.pdf}
\end{sourceblock}
过程 u(t) 可以表示为傅立叶级数，

\begin{sourceblock}
\includegraphics[page=712,pagebox=mediabox,trim={64.39bp 854.89bp 64.39bp 1056.87bp},clip,width=0.94\textwidth,max height=0.38\textheight]{Pope.pdf}
\end{sourceblock}
其中 \{a,,5,\} 是实数，\{c,\} 是复傅立叶系数。由于时间平均值 (u(t))7 为零，因此由等式得出： (E.7) co 也为零。展开正弦和余弦，方程。 (E.9) 是
u(t) = + an) + iby +

\begin{sourceblock}
\includegraphics[page=712,pagebox=mediabox,trim={64.39bp 452.94bp 64.39bp 1455.82bp},clip,width=0.94\textwidth,max height=0.38\textheight]{Pope.pdf}
\end{sourceblock}
由于 u(t) 是实数，c 满足共轭对称性，
(i.€., a, = a\_, and b, = -b\_,) 使得等式右边的虚数项（E.10）消失。因此，傅立叶级数（方程（E.10））变为

\begin{sourceblock}
\includegraphics[page=712,pagebox=mediabox,trim={64.39bp 119.39bp 64.39bp 1805.77bp},clip,width=0.94\textwidth,max height=0.38\textheight]{Pope.pdf}
\end{sourceblock}
也可以写成

\begin{sourceblock}
\includegraphics[page=713,pagebox=mediabox,trim={64.39bp 1679.79bp 64.39bp 236.97bp},clip,width=0.94\textwidth,max height=0.38\textheight]{Pope.pdf}
\end{sourceblock}
其中第 n 阶傅里叶模式的幅度为

\begin{sourceblock}
\includegraphics[page=713,pagebox=mediabox,trim={64.39bp 1549.80bp 64.39bp 388.95bp},clip,width=0.94\textwidth,max height=0.38\textheight]{Pope.pdf}
\end{sourceblock}
其相位为

\begin{sourceblock}
\includegraphics[page=713,pagebox=mediabox,trim={64.39bp 1429.82bp 64.39bp 510.94bp},clip,width=0.94\textwidth,max height=0.38\textheight]{Pope.pdf}
\end{sourceblock}
傅里叶系数的显式表达式可以通过乘以等式获得。 (E.9) 通过第 -mth 模式并求平均值：

\begin{sourceblock}
\includegraphics[page=713,pagebox=mediabox,trim={64.39bp 1071.86bp 64.39bp 682.91bp},clip,width=0.94\textwidth,max height=0.38\textheight]{Pope.pdf}
\end{sourceblock}
引入由以下定义的运算符 Fun\{\} 很方便

\begin{sourceblock}
\includegraphics[page=713,pagebox=mediabox,trim={64.39bp 899.89bp 64.39bp 996.87bp},clip,width=0.94\textwidth,max height=0.38\textheight]{Pope.pdf}
\end{sourceblock}
所以方程。 (E.16) 可以写成

\begin{sourceblock}
\includegraphics[page=713,pagebox=mediabox,trim={64.39bp 772.90bp 64.39bp 1170.85bp},clip,width=0.94\textwidth,max height=0.38\textheight]{Pope.pdf}
\end{sourceblock}
因此，算子3 \} 确定频率为an 的模式的傅立叶系数。

方程 (E.9) 是 u(t) 的频谱表示，将 u(t) 给出为离散傅里叶模式 ei ont 的总和，并用傅里叶系数 c 加权。考虑到非周期情况的扩展，谱表示也可以写成

\begin{sourceblock}
\includegraphics[page=713,pagebox=mediabox,trim={64.39bp 404.95bp 64.39bp 1496.81bp},clip,width=0.94\textwidth,max height=0.38\textheight]{Pope.pdf}
\end{sourceblock}
\begin{sourceblock}
\includegraphics[page=713,pagebox=mediabox,trim={64.39bp 265.97bp 64.39bp 1665.79bp},clip,width=0.94\textwidth,max height=0.38\textheight]{Pope.pdf}
\end{sourceblock}
其中a是连续频率。方程中的积分。 (E.19) 是傅立叶逆变换（参见.Ed q.(D.2)），因此 z ( a ) 被识别为

\section*{E.2 周期性随机过程}\phantomsection\addcontentsline{toc}{section}{E.2 周期性随机过程}
我们现在认为 u(t) 是一个统计平稳的周期性随机过程。 A11 上面获得的结果对于该过程的每个实现都有效。特别地，傅立叶系数 c 由等式给出。 （E.1 6）。然而，由于 u(t) 是随机的，因此傅里叶系数 c 是随机变量。我们现在展示均值 (c,) 是如何为零，并且对应于不同频率的系数是不相关的。

方程的平均值。 (E.9) 是

\begin{sourceblock}
\includegraphics[page=714,pagebox=mediabox,trim={64.39bp 1377.83bp 64.39bp 534.93bp},clip,width=0.94\textwidth,max height=0.38\textheight]{Pope.pdf}
\end{sourceblock}
回想一下，co 为零，对于 n 0，(w(t)) 独立于 t 的平稳性条件显然意味着 (c,) 为零。

傅里叶模态的协方差为

e (u(t)u(tl)) dt' dt
-铁-iwmt

\begin{sourceblock}
\includegraphics[page=714,pagebox=mediabox,trim={64.39bp 884.89bp 64.39bp 960.88bp},clip,width=0.94\textwidth,max height=0.38\textheight]{Pope.pdf}
\end{sourceblock}
第三行来自替换 t' = t s 以及自协方差的定义

\begin{sourceblock}
\includegraphics[page=714,pagebox=mediabox,trim={64.39bp 707.91bp 64.39bp 1237.84bp},clip,width=0.94\textwidth,max height=0.38\textheight]{Pope.pdf}
\end{sourceblock}
其中（由于平稳性）独立于 t。被积函数 eiQmSR(\textasciitilde{}) 在 s 内呈周期性，周期为 T，因此大括号中的积分与 t 无关。最后一行来自等式。 （E.8）和（E.17）。

从方程式可以立即看出。 (E.22) 除非 m 等于 -n，否则协方差 (c,c,) 为零：也就是说，对应于不同频率的傅立叶系数是不相关的。对于 rn = -n，方程。 (E.22) 变为
因此方差 ( ) 是 R(s) 的傅里叶系数，因此可以表示为

\begin{sourceblock}
\includegraphics[page=714,pagebox=mediabox,trim={64.39bp 125.98bp 64.39bp 1786.77bp},clip,width=0.94\textwidth,max height=0.38\textheight]{Pope.pdf}
\end{sourceblock}
可以观察到，R(s) 是实数并且是 s 的偶函数，并且它仅取决于幅度 Ic，与相位 8i 无关。

再次考虑到非周期性情况的扩展，我们通过以下方式定义频谱

\begin{sourceblock}
\includegraphics[page=715,pagebox=mediabox,trim={64.39bp 1573.80bp 64.39bp 341.96bp},clip,width=0.94\textwidth,max height=0.38\textheight]{Pope.pdf}
\end{sourceblock}
这样自协方差就可以写成

\begin{sourceblock}
\includegraphics[page=715,pagebox=mediabox,trim={64.39bp 1318.83bp 64.39bp 510.94bp},clip,width=0.94\textwidth,max height=0.38\textheight]{Pope.pdf}
\end{sourceblock}
从其定义可知，E(o)是实数偶函数
( 即 ( - ) ) 那么，通过以下方式定义（替代）ci 频谱是很方便的

\begin{sourceblock}
\includegraphics[page=715,pagebox=mediabox,trim={64.39bp 1124.86bp 64.39bp 822.90bp},clip,width=0.94\textwidth,max height=0.38\textheight]{Pope.pdf}
\end{sourceblock}
并重写方程。 （E.27）作为反余弦变换（方程（D.9））

\begin{sourceblock}
\includegraphics[page=715,pagebox=mediabox,trim={64.39bp 958.88bp 64.39bp 940.88bp},clip,width=0.94\textwidth,max height=0.38\textheight]{Pope.pdf}
\end{sourceblock}
将上式中的 s 设为 0，我们得到

\begin{sourceblock}
\includegraphics[page=715,pagebox=mediabox,trim={64.39bp 660.92bp 64.39bp 1202.85bp},clip,width=0.94\textwidth,max height=0.38\textheight]{Pope.pdf}
\end{sourceblock}
因此，(c c*)r 表示第 n 个模式对方差的贡献，类似地
是对 ( \textasciitilde{} ( t )f\textasciitilde{}ro)m 频率范围 o, 5 lo1 ob 的贡献。从方程式可以清楚地看出。 (E.26) 与 R(s) 一样，频谱 E ( o ) 与相位无关。

可以观察到，方程。 (E.27) 将 R(s) 识别为 \textasciitilde{} ( o(c)f. Eq. (D.2)) 的傅里叶逆变换。因此，可以直接从方程式验证。 (E.25)，E ( o ) 是 R(s) 的傅里叶变换：

\begin{sourceblock}
\includegraphics[page=715,pagebox=mediabox,trim={64.39bp 119.39bp 64.39bp 1810.77bp},clip,width=0.94\textwidth,max height=0.38\textheight]{Pope.pdf}
\end{sourceblock}
类似地，E(w) 是 R(s) 的傅里叶变换的两倍：

\begin{sourceblock}
\includegraphics[page=716,pagebox=mediabox,trim={64.39bp 1706.78bp 64.39bp 239.97bp},clip,width=0.94\textwidth,max height=0.38\textheight]{Pope.pdf}
\end{sourceblock}
k ( o ) 确定了 R(s) 并将其作为傅立叶变换对后，我们现在采用等式： (E.31) 作为 E ( o ) 的定义（而不是方程 ( .26)）。

频谱 \textasciitilde{} ( oca)n 也可以用傅里叶变换 z(o) 来表示，如方程 1 中定义。 （E.20）。考虑无穷小区间 (o, o d o )，其中包含零个或一个离散频率 a,。如果它不包含任何离散频率，则

\begin{sourceblock}
\includegraphics[page=716,pagebox=mediabox,trim={64.39bp 1320.83bp 64.39bp 625.92bp},clip,width=0.94\textwidth,max height=0.38\textheight]{Pope.pdf}
\end{sourceblock}
另一方面，如果它包含离散频率 a，则

\begin{sourceblock}
\includegraphics[page=716,pagebox=mediabox,trim={64.39bp 1165.85bp 64.39bp 763.90bp},clip,width=0.94\textwidth,max height=0.38\textheight]{Pope.pdf}
\end{sourceblock}
因此，一般来说，
频谱的基本属性是

( i ) 是非负的 ( i ( w )' 2 0) （参见方程 (E.35)）；

(ii) i ( o ) 是实数（因为 R(s) 是偶数，即 R(s) R(-s)）；和
(iii) 是偶数，即 i ( w ) i(-a) （因为 R(s) 是实数）。

表 E.1 总结了 u(t)、c、z(w)、R(s) 和 E(a) 之间的关系。

通过对方程进行傅立叶变换。 (E.9)，表明 z(w) 给定
由方程式(E.20) 是 u(t) 的傅立叶变换。 （提示：参见方程（D.l2）。）

显示统计上的傅立叶系数 c、a、ib
平稳、周期性随机过程满足

\begin{sourceblock}
\includegraphics[page=716,pagebox=mediabox,trim={64.39bp 246.97bp 64.39bp 1674.79bp},clip,width=0.94\textwidth,max height=0.38\textheight]{Pope.pdf}
\end{sourceblock}
\section*{E.3 非周期随机过程}\phantomsection\addcontentsline{toc}{section}{E.3 非周期随机过程}
\begin{center}\small 表 E.1。周期性和非周期性统计平稳随机过程的谱特性\end{center}
\clearpage
\begin{sourceblock}
\includegraphics[page=717,pagebox=mediabox,trim={64.39bp 186.98bp 64.39bp 257.97bp},clip,width=\textwidth,max height=0.72\textheight]{Pope.pdf}
\end{sourceblock}
在近似意义上，tEc 非周期性情况' 可以被视为
相邻离散频率之间是

\begin{sourceblock}
\includegraphics[page=718,pagebox=mediabox,trim={64.39bp 1723.78bp 64.39bp 222.97bp},clip,width=0.94\textwidth,max height=0.38\textheight]{Pope.pdf}
\end{sourceblock}
当 T 趋于无穷大时，它趋于零。因此，在任何给定的频率范围（a，Io cob）内，离散频率的数量
A ) 趋于无穷大，因此（在近似意义上）(cob 谱在 om 上变得连续
Monin 和 Yaglom (1975) 以及 Priestley (1981) 对非周期情况进行了严格的数学处理。简而言之，虽然非周期过程 u(t) 没有傅里叶变换，但它确实具有傅里叶-斯蒂尔切斯积分的频谱表示

\begin{sourceblock}
\includegraphics[page=718,pagebox=mediabox,trim={64.39bp 1246.84bp 64.39bp 654.92bp},clip,width=0.94\textwidth,max height=0.38\textheight]{Pope.pdf}
\end{sourceblock}
其中 Z ( o ) 是不可微的复随机函数。可以观察到，dZ(o)（在非周期性情况下）对应于z(o)do（在周期性情况下，方程（E.20））。非周期性情况的频谱对应于方程。 (E.35) 对于周期性情况是

\begin{sourceblock}
\includegraphics[page=718,pagebox=mediabox,trim={64.39bp 978.88bp 64.39bp 950.88bp},clip,width=0.94\textwidth,max height=0.38\textheight]{Pope.pdf}
\end{sourceblock}
表 E.1 比较了周期性和非周期性情况的频谱表示。

\section*{E.4 过程的导数}\phantomsection\addcontentsline{toc}{section}{E.4 过程的导数}
对于周期性情况，过程 u(t) 具有频谱表示式 Eq. （E.9）。对时间进行微分，我们得到 duldt 的谱表示：

\begin{sourceblock}
\includegraphics[page=718,pagebox=mediabox,trim={64.39bp 456.94bp 64.39bp 1429.82bp},clip,width=0.94\textwidth,max height=0.38\textheight]{Pope.pdf}
\end{sourceblock}
类似地，k阶导数的谱表示为

\begin{sourceblock}
\includegraphics[page=718,pagebox=mediabox,trim={64.39bp 279.96bp 64.39bp 1616.80bp},clip,width=0.94\textwidth,max height=0.38\textheight]{Pope.pdf}
\end{sourceblock}
因为 u (k) (t) 由 u(t) 的傅立叶系数（即方程 (E.42) 中的 c）确定，所以 d k ) ( t ) 的自协方差和谱由下式确定：

它遵循与得出方程相同的过程。 (E.25) 表示 u(k) 的自相关

\begin{sourceblock}
\includegraphics[page=719,pagebox=mediabox,trim={64.39bp 1589.80bp 64.39bp 258.97bp},clip,width=0.94\textwidth,max height=0.38\textheight]{Pope.pdf}
\end{sourceblock}
通过将其与方程的 (2k) 阶导数进行比较。 （E.25），

\begin{sourceblock}
\includegraphics[page=719,pagebox=mediabox,trim={64.39bp 1412.82bp 64.39bp 491.94bp},clip,width=0.94\textwidth,max height=0.38\textheight]{Pope.pdf}
\end{sourceblock}
我们得到结果

\begin{sourceblock}
\includegraphics[page=719,pagebox=mediabox,trim={64.39bp 1256.84bp 64.39bp 646.92bp},clip,width=0.94\textwidth,max height=0.38\textheight]{Pope.pdf}
\end{sourceblock}
d k ) ( t )i s 的频谱（参见方程 (E.26)）

\begin{sourceblock}
\includegraphics[page=719,pagebox=mediabox,trim={64.39bp 879.89bp 64.39bp 1068.87bp},clip,width=0.94\textwidth,max height=0.38\textheight]{Pope.pdf}
\end{sourceblock}
或者，可以通过对方程进行傅立叶变换来获得结果。 （E.45）。

总之，对于过程 \textasciitilde{} ( t )t 的 k 阶导数 u (k ) (t)，h e 自协方差 \&(s) 由等式给出： （E.45），而频谱是

\begin{sourceblock}
\includegraphics[page=719,pagebox=mediabox,trim={64.39bp 617.92bp 64.39bp 1321.83bp},clip,width=0.94\textwidth,max height=0.38\textheight]{Pope.pdf}
\end{sourceblock}
这两个结果既适用于周期性情况，也适用于非周期性情况。

\chapter{离散傅里叶变换}
我们考虑周期函数 u(t)，周期为 T，在周期内以 N 等间隔的时间采样，其中 N 是偶数。基于这些样本，离散傅里叶变换定义了 N 个傅里叶系数（与傅里叶级数的系数相关），从而提供了 u(t) 的离散谱表示。

快速傅里叶变换 (FFT) 是离散傅里叶变换的有效实现。在数值方法中，FFT及其逆可用于时域和频域之间、物理空间和波数空间之间的变换。在具有一个或多个统计同质方向的流的 DNS 和 LES 中，通常使用伪谱方法（在同质方向），并广泛使用 FFT 在物理空间和波数空间之间进行转换。

时间间隔 A t 定义为
采样时间是

\begin{sourceblock}
\includegraphics[page=720,pagebox=mediabox,trim={64.39bp 582.93bp 64.39bp 1361.83bp},clip,width=0.94\textwidth,max height=0.38\textheight]{Pope.pdf}
\end{sourceblock}
样本表示为

然后将离散傅立叶变换的复系数定义为 1 - 1 \textasciitilde{} < k < l \textasciitilde{} b \textasciitilde{}

\begin{sourceblock}
\includegraphics[page=720,pagebox=mediabox,trim={64.39bp 177.98bp 64.39bp 1701.79bp},clip,width=0.94\textwidth,max height=0.38\textheight]{Pope.pdf}
\end{sourceblock}
其中（与傅立叶级数一样）频率定义为

\begin{sourceblock}
\includegraphics[page=721,pagebox=mediabox,trim={64.39bp 1696.79bp 64.39bp 250.97bp},clip,width=0.94\textwidth,max height=0.38\textheight]{Pope.pdf}
\end{sourceblock}
如下所示，逆变换为
为了确认逆变换的形式，我们考虑数量
从复平面来看，ZjYN 是 N 个点的质心
2rrijkl N k的N个值。当 j 为零或 N 的整数倍时，每个点都位于 ( 1 , ) 处，因此 Zj,N 是单位。当j不是N的整数倍时，点关于原点对称分布，因此ZjT为零。因此
对于 j / N 整数，
有了这个结果，方程的右侧就可以得到。 (F.6) 可以写成
在最终的总和中，唯一的非零贡献是 j，这验证了逆变换，方程： （F.6）。

研究各系数之间的关系具有重要意义

这些量的定义（方程（FA）和（E.9））我们有

\begin{sourceblock}
\includegraphics[page=722,pagebox=mediabox,trim={64.39bp 1641.79bp 64.39bp 215.97bp},clip,width=0.94\textwidth,max height=0.38\textheight]{Pope.pdf}
\end{sourceblock}
在考虑一般情况之前，我们考虑更简单的情况，其中对于具有 (ak( 2 amax 的所有模式，傅里叶系数 ck 为零，其中 amaisx 是离散傅里叶变换中表示的最高频率，
在这种情况下，方程中的总和。 (FJO) 均有效地从 - ( 1) 到 ( N I)，因此系数 Ek 和 ck 相同。

对于一般情况，我们需要考虑高于 a,,, 的频率。对于 -(LN - 1) 5 k 5 LN 范围内的 k，并且对于非零整数 m，第 (k mN) 模式的频率为

\begin{sourceblock}
\includegraphics[page=722,pagebox=mediabox,trim={64.39bp 1071.86bp 64.39bp 869.89bp},clip,width=0.94\textwidth,max height=0.38\textheight]{Pope.pdf}
\end{sourceblock}
\begin{sourceblock}
\includegraphics[page=722,pagebox=mediabox,trim={64.39bp 955.88bp 64.39bp 992.87bp},clip,width=0.94\textwidth,max height=0.38\textheight]{Pope.pdf}
\end{sourceblock}
在采样时间 tj 处，第 (k mN) 个模式与第 k 个模式无法区分，因为

\begin{sourceblock}
\includegraphics[page=722,pagebox=mediabox,trim={64.39bp 777.90bp 64.39bp 1160.85bp},clip,width=0.94\textwidth,max height=0.38\textheight]{Pope.pdf}
\end{sourceblock}
第(k mN) 模式被称为第k模式。

系数 Zk 可以根据其定义（方程（F.4））确定，并用傅里叶级数代替

\begin{sourceblock}
\includegraphics[page=722,pagebox=mediabox,trim={64.39bp 397.95bp 64.39bp 1377.83bp},clip,width=0.94\textwidth,max height=0.38\textheight]{Pope.pdf}
\end{sourceblock}
因此，系数 Zk 是与第 k 个模式混叠的所有模式的傅里叶系数之和。

鉴于共轭对称性，N 个复数系数 Zk 可以用 N 个实数表示（例如，\%\{Ek\} for k = 0, 1, . , A N 和
对于 k = 1,2,。。 , LN- 1, 参见练习 El). 离散傅立叶变换 3 \{ \textasciitilde{} k \}
需要 N 2 次运算才能直接根据等式 1 中的总和计算 Tk。 (F.4)。然而，通过使用快速傅立叶变换(FFT)，可以在大约N log N运算中获得相同的结果(参见例如Brigham (1974))。因此，对于以足够小的时间间隔采样的周期性数据，FFT 是一种有效的工具，用于评估傅里叶系数、频谱、自相关（如频谱的傅里叶逆变换）、卷积和导数：

\begin{sourceblock}
\includegraphics[page=723,pagebox=mediabox,trim={64.39bp 1453.82bp 64.39bp 478.94bp},clip,width=0.94\textwidth,max height=0.38\textheight]{Pope.pdf}
\end{sourceblock}
与傅里叶变换一样，离散傅里叶变换有多种定义。这里使用的定义与傅里叶级数有最直接的联系。在数值实现中，通常使用练习中给出的替代定义。

证明，对于实数 u(t)，系数 ZIk 满足

\begin{sourceblock}
\includegraphics[page=723,pagebox=mediabox,trim={64.39bp 1068.87bp 64.39bp 870.89bp},clip,width=0.94\textwidth,max height=0.38\textheight]{Pope.pdf}
\end{sourceblock}
那 Fo 和 ?我是真实的。

\begin{sourceblock}
\includegraphics[page=723,pagebox=mediabox,trim={64.39bp 909.89bp 64.39bp 984.88bp},clip,width=0.94\textwidth,max height=0.38\textheight]{Pope.pdf}
\end{sourceblock}
\begin{sourceblock}
\includegraphics[page=723,pagebox=mediabox,trim={64.39bp 829.90bp 64.39bp 1112.86bp},clip,width=0.94\textwidth,max height=0.38\textheight]{Pope.pdf}
\end{sourceblock}
离散傅里叶变换的另一种定义是

\begin{sourceblock}
\includegraphics[page=723,pagebox=mediabox,trim={64.39bp 648.92bp 64.39bp 1232.84bp},clip,width=0.94\textwidth,max height=0.38\textheight]{Pope.pdf}
\end{sourceblock}
证明其逆函数是

系数 Zk 和 EL 之间有什么关系？

\chapter{幂律谱}
鲍威\textasciitilde{}法佩克特拉
在湍流研究中，高频（或大波数）光谱的形状受到相当多的关注。附录的目的是展示幂律谱 ( E ( o )

对于大型 co)、底层随机过程 u(t) 和二阶结构函数 D(s)。

我们考虑具有有限方差 (u2) 和积分时间尺度的统计平稳过程 u(t) 自相关函数

\begin{sourceblock}
\includegraphics[page=724,pagebox=mediabox,trim={64.39bp 1021.87bp 64.39bp 923.88bp},clip,width=0.94\textwidth,max height=0.38\textheight]{Pope.pdf}
\end{sourceblock}
和一半频谱形成傅里叶余弦变换对：

\begin{sourceblock}
\includegraphics[page=724,pagebox=mediabox,trim={64.39bp 850.89bp 64.39bp 1051.87bp},clip,width=0.94\textwidth,max height=0.38\textheight]{Pope.pdf}
\end{sourceblock}
\begin{sourceblock}
\includegraphics[page=724,pagebox=mediabox,trim={64.39bp 732.91bp 64.39bp 1163.85bp},clip,width=0.94\textwidth,max height=0.38\textheight]{Pope.pdf}
\end{sourceblock}
第三个感兴趣的量是二阶结构函数

\begin{sourceblock}
\includegraphics[page=724,pagebox=mediabox,trim={64.39bp 464.94bp 64.39bp 1455.82bp},clip,width=0.94\textwidth,max height=0.38\textheight]{Pope.pdf}
\end{sourceblock}
根据定义，幂律谱变化为

\begin{sourceblock}
\includegraphics[page=724,pagebox=mediabox,trim={64.39bp 309.96bp 64.39bp 1635.79bp},clip,width=0.94\textwidth,max height=0.38\textheight]{Pope.pdf}
\end{sourceblock}
而幂律结构函数变化为

\begin{sourceblock}
\includegraphics[page=724,pagebox=mediabox,trim={64.39bp 177.98bp 64.39bp 1765.78bp},clip,width=0.94\textwidth,max height=0.38\textheight]{Pope.pdf}
\end{sourceblock}
这里的目的是理解 p 和 q 的特定值的意义以及它们之间的联系。

通过在方程中设置 s 0 获得的第一个观测值。 (G.2) 的方差是
假设 (u2) 是有限的。因此，如果 E ( o ) 是幂律谱，则积分收敛的要求规定 p > 1。

u(t) 导数的光谱得出一系列类似的结果。假设 u(t) 的 n 次导数存在，并表示为

\begin{sourceblock}
\includegraphics[page=725,pagebox=mediabox,trim={64.39bp 1282.84bp 64.39bp 624.92bp},clip,width=0.94\textwidth,max height=0.38\textheight]{Pope.pdf}
\end{sourceblock}
u(")(t)i 的自相关性
（见附录E，方程（E.45）），其频谱为

\begin{sourceblock}
\includegraphics[page=725,pagebox=mediabox,trim={64.39bp 934.88bp 64.39bp 1014.87bp},clip,width=0.94\textwidth,max height=0.38\textheight]{Pope.pdf}
\end{sourceblock}
（参见方程式（E.47））。因此我们得到

\begin{sourceblock}
\includegraphics[page=725,pagebox=mediabox,trim={64.39bp 642.92bp 64.39bp 1259.84bp},clip,width=0.94\textwidth,max height=0.38\textheight]{Pope.pdf}
\end{sourceblock}
如果 u(t) 可微分 n 次（在均方意义上），则左侧是有限的。那么，如果E(co)是幂律谱，则要求方程中的积分： ((3.1 1) 收敛规定

\begin{sourceblock}
\includegraphics[page=725,pagebox=mediabox,trim={64.39bp 427.95bp 64.39bp 1519.81bp},clip,width=0.94\textwidth,max height=0.38\textheight]{Pope.pdf}
\end{sourceblock}
对于无限可微的过程，例如由纳维-斯托克斯方程演化的速度，可从方程 1 得出。 （G.12）（对于大 co）光谱衰减比 co 的任何幂都更快：例如，它可能会衰减为 exp(-co) 或 exp(-co2)。然而，在很大的频率范围内（比如 ol < co < wh），可能会出现幂律谱，其中

\begin{sourceblock}
\includegraphics[page=726,pagebox=mediabox,trim={64.39bp 1245.84bp 64.39bp 109.44bp},clip,width=\textwidth,max height=0.72\textheight]{Pope.pdf}
\par\vspace{0.45\baselineskip}
\small 图 G.1。无量纲幂律谱 E(w)： (G.14) 对于 v = i。
\end{sourceblock}
如果过程 u(t) 至少一次连续可微，则对于小 s，结构函数为

\begin{sourceblock}
\includegraphics[page=726,pagebox=mediabox,trim={64.39bp 942.88bp 64.39bp 934.88bp},clip,width=0.94\textwidth,max height=0.38\textheight]{Pope.pdf}
\end{sourceblock}
即，q = 2 的幂律（方程（G.6））。

检查无量纲幂律谱是有启发性的

\begin{sourceblock}
\includegraphics[page=726,pagebox=mediabox,trim={64.39bp 722.91bp 64.39bp 1164.85bp},clip,width=0.94\textwidth,max height=0.38\textheight]{Pope.pdf}
\end{sourceblock}
对于正参数 v 的各种值。无量纲积分时间尺度是统一的，即

\begin{sourceblock}
\includegraphics[page=726,pagebox=mediabox,trim={64.39bp 518.93bp 64.39bp 1383.83bp},clip,width=0.94\textwidth,max height=0.38\textheight]{Pope.pdf}
\end{sourceblock}
而（对于给定 v ）a 是指定的（参见下面的等式（G.18）），以便方差为单位。图 G.1 显示了 1 和 2 之间 v 的 E ( o )。对于大 o，双对数图上的直线清楚地显示了幂律行为

\begin{sourceblock}
\includegraphics[page=726,pagebox=mediabox,trim={64.39bp 306.96bp 64.39bp 1641.79bp},clip,width=0.94\textwidth,max height=0.38\textheight]{Pope.pdf}
\end{sourceblock}
相应的自相关（通过逆变换获得

\begin{sourceblock}
\includegraphics[page=726,pagebox=mediabox,trim={64.39bp 128.98bp 64.39bp 1747.78bp},clip,width=0.94\textwidth,max height=0.38\textheight]{Pope.pdf}
\end{sourceblock}
\begin{sourceblock}
\includegraphics[page=727,pagebox=mediabox,trim={64.39bp 1227.85bp 64.39bp 109.44bp},clip,width=\textwidth,max height=0.72\textheight]{Pope.pdf}
\par\vspace{0.45\baselineskip}
\small 图 G.2。自相关函数 R(s)，方程。 (G.19)，对于 v = 4，ala 密封 3
\end{sourceblock}
其中 K, 是第二类修正贝塞尔函数。归一化条件 (u2) R(0) 1 得出

\begin{sourceblock}
\includegraphics[page=727,pagebox=mediabox,trim={64.39bp 978.88bp 64.39bp 959.88bp},clip,width=0.94\textwidth,max height=0.38\textheight]{Pope.pdf}
\end{sourceblock}
所以方程。 （G.17）可以重写

\begin{sourceblock}
\includegraphics[page=727,pagebox=mediabox,trim={64.39bp 833.89bp 64.39bp 1113.86bp},clip,width=0.94\textwidth,max height=0.38\textheight]{Pope.pdf}
\end{sourceblock}
这些自相关如图 G.2 所示。 （对于 v，方程 (G.19) 给出的自相关性就是 R(s) exp(-IS\textbackslash\{\})。）

自相关的表达式还远未揭示出这一点。然而，K,,(sa)（对于小参数）的扩展导致结构函数（对于小 s）的信息丰富的表达式：

\begin{sourceblock}
\includegraphics[page=727,pagebox=mediabox,trim={64.39bp 423.95bp 64.39bp 1434.82bp},clip,width=0.94\textwidth,max height=0.38\textheight]{Pope.pdf}
\end{sourceblock}
因此，结构函数随着指数的幂律变化
这种行为在图 G.3 中很明显。

由幂律表达式得出的结论

\begin{sourceblock}
\includegraphics[page=728,pagebox=mediabox,trim={64.39bp 1257.84bp 64.39bp 109.44bp},clip,width=\textwidth,max height=0.72\textheight]{Pope.pdf}
\par\vspace{0.45\baselineskip}
\small 图 G.3。二阶结构函数 D(s)、Bq。 (G.20)，对于 v = 4,4,...,13,2。观察到，对于 v > 1 且 s 小，所有结构函数都随着 s? 变化。
\end{sourceblock}
\begin{center}\small 表 G.1。幂律谱方程的谱指数 p、结构函数指数 q 和基础过程 u(t) 的可微性之间的关系。 （G.14）\end{center}
\begin{sourceblock}
\includegraphics[page=728,pagebox=mediabox,trim={64.39bp 216.97bp 64.39bp 994.88bp},clip,width=\textwidth,max height=0.72\textheight]{Pope.pdf}
\end{sourceblock}
对于 v 1，对于较大的 a，幂律谱为

\begin{sourceblock}
\includegraphics[page=728,pagebox=mediabox,trim={64.39bp 119.39bp 64.39bp 1840.77bp},clip,width=0.94\textwidth,max height=0.38\textheight]{Pope.pdf}
\end{sourceblock}
而对于小 s，结构函数是

\begin{sourceblock}
\includegraphics[page=729,pagebox=mediabox,trim={64.39bp 1720.78bp 64.39bp 225.97bp},clip,width=0.94\textwidth,max height=0.38\textheight]{Pope.pdf}
\end{sourceblock}
这些关系源于等式。 （（3.14）和（G.20），从中可以推导出 C1 和 C2（取决于 q ）。这是一个代数问题（参见练习 G.1），表明 C1 和 C2 通过以下关系相关：

\begin{sourceblock}
\includegraphics[page=729,pagebox=mediabox,trim={64.39bp 1474.81bp 64.39bp 436.94bp},clip,width=0.94\textwidth,max height=0.38\textheight]{Pope.pdf}
\end{sourceblock}
对于特定情况 q 3，该比率为 0.2489；或者，到一个优秀的

\begin{sourceblock}
\includegraphics[page=729,pagebox=mediabox,trim={64.39bp 1311.83bp 64.39bp 576.93bp},clip,width=0.94\textwidth,max height=0.38\textheight]{Pope.pdf}
\end{sourceblock}
尽管我们考虑了幂律谱的具体示例（e Eq. (G14)），但得出的结论是一般性的。如果光谱表现出
Cleo-P 在相当大的频率范围内具有幂律行为 E(co)，则对于具有 q min(p- 1 2) 的结构函数，存在相应的幂律行为 D(s) C2sq。 （这一断言可以通过分析方程（G.4）来验证，参见 Monin 和 Yaglom (1975)。）对于 q 2，C1 和 Cz 通过方程（G.4）相关。 （G.25）。

确定方程式中的 C1 和 Cz。 （G.23）和（G.24）。随着使用
伽马函数的以下属性：

\begin{sourceblock}
\includegraphics[page=729,pagebox=mediabox,trim={64.39bp 750.91bp 64.39bp 1197.85bp},clip,width=0.94\textwidth,max height=0.38\textheight]{Pope.pdf}
\end{sourceblock}
\begin{sourceblock}
\includegraphics[page=729,pagebox=mediabox,trim={64.39bp 672.91bp 64.39bp 1273.84bp},clip,width=0.94\textwidth,max height=0.38\textheight]{Pope.pdf}
\end{sourceblock}
\begin{sourceblock}
\includegraphics[page=729,pagebox=mediabox,trim={64.39bp 588.93bp 64.39bp 1350.83bp},clip,width=0.94\textwidth,max height=0.38\textheight]{Pope.pdf}
\end{sourceblock}
验证方程（G.25）。

\chapter{欧拉 PDF 方程的推导}
在本附录中，速度 f (V; x, t) 的欧拉 PDF 的输运方程是从纳维-斯托克斯方程导出的。此推导基于 *细粒度 PDF 的属性； Dopazo (1994) 描述了一些其他技术。

细粒度 PDF 对于给定的流实现，速度的（单点、一次性、欧拉）细粒度 PDF 定义为

\begin{sourceblock}
\includegraphics[page=730,pagebox=mediabox,trim={64.39bp 933.88bp 64.39bp 946.88bp},clip,width=0.94\textwidth,max height=0.38\textheight]{Pope.pdf}
\end{sourceblock}
在每个点 x 和时间 t，f ' 是速度空间中位于 V = U(x,t) 的三维 delta 函数。 （附录 C 回顾了 delta 函数的属性。）

细粒度 PDF 在获取和操作 PDF 方程方面非常有用，因为它具有以下两个属性：

\begin{sourceblock}
\includegraphics[page=730,pagebox=mediabox,trim={64.39bp 640.92bp 64.39bp 1305.84bp},clip,width=0.94\textwidth,max height=0.38\textheight]{Pope.pdf}
\end{sourceblock}
\begin{sourceblock}
\includegraphics[page=730,pagebox=mediabox,trim={64.39bp 551.93bp 64.39bp 1373.83bp},clip,width=0.94\textwidth,max height=0.38\textheight]{Pope.pdf}
\end{sourceblock}
第一个关系是从均值 (Q(U(x,t ))) 的一般表达式获得的，方程：第 464 页上的 (12.2)，通过用 V' 替换积分变量 V，进行替换 Q(U(x,t )) = 6(U(x,t) - V)，并且

\begin{sourceblock}
\includegraphics[page=730,pagebox=mediabox,trim={64.39bp 178.98bp 64.39bp 1572.80bp},clip,width=0.94\textwidth,max height=0.38\textheight]{Pope.pdf}
\end{sourceblock}
最后一步源自 delta 函数的筛选特性（方程（C.39））。

第二个关系式， (H.3)，以类似的方式获得

\begin{sourceblock}
\includegraphics[page=731,pagebox=mediabox,trim={64.39bp 1313.83bp 64.39bp 609.92bp},clip,width=0.94\textwidth,max height=0.38\textheight]{Pope.pdf}
\end{sourceblock}
这五个步骤遵循 f' 的定义，方程 1。 ( 1 U 和 函数的无条件均值的定义；筛选性质；方程 (12.3) 的代入；以及条件均值 的定义，方程 (12.4)。

细粒度 PDF 的衍生词
推导 PDF 输运方程时需要细粒度 PDF f '( V ;x , t) 的时间和空间导数。作为一个更简单的预备，考虑具有细粒度 PDF 的标量值可微过程 u(t)，

\begin{sourceblock}
\includegraphics[page=731,pagebox=mediabox,trim={64.39bp 745.91bp 64.39bp 1198.85bp},clip,width=0.94\textwidth,max height=0.38\textheight]{Pope.pdf}
\end{sourceblock}
其中 v 是样本空间变量。 δ函数 6(v - a)（u 为常数）的导数表示为 6 (1) (v - a) (Eq. (C.21))，它是一个奇函数（Eq. ((2.22))）：

\begin{sourceblock}
\includegraphics[page=731,pagebox=mediabox,trim={64.39bp 471.94bp 64.39bp 1434.82bp},clip,width=0.94\textwidth,max height=0.38\textheight]{Pope.pdf}
\end{sourceblock}
因此，微分方程。 (H.6) 对于 t（使用链式法则），我们得到

\begin{sourceblock}
\includegraphics[page=731,pagebox=mediabox,trim={64.39bp 165.98bp 64.39bp 1648.79bp},clip,width=0.94\textwidth,max height=0.38\textheight]{Pope.pdf}
\end{sourceblock}
对于速度的细粒度 PDF，相同的过程产生重新

\begin{sourceblock}
\includegraphics[page=732,pagebox=mediabox,trim={64.39bp 1640.79bp 64.39bp 199.97bp},clip,width=0.94\textwidth,max height=0.38\textheight]{Pope.pdf}
\end{sourceblock}
\begin{sourceblock}
\includegraphics[page=732,pagebox=mediabox,trim={64.39bp 1511.81bp 64.39bp 377.95bp},clip,width=0.94\textwidth,max height=0.38\textheight]{Pope.pdf}
\end{sourceblock}
下面的推导所需的最终结果是

\begin{sourceblock}
\includegraphics[page=732,pagebox=mediabox,trim={64.39bp 1257.84bp 64.39bp 645.92bp},clip,width=0.94\textwidth,max height=0.38\textheight]{Pope.pdf}
\end{sourceblock}
第一步依赖于不可压缩性（V U = 0），第二步依赖于筛选特性（方程（C.15）），第三步依赖于 是一个自变量。

PDF f ( V ;x , t) 输运方程 有了已经获得的结果，PDF f ( V ;x , t) 输运方程的推导就很简单了。 f'(V;x,t) 的基本导数由等式获得。 ( 9 - ( H l l ) 是

\begin{sourceblock}
\includegraphics[page=732,pagebox=mediabox,trim={64.39bp 771.90bp 64.39bp 1158.85bp},clip,width=0.94\textwidth,max height=0.38\textheight]{Pope.pdf}
\end{sourceblock}
这个方程的平均值得出

\begin{sourceblock}
\includegraphics[page=732,pagebox=mediabox,trim={64.39bp 581.93bp 64.39bp 1307.84bp},clip,width=0.94\textwidth,max height=0.38\textheight]{Pope.pdf}
\end{sourceblock}
在等式中。 (H.11) 假设速度场不可压缩；但除此之外，这个 PDF 方程（Eq. (H.13)）非常笼统，并且不包含任何物理学。当使用 Navier-Stokes 方程（第 1-7 页的方程（2.35））代替 DUi/Dt 时，物理现象就出现了：

\begin{sourceblock}
\includegraphics[page=732,pagebox=mediabox,trim={64.39bp 271.97bp 64.39bp 1617.80bp},clip,width=0.94\textwidth,max height=0.38\textheight]{Pope.pdf}
\end{sourceblock}
压力的雷诺分解 (p (p) p') 导致

\begin{sourceblock}
\includegraphics[page=732,pagebox=mediabox,trim={64.39bp 119.39bp 64.39bp 1796.77bp},clip,width=0.94\textwidth,max height=0.38\textheight]{Pope.pdf}
\end{sourceblock}
请注意，d(p)/axi 是非随机的，因此不受均值和条件均值运算的影响；即

\begin{sourceblock}
\includegraphics[page=733,pagebox=mediabox,trim={64.39bp 1630.79bp 64.39bp 260.97bp},clip,width=0.94\textwidth,max height=0.38\textheight]{Pope.pdf}
\end{sourceblock}
分解方程的代入。 （H.15）代入等式。 (H.14) 产生文本中给出的 PDF 方程，方程。 （12.9）。

以下练习讨论脉动速度 PDF g(v;x, t) 输运方程的推导，并对 f 和 g 方程中出现的项进行各种分解。

通过遵循与用于导出传输相同的过程

f ( V ; x, t) 的方程， (H.14)，证明 g(vex
（脉动速度的 PDF）是

\begin{sourceblock}
\includegraphics[page=733,pagebox=mediabox,trim={64.39bp 1044.87bp 64.39bp 844.89bp},clip,width=0.94\textwidth,max height=0.38\textheight]{Pope.pdf}
\end{sourceblock}
使用方程式。第 126 页上的 (5.138) 来替代 Dui/Dt，因此
得到结果

从关系来看

\begin{sourceblock}
\includegraphics[page=733,pagebox=mediabox,trim={64.39bp 583.93bp 64.39bp 1361.83bp},clip,width=0.94\textwidth,max height=0.38\textheight]{Pope.pdf}
\end{sourceblock}
u(x, t) 和 U(x, t) 的 PDF 之间的关系表明它们的演化
方程的关系为

AF

在轴轴处，
以 gf(v;x , t) 为细粒度 PDF \textasciitilde{} ( u ( xt,) - v)，得到

\begin{sourceblock}
\includegraphics[page=734,pagebox=mediabox,trim={64.39bp 1664.79bp 64.39bp 198.97bp},clip,width=0.94\textwidth,max height=0.38\textheight]{Pope.pdf}
\end{sourceblock}
因此表明方程中 dp'/axj 中的项。 (H.18)可以重新表达
- 2 2vi i?vj [gxri(')]

dvi(g(p axi v))=axi dvi(g(\$-..-

\begin{sourceblock}
\includegraphics[page=734,pagebox=mediabox,trim={64.39bp 1431.82bp 64.39bp 516.93bp},clip,width=0.94\textwidth,max height=0.38\textheight]{Pope.pdf}
\end{sourceblock}
其中条件压力-速率-qf-应变张量 RC(vx,, t) 是 de-
按方程式罚款。 （12.20）。

获取关系

\begin{sourceblock}
\includegraphics[page=734,pagebox=mediabox,trim={64.39bp 1200.85bp 64.39bp 694.91bp},clip,width=0.94\textwidth,max height=0.38\textheight]{Pope.pdf}
\end{sourceblock}
因此表明方程中的粘性项。 （H.1 -8）可以重新

\begin{sourceblock}
\includegraphics[page=734,pagebox=mediabox,trim={64.39bp 1002.87bp 64.39bp 845.89bp},clip,width=0.94\textwidth,max height=0.38\textheight]{Pope.pdf}
\end{sourceblock}
其中条件耗散张量 \textasciitilde{} " ( v ,x , t) 定义为 bv
通过使用练习 H.3 和 H.4 的技术，证明 PDF
传输方程，Eq. (H.1 8)，可以重新表示为

\begin{sourceblock}
\includegraphics[page=734,pagebox=mediabox,trim={64.39bp 599.92bp 64.39bp 1286.84bp},clip,width=0.94\textwidth,max height=0.38\textheight]{Pope.pdf}
\end{sourceblock}
\chapter{特征函数}
对于随机变量 U，特征函数定义为

\begin{sourceblock}
\includegraphics[page=735,pagebox=mediabox,trim={64.39bp 1276.84bp 64.39bp 669.92bp},clip,width=0.94\textwidth,max height=0.38\textheight]{Pope.pdf}
\end{sourceblock}
它是实数变量 s （其维度为 u ）的无量纲复函数。与 (V) 的 PDF 一样，特征函数充分表征了随机变量 U。

在某些情况下，用特征函数比用 PDF 更容易获得结果。这些情况包括矩的计算、随机变量的变换、独立随机变量的总和、正态分布和联合正态分布以及模型 P D F 方程的解。现在给出本文中使用的特征函数的基本性质。

与 P D F 的关系 使用均值的定义作为 P D F 的积分（第 41 页上的方程（3.20）），定义特征函数的方程（方程（1.1））可以是

\begin{sourceblock}
\includegraphics[page=735,pagebox=mediabox,trim={64.39bp 539.93bp 64.39bp 1317.83bp},clip,width=0.94\textwidth,max height=0.38\textheight]{Pope.pdf}
\end{sourceblock}
这可以被识别为傅里叶逆变换（参见第 678 页的方程 (D.2)），以便 f ( V ) 和 Y(s) 形成傅里叶变换对：

\begin{sourceblock}
\includegraphics[page=735,pagebox=mediabox,trim={64.39bp 331.96bp 64.39bp 1576.80bp},clip,width=0.94\textwidth,max height=0.38\textheight]{Pope.pdf}
\end{sourceblock}
\begin{sourceblock}
\includegraphics[page=735,pagebox=mediabox,trim={64.39bp 241.97bp 64.39bp 1699.79bp},clip,width=0.94\textwidth,max height=0.38\textheight]{Pope.pdf}
\end{sourceblock}
因为 f ( V ) 是实数，所以 Y(s) 具有共轭对称性：Y(s) Y*(-s)。

原点处的行为 在方程式中将 s 设置为零。 (1.1) 得出重要结果

\begin{sourceblock}
\includegraphics[page=736,pagebox=mediabox,trim={64.39bp 1677.79bp 64.39bp 267.97bp},clip,width=0.94\textwidth,max height=0.38\textheight]{Pope.pdf}
\end{sourceblock}
这对应于 PDF 的标准化条件。 Y(s) 的 k 阶导数是
dsk 因此，U 的矩（关于原点）由下式给出
线性变换 当 a 和 b 为常数时，随机变量
具有特征函数

独立随机变量之和 如果 UI 和 U2 是具有特征函数 Yl(s) 和 Y2(s) 的独立随机变量，则它们的和 U U1 U2 的特征函数

\begin{sourceblock}
\includegraphics[page=736,pagebox=mediabox,trim={64.39bp 665.92bp 64.39bp 1215.85bp},clip,width=0.94\textwidth,max height=0.38\textheight]{Pope.pdf}
\end{sourceblock}
即特征函数的乘积。

正态分布 如果 U 是标准化正态随机变量，则其 PDF 为

\begin{sourceblock}
\includegraphics[page=736,pagebox=mediabox,trim={64.39bp 376.95bp 64.39bp 1517.81bp},clip,width=0.94\textwidth,max height=0.38\textheight]{Pope.pdf}
\end{sourceblock}
然后进行傅立叶逆变换（参见第 679 页的表 D.1），得出

\begin{sourceblock}
\includegraphics[page=736,pagebox=mediabox,trim={64.39bp 119.39bp 64.39bp 1680.79bp},clip,width=0.94\textwidth,max height=0.38\textheight]{Pope.pdf}
\end{sourceblock}
细粒度特征函数 正如定义细粒度 PDF f ( V ) G(U - V) 使其平均值为 PDF ( V ) 一样，细粒度特征函数也定义为
使得其均值一致为 ( s ) C， ( V ) 和 ( s ) 构成傅里叶变换对。当 Q 是第二个随机变量时，我们有（参见

\begin{sourceblock}
\includegraphics[page=737,pagebox=mediabox,trim={64.39bp 1388.82bp 64.39bp 485.94bp},clip,width=0.94\textwidth,max height=0.38\textheight]{Pope.pdf}
\end{sourceblock}
相应地，对于细粒度的特征函数我们有

\begin{sourceblock}
\includegraphics[page=737,pagebox=mediabox,trim={64.39bp 1080.86bp 64.39bp 772.90bp},clip,width=0.94\textwidth,max height=0.38\textheight]{Pope.pdf}
\end{sourceblock}
(Qff) n l u s 和 (QYf)a 也形成傅立叶变换对。

总结 表 1.1 总结了特征函数的几个有用属性及其与 PDF 的关系。

基于傅里叶变换方程。 (1.3) 和 (1.4) 验证

\begin{sourceblock}
\includegraphics[page=737,pagebox=mediabox,trim={64.39bp 412.95bp 64.39bp 1483.81bp},clip,width=0.94\textwidth,max height=0.38\textheight]{Pope.pdf}
\end{sourceblock}
使用方程式。 (1.7) 和 (1.12) 确定前六个矩
标准化正态分布。

让UI、U2、...，UN 是独立的，同分布的，标准的
化随机变量。表明

\begin{center}\small 表 I.1。特征函数和概率密度函数之间的关系（右侧的量是左侧量的傅里叶变换）\end{center}
\clearpage
\begin{sourceblock}
\includegraphics[page=738,pagebox=mediabox,trim={64.39bp 762.90bp 64.39bp 306.96bp},clip,width=\textwidth,max height=0.72\textheight]{Pope.pdf}
\end{sourceblock}
是标准化随机变量。表明特征函数
XN,YN(\textasciitilde{}i)s 的重刑，与 U S 的重刑有关）

\begin{sourceblock}
\includegraphics[page=738,pagebox=mediabox,trim={64.39bp 600.92bp 64.39bp 1339.83bp},clip,width=0.94\textwidth,max height=0.38\textheight]{Pope.pdf}
\end{sourceblock}
因此

\begin{sourceblock}
\includegraphics[page=738,pagebox=mediabox,trim={64.39bp 399.95bp 64.39bp 1482.81bp},clip,width=0.94\textwidth,max height=0.38\textheight]{Pope.pdf}
\end{sourceblock}
从而得到结果

\begin{sourceblock}
\includegraphics[page=738,pagebox=mediabox,trim={64.39bp 255.97bp 64.39bp 1669.79bp},clip,width=0.94\textwidth,max height=0.38\textheight]{Pope.pdf}
\end{sourceblock}
也就是说，XN 趋于标准化正态随机变量。 （这是
一个简单的版本

中心极限定理。）

分别为 U 的 PDF 和细粒度特征函数。显示
U Q 的特征函数为

\begin{sourceblock}
\includegraphics[page=739,pagebox=mediabox,trim={64.39bp 1625.80bp 64.39bp 261.97bp},clip,width=0.94\textwidth,max height=0.38\textheight]{Pope.pdf}
\end{sourceblock}
相应地，U Q 的 PDF 为
联合随机变量设 U \{ U , U 2 7。。 .U D) 是一组 D 随机变量。相应的特征函数定义为

\begin{sourceblock}
\includegraphics[page=739,pagebox=mediabox,trim={64.39bp 1147.86bp 64.39bp 794.90bp},clip,width=0.94\textwidth,max height=0.38\textheight]{Pope.pdf}
\end{sourceblock}
其中 s \{s1,s2。 , +, s D) 是自变量。

设 A 为常数 D 向量，B 为常数 D x D 矩阵，用于定义线性变换

\begin{sourceblock}
\includegraphics[page=739,pagebox=mediabox,trim={64.39bp 923.88bp 64.39bp 1023.87bp},clip,width=0.94\textwidth,max height=0.38\textheight]{Pope.pdf}
\end{sourceblock}
则U的特征函数Y(s)为

\begin{sourceblock}
\includegraphics[page=739,pagebox=mediabox,trim={64.39bp 780.90bp 64.39bp 1148.86bp},clip,width=0.94\textwidth,max height=0.38\textheight]{Pope.pdf}
\end{sourceblock}
通过设置 sk 0, k f j 可以从 Y(s) 中获得单个随机变量 U 的特征函数 Y ( j ) ( \textasciitilde{}o)f （从式（1.25）可以看出）。

如果 D 个随机变量相互独立，则

\begin{sourceblock}
\includegraphics[page=739,pagebox=mediabox,trim={64.39bp 493.94bp 64.39bp 1382.83bp},clip,width=0.94\textwidth,max height=0.38\textheight]{Pope.pdf}
\end{sourceblock}
其中 \#brackets 中的后缀不包括在求和约定中。它由等式得出。由(1.12)和(1.28)可知D个联合正态标准化随机变量U的特征函数为

\begin{sourceblock}
\includegraphics[page=739,pagebox=mediabox,trim={64.39bp 277.96bp 64.39bp 1668.79bp},clip,width=0.94\textwidth,max height=0.38\textheight]{Pope.pdf}
\end{sourceblock}
如果这些随机变量经过线性变换方程： (I.26)7 那么所得的随机 v e c t \textasciitilde{} rU 具有平均值

\begin{sourceblock}
\includegraphics[page=739,pagebox=mediabox,trim={64.39bp 119.39bp 64.39bp 1849.77bp},clip,width=0.94\textwidth,max height=0.38\textheight]{Pope.pdf}
\end{sourceblock}
\begin{sourceblock}
\includegraphics[page=740,pagebox=mediabox,trim={64.39bp 1728.78bp 64.39bp 162.98bp},clip,width=0.94\textwidth,max height=0.38\textheight]{Pope.pdf}
\end{sourceblock}
和特征函数（来自方程（1.27）和（1.29））

\begin{sourceblock}
\includegraphics[page=740,pagebox=mediabox,trim={64.39bp 1545.81bp 64.39bp 334.96bp},clip,width=0.94\textwidth,max height=0.38\textheight]{Pope.pdf}
\end{sourceblock}
实际上，该方程通常被视为联合正态随机变量的定义。

\chapter{扩散过程}
本附录简要介绍了扩散过程以及在其描述和分析中使用的一些数学技术。加德纳（Gardiner，1985）、吉莱斯皮（Gillespie，1992）、阿诺德（Arnold，1974）以及卡林和泰勒（Karlin and Taylor，1981）提供了更全面和严格的描述。

扩散过程是一种特殊的随机过程。它是一个具有连续样本路径（以及下面描述的其他属性）的连续时间马尔可夫过程。

马尔可夫过程

令 t 2 的 U(t) 为具有一次性 PDF f ( V t) 的随机过程。我们引入 N 次 t1 < t2。。。 < tN, （其中 tl > to），并考虑早期时间以 U(t) 为条件的 U(tN) 的 PDF \{U(tN-1),U (tN-2),- u(ti)\}？表示为
以单个过去时间为条件的 U(t) 的 PDF 表示为，例如，
根据定义，如果 U(t) 是马尔可夫过程，那么这些条件 PDF 是相等的：

这意味着，给定 U(tN-l) 先前值的知识
, U ( t l ) 没有提供有关未来 U(\textasciitilde{}N-2)U, (tN-3)7 值 U(tN) 的进一步信息。

Chapman-Kolmogorov 方程 对于任何过程，根据条件 PDF 的定义，我们有

\begin{sourceblock}
\includegraphics[page=742,pagebox=mediabox,trim={64.39bp 1633.79bp 64.39bp 270.97bp},clip,width=0.94\textwidth,max height=0.38\textheight]{Pope.pdf}
\end{sourceblock}
（参见练习 J.1）。对于马尔可夫过程，方程。 (J.1) 可以用 f1(V3;5 t) 代替，得到 Chapman-Kolmogorov 方程

\begin{sourceblock}
\includegraphics[page=742,pagebox=mediabox,trim={64.39bp 1421.82bp 64.39bp 481.94bp},clip,width=0.94\textwidth,max height=0.38\textheight]{Pope.pdf}
\end{sourceblock}
一个有用的概念是过程中的增量：正时间间隔 h 中的增量定义为

\begin{sourceblock}
\includegraphics[page=742,pagebox=mediabox,trim={64.39bp 1178.85bp 64.39bp 767.90bp},clip,width=0.94\textwidth,max height=0.38\textheight]{Pope.pdf}
\end{sourceblock}
值得注意的是，h 是正数，并且增量是在时间上向前定义的。一个过程可以被视为其增量的总和，例如，
U(ty) = U(to\} + Arp U (to) + Ary 1, U(t1) + 任意 ty, U (tw)。

增量 A,U(t) 的 PDF，以 U(r) = V 为条件，用 g(V;h, V,t) 表示。若将h取为t； - tf，则 U(t,) 可以重新表示为

\begin{sourceblock}
\includegraphics[page=742,pagebox=mediabox,trim={64.39bp 810.90bp 64.39bp 1136.86bp},clip,width=0.94\textwidth,max height=0.38\textheight]{Pope.pdf}
\end{sourceblock}
以及等式右侧的第一个条件 PDF。 (J.3) 是

\begin{sourceblock}
\includegraphics[page=742,pagebox=mediabox,trim={64.39bp 669.92bp 64.39bp 1268.84bp},clip,width=0.94\textwidth,max height=0.38\textheight]{Pope.pdf}
\end{sourceblock}
因此查普曼-柯尔莫哥洛夫方程可以重写为

\begin{sourceblock}
\includegraphics[page=742,pagebox=mediabox,trim={64.39bp 492.94bp 64.39bp 1408.82bp},clip,width=0.94\textwidth,max height=0.38\textheight]{Pope.pdf}
\end{sourceblock}
扩散过程 连续时间马尔可夫过程有多种性质不同的类型，它们之间的区别在于它们在 h 趋于零时极限内的增量 A,U(t) 的行为。扩散过程的一个定义属性是其样本路径是连续的。更准确地说，对于每一个
如果存在，则进程的 i\textasciitilde{}finitcsimalp 参数定义为
( V , t) lim -([AI,U(t)lHI U(t) = V),

\begin{sourceblock}
\includegraphics[page=743,pagebox=mediabox,trim={64.39bp 1613.80bp 64.39bp 316.96bp},clip,width=0.94\textwidth,max height=0.38\textheight]{Pope.pdf}
\end{sourceblock}
对于 n l , 2 , .... 除了方程。 （5.9），扩散过程的定义属性是漂移系数，

\begin{sourceblock}
\includegraphics[page=743,pagebox=mediabox,trim={64.39bp 1420.82bp 64.39bp 524.93bp},clip,width=0.94\textwidth,max height=0.38\textheight]{Pope.pdf}
\end{sourceblock}
和 d\$Tusion 系数，

\begin{sourceblock}
\includegraphics[page=743,pagebox=mediabox,trim={64.39bp 1270.84bp 64.39bp 658.92bp},clip,width=0.94\textwidth,max height=0.38\textheight]{Pope.pdf}
\end{sourceblock}
存在，并且剩余的无穷小参数为零：

\begin{sourceblock}
\includegraphics[page=743,pagebox=mediabox,trim={64.39bp 1137.86bp 64.39bp 809.90bp},clip,width=0.94\textwidth,max height=0.38\textheight]{Pope.pdf}
\end{sourceblock}
由常微分控制的可微确定性过程

\begin{sourceblock}
\includegraphics[page=743,pagebox=mediabox,trim={64.39bp 932.88bp 64.39bp 924.88bp},clip,width=0.94\textwidth,max height=0.38\textheight]{Pope.pdf}
\end{sourceblock}
是简并扩散过程，漂移 a(V, t) 和扩散系数 b(V, t)2 0。非简并扩散过程 ( . e b (V, t) > 0) 显然无处可微，因为 ([AJ.J(t)12 /h) 趋于正极限这一事实意味着 ( l A h U ( t ) / h l L ) 趋于无穷大。

Kramers-Moyal 方程 在 Chapman-Kolmogorov 方程（方程（J.8））中，右侧的 g 和 f 都涉及参数 V 将这些量展开为关于 V 的泰勒级数，得出

\begin{sourceblock}
\includegraphics[page=743,pagebox=mediabox,trim={64.39bp 325.96bp 64.39bp 1511.81bp},clip,width=0.94\textwidth,max height=0.38\textheight]{Pope.pdf}
\end{sourceblock}
Bv 除以 h，取极限 h 0，并使用等式： (J.10)，我们得到Kramers-Moyul方程

\begin{sourceblock}
\includegraphics[page=743,pagebox=mediabox,trim={64.39bp 119.39bp 64.39bp 1794.77bp},clip,width=0.94\textwidth,max height=0.38\textheight]{Pope.pdf}
\end{sourceblock}
该方程适用于参数 (V t) 存在的过程，并且对于 2 适当的初始条件为

\begin{sourceblock}
\includegraphics[page=744,pagebox=mediabox,trim={64.39bp 1684.79bp 64.39bp 260.97bp},clip,width=0.94\textwidth,max height=0.38\textheight]{Pope.pdf}
\end{sourceblock}
Fokker-Planck 方程 对于扩散过程，除了漂移 B1 a 和扩散 B2 b2 之外，所有参数 B 均为零。在这种情况下，等式。 (5.16) 简化为 Fokker-Planck 或前向 Kolmogorov 方程
-fl(v;tlvl?tl) = a v ti tl)l
2 [ b ( V , t ) 2 fl ( \textasciitilde{} ; t l v l , t l ) l \textasciitilde{} 该方程决定了条件 PDF 的演化。

边际 PDF f ( V ; t) 的相应方程是通过乘以 f (V1; t i ) 并在 VimS ince 上积分获得的，在方程 1 中。 (5. 18)，只有 f 与 V 相关，结果很简单

\begin{sourceblock}
\includegraphics[page=744,pagebox=mediabox,trim={64.39bp 995.87bp 64.39bp 906.89bp},clip,width=0.94\textwidth,max height=0.38\textheight]{Pope.pdf}
\end{sourceblock}
对于由常微分方程 Eq 控制的确定性过程。 （5.14），扩散系数为零，因此方程中的最后一项。 （5.18）和（5.19）消失。由此产生的方程称为刘维尔方程。

平稳分布 如果系数 a 和 b 与时间无关，则扩散过程可能在统计上是平稳的。在这种情况下，福克-普朗克方程（方程（5.19））简化为

\begin{sourceblock}
\includegraphics[page=744,pagebox=mediabox,trim={64.39bp 447.94bp 64.39bp 1452.82bp},clip,width=0.94\textwidth,max height=0.38\textheight]{Pope.pdf}
\end{sourceblock}
哪个有解决方案

其中，为方便起见，可以选择下限Vo，常数C由归一化条件确定。 （如果方程（5.21）右侧对所有 V 的积分不收敛，则 t T ( t ) 不具有

\begin{sourceblock}
\includegraphics[page=745,pagebox=mediabox,trim={64.39bp 1157.85bp 64.39bp 109.44bp},clip,width=\textwidth,max height=0.72\textheight]{Pope.pdf}
\par\vspace{0.45\baselineskip}
\small 图 J.1。维纳过程的三个样本路径。
\end{sourceblock}
维纳过程 最基本的扩散过程是维纳过程，用 W(t) 表示，所有其他过程都可以从中导出。这是由初始条件 W(0) 0 以及漂移和扩散系数的规范定义的（对于 t 2 0），

\begin{sourceblock}
\includegraphics[page=745,pagebox=mediabox,trim={64.39bp 777.90bp 64.39bp 1169.85bp},clip,width=0.94\textwidth,max height=0.38\textheight]{Pope.pdf}
\end{sourceblock}
W(t) 的一些样本路径如图 5.1 所示。

可以很容易地验证，从初始条件方程（5.18）可以解出 a 0 和 b2 1 的 Fokker-Planck 方程（方程（5.18））。 (5.17) 是

\begin{sourceblock}
\includegraphics[page=745,pagebox=mediabox,trim={64.39bp 496.94bp 64.39bp 1387.83bp},clip,width=0.94\textwidth,max height=0.38\textheight]{Pope.pdf}
\end{sourceblock}
i-ema，均值 Vl 和方差 t - tl 的正态分布。 \textasciitilde{} h u s f,or all h > 0，增量 AhW (t) 服从均值为零的正态分布，

\begin{sourceblock}
\includegraphics[page=745,pagebox=mediabox,trim={64.39bp 298.96bp 64.39bp 1588.80bp},clip,width=0.94\textwidth,max height=0.38\textheight]{Pope.pdf}
\end{sourceblock}
（该符号读作“分布等于”，并且 a2）d 表示平均值为 p 且方差为 a2,E q 的正态分布。 （3.41）。）。

维纳过程增量的一些重要属性（并非全部独立）

(v) W(t2) 对于 t 5 t l，W ( t l ) 独立于 W (t)，
(vi) ( [ W ( \textasciitilde{} \textasciitilde{} ) - W ( \textasciitilde{} : ) I [ W ( - \textasciitilde{} WZ ) t ) ) = 0 - 非重叠增量
时间间隔是独立的，
(vii) ([W(t4) - W ( t \textasciitilde{} ) l [ W ( h - ) W(tl)]) = - 增量的协方差
ments 等于时间间隔重叠的持续时间，
zN jviii) ([W(t,) - \textasciitilde{} ( t \textasciitilde{} )= ]tN\textasciitilde{} - )to，并且
(ix) W(t) 是高斯过程：W(t1)、W(t2) 的联合 PDF。 ..，W (\textasciitilde{}N)

是关节正常值。

练习 5.2 中推导了维纳过程的其他几个有趣的性质。

随机微分方程 由于扩散过程不可微，因此无法应用微分计算的标准工具。合适的方法是伊藤微积分，而不是微积分；并且，扩散过程不是用常微分方程描述，而是用随机微分方程描述。

过程 U(t) 的无穷小增量定义为
其中 dt 是正无穷小时间间隔。特别是对于维纳过程，我们有

\begin{sourceblock}
\includegraphics[page=746,pagebox=mediabox,trim={64.39bp 603.92bp 64.39bp 1332.83bp},clip,width=0.94\textwidth,max height=0.38\textheight]{Pope.pdf}
\end{sourceblock}
现在考虑由初始条件 U(to) = Uo 和增量定义的过程 U(t)

\begin{sourceblock}
\includegraphics[page=746,pagebox=mediabox,trim={64.39bp 426.95bp 64.39bp 1518.81bp},clip,width=0.94\textwidth,max height=0.38\textheight]{Pope.pdf}
\end{sourceblock}
对于给定函数 a ( V, t) 和 b ( V t ) It t 很容易验证由该随机微分方程定义的过程 U(t) 是一个扩散过程；并且，如符号所示，漂移和扩散系数为 a ( V , t ) 和 b(V,
随机变量完全由其 PDF 表征；具有相同 PDF 的两个随机变量在统计上是相同的。同样，扩散
具有相同系数的扩散过程在统计上是相同的。因此随机微分方程 Eq。方程（5.27）提供了扩散过程的一般表达式。

随机微分方程(5.27) 表明扩散过程的无穷小增量是高斯分布的，i-em，

\begin{sourceblock}
\includegraphics[page=747,pagebox=mediabox,trim={64.39bp 1543.81bp 64.39bp 398.95bp},clip,width=0.94\textwidth,max height=0.38\textheight]{Pope.pdf}
\end{sourceblock}
这种高斯性不是扩散过程的定义属性，而是其定义的推论。

白噪声 在随机微分方程理论发展之前，扩散过程通常被表示为涉及白噪声的常微分方程。关于除法方程。 (5.27) 由 d t 我们得到

\begin{sourceblock}
\includegraphics[page=747,pagebox=mediabox,trim={64.39bp 1098.86bp 64.39bp 809.90bp},clip,width=0.94\textwidth,max height=0.38\textheight]{Pope.pdf}
\end{sourceblock}
其中白噪声 \textasciitilde{} ( t 是) d W(t)/dt。由于 d U / d t 和 d W / d t 都不存在，因此不能以通常的方式解释该方程。因此，最好不要使用白噪声的概念，而是将扩散过程表示为随机微分方程。

矩的演化 无条件矩 (U(t)") 的演化方程可以从 Fokker-Planck 方程 (Eq. (5 . 9) ) 导出，或者从随机微分方程 (Eq. (5.27)) 导出 - 后一种方法具有指导意义。

取方程的平均值。 （5.25）并代入等式。 （5.27），我们得到

\begin{sourceblock}
\includegraphics[page=747,pagebox=mediabox,trim={64.39bp 494.94bp 64.39bp 1393.82bp},clip,width=0.94\textwidth,max height=0.38\textheight]{Pope.pdf}
\end{sourceblock}
现在 d W (t) 的均值为零，并且对于 t' 5 t，它独立于 U(t')。因此最后一项消失，导致

\begin{sourceblock}
\includegraphics[page=747,pagebox=mediabox,trim={64.39bp 279.96bp 64.39bp 1627.79bp},clip,width=0.94\textwidth,max height=0.38\textheight]{Pope.pdf}
\end{sourceblock}
类似地，U(t dt) 的平方平均值为

\begin{sourceblock}
\includegraphics[page=747,pagebox=mediabox,trim={64.39bp 119.39bp 64.39bp 1798.77bp},clip,width=0.94\textwidth,max height=0.38\textheight]{Pope.pdf}
\end{sourceblock}
交叉项为 2(Ua) dt，最终项为

\begin{sourceblock}
\includegraphics[page=748,pagebox=mediabox,trim={64.39bp 1632.79bp 64.39bp 256.97bp},clip,width=0.94\textwidth,max height=0.38\textheight]{Pope.pdf}
\end{sourceblock}
其中 o(h) 表示一个量，使得
（例如，h o(h)，对于所有 0）。因此 U(t) 的均方演化为

\begin{sourceblock}
\includegraphics[page=748,pagebox=mediabox,trim={64.39bp 1315.83bp 64.39bp 591.93bp},clip,width=0.94\textwidth,max height=0.38\textheight]{Pope.pdf}
\end{sourceblock}
请注意，对于可微过程 ( d \textasciitilde{} ( t ) \textasciitilde{} ) /isd zte ro，但是对于扩散过程，它是 (b2)，这导致方程 1 中的最后一项。 （J.35）。

Ornstein-Uhlenbeck (OU) 过程 OU 过程是最简单的统计平稳扩散过程。它由线性漂移系数定义

\begin{sourceblock}
\includegraphics[page=748,pagebox=mediabox,trim={64.39bp 966.88bp 64.39bp 979.88bp},clip,width=0.94\textwidth,max height=0.38\textheight]{Pope.pdf}
\end{sourceblock}
恒定扩散系数

\begin{sourceblock}
\includegraphics[page=748,pagebox=mediabox,trim={64.39bp 804.90bp 64.39bp 1099.86bp},clip,width=0.94\textwidth,max height=0.38\textheight]{Pope.pdf}
\end{sourceblock}
和初始条件

\begin{sourceblock}
\includegraphics[page=748,pagebox=mediabox,trim={64.39bp 693.91bp 64.39bp 1242.84bp},clip,width=0.94\textwidth,max height=0.38\textheight]{Pope.pdf}
\end{sourceblock}
其中 T 是正时间尺度，a 是常数。对应的随机微分方程为朗之万方程
对于 OU 过程，U(t) 的 PDF 的 Fokker-Planck 方程（方程 (5.18)）以 U(tl) = V1（对于 t > tl）为条件，( V ; t V t ) 为

\begin{sourceblock}
\includegraphics[page=748,pagebox=mediabox,trim={64.39bp 271.97bp 64.39bp 1632.79bp},clip,width=0.94\textwidth,max height=0.38\textheight]{Pope.pdf}
\end{sourceblock}
具有确定性初始条件方程。 (5.17)，该方程的解是正态分布

\begin{sourceblock}
\includegraphics[page=748,pagebox=mediabox,trim={64.39bp 119.39bp 64.39bp 1832.77bp},clip,width=0.94\textwidth,max height=0.38\textheight]{Pope.pdf}
\end{sourceblock}
（参见练习 5.4）。该解充分表征了该过程，表明条件均值 (U(t)l V1) 在时间尺度 T 上从 V1 衰减到 0；而条件方差在时间尺度 L T 上从 0 增加到 o2。在很大的时候，条件 PDF 趋于平稳分布
等式的一个重要推论。 （5.41）式表明OU过程是高斯过程。马尔可夫性质的结果是 U(t) 在 N 1 次 \{to = 0 , t l , t 2, ..., t \} 处的联合 PDF 可以写为 to 处的边际 PDF 与 N 个条件 PDF fl(V'; t.

t ) for n = 1,2,..., N。由于这些 PDF 中的每一个都是正态的（参见 - 1 .-I 方程 (5.38) 和 ( 5 . 4 ) ），因此 (N 1) 时间联合 PDF 是联合正态的，满足高斯过程的定义。

由于均值 ( \textasciitilde{} ( t ) i ) s 为零，因此自协方差（对于 s 2 0）为

\begin{sourceblock}
\includegraphics[page=749,pagebox=mediabox,trim={64.39bp 1202.85bp 64.39bp 741.91bp},clip,width=0.94\textwidth,max height=0.38\textheight]{Pope.pdf}
\end{sourceblock}
这可以通过条件平均值 (U ( t l s)/V 1) V1 e 轻松评估：

\begin{sourceblock}
\includegraphics[page=749,pagebox=mediabox,trim={64.39bp 965.88bp 64.39bp 923.88bp},clip,width=0.94\textwidth,max height=0.38\textheight]{Pope.pdf}
\end{sourceblock}
对于任何统计平稳过程，自协方差 R(s) 和自相关函数 p(s) 都是偶函数。因此，对于 OU 流程，

\begin{sourceblock}
\includegraphics[page=749,pagebox=mediabox,trim={64.39bp 756.90bp 64.39bp 1134.86bp},clip,width=0.94\textwidth,max height=0.38\textheight]{Pope.pdf}
\end{sourceblock}
\begin{sourceblock}
\includegraphics[page=749,pagebox=mediabox,trim={64.39bp 672.91bp 64.39bp 1269.84bp},clip,width=0.94\textwidth,max height=0.38\textheight]{Pope.pdf}
\end{sourceblock}
请注意，p(s) ds 定义的积分时间尺度为 T。

二阶结构函数D2(s)和频谱E(o)包含与自协方差R(s)相同的信息。对于OU进程，结构函数为

\begin{sourceblock}
\includegraphics[page=749,pagebox=mediabox,trim={64.39bp 315.96bp 64.39bp 1531.81bp},clip,width=0.94\textwidth,max height=0.38\textheight]{Pope.pdf}
\end{sourceblock}
而频谱（R(s) 的傅里叶变换的两倍）是

\begin{sourceblock}
\includegraphics[page=749,pagebox=mediabox,trim={64.39bp 156.98bp 64.39bp 1753.78bp},clip,width=0.94\textwidth,max height=0.38\textheight]{Pope.pdf}
\end{sourceblock}
p(s) 在原点处以及变体 D2(s) s（对于小 s）和 E(o) wd2（对于大 o）的斜率。

总而言之； Ornstein-Uhlenbeck 过程（由 Eangevin 方程，Eq. (5.39) 生成）是统计平稳高斯过程。因此，它完全由均值（为零）、方差 a2 和自相关函数 p(s) = e 来表征。

伊藤变换 考虑由下式定义的过程 q(t)

\begin{sourceblock}
\includegraphics[page=750,pagebox=mediabox,trim={64.39bp 1392.82bp 64.39bp 554.93bp},clip,width=0.94\textwidth,max height=0.38\textheight]{Pope.pdf}
\end{sourceblock}
其中 U(t) 是扩散过程（具有漂移 a ( V , t) 和扩散 b(V, Q2)，Q ( V ) 是可微函数，具有导数 Qf(V)、Q U ( V ,) 等。 q 的无限小增量为

\begin{sourceblock}
\includegraphics[page=750,pagebox=mediabox,trim={64.39bp 1163.85bp 64.39bp 767.90bp},clip,width=0.94\textwidth,max height=0.38\textheight]{Pope.pdf}
\end{sourceblock}
通过在关于 U(t) 的泰勒级数中展开 Q(U d U )，并代入方程： (5.27) 对于 dU，我们得到

\begin{sourceblock}
\includegraphics[page=750,pagebox=mediabox,trim={64.39bp 939.88bp 64.39bp 989.88bp},clip,width=0.94\textwidth,max height=0.38\textheight]{Pope.pdf}
\end{sourceblock}
因此，q(t) 本身就是一个扩散过程，具有随机微分方程

\begin{sourceblock}
\includegraphics[page=750,pagebox=mediabox,trim={64.39bp 813.90bp 64.39bp 1129.86bp},clip,width=0.94\textwidth,max height=0.38\textheight]{Pope.pdf}
\end{sourceblock}
其中系数是

\begin{sourceblock}
\includegraphics[page=750,pagebox=mediabox,trim={64.39bp 682.91bp 64.39bp 1260.84bp},clip,width=0.94\textwidth,max height=0.38\textheight]{Pope.pdf}
\end{sourceblock}
\begin{sourceblock}
\includegraphics[page=750,pagebox=mediabox,trim={64.39bp 605.92bp 64.39bp 1338.83bp},clip,width=0.94\textwidth,max height=0.38\textheight]{Pope.pdf}
\end{sourceblock}
这些关系形成了伊藤变换。与常微分方程的变换规则相比，本质的区别在于附加的漂移-1Q n b 2。

矢量值扩散过程 上面提出的标量值扩散过程的发展直接延伸到矢量值过程 U(t) \{U,(t), u2(t), - UD(\textasciitilde{})\}O。这里仅给出主要结果。

漂移系数是向量

\begin{sourceblock}
\includegraphics[page=750,pagebox=mediabox,trim={64.39bp 119.39bp 64.39bp 1805.77bp},clip,width=0.94\textwidth,max height=0.38\textheight]{Pope.pdf}
\end{sourceblock}
而扩散系数是 D x D 矩阵 B，其中元素

\begin{sourceblock}
\includegraphics[page=751,pagebox=mediabox,trim={64.39bp 1698.79bp 64.39bp 229.97bp},clip,width=0.94\textwidth,max height=0.38\textheight]{Pope.pdf}
\end{sourceblock}
从其定义可知，B 是对称正半定的。

U(t),f ( V ; t) 的 PDF 以及条件 PDF 的 Fokker-Planck 方程为

\begin{sourceblock}
\includegraphics[page=751,pagebox=mediabox,trim={64.39bp 1455.82bp 64.39bp 444.94bp},clip,width=0.94\textwidth,max height=0.38\textheight]{Pope.pdf}
\end{sourceblock}
参见等式。 （J.19）。

向量值维纳过程 W (t) = \{ W1 (t), W2(t), ...。 , WD(t)\}i 只是由独立的标量过程 Wi(t) 组成。增量 d W(t) 是联合正态，均值为零，协方差

\begin{sourceblock}
\includegraphics[page=751,pagebox=mediabox,trim={64.39bp 1199.85bp 64.39bp 747.91bp},clip,width=0.94\textwidth,max height=0.38\textheight]{Pope.pdf}
\end{sourceblock}
该过程是统计各向同性的：如果 A 是酉矩阵，则

\textasciitilde{} ( t ) A W(t) 也是一个向量值维纳过程。

U(t) 的随机微分方程可写为

\begin{sourceblock}
\includegraphics[page=751,pagebox=mediabox,trim={64.39bp 984.88bp 64.39bp 960.88bp},clip,width=0.94\textwidth,max height=0.38\textheight]{Pope.pdf}
\end{sourceblock}
其中（为了与方程（J.55）一致）系数 b n -sa tsfy

\begin{sourceblock}
\includegraphics[page=751,pagebox=mediabox,trim={64.39bp 855.89bp 64.39bp 1088.86bp},clip,width=0.94\textwidth,max height=0.38\textheight]{Pope.pdf}
\end{sourceblock}
请注意，非对称矩阵 b 并非由对称矩阵 B 唯一确定。b 的两种可能选择是 B 的对称平方根和 B 的 Cholesky 分解给出的下三角矩阵。b 的所有选择（与方程 (J.59) 一致）都会导致统计上相同的幻觉过程。

考虑三个随机变量 U1、U2 和 Ui。有条件的 PDF
由联合 PDF 定义，例如，

\begin{sourceblock}
\includegraphics[page=751,pagebox=mediabox,trim={64.39bp 381.95bp 64.39bp 1563.80bp},clip,width=0.94\textwidth,max height=0.38\textheight]{Pope.pdf}
\end{sourceblock}
\begin{sourceblock}
\includegraphics[page=751,pagebox=mediabox,trim={64.39bp 306.96bp 64.39bp 1638.79bp},clip,width=0.94\textwidth,max height=0.38\textheight]{Pope.pdf}
\end{sourceblock}
获取结果

\begin{sourceblock}
\includegraphics[page=751,pagebox=mediabox,trim={64.39bp 149.98bp 64.39bp 1752.78bp},clip,width=0.94\textwidth,max height=0.38\textheight]{Pope.pdf}
\end{sourceblock}
将时间间隔 (0, T) 分为 M 个相等的子间隔
持续时间 h TIM。对于离散时间 nh (n 0,1,. .., M)，
进程由 0 和定义

\begin{sourceblock}
\includegraphics[page=752,pagebox=mediabox,trim={64.39bp 1578.80bp 64.39bp 328.96bp},clip,width=0.94\textwidth,max height=0.38\textheight]{Pope.pdf}
\end{sourceblock}
其中 <1,\{2,-.., 是独立的标准化正态随机
(a) 显示在统计上与维纳过程相同
同时采样。

(b) 令 SM 表示增量的平方和

\begin{sourceblock}
\includegraphics[page=752,pagebox=mediabox,trim={64.39bp 1136.86bp 64.39bp 687.91bp},clip,width=0.94\textwidth,max height=0.38\textheight]{Pope.pdf}
\end{sourceblock}
获得结果 (SM) T 和 var(SM) 2 h \textasciitilde{} 2 \textasciitilde{}\textasciitilde{} 1 \textasciitilde{}。

因此认为随机变量 SMh 作为非随机变量
(c) 考虑 t 的图，其中连续值是
由直线段连接。表明预期的
每条线段的长度超过 ,/2h/7r，并且
长度总和的期望超过 ～ M T / z。

因此认为，在每个正区间中，样本路径
维纳过程的弧长无限大。

(d) W(nh) 和 W(nh) 的联合 PDF f,(V, ; h) 是多少？对于
n 2 1，考虑事件 C，由 W(nh) 和 VV [(n l)h] 定义
具有相反的符号。 V样本空间的哪个区域
对应事件C，？证明 C 的概率，

\begin{sourceblock}
\includegraphics[page=752,pagebox=mediabox,trim={64.39bp 377.95bp 64.39bp 1514.81bp},clip,width=0.94\textwidth,max height=0.38\textheight]{Pope.pdf}
\end{sourceblock}
（提示：考虑变换 f 的简单变换，
到标准化关节法线。）的预期数量
改变符号的次数为 NM P(Cn)。怎么样
当M增加时NM表现如何？因此认为，在任何积极的
间隔 (tl,t 2)，维纳过程取之间的每个值
(e) 对于任意正常数 c，证明在变换后的
时间t=ct，过程

\begin{sourceblock}
\includegraphics[page=753,pagebox=mediabox,trim={64.39bp 1662.79bp 64.39bp 276.96bp},clip,width=0.94\textwidth,max height=0.38\textheight]{Pope.pdf}
\end{sourceblock}
是一个维纳过程。

对于方程定义的扩散过程。 (J.27)，以下哪一项
数量是相关的：W(t)、dW(t)、U(t)、dU(t) 和 U(t dt)？

令 Y(s, t)（对于 t 2 t l ）为 OU 的特征函数
从确定性初始条件 U(tl) V1 开始处理；这样
其傅里叶变换是条件 PDF f ( V ; t V , t ) 从
f 的 Fokker-Planck 方程（方程 (5.40)）表明 Y(s, t) 演化

\begin{sourceblock}
\includegraphics[page=753,pagebox=mediabox,trim={64.39bp 1153.85bp 64.39bp 690.91bp},clip,width=0.94\textwidth,max height=0.38\textheight]{Pope.pdf}
\end{sourceblock}
利用特征法求得解

\begin{sourceblock}
\includegraphics[page=753,pagebox=mediabox,trim={64.39bp 1020.87bp 64.39bp 926.88bp},clip,width=0.94\textwidth,max height=0.38\textheight]{Pope.pdf}
\end{sourceblock}
\begin{sourceblock}
\includegraphics[page=753,pagebox=mediabox,trim={64.39bp 888.89bp 64.39bp 1053.87bp},clip,width=0.94\textwidth,max height=0.38\textheight]{Pope.pdf}
\end{sourceblock}
证明，在初始条件 ( s t ) = v1 的情况下，该解
对应于平均值的正态分布

\begin{sourceblock}
\includegraphics[page=753,pagebox=mediabox,trim={64.39bp 700.91bp 64.39bp 1244.84bp},clip,width=0.94\textwidth,max height=0.38\textheight]{Pope.pdf}
\end{sourceblock}
和方差

让速度 U(t) \{Ul(t),U 2(t),U 3(t)\} 的每个分量演化
通过独立的朗之万方程，即

\begin{sourceblock}
\includegraphics[page=753,pagebox=mediabox,trim={64.39bp 416.95bp 64.39bp 1472.81bp},clip,width=0.94\textwidth,max height=0.38\textheight]{Pope.pdf}
\end{sourceblock}
U 的联合 PDF 是什么？速度 q(t) 定义为

\begin{sourceblock}
\includegraphics[page=753,pagebox=mediabox,trim={64.39bp 282.96bp 64.39bp 1662.79bp},clip,width=0.94\textwidth,max height=0.38\textheight]{Pope.pdf}
\end{sourceblock}
证明 q(t) 由随机微分方程演化

\begin{sourceblock}
\includegraphics[page=753,pagebox=mediabox,trim={64.39bp 119.39bp 64.39bp 1797.77bp},clip,width=0.94\textwidth,max height=0.38\textheight]{Pope.pdf}
\end{sourceblock}
证明 q 的平稳分布为

\begin{sourceblock}
\includegraphics[page=754,pagebox=mediabox,trim={64.39bp 1698.79bp 64.39bp 207.97bp},clip,width=0.94\textwidth,max height=0.38\textheight]{Pope.pdf}
\end{sourceblock}
令 Y(s, t ) 为一般向量值的特征函数
扩散过程由方程给出。 （5.58）。通过展开 (exp(is.[ U(t)

dU(t)])) (见式(3)) 得到结果
Y(s, t dt) = Y(s, t) 是,(a,e d - IS (bjCbkee dt \textasciitilde{} ( d t ) ,

\begin{sourceblock}
\includegraphics[page=754,pagebox=mediabox,trim={64.39bp 1445.82bp 64.39bp 502.94bp},clip,width=0.94\textwidth,max height=0.38\textheight]{Pope.pdf}
\end{sourceblock}
其中系数例如为 a, = n,[U(t), t]。表明
该方程的傅里叶变换（除以 dt 时）是
福克-普朗克方程，Eq. （5.56）。

\clearpage\chapter*{参考文献}\addcontentsline{toc}{chapter}{参考文献}
\markboth{参考文献}{参考文献}
\clearpage
\begin{sourceblock}
\includegraphics[page=755,pagebox=mediabox,trim={70.83bp 278.57bp 70.83bp 198.98bp},clip,width=\textwidth,max height=0.82\textheight]{Pope.pdf}
\end{sourceblock}
\clearpage
\begin{sourceblock}
\includegraphics[page=756,pagebox=mediabox,trim={70.83bp 278.57bp 70.83bp 198.98bp},clip,width=\textwidth,max height=0.82\textheight]{Pope.pdf}
\end{sourceblock}
\clearpage
\begin{sourceblock}
\includegraphics[page=757,pagebox=mediabox,trim={70.83bp 278.57bp 70.83bp 198.98bp},clip,width=\textwidth,max height=0.82\textheight]{Pope.pdf}
\end{sourceblock}
\clearpage
\begin{sourceblock}
\includegraphics[page=758,pagebox=mediabox,trim={70.83bp 278.57bp 70.83bp 198.98bp},clip,width=\textwidth,max height=0.82\textheight]{Pope.pdf}
\end{sourceblock}
\clearpage
\begin{sourceblock}
\includegraphics[page=759,pagebox=mediabox,trim={70.83bp 278.57bp 70.83bp 198.98bp},clip,width=\textwidth,max height=0.82\textheight]{Pope.pdf}
\end{sourceblock}
\clearpage
\begin{sourceblock}
\includegraphics[page=760,pagebox=mediabox,trim={70.83bp 278.57bp 70.83bp 198.98bp},clip,width=\textwidth,max height=0.82\textheight]{Pope.pdf}
\end{sourceblock}
\clearpage
\begin{sourceblock}
\includegraphics[page=761,pagebox=mediabox,trim={70.83bp 278.57bp 70.83bp 198.98bp},clip,width=\textwidth,max height=0.82\textheight]{Pope.pdf}
\end{sourceblock}
\clearpage
\begin{sourceblock}
\includegraphics[page=762,pagebox=mediabox,trim={70.83bp 278.57bp 70.83bp 198.98bp},clip,width=\textwidth,max height=0.82\textheight]{Pope.pdf}
\end{sourceblock}
\clearpage
\begin{sourceblock}
\includegraphics[page=763,pagebox=mediabox,trim={70.83bp 278.57bp 70.83bp 198.98bp},clip,width=\textwidth,max height=0.82\textheight]{Pope.pdf}
\end{sourceblock}
\clearpage
\begin{sourceblock}
\includegraphics[page=764,pagebox=mediabox,trim={70.83bp 278.57bp 70.83bp 198.98bp},clip,width=\textwidth,max height=0.82\textheight]{Pope.pdf}
\end{sourceblock}
\clearpage
\begin{sourceblock}
\includegraphics[page=765,pagebox=mediabox,trim={70.83bp 278.57bp 70.83bp 198.98bp},clip,width=\textwidth,max height=0.82\textheight]{Pope.pdf}
\end{sourceblock}
\clearpage
\begin{sourceblock}
\includegraphics[page=766,pagebox=mediabox,trim={70.83bp 278.57bp 70.83bp 198.98bp},clip,width=\textwidth,max height=0.82\textheight]{Pope.pdf}
\end{sourceblock}
\clearpage
\begin{sourceblock}
\includegraphics[page=767,pagebox=mediabox,trim={70.83bp 278.57bp 70.83bp 198.98bp},clip,width=\textwidth,max height=0.82\textheight]{Pope.pdf}
\end{sourceblock}
\clearpage
\begin{sourceblock}
\includegraphics[page=768,pagebox=mediabox,trim={70.83bp 278.57bp 70.83bp 198.98bp},clip,width=\textwidth,max height=0.82\textheight]{Pope.pdf}
\end{sourceblock}
\clearpage
\begin{sourceblock}
\includegraphics[page=769,pagebox=mediabox,trim={70.83bp 278.57bp 70.83bp 198.98bp},clip,width=\textwidth,max height=0.82\textheight]{Pope.pdf}
\end{sourceblock}
\clearpage
\begin{sourceblock}
\includegraphics[page=770,pagebox=mediabox,trim={70.83bp 278.57bp 70.83bp 198.98bp},clip,width=\textwidth,max height=0.82\textheight]{Pope.pdf}
\end{sourceblock}
\clearpage
\begin{sourceblock}
\includegraphics[page=771,pagebox=mediabox,trim={70.83bp 278.57bp 70.83bp 198.98bp},clip,width=\textwidth,max height=0.82\textheight]{Pope.pdf}
\end{sourceblock}
\clearpage
\begin{sourceblock}
\includegraphics[page=772,pagebox=mediabox,trim={70.83bp 278.57bp 70.83bp 198.98bp},clip,width=\textwidth,max height=0.82\textheight]{Pope.pdf}
\end{sourceblock}
\clearpage
\begin{sourceblock}
\includegraphics[page=773,pagebox=mediabox,trim={70.83bp 278.57bp 70.83bp 198.98bp},clip,width=\textwidth,max height=0.82\textheight]{Pope.pdf}
\end{sourceblock}
\clearpage
\begin{sourceblock}
\includegraphics[page=774,pagebox=mediabox,trim={70.83bp 278.57bp 70.83bp 198.98bp},clip,width=\textwidth,max height=0.82\textheight]{Pope.pdf}
\end{sourceblock}
\clearpage
\begin{sourceblock}
\includegraphics[page=775,pagebox=mediabox,trim={70.83bp 278.57bp 70.83bp 198.98bp},clip,width=\textwidth,max height=0.82\textheight]{Pope.pdf}
\end{sourceblock}
\clearpage
\begin{sourceblock}
\includegraphics[page=776,pagebox=mediabox,trim={70.83bp 278.57bp 70.83bp 198.98bp},clip,width=\textwidth,max height=0.82\textheight]{Pope.pdf}
\end{sourceblock}
\clearpage\chapter*{作者索引}\addcontentsline{toc}{chapter}{作者索引}
\markboth{作者索引}{作者索引}
\clearpage
\begin{sourceblock}
\begin{tikzpicture}
\node[inner sep=0] (indeximage) {\begin{minipage}{\textwidth}\noindent\begin{minipage}{0.49\textwidth}\centering
\includegraphics[page=777,pagebox=mediabox,trim={70.83bp 198.97bp 656.75bp 198.98bp},clip,width=\linewidth,max height=0.82\textheight]{Pope.pdf}
\end{minipage}\hfill\begin{minipage}{0.49\textwidth}\centering
\includegraphics[page=777,pagebox=mediabox,trim={656.75bp 198.97bp 70.83bp 198.98bp},clip,width=\linewidth,max height=0.82\textheight]{Pope.pdf}
\end{minipage}
\end{minipage}};
\fill[white] ([xshift=-1cm,yshift=-0.05cm]indeximage.south) rectangle ([xshift=1cm,yshift=0.75cm]indeximage.south);
\end{tikzpicture}
\end{sourceblock}
\end{document}
